% =====================================================================
%  ANGLAIS — LE GRAND LIVRE D'AUTO-FORMATION
%  12 mois pour parler anglais comme vous parlez français
%  Méthode complète pour francophones : A1 -> B2 (avec passerelles C1)
%
%  FICHIER PRINCIPAL : compilez CE fichier (main.tex), deux fois.
%    - Overleaf : New Project -> Upload -> glisser tout le dossier,
%                 compiler main.tex (Menu -> Compiler : pdfLaTeX).
%    - En local : pdflatex main.tex && pdflatex main.tex
%                 (ou : latexmk -pdf main.tex, ou la commande make)
% =====================================================================
\documentclass[11pt,a4paper,oneside]{book}

% =====================================================================
%  PRÉAMBULE — « ANGLAIS : LE GRAND LIVRE D'AUTO-FORMATION »
%  Méthode complète pour francophones — 12 mois, A1 -> B2/C1
%
%  > Moteur conseillé : pdfLaTeX (2 passes pour la table des matières)
%  > Fonctionne sur Overleaf (TeX Live complet) et en local
%    (TeX Live / MiKTeX). Aucun fichier image externe n'est requis.
%  > Si un paquet manque, MiKTeX le télécharge automatiquement ;
%    sur TeX Live :  sudo tlmgr install tcolorbox titlesec pifont
% =====================================================================

% ---------- Encodage et polices -------------------------------------
\usepackage[T1]{fontenc}
\usepackage[utf8]{inputenc}
\usepackage{lmodern}
\usepackage{microtype}
\usepackage{textcomp}
\usepackage{pifont}

% ---------- Langues : français (principal) + anglais ----------------
\usepackage[english,french]{babel}
% Typographie française correcte ; l'anglais est inséré via \en{} ou
% \selectlanguage{english} pour les longs passages.
\frenchbsetup{
  AutoSpacePunctuation = true,
  StandardItemizeEnv   = true,
  StandardEnumerateEnv = true,
  ThinSpaceInFrenchNumbers = true,
  FrenchFootnotes      = false,
  IndentFirst          = false
}

% ---------- Mise en page --------------------------------------------
\usepackage[a4paper,top=2.2cm,bottom=2.2cm,left=2.1cm,right=2.1cm,
            headheight=14pt,footskip=1.1cm]{geometry}
\usepackage{fancyhdr}
\usepackage{titlesec}
\usepackage{multicol}
\usepackage{longtable}
\usepackage{array}
\usepackage{tabularx}
\usepackage{colortbl}
\usepackage{ragged2e}
\usepackage[usenames,dvipsnames,table]{xcolor}

% ---------- Couleurs du livre ---------------------------------------
\definecolor{bleuNuit}{HTML}{0F3D6E}
\definecolor{bleuClair}{HTML}{E9F1FA}
\definecolor{bleuMoyen}{HTML}{3D7EBF}
\definecolor{vert}{HTML}{1B7A43}
\definecolor{vertClair}{HTML}{E8F6EE}
\definecolor{rouge}{HTML}{AE2B2B}
\definecolor{rougeClair}{HTML}{FCECEC}
\definecolor{orange}{HTML}{C4720A}
\definecolor{orangeClair}{HTML}{FFF3E1}
\definecolor{violet}{HTML}{5B3A8E}
\definecolor{violetClair}{HTML}{F2ECFA}
\definecolor{gris}{HTML}{5A5A5A}
\definecolor{grisClair}{HTML}{F3F3F3}
\definecolor{grisTresClair}{HTML}{FAFAFA}
\definecolor{tabHead}{HTML}{0F3D6E}
\definecolor{tabAlt}{HTML}{F4F7FB}

% ---------- Symboles utilitaires ------------------------------------
\newcommand{\ok}{\ding{51}\,}
\newcommand{\ko}{\ding{55}\,}
\newcommand{\etoile}{\ding{72}\,}
\newcommand{\ra}{\textrightarrow\,}
\newcommand{\Ra}{\ensuremath{\Rightarrow}\,}
\newcommand{\cible}{\ding{119}\,}

% ---------- Phonétique : IPA (si tipa est installé) ------------------
\newif\ifavecTIPA
\IfFileExists{tipa.sty}{\avecTIPAtrue}{\avecTIPAfalse}
%  Voyelles :   \ipaee iː  \ipaiot ɪ  \ipaE ɛ/e  \ipae æ  \ipaa ɑː  \ipao ɔː
%               \ipaups ʊ  \ipaoo uː  \ipalv ʌ  \ipaer ɜː  \ipeschwa ə
%  Diphtongues: \ipaei eɪ  \ipaai aɪ  \ipaOi ɔɪ  \ipaau aʊ  \ipaou əʊ/oʊ
%  Consonnes : \ipatheta θ  \ipadh ð  \ipash ʃ  \ipazh ʒ  \ipang ŋ
%  Divers   :  \ipalong ː   \ipa{...} séquence libre (notation TIPA)
\ifavecTIPA
  \usepackage[safe]{tipa}
  \newcommand{\ipash}{\textipa{S}}        % ʃ
  \newcommand{\ipazh}{\textipa{Z}}        % ʒ
  \newcommand{\ipang}{\textipa{N}}        % ŋ
  \newcommand{\ipeschwa}{\textipa{@}}     % ə
  \newcommand{\ipatheta}{\textipa{T}}     % θ
  \newcommand{\ipadh}{\textipa{D}}        % ð
  \newcommand{\ipae}{\ae}                 % æ
  \newcommand{\ipaE}{\textipa{E}}         % ɛ
  \newcommand{\ipaa}{\textipa{A}}         % ɑ
  \newcommand{\ipao}{\textipa{O}}         % ɔ
  \newcommand{\ipaer}{\textipa{3}}        % ɜ
  \newcommand{\ipaups}{\textipa{U}}       % ʊ
  \newcommand{\ipaiot}{\textipa{I}}       % ɪ
  \newcommand{\ipalv}{\textipa{V}}        % ʌ
  \newcommand{\ipaoo}{\textipa{u}}        % u
  \newcommand{\ipaee}{\textipa{i}}        % i
  \newcommand{\ipalong}{\textipa{:}}      % ː
  \newcommand{\ipaei}{\textipa{eI}}       % eɪ
  \newcommand{\ipaai}{\textipa{aI}}       % aɪ
  \newcommand{\ipaOi}{\textipa{OI}}       % ɔɪ
  \newcommand{\ipaau}{\textipa{aU}}       % aʊ
  \newcommand{\ipaou}{\textipa{@U}}       % əʊ
  \newcommand{\ipaoou}{\textipa{oU}}      % oʊ (AmE)
  \newcommand{\ipa}[1]{\textipa{#1}}
\else
  % ---- Repli : transcription lisible si le paquet tipa est absent ----
  \newcommand{\ipash}{ch}
  \newcommand{\ipazh}{j}
  \newcommand{\ipang}{ng}
  \newcommand{\ipeschwa}{\textsf{@}}
  \newcommand{\ipatheta}{\ensuremath{\theta}}
  \newcommand{\ipadh}{dh}
  \newcommand{\ipae}{\ae}
  \newcommand{\ipaE}{\`e}
  \newcommand{\ipaa}{a:}
  \newcommand{\ipao}{o:}
  \newcommand{\ipaer}{eu:}
  \newcommand{\ipaups}{ou}
  \newcommand{\ipaiot}{i}
  \newcommand{\ipalv}{e}
  \newcommand{\ipaoo}{ou:}
  \newcommand{\ipaee}{i:}
  \newcommand{\ipalong}{:}
  \newcommand{\ipaei}{\'e-i}
  \newcommand{\ipaai}{a-i}
  \newcommand{\ipaOi}{o-i}
  \newcommand{\ipaau}{a-ou}
  \newcommand{\ipaou}{o-ou}
  \newcommand{\ipaoou}{o-ou}
  \newcommand{\ipa}[1]{#1}
\fi

% Prononciation « à la française » : toujours disponible, sans paquet
\newcommand{\pron}[1]{{\color{violet}\sffamily\small[#1]}}

% ---------- Anglais en ligne et blocs bilingues ---------------------
% \en{}  : passage en anglais (typographie anglaise, pas d'espace avant ! ? :)
% \enb{} : idem en gras   \eni{} : idem en italique
\newcommand{\en}[1]{\foreignlanguage{english}{#1}}
\newcommand{\enb}[1]{\foreignlanguage{english}{\textbf{#1}}}
\newcommand{\eni}[1]{\foreignlanguage{english}{\textit{#1}}}

% Exemple bilingue : anglais + traduction française en retrait
\newcommand{\ex}[2]{%
  \par\smallskip
  {\selectlanguage{english}#1\par}%
  \hspace*{1.2em}{\footnotesize\color{gris}\selectlanguage{french}\itshape #2}\par}

% Phrase utile (tableaux) : anglais | français
\newcommand{\phr}[2]{\en{#1} & \selectlanguage{french}#2 \tabularnewline}

% Mot de vocabulaire : mot | prononciation française | sens
\newcommand{\mot}[3]{\textbf{\en{#1}} & \pron{#2} & \selectlanguage{french}#3 \tabularnewline}

% Ligne de dialogue : personnage, réplique anglaise, traduction
\newcommand{\dial}[3]{%
  \par\smallskip
  {\bfseries\color{bleuNuit}\selectlanguage{french}#1}\enspace\textcolor{orange}{\textbullet}\enspace
  {\selectlanguage{english}#2\par}%
  \hspace*{1.6em}{\footnotesize\color{gris}\selectlanguage{french}\itshape #3}\par}

% Réponse courte / variante dans un dialogue
% Flèches d'intonation : sûres en mode texte comme en mode maths
\newcommand{\monte}{\ensuremath{\nearrow}}
\newcommand{\descend}{\ensuremath{\searrow}}

\newcommand{\variante}[1]{\par\hspace*{1.6em}{\footnotesize\selectlanguage{english}\textcolor{bleuMoyen}{#1}}\par}

% Trou à compléter dans les exercices
\newcommand{\trou}[1][2.2cm]{\rule{#1}{0.4pt}}
\newcommand{\troul}[1]{\underline{\hspace{#1}}}

% ---------- Boîtes colorées (tcolorbox si disponible) ----------------
\IfFileExists{tcolorbox.sty}{%
  \usepackage[skins,breakable]{tcolorbox}
  \tcbset{
    boite/.style={enhanced jigsaw, breakable, boxrule=0pt,
      leftrule=3pt, arc=0pt, outer arc=0pt, left=7pt, right=7pt,
      top=6pt, bottom=6pt, toptitle=4pt, bottomtitle=4pt,
      fonttitle=\bfseries\selectlanguage{french},
      coltitle=white, fontupper=\small, fontlower=\small},
    sansTitre/.style={enhanced jigsaw, breakable, boxrule=0.5pt, arc=2pt,
      left=6pt, right=6pt, top=5pt, bottom=5pt, fontupper=\small}
  }
  \newtcolorbox{regle}[1][]{boite, colback=bleuClair, colframe=bleuClair,
    borderline west={3pt}{0pt}{bleuNuit}, coltitle=white,
    title={\selectlanguage{french}\ok RÈGLE #1}, colbacktitle=bleuNuit}
  \newtcolorbox{piege}[1][]{boite, colback=rougeClair, colframe=rougeClair,
    borderline west={3pt}{0pt}{rouge},
    title={\selectlanguage{french}\ko PIÈGE FRANCOPHONE #1}, colbacktitle=rouge}
  \newtcolorbox{astuce}[1][]{boite, colback=vertClair, colframe=vertClair,
    borderline west={3pt}{0pt}{vert},
    title={\selectlanguage{french}\etoile ASTUCE #1}, colbacktitle=vert}
  \newtcolorbox{mission}[1][]{boite, colback=orangeClair, colframe=orangeClair,
    borderline west={3pt}{0pt}{orange},
    title={\selectlanguage{french}\cible MISSION DU JOUR #1}, colbacktitle=orange}
  \newtcolorbox{memo}[1][]{boite, colback=violetClair, colframe=violetClair,
    borderline west={3pt}{0pt}{violet},
    title={\selectlanguage{french}\ding{46}\ À RETENIR #1}, colbacktitle=violet}
  \newtcolorbox{dialoguebox}[1][]{sansTitre, colback=grisTresClair, colframe=bleuMoyen,
    title={\selectlanguage{french}\bfseries DIALOGUE #1}, coltitle=white,
    colbacktitle=bleuMoyen}
  \newtcolorbox{exo}[1][]{sansTitre, colback=white, colframe=bleuNuit,
    title={\selectlanguage{french}\bfseries EXERCICE #1}, coltitle=white,
    colbacktitle=bleuNuit}
  \newtcolorbox{corrige}[1][]{sansTitre, colback=vertClair, colframe=vert,
    title={\selectlanguage{french}\bfseries CORRIGÉ #1}, coltitle=white,
    colbacktitle=vert}
  \newtcolorbox{tableau}[1][]{sansTitre, colback=white, colframe=gris,
    title={\selectlanguage{french}\bfseries #1}, coltitle=white,
    colbacktitle=gris}
}{%
  % -------- Repli sans tcolorbox : contenu toujours lisible ----------
  \newenvironment{regle}[1][]{\par\medskip\noindent\begin{minipage}{\linewidth}
      \color{bleuNuit}\bfseries\selectlanguage{french}\ok RÈGLE #1\par\normalcolor\small\selectlanguage{french}}
      {\end{minipage}\par\medskip}
  \newenvironment{piege}[1][]{\par\medskip\noindent\begin{minipage}{\linewidth}
      \color{rouge}\bfseries\selectlanguage{french}\ko PIÈGE FRANCOPHONE #1\par\normalcolor\small\selectlanguage{french}}
      {\end{minipage}\par\medskip}
  \newenvironment{astuce}[1][]{\par\medskip\noindent\begin{minipage}{\linewidth}
      \color{vert}\bfseries\selectlanguage{french}\etoile ASTUCE #1\par\normalcolor\small\selectlanguage{french}}
      {\end{minipage}\par\medskip}
  \newenvironment{mission}[1][]{\par\medskip\noindent\begin{minipage}{\linewidth}
      \color{orange}\bfseries\selectlanguage{french}\cible MISSION DU JOUR #1\par\normalcolor\small\selectlanguage{french}}
      {\end{minipage}\par\medskip}
  \newenvironment{memo}[1][]{\par\medskip\noindent\begin{minipage}{\linewidth}
      \color{violet}\bfseries\selectlanguage{french}\ding{46}\ À RETENIR #1\par\normalcolor\small\selectlanguage{french}}
      {\end{minipage}\par\medskip}
  \newenvironment{dialoguebox}[1][]{\par\medskip\noindent\begin{minipage}{\linewidth}
      \bfseries\selectlanguage{french}DIALOGUE #1\par\mdseries\small}
      {\end{minipage}\par\medskip}
  \newenvironment{exo}[1][]{\par\medskip\noindent\begin{minipage}{\linewidth}
      \bfseries\selectlanguage{french}EXERCICE #1\par\mdseries\small}
      {\end{minipage}\par\medskip}
  \newenvironment{corrige}[1][]{\par\medskip\noindent\begin{minipage}{\linewidth}
      \bfseries\selectlanguage{french}CORRIGÉ #1\par\mdseries\small}
      {\end{minipage}\par\medskip}
  \newenvironment{tableau}[1][]{\par\medskip\noindent\begin{minipage}{\linewidth}
      \bfseries\selectlanguage{french}#1\par\mdseries\small}
      {\end{minipage}\par\medskip}
}

% ---------- Environnements simples (sans paquet) ---------------------
% Bloc 100 % anglais (dialogues, textes, listes en anglais)
\newenvironment{anglais}{\selectlanguage{english}\par\smallskip}{\par\smallskip\selectlanguage{french}}
% Bloc de vocabulaire en trois colonnes
\newenvironment{tabvocab}
  {\par\smallskip\selectlanguage{french}\noindent
   \begin{longtable}{@{}p{4.1cm}p{3.1cm}p{8.4cm}@{}}
   \rowcolor{tabHead}\color{white}\bfseries Mot & \bfseries Prononciation & \bfseries Sens, exemple \\[2pt]
   \endfirsthead
   \rowcolor{tabHead}\color{white}\bfseries Mot & \bfseries Prononciation & \bfseries Sens, exemple \\[2pt]
   \endhead}
  {\end{longtable}\par\smallskip}

% ---------- Titres (parties, chapitres, sections) --------------------
\titleformat{\part}[display]
  {\normalfont\sffamily\color{bleuNuit}}
  {\filleft\huge\bfseries\partname\ \thepart}{12pt}
  {\filcenter\Huge\bfseries}
  [\vspace{6pt}{\color{orange}\titlerule[2pt]}]
\titleformat{\chapter}[hang]
  {\normalfont\sffamily\Huge\bfseries\color{bleuNuit}}
  {\thechapter\enspace\textcolor{orange}{\vrule width 4pt height 24pt}\enspace}{0pt}{\Huge}
\titlespacing*{\chapter}{0pt}{-18pt}{22pt}
\titleformat{\section}
  {\normalfont\sffamily\Large\bfseries\color{bleuNuit}}{\thesection}{0.8em}{}
\titlespacing*{\section}{0pt}{16pt}{8pt}
\titleformat{\subsection}
  {\normalfont\sffamily\large\bfseries\color{orange}}{\thesubsection}{0.8em}{}
\titlespacing*{\subsection}{0pt}{12pt}{6pt}
\titleformat{\subsubsection}
  {\normalfont\sffamily\normalsize\bfseries\color{gris}}{\thesubsubsection}{0.8em}{}
\titlespacing*{\subsubsection}{0pt}{10pt}{4pt}

% ---------- En-têtes et pieds de page --------------------------------
\pagestyle{fancy}
\fancyhf{}
\fancyhead[L]{\footnotesize\sffamily\color{gris}\nouppercase{\rightmark}}
\fancyhead[R]{\footnotesize\sffamily\color{bleuNuit}\thepage}
\fancyfoot[L]{\footnotesize\sffamily\color{gris}Anglais --- le grand livre d'auto-formation}
\renewcommand{\headrulewidth}{0.4pt}
\renewcommand{\headrule}{\hbox to\headwidth{\color{bleuMoyen}\leaders\hrule height \headrulewidth\hfill}}
\renewcommand{\footrulewidth}{0pt}
\fancypagestyle{plain}{%
  \fancyhf{}\fancyfoot[C]{\footnotesize\sffamily\thepage}%
  \renewcommand{\headrulewidth}{0pt}}

% ---------- Listes plus aérées ---------------------------------------
\usepackage{enumitem}
\setlist[itemize]{leftmargin=1.4em,itemsep=2pt,topsep=4pt,parsep=0pt}
\setlist[enumerate]{leftmargin=1.8em,itemsep=2pt,topsep=4pt,parsep=0pt}

% ---------- Tableaux : style par défaut ------------------------------
\setlength{\tabcolsep}{5pt}
\renewcommand{\arraystretch}{1.25}
\newcolumntype{L}[1]{>{\RaggedRight\arraybackslash}p{#1}}
\newcolumntype{C}[1]{>{\Centering\arraybackslash}p{#1}}
\newcommand{\entete}[1]{\rowcolor{tabHead}\color{white}\bfseries #1}

% ---------- Liens internes -------------------------------------------
\usepackage[colorlinks=true,linkcolor=bleuNuit,urlcolor=bleuMoyen,
            citecolor=vert,pdftitle={Anglais - le grand livre d'auto-formation},
            pdfauthor={},pdfsubject={Methode complete d'anglais pour francophones},
            bookmarksopen=true,bookmarksnumbered=true]{hyperref}

% ---------- Divers ---------------------------------------------------
\setlength{\parskip}{4pt}
\setlength{\parindent}{0pt}
\raggedbottom
\sloppy
\newcommand{\pagevierge}{\clearpage}


\begin{document}

% ------------------------------------------------------------------
%  PAGES LIMINAIRES
% ------------------------------------------------------------------
\frontmatter
% =====================================================================
%  PAGES LIMINAIRES : titre, promesse, mode d'emploi, feuille de route
% =====================================================================
\begin{titlepage}
\thispagestyle{empty}
\begin{center}

\vspace*{0.6cm}
{\color{orange}\rule{\linewidth}{2.5pt}}\\[0.9cm]

{\sffamily\bfseries\color{bleuNuit}\fontsize{34}{40}\selectfont ANGLAIS\par}
\vspace{0.35cm}
{\sffamily\bfseries\color{bleuNuit}\fontsize{21}{26}\selectfont Le grand livre d'auto-formation\par}
\vspace{0.5cm}
{\color{gris}\large\itshape 12 mois pour parler anglais\\ comme vous parlez français\par}

\vspace{0.7cm}
{\color{orange}\rule{0.45\linewidth}{1.5pt}}

\vspace{0.9cm}
{\sffamily\large\color{bleuNuit}\bfseries Méthode complète pour francophones\par}
\vspace{0.25cm}
{\large Niveau de départ A1/A2 \quad\ra\quad Objectif B2 (conversation courante)\par}

\vspace{1.1cm}

\begin{minipage}{0.92\linewidth}
\centering
\color{bleuNuit}\large\bfseries
Ce que contient ce livre
\end{minipage}

\vspace{0.35cm}

\begin{minipage}{0.94\linewidth}
\small
\begin{tabular}{@{}p{7.3cm}p{7.3cm}@{}}
\ok Un plan de 12 mois, semaine par semaine & \ok Plus de 3\,000 mots et 800 phrases prêtes à l'emploi \\[4pt]
\ok Toute la grammaire, expliquée en français & \ok 40 situations de conversation, dialogues modèles et jeux de rôle \\[4pt]
\ok La prononciation : 44 sons, accent, intonation & \ok 150 \en{phrasal verbs}, 120 expressions idiomatiques \\[4pt]
\ok Les 100 erreurs typiques des francophones & \ok Des centaines d'exercices, tous corrigés \\[4pt]
\ok La méthode d'écoute, de lecture et d'écriture & \ok Tests de niveau et grilles d'auto-évaluation \\
\end{tabular}
\end{minipage}

\vfill

{\color{gris}\small Édition 1.0 --- Document LaTeX librement modifiable et imprimable.\par}
\vspace{0.3cm}
{\color{orange}\rule{\linewidth}{2.5pt}}

\end{center}
\end{titlepage}

% ---------------------------------------------------------------------
\chapter*{Avant de commencer}
\addcontentsline{toc}{chapter}{Avant de commencer}
\markboth{Avant de commencer}{Avant de commencer}

Vous avez déjà des notions d'anglais. Vous comprenez peut-être un peu, vous lisez des mots
isolés, mais dès qu'il faut \emph{parler}, tout se bloque : les mots ne viennent pas, la
prononciation vous gêne, vous traduisez du français dans votre tête et vous arrivez deux
secondes trop tard.

Ce livre a été écrit pour exactement ce problème. Il ne s'agit pas d'un manuel scolaire de
plus, mais d'un \textbf{programme d'entraînement complet} : une méthode, un calendrier, des
contenus et des exercices, organisés pour qu'en douze mois vous passiez de « je connais
un peu l'anglais » à « je parle anglais, je comprends ce qu'on me dit et je réponds sans
hésiter ».

\begin{memo}[La promesse de ce livre]
Si vous suivez le plan du chapitre 2 --- \textbf{90 minutes par jour, six jours sur sept} ---
pendant douze mois, vous aurez accumulé environ \textbf{470 heures d'entraînement dirigé}
et \textbf{3\,000 heures de contact} avec l'anglais si vous ajoutez l'immersion passive
recommandée. C'est le volume qui sépare un francophone débutant d'un locuteur \textbf{B2
solide}, capable de tenir une conversation de 30 minutes avec un anglophone, de suivre un
film sans sous-titres français et de travailler en anglais.
\end{memo}

\section*{Pourquoi un francophone peut réussir vite}
Trois bonnes nouvelles, à lire les jours de découragement :

\begin{enumerate}
  \item \textbf{L'anglais et le français partagent environ 30 à 45\,\% de leur vocabulaire.}
  Depuis 1066, le français a massivement irrigué l'anglais : \en{government}, \en{language},
  \en{important}, \en{possible}, \en{different}, \en{information}, \en{nature}, \en{culture},
  \en{justice}, \en{village}, \en{table}, \en{fruit}, \en{hotel}. Vous connaissez déjà des
  milliers de mots anglais sans le savoir. Le chapitre 12 vous apprend à exploiter ce trésor.
  \item \textbf{La grammaire anglaise est plus simple que la grammaire française.} Pas de
  genre sur les noms, pas d'accord des adjectifs, pas de subjonctif à mémoriser, une seule
  forme de verbe au présent (\en{-s} à la 3\textsuperscript{e} personne), pas de vouvoiement.
  Ce qui est difficile en anglais, ce sont \emph{les temps} (le \en{present perfect} !)
  et \emph{la prononciation}. C'est précisément là que ce livre insiste.
  \item \textbf{L'anglais est partout.} Musique, séries, football, technologie, réseaux
  sociaux, religion, commerce : vous pouvez trouver de l'anglais gratuit toute la journée.
  L'immersion ne vous coûtera presque rien --- seulement de la discipline.
\end{enumerate}

\section*{Les trois pièges qui font échouer (et comment les éviter)}
\begin{tableau}[Ce qui fait échouer 90 \% des autodidactes]
\begin{tabularx}{\linewidth}{@{}p{3.9cm}X@{}}
\entete{Erreur} & \entete{Correction apportée par ce livre} \tabularnewline
\hline
\textbf{Étudier la grammaire sans jamais parler} &
Chaque unité se termine par une \textbf{mission orale} : dire 10 phrases à voix haute,
se raconter quelque chose, imiter un dialogue. La parole est entraînée dès le jour 1. \tabularnewline
\hline
\textbf{Apprendre des listes de mots isolés} &
Vous apprendrez des \textbf{blocs} (\en{I'm looking forward to it}) et non des mots
seuls, avec la prononciation, et vous les reverrez grâce à la répétition espacée
(chapitre 3). \tabularnewline
\hline
\textbf{Recommencer dix fois le niveau débutant} &
Un \textbf{calendrier daté} avec des jalons : mois 1, mois 3, mois 6, mois 12. Vous savez
toujours où vous en êtes, et les tests du chapitre 19 le prouvent. \tabularnewline
\end{tabularx}
\end{tableau}

\section*{Comment lire ce livre}
Ce livre utilise un code couleur constant. En voici la légende --- gardez cette page en tête
pendant toute votre année :

\begin{regle}
Une règle de grammaire ou de prononciation, formulée simplement, avec sa structure.
\end{regle}

\begin{piege}
Un endroit où les francophones se trompent presque toujours. Ce sont les encadrés les plus
rentables du livre : chaque piège évité, c'est une faute que vous ne ferez plus jamais.
\end{piege}

\begin{astuce}
Un raccourci de mémorisation, une astuce de prononciation, un moyen de gagner du temps.
\end{astuce}

\begin{mission}
Ce que vous devez \textbf{faire} aujourd'hui, à voix haute. Un livre de langue ne s'apprend
pas avec les yeux : une mission non faite est une leçon non apprise.
\end{mission}

\begin{memo}
La synthèse de l'unité : à relire la veille d'un test ou dans le bus.
\end{memo}

\begin{exo}[X --- modèle]
Les exercices sont numérotés par unité. Le \textbf{corrigé suit immédiatement} l'exercice
(encadré vert), pour que vous puissiez travailler seul sans attendre. Cachez-le avec une
feuille si vous préférez chercher d'abord.
\end{exo}

\begin{corrige}
Voici à quoi ressemble un corrigé.
\end{corrige}

% Conventions d'écriture
\begin{tableau}[Conventions typographiques]
\begin{tabularx}{\linewidth}{@{}p{5.4cm}X@{}}
\entete{Ce que vous voyez} & \entete{Ce que cela veut dire} \tabularnewline
\hline
\en{I have been working.} & Texte en \textbf{anglais} (à lire à voix haute). \tabularnewline
\hline
\textit{\color{gris}J'ai travaillé / je travaille depuis\ldots} & Traduction française, en petit et en gris. \tabularnewline
\hline
\pron{aïe have bin oueur-king} & \textbf{Prononciation approchée à la française}.
Indispensable les premiers mois ; passez ensuite à l'écoute réelle (chapitre 16). \tabularnewline
\hline
\trou \quad ou \quad \trou[1.2cm] & Trou à compléter dans un exercice. \tabularnewline
\hline
\en{*I have 25 years.} & L'astérisque signale une phrase \textbf{incorrecte}, à ne pas imiter. \tabularnewline
\hline
\en{BrE} / \en{AmE} & Britannique / américain. Les deux sont donnés quand ils diffèrent. \tabularnewline
\end{tabularx}
\end{tableau}

\begin{astuce}[Le mot \en{to be} ou \en{to have} ?]
Dans ce livre, tout exemple anglais est suivi de sa traduction. Ne lisez \textbf{jamais}
un exemple sans le dire à voix haute au moins une fois. Votre bouche doit s'habituer, pas
seulement vos yeux.
\end{astuce}

% ---------------------------------------------------------------------
\chapter*{Feuille de route en un coup d'\oe il}
\addcontentsline{toc}{chapter}{Feuille de route en un coup d'\oe il}
\markboth{Feuille de route}{Feuille de route}

Voici votre année, résumée sur une page. Le détail jour par jour se trouve au chapitre 2.

\begin{tableau}[Les 12 mois --- objectifs, contenu, jalon de validation]
\footnotesize
\begin{longtable}{@{}p{1.7cm}p{2.1cm}p{6.3cm}p{4.6cm}@{}}
\entete{Mois} & \entete{Niveau visé} & \entete{Ce que vous travaillez} & \entete{Vous devez savoir faire} \tabularnewline
\endfirsthead
\entete{Mois} & \entete{Niveau visé} & \entete{Ce que vous travaillez} & \entete{Vous devez savoir faire} \tabularnewline
\endhead
\hline
\textbf{1} & A1 --- démarrage &
Sons de l'anglais, accent tonique, \en{be}/\en{have}, présent simple, se présenter,
300 mots, écoute quotidienne 15 min &
Se présenter en 2 min, dire 200 phrases de base, épeler son nom \tabularnewline
\hline
\textbf{2} & A1 --- consolidation &
Questions, négation, articles, pluriels, \en{there is}, prépositions de lieu/temps,
vie quotidienne, 600 mots &
Raconter sa journée, poser 30 questions types \tabularnewline
\hline
\textbf{3} & A1\ra A2 &
Passé simple (verbes irréguliers), \en{past continuous}, futur \en{going to}/\en{will},
comparatifs, 900 mots, 1\textsuperscript{er} dialogue réel &
Raconter un souvenir, décrire un projet, tenir 5 min de conversation guidée \tabularnewline
\hline
\textbf{4} & A2 &
\en{Present perfect} (le grand chapitre), \en{for}/\en{since}, \en{used to},
modaux de base, ville/transports/santé, 1\,200 mots &
Parler de son expérience, réserver, demander son chemin, chez le médecin \tabularnewline
\hline
\textbf{5} & A2 &
Tous les temps en revue, voix passive simple, \en{should}/\en{must}/\en{have to},
travail et métiers, 1\,500 mots, lecture d'un premier roman simplifié &
Décrire un emploi, une règle, un processus ; lire un texte A2 sans dictionnaire \tabularnewline
\hline
\textbf{6} & A2\ra B1 &
\textbf{Bilan à mi-parcours.} Conditionnels 0/1/2, relatifs, gérondif/infinitif,
\en{phrasal verbs} (1\textsuperscript{re} série), 1\,800 mots &
Donner son opinion, nuancer, comprendre une série avec sous-titres anglais \tabularnewline
\hline
\textbf{7} & B1 &
Conditionnel 3, discours rapporté, \en{wish}, inversion, quantifieurs,
médias et actualités, 2\,100 mots &
Raconter une anecdote complexe, rapporter des propos, débattre 10 min \tabularnewline
\hline
\textbf{8} & B1 &
Voix passive complète, \en{have something done}, \en{cleft sentences},
modaux du passé, voyages/aéroport/hôtel, 2\,400 mots &
Se débrouiller en voyage de A à Z, présenter un exposé de 5 min \tabularnewline
\hline
\textbf{9} & B1\ra B2 &
Connecteurs avancés, collocations, expressions idiomatiques, registre
(formel/familier), 2\,700 mots, écriture d'e-mails et de lettres &
Écrire un e-mail professionnel, adapter son registre, comprendre un podcast clair \tabularnewline
\hline
\textbf{10} & B2 &
Anglais professionnel : réunion, entretien d'embauche, téléphone, négociation,
présentation ; 3\,000 mots &
Passer un entretien, animer une réunion, téléphoner sans panique \tabularnewline
\hline
\textbf{11} & B2 &
Fluidité : débit, \emph{fillers}, reformulation, humour, accents (britannique,
américain, africain, indien), 30 min de conversation libre par semaine &
Parler 20 min sans préparation sur un sujet inconnu \tabularnewline
\hline
\textbf{12} & B2\ra C1 &
Révision générale, test blanc complet, certification éventuelle (TOEIC/IELTS),
consolidation des points faibles &
Tenir une conversation d'une heure, regarder un film sans sous-titres \tabularnewline
\end{longtable}
\end{tableau}

\begin{mission}[Avant même de commencer]
\begin{enumerate}
  \item Passez le \textbf{test de positionnement} du chapitre 19 (40 min). Il vous dira si
  vous démarrez au mois 1 ou au mois 3.
  \item Choisissez votre \textbf{créneau quotidien} et notez-le : \trou[3cm] h \ra \trou[2cm] h.
  \item Installez \textbf{une application de répétition espacée} (Anki, gratuite) et
  \textbf{un dictionnaire} (Cambridge ou WordReference).
  \item Mettez votre \textbf{téléphone en anglais} aujourd'hui (Réglages \ra Langue).
  Cinq minutes de courage, un an de bénéfices.
  \item Dites à voix haute, maintenant : \en{``In one year, I will speak English every day
  without fear.''} --- \textit{Dans un an, je parlerai anglais tous les jours sans peur.}
\end{enumerate}
\end{mission}


\tableofcontents
\pagevierge

% ------------------------------------------------------------------
%  CORPS DU LIVRE
% ------------------------------------------------------------------
\mainmatter

% ===== PARTIE I — LE PLAN DE BATAILLE ==============================
\part{Le plan de bataille : devenir vraiment bilingue en 12 mois}
% =====================================================================
%  CHAPITRES 1 ET 2 — Comment on devient bilingue + le plan de 12 mois
% =====================================================================
\chapter{Comment on devient vraiment capable de parler}

\begin{quote}
\itshape
On ne devient pas bilingue en \textbf{comprenant} une langue. On devient bilingue en
\textbf{produisant} une langue, beaucoup, souvent, sans attendre d'être parfait.
\end{quote}

Avant de toucher à la grammaire, il faut comprendre comment le cerveau fabrique une langue.
Cinq minutes de lecture ici vous feront gagner des mois.

\section{Ce que dit la recherche (en quatre idées simples)}

\subsection{Idée 1 --- L'input compréhensible : comprendre à 90\,\%}
Le linguiste Stephen Krashen a montré qu'on acquiert une langue quand on comprend des
messages \emph{légèrement} au-dessus de son niveau. Ni trop facile (on n'apprend rien),
ni trop difficile (on décroche). La règle pratique :

\begin{memo}[La règle des 90\,\%]
Choisissez un texte ou un audio où vous comprenez \textbf{9 mots sur 10}. Si vous comprenez
moins de 8 mots sur 10, c'est trop dur : descendez d'un niveau. Si vous comprenez
10 mots sur 10, c'est trop facile : montez.
\end{memo}

\subsection{Idée 2 --- L'output : parler n'est pas la conséquence de l'apprentissage, c'en est le moteur}
Pendant longtemps on a cru : « d'abord je stocke, ensuite je parle ». Faux. Parler
\emph{fabrique} la compétence : quand vous cherchez un mot et ne le trouvez pas, votre
cerveau enregistre un « trou » qu'il comblera durablement la prochaine fois qu'il rencontrera
le mot. C'est ce qu'on appelle l'\emph{effet de production} et l'\emph{hypothèse de
l'output} (Merrill Swain). Conséquence pour vous : \textbf{parlez dès la première semaine},
même mal, même seul.

\subsection{Idée 3 --- La répétition espacée : revoir juste avant d'oublier}
Vous oubliez 70\,\% d'un mot nouveau en 24 heures si vous ne le revoyez pas. Mais si vous le
revoyez à J+1, J+3, J+7, J+15, J+30, il reste. C'est le principe des logiciels de
répétition espacée (\en{spaced repetition systems}, SRS) comme \textbf{Anki}. Un mot appris
sans SRS coûte 10 fois plus cher en temps qu'un mot appris avec SRS.

\subsection{Idée 4 --- Le sommeil et la régularité battent l'intensité}
Deux heures par jour pendant un an (730 h) valent infiniment mieux que dix heures par jour
pendant trois semaines. La consolidation de la mémoire se fait la nuit ; le cerveau a besoin
de contacts \emph{répétés}. Visez la \textbf{chaîne ininterrompue} : 365 jours d'affilée
valent mieux que 250 jours parfaits.

\section{Combien d'heures faut-il ? (le calcul honnête)}
Le Cadre européen commun de référence pour les langues (CECRL) et les estimations de
l'institut américain FSI donnent des ordres de grandeur fiables pour un francophone :

\begin{tableau}[Heures d'étude pour un francophone --- anglais]
\begin{tabularx}{\linewidth}{@{}p{2.2cm}p{3.4cm}X@{}}
\entete{Niveau} & \entete{Volume cumulé} & \entete{Ce que vous savez faire} \tabularnewline
\hline
\textbf{A1} & 80 à 120 h & Se présenter, formules de survie, phrases très simples. \tabularnewline
\textbf{A2} & 200 à 300 h & Échanges courants : achats, famille, travail, trajets, passé et futur simples. \tabularnewline
\textbf{B1} & 400 à 550 h & Se débrouiller en voyage, raconter, expliquer, donner un avis, suivre une série avec sous-titres anglais. \tabularnewline
\textbf{B2} & 700 à 900 h & Parler avec aisance et spontanéité, argumenter, comprendre les informations et la plupart des films. \textbf{C'est l'objectif de ce livre.} \tabularnewline
\textbf{C1} & 1\,000 à 1\,400 h & Nuancer, plaisanter, rédiger des textes complexes, travailler en anglais sans effort. \tabularnewline
\end{tabularx}
\end{tableau}

\begin{astuce}[Votre budget temps]
90 minutes par jour $\times$ 6 jours $\times$ 52 semaines = \textbf{468 heures d'étude active}.
Ajoutez l'immersion passive (musique, séries, téléphone en anglais, radio) : environ
\textbf{2 heures par jour}, soit 730 heures de contact. Total annuel : plus de
\textbf{1\,100 heures d'exposition}. Avec un départ A1/A2, cela vous place \textbf{entre B2
et C1}. Le chiffre n'est pas une promesse : c'est la seule chose que vous contrôlez.
\end{astuce}

\section{Les cinq systèmes à faire tourner en parallèle}
Une langue n'est pas une matière scolaire, c'est cinq compétences qui progressent à des
vitesses différentes. Voici les cinq « moteurs » et la part de votre temps quotidien à leur
consacrer.

\begin{tableau}[Les 5 moteurs et leur poids dans votre journée de 90 min]
\begin{tabularx}{\linewidth}{@{}p{3.6cm}p{2.0cm}X@{}}
\entete{Moteur} & \entete{Temps} & \entete{Activités concrètes} \tabularnewline
\hline
\textbf{1. Oreille} (compréhension orale) & 25 min &
Écoute intensive (5 passes sur 2 min), dictée, podcasts, séries, chansons. \tabularnewline
\textbf{2. Bouche} (production orale) & 25 min &
Shadowing, lecture à voix haute, auto-conversation, dialogues du livre, appels avec
partenaire. \tabularnewline
\textbf{3. Vocabulaire} & 15 min &
Anki (révisions + 15 mots nouveaux), blocs de phrases, familles de mots. \tabularnewline
\textbf{4. Grammaire / lecture} & 15 min &
Une unité de ce livre + 5 exercices + un texte court. \tabularnewline
\textbf{5. Écriture} & 10 min &
Journal : 5 à 10 phrases sur votre journée, puis correction. \tabularnewline
\end{tabularx}
\end{tableau}

\begin{piege}[Les francophones négligent le moteur n°2]
L'école africaine et française enseigne l'anglais \emph{par l'écrit} : on lit, on conjugue,
on fait des thèmes. Résultat : des adultes qui savent accorder un \en{present perfect} mais
qui restent muets. Ce livre inverse la proportion : \textbf{un tiers de votre temps doit
sortir de votre bouche}, chaque jour, dès le premier jour.
\end{piege}

\section{Les six mythes qui font perdre du temps}
\begin{enumerate}
  \item \textbf{« Je suis nul en langues.} » Faux : vous avez appris le français, le wolof,
  ou une autre langue sans manuel, en 3 ans, en écoutant et en osant. Vous savez faire.
  \item \textbf{« Il faut un professeur particulier.} » Utile, pas indispensable. Un
  partenaire de conversation gratuit (application, club, église, camarade) 2 fois par
  semaine suffit si vous travaillez sérieusement seul le reste du temps.
  \item \textbf{« Il faut vivre dans un pays anglophone.} » L'immersion aide énormément,
  mais on peut fabriquer une immersion artificielle chez soi (chapitre 4). Des millions de
  personnes ont atteint C1 sans quitter leur pays.
  \item \textbf{« Je dois connaître tous les mots.} » 3\,000 mots couvrent 95\,\% d'une
  conversation quotidienne. 1\,000 mots en couvrent 85\,\%. La quantité ne sert à rien sans
  les mots \emph{utiles}.
  \item \textbf{« Je parlerai quand je serai prêt.} » On n'est jamais prêt. On parle, donc
  on devient prêt.
  \item \textbf{« L'accent parfait ou rien.} » Votre accent français (ou sénégalais) n'est
  pas un défaut : c'est votre identité. L'objectif est l'\textbf{intelligibilité} --- qu'on
  vous comprenne du premier coup --- pas l'imitation d'un acteur de Hollywood.
\end{enumerate}

\section{Les trois leviers de la motivation sur douze mois}
\begin{itemize}
  \item \textbf{Le levier visuel} : imprimez l'annexe F (journal de progression) et collez-la
  au mur. Une case cochée par jour. Une chaîne de 365 cases est psychologiquement
  imparable : on ne veut pas la casser.
  \item \textbf{Le levier social} : annoncez votre objectif à deux personnes. Mieux :
  trouvez un partenaire d'apprentissage (même à distance) et fixez un rendez-vous hebdomadaire.
  \item \textbf{Le levier de la preuve} : enregistrez-vous le jour 1 (\en{``My name is\ldots''})
  et conservez le fichier. Réenregistrez-vous chaque mois. Écouter sa propre voix à trois
  mois d'intervalle est la meilleure motivation qui existe.
\end{itemize}

\begin{mission}[Jour 0 --- aujourd'hui, 20 minutes]
\begin{enumerate}
  \item Enregistrez-vous sur votre téléphone en anglais pendant 2 minutes : présentez-vous,
  dites ce que vous faites dans la vie, ce que vous voulez dans un an. Nommez le fichier
  \en{jour-000.m4a}. Ne l'effacez jamais.
  \item Écrivez en dessous de l'enregistrement, en français, \textbf{pourquoi} vous voulez
  parler anglais (travail, religion, études, famille, voyages, Internet). Soyez précis.
  Relisez cette phrase chaque lundi.
  \item Cochez la première case de votre journal de progression.
\end{enumerate}
\end{mission}


% =====================================================================
\chapter{Le plan des 12 mois, mois par mois}

Ce chapitre est le c\oe ur opérationnel du livre. Il vous dit \textbf{quoi faire, quand,
combien de temps et comment vérifier}. Tout le reste du livre est au service de ce plan.

\section{La journée type (90 minutes)}
Voici le squelette quotidien. Adaptez les horaires, gardez les proportions.

\begin{tableau}[Votre séance quotidienne --- 90 minutes]
\begin{tabularx}{\linewidth}{@{}p{1.6cm}p{4.4cm}X@{}}
\entete{Ordre} & \entete{Bloc} & \entete{Détail} \tabularnewline
\hline
1 & \textbf{Échauffement} (5 min) &
Lecture à voix haute d'un paragraphe déjà connu (du chapitre en cours ou d'un dialogue
appris). Objectif : remettre la bouche en mouvement, sans réfléchir. \tabularnewline
2 & \textbf{Écoute intensive} (20 min) &
Un audio court (1 à 3 min) travaillé en 5 passes : écoute globale \ra écoute avec le texte
\ra écoute sans le texte \ra dictée d'une phrase \ra shadowing. Protocole complet au
chapitre 16. \tabularnewline
3 & \textbf{Vocabulaire SRS} (15 min) &
Révisions Anki du jour + 15 mots/blocs nouveaux issus de l'unité du livre. Toujours à voix
haute, toujours en phrase. \tabularnewline
4 & \textbf{Unité du livre} (25 min) &
Lecture de l'unité (règles, tableaux, pièges) + exercices écrits + vérification des corrigés. \tabularnewline
5 & \textbf{Production orale} (15 min) &
La « mission orale » de l'unité : répéter 10 phrases, se raconter sa journée, rejouer le
dialogue en changeant les détails, s'enregistrer. \tabularnewline
6 & \textbf{Journal écrit} (10 min) &
5 à 10 phrases sur votre journée, en anglais, sans dictionnaire d'abord, puis correction
avec un dictionnaire ou un correcteur. \tabularnewline
\end{tabularx}
\end{tableau}

\begin{astuce}[Vous n'avez que 30 minutes ?]
Ne supprimez pas au hasard. Gardez toujours, dans l'ordre : \textbf{(1)} les 15 min de
vocabulaire SRS, \textbf{(2)} 10 min de parole à voix haute, \textbf{(3)} 5 min d'écoute.
La grammaire peut attendre un jour ; la répétition espacée et la parole, jamais. Une
journée à 30 minutes vaut mieux qu'une journée à zéro.
\end{astuce}

\begin{astuce}[Le temps mort devient du temps d'anglais]
Vous avez chaque jour 1 à 2 heures de temps mort (transport, cuisine, marche, attente,
ménage). Transformez-les en immersion : écouteurs, un podcast ou une playlist. Cela double
votre exposition sans effort supplémentaire. C'est la différence entre ceux qui réussissent
en 12 mois et ceux qui échouent en 3 ans.
\end{astuce}

\section{La semaine type}
\begin{tableau}[Semaine de 6 jours + 1 jour de bilan]
\begin{tabularx}{\linewidth}{@{}p{2.0cm}p{4.0cm}X@{}}
\entete{Jour} & \entete{Accent principal} & \entete{Contenu} \tabularnewline
\hline
\textbf{Lundi} & Grammaire neuve & Séance complète ; nouvelle unité du livre. Fixez l'objectif de la semaine. \tabularnewline
\textbf{Mardi} & Vocabulaire + écoute & Séance complète ; 1\textsuperscript{re} révision de l'unité de lundi à voix haute (5 min). \tabularnewline
\textbf{Mercredi} & \textbf{Parole} & Séance + 30 min de conversation (partenaire, application, ou monologue enregistré sur un sujet du livre). \tabularnewline
\textbf{Jeudi} & Grammaire + lecture & Séance complète ; lecture extensive 15 min (texte facile, plaisir). \tabularnewline
\textbf{Vendredi} & Vocabulaire + écoute & Séance + dictée de 5 phrases + chanson travaillée. \tabularnewline
\textbf{Samedi} & \textbf{Intégration} & Séance allégée + 45 min de « projet » : série, appel, présentation, écriture d'un texte long. \tabularnewline
\textbf{Dimanche} & Bilan (45 min) & Test rapide de la semaine, relecture des corrigés, préparation du matériel de la semaine, enregistrement vocal de 3 min. \tabularnewline
\end{tabularx}
\end{tableau}

\section{Les douze mois en détail}

\subsection{Mois 1 --- Démarrer fort (A1)}
\textbf{Priorités :} les sons (chapitre 5), les 300 mots les plus fréquents, \en{be},
\en{have}, le présent simple, se présenter, les nombres, l'alphabet.

\begin{itemize}
  \item \textbf{Chapitres du livre :} 1 et 2 (la méthode, le plan), 3 et 4 (les huit techniques, les outils), 5 (les 44 sons), 8 (unités 1 à 4 : \en{be}, \en{have}, le présent simple, les articles), 12 (thèmes 1 à 4).
  \item \textbf{Vocabulaire :} 15 mots/jour $\times$ 26 jours = \textbf{390 mots}.
  \item \textbf{Écoute :} 15 min/jour d'un podcast pour débutants (voir annexe E), vitesse 0,75×.
  \item \textbf{Parole :} chaque jour, 10 phrases à voix haute + présentation de soi
  (2 min) une fois par semaine, enregistrée.
\end{itemize}
\textbf{Jalon de fin de mois :} vous vous présentez pendant 2 minutes sans lire ;
vous comprenez et dites 200 phrases de survie ; vous reconnaissez les 44 sons de l'anglais.

\subsection{Mois 2 --- Construire la phrase (A1)}
\textbf{Priorités :} poser des questions, nier, les articles, le pluriel, \en{there is/are},
les prépositions de lieu et de temps, les possessifs, décrire.

\begin{itemize}
  \item \textbf{Chapitres :} 6 (accent, rythme, \en{connected speech}), 8 (unités 5 à 10 : questions, négation, pluriels, \en{there is/are}, prépositions, possessifs), 12 (thèmes 5 à 8), 14 (situations 1 à 8).
  \item \textbf{Vocabulaire :} 600 mots cumulés.
  \item \textbf{Parole :} décrire une pièce, une personne, sa famille, son quartier (3 min chacun).
\end{itemize}
\textbf{Jalon :} vous posez 30 questions types sans hésiter ; vous décrivez une image
pendant 2 minutes.

\subsection{Mois 3 --- Le passé et le futur (A1\ra A2)}
\textbf{Priorités :} \en{past simple} (et les 60 verbes irréguliers essentiels),
\en{past continuous}, \en{going to}, \en{will}, comparatifs, raconter.

\begin{itemize}
  \item \textbf{Chapitres :} 9 (unités 1 à 4 : \en{past simple}, \en{past continuous},
  \en{going to}, \en{will}), annexe A (verbes irréguliers, groupes 1 et 2),
  12 (thèmes 9 à 12), 14 (situations 9 à 14).
  \item \textbf{Vocabulaire :} 900 mots cumulés.
  \item \textbf{Parole :} raconter sa journée d'hier, un souvenir d'enfance, ses projets
  de l'année (3 min chacun).
\end{itemize}
\textbf{Jalon :} vous racontez une histoire de 2 minutes au passé ; vous connaissez
60 verbes irréguliers par c\oe ur.

\subsection{Mois 4 --- Le mur du \en{present perfect} (A2)}
\textbf{Priorités :} \en{present perfect} contre \en{past simple} (le plus gros piège
francophone), \en{for}/\en{since}, \en{just}/\en{already}/\en{yet}, \en{used to},
modaux de base.

\begin{itemize}
  \item \textbf{Chapitres :} 9 (unités 5 à 8 : \en{present perfect}, \en{for}/\en{since}, \en{used to}, modaux de base), 18 (erreurs 1 à 30), 14 (situations 15 à 20).
  \item \textbf{Vocabulaire :} 1\,200 mots cumulés --- santé, ville, transports, courses.
  \item \textbf{Parole :} parler de son expérience (\en{I have never\ldots}, \en{I have been\ldots}),
  situations de voyage.
\end{itemize}
\textbf{Jalon :} dans un test de 20 phrases, vous choisissez correctement \en{past simple}
ou \en{present perfect} dans 18 cas sur 20.

\subsection{Mois 5 --- Consolider et travailler (A2)}
\textbf{Priorités :} révision de tous les temps, voix passive simple, \en{should}/\en{must}/
\en{have to}, vocabulaire du travail et des métiers, premier roman simplifié.

\begin{itemize}
  \item \textbf{Chapitres :} 10 (unités 1 à 3 : passif simple, \en{should}/\en{must}/\en{have to}),
  12 (thèmes 13 à 16), 17 (partie A : la méthode de lecture).
  \item \textbf{Lecture :} un \en{graded reader} niveau A2 (60 à 80 pages), 10 pages/jour.
  \item \textbf{Vocabulaire :} 1\,500 mots cumulés.
\end{itemize}
\textbf{Jalon :} vous lisez un texte A2 inconnu et en comprenez l'essentiel sans dictionnaire ;
vous parlez de votre métier pendant 5 minutes.

\subsection{Mois 6 --- Bilan de mi-parcours (A2\ra B1)}
\textbf{Priorités :} test complet, conditionnels 0/1/2, propositions relatives,
gérondif contre infinitif, première série de \en{phrasal verbs}, première conversation
libre de 20 minutes.

\begin{itemize}
  \item \textbf{Chapitres :} 10 (unités 4 à 5 : conditionnels 0/1/2, propositions relatives,
  gérondif contre infinitif), 13 (les 75 premiers \emph{phrasal verbs}),
  16 (le protocole des 5 passes), 19 (test du mois 6).
  \item \textbf{Parole :} \textbf{deux conversations libres de 30 minutes par semaine} avec un
  partenaire (application, club, camarade). Sujet imposé : « mon pays, ma famille, mes projets ».
  \item \textbf{Vocabulaire :} 1\,800 mots cumulés.
\end{itemize}
\textbf{Jalon :} test de mi-parcours réussi (chapitre 19) ; vous comprenez une série
avec sous-titres \emph{anglais} ; vous tenez 10 minutes de conversation.

\subsection{Mois 7 --- La nuance (B1)}
\textbf{Priorités :} conditionnel 3, discours rapporté, \en{wish}/\en{if only},
quantifieurs, vocabulaire des médias et de l'actualité, débat.

\begin{itemize}
  \item \textbf{Chapitres :} 11 (unités 1 à 3 : conditionnel 3, discours rapporté,
  \en{wish}/\en{if only}), 13 (expressions idiomatiques 1 à 60), 14 (situations 21 à 28).
  \item \textbf{Écoute :} passage aux podcasts « naturels lents » et aux journaux télévisés.
  \item \textbf{Vocabulaire :} 2\,100 mots cumulés.
\end{itemize}
\textbf{Jalon :} vous donnez et défendez une opinion pendant 10 minutes ; vous comprenez
un journal télévisé à 70\,\%.

\subsection{Mois 8 --- Voyager et présenter (B1)}
\textbf{Priorités :} voix passive complète, \en{have something done}, structures emphatiques,
modaux du passé, situations de voyage (aéroport, hôtel, douane, urgence).

\begin{itemize}
  \item \textbf{Chapitres :} 11 (unités 4 à 7 : passif complet, \en{have something done}, structures emphatiques, modaux du passé), 14 (situations 29 à 36), 16 (compréhension des accents).
  \item \textbf{Projet :} préparer et enregistrer une \textbf{présentation de 5 minutes}
  (votre ville, votre métier, un sujet qui vous passionne).
  \item \textbf{Vocabulaire :} 2\,400 mots cumulés.
\end{itemize}
\textbf{Jalon :} vous gérez seul un trajet complet en anglais (aéroport, hôtel, taxi,
réclamation) ; votre présentation de 5 minutes est fluide.

\subsection{Mois 9 --- Les registres (B1\ra B2)}
\textbf{Priorités :} connecteurs, collocations, idiomes, registre formel/familier,
écriture (e-mails, lettres, CV).

\begin{itemize}
  \item \textbf{Chapitres :} 13 (expressions 61 à 120, argot, les 150 faux amis), 17 (parties B à E : orthographe, ponctuation, les 8 modèles, le journal), annexe D (connecteurs et prépositions).
  \item \textbf{Écriture :} 3 e-mails par semaine (formel, informel, réclamation) + un
  texte de 250 mots chaque dimanche.
  \item \textbf{Vocabulaire :} 2\,700 mots cumulés.
\end{itemize}
\textbf{Jalon :} vous écrivez un e-mail professionnel correct du premier coup ; vous
reconnaissez et utilisez 100 expressions idiomatiques.

\subsection{Mois 10 --- L'anglais professionnel (B2)}
\textbf{Priorités :} entretien d'embauche, réunion, téléphone, négociation, présentation
avec supports.

\begin{itemize}
  \item \textbf{Chapitres :} 14 (situations 37 à 40 et blocs professionnels), 17 (e-mail professionnel, CV, lettre de motivation), 18 (erreurs 31 à 100).
  \item \textbf{Projet :} simuler un entretien d'embauche complet (45 min) avec un
  partenaire ou en auto-questionnement enregistré.
  \item \textbf{Vocabulaire :} 3\,000 mots cumulés.
\end{itemize}
\textbf{Jalon :} vous passez un entretien de 30 minutes ; vous animez une réunion simple.

\subsection{Mois 11 --- La fluidité (B2)}
\textbf{Priorités :} débit, \emph{fillers} naturels, reformulation, humour, compréhension
des accents (britannique, américain, nigérian, ghanéen, indien, sud-africain).

\begin{itemize}
  \item \textbf{Chapitres :} 15 (les 7 drills de fluidité, les \emph{fillers}, le plan
  anti-blocage), 16 (les accents du monde), 6 (révision intonation et \en{connected speech}).
  \item \textbf{Parole :} 30 minutes de conversation libre par jour si possible.
\end{itemize}
\textbf{Jalon :} vous parlez 20 minutes sur un sujet inconnu, sans préparation, sans
repasser au français.

\subsection{Mois 12 --- Certification et autonomie (B2\ra C1)}
\textbf{Priorités :} révision générale, test blanc complet, préparation éventuelle d'un
examen (TOEIC, IELTS, Cambridge B2 First), construction de votre plan pour l'année 2.

\begin{itemize}
  \item \textbf{Chapitres :} 19 (tous les tests et la préparation aux certifications), annexe C (les 14 fiches de révision), annexe F (le bilan des 365 jours).
  \item \textbf{Projet final :} trois conversations d'une heure, un film entier sans
  sous-titres, une présentation de 10 minutes.
\end{itemize}
\textbf{Jalon final :} comparez votre enregistrement du jour 0 à celui du jour 365.
C'est là que vous saurez.

\section{Que faire quand vous prenez du retard}
Vous \emph{allez} prendre du retard : maladie, travail, voyage, mariage, Ramadan, examens,
enfants. C'est normal et ce n'est pas grave --- à une condition : ne jamais transformer un
retard en abandon.

\begin{regle}[La règle des deux jours]
Vous pouvez rater \textbf{un} jour. Jamais \textbf{deux} de suite. Le deuxième jour manqué
est celui qui crée l'habitude de manquer. Si vous ratez deux jours, reprenez le troisième
avec une séance de 15 minutes seulement : l'important est de recoller la chaîne.
\end{regle}

\begin{itemize}
  \item \textbf{Retard de 1 à 2 semaines :} ne « rattrapez » pas en doublant les séances.
  Reprenez le fil, et décalez simplement la fin du plan de deux semaines.
  \item \textbf{Retard d'un mois ou plus :} refaites le test de positionnement, et reprenez
  au mois correspondant, en gardant le vocabulaire SRS à jour (c'est lui qui vous a permis
  de ne rien perdre).
  \item \textbf{Découragement :} relisez votre enregistrement du mois précédent. Les progrès
  sont invisibles au jour le jour, évidents à trois mois.
\end{itemize}

\begin{memo}[Le contrat avec vous-même]
Écrivez et signez, aujourd'hui :\\[4pt]
\textit{Je, \trou[5cm], m'engage à étudier l'anglais \trou[1.5cm] minutes par jour,
\trou[1cm] jours par semaine, du \trou[2.5cm] au \trou[2.5cm], et à parler à voix haute
chaque jour. Mon objectif au \trou[2.5cm] (dans 12 mois) : tenir une conversation
de 30 minutes en anglais sans préparation.}\\[4pt]
Signature : \trou[5cm] \hfill Date : \trou[3cm]
\end{memo}

% =====================================================================
%  CHAPITRES 3 ET 4 — Les techniques d'entraînement + les outils
% =====================================================================
\chapter{Les huit techniques qui font parler}

Une méthode de langue est une boîte à outils. Voici les huit outils qui produisent
l'essentiel des résultats. Chacun a son protocole précis --- appliqué sérieusement, pas
approximativement. Maîtrisez-les tous dans le premier mois : ensuite, votre année n'est plus
qu'une question de contenu.

\section{Technique 1 --- Le \emph{shadowing} (l'ombre)}
\textbf{Ce que c'est :} écouter un audio et le répéter \emph{presque en même temps} que le
locuteur, comme une ombre qui suit, avec un décalage d'une demi-seconde. C'est l'exercice le
plus puissant pour la prononciation, le rythme, l'intonation et la vitesse de parole.

\begin{tableau}[Protocole du shadowing --- 7 étapes, 15 minutes]
\begin{tabularx}{\linewidth}{@{}p{1.2cm}p{4.2cm}X@{}}
\entete{Étape} & \entete{Action} & \entete{Détail} \tabularnewline
\hline
1 & Choisir un audio court & 30 à 90 secondes, niveau compréhensible à 90\,\%, voix que vous aimez. Toujours le même pendant 3 jours. \tabularnewline
2 & Écouter sans texte & 2 fois. Comprendre le sens global. \tabularnewline
3 & Lire le texte & 1 fois, à voix haute lentement, en cherchant les mots inconnus. \tabularnewline
4 & Écouter en lisant & 2 fois, en suivant du doigt. Repérer les liaisons et les mots « mangés ». \tabularnewline
5 & \textbf{Shadower avec le texte} & 3 fois : vous parlez en même temps que l'audio, les yeux sur le texte. Imiter \emph{tout} : mélodie, pauses, énergie. \tabularnewline
6 & \textbf{Shadower sans le texte} & 3 fois. C'est là que la mémoire et l'oreille travaillent. \tabularnewline
7 & S'enregistrer et comparer & 1 fois. Écoutez votre version puis l'original. Notez \textbf{un} défaut à corriger demain. \tabularnewline
\end{tabularx}
\end{tableau}

\begin{astuce}[Variantes utiles]
\textbf{Vitesse réduite} (0,75×) les deux premières semaines, puis vitesse normale, puis
1,1× pour l'agilité. \textbf{Shadowing en marchant} : marcher en shadowant active le rythme
corporel et améliore nettement la prosodie. \textbf{Shadowing chuchoté} dans le bus : vous
travaillez sans gêner personne.
\end{astuce}

\begin{piege}[Ne shadowez pas votre propre accent]
Le but n'est pas de « dire le texte », mais de \textbf{copier le son}. Si vous lisez avec
votre accent français habituel, l'exercice ne sert presque à rien. Exagérez l'imitation,
comme un acteur : c'est la seule façon de casser vos automatismes.
\end{piege}

\section{Technique 2 --- La répétition espacée (Anki)}
\textbf{Ce que c'est :} un logiciel qui vous montre chaque carte au moment exact où vous
alliez l'oublier. C'est le multiplicateur de vocabulaire.

\subsection{Régler Anki en 5 minutes}
\begin{enumerate}
  \item Créez un paquet nommé \en{Anglais --- 12 mois}.
  \item Type de carte : \textbf{Phrase à trou} (\en{Cloze}) de préférence --- bien plus
  efficace que « mot isolé ».
  \item Cartes nouvelles par jour : \textbf{15} (au début), jusqu'à 25 à partir du mois 3.
  \item Intervalle maximal : 180 jours. Pas de limite de révisions par jour.
  \item Activez l'audio automatique si disponible (ou ajoutez le son vous-même depuis Forvo).
\end{enumerate}

\subsection{Comment fabriquer une bonne carte}
\begin{tableau}[Mauvaise carte \ra bonne carte]
\begin{tabularx}{\linewidth}{@{}p{7.6cm}X@{}}
\entete{\ko À éviter} & \entete{\ok À faire} \tabularnewline
\hline
Recto : \en{to achieve} \quad Verso : \en{réussir, accomplir} &
Recto : \en{She finally \{\{c1::achieved\}\} her goal.} \\
& Verso : \en{achieve} = atteindre, accomplir --- \pron{a-tchi-v} \\
& + la phrase entière lue à voix haute. \tabularnewline
\hline
Un mot seul, hors contexte & Toujours \textbf{une phrase complète}, personnelle si possible (votre vie, votre travail). \tabularnewline
\hline
Une carte = 5 sens possibles & \textbf{Une seule réponse attendue} par carte. \tabularnewline
\hline
Traduire FR \ra EN uniquement & Les deux sens : EN \ra FR (compréhension) et FR \ra EN (production). \tabularnewline
\end{tabularx}
\end{tableau}

\begin{memo}[Les 3 règles d'or d'Anki]
\textbf{1.} Faites vos révisions \textbf{tous les jours}, à la même heure (15 min suffisent).
\textbf{2.} Dites chaque carte \textbf{à voix haute}, en entier, avec l'intonation.
\textbf{3.} Ne créez que des cartes pour des mots que vous avez \textbf{réellement rencontrés}
(dans ce livre, un podcast, une série) --- jamais des listes apprises « au cas où ».
\end{memo}

\section{Technique 3 --- L'auto-conversation (\emph{self-talk})}
\textbf{Ce que c'est :} parler seul en anglais, à voix haute ou dans sa tête, en commentant
ce que vous faites. C'est gratuit, disponible 24 h/24, et cela entraîne exactement le muscle
qui vous manque : \textbf{produire sans interlocuteur pour vous aider}.

\begin{tableau}[Les 5 niveaux du self-talk --- progressez dans l'ordre]
\begin{tabularx}{\linewidth}{@{}p{1.4cm}p{3.6cm}X@{}}
\entete{Niveau} & \entete{Exercice} & \entete{Exemple} \tabularnewline
\hline
\textbf{1} & Nommer & Vous regardez autour de vous et nommez tout ce que vous voyez, à voix haute, 2 minutes. \en{a table, a phone, the window, my shoes\ldots} \tabularnewline
\textbf{2} & Commenter ses actes & \en{I'm opening the fridge. I'm looking for some water. It's cold.} (en cuisinant, en vous habillant, en marchant) \tabularnewline
\textbf{3} & Raconter sa journée & Chaque soir, 3 minutes : \en{This morning I woke up at six. Then I\ldots} \tabularnewline
\textbf{4} & Anticiper & Le matin : \en{Today I'm going to\ldots\ First I will\ldots\ Then\ldots} \tabularnewline
\textbf{5} & Argumenter / débatteur & Choisissez un sujet (\en{Is social media good for young people?}) et parlez 3 minutes \emph{pour}, puis 3 minutes \emph{contre}. \tabularnewline
\end{tabularx}
\end{tableau}

\begin{astuce}[Le « trou de mot » est votre meilleur professeur]
Quand un mot manque en self-talk, \textbf{ne cherchez pas tout de suite}. Contournez-le :
\en{``It's the thing you use to open a bottle\ldots''}. Puis notez le mot manquant sur votre
téléphone et créez une carte Anki le soir. Vous ne l'oublierez plus jamais, parce que votre
cerveau l'a \emph{cherché}.
\end{astuce}

\section{Technique 4 --- Le journal en anglais}
Dix minutes par jour, cinq à dix phrases. Aucun dictionnaire au premier jet --- écrivez avec
ce que vous avez, puis corrigez.

\begin{tableau}[Prompts de journal --- 30 sujets pour le premier mois]
\footnotesize
\begin{multicols}{2}
\begin{enumerate}
\item Ce que j'ai fait hier.
\item Ma famille.
\item Mon quartier / ma ville.
\item Mon travail ou mes études.
\item Ce que j'ai mangé aujourd'hui.
\item Mes projets pour cette année.
\item La personne que j'admire le plus.
\item Mon meilleur souvenir d'enfance.
\item Ce qui m'a rendu heureux cette semaine.
\item Ce qui m'a énervé cette semaine.
\item Ma journée idéale.
\item Le temps qu'il fait et ce que je fais avec.
\item Un film ou une série que j'ai vu.
\item Ce que je veux dans cinq ans.
\item Mon téléphone : à quoi il me sert.
\item Une habitude que je veux changer.
\item Un voyage que j'ai fait.
\item La description de ma chambre.
\item Ce que j'ai appris cette semaine en anglais.
\item Un problème que j'ai résolu.
\item Mon plat préféré, comment on le cuisine.
\item Une conversation que j'ai eue aujourd'hui.
\item Ce dont j'ai peur.
\item Mes trois objectifs du mois.
\item Une personne que je veux remercier.
\item Ce que je ferais avec un million.
\item Ma routine du matin, minute par minute.
\item Une chose nouvelle que j'ai essayée.
\item Ce qui est difficile en anglais pour moi.
\item Bilan de mon mois d'anglais.
\end{enumerate}
\end{multicols}
\end{tableau}

\section{Technique 5 --- La dictée (\en{dictation})}
Écrire ce que l'on entend : c'est le seul exercice qui prouve si vous entendez \emph{vraiment}.
Protocole (10 minutes) :
\begin{enumerate}
  \item Écoutez une phrase (5 à 10 mots) une fois, sans écrire.
  \item Écoutez-la à nouveau et écrivez ce que vous entendez, en mettant un tiret pour chaque
  trou : \en{I --- to the --- yesterday}.
  \item Réécoutez jusqu'à 5 fois en complétant.
  \item Comparez avec le texte exact. Comptez les trous.
  \item \textbf{Analysez chaque erreur} : était-ce un mot inconnu ? un mot connu mais « mangé »
  par la liaison (\en{going to} entendu comme \en{gonna}) ? un son non perçu (\en{th}, \en{h} final) ?
  \item Relisez la phrase à voix haute trois fois, en imitant l'audio.
\end{enumerate}
Cinq phrases par jour, cinq jours par semaine : au bout de trois mois, votre oreille n'a plus
rien à voir.

\section{Technique 6 --- La lecture à voix haute}
Prenez n'importe quel texte de ce livre (dialogue, exemple, paragraphe) et lisez-le à voix
haute, lentement, en articulant, 5 minutes par jour. Bénéfices : votre bouche s'habitue aux
enchaînements anglais ; vous mémorisez des structures par la voie motrice ; vous repérez
immédiatement les mots que vous « savez » mais ne savez pas dire.

\begin{astuce}[La règle des trois lectures]
\textbf{1\textsuperscript{re} lecture :} lente, mot à mot, en corrigeant la prononciation.
\textbf{2\textsuperscript{e} lecture :} normale, en respectant les pauses et l'accent tonique.
\textbf{3\textsuperscript{e} lecture :} rapide et expressive, comme un acteur --- en levant
les yeux du texte à la fin de chaque phrase.
\end{astuce}

\section{Technique 7 --- L'immersion fabriquée}
Vous ne vivez pas à Londres ? Fabriquez Londres chez vous. Voici la liste à mettre en place
\textbf{dans la première semaine}, une fois pour toutes :

\begin{tableau}[Votre environnement d'immersion --- à installer une fois]
\begin{tabularx}{\linewidth}{@{}p{4.4cm}X@{}}
\entete{Élément} & \entete{Action concrète} \tabularnewline
\hline
\textbf{Téléphone et ordinateur} & Langue du système en anglais. Vous apprendrez 200 mots utiles en trois jours sans effort. \tabularnewline
\textbf{Réseaux sociaux} & Abonnez-vous à 10 comptes anglophones dans des domaines que vous aimez déjà (football, cuisine, religion, technologie, humour, actualité). L'algorithme vous servira de l'anglais toute la journée. \tabularnewline
\textbf{Musique} & Playlist anglaise pour la cuisine, le transport, le sport. Une chanson travaillée par semaine (paroles + sens + chant). \tabularnewline
\textbf{Séries et films} & Sous-titres \textbf{anglais} (jamais français) à partir du mois 3. Avant : sous-titres français, mais 10 minutes par épisode en sous-titres anglais. \tabularnewline
\textbf{Radio / podcast} & Un flux anglais en fond sonore pendant les tâches ménagères, même sans comprendre. Le cerveau s'habitue au rythme. \tabularnewline
\textbf{YouTube} & Historique de recommandations à « nettoyer » : regardez 20 vidéos en anglais, l'algorithme basculera. \tabularnewline
\textbf{Étiquettes} & Collez des post-it en anglais sur 20 objets de la maison (\en{fridge, mirror, cupboard, kettle, shelf\ldots}). \tabularnewline
\textbf{Pensée} & Deux « alarmes » par jour (11 h et 17 h) : au bip, pensez en anglais pendant 2 minutes, quoi qu'il arrive. \tabularnewline
\textbf{Personnes} & Un partenaire de langue identifié (application, groupe, club, collègue, camarade de classe) --- rendez-vous hebdomadaire fixé dans l'agenda. \tabularnewline
\end{tabularx}
\end{tableau}

\section{Technique 8 --- L'enregistrement de soi}
Une fois par semaine, enregistrez 3 minutes de parole libre sur un sujet donné. Écoutez
le lendemain --- pas immédiatement (vous seriez trop sévère) --- avec une grille de trois
questions :
\begin{enumerate}
  \item \textbf{Fluidité :} combien d'hésitations de plus de 3 secondes ? combien de retours
  au français ?
  \item \textbf{Prononciation :} quels sons ont dérapé (\en{th}, \en{h}, voyelles, accent tonique) ?
  \item \textbf{Grammaire :} quelles erreurs reviennent ? (Notez-les : trois maximum, pas plus.
  Travaillez-les la semaine suivante.)
\end{enumerate}

\begin{memo}[Le cycle hebdomadaire complet]
\textbf{Lundi :} nouvelle unité. \textbf{Mardi :} révision + shadowing.
\textbf{Mercredi :} conversation avec partenaire. \textbf{Jeudi :} lecture + grammaire.
\textbf{Vendredi :} dictée + chanson. \textbf{Samedi :} projet (série, présentation).
\textbf{Dimanche :} enregistrement de 3 min + test rapide + bilan écrit.
\end{memo}


% =====================================================================
\chapter{Outils, ressources et partenaires}

Tout ce dont vous avez besoin est \textbf{gratuit} ou presque. Voici l'équipement minimal,
puis les ressources classées par niveau. (Les liens détaillés sont en annexe E.)

\section{L'équipement minimal (à réunir en une soirée)}
\begin{enumerate}
  \item \textbf{Un téléphone} avec un enregistreur vocal et des écouteurs. C'est l'outil
  central : il remplace un laboratoire de langue.
  \item \textbf{Anki} (ordinateur ou téléphone) --- répétition espacée. Gratuit sur Android
  (\en{AnkiDroid}) et ordinateur ; payant sur iPhone (\en{AnkiMobile}), où \en{AnkiWeb}
  gratuit suffit.
  \item \textbf{Un dictionnaire anglais-anglais} : Cambridge Dictionary ou Oxford Learner's
  Dictionaries. Passer au dictionnaire unilingue dès le mois 3 est un accélérateur majeur.
  \item \textbf{Un dictionnaire bilingue de secours} : WordReference (avec ses forums, très
  utiles pour les nuances) ou Linguee/DeepL (pour voir des traductions en contexte réel).
  \item \textbf{Un correcteur} : LanguageTool ou le correcteur de Google Docs, pour vos
  journaux écrits. Attention : lisez la correction, ne la copiez pas bêtement.
  \item \textbf{Ce livre}, imprimé ou sur tablette, plus un cahier de 100 pages.
\end{enumerate}

\section{Ressources d'écoute classées par niveau}
\begin{tableau}[Où écouter, selon votre mois]
\begin{tabularx}{\linewidth}{@{}p{2.0cm}p{4.6cm}X@{}}
\entete{Niveau} & \entete{Type de ressource} & \entete{Exemples et usage} \tabularnewline
\hline
\textbf{A1} (mois 1-2) & Podcasts pour apprenants, débit très lent &
Chaînes et podcasts de type « English for beginners », \en{BBC Learning English}
(rubriques débutants), \en{VOA Learning English} (débit lent, vocabulaire contrôlé).
15 min/jour, vitesse 0,75×. \tabularnewline
\textbf{A2} (mois 3-5) & Histoires courtes, dialogues lents, chansons claires &
Podcasts d'histoires pour apprenants (\en{storytelling} lent), \en{6 Minute English} (BBC),
dessins animés pour enfants (vocabulaire simple, articulation exagérée --- ne les méprisez
pas : c'est un excellent matériau). \tabularnewline
\textbf{B1} (mois 6-8) & Podcasts « naturels mais articulés », TED-Ed, séries &
Conférences TED courtes sous-titrées, documentaires, sitcoms (\en{Friends},
\en{Modern Family}), chaînes YouTube de vulgarisation. \tabularnewline
\textbf{B2} (mois 9-12) & Contenu natif : actualités, interviews, humour &
Journaux télévisés (BBC, CNN, Al Jazeera English), podcasts natifs sur vos centres d'intérêt,
débats, stand-up, films sans sous-titres. \tabularnewline
\textbf{Accents} (mois 11) & Variétés d'anglais &
Cherchez volontairement un accent britannique, un américain, un nigérian, un ghanéen, un
kényan, un sud-africain, un indien. Votre oreille doit être mondiale, pas seulement
hollywoodienne. \tabularnewline
\end{tabularx}
\end{tableau}

\section{Trouver des gens à qui parler (le point le plus important du chapitre)}
\begin{tableau}[Huit façons de parler à un humain, de la plus facile à la plus efficace]
\begin{tabularx}{\linewidth}{@{}p{4.6cm}X@{}}
\entete{Solution} & \entete{Comment faire} \tabularnewline
\hline
\textbf{Applications d'échange linguistique} &
HelloTalk, Tandem : vous aidez quelqu'un en français/wolof, il vous aide en anglais.
Messages vocaux d'abord, appels ensuite. 20 min/jour possible. \tabularnewline
\textbf{Serviteurs Discord / groupes WhatsApp} &
De nombreux serveurs d'apprentissage de l'anglais ont des salons vocaux ouverts 24 h/24.
Cherchez « English learning voice chat ». Gratuit, immédiat, anonyme --- donc sans honte. \tabularnewline
\textbf{Clubs de conversation locaux} &
Instituts culturels, centres américains/britanniques, bibliothèques, universités, cafés
langues. Une séance par semaine, souvent gratuite. \tabularnewline
\textbf{Communautés religieuses} &
Dans beaucoup de villes d'Afrique de l'Ouest, des églises, des groupes d'étude ou des
associations tiennent des cultes ou séances en anglais. Renseignez-vous : c'est une source
régulière et bienveillante de pratique. \tabularnewline
\textbf{Un camarade de même niveau} &
La solution la plus sous-estimée : trouvez une personne de votre entourage qui veut aussi
apprendre, et imposez-vous \textbf{une heure par semaine en anglais uniquement}, même
imparfait. Vous progresserez plus vite qu'avec un professeur que vous voyez une fois par mois. \tabularnewline
\textbf{Professeur particulier en ligne} &
Plateformes de tutorat (payantes) : 30 min, une à deux fois par semaine. Utile surtout pour
corriger la prononciation et débloquer l'oral à partir du mois 4. \tabularnewline
\textbf{Parler à un assistant vocal} &
L'assistant de votre téléphone, réglé en anglais, comprend et répond. Excellent pour
s'entraîner sans jugement : \en{``What's the weather today?'', ``Set a timer for ten
minutes'', ``Tell me a joke.''} \tabularnewline
\textbf{Parler seul (self-talk)} &
Toujours disponible, jamais en panne, jamais moqueur. C'est 80\,\% du volume de parole
de la plupart des autodidactes qui réussissent. \tabularnewline
\end{tabularx}
\end{tableau}

\begin{piege}[La honte est votre seul véritable ennemi]
Les francophones ne manquent ni d'intelligence ni de ressources : ils manquent
d'\textbf{audace}. Vous ferez des fautes. Vous aurez l'accent. Vous chercherez vos mots.
Les anglophones s'en moquent --- ils sont même admiratifs, parce qu'eux ne parlent qu'une
langue. Répétez-vous cette phrase chaque fois que vous hésitez à ouvrir la bouche :
\en{``Done is better than perfect.''} \textit{(Fait vaut mieux que parfait.)}
\end{piege}

\section{Trois routines à copier}

\subsection{Routine « trajet » (20 min, deux fois par jour)}
\begin{enumerate}
  \item 5 min : révisions Anki (téléphone en main).
  \item 10 min : shadowing d'un épisode de podcast déjà écouté.
  \item 5 min : self-talk --- raconter à voix basse ce que vous allez faire aujourd'hui.
\end{enumerate}

\subsection{Routine « soir calme » (45 min)}
\begin{enumerate}
  \item 15 min : unité de grammaire + exercices.
  \item 15 min : un épisode de série, 3 passes (voir chapitre 16).
  \item 10 min : journal écrit.
  \item 5 min : relecture à voix haute de 10 phrases de la journée.
\end{enumerate}

\subsection{Routine « dimanche bilan » (60 min)}
\begin{enumerate}
  \item 10 min : test rapide de la semaine (questions du livre).
  \item 15 min : enregistrement de 3 min + écoute + 3 défauts notés.
  \item 15 min : nettoyage des cartes Anki (supprimer les cartes inutiles, corriger les
  mauvaises).
  \item 10 min : choix du matériel de la semaine prochaine.
  \item 10 min : conversation avec un partenaire (même 10 minutes comptent).
\end{enumerate}

\begin{mission}[Semaine 1 --- installer la machine]
\begin{enumerate}
  \item Jour 1 : téléphone en anglais + Anki installé + 20 post-it collés.
  \item Jour 2 : premier shadowing (15 min) sur un audio de 1 minute.
  \item Jour 3 : premier journal (5 phrases) + première dictée (3 phrases).
  \item Jour 4 : auto-conversation niveau 1 et 2 (5 min).
  \item Jour 5 : trouver et contacter un partenaire de langue (un seul message suffit).
  \item Jour 6 : enregistrement de 2 minutes sur « pourquoi j'apprends l'anglais ».
  \item Jour 7 : bilan + cases cochées + récompense (vous en aurez besoin toute l'année :
  prévoyez une petite récompense chaque dimanche).
\end{enumerate}
\end{mission}


% ===== PARTIE II — LES SONS DE L'ANGLAIS ===========================
\part{Les sons de l'anglais : comprendre et être compris}
% =====================================================================
%  CHAPITRES 5, 6 ET 7 — Les sons, le rythme, la clinique francophone
% =====================================================================
\chapter{Les 44 sons de l'anglais}

\section{Pourquoi commencer par les sons ?}
Trois raisons, toutes décisives :

\begin{enumerate}
  \item \textbf{L'anglais n'est pas une langue phonétique.} En français, une lettre = un son,
  ou presque. En anglais, \en{ough} se prononce de six façons différentes :
  \en{through} \pron{throu}, \en{though} \pron{dzo}, \en{thought} \pron{thoot},
  \en{tough} \pron{teuf}, \en{cough} \pron{kof}, \en{bough} \pron{baou}. Si vous lisez
  « à la française », vous prononcez faux --- et vous ne reconnaissez pas le mot quand
  quelqu'un d'autre le dit.
  \item \textbf{L'anglais possède 9 sons qui n'existent pas en français.} Tant que votre
  oreille ne les distingue pas, vous entendez \en{ship} et \en{sheep} comme le même mot,
  et vous produisez \en{think} comme \en{sink}. Ce n'est pas un problème de gorge, c'est un
  problème d'\textbf{oreille} : on ne peut pas produire un son qu'on n'entend pas.
  \item \textbf{C'est le premier obstacle ressenti à l'oral.} Beaucoup de francophones
  connaissent les mots mais n'osent pas parler, parce qu'ils « sonnent faux ». Deux semaines
  de travail phonétique suppriment 80\,\% de cette gêne.
\end{enumerate}

\begin{astuce}[Comment lire les tableaux qui suivent]
Chaque son est donné sous trois formes :
\textbf{(1)} le symbole phonétique international (API) entre barres obliques ;
\textbf{(2)} un \textbf{mot anglais modèle} --- apprenez-le, il vaut mille explications ;
\textbf{(3)} une \textbf{prononciation approchée en lettres françaises}, en violet, entre
crochets. La respelling française est une béquille : elle vous dépanne les premiers mois,
mais elle est forcément imprécise. Dès que vous pouvez, remplacez-la par l'écoute réelle du
mot (Forvo, dictionnaire Cambridge avec haut-parleur, YouGlish).
\end{astuce}

\section{Les 12 voyelles simples}
C'est ici que les francophones perdent le plus. Le français a 12 à 16 voyelles selon les
analyses, l'anglais en a 12 --- mais \textbf{différentes}, et surtout l'anglais oppose des
voyelles \textbf{longues} et \textbf{brèves} qui changent le sens des mots.

\begin{tableau}[Les 12 voyelles simples (monohtongues)]
\footnotesize
\begin{longtable}{@{}p{1.7cm}p{3.6cm}p{2.6cm}p{7.2cm}@{}}
\entete{API} & \entete{Mot modèle (et autres exemples)} & \entete{À la française} & \entete{Comment la produire / piège francophone} \tabularnewline
\endfirsthead
\entete{API} & \entete{Mot modèle} & \entete{À la française} & \entete{Comment la produire / piège francophone} \tabularnewline
\endhead
\hline
/\ipaee\ipalong/ & \en{sheep}, \en{see}, \en{tea}, \en{free}, \en{easy} & \pron{chiiip} &
Long « i » tendu, lèvres étirées, son tenu. \textbf{Tenez-le} : \en{sheeeep}. Le français ne
tient pas ses voyelles --- c'est la première chose à changer. \tabularnewline
\hline
/\ipaiot/ & \en{ship}, \en{sit}, \en{big}, \en{English}, \en{minute} & \pron{chip} &
Bref, relâché, entre le « i » et le « é ». Bouche presque fermée mais muscles mous.
\textbf{La paire \en{sheep}/\en{ship} est le test classique.} \tabularnewline
\hline
/\ipaE/ & \en{bed}, \en{ten}, \en{said}, \en{many}, \en{friend} & \pron{bèd} &
« è » bref. Attention : \en{many}, \en{any}, \en{said}, \en{friend} ne se prononcent pas
comme ils s'écrivent. \tabularnewline
\hline
/\ipae/ & \en{cat}, \en{bad}, \en{man}, \en{apple}, \en{happy} & \pron{katt} &
\textbf{N'existe pas en français.} Ouvrez grand la bouche, mâchoire descendue, comme chez le
dentiste, et dites un « a » en tirant vers le « è ». Son très ouvert et nasal. \tabularnewline
\hline
/\ipaa\ipalong/ & \en{father}, \en{car}, \en{start}, \en{class} & \pron{faa-dheu} &
« a » long, arrière de la bouche, lèvres neutres. \en{AmE} \en{class} = \pron{klass},
\en{BrE} = \pron{klaas}. \tabularnewline
\hline
/\ipao/ (bref) & \en{hot}, \en{dog}, \en{stop}, \en{box} & \pron{hot} &
« o » bref et ouvert, lèvres arrondies. Américain : souvent plus proche de « aa »
(\en{hot} \pron{haat}). \tabularnewline
\hline
/\ipao\ipalong/ & \en{four}, \en{thought}, \en{law}, \en{door}, \en{more} & \pron{foo} &
« ô » long, lèvres bien arrondies et projetées. \tabularnewline
\hline
/\ipaups/ & \en{book}, \en{put}, \en{good}, \en{could}, \en{woman} & \pron{bouk} &
« ou » très bref et relâché. Ne le confondez jamais avec le « ou » long ci-dessous :
\en{full} \pron{foul} $\neq$ \en{fool} \pron{fouul}. \tabularnewline
\hline
/\ipaoo\ipalong/ & \en{too}, \en{food}, \en{blue}, \en{group}, \en{you} & \pron{touu} &
« ou » long, lèvres très arrondies, son tenu. \tabularnewline
\hline
/\ipalv/ & \en{cup}, \en{love}, \en{money}, \en{country}, \en{sun} & \pron{keup} &
\textbf{N'existe pas en français.} Bref, central, entre « a » et « eu », mâchoire relâchée.
C'est le \en{u} de \en{but}. Les francophones disent « keup » à la française : trop long et
trop rond. Soyez bref et mou. \tabularnewline
\hline
/\ipaer\ipalong/ & \en{bird}, \en{her}, \en{work}, \en{learn}, \en{early} & \pron{beuud} &
« eu » long + \textbf{r} anglais prononcé (contrairement au français où le « euh » final
disparaît). \en{bird}, \en{word}, \en{world}, \en{girl}, \en{heard} : cinq orthographes,
un seul son. \tabularnewline
\hline
/\ipeschwa/ & \en{about}, \en{teacher}, \en{banana}, \en{problem} & \pron{e} très faible &
\textbf{Le son le plus fréquent de l'anglais (environ 30\,\% des voyelles).} Voyelle neutre,
à peine prononcée, dans toutes les syllabes \emph{non accentuées}. \en{about} = \pron{e-baout},
\en{banana} = \pron{be-na-neu}, \en{teacher} = \pron{tii-tcheu}. \tabularnewline
\end{longtable}
\end{tableau}

\begin{piege}[Le français n'a pas de voyelles longues --- l'anglais si]
\en{sheep} et \en{ship}, \en{fool} et \en{full}, \en{cart} et \en{cut}, \en{beach} et
\en{bitch} : quatre paires où la longueur change le sens. Un francophone qui abrège toutes
ses voyelles produit des contresens --- parfois embarrassants. Entraînement : tenez la
voyelle longue une seconde entière, puis dites la brève en un éclair.
\end{piege}

\section{Les 8 diphtongues}
Une diphtongue est un glissement d'une voyelle vers une autre, \emph{dans la même syllabe}.
Le français en a peu (\en{oui}, \en{bien}) ; l'anglais en est plein. Le secret : le premier
élément est fort et long, le second est faible et court.

\begin{tableau}[Les 8 diphtongues]
\footnotesize
\begin{longtable}{@{}p{1.9cm}p{4.3cm}p{2.6cm}p{6.3cm}@{}}
\entete{API} & \entete{Mots modèles} & \entete{À la française} & \entete{Remarque} \tabularnewline
\endfirsthead
\entete{API} & \entete{Mots modèles} & \entete{À la française} & \entete{Remarque} \tabularnewline
\endhead
\hline
/\ipaei/ & \en{say}, \en{day}, \en{eight}, \en{name}, \en{they} & \pron{dé-i} &
Le « ai » français est \emph{une} voyelle ; l'anglais glisse. \en{day} = \pron{dé-i},
jamais \pron{dè}. \tabularnewline
\hline
/\ipaai/ & \en{my}, \en{time}, \en{five}, \en{eye}, \en{idea} & \pron{ma-i} &
Glissement net « a \ra i ». \en{time} = \pron{ta-im}, pas \pron{taïm} mou. \tabularnewline
\hline
/\ipaOi/ & \en{boy}, \en{oil}, \en{voice}, \en{enjoy} & \pron{bo-i} &
Identique au français « oi », mais plus appuyé. \tabularnewline
\hline
/\ipaau/ & \en{now}, \en{house}, \en{out}, \en{about} & \pron{na-ou} &
Commence bouche grande ouverte. \en{now} = \pron{na-ou}, pas \pron{no}. \tabularnewline
\hline
/\ipaou/ \en{(BrE)} / \ipaoou \en{(AmE)} & \en{go}, \en{home}, \en{no}, \en{show} & \pron{gou-ou} &
\textbf{Erreur n°1 des francophones :} dire \en{no} = \pron{no}. C'est \pron{nou-ou}.
Le « o » anglais est toujours une diphtongue. \tabularnewline
\hline
/\ipaiot\ipeschwa/ & \en{here}, \en{near}, \en{idea}, \en{ear} & \pron{hi-eu} &
\en{here} = \pron{hi-eu}, pas \pron{hèr}. \tabularnewline
\hline
/\ipaE\ipeschwa/ & \en{hair}, \en{care}, \en{where}, \en{there} & \pron{hè-eu} &
\en{hair} = \pron{hè-eu}. Attention : \en{here} et \en{hair} sont différents. \tabularnewline
\hline
/\ipaups\ipeschwa/ & \en{tour}, \en{pure}, \en{sure}, \en{cure} & \pron{tou-eu} &
De plus en plus prononcé \pron{too} (\en{sure} = \pron{choo}). \tabularnewline
\end{longtable}
\end{tableau}

\section{Les 24 consonnes}
Bonne nouvelle : 19 consonnes anglaises existent en français et ne posent aucun problème
(\en{p, b, t, d, k, g, f, v, s, z, m, n, l} --- avec quelques nuances). Cinq demandent un
vrai travail, signalées par \etoile.

\begin{tableau}[Les 24 consonnes]
\footnotesize
\begin{longtable}{@{}p{1.8cm}p{4.0cm}p{2.4cm}p{6.9cm}@{}}
\entete{API} & \entete{Mots modèles} & \entete{À la française} & \entete{Remarque pour francophone} \tabularnewline
\endfirsthead
\entete{API} & \entete{Mots modèles} & \entete{À la française} & \entete{Remarque pour francophone} \tabularnewline
\endhead
\hline
/p/ & \en{pen}, \en{top}, \en{apple} & \pron{pèn} &
Comme en français, mais \textbf{aspiré} en début de mot : un petit souffle
(\en{pen} = p\,+\,souffle). Devant \en{s}, pas d'aspiration : \en{spin} = \pron{spine}. \tabularnewline
\hline
/b/ & \en{bad}, \en{job}, \en{about} & \pron{bad} & Identique au français. \tabularnewline
\hline
/t/ & \en{tea}, \en{stop}, \en{water} & \pron{tii} &
Aspiré en début de mot. Américain : entre deux voyelles, il devient un « d » bref
(\en{water} = \pron{ouo-deu}, \en{city} = \pron{si-di}). \tabularnewline
\hline
/d/ & \en{day}, \en{did}, \en{red} & \pron{dé-i} & Identique. \tabularnewline
\hline
/k/ & \en{cat}, \en{school}, \en{back} & \pron{katt} &
Aspiré en début de mot. \en{school} = \pron{skoul} (pas de souffle après \en{s}). \tabularnewline
\hline
/g/ & \en{go}, \en{big}, \en{again} & \pron{gou} &
Toujours dur, jamais « j » : \en{get} = \pron{guèt}, \en{girl} = \pron{gueul}. \tabularnewline
\hline
/f/ & \en{fat}, \en{life}, \en{phone} & \pron{fatt} &
\en{ph} = \en{f} : \en{phone}, \en{photo}, \en{phrase}. \tabularnewline
\hline
/v/ & \en{van}, \en{love}, \en{very} & \pron{van} &
Ne confondez jamais \en{v} et \en{w} (voir plus bas). \tabularnewline
\hline
\etoile /\ipatheta/ & \en{think}, \en{three}, \en{thing}, \en{math} & \pron{thi-nk} &
\textbf{N'existe pas en français.} Langue \emph{entre les dents}, souffle. Jamais « s »,
jamais « f ». \en{think} $\neq$ \en{sink}. \tabularnewline
\hline
\etoile /\ipadh/ & \en{this}, \en{that}, \en{mother}, \en{the} & \pron{dhis} &
Même position de la langue, mais \textbf{avec la voix} (cordes vocales qui vibrent).
\en{this} $\neq$ \en{dis}. Ce son apparaît dans les mots les plus fréquents de l'anglais
(\en{the, this, that, they, there, them, then, than, with, mother, father, other}). \tabularnewline
\hline
/s/ & \en{see}, \en{city}, \en{miss} & \pron{sii} &
\en{c} devant \en{e/i} = \en{s} : \en{city}, \en{nice}. \tabularnewline
\hline
/z/ & \en{zoo}, \en{is}, \en{music}, \en{was} & \pron{zou} &
\en{s} entre voyelles = \en{z} : \en{rose}, \en{busy}, \en{music}. \tabularnewline
\hline
/\ipash/ & \en{she}, \en{nation}, \en{sure}, \en{special} & \pron{chii} &
Comme le « ch » français. \en{ti}, \en{ci}, \en{ss} peuvent donner ce son :
\en{nation} \pron{né-i-cheun}, \en{special} \pron{spé-chl}, \en{pressure} \pron{pré-cheu}. \tabularnewline
\hline
/\ipazh/ & \en{measure}, \en{vision}, \en{usual}, \en{pleasure} & \pron{mé-jeu} &
Comme le « j » français. Rare (environ 150 mots). \tabularnewline
\hline
\etoile /h/ & \en{hat}, \en{hello}, \en{house}, \en{behind} & \pron{hètt} &
\textbf{N'existe pas en français} (notre « h » est muet). C'est un souffle, comme sur des
lunettes pour les nettoyer. \textbf{Ne l'avalez jamais} : \en{he, him, her, house, happy,
hand, head, hold, hope}. Et ne l'ajoutez jamais là où il n'y en a pas : \en{hour},
\en{honest}, \en{honor}, \en{heir} sont muets. \tabularnewline
\hline
/t\ipash/ & \en{chair}, \en{watch}, \en{teach}, \en{nature} & \pron{tchè-eu} &
Un seul son, pas « t » + « ch » séparés. \tabularnewline
\hline
/d\ipazh/ & \en{joy}, \en{age}, \en{judge}, \en{gym} & \pron{djo-i} &
Un seul son. \en{j} = \en{d\ipazh} presque toujours (\en{jam}, \en{just}, \en{job}). \tabularnewline
\hline
/m/ & \en{man}, \en{time}, \en{come} & \pron{mann} & Identique. \tabularnewline
\hline
/n/ & \en{no}, \en{ten}, \en{when} & \pron{nou} & Identique. \tabularnewline
\hline
\etoile /\ipang/ & \en{sing}, \en{long}, \en{thing}, \en{English} & \pron{si-ng} &
\textbf{N'existe pas en français en fin de mot.} Arrière de la langue contre le palais,
\emph{sans} prononcer un « g » final. \en{sing} = \pron{si-ng} (pas \pron{si-ngue}).
Mais \en{finger} = \pron{fi-ngueu} (le \en{g} se fait entendre), \en{singer} = \pron{si-ngheu}. \tabularnewline
\hline
/l/ & \en{leg}, \en{light}, \en{all}, \en{feel} & \pron{lèg} &
Deux \en{l} anglais : clair devant voyelle (\en{light}), « sombre » --- langue en arrière ---
en fin de mot (\en{all}, \en{feel}, \en{milk}). Ne l'oubliez jamais en fin de mot. \tabularnewline
\hline
\etoile /r/ & \en{red}, \en{right}, \en{very}, \en{around} & \pron{rèd} &
\textbf{Très différent du \en{r} français} (qui vient de la gorge). Le \en{r} anglais :
langue recroquevillée \emph{au milieu} de la bouche, sans toucher le palais, lèvres un peu
arrondies. Aucun grattement. Et il se prononce \textbf{partout}, y compris avant consonne
(\en{car} \pron{kaa-r}, \en{four} \pron{foo-r}). \tabularnewline
\hline
\etoile /w/ & \en{we}, \en{water}, \en{work}, \en{want}, \en{now} & \pron{ouii} &
\textbf{N'existe pas comme phonème en français.} Lèvres très arrondies puis ouverture
rapide : c'est un « ou » qui glisse vers la voyelle suivante. \textbf{Jamais « v »} :
\en{we} $\neq$ \en{ve}, \en{wine} $\neq$ \en{vine}, \en{west} $\neq$ \en{vest}. \tabularnewline
\hline
/j/ & \en{yes}, \en{you}, \en{year}, \en{use} (verbe) & \pron{ièss} &
Le « y » anglais = « i » français de \en{bien}. Attention : \en{you} = \pron{iou},
\en{year} = \pron{ii-eu}. \tabularnewline
\end{longtable}
\end{tableau}

\begin{memo}[Les 6 sons à travailler en priorité absolue]
\begin{enumerate}
  \item \ipatheta\ et \ipadh\ (\en{think}, \en{this}) --- langue entre les dents.
  \item /h/ aspiré (\en{he}, \en{house}) --- ne jamais l'avaler.
  \item /w/ (\en{we}, \en{water}) --- ne jamais le dire « v ».
  \item /r/ anglais (\en{red}, \en{car}) --- jamais le \en{r} de gorge français.
  \item /\ipaiot/ contre /\ipaee\ipalong/ (\en{ship}/\en{sheep}) --- bref contre long.
  \item /\ipeschwa/ (le schwa) --- la clé du rythme anglais (chapitre 6).
\end{enumerate}
\end{memo}

\section{L'orthographe anglaise : ce qui se prononce et ce qui ne se prononce pas}

\subsection{Les lettres muettes (environ 60\,\% des mots anglais en contiennent)}
\begin{tableau}[Lettres muettes --- à mémoriser par groupes]
\footnotesize
\begin{longtable}{@{}p{2.2cm}p{5.2cm}p{7.2cm}@{}}
\entete{Lettre muette} & \entete{Mots} & \entete{Prononciation} \tabularnewline
\endfirsthead
\entete{Lettre muette} & \entete{Mots} & \entete{Prononciation} \tabularnewline
\endhead
\hline
\en{k} devant \en{n} & \en{know, knee, knife, knock, knight, knot} & \pron{nou}, \pron{nii}, \pron{naif}, \pron{nok}, \pron{naït} \tabularnewline
\en{w} devant \en{r} & \en{write, wrong, wrist, wrap, wreck} & \pron{raït}, \pron{rong} \tabularnewline
\en{b} final après \en{m} & \en{climb, bomb, thumb, comb, dumb, lamb} & \pron{klaïm}, \pron{bom}, \pron{theum} \tabularnewline
\en{t} dans certains mots & \en{listen, often, castle, whistle, Christmas} & \pron{li-s'n}, \pron{o-f'n}, \pron{kaa-s'l} \tabularnewline
\en{l} dans certains mots & \en{walk, talk, half, calm, would, should, could} & \pron{ouook}, \pron{toook}, \pron{haaf}, \pron{kaam}, \pron{woud} \tabularnewline
\en{g} devant \en{n} & \en{sign, design, foreign, gnaw} & \pron{saïn}, \pron{di-zaïn}, \pron{fo-rin} \tabularnewline
\en{h} initial & \en{hour, honest, honor, heir, herb (AmE)} & \pron{a-ou-eu}, \pron{o-nèst} \tabularnewline
\en{u} après \en{g} & \en{guess, guitar, guide, guest, build} & \pron{guèss}, \pron{gi-taa}, \pron{bild} \tabularnewline
\en{i} dans \en{-tion} & \en{nation, station, question, information} & \pron{né-i-cheun}, \en{question} = \pron{kwès-tchoun} \tabularnewline
\en{e} final & \en{make, time, come, house, name} & Le \en{e} est muet mais \textbf{allonge la voyelle} : \en{make} = \pron{mé-ik} et non \pron{mak}. \tabularnewline
\end{longtable}
\end{tableau}

\subsection{Un son, plusieurs orthographes}
\begin{itemize}
  \item \textbf{/\ipaee\ipalong/} : \en{see, sea, key, machine, people, believe, receive, field, these}
  \item \textbf{/\ipash/} : \en{she, nation, special, sure, ocean, conscience, sugar}
  \item \textbf{/t\ipash/} : \en{chair, watch, teach, nature, question, cello}
  \item \textbf{/\ipalv/} : \en{cup, love, money, does, blood, other, brother, enough}
  \item \textbf{/\ipaer\ipalong/} : \en{bird, work, learn, heard, girl, journey, early}
  \item \textbf{/\ipaei/} : \en{say, day, eight, name, they, rain, great, weigh, baby}
\end{itemize}

\subsection{Même orthographe, sons différents (les \en{heteronyms})}
\begin{tableau}[Mots qui changent de prononciation selon le sens]
\footnotesize
\begin{longtable}{@{}p{3.0cm}p{5.4cm}p{5.4cm}@{}}
\entete{Mot} & \entete{Nom (accent sur la 1\textsuperscript{re} syllabe)} & \entete{Verbe (accent sur la 2\textsuperscript{e} syllabe)} \tabularnewline
\endfirsthead
\entete{Mot} & \entete{Nom} & \entete{Verbe} \tabularnewline
\endhead
\hline
\en{record} & \pron{rè-kood} un disque, un record & \pron{ri-kood} enregistrer \tabularnewline
\en{present} & \pron{prè-z'nt} un cadeau & \pron{pri-zènt} présenter \tabularnewline
\en{object} & \pron{o-bdject} un objet & \pron{eub-dject} s'opposer \tabularnewline
\en{produce} & \pron{pro-diouss} des produits & \pron{preu-diouss} produire \tabularnewline
\en{conduct} & \pron{kon-deukt} le comportement & \pron{keun-deukt} diriger \tabularnewline
\en{permit} & \pron{peu-mitt} un permis & \pron{peu-mitt} permettre \tabularnewline
\en{increase} & \pron{in-kriiss} une augmentation & \pron{in-kriiss} augmenter \tabularnewline
\end{longtable}
\end{tableau}

\section{60 paires minimales : l'entraînement de l'oreille}
Une \emph{paire minimale} est un couple de mots qui ne diffèrent que par un son. C'est
l'exercice fondamental pour installer les sons anglais dans votre oreille.

\textbf{Protocole :} faites écouter les deux mots par un outil audio (Cambridge, Forvo),
puis répétez-les 5 fois chacun, puis demandez à quelqu'un (ou à une application) de vous en
dire un seul au hasard : vous devez identifier lequel. 5 paires par jour suffisent.

\begin{tableau}[60 paires minimales --- par son]
\footnotesize
\begin{longtable}{@{}p{2.6cm}p{4.0cm}p{4.0cm}p{3.6cm}@{}}
\entete{Contraste} & \entete{Mot 1} & \entete{Mot 2} & \entete{Sens} \tabularnewline
\endfirsthead
\entete{Contraste} & \entete{Mot 1} & \entete{Mot 2} & \entete{Sens} \tabularnewline
\endhead
\hline
\ipaiot\ / \ipaee\ipalong & \en{ship} & \en{sheep} & bateau / mouton \tabularnewline
\ipaiot\ / \ipaee\ipalong & \en{sit} & \en{seat} & s'asseoir / siège \tabularnewline
\ipaiot\ / \ipaee\ipalong & \en{live} & \en{leave} & habiter / quitter \tabularnewline
\ipaiot\ / \ipaee\ipalong & \en{fill} & \en{feel} & remplir / sentir \tabularnewline
\ipaiot\ / \ipaee\ipalong & \en{slip} & \en{sleep} & glisser / dormir \tabularnewline
\ipaiot\ / \ipaee\ipalong & \en{bit} & \en{beat} & un morceau / battre \tabularnewline
\ipaiot\ / \ipaee\ipalong & \en{hit} & \en{heat} & frapper / chaleur \tabularnewline
\ipaiot\ / \ipaee\ipalong & \en{fit} & \en{feet} & convenir / pieds \tabularnewline
\hline
\ipae\ / \ipaE & \en{cat} & \en{cut} & chat / coupure \tabularnewline
\ipae\ / \ipaE & \en{bad} & \en{bed} & mauvais / lit \tabularnewline
\ipae\ / \ipaE & \en{man} & \en{men} & homme / hommes \tabularnewline
\ipae\ / \ipaE & \en{hat} & \en{hot} & chapeau / chaud \tabularnewline
\ipae\ / \ipaE & \en{pan} & \en{pen} & poêle / stylo \tabularnewline
\ipae\ / \ipaE & \en{bag} & \en{beg} & sac / mendier \tabularnewline
\ipae\ / \ipaE & \en{sat} & \en{set} & assis / un ensemble \tabularnewline
\hline
\ipalv\ / \ipaups & \en{cup} & \en{coop} & tasse / poulailler \tabularnewline
\ipalv\ / \ipaoo\ipalong & \en{full} & \en{fool} & plein / fou \tabularnewline
\ipalv\ / \ipaa\ipalong & \en{cut} & \en{cart} & couper / chariot \tabularnewline
\ipalv\ / \ipaer\ipalong & \en{fun} & \en{fern} & amusement / fougère \tabularnewline
\hline
\ipatheta\ / s & \en{think} & \en{sink} & penser / évier \tabularnewline
\ipatheta\ / s & \en{thick} & \en{sick} & épais / malade \tabularnewline
\ipatheta\ / f & \en{three} & \en{free} & trois / libre \tabularnewline
\ipatheta\ / t & \en{thing} & \en{ting} & chose / tintement \tabularnewline
\ipadh\ / d & \en{they} & \en{day} & ils / jour \tabularnewline
\ipadh\ / d & \en{then} & \en{den} & alors / tanière \tabularnewline
\ipadh\ / z & \en{this} & \en{zoo} & ceci / zoo \tabularnewline
\ipadh\ / d & \en{those} & \en{doze} & ceux-là / somnoler \tabularnewline
\hline
/w/ / v & \en{west} & \en{vest} & ouest / gilet \tabularnewline
/w/ / v & \en{wine} & \en{vine} & vin / vigne \tabularnewline
/w/ / v & \en{wale} & \en{veil} & --- / voile \tabularnewline
/w/ / v & \en{wet} & \en{vet} & mouillé / vétérinaire \tabularnewline
\hline
/b/ / p & \en{bat} & \en{pat} & chauve-souris / tape \tabularnewline
/b/ / p & \en{berry} & \en{perry} & baie / --- \tabularnewline
/b/ / p & \en{ban} & \en{pan} & interdiction / poêle \tabularnewline
/b/ / p & \en{boat} & \en{coat} & bateau / manteau \tabularnewline
/b/ / p & \en{back} & \en{pack} & dos / paquet \tabularnewline
\hline
/d/ / t & \en{doe} & \en{toe} & biche / orteil \tabularnewline
/d/ / t & \en{dear} & \en{tear} & cher / larme \tabularnewline
/d/ / t & \en{dime} & \en{time} & dix cents / temps \tabularnewline
\hline
/g/ / k & \en{goat} & \en{coat} & chèvre / manteau \tabularnewline
/g/ / k & \en{gift} & \en{kift} & cadeau / --- \tabularnewline
/g/ / k & \en{gale} & \en{kale} & grand vent / chou \tabularnewline
\hline
/n/ /\ipang/ & \en{sin} & \en{sing} & péché / chanter \tabularnewline
/n/ /\ipang/ & \en{thin} & \en{thing} & mince / chose \tabularnewline
/n/ /\ipang/ & \en{ran} & \en{rang} & courut / sonna \tabularnewline
/n/ /\ipang/ & \en{ban} & \en{bang} & interdiction / bang \tabularnewline
\hline
/l/ / r & \en{light} & \en{right} & léger / droite \tabularnewline
/l/ / r & \en{lead} & \en{read} & mener / lire \tabularnewline
/l/ / r & \en{love} & \en{rove} & amour / errer \tabularnewline
/l/ / r & \en{fly} & \en{fry} & voler / frire \tabularnewline
/l/ / r & \en{collect} & \en{correct} & collecter / corriger \tabularnewline
/l/ / r & \en{low} & \en{row} & bas / ramer \tabularnewline
\hline
/z/ / s & \en{zip} & \en{sip} & fermeture éclair / gorgée \tabularnewline
/z/ / s & \en{bus} & \en{buzz} & bus / bourdonner \tabularnewline
/z/ / s & \en{rice} & \en{rise} & riz / montée \tabularnewline
/z/ / s & \en{place} & \en{plays} & lieu / joue \tabularnewline
\hline
/\ipash/ / s & \en{she} & \en{sea} & elle / mer \tabularnewline
/\ipash/ / t\ipash/ & \en{share} & \en{chair} & partager / chaise \tabularnewline
/\ipash/ / t\ipash/ & \en{shore} & \en{chore} & rivage / corvée \tabularnewline
/\ipazh/ / d\ipazh/ & \en{vision} & \en{Vishnu} & vision / --- \tabularnewline
\hline
\ipaou / \ipao\ipalong & \en{goat} & \en{got} & chèvre / obtenu \tabularnewline
\ipaou / \ipao & \en{no} & \en{now} & non / maintenant \tabularnewline
\ipaou / \ipaoo\ipalong & \en{hope} & \en{hoop} & espoir / cerceau \tabularnewline
/\ipaa\ipalong/ / \ipae/ & \en{cart} & \en{cat} & chariot / chat \tabularnewline
/\ipao\ipalong/ / \ipao\ (bref)/ & \en{caught} & \en{cot} & attrapé / lit d'enfant \tabularnewline
/\ipaiot\ipeschwa/ / \ipaE\ipeschwa/ & \en{here} & \en{hair} & ici / cheveux \tabularnewline
/\ipaiot\ipeschwa/ / \ipaE\ipeschwa/ & \en{near} & \en{nair} & près / --- \tabularnewline
/h/ / voyelle & \en{hill} & \en{ill} & colline / malade \tabularnewline
/h/ / voyelle & \en{heat} & \en{eat} & chaleur / manger \tabularnewline
/h/ / voyelle & \en{hat} & \en{at} & chapeau / à \tabularnewline
\end{longtable}
\end{tableau}

\section{Exercices}

\begin{exo}[5.1 --- Identifiez le son]
Écoutez (ou imaginez la prononciation correcte) et classez chaque mot dans la bonne colonne :
\en{ship -- sheep -- live -- leave -- feel -- fill -- sit -- seat -- heat -- hit -- rich -- reach}
\begin{center}
\begin{tabular}{@{}p{5.5cm}p{5.5cm}@{}}
\entete{/\ipaiot/ bref} & \entete{/\ipaee\ipalong/ long} \tabularnewline
\hline
\trou[4.5cm] & \trou[4.5cm] \\
\trou[4.5cm] & \trou[4.5cm] \\
\end{tabular}
\end{center}
\end{exo}
\begin{corrige}[5.1]
\textbf{Bref} /\ipaiot/ : \en{ship, live, fill, sit, hit, rich}.
\textbf{Long} /\ipaee\ipalong/ : \en{sheep, leave, feel, seat, heat, reach}.
\end{corrige}

\begin{exo}[5.2 --- Transcrivez à la française]
Donnez la prononciation approchée des douze mots suivants :
\en{water, thought, three, this, house, know, right, work, book, food, cup, go}.
\end{exo}
\begin{corrige}[5.2]
\en{water} \pron{ouo-teu} \quad \en{thought} \pron{thoot} \quad \en{three} \pron{thrii} \quad
\en{this} \pron{dhis} \quad \en{house} \pron{ha-ouss} \quad \en{know} \pron{nou} \quad
\en{right} \pron{raït} \quad \en{work} \pron{oueu-k} \quad \en{book} \pron{bouk} \quad
\en{food} \pron{fouud} \quad \en{cup} \pron{keup} \quad \en{go} \pron{gou-ou}.
\end{corrige}

\begin{exo}[5.3 --- Quelle paire ?]
Pour chaque phrase, choisissez le mot correct :
\begin{enumerate}
  \item I want to (sit / seat) down.
  \item The (ship / sheep) is very big on the sea.
  \item I (think / sink) you are right.
  \item This is my (wine / vine) --- it's from France.
  \item She has beautiful (hair / here).
  \item Don't (fill / feel) bad about it.
  \item The (cart / cut) is in the garden.
  \item We need a (bag / beg) for the shopping.
  \item He will (live / leave) tomorrow morning.
  \item My (dad / dead) works in a bank.
\end{enumerate}
\end{exo}
\begin{corrige}[5.3]
1. sit \quad 2. ship \quad 3. think \quad 4. wine \quad 5. hair \quad
6. feel \quad 7. cart \quad 8. bag \quad 9. leave \quad 10. dad
\end{corrige}

\begin{exo}[5.4 --- Les lettres muettes]
Dans chaque mot, entourez la lettre qui ne se prononce pas :
\en{knee, write, climb, listen, walk, hour, sign, guess, half, know, comb, island, talk, honest, wrist}.
\end{exo}
\begin{corrige}[5.4]
\en{knee (k), write (w), climb (b), listen (t), walk (l), hour (h), sign (g), guess (u),
half (l), know (k), comb (b), island (s), talk (l), honest (h), wrist (w)}.
\end{corrige}

\begin{mission}[Chapitre 5 --- 10 jours de sons]
\textbf{Jour 1-2 :} les voyelles longues et brèves (paires \en{sheep/ship, full/fool, cart/cut}).
\textbf{Jour 3-4 :} \ipatheta\ et \ipadh\ (30 mots, 5 fois chacun devant un miroir, langue visible entre les dents).
\textbf{Jour 5 :} /h/ aspiré (30 mots) --- expirez avant de parler.
\textbf{Jour 6 :} /w/ contre /v/ (20 paires).
\textbf{Jour 7 :} /r/ anglais (30 mots devant un miroir : la langue ne touche rien).
\textbf{Jour 8 :} \ipang\ final (\en{sing, long, thing, bring, wrong, song, young, strong}).
\textbf{Jour 9 :} les lettres muettes (tableau ci-dessus, 60 mots).
\textbf{Jour 10 :} test --- lisez à voix haute les 60 paires minimales, enregistrez-vous,
comparez avec un audio de référence, notez vos 3 sons les plus faibles. Ce sont vos
priorités du chapitre 7.
\end{mission}


% =====================================================================
\chapter{Accent tonique, rythme, intonation et \en{connected speech}}

Ce chapitre explique \textbf{pourquoi les francophones ne comprennent pas les anglophones
natifs} --- et ce n'est pas une question de vocabulaire.

\section{Deux rythmes incompatibles}
\begin{tableau}[Français contre anglais : la différence fondamentale]
\begin{tabularx}{\linewidth}{@{}p{4.0cm}X@{}}
\entete{} & \entete{Caractéristique} \tabularnewline
\hline
\textbf{Français} & Langue à \textbf{rythme syllabique} : chaque syllabe a à peu près la
même durée et la même force. \en{Je-ne-sais-pas} = quatre coups réguliers. L'accent tombe
sur la dernière syllabe du groupe. \tabularnewline
\hline
\textbf{Anglais} & Langue à \textbf{accent tonique} : les syllabes accentuées sont
\textbf{fortes, longues et hautes}, les autres sont \textbf{écrasées, courtes et réduites}
au schwa. Le temps entre deux accents est à peu près constant, quel que soit le nombre de
syllabes faibles entre elles. \tabularnewline
\hline
\textbf{Conséquence} & Un anglophone qui entend du français régulier le trouve « plat,
rapide, monotone ». Un francophone qui entend de l'anglais accentué trouve que les mots
« disparaissent » : il cherche \en{I want to go to the shop} et entend \en{I wanna go t'the
shop}. \textbf{Vous ne comprenez pas parce que vous cherchez des syllabes françaises là où
il n'y a que des accents anglais.} \tabularnewline
\end{tabularx}
\end{tableau}

\begin{memo}[La règle d'or du chapitre]
En anglais, ce qui compte n'est pas ce que vous dites, c'est \textbf{où vous frappez}.
Une syllabe accentuée mal placée rend un mot incompréhensible, même si tous ses sons sont
justes. Exemple : \en{banana} \pron{be-NA-neu} est compris ; \pron{ba-na-na} (à la
française) ne l'est pas.
\end{memo}

\section{L'accent tonique dans les mots}

\subsection{Les quatre principes}
\begin{enumerate}
  \item Tout mot anglais de deux syllabes ou plus a \textbf{une} syllabe accentuée
  (éventuellement une secondaire).
  \item Une syllabe accentuée est \textbf{plus longue, plus forte et plus aiguë}.
  \item Toute syllabe non accentuée tend vers le \textbf{schwa} /\ipeschwa/.
  \item L'accent est \textbf{imprévisible} --- il faut donc l'apprendre avec le mot, comme
  une partie du mot. Dans votre Anki, notez toujours la syllabe accentuée (par exemple en
  MAJUSCULES : \en{baNAna}, \en{COMfortable}, \en{deVELopment}).
\end{enumerate}

\subsection{Les règles qui marchent 80\,\% du temps}
\begin{tableau}[Règles de placement de l'accent]
\begin{tabularx}{\linewidth}{@{}p{5.6cm}X@{}}
\entete{Règle} & \entete{Exemples} \tabularnewline
\hline
\textbf{Nom de 2 syllabes} : accent sur la 1\textsuperscript{re} &
\en{TAble, WAter, MOney, CIty, PREsent, EXport, CHIna} \tabularnewline
\textbf{Verbe de 2 syllabes} : accent sur la 2\textsuperscript{e} &
\en{to beGIN, to preSENT, to exPORT, to deCIDE, to reLAX, to rePEAT} \tabularnewline
\textbf{Mots en \en{-tion}, \en{-sion}, \en{-cian}} : accent sur l'avant-dernière syllabe &
\en{naTION, inforMAtion, deciSIon, musIcian, televiSIon} \tabularnewline
\textbf{Mots en \en{-ic}} : accent sur l'avant-dernière &
\en{scienTIfic, draMAtic, econoMIc, fanTAStic, reaLIStic} \tabularnewline
\textbf{Mots en \en{-ity}, \en{-ify}, \en{-ise}} : accent 3 syllabes avant la fin &
\en{uniVERsity, possiBIlity, CLAssify, adVERtise, ORganise} \tabularnewline
\textbf{Mots en \en{-ee}, \en{-eer}, \en{-esque}} : accent \textbf{sur} la finale &
\en{emPLOyee, volunTEER, engiNEER, pictuRESque} \tabularnewline
\textbf{Composés noms} : accent sur le 1\textsuperscript{er} élément &
\en{GREENhouse, FOOTball, POSTman, BATHroom} \tabularnewline
\textbf{Nombres en \en{-teen}} : accent sur la 2\textsuperscript{e} ; \en{-ty} : sur la 1\textsuperscript{re} &
\en{thirTEEN} contre \en{THIRty} --- \textbf{différence capitale à l'oral} \tabularnewline
\end{tabularx}
\end{tableau}

\begin{piege}[\en{thirteen} ou \en{thirty} ?]
C'est la confusion la plus coûteuse de la vie quotidienne (prix, dates, numéros, âges).
\en{thirteen} = \pron{theu-tiin} (accent final, \en{ee} long) ; \en{thirty} =
\pron{THEU-di} (accent initial, \en{ir} bref). En cas de doute, l'anglophone demande :
\en{``Sorry, thirteen or thirty?''} --- entraînez la paire 13/30, 14/40, 15/50, 16/60,
17/70, 18/80, 19/90 jusqu'à l'automatisme.
\end{piege}

\subsection{Accent de mot et sens}
Certains mots changent de \textbf{sens} selon l'accent (voir tableau des \en{heteronyms},
chapitre 5). Retenez les dix plus utiles : \en{record, present, object, produce, conduct,
permit, increase, decrease, protest, suspect}.

\section{Le schwa : le son qui change tout}
/\ipeschwa/ est une voyelle neutre, à peine articulée, qui remplace \textbf{toute voyelle
d'une syllabe non accentuée}. C'est le son le plus fréquent de l'anglais.

\begin{tableau}[Le schwa se cache partout --- mêmes lettres, son réduit]
\footnotesize
\begin{longtable}{@{}p{3.4cm}p{3.0cm}p{8.0cm}@{}}
\entete{Mot} & \entete{Prononciation} & \entete{Où sont les schwas ?} \tabularnewline
\endfirsthead
\entete{Mot} & \entete{Prononciation} & \entete{Où sont les schwas ?} \tabularnewline
\endhead
\hline
\en{about} & \pron{e-baout} & Le \en{a} initial devient schwa. \tabularnewline
\en{banana} & \pron{be-NA-neu} & Les deux \en{a} non accentués deviennent schwa. \tabularnewline
\en{problem} & \pron{PRO-beulm} & Le \en{e} de la 2\textsuperscript{e} syllabe. \tabularnewline
\en{support} & \pron{se-POOT} & Le \en{u} initial. \tabularnewline
\en{photograph} & \pron{FOU-teu-graaf} & Le \en{o} de la 2\textsuperscript{e} syllabe. \tabularnewline
\en{photographer} & \pron{fe-TO-gra-feu} & Les trois voyelles non accentuées. \tabularnewline
\en{computer} & \pron{keum-Piou-teu} & \en{o} et \en{e} finaux. \tabularnewline
\en{interesting} & \pron{IN-tres-ting} & Souvent réduit à 3 syllabes, pas 5. \tabularnewline
\en{comfortable} & \pron{KOMF-teu-bl} & 3 syllabes, pas 5. \tabularnewline
\en{vegetable} & \pron{VEDJ-teu-bl} & 3 syllabes, pas 4. \tabularnewline
\en{today} & \pron{t'e-dé-i} & Le \en{o} devient \pron{t'}. \tabularnewline
\en{police} & \pron{pe-Liiss} & Le \en{o} initial. \tabularnewline
\end{longtable}
\end{tableau}

\begin{astuce}[Exercice du schwa --- 5 min/jour]
Prenez n'importe quel mot de trois syllabes ou plus, et dites-le de deux façons :
\textbf{(1)} à la française, toutes les syllabes égales ; \textbf{(2)} à l'anglaise, en
frappant la syllabe accentuée et en écrasant les autres. Enregistrez les deux versions.
Vous entendrez immédiatement la différence --- et vous comprendrez pourquoi les natifs ne
vous comprenaient pas.
\end{astuce}

\section{L'accent de phrase : les mots forts et les mots faibles}
Dans une phrase anglaise, on accentue les \textbf{mots de contenu} et on écrase les
\textbf{mots de forme}.

\begin{tableau}[Mots accentués et mots réduits]
\begin{tabularx}{\linewidth}{@{}p{7.5cm}X@{}}
\entete{\textbf{Accentués} (portent le sens)} & \entete{\textbf{Réduits} (structure)} \tabularnewline
\hline
Noms, verbes pleins, adjectifs, adverbes ; pronoms démonstratifs (\en{this, that}) ;
négations (\en{not, never}) ; mots interrogatifs (\en{what, where, why, how}) ;
auxiliaires en fin de phrase ou contractés avec emphase ; prépositions suivies d'un accent
de contraste. &
Articles (\en{a, an, the}) ; prépositions (\en{of, to, for, at, from, by}) ; conjonctions
(\en{and, but, or, as, that}) ; auxiliaires (\en{be, have, do, can, will, would}) ;
pronoms personnels (\en{he, him, her, them, us, you}) ; possessifs (\en{my, your, his}). \tabularnewline
\end{tabularx}
\end{tableau}

\ex{I want to go to the shop to buy some bread.}{Je veux aller à la boutique acheter du pain.}
\textbf{Version réelle :} \pron{aï WANT t' g' t' the SHOP t' BUY s'm BRED} ---
cinq accents forts, tout le reste écrasé.

\ex{She can speak French and English.}{Elle sait parler français et anglais.}
\textbf{Version réelle :} \pron{ch' k'n SPIIK French 'n INGLISH}.

\begin{piege}[La négation et le contraste doivent être frappés fort]
\en{I \textbf{didn't} say that} : si vous réduisez \en{didn't}, on entend \en{I said that}
--- l'inverse du sens. De même : \en{I can do it} (\en{can} faible \pron{k'n}) contre
\en{I \textbf{can't} do it} (\en{can't} fort \pron{kaant}). Entraînez-vous à frapper
\en{not, never, no, can't, don't, isn't, wasn't}.
\end{piege}

\section{L'intonation : la mélodie des phrases}
\begin{tableau}[Les 5 courbes d'intonation essentielles]
\begin{tabularx}{\linewidth}{@{}p{4.2cm}p{2.6cm}X@{}}
\entete{Type de phrase} & \entete{Mélodie} & \entete{Exemple} \tabularnewline
\hline
Question fermée (oui/non) & Montante \ra & \en{Are you coming?} \pron{a y' CO-ming $\monte$} \tabularnewline
Question ouverte (\en{wh-}) & \textbf{Descendante} \descend & \en{Where do you live?} --- \textbf{erreur fréquente : les francophones la font monter.} \tabularnewline
Affirmation & Descendante \descend & \en{I live in Dakar.} \pron{i LIVE in DA-kaar $\descend$} \tabularnewline
Liste & Montante, montante\ldots\ puis descendante & \en{I bought bread$\monte$, milk$\monte$ and eggs$\descend$.} \tabularnewline
Surprise, doute, politesse & Très montante & \en{Really?} \pron{rii-AL-lii $\monte$} \quad \en{Sorry?} \tabularnewline
\end{tabularx}
\end{tableau}

\begin{astuce}[L'intonation porte le message émotionnel]
La même phrase dite avec trois mélodies donne trois sens :\\
\en{``You're coming.''} \descend\ = constatation. \quad
\en{``You're coming?''} \monte\ = question étonnée. \quad
\en{``\textbf{You're} coming?''} \monte\ avec accent sur \en{you} = « \emph{toi}, tu viens ? ».
\end{astuce}

\section{\en{Connected speech} : pourquoi les mots disparaissent}
Les anglophones ne disent pas les mots un par un : ils les enchaînent en un seul flux sonore.
Voici les sept mécanismes --- chacun avec un exemple que vous devez savoir reconnaître à l'oreille.

\subsection{1. Le \emph{linking} (consonne finale + voyelle initiale)}
La consonne finale d'un mot colle à la voyelle du mot suivant, comme une liaison française.
\begin{itemize}
  \item \en{an apple} \pron{e-NA-pl} --- le \en{n} passe dans \en{apple}.
  \item \en{turn it off} \pron{teu-ri-NI-tof}.
  \item \en{not at all} \pron{no-ta-TAAL}.
  \item \en{pick it up} \pron{pi-ki-TUP}.
  \item \en{an hour ago} \pron{e-nau-e-ROU}.
\end{itemize}

\subsection{2. Le \en{r} de liaison (britannique)}
Le \en{r} final, muet en britannique, se prononce devant une voyelle :
\en{far away} \pron{faa-re-wé-i}, \en{for example} \pron{fo-rig-Zaampl},
\en{here it is} \pron{hi-e-ri-TIZ}.

\subsection{3. Les intrusions /w/ et /j/}
Entre deux voyelles, l'anglais insère un son de transition :
\begin{itemize}
  \item \en{go out} \pron{gou-Waout} ; \en{do it} \pron{dou-WIT} ; \en{who is it?} \pron{hou-WI-zit}.
  \item \en{I am} \pron{aï-YAM} ; \en{see it} \pron{sii-YIT} ; \en{the end} \pron{dhi-YEND}.
\end{itemize}

\subsection{4. Les élisions (sons supprimés)}
\begin{itemize}
  \item \en{next day} \pron{nèks-dé-i} (le \en{t} saute).
  \item \en{last night} \pron{laas-naït}.
  \item \en{handbag} \pron{hann-bag}.
  \item \en{comfortable} \pron{KOMF-teu-bl} ; \en{family} \pron{FAM-li} ; \en{camera} \pron{KAM-reu}.
  \item \en{probably} \pron{PROB-li} ; \en{actually} \pron{AK-chou-a-li} ; \en{chocolate} \pron{TCHOK-l't}.
\end{itemize}

\subsection{5. Les assimilations (sons qui se transforment)}
\begin{itemize}
  \item \en{t} + \en{y} \ra \en{ch} : \en{don't you} \pron{doun-tchou} ; \en{meet you} \pron{mii-tchou}.
  \item \en{d} + \en{y} \ra \en{j} : \en{did you} \pron{di-djou} ; \en{would you} \pron{wou-djou}.
  \item \en{s} + \en{y} \ra \en{ch} : \en{miss you} \pron{mi-chou}.
  \item \en{z} + \en{y} \ra \en{j} : \en{was your} \pron{ouo-djou}.
\end{itemize}

\subsection{6. Le \emph{glottal stop} (coup de glotte)}
Le \en{t} entre voyelles ou devant consonne devient un arrêt net de la gorge (très britannique) :
\en{butter} \pron{ba-eu}, \en{football} \pron{fou-bo-l}, \en{not now} \pron{no-naou},
\en{water} \pron{ouo-eu}. En américain, ce \en{t} devient plutôt un \en{d} bref
(\en{water} = \pron{oua-deu}).

\subsection{7. Les formes faibles (\emph{weak forms})}
\begin{tableau}[Formes faibles des petits mots --- à reconnaître absolument]
\footnotesize
\begin{longtable}{@{}p{2.4cm}p{3.2cm}p{2.6cm}p{5.8cm}@{}}
\entete{Mot} & \entete{Forme forte} & \entete{Forme faible} & \entete{Exemple} \tabularnewline
\endfirsthead
\entete{Mot} & \entete{Forte} & \entete{Faible} & \entete{Exemple} \tabularnewline
\endhead
\hline
\en{and} & \pron{and} & \pron{'n} & \en{fish and chips} = \pron{fi-ch'n tchips} \tabularnewline
\en{of} & \pron{ov} & \pron{euv} (schwa) & \en{a cup of tea} = \pron{e-keup-euv-tii} \tabularnewline
\en{to} & \pron{tou} & \pron{t'} & \en{go to bed} = \pron{gou-t'-bèd} \tabularnewline
\en{for} & \pron{foo} & \pron{fe} & \en{for me} = \pron{fe-mii} \tabularnewline
\en{can} & \pron{kann} & \pron{k'n} & \en{I can swim} = \pron{aï k'n SWIM} \tabularnewline
\en{have} & \pron{hav} & \pron{hev/'v} & \en{should have} = \pron{chou-dev} \ra \en{shoulda} \tabularnewline
\en{do} & \pron{dou} & \pron{de} & \en{do you know?} = \pron{d'ye NOU} \tabularnewline
\en{at} & \pron{at} & \pron{et} & \en{at work} = \pron{et-oueu-k} \tabularnewline
\en{from} & \pron{from} & \pron{freum} & \en{from Dakar} = \pron{freum-da-kaar} \tabularnewline
\en{them} & \pron{dhèm} & \pron{dhem/'em} & \en{tell them} = \pron{tè-l'em} \tabularnewline
\en{us} & \pron{euss} & \pron{es} & \en{let us go} = \en{let's go} = \pron{lèts-gou} \tabularnewline
\en{was} & \pron{ouoz} & \pron{weu-z} & \en{it was good} = \pron{i-teu-z-GOUD} \tabularnewline
\en{your} & \pron{yoo} & \pron{yeu} & \en{your name} = \pron{yeu-né-im} \tabularnewline
\en{he/him/her} & \pron{hii/im/heu} & \pron{i/im/eu} & \en{tell him} = \pron{té-lim} ; \en{I saw her} = \pron{aï soo-weu} \tabularnewline
\end{longtable}
\end{tableau}

\subsection{Les réductions familières (à comprendre, à utiliser avec prudence)}
\begin{tableau}[Réductions courantes de l'anglais parlé]
\footnotesize
\begin{longtable}{@{}p{3.2cm}p{3.0cm}p{8.0cm}@{}}
\entete{Forme réduite} & \entete{Forme complète} & \entete{Exemple et traduction} \tabularnewline
\endfirsthead
\entete{Réduite} & \entete{Complète} & \entete{Exemple} \tabularnewline
\endhead
\hline
\en{gonna} & \en{going to} & \en{I'm gonna call you.} --- \textit{Je vais t'appeler.} \tabularnewline
\en{wanna} & \en{want to} & \en{Do you wanna eat?} --- \textit{Tu veux manger ?} \tabularnewline
\en{gotta} & \en{got to / have to} & \en{I gotta go.} --- \textit{Je dois y aller.} \tabularnewline
\en{hafta} & \en{have to} & \en{We hafta leave now.} --- \textit{On doit partir maintenant.} \tabularnewline
\en{oughta} & \en{ought to} & \en{You oughta see a doctor.} --- \textit{Tu devrais voir un médecin.} \tabularnewline
\en{supposta} & \en{supposed to} & \en{What was I supposta do?} --- \textit{Qu'est-ce que j'étais censé faire ?} \tabularnewline
\en{used to} & \en{used to} & \pron{iou-s-tou} --- \textit{avant, j'avais l'habitude de\ldots} \tabularnewline
\en{lemme} & \en{let me} & \en{Lemme see.} --- \textit{Laisse-moi voir.} \tabularnewline
\en{gimme} & \en{give me} & \en{Gimme a minute.} --- \textit{Donne-moi une minute.} \tabularnewline
\en{dunno} & \en{don't know} & \en{I dunno.} --- \textit{Je sais pas.} \tabularnewline
\en{kinda / sorta} & \en{kind of / sort of} & \en{It's kinda cold.} --- \textit{C'est un peu froid.} \tabularnewline
\en{d'you / dya} & \en{do you} & \en{D'you want some?} --- \textit{Tu en veux ?} \tabularnewline
\en{whaddaya} & \en{what do you} & \en{Whaddaya mean?} --- \textit{Qu'est-ce que tu veux dire ?} \tabularnewline
\en{c'mon} & \en{come on} & \en{C'mon, let's go!} --- \textit{Allez, on y va !} \tabularnewline
\en{y'know} & \en{you know} & \en{It's hard, y'know.} --- \textit{C'est dur, tu sais.} \tabularnewline
\en{innit} \en{(BrE)} & \en{isn't it} & \en{It's cold, innit?} --- \textit{Il fait froid, non ?} \tabularnewline
\end{longtable}
\end{tableau}

\begin{astuce}[Les comprendre, oui ; les utiliser, plus tard]
Vous devez \textbf{reconnaître} toutes ces réductions à l'oreille dès le mois 3 --- sinon les
séries resteront incompréhensibles. Mais ne les \textbf{employez} qu'à partir du mois 8 ou 9,
et seulement dans un contexte familier : un \en{gonna} dans un entretien d'embauche est
maladroit, dans une discussion entre amis il est naturel.
\end{astuce}

\section{Exercices de rythme et d'intonation}

\begin{exo}[6.1 --- Marquez les accents]
Écrivez les mots suivants en majuscules sur la syllabe accentuée :
\en{banana, information, comfortable, development, photographer, university, employee,
interesting, television, tomorrow, yesterday, computer}.
\end{exo}
\begin{corrige}[6.1]
\en{baNAna, inforMAtion, KOMfortable, deVElopment, phoTOgrapher, uniVERsity, employEE,
INteresting, TELevision, toMOrrow, YEStarday, comPUter}.
\end{corrige}

\begin{exo}[6.2 --- Nom ou verbe ?]
Prononcez correctement selon le sens, puis traduisez :
\begin{enumerate}
  \item I bought a new \en{record}. / Please \en{record} my voice.
  \item She gave me a \en{present}. / He will \en{present} the project.
  \item There was an \en{increase} in prices. / Prices \en{increase} every year.
  \item Do you \en{object}? / It's a strange \en{object}.
\end{enumerate}
\end{exo}
\begin{corrige}[6.2]
1. \pron{RÈ-kood} un disque / \pron{ri-KOOD} enregistrer. 2. \pron{PRÈ-z'nt} un cadeau /
\pron{pri-ZÈNT} présenter. 3. \pron{IN-kriiss} une augmentation / \pron{in-KRISS} augmenter.
4. \pron{eub-DJECT} s'opposer / \pron{OB-dject} un objet.
\end{corrige}

\begin{exo}[6.3 --- Écrivez ce que vous entendez]
Remettez ces phrases réduites en anglais correct et complet :
\begin{enumerate}
  \item \en{``I'm gonna call ya later.''}
  \item \en{``D'you wanna cuppa tea?''}
  \item \en{``He shoulda told us.''}
  \item \en{``Whaddaya doin'?''}
  \item \en{``Lemme see it, gimme a sec.''}
  \item \en{``I dunno what he's talkin' about.''}
  \item \en{``We gotta go, it's kinda late.''}
  \item \en{``Tell 'em I said hi.''}
\end{enumerate}
\end{exo}
\begin{corrige}[6.3]
\begin{enumerate}
  \item \en{I am going to call you later.} --- \textit{Je t'appellerai plus tard.}
  \item \en{Do you want a cup of tea?} --- \textit{Tu veux une tasse de thé ?}
  \item \en{He should have told us.} --- \textit{Il aurait dû nous le dire.}
  \item \en{What are you doing?} --- \textit{Qu'est-ce que tu fais ?}
  \item \en{Let me see it, give me a second.} --- \textit{Laisse-moi voir, donne-moi une seconde.}
  \item \en{I don't know what he is talking about.} --- \textit{Je ne sais pas de quoi il parle.}
  \item \en{We have got to go, it is kind of late.} --- \textit{On doit y aller, il est un peu tard.}
  \item \en{Tell them I said hello.} --- \textit{Dis-leur que je leur ai dit bonjour.}
\end{enumerate}
\end{corrige}

\begin{exo}[6.4 --- Intonation]
Indiquez la mélodie (montante \monte\ ou descendante \descend) :
\begin{enumerate}
  \item \en{Do you speak English?}
  \item \en{Where do you come from?}
  \item \en{My name is Ali.}
  \item \en{I need bread, milk and eggs.}
  \item \en{You're joking!}
  \item \en{Can you help me?}
  \item \en{What time is it?}
  \item \en{It's half past nine.}
\end{enumerate}
\end{exo}
\begin{corrige}[6.4]
1. \monte\ \ 2. \descend\ \ 3. \descend\ \ 4. \monte\monte\descend\ (sur \en{eggs}) \
5. \monte\ \ 6. \monte\ \ 7. \descend\ \ 8. \descend
\end{corrige}

\begin{mission}[Chapitre 6 --- le drill des 5 accents]
Pendant 7 jours, prenez chaque jour 5 phrases de ce livre (n'importe lesquelles) et :
\textbf{(1)} marquez au crayon les syllabes accentuées ; \textbf{(2)} lisez-les en exagérant
les accents (fort et haut) et en écrasant le reste ; \textbf{(3)} enregistrez-vous ;
\textbf{(4)} comparez à un natif (YouGlish : tapez la phrase, écoutez 5 personnes la dire).
À la fin de la semaine, votre rythme aura changé --- les gens vous diront que vous
« sonnez mieux », sans savoir pourquoi.
\end{mission}


% =====================================================================
\chapter{Clinique de prononciation pour francophones}

Voici les quinze erreurs systématiques des francophones (dont les Africains francophones),
avec pour chacune : la cause, la correction et un entraînement. C'est le chapitre le plus
rentable du livre : chacune de ces erreurs corrigée améliore immédiatement votre
intelligibilité.

\section{Les 15 erreurs et leur correction}

\subsection{Erreur 1 --- Avaler le \en{h}}
\textbf{Ce que vous dites :} \en{he} \pron{i}, \en{house} \pron{ouss}, \en{happy} \pron{api}.
\textbf{Ce qu'on entend :} \en{eat, ounce, appy} --- ou simplement « un accent étrange ».
\textbf{Cause :} le \en{h} français est muet, votre bouche a pris l'habitude de ne pas souffler.
\textbf{Correction :} le \en{h} anglais est un vrai souffle. Expirez avant de prononcer la
voyelle, comme pour faire de la buée sur une vitre.
\textbf{Drill :} \en{he, him, her, his, here, house, home, hope, hand, head, hair, hat, hot,
heart, help, hold, how, hello, behind, ahead, perhaps} --- 3 fois chacun, main devant la
bouche pour sentir le souffle.

\subsection{Erreur 2 --- Transformer le \en{th} en \en{s}, \en{z}, \en{f} ou \en{d}}
\textbf{Ce que vous dites :} \en{think} \pron{sink}, \en{this} \pron{dis}, \en{three}
\pron{free}, \en{brother} \pron{brozeu}.
\textbf{Correction :} langue \emph{entre} les dents --- vraiment entre, on doit la voir.
\ipatheta\ = souffle sans voix (\en{think}) ; \ipadh\ = souffle avec voix (\en{this}).
\textbf{Drill :} \en{think, thing, three, thank, thought, through, thousand, thumb, bath,
mouth, month, birthday} --- puis \en{this, that, these, those, there, they, them, then,
than, mother, father, brother, other, weather, with}.
\textbf{Astuce visuelle :} répétez devant un miroir jusqu'à voir votre langue dépasser.

\subsection{Erreur 3 --- Le \en{w} prononcé \en{v}}
\textbf{Ce que vous dites :} \en{we} \pron{vi}, \en{water} \pron{vateu}, \en{work} \pron{veuk}.
\textbf{Correction :} /w/ = lèvres très arrondies et projetées en avant, puis ouverture
rapide --- comme un « ou » qui glisse. Aucune lèvre sur les dents (c'est cela, le \en{v}).
\textbf{Drill :} \en{we, will, would, water, work, world, want, when, where, why, which,
woman, window, week, way, well, always, now, know, grow, show} --- en alternance avec
\en{ve, veal, very, vest, vet, vine, visit, voice}.

\subsection{Erreur 4 --- Le \en{r} de gorge français}
\textbf{Correction :} le \en{r} anglais ne vient pas de la gorge. Recroquevillez la langue
au milieu de la bouche, sans contact, lèvres légèrement arrondies. C'est un son « mou »,
proche d'un \en{w} très bref.
\textbf{Drill :} \en{red, right, run, rain, road, room, read, write, rice, ring, round,
river, around, sorry, really, very, every, remember, important, world, work, girl, bird}.
\textbf{Test :} dites \en{``around the world''} dix fois : si votre gorge gratte, recommencez.

\subsection{Erreur 5 --- Ajouter un \en{e} final (le fameux « euh »)}
\textbf{Ce que vous dites :} \en{book} \pron{bou-keu}, \en{dog} \pron{do-gueu}, \en{stop}
\pron{sto-peu}, \en{work} \pron{oueu-keu}.
\textbf{Cause :} le français n'aime pas finir sur une consonne ; votre bouche ajoute une
voyelle de soutien.
\textbf{Correction :} arrêtez le son net. Fermez la bouche à la fin de \en{book},
coupez le souffle à la fin de \en{stop}.
\textbf{Drill :} \en{book, look, take, make, work, walk, talk, stop, jump, help, want, need,
good, bad, big, small, cat, dog, map, cup, black, white, red, green, blue} --- en terminant
chaque mot par un silence franc.

\subsection{Erreur 6 --- Ne pas distinguer \en{ship} et \en{sheep}}
Voir chapitre 5 : bref contre long, tendu contre relâché. Drill : les 20 paires en
/\ipaiot/ -- /\ipaee\ipalong/ de la liste des paires minimales.

\subsection{Erreur 7 --- \en{no} prononcé \pron{no}}
\textbf{Correction :} le \en{o} anglais final est une diphtongue : \pron{nou-ou}.
Idem pour \en{go, so, home, phone, only, open, close, cold, both, note}.
\textbf{Drill :} \en{no, go, so, home, phone, hope, note, both, cold, only, open, close,
over, road, show, know, slow, grow, throw, alone}.

\subsection{Erreur 8 --- L'intonation plate et le rythme syllabique}
Vous parlez anglais avec le rythme du français : toutes les syllabes égales. Résultat : les
natifs doivent faire un effort pour vous comprendre, et vous fatiguez vite.
\textbf{Correction :} exagérez. Frappez les mots de contenu, écrasez les petits mots.
C'est inconfortable au début --- c'est exactement ce qu'il faut.
\textbf{Drill :} shadowing quotidien (chapitre 3) sur des phrases courtes.

\subsection{Erreur 9 --- Le \en{-s} final oublié ou mal sonorisé}
\textbf{Ce que vous dites :} \en{he work} au lieu de \en{he works}, \en{two book}.
\textbf{Et le pire :} le \en{-s} du pluriel se prononce /s/, /z/ ou /\ipaiot z/ selon le mot :
\en{cats} \pron{kats}, \en{dogs} \pron{dogz}, \en{boxes} \pron{bok-siz}.
\textbf{Correction :} après une consonne sourde \ra /s/ ; après une voyelle ou une consonne
sonore \ra /z/ ; après \en{s, z, sh, ch, x} \ra /iz/.
\textbf{Drill :} \en{books / bags / buses ; cats / dogs / houses ; maps / pens / watches}.

\subsection{Erreur 10 --- Le \en{ng} final prononcé \pron{ingue}}
\en{sing} = \pron{si-ng}, pas \pron{si-ngue}. La langue touche le palais, le son est nasal,
et \textbf{rien ne suit}. Drill : \en{sing, song, long, strong, bring, thing, wrong, young,
king, ring, morning, evening, nothing, everything, interesting, English}.

\subsection{Erreur 11 --- Les voyelles longues abrégées}
\en{car} = \pron{kaa}, pas \pron{kar}. \en{four} = \pron{foo-r}. \en{bird} = \pron{beuud}.
Le français n'a pas de voyelles longues ; l'anglais en tire des distinctions de sens.
Drill : tenez chaque voyelle une seconde : \en{car, park, heart, start, father} ---
\en{four, more, door, floor, for} --- \en{bird, word, world, girl, third}.

\subsection{Erreur 12 --- Le \en{l} final disparu}
\en{all} \pron{oo} au lieu de \pron{ool}. \en{milk} \pron{mi-k}. \en{people} \pron{piip}.
\textbf{Correction :} le \en{l} final anglais se fait en levant la pointe de la langue vers
les dents du haut, à la fin du mot. Drill : \en{all, call, small, well, tell, feel, milk,
help, people, table, ball, full, cool, school, world, girl}.

\subsection{Erreur 13 --- La liaison à la française là où l'anglais ne lie pas}
En anglais, on ne lie pas un \en{h} (\en{the house} ne devient pas \pron{dhi-zouse}) et on
ne lie pas à travers un mot accentué. En revanche, on lie volontiers consonne \ra voyelle.
\textbf{Correction :} liez seulement quand un mot finit par une consonne et que le suivant
commence par une voyelle : \en{an apple, my old car, turn it on}.

\subsection{Erreur 14 --- Le \en{dark} et le \en{clear} \en{l}, le \en{t} américain}
\en{water} : britannique \pron{ouo-teu} (coup de glotte), américain \pron{oua-deu}.
Choisissez \textbf{une} variété et tenez-vous-y les premiers mois --- mélanger les deux
produit un accent flou. Recommandation : l'américain est plus présent dans les médias,
le britannique est plus enseigné en Afrique de l'Ouest. Les deux sont parfaitement acceptés.

\subsection{Erreur 15 --- Parler trop doucement par peur, ou trop vite par panique}
Trop doucement : on ne vous entend pas, on vous fait répéter, vous perdez confiance.
Trop vite : vous mangez les finales et les accents. \textbf{Correction :} visez un débit
\emph{moyen, régulier}, avec des pauses claires entre les groupes de mots. La pause est un
outil de puissance, pas un aveu de faiblesse : les bons orateurs s'arrêtent.

\section{30 \en{tongue twisters} (virelangues) --- votre salle de sport}
Travaillez-les \textbf{lentement d'abord, puis vite}, trois fois de suite chacun, sans erreur.
Cinq minutes par jour suffisent.

\begin{tableau}[30 virelangues classés par son]
\footnotesize
\begin{longtable}{@{}p{2.6cm}p{7.6cm}p{4.0cm}@{}}
\entete{Son visé} & \entete{Virelangue} & \entete{Traduction} \tabularnewline
\endfirsthead
\entete{Son} & \entete{Virelangue} & \entete{Traduction} \tabularnewline
\endhead
\hline
\ipatheta & \en{The thirty-three thieves thought that they thrilled the throne throughout Thursday.} & Les trente-trois voleurs pensaient avoir ému le trône tout au long de jeudi. \tabularnewline
\ipatheta & \en{Three thin thinkers thought thick thoughts.} & Trois minces penseurs pensaient d'épaisses pensées. \tabularnewline
\ipadh & \en{This is the sixth sick sheikh's sixth sheep.} & C'est la sixième brebis du sixième cheikh malade. \tabularnewline
\ipatheta/\ipadh & \en{They threw three thick things through the thin things.} & Ils ont jeté trois choses épaisses à travers les choses minces. \tabularnewline
/s/ / \ipash & \en{She sells seashells by the seashore.} & Elle vend des coquillages au bord de la mer. \tabularnewline
/s/ / \ipash & \en{Surely Sylvia swims slowly, silently.} & Sûrement Sylvia nage lentement, silencieusement. \tabularnewline
/p/ / b/ & \en{Peter Piper picked a peck of pickled peppers.} & Peter Piper a pris une pincée de poivrons au vinaigre. \tabularnewline
/b/ & \en{A big black bug bit a big black bear.} & Un gros insecte noir a mordu un gros ours noir. \tabularnewline
/f/ / v/ & \en{Four fine fresh figs for you.} & Quatre belles figues fraîches pour toi. \tabularnewline
/w/ & \en{Which witch is which?} & Quelle sorcière est laquelle ? \tabularnewline
/w/ / v/ & \en{Very well, Will we visit the west village Wednesday?} & Très bien, Will, visiterons-nous le village de l'ouest mercredi ? \tabularnewline
/w/ & \en{We surely shall see the sun shine shortly.} & Nous verrons sûrement bientôt briller le soleil. \tabularnewline
/r/ / l/ & \en{Red lorry, yellow lorry.} \quad (×10, vite) & Camion rouge, camion jaune. \tabularnewline
/r/ / l/ & \en{Really leathery, rarely leathery.} & Vraiment coriace, rarement coriace. \tabularnewline
/r/ / l/ & \en{Rolling red wagons, roaring red lions.} & Wagons rouges qui roulent, lions rouges qui rugissent. \tabularnewline
/r/ / w/ & \en{Round the rugged rock the ragged rascal ran.} & Autour du rocher rugueux, le coquin déguenailé courut. \tabularnewline
/l/ & \en{Lovely lemon liniment, lively little lady.} & Joli liniment au citron, petite dame vive. \tabularnewline
/l/ / n/ & \en{Loyal, royal, loyal, royal.} & Loyal, royal, loyal, royal. \tabularnewline
/t/ / d/ & \en{Two tiny tigers take two taxis to town.} & Deux petits tigres prennent deux taxis en ville. \tabularnewline
/d/ & \en{Double Dutch dancers dance daily downtown.} & Les danseurs de Double Dutch dansent tous les jours au centre-ville. \tabularnewline
/k/ / g/ & \en{Good blood, bad blood.} \quad (×10) & Bon sang, mauvais sang. \tabularnewline
/k/ & \en{Can you can a can as a canner can can a can?} & Peux-tu mettre en boîte comme un metteur en boîte peut mettre en boîte ? \tabularnewline
/m/ / n/ & \en{Many mumbling mice are making merry music in the moonlight.} & Beaucoup de souris marmonnantes font de la joyeuse musique au clair de lune. \tabularnewline
/n/ & \en{Nine nice night nurses nursing nicely.} & Neuf gentilles infirmières de nuit soignent gentiment. \tabularnewline
\ipang & \en{The king and the queen are singing in the morning.} & Le roi et la reine chantent le matin. \tabularnewline
\ipae/ \ipaE & \en{A sad man sat on a black cat's mat.} & Un homme triste s'est assis sur le tapis d'un chat noir. \tabularnewline
\ipaee/ \ipaiot & \en{Six sheep sleep in a ship's sheet.} & Six moutons dorment dans le drap d'un navire. \tabularnewline
\ipaou & \en{No one knows how old the old oak grows.} & Personne ne sait quel âge a le vieux chêne. \tabularnewline
/h/ & \en{How much wood would a woodchuck chuck if a woodchuck could chuck wood?} & Combien de bois une marmotte jetterait si une marmotte pouvait jeter du bois ? \tabularnewline
Tout & \en{I thought a thought, but the thought I thought wasn't the thought I thought I thought.} & J'ai pensé une pensée, mais la pensée que j'ai pensée n'était pas celle que je croyais penser. \tabularnewline
\end{longtable}
\end{tableau}

\section{50 phrases-drills : la prononciation en contexte}
Lisez chaque phrase trois fois (lentement, normalement, vite et expressif), puis enregistrez
les 10 premières et comparez à un natif.

\begin{tableau}[50 phrases à maîtriser phonétiquement]
\footnotesize
\begin{longtable}{@{}p{0.7cm}p{7.4cm}p{6.1cm}@{}}
\entete{\#} & \entete{Phrase} & \entete{Point travaillé} \tabularnewline
\endfirsthead
\entete{\#} & \entete{Phrase} & \entete{Point travaillé} \tabularnewline
\endhead
1 & \en{I think this is the third thing they thought about.} & \ipatheta, \ipadh \tabularnewline
2 & \en{He has a house in the heart of the city.} & /h/ aspiré \tabularnewline
3 & \en{We will work with water and wood every week.} & /w/ \tabularnewline
4 & \en{The red car around the world really runs right.} & /r/ anglais \tabularnewline
5 & \en{She should have looked at the book on the shelf.} & finales consonantiques \tabularnewline
6 & \en{The sheep on the ship sleep six sheets deep.} & \ipaee\ contre \ipaiot \tabularnewline
7 & \en{I want a hot cup of tea, not a cold one.} & \ipae, \ipao, \ipalv \tabularnewline
8 & \en{My brother's mother and father live in the north.} & \ipadh, \ipatheta \tabularnewline
9 & \en{No one knows how old the old oak grows.} & diphtongue \ipaou \tabularnewline
10 & \en{The king is singing a long song in the morning.} & \ipang \tabularnewline
11 & \en{It's a beautiful day, isn't it?} & schwa, question tag \tabularnewline
12 & \en{What do you want to do today?} & réductions \en{whaddaya wanna} \tabularnewline
13 & \en{I'm going to have to go now.} & \en{gonna hafta} \tabularnewline
14 & \en{Could you tell him I called?} & \en{couldja tellim} \tabularnewline
15 & \en{I don't know what she's talking about.} & \en{dunno whashe talkinabout} \tabularnewline
16 & \en{Would you like a cup of coffee or a glass of water?} & \en{wouja}, formes faibles \tabularnewline
17 & \en{The weather is better than yesterday.} & \ipadh, comparatif \tabularnewline
18 & \en{He's been working here for about three years.} & \en{present perfect continuous} \tabularnewline
19 & \en{I've never seen anything like that before.} & \en{I've} \pron{aïv}, \en{anything} \tabularnewline
20 & \en{She shouldn't have said that to him.} & \en{shouldn't have} = \pron{chou-dent-ev} \tabularnewline
21 & \en{Where were you when I was waiting for you?} & /w/ + \ipaer\ipalong \tabularnewline
22 & \en{My father's favourite film is a French film.} & \ipaer\ipalong, /f/ \tabularnewline
23 & \en{Actually, I'd rather have a vegetable salad.} & réductions de syllabes \tabularnewline
24 & \en{The library is behind the bakery, near the river.} & \en{library} \pron{LAÏ-bre-ri} \tabularnewline
25 & \en{February is the shortest month of the year.} & \en{February} \pron{FÉ-brou-é-ri} \tabularnewline
26 & \en{I usually go to the university by bus.} & /\ipazh/ dans \en{usually} \tabularnewline
27 & \en{The children are playing in the garden.} & \en{children} \pron{TCHIL-dren} \tabularnewline
28 & \en{Question after question, he never answered.} & \en{question} \pron{KWÈS-tchoun} \tabularnewline
29 & \en{She's a photographer and he's an engineer.} & accents mobiles \tabularnewline
30 & \en{It's not very interesting, to be honest.} & \en{honest} sans /h/ \tabularnewline
31 & \en{Half of the class was absent, which is unusual.} & \en{half} sans /l/ \tabularnewline
32 & \en{Could I have a knife and a fork, please?} & \en{knife} sans /k/ \tabularnewline
33 & \en{He wrote a wrong address on the envelope.} & \en{wrote, wrong} sans /w/ \tabularnewline
34 & \en{Climbing the mountain was hard but worth it.} & \en{climb} sans /b/, \en{worth} \tabularnewline
35 & \en{I'm looking forward to hearing from you.} & bloc social essentiel \tabularnewline
36 & \en{Would you mind opening the window?} & politesse + \en{would you} \tabularnewline
37 & \en{I'm not sure whether he'll come or not.} & \en{whether} \pron{WÈ-dheu} \tabularnewline
38 & \en{The problem is that we don't have enough time.} & \en{enough} \pron{i-NAF} \tabularnewline
39 & \en{She's used to working late, so it's fine.} & \en{used to} \pron{iou-s-tou} \tabularnewline
40 & \en{What would you have done in my place?} & conditionnel 3, \en{would you} \tabularnewline
41 & \en{I've been living here since I was a child.} & \en{since} + voyelles longues \tabularnewline
42 & \en{There's a lot of traffic at this hour.} & \en{hour} sans /h/, \en{a lot of} \tabularnewline
43 & \en{Although it's cold, I'm going out for a walk.} & \en{although} \pron{o-l-DZOOU} \tabularnewline
44 & \en{He asked me whether I had heard the news.} & \en{asked} \pron{aaskt}, discours rapporté \tabularnewline
45 & \en{You should have told me earlier, honestly.} & \en{should have} + \en{honestly} \tabularnewline
46 & \en{Everybody knows that nobody knows everything.} & \en{everybody} \pron{È-vri-ba-di} \tabularnewline
47 & \en{I'd like a glass of water without ice, please.} & \en{water}, \en{without} \tabularnewline
48 & \en{The government has decided to raise taxes.} & \en{government} \pron{GOU-veun-m'nt} \tabularnewline
49 & \en{It's probably going to rain this afternoon.} & \en{probably} \pron{PROB-li} \tabularnewline
50 & \en{Thank you very much for your help, I really appreciate it.} & remerciement complet \tabularnewline
\end{longtable}
\end{tableau}

\section{Programme de correction en 30 jours}
\begin{tableau}[30 jours pour réparer votre accent]
\begin{tabularx}{\linewidth}{@{}p{2.0cm}X@{}}
\entete{Jours} & \entete{Travail quotidien (10 à 15 minutes)} \tabularnewline
\hline
\textbf{1 à 5} & \ipatheta\ et \ipadh\ : 40 mots par jour devant un miroir + 3 virelangues. \tabularnewline
\textbf{6 à 9} & /h/ et /w/ : 40 mots + paires minimales \en{west/vest} + 3 virelangues. \tabularnewline
\textbf{10 à 13} & /r/ et /l/ : 40 mots + \en{red lorry yellow lorry} ×10 + lecture à voix haute de 5 phrases-drills. \tabularnewline
\textbf{14 à 17} & Voyelles : paires \en{sheep/ship}, \en{full/fool}, \en{cart/cut}, \en{cat/cut}, \en{no/now} + drill des 50 phrases (1 à 15). \tabularnewline
\textbf{18 à 21} & Finales consonantiques et \ipang\ : aucun « euh » ajouté + drill des 50 phrases (16 à 30). \tabularnewline
\textbf{22 à 25} & Rythme et accents : shadowing quotidien d'un dialogue de ce livre + marquage des accents au crayon. \tabularnewline
\textbf{26 à 28} & \en{Connected speech} : réductions (\en{gonna, wanna, gotta}) et formes faibles + drill des 50 phrases (31 à 50). \tabularnewline
\textbf{29 à 30} & \textbf{Bilan :} enregistrez 5 minutes de parole libre, comparez avec votre enregistrement du jour 1, listez vos 3 sons restants à travailler en continu. \tabularnewline
\end{tabularx}
\end{tableau}

\begin{memo}[Ce qu'il faut retenir de la Partie II]
\begin{enumerate}
  \item L'anglais a \textbf{9 sons inexistants en français} : \ipatheta, \ipadh, /h/, /w/,
  /r/, \ipang, \ipae, \ipalv, et le schwa. Travaillez-les un par un, dans cet ordre.
  \item \textbf{L'accent tonique} est plus important que les sons : \en{baNAna}, pas
  \en{banana}.
  \item Le \textbf{schwa} est la clé du rythme : écrasez les syllabes non accentuées.
  \item Le \textbf{\en{connected speech}} explique 80\,\% de vos difficultés d'écoute :
  les natifs ne prononcent pas les mots séparément.
  \item Choisissez \textbf{une} variété (américaine ou britannique) pour les 6 premiers mois.
  \item Enregistrez-vous \textbf{chaque semaine} : c'est le seul miroir honnête.
\end{enumerate}
\end{memo}


% ===== PARTIE III — GRAMMAIRE COMPLÈTE =============================
\part{Grammaire complète, du niveau A1 au niveau B2}
% =====================================================================
%  CHAPITRE 8 — GRAMMAIRE A1 : LES FONDATIONS
% =====================================================================
\chapter{Les fondations : pronoms, \en{be}, \en{have}, présent simple}

\begin{quote}\itshape
Vous n'avez pas besoin de connaître toute la grammaire pour parler. Vous avez besoin de
connaître \textbf{parfaitement} les dix structures qui représentent 70\,\% des phrases
réelles. Ce chapitre et le suivant vous les donnent.
\end{quote}

Chaque unité suit le même format : \textbf{règle} \ra \textbf{tableau} \ra \textbf{exemples
bilingues} \ra \textbf{piège francophone} \ra \textbf{exercices corrigés} \ra
\textbf{mission orale}. Ne sautez jamais la mission orale.

\section{Unité 1 --- Les pronoms personnels et le verbe \en{to be}}

\subsection{Les pronoms personnels}
\begin{tableau}[Pronoms personnels sujets et compléments]
\begin{tabularx}{\linewidth}{@{}p{3.2cm}p{3.2cm}p{3.0cm}X@{}}
\entete{Sujet} & \entete{Complément} & \entete{Possessif} & \entete{Exemple} \tabularnewline
\hline
\en{I} \pron{aï} --- je & \en{me} --- me, moi & \en{my / mine} & \en{I love you. She loves me.} \tabularnewline
\en{you} --- tu, vous & \en{you} --- te, toi, vous & \en{your / yours} & \en{You are late. I call you.} \tabularnewline
\en{he} --- il & \en{him} --- le, lui & \en{his / his} & \en{He is my brother. I see him.} \tabularnewline
\en{she} --- elle & \en{her} --- la, elle, lui & \en{her / hers} & \en{She is a doctor. I know her.} \tabularnewline
\en{it} --- il, elle, ça (choses/animaux) & \en{it} --- le, ça & \en{its / ---} & \en{It is cold. I like it.} \tabularnewline
\en{we} --- nous & \en{us} --- nous & \en{our / ours} & \en{We are ready. Help us.} \tabularnewline
\en{you} --- vous & \en{you} --- vous & \en{your / yours} & \en{You are my friends.} \tabularnewline
\en{they} --- ils, elles & \en{them} --- les, eux, elles & \en{their / theirs} & \en{They are here. I told them.} \tabularnewline
\end{tabularx}
\end{tableau}

\begin{piege}[L'anglais n'a pas de vouvoiement --- et c'est une bonne nouvelle]
\en{you} veut dire \textbf{tu} et \textbf{vous}, au singulier et au pluriel. Pas d'hésitation,
pas de faute possible. En revanche, on s'adresse presque toujours aux gens par leur prénom,
même au travail. Dire \en{Mr Diop} à un collègue peut paraître distant ; \en{Moussa} est normal.
\end{piege}

\begin{piege}[\en{it} : le pronom obligatoire]
Le français dit « il pleut », « c'est froid », « il est 8 heures ». L'anglais exige toujours
un sujet explicite : \en{it rains}, \en{it is cold}, \en{it is eight o'clock}.
\textbf{Ne dites jamais} \en{*Is cold} ou \en{*Rains}. C'est la faute n°1 des francophones
débutants.
\end{piege}

\subsection{Le verbe \en{to be} (être)}
\begin{tableau}[\en{to be} au présent --- avec contractions]
\begin{tabularx}{\linewidth}{@{}p{2.6cm}p{3.4cm}p{3.4cm}X@{}}
\entete{Forme} & \entete{Complète} & \entete{Contractée} & \entete{Traduction} \tabularnewline
\hline
Affirmation & \en{I am} & \en{I'm} \pron{aïm} & je suis \tabularnewline
 & \en{you are} & \en{you're} \pron{iou-ar} & tu es / vous êtes \tabularnewline
 & \en{he / she / it is} & \en{he's / she's / it's} & il / elle / c'est \tabularnewline
 & \en{we are} & \en{we're} \pron{ouii-ar} & nous sommes \tabularnewline
 & \en{they are} & \en{they're} \pron{dhè-ar} & ils / elles sont \tabularnewline
\hline
Négation & \en{I am not} & \en{I'm not} & je ne suis pas \tabularnewline
 & \en{you are not} & \en{you aren't} / \en{you're not} & tu n'es pas \tabularnewline
 & \en{he is not} & \en{he isn't} / \en{he's not} & il n'est pas \tabularnewline
\hline
Question & \en{Am I\ldots?} & --- & suis-je\ldots ? \tabularnewline
 & \en{Are you\ldots?} & --- & es-tu\ldots ? \tabularnewline
 & \en{Is he\ldots?} & --- & est-il\ldots ? \tabularnewline
\end{tabularx}
\end{tableau}

\begin{astuce}[À l'oral, utilisez toujours les formes contractées]
\en{I am French} sonne comme un robot ou comme une insistance (\textit{MOI, je suis français !}).
La forme naturelle est \en{I'm French.} En conversation, les contractées représentent 95\,\%
des cas. Les formes complètes s'utilisent en fin de phrase (\en{Yes, I am.}) et pour insister.
\end{astuce}

\subsection{Les six emplois de \en{be} que le français rend autrement}
\begin{tableau}[\en{be} là où le français utilise \en{avoir}]
\begin{tabularx}{\linewidth}{@{}p{6.6cm}X@{}}
\entete{Anglais} & \entete{Français} \tabularnewline
\hline
\en{I \textbf{am} 25 years old.} & J'\textbf{ai} 25 ans. \tabularnewline
\en{She \textbf{is} hungry / thirsty / cold.} & Elle \textbf{a} faim / soif / froid. \tabularnewline
\en{He \textbf{is} afraid of dogs.} & Il \textbf{a} peur des chiens. \tabularnewline
\en{We \textbf{are} late.} & Nous \textbf{sommes} en retard. (identique) \tabularnewline
\en{They \textbf{are} right.} & Ils \textbf{ont} raison. \tabularnewline
\en{How old \textbf{are} you?} & Quel âge \textbf{as}-tu ? \tabularnewline
\end{tabularx}
\end{tableau}

\begin{piege}[\en{*I have 25 years}]
Erreur systématique. Retenez le bloc entier : \en{I am X years old} --- ou, en conversation
courante, simplement \en{I'm 25}. De même : \en{*I have cold} \ra \en{I'm cold} ;
\en{*I have fear} \ra \en{I'm afraid}.
\end{piege}

\begin{exo}[1.1 --- Complétez avec \en{be}]
\begin{enumerate}
  \item I \trou[1.2cm] from Senegal. \item She \trou[1.2cm] my sister.
  \item They \trou[1.2cm] not at home. \item \trou[1.2cm] you married?
  \item We \trou[1.2cm] hungry. \item He \trou[1.2cm] 30 years old.
  \item It \trou[1.2cm] very hot today. \item Where \trou[1.2cm] your parents?
  \item I \trou[1.2cm] tired but I \trou[1.2cm] happy. \item \trou[1.2cm] she a doctor? --- No, she \trou[1.2cm].
\end{enumerate}
\end{exo}
\begin{corrige}[1.1]
1. \en{am} ('m) \quad 2. \en{is} ('s) \quad 3. \en{are} \quad 4. \en{Are} \quad
5. \en{are} \quad 6. \en{is} \quad 7. \en{is} \quad 8. \en{are} \quad
9. \en{am} / \en{am} \quad 10. \en{Is} --- \en{she isn't}
\end{corrige}

\begin{exo}[1.2 --- Traduisez en anglais]
\begin{enumerate}
  \item J'ai 32 ans. \item Elle a faim. \item Nous avons raison.
  \item Il a peur de la mer. \item Ils ne sont pas prêts. \item Quel âge as-tu ?
  \item C'est mon ami. \item Je ne suis pas fatigué.
\end{enumerate}
\end{exo}
\begin{corrige}[1.2]
1. \en{I'm 32 (years old).} 2. \en{She's hungry.} 3. \en{We're right.}
4. \en{He's afraid of the sea.} 5. \en{They aren't ready.} 6. \en{How old are you?}
7. \en{He's / This is my friend.} 8. \en{I'm not tired.}
\end{corrige}

\begin{mission}[Unité 1]
Dites à voix haute, en vous regardant dans un miroir, 10 fois de suite, en variant le
pronom : \en{I'm happy --- you're happy --- he's happy --- she's happy --- it's good ---
we're happy --- they're happy}. Puis présentez-vous en 5 phrases :
\en{Hello, my name is\ldots\ I'm from\ldots\ I'm \ldots years old. I'm a \ldots\ I'm not
married / I'm married.}
\end{mission}


\section{Unité 2 --- \en{have}, \en{have got} et la possession}

\begin{regle}[\en{have} = avoir]
\en{I / you / we / they have} \quad---\quad \en{he / she / it has}\\
Contractions : \en{I've, you've, we've, they've, he's, she's, it's}.\\
Négation : \en{I don't have} / \en{he doesn't have} (américain courant) \quad ou \quad
\en{I haven't got} / \en{he hasn't got} (britannique).\\
Question : \en{Do you have\ldots?} / \en{Have you got\ldots?}
\end{regle}

\begin{tableau}[\en{have} ou \en{have got} ?]
\begin{tabularx}{\linewidth}{@{}p{4.0cm}p{5.2cm}X@{}}
\entete{} & \entete{Américain (recommandé)} & \entete{Britannique} \tabularnewline
\hline
J'ai une voiture. & \en{I have a car.} / \en{I've got a car.} & \en{I've got a car.} \tabularnewline
Tu as un stylo ? & \en{Do you have a pen?} & \en{Have you got a pen?} \tabularnewline
Il n'a pas d'argent. & \en{He doesn't have any money.} & \en{He hasn't got any money.} \tabularnewline
\end{tabularx}
\end{tableau}

\begin{piege}[\en{He doesn't \textbf{has}}]
Après \en{does / doesn't / did / didn't}, le verbe revient à sa forme \textbf{de base} :
\en{He doesn't have} (et non \en{*doesn't has}). Même chose aux questions :
\en{Does she have a car?} (et non \en{*Does she has}).
\end{piege}

\subsection{Adjectifs et pronoms possessifs}
\begin{tableau}[Les possessifs]
\begin{tabularx}{\linewidth}{@{}p{2.6cm}p{3.4cm}p{3.4cm}X@{}}
\entete{Personne} & \entete{Adjectif + nom} & \entete{Pronom (seul)} & \entete{Exemple} \tabularnewline
\hline
je & \en{my} --- mon & \en{mine} --- le mien & \en{my book / the book is mine} \tabularnewline
tu / vous & \en{your} --- ton, votre & \en{yours} & \en{your name / it's yours} \tabularnewline
il & \en{his} --- son (à lui) & \en{his} & \en{his car} \tabularnewline
elle & \en{her} --- son (à elle) & \en{hers} & \en{her bag} \tabularnewline
il/elle (chose) & \en{its} --- son (à cela) & --- & \en{its colour} \tabularnewline
nous & \en{our} --- notre & \en{ours} & \en{our house} \tabularnewline
ils/elles & \en{their} --- leur & \en{theirs} & \en{their children} \tabularnewline
\end{tabularx}
\end{tableau}

\begin{memo}[Le possessif suit le \textbf{possesseur}, pas l'objet]
En français, « son » dépend du genre de l'objet possédé (\emph{son} livre, \emph{sa} voiture).
En anglais, il dépend de \textbf{qui possède} : \en{his book} (le livre d'un homme),
\en{her book} (le livre d'une femme) --- même si c'est le même objet. Et \en{his} ou \en{her}
ne changent jamais au pluriel : \en{his books}, \en{her cars}.
\end{memo}

\subsection{Le génitif : \en{'s} (appartenance)}
\begin{regle}
Nom + \en{'s} = « de » : \en{Moussa's car} (la voiture de Moussa), \en{my brother's wife}
(la femme de mon frère), \en{the children's room}. Au pluriel en \en{-s}, on met seulement
l'apostrophe : \en{the students' books}. Pour les choses, on préfère \en{of} :
\en{the door of the car}, \en{the end of the film}.
\end{regle}

\begin{exo}[2.1 --- Traduisez]
\begin{enumerate}
  \item J'ai deux frères. \item Elle a une belle maison. \item Nous n'avons pas de temps.
  \item Est-ce qu'il a une voiture ? \item Le téléphone de Fatou. \item Le sac des enfants.
  \item C'est le mien. \item Leur maison est grande.
\end{enumerate}
\end{exo}
\begin{corrige}[2.1]
1. \en{I have two brothers.} 2. \en{She has a beautiful house.}
3. \en{We don't have (any) time.} 4. \en{Does he have a car?} / \en{Has he got a car?}
5. \en{Fatou's phone.} 6. \en{The children's bag.} 7. \en{It's mine.}
8. \en{Their house is big.}
\end{corrige}

\section{Unité 3 --- Le présent simple}

\begin{regle}[Formation]
\textbf{Toutes les personnes : base verbale}, sauf la 3\textsuperscript{e} personne du
singulier qui prend \en{-s}.\\
\en{I work, you work, he work\textbf{s}, she work\textbf{s}, it work\textbf{s}, we work,
they work.}\\
Négation : \en{do not (don't)} + base --- \en{I don't work}, \en{he doesn't work}.\\
Question : \en{do / does} + sujet + base --- \en{Do you work?}, \en{Does she work?}
\end{regle}

\begin{tableau}[Le \en{-s} de la 3\textsuperscript{e} personne : orthographe]
\begin{tabularx}{\linewidth}{@{}p{5.6cm}X@{}}
\entete{Cas} & \entete{Règle et exemples} \tabularnewline
\hline
Règle générale & + \en{s} : \en{works, plays, reads, eats} \tabularnewline
Verbes en \en{-o}, \en{-ch}, \en{-sh}, \en{-ss}, \en{-x} & + \en{es} : \en{goes, watches, washes, passes, fixes} \tabularnewline
Consonne + \en{y} & \en{y} \ra \en{ies} : \en{study} \ra \en{studies}, \en{try} \ra \en{tries}, \en{cry} \ra \en{cries} \tabularnewline
Voyelle + \en{y} & + \en{s} simplement : \en{plays, buys, says, enjoys} \tabularnewline
Irréguliers & \en{have} \ra \en{has} \tabularnewline
\end{tabularx}
\end{tableau}

\subsection{Quand employer le présent simple ?}
\begin{enumerate}
  \item \textbf{Habitudes et routines} : \en{I wake up at six every day.}
  \item \textbf{Vérités générales} : \en{Water boils at 100 degrees.} \en{The sun rises in the east.}
  \item \textbf{États permanents} : \en{She lives in Thiès.} \en{He works in a bank.}
  \item \textbf{Horaires officiels} (futur) : \en{The train leaves at 8 pm.}
  \item \textbf{Avec les verbes d'état} (voir, entendre, aimer, savoir\ldots) :
  \en{I know, I understand, I want, I need, I like, I love, I prefer, I believe, I remember}.
\end{enumerate}

\begin{piege}[Les verbes d'état ne se mettent pas au \en{-ing}]
\en{*I am knowing the answer} \ra \en{I know the answer}.
\en{*I am wanting a coffee} \ra \en{I want a coffee}.
Liste à connaître : \en{know, want, need, like, love, hate, prefer, believe, understand,
remember, forget, mean, belong, own, seem, sound, smell, taste, feel (sens), hope, wish}.
\end{piege}

\subsection{Les adverbes de fréquence et leur place}
\begin{tableau}[Échelle de fréquence + placement]
\begin{tabularx}{\linewidth}{@{}p{3.4cm}p{2.0cm}X@{}}
\entete{Adverbe} & \entete{\%} & \entete{Exemple (placement)} \tabularnewline
\hline
\en{always} & 100 & \en{I \textbf{always} drink tea.} \tabularnewline
\en{usually} & 90 & \en{She \textbf{usually} takes the bus.} \tabularnewline
\en{often} & 70 & \en{We \textbf{often} watch football.} \tabularnewline
\en{sometimes} & 50 & \en{They \textbf{sometimes} cook fish.} (aussi en début ou fin de phrase) \tabularnewline
\en{rarely / seldom} & 10 & \en{He \textbf{rarely} goes out.} \tabularnewline
\en{never} & 0 & \en{I \textbf{never} eat pork.} \tabularnewline
\en{hardly ever} & 2 & \en{We \textbf{hardly ever} travel.} \tabularnewline
\end{tabularx}
\end{tableau}

\begin{regle}[Placement]
\textbf{Avant} le verbe principal : \en{I \textbf{never} smoke.}\\
\textbf{Après} le verbe \en{be} : \en{He is \textbf{always} late.}\\
\textbf{Entre} l'auxiliaire et le verbe : \en{I have \textbf{never} been to London.}
\en{She doesn't \textbf{often} call.}
\end{regle}

\begin{piege}[Jamais de double négation]
Le français dit « je ne sais \emph{pas} » ; le wolof et d'autres langues autorisent la double
négation. L'anglais \textbf{non} : \en{never} et \en{not} ne cohabitent jamais.\\
\en{*I don't never go there.} \ra \en{I never go there.} ou \en{I don't ever go there.}
\end{piege}

\begin{exo}[3.1 --- Mettez le verbe à la bonne forme]
\begin{enumerate}
  \item She (work) \trou[2cm] in a hospital.
  \item They (not / like) \trou[2cm] coffee.
  \item (you / speak) \trou[2.5cm] English?
  \item My brother (study) \trou[2cm] every evening.
  \item He (go) \trou[2cm] to the mosque on Friday.
  \item We (watch) \trou[2cm] the news at 8 pm.
  \item (she / have) \trou[2.5cm] a car?
  \item I (not / understand) \trou[2.5cm] this word.
  \item The bus (leave) \trou[2cm] at 7 am.
  \item My parents (live) \trou[2cm] in Dakar.
\end{enumerate}
\end{exo}
\begin{corrige}[3.1]
1. \en{works} 2. \en{don't like} 3. \en{Do you speak} 4. \en{studies} 5. \en{goes}
6. \en{watch} 7. \en{Does she have} 8. \en{don't understand} 9. \en{leaves} 10. \en{live}
\end{corrige}

\begin{exo}[3.2 --- Placez l'adverbe de fréquence]
\begin{enumerate}
  \item I drink tea. (always) \item She is late. (never)
  \item We go to the beach. (sometimes) \item He watches TV. (hardly ever)
  \item They arrive on time. (usually) \item I have eaten sushi. (never)
\end{enumerate}
\end{exo}
\begin{corrige}[3.2]
1. \en{I always drink tea.} 2. \en{She is never late.} 3. \en{We sometimes go to the beach.}
4. \en{He hardly ever watches TV.} 5. \en{They usually arrive on time.}
6. \en{I have never eaten sushi.}
\end{corrige}

\begin{mission}[Unité 3 --- votre routine, 3 minutes]
Décrivez votre journée habituelle en 12 phrases au présent simple, à voix haute, avec des
adverbes de fréquence : \en{I always wake up at\ldots\ I usually drink\ldots\ I never\ldots\
Sometimes I\ldots} Enregistrez-vous. Réécoutez le lendemain et corrigez vos \en{-s} oubliés
--- c'est l'erreur la plus fréquente à l'oral.
\end{mission}


\section{Unité 4 --- Les articles : \en{a}, \en{an}, \en{the}, ou rien}

C'est \textbf{le} point où les francophones font des fautes à vie. La raison : le français
met un article presque partout, l'anglais souvent pas.

\begin{regle}[Les trois cas]
\begin{enumerate}
  \item \en{a / an} = \textbf{un, une} (singulier, non précis, première mention).
  \en{a} devant consonne \emph{sonore}, \en{an} devant voyelle \emph{sonore} :
  \en{a book, an apple, a university} (\pron{iou} = consonne), \en{an hour} (\en{h} muet).
  \item \en{the} = \textbf{le, la, les} (chose précise, déjà mentionnée, unique).
  \item \textbf{Aucun article} (article « zéro ») : généralités, pluriels et indénombrables
  pris en général, noms de pays (sauf exceptions), repas, langues, matières, sports.
\end{enumerate}
\end{regle}

\begin{tableau}[Article zéro : les 10 cas à connaître]
\begin{tabularx}{\linewidth}{@{}p{5.4cm}p{5.0cm}X@{}}
\entete{Règle} & \entete{Anglais} & \entete{Français (avec article !)} \tabularnewline
\hline
Généralité avec un pluriel & \en{Dogs are friendly.} & \textbf{Les} chiens sont gentils. \tabularnewline
Généralité avec un indénombrable & \en{Water is essential.} & \textbf{L'}eau est essentielle. \tabularnewline
Repas & \en{I had breakfast at 7.} & J'ai pris \textbf{le} petit-déjeuner. \tabularnewline
Langues & \en{She speaks English.} & Elle parle \textbf{l'}anglais. \tabularnewline
Matières & \en{This table is wood.} / \en{made of wood} & Cette table est \textbf{en} bois. \tabularnewline
Sports et jeux & \en{I play football.} & Je joue \textbf{au} football. \tabularnewline
Noms de pays (singulier sans \en{States/Kingdom/Republic}) & \en{I live in Senegal.} & Je vis \textbf{au} Sénégal. \tabularnewline
Villes & \en{She lives in Thiès.} & Elle vit \textbf{à} Thiès. \tabularnewline
Années, mois, jours & \en{in 2026, on Monday, in May} & \textbf{en} 2026, \textbf{le} lundi \tabularnewline
Institutions au sens général & \en{go to school, go to bed, go to church, go to work, go home} & aller à \textbf{l'}école, se coucher \tabularnewline
\end{tabularx}
\end{tableau}

\begin{tableau}[Les exceptions qui contredisent la règle des pays]
\begin{tabularx}{\linewidth}{@{}p{6.4cm}X@{}}
\entete{Avec \en{the}} & \entete{Exemples} \tabularnewline
\hline
Pays au pluriel ou contenant \en{States, Kingdom, Republic, Union} &
\en{the United States, the United Kingdom, the Netherlands, the Philippines,
the Democratic Republic of Congo, the European Union} \tabularnewline
Fleuves, mers, océans, déserts, chaînes de montagnes &
\en{the Senegal River, the Nile, the Atlantic, the Sahara, the Alps} \tabularnewline
Journaux, hôtels, musées, cinémas & \en{the Times, the Hilton, the Louvre} \tabularnewline
Adjectif + nationalité au pluriel & \en{the French, the British, the Senegalese} \tabularnewline
\end{tabularx}
\end{tableau}

\begin{piege}[\en{I like the music} ou \en{I like music} ?]
\en{I like music.} = j'aime la musique en général (le bon choix 9 fois sur 10).\\
\en{I like the music} = j'aime \emph{cette} musique-là (celle qu'on entend maintenant).\\
Même logique : \en{Life is beautiful} (la vie en général) / \en{The life of Mandela}
(la vie d'une personne précise). Les francophones ajoutent \en{the} par réflexe : c'est
la faute la plus tenace de toute la langue.
\end{piege}

\begin{exo}[4.1 --- \en{a}, \en{an}, \en{the} ou rien ?]
\begin{enumerate}
  \item I bought \trou[1cm] new phone yesterday. \trou[1cm] phone is blue.
  \item She is \trou[1cm] engineer. She works in \trou[1cm] Dakar.
  \item \trou[1cm] life is short.
  \item We had \trou[1cm] lunch at \trou[1cm] one o'clock.
  \item He speaks \trou[1cm] French and \trou[1cm] English.
  \item \trou[1cm] sun rises in \trou[1cm] east.
  \item I go to \trou[1cm] work by \trou[1cm] bus.
  \item \trou[1cm] United States is a big country.
  \item This is \trou[1cm] hour of my life I will never forget.
  \item She studies at \trou[1cm] university in \trou[1cm] Senegal.
\end{enumerate}
\end{exo}
\begin{corrige}[4.1]
1. \en{a} --- \en{The} \quad 2. \en{an} --- rien \quad 3. rien \quad 4. rien --- rien \quad
5. rien --- rien \quad 6. \en{The} --- \en{the} \quad 7. rien --- rien \quad
8. \en{The} \quad 9. \en{an} \quad 10. rien (ou \en{a} si on précise laquelle) --- rien
\end{corrige}

\section{Unité 5 --- Le pluriel, les dénombrables et \en{some / any}}

\begin{tableau}[Formation du pluriel]
\begin{tabularx}{\linewidth}{@{}p{5.0cm}X@{}}
\entete{Cas} & \entete{Exemples} \tabularnewline
\hline
Règle générale : + \en{s} & \en{book} \ra \en{books}, \en{car} \ra \en{cars} \tabularnewline
\en{-s, -sh, -ch, -x, -o} : + \en{es} & \en{bus} \ra \en{buses}, \en{box} \ra \en{boxes}, \en{watch} \ra \en{watches}, \en{tomato} \ra \en{tomatoes}, \en{potato} \ra \en{potatoes} (mais \en{photo} \ra \en{photos}, \en{video} \ra \en{videos}) \tabularnewline
Consonne + \en{y} : \en{y} \ra \en{ies} & \en{city} \ra \en{cities}, \en{baby} \ra \en{babies}, \en{country} \ra \en{countries} \tabularnewline
\en{-f / -fe} : \ra \en{-ves} & \en{leaf} \ra \en{leaves}, \en{wife} \ra \en{wives}, \en{knife} \ra \en{knives}, \en{life} \ra \en{lives}, \en{shelf} \ra \en{shelves} (mais \en{roof} \ra \en{roofs}) \tabularnewline
\textbf{Irréguliers} & \en{man} \ra \en{men} ; \en{woman} \ra \en{women} ; \en{child} \ra \en{children} ; \en{person} \ra \en{people} ; \en{tooth} \ra \en{teeth} ; \en{foot} \ra \en{feet} ; \en{mouse} \ra \en{mice} ; \en{goose} \ra \en{geese} \tabularnewline
\textbf{Invariables} & \en{sheep, fish, deer, series, species, aircraft} \tabularnewline
\textbf{Toujours pluriel} & \en{trousers, jeans, glasses, scissors, clothes, police, people} \tabularnewline
\end{tabularx}
\end{tableau}

\subsection{Dénombrables et indénombrables}
\begin{tableau}[Les indénombrables les plus utiles]
\begin{tabularx}{\linewidth}{@{}p{4.4cm}p{4.4cm}X@{}}
\entete{Alimentation} & \entete{Abstraits} & \entete{Matières / autres} \tabularnewline
\hline
\en{bread, rice, water, milk, sugar, salt, meat, fish, fruit, butter, cheese, tea, coffee, oil} &
\en{information, advice, news, money, work, time, help, knowledge, progress, research, furniture, luggage, equipment, traffic, homework} &
\en{paper, wood, gold, iron, hair, grass, electricity, weather, rain, snow, music, English, sleep} \tabularnewline
\end{tabularx}
\end{tableau}

\begin{piege}[Ces mots sont indénombrables en anglais et dénombrables en français]
\en{*an information} \ra \en{a piece of information}\\
\en{*an advice} \ra \en{a piece of advice}\\
\en{*a news} \ra \en{a piece of news} (\en{news} est singulier : \en{The news \textbf{is} bad.})\\
\en{*a furniture} \ra \en{a piece of furniture}\\
\en{*a luggage} \ra \en{a piece of luggage / a bag}\\
\en{*a hair} (= la chevelure) \ra \en{a hair} existe mais veut dire \emph{un poil} !\\
\en{*two breads} \ra \en{two loaves of bread} --- \en{*a money} \ra \en{some money}
\end{piege}

\subsection{\en{some}, \en{any}, \en{no}, \en{much}, \en{many}, \en{a lot of}}
\begin{tableau}[Quantifieurs de base]
\begin{tabularx}{\linewidth}{@{}p{2.8cm}p{4.0cm}X@{}}
\entete{Mot} & \entete{Emploi} & \entete{Exemples} \tabularnewline
\hline
\en{some} & Phrase affirmative ; offres et demandes polies & \en{I need some water.} / \en{Would you like some tea?} / \en{Can I have some bread?} \tabularnewline
\en{any} & Négation et question & \en{I don't have any money.} / \en{Do you have any brothers?} \tabularnewline
\en{no} & Négation sans \en{not} & \en{I have no time.} (= \en{I don't have any time.}) \tabularnewline
\en{much} & Indénombrables (négation, question) & \en{How much sugar?} / \en{I don't have much time.} \tabularnewline
\en{many} & Dénombrables & \en{How many children?} / \en{There aren't many cars.} \tabularnewline
\en{a lot of / lots of} & Les deux, à l'affirmative (le plus naturel) & \en{I have a lot of friends.} / \en{She drinks a lot of water.} \tabularnewline
\en{a few / few} & Dénombrables : quelques / très peu & \en{a few friends} (positif) / \en{few friends} (presque aucun) \tabularnewline
\en{a little / little} & Indénombrables : un peu / très peu & \en{a little sugar} / \en{little hope} \tabularnewline
\end{tabularx}
\end{tableau}

\begin{exo}[5.1 --- Pluriels]
Écrivez le pluriel : \en{child, woman, foot, tooth, person, mouse, bus, city, wife, leaf,
sheep, potato, box, baby, man, fish, glass, country, knife, life}.
\end{exo}
\begin{corrige}[5.1]
\en{children, women, feet, teeth, people, mice, buses, cities, wives, leaves, sheep,
potatoes, boxes, babies, men, fish, glasses, countries, knives, lives}.
\end{corrige}

\begin{exo}[5.2 --- \en{some}, \en{any}, \en{much}, \en{many}, \en{a lot of}]
\begin{enumerate}
  \item I don't have \trou[1.5cm] time.
  \item How \trou[1.5cm] brothers do you have?
  \item There is \trou[1.5cm] milk in the fridge.
  \item Do you want \trou[1.5cm] tea?
  \item She has \trou[1.5cm] of friends.
  \item How \trou[1.5cm] money do you need?
  \item We don't have \trou[1.5cm] bread left.
  \item There are \trou[1.5cm] people in the room.
\end{enumerate}
\end{exo}
\begin{corrige}[5.2]
1. \en{any / much} 2. \en{many} 3. \en{some} 4. \en{some} (offre) 5. \en{a lot / lots}
6. \en{much} 7. \en{any / much} 8. \en{a lot of / many}
\end{corrige}

\begin{exo}[5.3 --- Corrigez l'erreur]
\begin{enumerate}
  \item \en{*Can you give me an advice?}
  \item \en{*I have many informations for you.}
  \item \en{*The news are good today.}
  \item \en{*I bought two breads and a milk.}
  \item \en{*She has a beautiful hairs.}
  \item \en{*My luggages are heavy.}
\end{enumerate}
\end{exo}
\begin{corrige}[5.3]
1. \en{Can you give me a piece of advice?} (ou \en{some advice})
2. \en{I have a lot of information for you.}
3. \en{The news is good today.}
4. \en{I bought two loaves of bread and some milk.}
5. \en{She has beautiful hair.}
6. \en{My luggage is heavy.} / \en{My bags are heavy.}
\end{corrige}


\section{Unité 6 --- Poser des questions : la mécanique complète}

\begin{regle}[L'ordre des mots : AUX + SUJET + VERBE]
L'anglais n'a pas d'inversion sujet-verbe comme le français. On place un \textbf{auxiliaire}
devant le sujet :\\
\en{You work here.} \ra \en{\textbf{Do} you work here?}\\
\en{She is tired.} \ra \en{\textbf{Is} she tired?}\\
\en{They can swim.} \ra \en{\textbf{Can} they swim?}\\
\en{He went out.} \ra \en{\textbf{Did} he go out?} (le verbe redevient base !)
\end{regle}

\begin{tableau}[Les mots interrogatifs (\en{Wh-} words)]
\begin{tabularx}{\linewidth}{@{}p{2.8cm}p{3.4cm}X@{}}
\entete{Mot} & \entete{Sens} & \entete{Exemple} \tabularnewline
\hline
\en{What} & que, quoi, quel & \en{What do you want?} --- \textit{Que veux-tu ?} \tabularnewline
\en{What\ldots like} & comment est (description) & \en{What's he like?} --- \textit{Comment est-il ? (caractère)} \tabularnewline
\en{Where} & où & \en{Where do you live?} \tabularnewline
\en{When} & quand & \en{When does the film start?} \tabularnewline
\en{Why} & pourquoi & \en{Why are you late?} --- réponse par \en{Because\ldots} \tabularnewline
\en{Who} & qui & \en{Who is that man?} / \en{Who called?} \tabularnewline
\en{Whom} & qui (complément, formel) & \en{To whom did you speak?} \tabularnewline
\en{Whose} & à qui & \en{Whose bag is this?} \tabularnewline
\en{Which} & lequel, quel (choix limité) & \en{Which colour do you prefer?} \tabularnewline
\en{How} & comment & \en{How does it work?} \tabularnewline
\en{How much} & combien (indénombrable) & \en{How much is it?} \tabularnewline
\en{How many} & combien (dénombrable) & \en{How many people came?} \tabularnewline
\en{How long} & combien de temps & \en{How long does it take?} \tabularnewline
\en{How often} & à quelle fréquence & \en{How often do you study?} \tabularnewline
\en{How far} & à quelle distance & \en{How far is the airport?} \tabularnewline
\en{How old} & quel âge & \en{How old are you?} \tabularnewline
\en{What time} & à quelle heure & \en{What time do you get up?} \tabularnewline
\end{tabularx}
\end{tableau}

\begin{piege}[Les questions \emph{sujet} : pas d'auxiliaire !]
Quand le mot interrogatif \textbf{est} le sujet, l'ordre reste normal et le verbe prend
\en{-s} :\\
\en{Who wants coffee?} (et non \en{*Who does want coffee?})\\
\en{What happened?} (et non \en{*What did happen?})\\
\en{Who called you?} \quad \en{Which bus goes to the centre?}
\end{piege}

\subsection{Les questions indirectes (la politesse)}
\begin{regle}
Devant \en{Can you tell me\ldots}, \en{Do you know\ldots}, \en{I wonder\ldots},
\en{Could you tell me\ldots}, l'ordre redevient \textbf{affirmatif} : sujet puis verbe,
et \en{do/does/did} disparaît.
\end{regle}
\ex{Where is the station? \ra\ Could you tell me where the station is?}{Où est la gare ? \ra\ Pourriez-vous me dire où est la gare ?}
\ex{What time does it start? \ra\ Do you know what time it starts?}{À quelle heure ça commence ? \ra\ Savez-vous à quelle heure ça commence ?}
\ex{Why did he leave? \ra\ I wonder why he left.}{Pourquoi est-il parti ? \ra\ Je me demande pourquoi il est parti.}

\begin{astuce}[Les formules qui font « anglais naturel »]
\en{Do you know if\ldots?} (Savez-vous si\ldots ?) \quad \en{I was wondering whether\ldots}
(Je me demandais si\ldots) \quad \en{Sorry to bother you, but\ldots} (Pardon de vous
déranger, mais\ldots) \quad \en{Would you mind telling me\ldots?} (Cela vous ennuierait-il
de me dire\ldots ?). Ces quatre formules transforment instantanément un anglais scolaire en
anglais poli d'adulte.
\end{astuce}

\subsection{Les 30 questions à connaître par c\oe ur}
\begin{tableau}[30 questions essentielles]
\footnotesize
\begin{longtable}{@{}p{6.6cm}p{7.0cm}@{}}
\entete{Anglais} & \entete{Français} \tabularnewline
\endfirsthead
\entete{Anglais} & \entete{Français} \tabularnewline
\endhead
\phr{What's your name?}{Comment vous appelez-vous ?}
\phr{Where are you from?}{D'où venez-vous ?}
\phr{Where do you live?}{Où habitez-vous ?}
\phr{How old are you?}{Quel âge avez-vous ?}
\phr{What do you do? (for a living)}{Que faites-vous dans la vie ?}
\phr{Are you married?}{Êtes-vous marié(e) ?}
\phr{Do you have children?}{Avez-vous des enfants ?}
\phr{Do you speak English?}{Parlez-vous anglais ?}
\phr{Can you repeat that, please?}{Pouvez-vous répéter, s'il vous plaît ?}
\phr{Can you speak more slowly, please?}{Pouvez-vous parler plus lentement ?}
\phr{What does ``\ldots'' mean?}{Que veut dire « \ldots » ?}
\phr{How do you say ``\ldots'' in English?}{Comment dit-on « \ldots » en anglais ?}
\phr{How do you spell that?}{Comment cela s'épelle-t-il ?}
\phr{What time is it?}{Quelle heure est-il ?}
\phr{How much is it?}{Combien ça coûte ?}
\phr{Where is the bathroom?}{Où sont les toilettes ?}
\phr{How do I get to\ldots?}{Comment aller à\ldots ?}
\phr{Is it far from here?}{C'est loin d'ici ?}
\phr{What are you doing?}{Que fais-tu ?}
\phr{What did you do yesterday?}{Qu'as-tu fait hier ?}
\phr{Where did you go?}{Où es-tu allé ?}
\phr{Why not?}{Pourquoi pas ?}
\phr{What's wrong?}{Qu'est-ce qui ne va pas ?}
\phr{What happened?}{Que s'est-il passé ?}
\phr{How was your day?}{Comment s'est passée ta journée ?}
\phr{Are you hungry?}{As-tu faim ?}
\phr{Do you mind if I\ldots?}{Cela t'ennuie si je\ldots ?}
\phr{What do you think?}{Qu'en pensez-vous ?}
\phr{Could you help me, please?}{Pourriez-vous m'aider ?}
\phr{When will you be ready?}{Quand serez-vous prêt ?}
\end{longtable}
\end{tableau}

\begin{exo}[6.1 --- Transformez en question]
\begin{enumerate}
  \item \en{You work in a bank.} \item \en{She lives in Dakar.}
  \item \en{They arrived yesterday.} \item \en{He can drive.}
  \item \en{You are thirsty.} \item \en{She wants a coffee.}
  \item \en{Someone called.} (avec \en{who}) \item \en{You need help.} (avec \en{what})
\end{enumerate}
\end{exo}
\begin{corrige}[6.1]
1. \en{Do you work in a bank?} 2. \en{Does she live in Dakar?}
3. \en{Did they arrive yesterday?} 4. \en{Can he drive?} 5. \en{Are you thirsty?}
6. \en{Does she want a coffee?} 7. \en{Who called?} 8. \en{What do you need?}
\end{corrige}

\begin{exo}[6.2 --- Questions indirectes]
Rendez ces questions polies :
\begin{enumerate}
  \item \en{Where is the bus stop?} (avec \en{Can you tell me\ldots})
  \item \en{What time does the shop close?} (avec \en{Do you know\ldots})
  \item \en{Why is he angry?} (avec \en{I wonder\ldots})
  \item \en{How much does it cost?} (avec \en{Could you tell me\ldots})
\end{enumerate}
\end{exo}
\begin{corrige}[6.2]
1. \en{Can you tell me where the bus stop is?} 2. \en{Do you know what time the shop closes?}
3. \en{I wonder why he is angry.} 4. \en{Could you tell me how much it costs?}
\end{corrige}

\section{Unité 7 --- La négation}

\begin{tableau}[Toutes les façons de nier]
\begin{tabularx}{\linewidth}{@{}p{3.6cm}p{4.6cm}X@{}}
\entete{Type} & \entete{Construction} & \entete{Exemples} \tabularnewline
\hline
Avec \en{be} & sujet + \en{be} + \en{not} & \en{I'm not tired. He isn't here.} \tabularnewline
Avec un auxiliaire/modal & auxiliaire + \en{not} & \en{I can't swim. She won't come. We haven't eaten.} \tabularnewline
Avec un verbe simple & \en{don't / doesn't / didn't} + base & \en{I don't know. He doesn't work. They didn't come.} \tabularnewline
Avec un adverbe négatif & \en{never, hardly, rarely, nothing, nobody} & \en{I never lie. Nobody called. Nothing happened.} \tabularnewline
Réponses courtes & \en{No, I don't. / No, she can't.} & \en{Do you like it?} --- \en{No, I don't.} \tabularnewline
\end{tabularx}
\end{tableau}

\begin{memo}[Les mots négatifs suffisent seuls]
\en{nothing} = \en{not anything} ; \en{nobody / no one} = \en{not anybody} ;
\en{nowhere} = \en{not anywhere} ; \en{never} = \en{not ever}.\\
Un seul par phrase : \en{I said nothing.} \textbf{ou} \en{I didn't say anything.}
Jamais les deux (\en{*I didn't say nothing}).
\end{memo}

\begin{piege}[Les réponses \en{yes} / \en{no} après une question négative]
\en{Don't you like coffee?} --- Si vous aimez le café, répondez \en{Yes, I do!} (= « si ! »).
Le français dit « si », l'anglais dit \en{yes}. Ne traduisez pas « non » littéralement :
répondez toujours selon la \textbf{réalité}, pas selon la forme de la question.
\end{piege}

\section{Unité 8 --- \en{There is / There are} et les prépositions de lieu}

\begin{regle}
\en{There is} + singulier / indénombrable \quad---\quad \en{There are} + pluriel\\
Négation : \en{There isn't / There aren't} \quad---\quad Question : \en{Is there\ldots? / Are there\ldots?}\\
Contracté : \en{There's a problem.} (\en{there're} s'écrit rarement)
\end{regle}

\begin{tableau}[Les prépositions de lieu essentielles]
\begin{tabularx}{\linewidth}{@{}p{2.6cm}p{3.6cm}X@{}}
\entete{Préposition} & \entete{Sens} & \entete{Exemple} \tabularnewline
\hline
\en{in} & dans (espace clos, ville, pays) & \en{in the box, in Dakar, in Senegal, in the car} \tabularnewline
\en{on} & sur (surface, étage, transport collectif) & \en{on the table, on the wall, on the second floor, on the bus} \tabularnewline
\en{at} & à (point précis, adresse, lieu d'activité) & \en{at the door, at home, at work, at school, at 25 Independence Avenue} \tabularnewline
\en{under} & sous & \en{under the bed} \tabularnewline
\en{next to / beside} & à côté de & \en{next to the bank} \tabularnewline
\en{between} & entre & \en{between the shop and the mosque} \tabularnewline
\en{in front of} & devant & \en{in front of the house} \tabularnewline
\en{behind} & derrière & \en{behind the car} \tabularnewline
\en{opposite} & en face de & \en{opposite the market} \tabularnewline
\en{near / close to} & près de & \en{near the beach} \tabularnewline
\en{above / below} & au-dessus / au-dessous (sans contact) & \en{above the door} \tabularnewline
\en{into / out of} & dans (mouvement) / hors de & \en{He went into the room.} \tabularnewline
\en{through} & à travers & \en{through the window, through the city} \tabularnewline
\en{across} & de l'autre côté, en traversant & \en{across the street} \tabularnewline
\en{along} & le long de & \en{along the river} \tabularnewline
\en{around} & autour de & \en{around the table} \tabularnewline
\end{tabularx}
\end{tableau}

\begin{astuce}[\en{in the car} mais \en{on the bus}]
On dit \en{in} pour les petits véhicules où l'on s'assoit (\en{in a car, in a taxi,
in a truck, in a helicopter}) et \en{on} pour ceux où l'on \emph{marche}
(\en{on a bus, on a train, on a plane, on a ship}) --- et aussi \en{on a bike, on a motorbike,
on a horse}.
\end{astuce}

\subsection{Les démonstratifs}
\begin{tableau}[\en{this / that / these / those}]
\begin{tabularx}{\linewidth}{@{}p{3.0cm}p{3.4cm}X@{}}
\entete{} & \entete{Près de moi} & \entete{Loin de moi} \tabularnewline
\hline
Singulier & \en{this} \pron{dhis} --- ceci, ce & \en{that} \pron{dhat} --- cela, ce \tabularnewline
Pluriel & \en{these} \pron{dhiz} --- ces & \en{those} \pron{dhouz} --- ces (là-bas) \tabularnewline
\end{tabularx}
\end{tableau}
\ex{This is my friend Ali. Those are my keys.}{Voici mon ami Ali. Ce sont mes clés (là-bas).}
\ex{I don't like that idea. What's this?}{Je n'aime pas cette idée. Qu'est-ce que c'est ?}

\begin{exo}[8.1 --- Traduisez]
\begin{enumerate}
  \item Il y a un problème. \item Il n'y a pas de chaises dans la salle.
  \item Y a-t-il un hôpital près d'ici ? \item Le chat est sous la table.
  \item La banque est entre la pharmacie et le marché. \item Ces livres sont à moi.
  \item Qu'est-ce que c'est ? \item Je suis au travail.
\end{enumerate}
\end{exo}
\begin{corrige}[8.1]
1. \en{There's a problem.} 2. \en{There aren't any chairs in the room.}
3. \en{Is there a hospital near here?} 4. \en{The cat is under the table.}
5. \en{The bank is between the pharmacy and the market.} 6. \en{These books are mine.}
7. \en{What is this?} / \en{What's this?} 8. \en{I'm at work.}
\end{corrige}

\section{Unité 9 --- Les adjectifs}

\begin{regle}[Trois règles d'or]
\begin{enumerate}
  \item L'adjectif se place \textbf{avant} le nom : \en{a big house} (une grande maison).
  \item L'adjectif est \textbf{invariable} : \en{a blue car, two blue cars} (jamais \en{*blues}).
  \item Avec \en{be}, l'adjectif reste après le verbe : \en{The house is big.}
\end{enumerate}
\end{regle}

\begin{tableau}[L'ordre des adjectifs : OSASCOMP]
\begin{tabularx}{\linewidth}{@{}p{2.2cm}p{3.4cm}X@{}}
\entete{Rang} & \entete{Type} & \entete{Exemple dans l'ordre} \tabularnewline
\hline
\textbf{O} & Opinion & \en{beautiful, nice, ugly, expensive, old-fashioned} \tabularnewline
\textbf{S} & Size / shape & \en{big, small, tall, round, long} \tabularnewline
\textbf{A} & Age & \en{new, old, young, ancient} \tabularnewline
\textbf{C} & Colour & \en{red, black, dark blue} \tabularnewline
\textbf{O} & Origin & \en{Senegalese, French, Chinese} \tabularnewline
\textbf{M} & Material & \en{wooden, metal, plastic, cotton} \tabularnewline
\textbf{P} & Purpose & \en{sleeping bag, running shoes, dining room} \tabularnewline
\end{tabularx}
\end{tableau}
\ex{a beautiful big old brown Senegalese wooden house}{une belle grande vieille maison sénégalaise en bois marron}
\textit{(En pratique, on met rarement plus de trois adjectifs.)}

\subsection{Adjectifs en \en{-ed} et \en{-ing}}
\begin{regle}
\en{-ed} = ce que \textbf{je ressens} (\emph{é} comme « je suis fatigué »).\\
\en{-ing} = ce qui \textbf{cause} le sentiment (\emph{ant} comme « c'est fatigant »).
\end{regle}
\ex{I'm bored. --- The film is boring.}{Je m'ennuie. --- Le film est ennuyeux.}
\ex{She's interested in history. --- History is interesting.}{Elle s'intéresse à l'histoire. --- L'histoire est intéressante.}

\begin{piege}[\en{*I am boring}]
Si vous dites \en{I'm boring}, vous dites « je suis quelqu'un d'ennuyeux ». Vous vouliez
dire « je m'ennuie » : \en{I'm bored}. Même piège avec \en{I'm excited} (j'ai hâte / je suis
enthousiaste) contre \en{I'm exciting} (je suis quelqu'un d'excitant), et
\en{I'm tired} contre \en{I'm tiring}.
\end{piege}

\subsection{Les 40 adjectifs de base}
\begin{tableau}[Adjectifs essentiels et leurs contraires]
\footnotesize
\begin{tabularx}{\linewidth}{@{}X X X X@{}}
\en{big} --- \en{small} (grand --- petit) & \en{long} --- \en{short} (long --- court) & \en{tall} --- \en{short} (grand --- petit, taille) & \en{old} --- \en{new / young} (vieux --- neuf / jeune) \\
\en{good} --- \en{bad} (bon --- mauvais) & \en{hot} --- \en{cold} (chaud --- froid) & \en{fast} --- \en{slow} (rapide --- lent) & \en{easy} --- \en{difficult / hard} (facile --- difficile) \\
\en{happy} --- \en{sad} (heureux --- triste) & \en{rich} --- \en{poor} (riche --- pauvre) & \en{clean} --- \en{dirty} (propre --- sale) & \en{full} --- \en{empty} (plein --- vide) \\
\en{strong} --- \en{weak} (fort --- faible) & \en{beautiful} --- \en{ugly} (beau --- laid) & \en{cheap} --- \en{expensive} (bon marché --- cher) & \en{open} --- \en{closed} (ouvert --- fermé) \\
\en{early} --- \en{late} (tôt --- tard) & \en{right} --- \en{wrong} (correct --- faux) & \en{true} --- \en{false} (vrai --- faux) & \en{safe} --- \en{dangerous} (sûr --- dangereux) \\
\en{heavy} --- \en{light} (lourd --- léger) & \en{wide} --- \en{narrow} (large --- étroit) & \en{high} --- \en{low} (haut --- bas) & \en{deep} --- \en{shallow} (profond --- peu profond) \\
\en{young} --- \en{old} (jeune --- âgé) & \en{kind} --- \en{mean} (gentil --- méchant) & \en{quiet} --- \en{noisy} (calme --- bruyant) & \en{free} --- \en{busy} (libre --- occupé) \\
\en{near} --- \en{far} (près --- loin) & \en{wet} --- \en{dry} (mouillé --- sec) & \en{dark} --- \en{bright} (sombre --- lumineux) & \en{interesting} --- \en{boring} (intéressant --- ennuyeux) \\
\end{tabularx}
\end{tableau}

\section{Unité 10 --- L'impératif, \en{can} et les formules de politesse}

\subsection{L'impératif}
\begin{regle}
Base verbale seule, sans sujet : \en{Come here!} \en{Don't touch that!} \en{Let's go!}
(Allons-y). Pour inclure quelqu'un : \en{Let's + base} ; négatif : \en{Let's not argue.}
\end{regle}
\ex{Sit down, please. Don't forget your bag. Let's start.}{Asseyez-vous, s'il vous plaît. N'oubliez pas votre sac. Commençons.}

\subsection{\en{can} : capacité, permission, possibilité}
\begin{tableau}[\en{can} --- invariable, jamais de \en{-s}, jamais de \en{to}]
\begin{tabularx}{\linewidth}{@{}p{3.6cm}p{4.4cm}X@{}}
\entete{Emploi} & \entete{Exemple} & \entete{Traduction} \tabularnewline
\hline
Capacité & \en{I can swim. She can't drive.} & Je sais nager. Elle ne sait pas conduire. \tabularnewline
Permission (informel) & \en{Can I open the window?} & Puis-je ouvrir la fenêtre ? \tabularnewline
Possibilité & \en{It can be dangerous.} & Ça peut être dangereux. \tabularnewline
Demande polie & \en{Can you help me, please?} & Pouvez-vous m'aider ? \tabularnewline
Plus poli : \en{could} & \en{Could you help me?} & Pourriez-vous m'aider ? \tabularnewline
Encore plus poli : \en{would you mind} & \en{Would you mind opening the door?} & Cela vous ennuierait-il d'ouvrir la porte ? \tabularnewline
\end{tabularx}
\end{tableau}

\begin{piege}[\en{*I can to speak English}]
Après un modal (\en{can, could, must, should, will, would, may, might}), le verbe est à la
base, \textbf{sans \en{to}} : \en{I can speak English.} \en{You must go.} \en{She will call you.}
\end{piege}

\subsection{Les formules de politesse (à apprendre comme des blocs)}
\begin{tableau}[La politesse anglaise, du plus direct au plus poli]
\footnotesize
\begin{tabularx}{\linewidth}{@{}X X@{}}
\en{Give me a coffee.} \textit{(ordre --- à éviter)} & \en{I want a coffee.} \textit{(trop direct)} \\
\en{Can I have a coffee, please?} \textit{(courant, parfait)} & \en{Could I have a coffee, please?} \textit{(plus poli)} \\
\en{I'd like a coffee, please.} \textit{(le plus naturel et sûr)} & \en{Would you mind passing me the sugar?} \textit{(très poli)} \\
\end{tabularx}
\end{tableau}

\begin{memo}[Les six blocs de politesse à automatiser]
\en{Could you\ldots, please?} \quad \en{Would you mind + \ldots ing?} \quad
\en{I'd like\ldots} \quad \en{May I\ldots?} \quad \en{Sorry to bother you, but\ldots} \quad
\en{Thanks a lot, I really appreciate it.}
\end{memo}

\section{Unité 11 --- Les prépositions de temps : \en{at}, \en{on}, \en{in}}

\begin{tableau}[La règle des trois prépositions de temps]
\begin{tabularx}{\linewidth}{@{}p{1.8cm}p{5.0cm}X@{}}
\entete{Mot} & \entete{S'emploie pour} & \entete{Exemples} \tabularnewline
\hline
\en{at} & Heures, moments précis, fêtes (sans \en{day}) & \en{at 7 o'clock, at noon, at midnight, at night, at the weekend (BrE), at Christmas, at the moment} \tabularnewline
\en{on} & Jours, dates, jours de fête & \en{on Monday, on 5 May, on my birthday, on Christmas Day, on Friday morning} \tabularnewline
\en{in} & Mois, années, saisons, siècles, durées, parties de la journée & \en{in May, in 2026, in summer, in the 21st century, in the morning, in two hours, in the past} \tabularnewline
\hline
\multicolumn{3}{@{}l}{\textbf{Sans préposition} : \en{today, tomorrow, yesterday, this week, last month, next year, every day, tonight}} \tabularnewline
\end{tabularx}
\end{tableau}

\begin{piege}[\en{*in Monday}, \en{*on the morning}, \en{*at July}]
Trois réflexes français à casser : \en{on Monday} (jamais \en{in}), \en{in the morning}
(mais \en{on Monday morning}), \en{at 8 o'clock} (jamais \en{in}). Et jamais de préposition
devant \en{next / last / this / every} : \en{I'll see you next week} (pas \en{*in next week}).
\end{piege}

\subsection{L'heure et les dates}
\begin{tableau}[Dire l'heure]
\begin{tabularx}{\linewidth}{@{}p{4.0cm}X@{}}
\entete{Heure} & \entete{Façons de la dire} \tabularnewline
\hline
8:00 & \en{It's eight o'clock.} \tabularnewline
8:05 & \en{It's eight oh five.} / \en{It's five past eight.} \tabularnewline
8:15 & \en{It's eight fifteen.} / \en{It's a quarter past eight.} \tabularnewline
8:30 & \en{It's eight thirty.} / \en{It's half past eight.} \tabularnewline
8:45 & \en{It's eight forty-five.} / \en{It's a quarter to nine.} \tabularnewline
8:50 & \en{It's eight fifty.} / \en{It's ten to nine.} \tabularnewline
\end{tabularx}
\end{tableau}

\begin{tableau}[Les dates]
\begin{tabularx}{\linewidth}{@{}p{4.0cm}X@{}}
\entete{Forme} & \entete{Écriture et prononciation} \tabularnewline
\hline
Britannique & \en{29/10/2026} = \en{29 October 2026} = \en{the twenty-ninth of October, twenty twenty-six} \tabularnewline
Américaine & \en{10/29/2026} = \en{October 29, 2026} = \en{October twenty-ninth} \tabularnewline
\textbf{Attention} & \en{10/11} = 10 novembre (BrE) ou 11 octobre (AmE) --- écrivez le mois en lettres dans les documents importants. \tabularnewline
\end{tabularx}
\end{tableau}

\subsection{Les nombres}
\begin{tableau}[Nombres de 0 à 1\,000\,000 --- pièges inclus]
\footnotesize
\begin{tabularx}{\linewidth}{@{}X X@{}}
0 \en{zero / oh} & 11 \en{eleven} \quad 12 \en{twelve} \\
1 \en{one} \quad 2 \en{two} \quad 3 \en{three} & 13 \en{thirteen} \quad 30 \en{thirty} \\
4 \en{four} \quad 5 \en{five} \quad 6 \en{six} & 14 \en{fourteen} \quad 40 \en{forty} (sans \en{u} !) \\
7 \en{seven} \quad 8 \en{eight} \quad 9 \en{nine} & 15 \en{fifteen} \quad 50 \en{fifty} \\
10 \en{ten} & 16 \en{sixteen} \quad 60 \en{sixty} \\
20 \en{twenty} \quad 21 \en{twenty-one} & 17 \en{seventeen} \quad 70 \en{seventy} \\
100 \en{a / one hundred} & 18 \en{eighteen} \quad 80 \en{eighty} \\
1\,000 \en{a / one thousand} & 19 \en{nineteen} \quad 90 \en{ninety} \\
1\,000\,000 \en{a million} & 2026 \en{twenty twenty-six} \tabularnewline
\end{tabularx}
\end{tableau}

\begin{astuce}[Les nombres ne prennent jamais de \en{-s}]
\en{two hundred people}, \en{three thousand cars}, \en{five million} --- jamais
\en{*two hundreds}. Mais : \en{hundreds of people}, \en{thousands of cars}
(avec \en{of}, le pluriel réapparaît).
\end{astuce}

\begin{exo}[11.1 --- Prépositions de temps]
Complétez avec \en{at}, \en{on}, \en{in} ou rien :
\begin{enumerate}
  \item \trou[1cm] Monday \item \trou[1cm] the morning \item \trou[1cm] 6 o'clock
  \item \trou[1cm] 2026 \item \trou[1cm] July \item \trou[1cm] tomorrow
  \item \trou[1cm] night \item \trou[1cm] my birthday \item \trou[1cm] last week
  \item \trou[1cm] the weekend \item \trou[1cm] two hours \item \trou[1cm] summer
\end{enumerate}
\end{exo}
\begin{corrige}[11.1]
1. \en{on} 2. \en{in} 3. \en{at} 4. \en{in} 5. \en{in} 6. rien 7. \en{at}
8. \en{on} 9. rien 10. \en{at} (BrE) / \en{on} (AmE) 11. \en{in} 12. \en{in}
\end{corrige}

\begin{exo}[11.2 --- Écrivez en toutes lettres]
\en{3:15, 7:45, 12:30, 14, 40, 19, 100, 2\,000, 15\,000, 05/12/2026 (BrE), 1987}.
\end{exo}
\begin{corrige}[11.2]
\en{three fifteen / a quarter past three ; seven forty-five / a quarter to eight ;
half past twelve ; fourteen ; forty ; nineteen ; a hundred ; two thousand ;
fifteen thousand ; the fifth of December twenty twenty-six ; nineteen eighty-seven}.
\end{corrige}

\begin{memo}[Révision express du chapitre 8]
\begin{enumerate}
  \item Sujet obligatoire toujours : \en{It is cold}, \en{It rains}, \en{There is a problem}.
  \item \en{be} pour l'âge, la faim, le froid, la peur, la raison.
  \item Présent simple : \en{-s} à la 3\textsuperscript{e} personne, \en{do/does} aux questions
  et négations, verbe de base après \en{do/does/did}.
  \item Article zéro pour les généralités, les repas, les langues, les sports, les pays.
  \item Verbes d'état jamais au \en{-ing}.
  \item Adjectif avant le nom, invariable.
  \item \en{at} + heure, \en{on} + jour, \en{in} + mois/année/période.
  \item Après un modal, base verbale sans \en{to}.
  \item Une seule négation par phrase.
  \item Question : auxiliaire + sujet + verbe.
\end{enumerate}
\end{memo}

% =====================================================================
%  CHAPITRE 9 — GRAMMAIRE A2 : LES TEMPS DU VERBE
% =====================================================================
\chapter{Les temps : parler du passé, du présent et de l'avenir}

Le système verbal anglais est \textbf{plus simple que le français} (pas de passé simple
littéraire, pas de subjonctif obligatoire, pas d'imparfait distinct), mais il fait une
distinction que le français ignore : \textbf{l'action est-elle terminée ou encore ouverte ?}
C'est tout le secret du \en{present perfect}.

\section{La carte générale des temps}
\begin{tableau}[Les 8 temps que vous utiliserez réellement]
\footnotesize
\begin{longtable}{@{}p{2.9cm}p{3.6cm}p{4.2cm}p{3.6cm}@{}}
\entete{Temps} & \entete{Construction} & \entete{Emploi principal} & \entete{Exemple} \tabularnewline
\endfirsthead
\entete{Temps} & \entete{Construction} & \entete{Emploi} & \entete{Exemple} \tabularnewline
\endhead
\hline
\textbf{Présent simple} & base + \en{s} & habitude, vérité générale & \en{I work here.} \tabularnewline
\textbf{Présent continu} & \en{be} + \en{-ing} & action en cours maintenant & \en{I am working.} \tabularnewline
\textbf{Passé simple} & \en{-ed} / irrégulier & action terminée à un moment précis du passé & \en{I worked yesterday.} \tabularnewline
\textbf{Passé continu} & \en{was/were} + \en{-ing} & décor, action en cours à un moment du passé & \en{I was working at 8 pm.} \tabularnewline
\textbf{\en{Present perfect}} & \en{have} + participe passé & lien passé-présent, expérience, résultat & \en{I have worked here for 3 years.} \tabularnewline
\textbf{\en{Perf. continuous}} & \en{have been} + \en{-ing} & durée d'une activité encore en cours & \en{I have been working since 7 am.} \tabularnewline
\textbf{\en{Past perfect}} & \en{had} + participe passé & passé antérieur à un autre passé & \en{I had already eaten when he called.} \tabularnewline
\textbf{Futur} & \en{will} / \en{going to} / présent continu & décision, intention, plan & \en{I'll call you. I'm going to study.} \tabularnewline
\end{longtable}
\end{tableau}

\section{Unité 1 --- Le présent continu (\en{be} + \en{-ing})}

\begin{regle}[Formation]
\en{am / is / are} + verbe en \en{-ing} : \en{I'm working, she's reading, they're coming}.\\
Négation : \en{I'm not working.} \quad Question : \en{Are you working?} \en{What are you doing?}
\end{regle}

\begin{tableau}[Orthographe du \en{-ing}]
\begin{tabularx}{\linewidth}{@{}p{5.4cm}X@{}}
\entete{Cas} & \entete{Exemples} \tabularnewline
\hline
Règle générale : + \en{ing} & \en{work} \ra \en{working}, \en{play} \ra \en{playing}, \en{read} \ra \en{reading} \tabularnewline
Verbe se terminant par \en{e} muet : \en{e} \ra \en{ing} & \en{make} \ra \en{making}, \en{write} \ra \en{writing}, \en{come} \ra \en{coming}, \en{live} \ra \en{living} \tabularnewline
\en{-ee}, \en{-ye}, \en{-oe} : on garde tout & \en{see} \ra \en{seeing}, \en{agree} \ra \en{agreeing}, \en{dye} \ra \en{dyeing} \tabularnewline
Une syllabe, voyelle + consonne : on double & \en{sit} \ra \en{sitting}, \en{run} \ra \en{running}, \en{get} \ra \en{getting}, \en{stop} \ra \en{stopping}, \en{begin} \ra \en{beginning} \tabularnewline
\en{-ie} \ra \en{-ying} & \en{die} \ra \en{dying}, \en{lie} \ra \en{lying}, \en{tie} \ra \en{tying} \tabularnewline
\end{tabularx}
\end{tableau}

\subsection{Quand l'employer ?}
\begin{enumerate}
  \item \textbf{Action en cours au moment où je parle} : \en{Be quiet, I'm working.}
  \item \textbf{Situation temporaire} : \en{I'm living with my brother this month.}
  \item \textbf{Tendance actuelle, changement} : \en{Prices are rising. My English is improving.}
  \item \textbf{Projet futur déjà décidé} : \en{I'm travelling to Dakar tomorrow.}
  \item \textbf{Action répétée avec \en{always}} (agacement) : \en{He's always losing his keys!}
\end{enumerate}

\begin{tableau}[Présent simple contre présent continu]
\begin{tabularx}{\linewidth}{@{}p{5.4cm}p{5.4cm}X@{}}
\entete{Simple (habituel, permanent)} & \entete{Continu (maintenant, temporaire)} & \entete{Français} \tabularnewline
\hline
\en{I work in a bank.} & \en{I'm working on a report.} & Je travaille dans une banque. / Je suis en train de travailler sur un rapport. \tabularnewline
\en{She drinks tea every morning.} & \en{She's drinking tea now.} & Elle boit du thé chaque matin. / Elle boit du thé maintenant. \tabularnewline
\en{It rains a lot in August.} & \en{It's raining.} & Il pleut beaucoup en août. / Il pleut (là). \tabularnewline
\en{They live in Thiès.} & \en{They're staying at a hotel.} & Ils habitent Thiès. / Ils logent à l'hôtel. \tabularnewline
\end{tabularx}
\end{tableau}

\begin{astuce}[Les mots qui déclenchent le continu]
\en{now, right now, at the moment, currently, today, this week, these days, look!, listen!,
still, temporarily} --- dès que vous les voyez, pensez \en{-ing}.
\end{astuce}

\begin{piege}[Le français « je travaille » est ambigu, l'anglais non]
« Je travaille dans une banque » = \en{I work in a bank} (permanent).\\
« Je travaille en ce moment » = \en{I'm working right now} (en cours).\\
Si vous dites \en{I'm working in a bank}, un anglophone comprend « temporairement, cette
période-ci ». Choisissez toujours selon le sens, jamais selon la traduction mot à mot.
\end{piege}

\begin{exo}[1.1 --- Simple ou continu ?]
\begin{enumerate}
  \item Look! It (rain) \trou[2cm].
  \item She usually (take) \trou[2cm] the bus, but today she (walk) \trou[2cm].
  \item I (not / understand) \trou[2.5cm] this exercise.
  \item They (build) \trou[2cm] a new bridge at the moment.
  \item He (know) \trou[2cm] the answer.
  \item We (study) \trou[2cm] English every day.
  \item Shh! The baby (sleep) \trou[2cm].
  \item My brother (be) \trou[2cm] a teacher, but this summer he (work) \trou[2cm] in a shop.
\end{enumerate}
\end{exo}
\begin{corrige}[1.1]
1. \en{is raining} 2. \en{takes} --- \en{is walking} 3. \en{don't understand}
4. \en{are building} 5. \en{knows} 6. \en{study} 7. \en{is sleeping}
8. \en{is} --- \en{is working}
\end{corrige}

\section{Unité 2 --- Le passé simple (\en{past simple})}

\begin{regle}[Formation]
\textbf{Verbes réguliers :} base + \en{-ed} --- \en{I worked, she played, we watched}.\\
\textbf{Verbes irréguliers :} 2\textsuperscript{e} forme à apprendre par c\oe ur
(\en{go} \ra \en{went}, \en{see} \ra \en{saw}, \en{eat} \ra \en{ate}). Voir l'annexe A.\\
\textbf{\en{be} :} \en{was} (\en{I/he/she/it}) / \en{were} (\en{you/we/they}).\\
\textbf{Négation :} \en{didn't} + \textbf{base} --- \en{I didn't go} (jamais \en{*didn't went}).\\
\textbf{Question :} \en{did} + sujet + \textbf{base} --- \en{Did you see him?}
\end{regle}

\begin{tableau}[Prononciation du \en{-ed} : trois sons, pas un]
\begin{tabularx}{\linewidth}{@{}p{2.4cm}p{4.4cm}X@{}}
\entete{Son} & \entete{Après} & \entete{Exemples} \tabularnewline
\hline
/t/ & consonne sourde (\en{p, k, f, s, sh, ch, x}) & \en{worked} \pron{oueu-kt}, \en{stopped} \pron{stopt}, \en{watched} \pron{ouo-tcht} \tabularnewline
/d/ & consonne sonore ou voyelle & \en{played} \pron{plé-id}, \en{called} \pron{kold}, \en{lived} \pron{livd} \tabularnewline
/\ipaiot d/ & \en{t} ou \en{d} final & \en{wanted} \pron{ouon-tid}, \en{needed} \pron{nii-did}, \en{decided} \pron{di-saï-did} \tabularnewline
\end{tabularx}
\end{tableau}

\subsection{Quand employer le passé simple ?}
\textbf{Une action terminée, à un moment déterminé du passé.} Les mots-signaux sont vos
amis : \en{yesterday, last night / week / year, two days ago, in 2010, when I was a child,
then, at that time}.

\ex{I visited Dakar last month. She didn't come yesterday. Did you call him?}{J'ai visité Dakar le mois dernier. Elle n'est pas venue hier. L'as-tu appelé ?}

\begin{piege}[\en{*I have seen him yesterday}]
Dès qu'un \textbf{moment passé précis} est mentionné (\en{yesterday}, \en{last week},
\en{in 2015}, \en{two hours ago}), c'est le \textbf{passé simple} obligatoire. Le
\en{present perfect} est alors une faute grave --- et les francophones la font pendant des
années. Comparez :\\
\en{I saw him yesterday.} \ok \qquad \en{*I have seen him yesterday.} \ko\\
\en{I have seen him.} \ok (sans date : le lien avec le présent compte)
\end{piege}

\begin{piege}[\en{was} / \en{were} et le subjonctif \en{were}]
\en{I was tired, they were happy}. Mais dans une hypothèse irréelle, on emploie
\en{were} pour toutes les personnes : \en{If I \textbf{were} you\ldots} (Si j'étais toi\ldots),
\en{I wish I \textbf{were} taller}. Forme familière acceptée : \en{If I was you\ldots}
\end{piege}

\begin{exo}[2.1 --- Mettez au passé simple]
\begin{enumerate}
  \item \en{go, see, eat, have, do, make, take, come, say, tell}
  \item \en{buy, bring, think, teach, catch, meet, feel, leave, find, know}
  \item \en{work, play, watch, study, stop, live, want, need, call, try}
\end{enumerate}
\end{exo}
\begin{corrige}[2.1]
1. \en{went, saw, ate, had, did, made, took, came, said, told}
2. \en{bought, brought, thought, taught, caught, met, felt, left, found, knew}
3. \en{worked, played, watched, studied, stopped, lived, wanted, needed, called, tried}
\end{corrige}

\begin{exo}[2.2 --- Traduisez]
\begin{enumerate}
  \item Je ne suis pas venu hier. \item Est-ce que tu as mangé ?
  \item Elle a acheté une voiture l'année dernière. \item Nous étions très fatigués.
  \item Ils ne savaient pas la réponse. \item Quand je vivais à Dakar, je travaillais dans un magasin.
\end{enumerate}
\end{exo}
\begin{corrige}[2.2]
1. \en{I didn't come yesterday.} 2. \en{Did you eat?} 3. \en{She bought a car last year.}
4. \en{We were very tired.} 5. \en{They didn't know the answer.}
6. \en{When I lived in Dakar, I worked in a shop.}
\end{corrige}

\section{Unité 3 --- Le passé continu et le récit au passé}

\begin{regle}
\en{was / were} + \en{-ing}. Emplois :
\textbf{(1)} décor d'une histoire : \en{The sun was shining and the birds were singing.}
\textbf{(2)} action en cours interrompue par une autre :
\en{I was sleeping when the phone rang.}
\textbf{(3)} deux actions simultanées : \en{While I was cooking, he was watching TV.}
\end{regle}

\begin{tableau}[\en{when} contre \en{while}]
\begin{tabularx}{\linewidth}{@{}p{5.6cm}X@{}}
\entete{Structure} & \entete{Exemple} \tabularnewline
\hline
\en{was/were + ing} + \en{when} + passé simple & \en{I was reading when she arrived.} (action longue interrompue) \tabularnewline
\en{passé simple} + \en{while} + \en{was/were + ing} & \en{She arrived while I was reading.} \tabularnewline
\en{while} + \en{ing} (sujet commun) & \en{While walking, I saw him.} \tabularnewline
Deux passés continus & \en{While I was working, they were playing.} \tabularnewline
\end{tabularx}
\end{tableau}

\begin{astuce}[Raconter une histoire : le cocktail des trois passés]
Un bon récit anglais mélange trois temps, toujours dans le même ordre logique :\\
\textbf{1. Passé continu} pour le décor : \en{It was raining and I was waiting for the bus.}\\
\textbf{2. Passé simple} pour les événements : \en{Suddenly a man came and spoke to me.}\\
\textbf{3. \en{past perfect}} pour l'antériorité : \en{I realised I had met him before.}\\
Maîtrisez ce trio et vos récits sonneront immédiatement « anglais ».
\end{astuce}

\begin{exo}[3.1 --- Passé simple ou passé continu ?]
\begin{enumerate}
  \item I (walk) \trou[2cm] home when I (see) \trou[2cm] an accident.
  \item While she (cook) \trou[2cm], the lights (go) \trou[2cm] out.
  \item What (you / do) \trou[2.5cm] at 9 pm yesterday?
  \item He (not / listen) \trou[2.5cm], so he (not / understand) \trou[2.5cm].
  \item We (drive) \trou[2cm] to Thiès when the car (break) \trou[2cm] down.
\end{enumerate}
\end{exo}
\begin{corrige}[3.1]
1. \en{was walking} --- \en{saw} 2. \en{was cooking} --- \en{went}
3. \en{were you doing} 4. \en{wasn't listening} --- \en{didn't understand}
5. \en{were driving} --- \en{broke}
\end{corrige}

\section{Unité 4 --- Le \en{present perfect} : le chapitre le plus important du livre}

\begin{regle}[Formation]
\en{have / has} + \textbf{participe passé} (3\textsuperscript{e} forme) :\\
\en{I have finished. She has gone. We have seen it. They haven't arrived. Have you eaten?}\\
Contractions : \en{I've, you've, he's, she's, we've, they've} ; \en{haven't, hasn't}.
\end{regle}

\subsection{L'idée en une image}
Le \textbf{passé simple} place l'action dans une boîte fermée du passé : le pont est coupé
avec le présent. Le \textbf{\en{present perfect}} est un \textbf{pont} : l'action est passée,
mais son \emph{effet} ou sa \emph{validité} touche le présent.

\begin{tableau}[Les cinq emplois, avec le test de traduction]
\begin{tabularx}{\linewidth}{@{}p{3.2cm}p{4.6cm}X@{}}
\entete{Emploi} & \entete{Exemple} & \entete{Ce que cela veut dire maintenant} \tabularnewline
\hline
\textbf{1. Résultat présent} & \en{I've lost my keys.} & Je ne les ai \emph{toujours} pas --- je suis bloqué maintenant. (\en{I lost my keys} = simple constat passé, je les ai peut-être retrouvées.) \tabularnewline
\textbf{2. Expérience de vie} & \en{I've been to London twice.} & Dans ma vie jusqu'à maintenant. Aucune date n'est donnée. \tabularnewline
\textbf{3. Durée jusqu'à présent} (avec \en{for/since}) & \en{I've lived here for ten years.} & J'y habite \emph{encore}. (\en{I lived here for ten years} = je n'y habite plus.) \tabularnewline
\textbf{4. Période non terminée} & \en{I've drunk three coffees today.} & La journée n'est pas finie, je peux en boire un quatrième. \tabularnewline
\textbf{5. Nouvelle récente} & \en{The president has arrived.} & On vient de me l'annoncer, c'est d'actualité. \tabularnewline
\end{tabularx}
\end{tableau}

\subsection{Les six mots-signaux}
\begin{tableau}[\en{just, already, yet, ever, never, for, since, recently, so far}]
\begin{tabularx}{\linewidth}{@{}p{2.6cm}p{4.4cm}X@{}}
\entete{Mot} & \entete{Position} & \entete{Exemple et traduction} \tabularnewline
\hline
\en{just} & entre \en{have} et le participe & \en{I've just finished.} --- \textit{Je viens de finir.} \tabularnewline
\en{already} & entre \en{have} et le participe & \en{She's already left.} --- \textit{Elle est déjà partie.} (affirmatif) \tabularnewline
\en{yet} & \textbf{en fin} de phrase & \en{Have you finished yet?} / \en{I haven't eaten yet.} --- \textit{déjà / encore} (question et négation) \tabularnewline
\en{ever} & dans les questions & \en{Have you ever eaten sushi?} --- \textit{As-tu déjà mangé des sushis ?} \tabularnewline
\en{never} & avant le participe & \en{I've never seen snow.} --- \textit{Je n'ai jamais vu la neige.} \tabularnewline
\en{for} & + \textbf{durée} & \en{for two hours, for ten years, for a long time} \tabularnewline
\en{since} & + \textbf{point de départ} & \en{since 2015, since Monday, since I was a child} \tabularnewline
\en{recently / lately} & en fin ou avant le participe & \en{Have you seen him recently?} --- \textit{récemment} \tabularnewline
\en{so far / up to now} & en début ou fin & \en{So far, everything is fine.} --- \textit{Jusqu'ici, tout va bien.} \tabularnewline
\end{tabularx}
\end{tableau}

\begin{piege}[\en{for} et \en{since} : la faute la plus fréquente]
\en{*I live here since 2010.} \ko \ra \en{I've lived here since 2010.} \ok\\
\en{*I am knowing him for five years.} \ko \ra \en{I've known him for five years.} \ok\\
Le français utilise le présent (« j'habite ici depuis 2010 ») ; l'anglais exige le
\en{present perfect}. C'est un automatisme à installer : dès que vous voyez
« depuis », votre cerveau doit déclencher \en{have}.
\end{piege}

\begin{piege}[\en{been} contre \en{gone}]
\en{She has \textbf{been} to Paris.} = elle y est allée (et elle est revenue --- expérience).\\
\en{She has \textbf{gone} to Paris.} = elle y est allée (et elle y est encore --- elle est absente).\\
\en{He's been out.} = il était sorti (il est rentré). \quad \en{He's gone out.} = il est sorti (il n'est pas là).
\end{piege}

\subsection{\en{Present perfect} contre \en{past simple} : le test décisif}
\begin{tableau}[Le seul critère qui compte : la période est-elle finie ?]
\begin{tabularx}{\linewidth}{@{}p{6.0cm}X@{}}
\entete{\en{past simple} (période terminée, date donnée)} & \entete{\en{present perfect} (période ouverte, pas de date)} \tabularnewline
\hline
\en{I saw him \textbf{yesterday}.} & \en{I've seen him \textbf{before}.} \tabularnewline
\en{She \textbf{lived} in Dakar for five years.} (elle n'y vit plus) & \en{She \textbf{has lived} in Dakar for five years.} (elle y vit encore) \tabularnewline
\en{We \textbf{had} breakfast at 7.} & \en{We \textbf{haven't had} breakfast \textbf{yet}.} \tabularnewline
\en{\textbf{Did} you \textbf{go} out last night?} & \en{\textbf{Have} you ever \textbf{been} to Paris?} \tabularnewline
\en{Shakespeare \textbf{wrote} 37 plays.} (il est mort) & \en{My favourite author \textbf{has written} 12 novels.} (il vit encore) \tabularnewline
\en{I \textbf{lost} my keys, but I found them.} & \en{I \textbf{have lost} my keys --- I can't get in.} \tabularnewline
\end{tabularx}
\end{tableau}

\begin{astuce}[L'astuce de la « conversation qui démarre »]
En anglais réel, une conversation commence souvent par un \en{present perfect}, puis
bascule au passé simple dès qu'on entre dans les détails :\\
\en{``Have you seen the new film?''} --- \en{``Yes, I have!''} --- \en{``When did you see it?''}
--- \en{``I saw it on Saturday. It was amazing.''}\\
C'est le schéma naturel : \textbf{ouverture au \en{perfect}, détails au passé simple}.
\end{astuce}

\subsection{Le \en{present perfect continuous}}
\begin{regle}
\en{have / has been} + \en{-ing}. Sert à mettre l'accent sur la \textbf{durée} ou sur
l'\textbf{activité} (souvent encore en cours ou tout juste arrêtée, avec un effet visible).
\end{regle}
\ex{I've been waiting for two hours.}{J'attends depuis deux heures (et j'attends encore).}
\ex{You look tired. --- I've been working all night.}{Tu as l'air fatigué. --- J'ai travaillé toute la nuit.}
\ex{It has been raining. The ground is wet.}{Il a plu (on voit le résultat).}
\ex{How long have you been learning English?}{Depuis combien de temps apprends-tu l'anglais ?}

\begin{tableau}[\en{perfect simple} contre \en{perfect continuous}]
\begin{tabularx}{\linewidth}{@{}p{6.0cm}X@{}}
\entete{Simple : résultat, quantité, achèvement} & \entete{Continu : durée, activité, processus} \tabularnewline
\hline
\en{I've read three books this month.} (quantité achevée) & \en{I've been reading a lot this month.} (activité) \tabularnewline
\en{She's written the report.} (c'est fini) & \en{She's been writing the report all day.} (elle y travaille encore) \tabularnewline
\en{I've known him for years.} (verbe d'état --- jamais de continu) & \en{I've been knowing him} \ko \textbf{interdit} \tabularnewline
\end{tabularx}
\end{tableau}

\begin{piege}[Verbes d'état : jamais au \en{-ing}, même au \en{perfect continuous}]
\en{*I have been knowing him for years.} \ra \en{I have known him for years.}\\
Idem : \en{*have been wanting, *have been believing, *have been owning, *have been meaning}.
\end{piege}

\begin{exo}[4.1 --- \en{present perfect} ou \en{past simple} ?]
\begin{enumerate}
  \item I (lose) \trou[2cm] my phone yesterday.
  \item I (lose) \trou[2cm] my phone --- can you help me find it?
  \item She (visit) \trou[2cm] Morocco three times in her life.
  \item She (visit) \trou[2cm] Morocco in 2019.
  \item We (know) \trou[2cm] each other since 2012.
  \item I (see) \trou[2cm] that film last week.
  \item (you / ever / eat) \trou[3cm] sushi?
  \item He just (arrive) \trou[2cm].
  \item My grandmother (write) \trou[2cm] six books before she died.
  \item I (not / finish) \trou[2.5cm] my homework yet.
\end{enumerate}
\end{exo}
\begin{corrige}[4.1]
1. \en{lost} 2. \en{have lost} 3. \en{has visited} 4. \en{visited} 5. \en{have known}
6. \en{saw} 7. \en{Have you ever eaten} 8. \en{has just arrived} 9. \en{wrote}
10. \en{haven't finished}
\end{corrige}

\begin{exo}[4.2 --- \en{for} ou \en{since} ?]
\begin{enumerate}
  \item \trou[1.5cm] two years \item \trou[1.5cm] 2023 \item \trou[1.5cm] I was a child
  \item \trou[1.5cm] three hours \item \trou[1.5cm] last Monday \item \trou[1.5cm] a long time
  \item \trou[1.5cm] we met \item \trou[1.5cm] ten minutes
\end{enumerate}
\end{exo}
\begin{corrige}[4.2]
\en{for} : 1, 4, 6, 8 (durées). \quad \en{since} : 2, 3, 5, 7 (points de départ).
\end{corrige}

\begin{exo}[4.3 --- Traduisez (pièges)]
\begin{enumerate}
  \item J'habite ici depuis dix ans. \item Je connais Fatou depuis l'école.
  \item Elle est allée au marché (elle n'est pas là). \item Elle est allée au marché trois fois cette semaine.
  \item J'attends depuis une heure. \item Tu as déjà mangé ? \item Je n'ai jamais vu ça.
  \item Nous avons acheté cette voiture en 2020. \item Il vient de sortir.
  \item Depuis combien de temps apprenez-vous l'anglais ?
\end{enumerate}
\end{exo}
\begin{corrige}[4.3]
1. \en{I've lived here for ten years.} 2. \en{I've known Fatou since school.}
3. \en{She's gone to the market.} 4. \en{She's been to the market three times this week.}
5. \en{I've been waiting for an hour.} 6. \en{Have you eaten yet?} / \en{Have you already eaten?}
7. \en{I've never seen that.} 8. \en{We bought this car in 2020.}
9. \en{He's just gone out.} 10. \en{How long have you been learning English?}
\end{corrige}

\section{Unité 5 --- Le \en{past perfect} : le passé du passé}

\begin{regle}
\en{had} + participe passé (toutes les personnes) : \en{I had eaten, she had left,
they hadn't come, had you seen it?}
\end{regle}
\textbf{Emploi :} marquer qu'une action est \textbf{antérieure} à une autre action passée.
\ex{When I arrived, the train had left.}{Quand je suis arrivé, le train était (déjà) parti.}
\ex{She was tired because she hadn't slept.}{Elle était fatiguée parce qu'elle n'avait pas dormi.}
\ex{I realised I had met him before.}{J'ai réalisé que je l'avais déjà rencontré.}
\ex{He said he had finished.}{Il a dit qu'il avait fini. (discours rapporté)}

\begin{astuce}[Quand il est facultatif]
Si l'ordre chronologique est clair grâce à \en{after} ou \en{before}, le passé simple suffit :
\en{After she finished, she left.} Mais avec \en{when}, le \en{past perfect} lève toute
ambiguïté : \en{When she had finished, she left.} (elle a fini, \emph{puis} elle est partie)
contre \en{When she finished, I left.} (au moment où elle a fini, je suis parti).
\end{astuce}

\section{Unité 6 --- Le futur : quatre façons, quatre nuances}

\begin{tableau}[Les quatre futurs et leurs nuances]
\begin{tabularx}{\linewidth}{@{}p{3.0cm}p{4.0cm}X@{}}
\entete{Forme} & \entete{Emploi} & \entete{Exemples} \tabularnewline
\hline
\en{will} + base & Décision prise \textbf{à l'instant}, promesse, prédiction, offre & \en{It's cold. --- I'll close the window.} / \en{I'll help you.} / \en{It will rain tomorrow.} \tabularnewline
\en{be going to} + base & Intention \textbf{déjà formée}, prévision fondée sur un signe visible & \en{I'm going to study tonight.} / \en{Look at those clouds --- it's going to rain.} \tabularnewline
\en{présent continu} & Plan \textbf{organisé, fixé} (rendez-vous, billet acheté) & \en{I'm meeting Ali at 5.} / \en{We're flying to Paris on Monday.} \tabularnewline
\en{présent simple} & Horaire officiel, programme & \en{The train leaves at 7 am.} / \en{The meeting starts at 9.} \tabularnewline
\end{tabularx}
\end{tableau}

\begin{piege}[Ne traduisez pas « aller + infinitif » systématiquement par \en{going to}]
« Je vais le faire » (décision immédiate) = \en{I'll do it.}\\
« Je vais le faire » (projet décidé avant) = \en{I'm going to do it.}\\
« Je le fais demain » (rendez-vous fixé) = \en{I'm doing it tomorrow.}\\
En conversation, les natifs disent très souvent \en{gonna} : \en{I'm gonna call you}.
\end{piege}

\subsection{\en{will} : les formes contractées et les réponses courtes}
\en{I'll, you'll, he'll, she'll, we'll, they'll} \pron{aï-l, iou-l, hi-l\ldots}\\
Négation : \en{won't} \pron{ouount} --- \en{I won't be late.}\\
Question : \en{Will you come?} \quad Réponse courte : \en{Yes, I will.} / \en{No, I won't.}\\
\en{shall} (surtout britannique, avec \en{I/we}) : \en{Shall we go?} (On y va ?),
\en{Shall I open the window?} (Je t'ouvre la fenêtre ?)

\subsection{Le futur continu et le futur antérieur (niveau B1-B2)}
\ex{This time tomorrow, I'll be flying to Dakar.}{Demain à cette heure-ci, je serai en train de voler vers Dakar.}
\ex{By 2028, I will have finished my degree.}{D'ici 2028, j'aurai terminé mes études.}
\ex{Will you have left by 6 pm?}{Seras-tu parti avant 18 h ?}

\subsection{Les subordonnées de temps : présent, pas futur !}
\begin{regle}
Après \en{when, as soon as, before, after, until, if, while}, on emploie le
\textbf{présent} pour parler de l'avenir.
\end{regle}
\ex{I'll call you when I \textbf{arrive}.}{Je t'appellerai quand j'arriverai. (\en{*when I will arrive})}
\ex{We'll eat as soon as he \textbf{comes}.}{Nous mangerons dès qu'il viendra.}
\ex{If it \textbf{rains}, we'll stay home.}{S'il pleut, nous resterons à la maison.}

\begin{exo}[6.1 --- Choisissez le futur]
\begin{enumerate}
  \item The phone is ringing. --- I (answer) \trou[2cm] it.
  \item I (visit) \trou[2.5cm] my uncle next weekend, I've already bought the ticket.
  \item Look out! You (fall) \trou[2cm]!
  \item The conference (begin) \trou[2cm] at 9 am tomorrow.
  \item I promise I (not / tell) \trou[2.5cm] anyone.
  \item We (study) \trou[2.5cm] harder next year, that's our plan.
  \item I'll send it as soon as I (get) \trou[2cm] the file.
  \item This time next week, I (lie) \trou[3cm] on the beach.
\end{enumerate}
\end{exo}
\begin{corrige}[6.1]
1. \en{will answer} / \en{'ll answer} (décision immédiate) 2. \en{am visiting} (plan fixé)
3. \en{are going to fall} (signe visible) 4. \en{begins} (horaire)
5. \en{won't tell} 6. \en{are going to study} (intention) 7. \en{get}
8. \en{will be lying} (futur continu)
\end{corrige}

\section{Unité 7 --- \en{used to} et les habitudes passées}

\begin{regle}
\en{used to} + base = une habitude ou un état \textbf{révolu} du passé.\\
\en{I used to smoke.} (je fumais, maintenant non)\\
Négation : \en{I didn't use to smoke.} \quad Question : \en{Did you use to smoke?}
\end{regle}

\begin{tableau}[Trois formes à ne jamais confondre]
\begin{tabularx}{\linewidth}{@{}p{3.4cm}p{4.0cm}X@{}}
\entete{Forme} & \entete{Sens} & \entete{Exemple} \tabularnewline
\hline
\en{used to} + base & habitude passée révolue & \en{I used to live in Dakar.} --- \textit{J'habitais Dakar (avant).} \tabularnewline
\en{be used to} + nom / \en{-ing} & être habitué à (état actuel) & \en{I'm used to the heat.} --- \textit{Je suis habitué à la chaleur.} \tabularnewline
\en{get used to} + nom / \en{-ing} & s'habituer à (processus) & \en{I'm getting used to waking up early.} --- \textit{Je m'habitue à me lever tôt.} \tabularnewline
\end{tabularx}
\end{tableau}

\begin{piege}[Prononciation et orthographe]
\en{used to} se prononce \pron{iou-s-tou}, jamais \pron{iouz-dou}. Et dans
\en{I'm used to \textbf{working}}, le verbe qui suit est en \en{-ing} --- car \en{to} y est
une préposition, pas un infinitif. \en{*I'm used to work} est une faute classique.
\end{piege}

\ex{When I was young, I used to play football every day.}{Quand j'étais jeune, je jouais au football tous les jours.}
\ex{He didn't use to like coffee, but now he drinks it every morning.}{Il n'aimait pas le café avant, mais maintenant il en boit chaque matin.}
\ex{She would sit by the window for hours.}{Elle s'asseyait près de la fenêtre pendant des heures. (\en{would} pour les habitudes passées répétées --- jamais pour les états : \en{*I would have a car})}

\section{Unité 8 --- Le tableau de bord de tous les temps}
Gardez ce tableau sous les yeux pendant deux mois : il vaut tous les résumés.

\begin{tableau}[Le verbe \en{to work} et le verbe \en{to go} à tous les temps]
\footnotesize
\begin{longtable}{@{}p{3.4cm}p{5.0cm}p{5.4cm}@{}}
\entete{Temps} & \entete{Affirmation} & \entete{Signal / sens} \tabularnewline
\endfirsthead
\entete{Temps} & \entete{Affirmation} & \entete{Signal / sens} \tabularnewline
\endhead
\hline
Présent simple & \en{I work. He works.} & \en{every day, usually, always} \tabularnewline
Présent continu & \en{I am working.} & \en{now, at the moment, today} \tabularnewline
Passé simple & \en{I worked. I went.} & \en{yesterday, last week, in 2010, ago} \tabularnewline
Passé continu & \en{I was working. I was going.} & \en{while, when, at 8 pm yesterday} \tabularnewline
\en{Present perfect} & \en{I have worked. I have gone.} & \en{ever, never, just, already, yet, for, since, so far} \tabularnewline
\en{Perf. continuous} & \en{I have been working.} & \en{for two hours, all day, how long, since} \tabularnewline
\en{Past perfect} & \en{I had worked. I had gone.} & \en{before, after, by the time, when} \tabularnewline
\en{Past perf. continuous} & \en{I had been working.} & durée avant un moment passé \tabularnewline
Futur \en{will} & \en{I will work. I'll go.} & décision, promesse, \en{tomorrow, next year} \tabularnewline
Futur \en{going to} & \en{I am going to work.} & intention, plan, évidence \tabularnewline
Futur continu & \en{I will be working.} & \en{this time tomorrow} \tabularnewline
Futur antérieur & \en{I will have worked.} & \en{by 2030, by the time} \tabularnewline
\end{longtable}
\end{tableau}

\begin{exo}[8.1 --- Le grand test des temps (20 phrases)]
Conjuguez correctement :
\begin{enumerate}
  \item Right now, I (read) \trou[2cm] a book about Africa.
  \item I (read) \trou[2cm] three books this year.
  \item I (read) \trou[2cm] that book last year.
  \item She (live) \trou[2cm] in Paris since 2018.
  \item She (live) \trou[2cm] in Paris from 2010 to 2015.
  \item We (wait) \trou[2.5cm] for you for an hour!
  \item When I (arrive) \trou[2cm], the meeting (already / start) \trou[2.5cm].
  \item Look at the sky --- it (rain) \trou[2.5cm].
  \item I (call) \trou[2cm] you as soon as I (finish) \trou[2cm].
  \item Every morning he (drink) \trou[2cm] tea, but yesterday he (drink) \trou[2cm] coffee.
  \item (you / ever / be) \trou[3cm] to the United States?
  \item I (not / see) \trou[2.5cm] him since Monday.
  \item By next June, they (build) \trou[3cm] the whole house.
  \item While I (walk) \trou[2.5cm] home, I (meet) \trou[2cm] an old friend.
  \item He (be) \trou[2cm] tired because he (work) \trou[3cm] all night.
  \item This time next year, I (study) \trou[3cm] at university.
  \item I (used to / play) \trou[2.5cm] the guitar when I was a teenager.
  \item She (get) \trou[2.5cm] used to the noise now, it's fine.
  \item The shop (open) \trou[2cm] at 9 and (close) \trou[2cm] at 6 every day.
  \item (you / watch) \trou[3cm] TV when I called?
\end{enumerate}
\end{exo}
\begin{corrige}[8.1]
1. \en{am reading} 2. \en{have read} 3. \en{read} \pron{rèd} 4. \en{has lived}
5. \en{lived} 6. \en{have been waiting} 7. \en{arrived} --- \en{had already started}
8. \en{is going to rain} 9. \en{will call / 'll call} --- \en{finish}
10. \en{drinks} --- \en{drank} 11. \en{Have you ever been} 12. \en{haven't seen}
13. \en{will have built} 14. \en{was walking} --- \en{met} 15. \en{was} ---
\en{had been working} 16. \en{will be studying} 17. \en{used to play}
18. \en{is getting} 19. \en{opens} --- \en{closes} 20. \en{Were you watching}
\end{corrige}

\begin{mission}[Chapitre 9 --- les trois récits]
Enregistrez-vous trois fois, une fois par jour :\\
\textbf{Jour 1 --- hier :} 2 minutes au passé simple et au passé continu
(\en{Yesterday I woke up at\ldots\ While I was\ldots, I\ldots}).\\
\textbf{Jour 2 --- votre vie :} 2 minutes au \en{present perfect}
(\en{I have lived\ldots\ I have never\ldots\ I have known\ldots for\ldots}).\\
\textbf{Jour 3 --- l'avenir :} 2 minutes avec \en{will}, \en{going to} et le présent continu
(\en{Next year I'm going to\ldots\ In ten years I'll\ldots\ On Saturday I'm meeting\ldots}).\\
Réécoutez le jour 4 : comptez vos erreurs de temps. Ce sont vos priorités de la semaine
suivante.
\end{mission}

% =====================================================================
%  CHAPITRE 10 — GRAMMAIRE B1 : MODAUX, CONDITIONNELS, PASSIF, RAPPORTÉ
% =====================================================================
\chapter{Structures intermédiaires : modaux, conditions, passif, récit}

Ce chapitre contient les structures qui font passer un discours de « scolaire » à « adulte ».
Chacune est courte à comprendre et longue à automatiser : apprenez-les en blocs, pas en règles.

\section{Unité 1 --- Les verbes modaux}

\begin{regle}[Ce qu'est un modal]
Un modal (\en{can, could, may, might, must, shall, should, will, would}) ne change
\textbf{jamais} : pas de \en{-s}, pas de \en{-ing}, pas de \en{to} après lui, et il forme
seul la négation et la question.\\
\en{She can swim.} \quad \en{She can't swim.} \quad \en{Can she swim?}\\
\textbf{Structure :} sujet + modal + \textbf{base verbale}.
\end{regle}

\begin{tableau}[Le grand tableau des modaux et de leurs sens]
\footnotesize
\begin{longtable}{@{}p{2.6cm}p{4.0cm}p{4.6cm}p{2.6cm}@{}}
\entete{Modal} & \entete{Sens} & \entete{Exemple} & \entete{Traduction} \tabularnewline
\endfirsthead
\entete{Modal} & \entete{Sens} & \entete{Exemple} & \entete{Traduction} \tabularnewline
\endhead
\hline
\en{can} & capacité, possibilité, permission informelle & \en{I can drive. Can I help?} & je peux, je sais \tabularnewline
\en{could} & capacité passée, demande polie, possibilité faible & \en{I could swim at 5. Could you help me?} & je pouvais, pourriez-vous \tabularnewline
\en{be able to} & capacité (remplace \en{can} aux temps manquants) & \en{I will be able to come. I haven't been able to sleep.} & je pourrai, je n'ai pas pu \tabularnewline
\en{may} & permission formelle, possibilité moyenne & \en{May I come in? It may rain.} & puis-je, il se peut \tabularnewline
\en{might} & possibilité faible (20-30\,\%) & \en{She might be late.} & elle sera peut-être en retard \tabularnewline
\en{must} & obligation forte venant de moi, déduction logique & \en{I must study. He must be tired.} & je dois, il doit être \tabularnewline
\en{have to} & obligation venant de l'extérieur (règle, loi) & \en{I have to wear a uniform.} & je dois (obligé par les circonstances) \tabularnewline
\en{must not} & \textbf{interdiction} & \en{You mustn't smoke here.} & tu ne dois surtout pas \tabularnewline
\en{don't have to} & \textbf{absence d'obligation} & \en{You don't have to pay.} & tu n'es pas obligé (mais tu peux) \tabularnewline
\en{should / ought to} & conseil, recommandation & \en{You should rest. You ought to call her.} & tu devrais \tabularnewline
\en{will} & futur, décision, promesse, habitude & \en{I'll do it. He will always help.} & je ferai \tabularnewline
\en{would} & conditionnel, habitude passée, politesse & \en{I would go if I could. Would you like tea?} & j'irais, voudriez-vous \tabularnewline
\en{shall} & offre ou suggestion (surtout \en{I/we}, britannique) & \en{Shall we dance? Shall I call him?} & veux-tu que\ldots ? \tabularnewline
\en{need} & nécessité (surtout négatif et interrogatif) & \en{You needn't worry. Need I come?} & tu n'as pas besoin \tabularnewline
\en{dare} & oser (littéraire) & \en{How dare you! I daren't tell her.} & comment oses-tu \tabularnewline
\end{longtable}
\end{tableau}

\begin{piege}[\en{mustn't} ne veut pas dire « ne pas être obligé »]
\en{You \textbf{mustn't} tell anyone.} = il est \textbf{interdit} de le dire.\\
\en{You \textbf{don't have to} tell anyone.} = tu n'es \textbf{pas obligé} de le dire.\\
Le français « tu ne dois pas » est ambigu ; l'anglais non. C'est une faute de sens très
coûteuse (consignes de sécurité, règles de travail, examens).
\end{piege}

\begin{piege}[\en{*I can to go}, \en{*he cans}, \en{*I am can}]
Trois fautes de débutant qui persistent : pas de \en{to} après un modal, pas de \en{-s},
pas de \en{be}. Et pour les temps où \en{can} n'existe pas, utilisez \en{be able to} :
\en{I will be able to} (futur), \en{I have been able to} (\en{perfect}),
\en{I would like to be able to} (conditionnel).
\end{piege}

\subsection{Les modaux de déduction : « ça doit être », « ça ne peut pas être »}
\begin{tableau}[Échelle de certitude]
\begin{tabularx}{\linewidth}{@{}p{3.2cm}p{2.4cm}X@{}}
\entete{Structure} & \entete{Certitude} & \entete{Exemple} \tabularnewline
\hline
\en{must} & 95\,\% (presque sûr) & \en{He must be at home --- his car is here.} \textit{(Il doit être chez lui.)} \tabularnewline
\en{will} & 90\,\% & \en{That will be the postman.} \textit{(Ce sera le facteur.)} \tabularnewline
\en{should} & 75\,\% & \en{She should be there by now.} \textit{(Elle devrait y être.)} \tabularnewline
\en{may} & 50\,\% & \en{It may be true.} \textit{(C'est peut-être vrai.)} \tabularnewline
\en{might / could} & 30\,\% & \en{He might be sleeping.} \textit{(Il dort peut-être.)} \tabularnewline
\en{can't} & 5\,\% (impossible) & \en{That can't be Ali --- he's in Paris!} \textit{(Ça ne peut pas être Ali.)} \tabularnewline
\end{tabularx}
\end{tableau}

\subsection{Les modaux au passé : \en{modal} + \en{have} + participe passé}
\begin{tableau}[Le regret, la déduction passée, la possibilité manquée]
\begin{tabularx}{\linewidth}{@{}p{3.6cm}p{4.4cm}X@{}}
\entete{Structure} & \entete{Sens} & \entete{Exemple} \tabularnewline
\hline
\en{must have done} & déduction sur le passé (c'est forcément arrivé) & \en{She must have forgotten.} \textit{(Elle a dû oublier.)} \tabularnewline
\en{can't have done} & impossibilité passée & \en{He can't have left already.} \textit{(Il ne peut pas être déjà parti.)} \tabularnewline
\en{may/might have done} & possibilité passée & \en{They might have missed the bus.} \textit{(Ils ont peut-être raté le bus.)} \tabularnewline
\en{should have done} & regret, reproche (j'aurais dû) & \en{I should have studied more.} \textit{(J'aurais dû étudier davantage.)} \tabularnewline
\en{shouldn't have done} & reproche inverse (je n'aurais pas dû) & \en{You shouldn't have said that.} \textit{(Tu n'aurais pas dû dire ça.)} \tabularnewline
\en{could have done} & possibilité non réalisée & \en{I could have won, but I fell.} \textit{(J'aurais pu gagner, mais je suis tombé.)} \tabularnewline
\en{needn't have done} & action faite alors qu'elle était inutile & \en{You needn't have cooked --- we ate out.} \textit{(Tu n'avais pas besoin de cuisiner.)} \tabularnewline
\end{tabularx}
\end{tableau}

\begin{astuce}[Les trois structures qui vous rendront immédiatement « naturel »]
\en{You should have told me!} (Tu aurais dû me le dire !) \quad
\en{It must have been hard.} (Ça a dû être dur.) \quad
\en{I could have done it.} (J'aurais pu le faire.)\\
Ces trois blocs, dits avec la bonne intonation, remplacent une heure de conversation
maladroite --- et expriment l'empathie, ce qui est le c\oe ur du contact humain.
\end{astuce}

\begin{exo}[1.1 --- Choisissez le modal]
\begin{enumerate}
  \item You \trou[1.8cm] smoke in the hospital. (interdiction)
  \item You \trou[1.8cm] pay, it's free. (absence d'obligation)
  \item \trou[1.8cm] I open the window? (permission polie)
  \item He \trou[1.8cm] be at home, the lights are on. (déduction forte)
  \item She \trou[1.8cm] come, but I'm not sure. (possibilité faible)
  \item You \trou[1.8cm] see a doctor. (conseil)
  \item When I was five, I \trou[1.8cm] already swim. (capacité passée)
  \item I \trou[1.8cm] to work every day --- it's the law. (obligation extérieure)
  \item I \trou[2.5cm] have called you yesterday, I'm sorry. (regret)
  \item They \trou[2.5cm] have missed the train --- they aren't here. (déduction passée)
\end{enumerate}
\end{exo}
\begin{corrige}[1.1]
1. \en{mustn't} 2. \en{don't have to} 3. \en{May / Can / Could} 4. \en{must}
5. \en{might / may / could} 6. \en{should} 7. \en{could} 8. \en{have}
9. \en{should have} 10. \en{must have}
\end{corrige}

\section{Unité 2 --- Les conditionnels}

\begin{tableau}[Les quatre conditionnels en un tableau]
\footnotesize
\begin{longtable}{@{}p{2.4cm}p{4.6cm}p{4.2cm}p{2.8cm}@{}}
\entete{Type} & \entete{Structure} & \entete{Exemple} & \entete{Emploi} \tabularnewline
\endfirsthead
\entete{Type} & \entete{Structure} & \entete{Exemple} & \entete{Emploi} \tabularnewline
\endhead
\hline
\textbf{0} & \en{if} + présent, présent & \en{If you heat water, it boils.} & vérité générale, fait scientifique \tabularnewline
\textbf{1} & \en{if} + présent, \en{will} + base & \en{If it rains, I'll stay home.} & situation réelle et probable dans l'avenir \tabularnewline
\textbf{2} & \en{if} + passé simple, \en{would} + base & \en{If I had time, I would travel.} & hypothèse improbable ou imaginaire au présent \tabularnewline
\textbf{3} & \en{if} + \en{past perfect}, \en{would have} + participe passé & \en{If I had studied, I would have passed.} & regret sur le passé, irréel \tabularnewline
\textbf{Mixte} & \en{if} + \en{past perfect}, \en{would} + base & \en{If I had taken that job, I would be rich now.} & cause passée, conséquence présente \tabularnewline
\end{longtable}
\end{tableau}

\begin{regle}[Le point qui fait toute la différence]
Dans la proposition en \en{if}, on emploie un temps \textbf{d'un cran en arrière} par rapport
au sens réel : présent pour du futur (type 1), passé pour du présent (type 2),
\en{past perfect} pour du passé (type 3). \textbf{Jamais de \en{will} après \en{if}}
(sauf dans \en{if you will} au sens de « si vous voulez bien »).
\end{regle}

\begin{piege}[\en{*If I will have time}]
\ko \en{If I will have time, I will call you.}\\
\ok \en{If I \textbf{have} time, I will call you.}\\
Même règle après \en{when, as soon as, unless, until, before, after} :
\en{When she arrives, we'll eat.} (pas \en{*when she will arrive}).
\end{piege}

\begin{astuce}[\en{If I were you\ldots} --- la formule de conseil]
\en{If I were you, I would take the job.} (Si j'étais toi, je prendrais le poste.)\\
C'est LE bloc à connaître pour donner un conseil avec élégance. \en{were} est la forme
correcte au subjonctif (toutes les personnes) ; \en{was} est accepté à l'oral.
\end{astuce}

\subsection{Les alternatives à \en{if}}
\begin{tableau}[Autres conjonctions conditionnelles]
\begin{tabularx}{\linewidth}{@{}p{3.4cm}p{4.4cm}X@{}}
\entete{Conjonction} & \entete{Sens} & \entete{Exemple} \tabularnewline
\hline
\en{unless} & à moins que (= \en{if\ldots not}) & \en{Unless you hurry, you'll be late.} \textit{(À moins que tu te dépêches, tu seras en retard.)} \tabularnewline
\en{as long as / so long as} & à condition que & \en{You can go, as long as you're back by ten.} \tabularnewline
\en{provided that / providing that} & à condition que (formel) & \en{We'll sign, provided that you agree.} \tabularnewline
\en{in case} & au cas où (précaution, \textbf{pas} une condition) & \en{Take an umbrella in case it rains.} \textit{(Prends un parapluie au cas où il pleuvrait.)} \tabularnewline
\en{otherwise / or} & sinon & \en{Hurry up, otherwise you'll miss it.} \tabularnewline
\en{even if} & même si & \en{Even if it rains, we'll go.} \tabularnewline
\en{suppose / supposing / what if} & et si\ldots ? & \en{Suppose you lost your job, what would you do?} \tabularnewline
\end{tabularx}
\end{tableau}

\begin{piege}[\en{in case} n'est pas « si »]
\en{I'll take an umbrella \textbf{in case} it rains.} = au cas où il pleuvrait (je me prépare).\\
\en{I'll take an umbrella \textbf{if} it rains.} = s'il pleut (condition).\\
Le français confond souvent les deux ; l'anglais les distingue nettement.
\end{piege}

\subsection{L'inversion conditionnelle (registre soutenu)}
\ex{Had I known, I would have come.}{Si j'avais su, je serais venu. (= \en{If I had known\ldots})}
\ex{Should you need help, please call me.}{Si vous avez besoin d'aide, appelez-moi. (= \en{If you should need\ldots})}
\ex{Were I richer, I would buy a house.}{Si j'étais plus riche, j'achèterais une maison.}

\begin{exo}[2.1 --- Mettez les verbes au bon temps]
\begin{enumerate}
  \item If it (rain) \trou[2cm] tomorrow, we (cancel) \trou[2.5cm] the match.
  \item If I (be) \trou[2cm] you, I (not / say) \trou[2.5cm] that.
  \item If she (study) \trou[2.5cm] harder, she (pass) \trou[2.5cm] the exam last year.
  \item I (buy) \trou[2cm] a car if I (have) \trou[2cm] enough money.
  \item If you heat ice, it (melt) \trou[2cm].
  \item Unless you (hurry) \trou[2cm], you (miss) \trou[2cm] the bus.
  \item If I (know) \trou[2.5cm] about the party, I (come) \trou[2.5cm].
  \item What (you / do) \trou[2.5cm] if you (win) \trou[2cm] the lottery?
  \item If he (not / forget) \trou[2.5cm] his passport, he (not / miss) \trou[2.5cm] the flight.
  \item Take a coat in case it (get) \trou[2cm] cold.
\end{enumerate}
\end{exo}
\begin{corrige}[2.1]
1. \en{rains} --- \en{will cancel} 2. \en{were / was} --- \en{wouldn't say}
3. \en{had studied} --- \en{would have passed} 4. \en{would buy} --- \en{had}
5. \en{melts} 6. \en{hurry} --- \en{will miss} 7. \en{had known} --- \en{would have come}
8. \en{would you do} --- \en{won} 9. \en{hadn't forgotten} --- \en{wouldn't have missed}
10. \en{gets}
\end{corrige}

\section{Unité 3 --- \en{wish}, \en{if only}, \en{would rather}, \en{it's time}}

\begin{tableau}[Exprimer le regret et le souhait]
\begin{tabularx}{\linewidth}{@{}p{3.6cm}p{4.2cm}X@{}}
\entete{Structure} & \entete{Sens} & \entete{Exemple} \tabularnewline
\hline
\en{wish} + passé simple & regret sur le \textbf{présent} & \en{I wish I \textbf{had} a car.} \textit{(J'aimerais avoir une voiture --- je n'en ai pas.)} \tabularnewline
\en{wish} + \en{past perfect} & regret sur le \textbf{passé} & \en{I wish I \textbf{had studied} more.} \textit{(J'aurais aimé étudier plus.)} \tabularnewline
\en{wish} + \en{would} & agacement, désir de changement chez l'autre & \en{I wish you \textbf{would stop} talking.} \textit{(J'aimerais que tu arrêtes de parler.)} \tabularnewline
\en{if only} + passé / \en{past perfect} & même sens, plus fort & \en{If only I \textbf{had} more time!} \tabularnewline
\en{would rather} + base & préférence (même sujet) & \en{I'd rather stay home.} \textit{(Je préférerais rester.)} \tabularnewline
\en{would rather} + sujet + passé & préférence (autre sujet) & \en{I'd rather you \textbf{didn't} smoke.} \textit{(Je préférerais que tu ne fumes pas.)} \tabularnewline
\en{it's (high) time} + passé & il serait temps que & \en{It's time we \textbf{left}.} \textit{(Il est temps qu'on parte.)} \tabularnewline
\en{hope} + présent / futur & espoir réel (pas de regret) & \en{I hope you \textbf{pass} the exam.} \textit{(J'espère que tu réussiras.)} \tabularnewline
\end{tabularx}
\end{tableau}

\begin{piege}[\en{wish} contre \en{hope}]
\en{I \textbf{hope} you feel better.} = j'espère que tu iras mieux (possible).\\
\en{I \textbf{wish} you felt better.} = j'aimerais que tu ailles mieux (mais ce n'est pas le cas).\\
\en{wish} implique toujours une dose d'irréel ou de regret. \en{hope}, lui, est tourné vers
l'avenir et s'emploie avec un temps normal. \en{*I wish you will come} est incorrect :
\en{I hope you will come} \ok\ ou \en{I wish you would come} \ok.
\end{piege}

\begin{exo}[3.1 --- Complétez]
\begin{enumerate}
  \item I wish I (be) \trou[2cm] taller. \item She wishes she (not / say) \trou[2.5cm] that.
  \item I wish it (stop) \trou[2cm] raining! \item I'd rather you (not / tell) \trou[2.5cm] anyone.
  \item It's time we (go) \trou[2cm] home. \item I hope you (have) \trou[2cm] a great trip.
  \item If only I (can) \trou[2.5cm] speak English perfectly.
  \item He wishes he (take) \trou[2.5cm] the other job last year.
\end{enumerate}
\end{exo}
\begin{corrige}[3.1]
1. \en{were / was} 2. \en{hadn't said} 3. \en{would stop} 4. \en{didn't tell}
5. \en{went} 6. \en{have / will have} 7. \en{could} 8. \en{had taken}
\end{corrige}

\section{Unité 4 --- La voix passive}

\begin{regle}[Formation]
\en{be} (au temps voulu) + \textbf{participe passé}. L'auteur de l'action passe au second plan
(avec \en{by}) ou disparaît.\\
\textbf{Actif :} \en{Somebody stole my bike.} \quad \textbf{Passif :} \en{My bike \textbf{was
stolen} (by somebody).}
\end{regle}

\begin{tableau}[Le passif à tous les temps]
\begin{tabularx}{\linewidth}{@{}p{3.6cm}p{5.0cm}X@{}}
\entete{Temps} & \entete{Passif} & \entete{Exemple} \tabularnewline
\hline
Présent simple & \en{am/is/are} + p.p. & \en{English is spoken here.} \tabularnewline
Présent continu & \en{am/is/are being} + p.p. & \en{The road is being repaired.} \tabularnewline
Passé simple & \en{was/were} + p.p. & \en{The house was built in 1920.} \tabularnewline
Passé continu & \en{was/were being} + p.p. & \en{The car was being washed.} \tabularnewline
\en{Present perfect} & \en{have/has been} + p.p. & \en{My phone has been stolen.} \tabularnewline
\en{Past perfect} & \en{had been} + p.p. & \en{The work had been finished.} \tabularnewline
Futur & \en{will be} + p.p. & \en{The results will be published tomorrow.} \tabularnewline
Modaux & \en{can/must/should be} + p.p. & \en{This must be signed. It can be done.} \tabularnewline
\en{going to} & \en{going to be} + p.p. & \en{A new bridge is going to be built.} \tabularnewline
Infinitif & \en{to be} + p.p. & \en{It needs to be cleaned.} \tabularnewline
\en{-ing} & \en{being} + p.p. & \en{I hate being interrupted.} \tabularnewline
\end{tabularx}
\end{tableau}

\subsection{Quand employer le passif ?}
\begin{enumerate}
  \item L'auteur est \textbf{inconnu} : \en{My wallet was stolen.}
  \item L'auteur est \textbf{évident ou sans importance} : \en{The rubbish is collected every Monday.}
  \item On veut rester \textbf{neutre ou poli} : \en{Mistakes were made.} (personne n'est accusé)
  \item Le résultat compte plus que l'acteur : \en{The bridge was opened in 2010.}
  \item Registre formel, administratif, scientifique, journalistique.
\end{enumerate}

\begin{piege}[Le français adore « on », l'anglais adore le passif]
\en{On m'a volé.} \ra \en{I was robbed.} (\en{*One stole me})\\
\en{On parle anglais ici.} \ra \en{English is spoken here.}\\
\en{On dit que c'est vrai.} \ra \en{It is said to be true.} / \en{They say it's true.}\\
\en{On a annulé la réunion.} \ra \en{The meeting was cancelled.}\\
Le pronom \en{one} existe en anglais mais il est littéraire et rare : ne l'utilisez pas
comme équivalent de « on ».
\end{piege}

\subsection{Verbes à deux compléments et passif impersonnel}
\ex{They gave him a prize. \ra\ He was given a prize. \ra\ A prize was given to him.}{On lui a donné un prix.}
\ex{They say that he is rich. \ra\ He is said to be rich. \ra\ It is said that he is rich.}{On dit qu'il est riche.}
\ex{People believe that the economy will improve. \ra\ The economy is believed to be improving.}{On pense que l'économie va s'améliorer.}

\subsection{\en{get} + participe passé : le passif familier}
\en{get married, get lost, get hurt, get fired, get arrested, get dressed, get broken,
get elected, get used to} --- très fréquent à l'oral, légèrement moins formel que \en{be}.
\ex{He got fired last week.}{Il s'est fait virer la semaine dernière.}

\subsection{\en{have something done} : faire faire quelque chose par un autre}
\begin{regle}
\en{have} + objet + \textbf{participe passé} = faire réaliser une action par quelqu'un d'autre.
\end{regle}
\ex{I had my hair cut.}{Je me suis fait couper les cheveux. (par un coiffeur)}
\ex{We're having the house painted next week.}{Nous faisons peindre la maison la semaine prochaine.}
\ex{I need to get my phone repaired.}{Je dois faire réparer mon téléphone. (\en{get} = variante familière)}

\begin{exo}[4.1 --- Mettez au passif]
\begin{enumerate}
  \item \en{They build these cars in Germany.}
  \item \en{Somebody has taken my pen.}
  \item \en{The storm destroyed the roof.}
  \item \en{People speak French in Senegal.}
  \item \en{They are repairing the road now.}
  \item \en{The company will launch a new product.}
  \item \en{You must sign this document.}
  \item \en{Someone told me the news.}
  \item \en{They say she is a genius.}
  \item \en{Somebody stole my bag on the bus.}
\end{enumerate}
\end{exo}
\begin{corrige}[4.1]
1. \en{These cars are built in Germany.} 2. \en{My pen has been taken.}
3. \en{The roof was destroyed (by the storm).} 4. \en{French is spoken in Senegal.}
5. \en{The road is being repaired now.} 6. \en{A new product will be launched.}
7. \en{This document must be signed.} 8. \en{I was told the news.}
9. \en{She is said to be a genius.} / \en{It is said that she is a genius.}
10. \en{My bag was stolen on the bus.}
\end{corrige}

\begin{exo}[4.2 --- \en{have something done}]
Traduisez :
\begin{enumerate}
  \item Je fais réparer ma voiture chaque année. \item Elle s'est fait voler son téléphone.
  \item Nous devons faire peindre la maison. \item Je me suis fait examiner les yeux hier.
\end{enumerate}
\end{exo}
\begin{corrige}[4.2]
1. \en{I have my car repaired every year.} 2. \en{She had her phone stolen.}
3. \en{We need to have the house painted.} 4. \en{I had my eyes tested yesterday.}
\end{corrige}

\section{Unité 5 --- Le discours rapporté (\en{reported speech})}

\begin{regle}[Le principe du « pas en arrière »]
Quand on rapporte des paroles au passé (\en{He said that\ldots}), le temps de la phrase
d'origine recule d'un cran, et les repères de temps et de lieu changent.
\end{regle}

\begin{tableau}[La transformation des temps]
\begin{tabularx}{\linewidth}{@{}p{5.4cm}X@{}}
\entete{Discours direct} & \entete{Discours rapporté} \tabularnewline
\hline
présent simple \ra & passé simple \tabularnewline
présent continu \ra & passé continu \tabularnewline
\en{present perfect} \ra & \en{past perfect} \tabularnewline
passé simple \ra & \en{past perfect} \tabularnewline
\en{will} \ra & \en{would} \tabularnewline
\en{can} \ra & \en{could} \tabularnewline
\en{may} \ra & \en{might} \tabularnewline
\en{must} / \en{have to} \ra & \en{had to} \tabularnewline
\en{should, could, might, would} \ra & inchangés \tabularnewline
\end{tabularx}
\end{tableau}

\begin{tableau}[Les changements de repères]
\begin{tabularx}{\linewidth}{@{}p{4.4cm}X@{}}
\entete{Direct} & \entete{Rapporté} \tabularnewline
\hline
\en{today / tonight} & \en{that day / that night} \tabularnewline
\en{yesterday} & \en{the day before / the previous day} \tabularnewline
\en{tomorrow} & \en{the next day / the following day} \tabularnewline
\en{now} & \en{then} \tabularnewline
\en{last week} & \en{the week before} \tabularnewline
\en{next month} & \en{the following month} \tabularnewline
\en{here} & \en{there} \tabularnewline
\en{this / these} & \en{that / those} \tabularnewline
\en{ago} & \en{before} \tabularnewline
\end{tabularx}
\end{tableau}

\subsection{\en{say} contre \en{tell}}
\begin{regle}
\en{say} + (proposition) : \en{He said (that) he was tired.}\\
\en{tell} + \textbf{personne} + (proposition) : \en{He told \textbf{me} (that) he was tired.}\\
Blocs figés : \en{tell the truth, tell a lie, tell a story, tell the time} ---
\en{say hello, say goodbye, say something, say a word}.
\end{regle}

\subsection{Rapporter une question}
\begin{regle}
\textbf{Question ouverte} (\en{wh-}) : on garde le mot interrogatif, puis ordre affirmatif.\\
\en{``Where do you live?''} \ra \en{He asked me where I lived.}\\
\textbf{Question fermée} : on introduit par \en{if} ou \en{whether}.\\
\en{``Are you married?''} \ra \en{She asked me if I was married.}
\end{regle}

\begin{piege}[\en{*He asked me where did I live}]
La faute la plus fréquente du discours rapporté : l'ordre interrogatif ne survit pas.
Après \en{asked}, c'est \textbf{sujet + verbe} : \en{where I lived}, \en{what she wanted},
\en{whether they had finished}.
\end{piege}

\subsection{Rapporter un ordre, une demande, un conseil}
\begin{regle}
Verbe rapporteur + \textbf{personne} + \en{to} + base (ou \en{not to} + base).
\end{regle}
\ex{``Close the door.'' \ra\ He told me to close the door.}{« Ferme la porte. » \ra Il m'a dit de fermer la porte.}
\ex{``Don't touch it.'' \ra\ She warned us not to touch it.}{« N'y touchez pas. » \ra Elle nous a prévenus de ne pas y toucher.}
\ex{``Please help me.'' \ra\ He asked me to help him.}{« Aide-moi, s'il te plaît. » \ra Il m'a demandé de l'aider.}

\subsection{Les verbes rapporteurs qui remplacent \en{say} (niveau B2)}
\begin{tableau}[Verbes rapporteurs et leurs structures]
\footnotesize
\begin{tabularx}{\linewidth}{@{}X X@{}}
\en{admit doing / that} --- admettre & \en{deny doing / that} --- nier \\
\en{promise to do} --- promettre & \en{refuse to do} --- refuser \\
\en{offer to do} --- proposer & \en{agree to do} --- accepter \\
\en{decide to do} --- décider & \en{plan to do} --- prévoir \\
\en{suggest doing / that} --- suggérer & \en{recommend doing} --- recommander \\
\en{apologise for doing} --- s'excuser de & \en{insist on doing} --- insister pour \\
\en{accuse someone of doing} --- accuser de & \en{blame someone for} --- reprocher à \\
\en{warn someone not to do} --- avertir de ne pas & \en{advise someone to do} --- conseiller de \\
\en{remind someone to do} --- rappeler de & \en{encourage someone to do} --- encourager à \\
\en{complain about} --- se plaindre de & \en{explain that / how} --- expliquer que \\
\end{tabularx}
\end{tableau}

\ex{``I'm sorry I'm late.'' \ra\ He apologised for being late.}{« Je suis désolé d'être en retard. » \ra Il s'est excusé d'être en retard.}
\ex{``You should see a doctor.'' \ra\ She advised me to see a doctor.}{« Tu devrais voir un médecin. » \ra Elle m'a conseillé de voir un médecin.}
\ex{``Don't forget to call me.'' \ra\ He reminded me to call him.}{« N'oublie pas de m'appeler. » \ra Il m'a rappelé de l'appeler.}
\ex{``Let's go out.'' \ra\ She suggested going out.}{« Sortons. » \ra Elle a proposé de sortir.}

\begin{astuce}[\en{that} est facultatif --- et souvent omis]
\en{He said (that) he was tired.} À l'oral, les natifs suppriment \en{that} la plupart du temps.
Le supprimer vous rend immédiatement plus naturel.
\end{astuce}

\begin{exo}[5.1 --- Rapportez ces phrases]
\begin{enumerate}
  \item \en{``I am very busy today.''} (he said)
  \item \en{``We have finished the work.''} (they told me)
  \item \en{``Will you help me?''} (she asked him)
  \item \en{``Where does she live?''} (I asked)
  \item \en{``Don't open the window.''} (the teacher told us)
  \item \en{``I will call you tomorrow.''} (he promised me)
  \item \en{``I saw her yesterday.''} (she said)
  \item \en{``Can you swim?''} (my mother asked me)
  \item \en{``You must leave now.''} (the guard told them)
  \item \en{``I'm sorry for the mistake.''} (he said)
\end{enumerate}
\end{exo}
\begin{corrige}[5.1]
1. \en{He said (that) he was very busy that day.}
2. \en{They told me (that) they had finished the work.}
3. \en{She asked him if / whether he would help her.}
4. \en{I asked where she lived.}
5. \en{The teacher told us not to open the window.}
6. \en{He promised to call me the next day.} / \en{He promised me that he would call me the following day.}
7. \en{She said (that) she had seen her the day before.}
8. \en{My mother asked me if / whether I could swim.}
9. \en{The guard told them (that) they had to leave then.}
10. \en{He apologised for the mistake.}
\end{corrige}

\begin{memo}[Chapitre 10 --- l'essentiel en 12 lignes]
\begin{enumerate}
  \item Après un modal : base verbale, jamais \en{to}, jamais \en{-s}.
  \item \en{mustn't} = interdit ; \en{don't have to} = pas obligé.
  \item \en{must be} = ça doit être ; \en{can't be} = c'est impossible.
  \item \en{should have done} = j'aurais dû.
  \item Type 1 : \en{if} + présent \ra \en{will}. Type 2 : \en{if} + passé \ra \en{would}.
  Type 3 : \en{if} + \en{past perfect} \ra \en{would have} + p.p.
  \item Jamais de \en{will} après \en{if}.
  \item \en{wish} + passé = regret présent ; \en{wish} + \en{past perfect} = regret passé.
  \item Passif = \en{be} + participe passé ; le français « on » se traduit souvent par le passif.
  \item \en{have something done} = faire faire par quelqu'un d'autre.
  \item Discours rapporté : un temps en arrière + changement des repères de temps.
  \item Après \en{asked}, ordre affirmatif : \en{where I lived} (pas \en{*where did I live}).
  \item Verbes rapporteurs riches : \en{promise to, admit doing, accuse of, advise to,
  apologise for, suggest doing}.
\end{enumerate}
\end{memo}

% =====================================================================
%  CHAPITRE 11 — GRAMMAIRE B2 : LA FINESSE
% =====================================================================
\chapter{Structures avancées : relatives, gérondif, comparaisons, emphase}

Vous savez former des phrases correctes. Ce chapitre vous apprend à former des phrases
\textbf{riches} --- celles qui font dire à un anglophone : « ton anglais est vraiment bon ».

\section{Unité 1 --- Les propositions relatives}

\begin{regle}[Les pronoms relatifs]
\en{who} (personnes, sujet) --- \en{whom} (personnes, complément, formel) ---
\en{whose} (possession) --- \en{which} (choses, animaux) --- \en{that} (personnes ou choses,
familier et très fréquent) --- \en{where} (lieu) --- \en{when} (temps) --- \en{why} (raison).
\end{regle}

\begin{tableau}[Relatives définissantes et non définissantes]
\begin{tabularx}{\linewidth}{@{}p{4.4cm}p{6.4cm}X@{}}
\entete{Type} & \entete{Règle} & \entete{Exemple} \tabularnewline
\hline
\textbf{Définissante} (essentielle) & Sans virgules. Elle identifie de qui ou de quoi on parle. \en{that} possible. Pronom complément omisssible. & \en{The man \textbf{who called} is my uncle.} / \en{The book \textbf{(that) I read} was great.} \tabularnewline
\textbf{Non définissante} (commentaire) & Entre virgules. Information supplémentaire. \en{that} \textbf{interdit}. Pronom jamais omis. & \en{My brother, \textbf{who lives in Dakar}, is a doctor.} \tabularnewline
\end{tabularx}
\end{tableau}

\begin{astuce}[Le pronom complément disparaît très souvent]
\en{The film (that) I saw yesterday} \ra \en{The film I saw yesterday}.
\en{The man (who) you met} \ra \en{The man you met}.
Cette omission est naturelle et fréquente --- mais impossible quand le pronom est sujet
(\en{*The man called is my uncle}).
\end{astuce}

\begin{tableau}[Les relatifs moins évidents]
\begin{tabularx}{\linewidth}{@{}p{3.4cm}p{4.4cm}X@{}}
\entete{Structure} & \entete{Emploi} & \entete{Exemple} \tabularnewline
\hline
\en{whose} & possession (personnes et choses) & \en{That's the girl \textbf{whose} father is a pilot.} \tabularnewline
\en{where} & lieu (= \en{in which}) & \en{This is the house \textbf{where} I was born.} \tabularnewline
\en{when} & temps (= \en{at which}) & \en{I remember the day \textbf{when} we met.} \tabularnewline
\en{why} & raison (après \en{reason}) & \en{That's the reason \textbf{why} I left.} \tabularnewline
\en{which} renvoyant à toute la phrase & commentaire sur ce qui précède & \en{He was late, \textbf{which} annoyed everyone.} \tabularnewline
\en{preposition + which/whom} & registre soutenu & \en{The city \textbf{in which} I live} = \en{The city \textbf{(that)} I live \textbf{in}} \tabularnewline
\end{tabularx}
\end{tableau}

\begin{piege}[Ne répétez pas le pronom]
\en{*The man who I met him yesterday} \ra \en{The man (who) I met yesterday}.
\en{*The book which I bought it} \ra \en{The book (which) I bought}.
Le relatif \emph{remplace} déjà le complément : en ajouter un second est la faute typique
des francophones qui traduisent mot à mot.
\end{piege}

\begin{exo}[1.1 --- Complétez avec \en{who}, \en{which}, \en{whose}, \en{where}, \en{when} ou rien]
\begin{enumerate}
  \item That's the woman \trou[1.5cm] helped me.
  \item The car \trou[1.5cm] I bought is blue.
  \item This is the school \trou[1.5cm] I studied.
  \item I'll never forget the day \trou[1.5cm] we met.
  \item He's the boy \trou[1.5cm] father is a minister.
  \item The film, \trou[1.5cm] was directed by a Senegalese, won a prize.
  \item The man \trou[1.5cm] you spoke to is my boss.
  \item She was late, \trou[1.5cm] made everyone angry.
\end{enumerate}
\end{exo}
\begin{corrige}[1.1]
1. \en{who / that} 2. \en{which / that / ---} (omis) 3. \en{where} 4. \en{when}
5. \en{whose} 6. \en{which} (non définissante : \en{that} interdit)
7. \en{who / whom / that / ---} 8. \en{which}
\end{corrige}

\section{Unité 2 --- Gérondif (\en{-ing}) ou infinitif (\en{to} + base) ?}

C'est l'un des points les plus subtils de l'anglais. Il existe trois situations.

\subsection{Cas 1 --- Verbes suivis toujours du gérondif}
\begin{tableau}[Verbes + \en{-ing} (liste essentielle)]
\footnotesize
\begin{tabularx}{\linewidth}{@{}X X X@{}}
\en{enjoy doing} & \en{avoid doing} & \en{finish doing} \\
\en{mind doing} & \en{keep doing} & \en{practise doing} \\
\en{suggest doing} & \en{consider doing} & \en{imagine doing} \\
\en{deny doing} & \en{admit doing} & \en{risk doing} \\
\en{miss doing} & \en{delay doing} & \en{can't help doing} \\
\en{give up doing} & \en{go on doing} & \en{feel like doing} \\
\en{it's no use doing} & \en{it's worth doing} & \en{look forward to doing} \\
\end{tabularx}
\end{tableau}
\ex{I enjoy learning English. She avoided answering. Do you mind waiting?}{J'aime apprendre l'anglais. Elle a évité de répondre. Ça t'ennuie d'attendre ?}
\ex{I'm looking forward to \textbf{seeing} you.}{J'attends de te voir avec impatience.}

\begin{piege}[\en{look forward to} + \en{-ing}]
Ici, \en{to} est une \textbf{préposition}, pas une marque d'infinitif : donc
\en{looking forward to meeting you} \ok, \en{*looking forward to meet you} \ko.
Même logique : \en{be used to doing}, \en{object to doing}, \en{dedicated to doing},
\en{when it comes to doing}.
\end{piege}

\subsection{Cas 2 --- Verbes suivis toujours de l'infinitif}
\begin{tableau}[Verbes + \en{to} + base (liste essentielle)]
\footnotesize
\begin{tabularx}{\linewidth}{@{}X X X@{}}
\en{want to do} & \en{decide to do} & \en{hope to do} \\
\en{plan to do} & \en{promise to do} & \en{refuse to do} \\
\en{agree to do} & \en{offer to do} & \en{manage to do} \\
\en{learn to do} & \en{would like to do} & \en{need to do} \\
\en{seem to do} & \en{appear to do} & \en{pretend to do} \\
\en{afford to do} & \en{choose to do} & \en{fail to do} \\
\en{expect to do} & \en{deserve to do} & \en{ask someone to do} \\
\end{tabularx}
\end{tableau}
\ex{I want to learn. He promised to come. We decided to leave.}{Je veux apprendre. Il a promis de venir. Nous avons décidé de partir.}

\subsection{Cas 3 --- Verbes qui changent de sens}
\begin{tableau}[Les huit verbes à double sens --- à connaître absolument]
\begin{tabularx}{\linewidth}{@{}p{2.4cm}p{4.2cm}p{4.2cm}X@{}}
\entete{Verbe} & \entete{+ \en{-ing}} & \entete{+ \en{to} + base} & \entete{Différence} \tabularnewline
\hline
\en{remember} & \en{I remember locking the door.} & \en{Remember to lock the door.} & souvenir du passé / tâche à faire \tabularnewline
\en{forget} & \en{I'll never forget meeting her.} & \en{Don't forget to call.} & souvenir / obligation \tabularnewline
\en{stop} & \en{He stopped smoking.} & \en{He stopped to smoke.} & arrêter une activité / s'arrêter pour faire autre chose \tabularnewline
\en{try} & \en{Try restarting the phone.} & \en{I tried to open it.} & expérimenter un moyen / faire un effort (souvent sans succès) \tabularnewline
\en{regret} & \en{I regret telling him.} & \en{I regret to inform you\ldots} & regretter un acte passé / annoncer poliment \tabularnewline
\en{go on} & \en{He went on talking.} & \en{He went on to explain.} & continuer la même chose / passer à autre chose \tabularnewline
\en{mean} & \en{Sorry, I didn't mean shouting.} & \en{I mean to finish tonight.} & --- / avoir l'intention \tabularnewline
\en{like} & \en{I like swimming.} (plaisir général) & \en{I like to swim in the morning.} (habitude choisie) & nuance très fine \tabularnewline
\end{tabularx}
\end{tableau}

\subsection{Après une préposition : toujours \en{-ing}}
\begin{regle}
Toute préposition (\en{about, after, before, for, of, without, instead of, because of,
interested in, good at, tired of, afraid of, sorry for, think about, dream of, succeed in,
believe in, insist on}) est suivie du gérondif.
\end{regle}
\ex{Thank you for coming. He left without saying goodbye. I'm interested in working with you.}{Merci d'être venu. Il est parti sans dire au revoir. Travailler avec vous m'intéresse.}

\subsection{L'infinitif de but et les autres constructions}
\begin{tableau}[Infinitif, \en{for}, \en{so that}]
\begin{tabularx}{\linewidth}{@{}p{4.2cm}X@{}}
\entete{Structure} & \entete{Exemple} \tabularnewline
\hline
\en{to} + base (but) & \en{I came \textbf{to help} you.} --- \textit{Je suis venu pour t'aider.} \tabularnewline
\en{in order to / so as to} (but, plus net) & \en{He studied hard \textbf{in order to pass}.} \tabularnewline
\en{so as not to} (but négatif) & \en{I whispered \textbf{so as not to wake} her.} \tabularnewline
\en{for} + nom (moyen) & \en{This knife is \textbf{for bread}.} \tabularnewline
\en{for} + personne + \en{to} + base & \en{It's hard \textbf{for me to understand}.} \tabularnewline
\en{so that} + proposition & \en{I called \textbf{so that} she \textbf{wouldn't} worry.} \tabularnewline
\en{too} + adjectif + \en{to} & \en{It's \textbf{too cold to swim}.} --- \textit{Trop froid pour nager.} \tabularnewline
\en{adjectif} + \en{enough} + \en{to} & \en{He's \textbf{old enough to drive}.} --- \textit{Assez grand pour conduire.} \tabularnewline
\end{tabularx}
\end{tableau}

\begin{piege}[\en{too} et \en{enough} : l'ordre change]
\en{too} \textbf{avant} l'adjectif : \en{too hot}. \en{enough} \textbf{après} l'adjectif :
\en{hot enough}. \en{*enough hot} est une faute immédiate et très visible.
\end{piege}

\begin{exo}[2.1 --- \en{-ing} ou \en{to} ?]
\begin{enumerate}
  \item I enjoy (listen) \trou[2cm] to music.
  \item She promised (help) \trou[2cm] me.
  \item Would you mind (open) \trou[2.5cm] the door?
  \item We decided (leave) \trou[2cm] early.
  \item He stopped (smoke) \trou[2cm] two years ago.
  \item He stopped (buy) \trou[2cm] some bread on his way home.
  \item Don't forget (send) \trou[2.5cm] the email.
  \item I'll never forget (visit) \trou[2cm] Paris.
  \item Thanks for (come) \trou[2cm] so quickly.
  \item It's worth (see) \trou[2cm].
  \item I'm looking forward to (hear) \trou[2cm] from you.
  \item She's very good at (explain) \trou[2cm] things.
  \item I tried (open) \trou[2cm] the window, but it was stuck.
  \item Try (press) \trou[2.5cm] the red button.
  \item I regret (say) \trou[2cm] those words.
\end{enumerate}
\end{exo}
\begin{corrige}[2.1]
1. \en{listening} 2. \en{to help} 3. \en{opening} 4. \en{to leave} 5. \en{smoking}
6. \en{to buy} 7. \en{to send} 8. \en{visiting} 9. \en{coming} 10. \en{seeing}
11. \en{hearing} 12. \en{explaining} 13. \en{to open} 14. \en{pressing}
15. \en{saying}
\end{corrige}

\section{Unité 3 --- Comparaison et superlatif}

\begin{tableau}[Formation régulière]
\begin{tabularx}{\linewidth}{@{}p{3.6cm}p{3.6cm}X@{}}
\entete{Type d'adjectif} & \entete{Comparatif} & \entete{Superlatif} \tabularnewline
\hline
1 syllabe & \en{-er} + \en{than} & \en{the} + \en{-est} \tabularnewline
& \en{taller than, faster than} & \en{the tallest, the fastest} \tabularnewline
1 syllabe, consonne-voyelle-consonne & on double & \en{bigger, the biggest} ; \en{hotter, the hottest} \tabularnewline
fin en \en{-e} & \en{-r} & \en{nicer, the nicest} \tabularnewline
consonne + \en{y} & \en{-ier} & \en{happier, the happiest} ; \en{easier, the easiest} \tabularnewline
2 syllabes et plus & \en{more} + adjectif + \en{than} & \en{the most} + adjectif \tabularnewline
& \en{more beautiful than} & \en{the most beautiful} \tabularnewline
Irréguliers & \en{good} \ra \en{better} ; \en{bad} \ra \en{worse} ; \en{far} \ra \en{further / farther} & \en{the best ; the worst ; the furthest} \tabularnewline
\end{tabularx}
\end{tableau}

\subsection{Les structures de comparaison utiles}
\begin{tableau}[Douze structures à maîtriser]
\begin{tabularx}{\linewidth}{@{}p{4.4cm}X@{}}
\entete{Structure} & \entete{Exemple et traduction} \tabularnewline
\hline
\en{as + adj + as} & \en{He is \textbf{as tall as} his father.} --- \textit{Il est aussi grand que son père.} \tabularnewline
\en{not as / not so + adj + as} & \en{It's \textbf{not as expensive as} I thought.} --- \textit{Ce n'est pas aussi cher que je pensais.} \tabularnewline
\en{much / far / a lot + comparatif} & \en{It's \textbf{much better}. She's \textbf{far more experienced}.} --- \textit{Bien meilleur. Bien plus expérimentée.} \tabularnewline
\en{a bit / slightly + comparatif} & \en{It's \textbf{a bit harder} than I expected.} --- \textit{Un peu plus difficile.} \tabularnewline
\en{the + comparatif\ldots, the + comparatif\ldots} & \en{\textbf{The more} you practise, \textbf{the better} you get.} --- \textit{Plus tu pratiques, mieux c'est.} \tabularnewline
\en{comparatif + and + comparatif} & \en{It's getting \textbf{colder and colder}.} --- \textit{Il fait de plus en plus froid.} \tabularnewline
\en{less + adjectif + than} & \en{This one is \textbf{less expensive than} that one.} \tabularnewline
\en{the least} & \en{the least interesting film} --- \textit{le film le moins intéressant} \tabularnewline
\en{the same as} & \en{Your phone is \textbf{the same as} mine.} --- \textit{Le même que le mien.} \tabularnewline
\en{different from} & \en{English is \textbf{different from} French.} (et non \en{*different than} en BrE) \tabularnewline
\en{as many / as much as} & \en{I don't eat \textbf{as much as} you.} / \en{There weren't \textbf{as many people as} last year.} \tabularnewline
\en{prefer A to B} & \en{I \textbf{prefer tea to coffee}.} / \en{I'd \textbf{rather stay than go}.} \tabularnewline
\end{tabularx}
\end{tableau}

\begin{piege}[\en{*more tall}, \en{*the most big}]
Un adjectif d'une syllabe prend \en{-er/-est}, jamais \en{more/most} :
\en{taller} \ok, \en{*more tall} \ko. Exception courante : on peut dire \en{more fun}
et \en{more clear} (acceptés), mais \en{funner} est évité.
\end{piege}

\subsection{Comparer des quantités et des verbes}
\ex{I earn more than you. She works harder than me. He drives more carefully than I do.}{Je gagne plus que toi. Elle travaille plus dur que moi. Il conduit plus prudemment que moi.}
\ex{There are fewer cars than yesterday but less traffic noise.}{Il y a moins de voitures qu'hier mais moins de bruit de circulation.}
\begin{memo}[\en{fewer} ou \en{less} ?]
\en{fewer} + dénombrable pluriel (\en{fewer people, fewer errors}) ;
\en{less} + indénombrable (\en{less water, less time, less money}).
\end{memo}

\begin{exo}[3.1 --- Comparez]
\begin{enumerate}
  \item Dakar is (big) \trou[2cm] than Thiès.
  \item This exercise is (difficult) \trou[2.5cm] than the last one.
  \item She is (good) \trou[2.5cm] at English than I am.
  \item It's (bad) \trou[2cm] day of my life.
  \item My car is not (expensive) \trou[3cm] as yours.
  \item The (far) \trou[2cm] you go, the (beautiful) \trou[2.5cm] it gets.
  \item He runs (fast) \trou[2.5cm] than his brother.
  \item I have (few / less) \trou[2cm] time than you.
  \item This is the (interesting) \trou[2.5cm] book I have ever read.
  \item The weather is getting (hot) \trou[2cm] and (hot) \trou[2cm].
\end{enumerate}
\end{exo}
\begin{corrige}[3.1]
1. \en{bigger} 2. \en{more difficult} 3. \en{better} 4. \en{the worst}
5. \en{as expensive} 6. \en{further} --- \en{more beautiful} 7. \en{faster}
8. \en{less} 9. \en{most interesting} 10. \en{hotter} --- \en{hotter}
\end{corrige}

\section{Unité 4 --- Les \en{question tags} (les petites questions en fin de phrase)}

\begin{regle}[Le mécanisme]
Phrase affirmative \ra tag \textbf{négatif} ; phrase négative \ra tag \textbf{affirmatif}.
On reprend \textbf{l'auxiliaire} de la phrase (ou \en{do/does/did} s'il n'y en a pas) et
le \textbf{pronom} correspondant au sujet.
\end{regle}

\begin{tableau}[Tags --- tableau complet]
\begin{tabularx}{\linewidth}{@{}p{5.4cm}X@{}}
\entete{Phrase} & \entete{Tag} \tabularnewline
\hline
\en{You're French,} & \en{aren't you?} \tabularnewline
\en{She isn't coming,} & \en{is she?} \tabularnewline
\en{They live here,} & \en{don't they?} \tabularnewline
\en{He works hard,} & \en{doesn't he?} \tabularnewline
\en{You went out,} & \en{didn't you?} \tabularnewline
\en{She can swim,} & \en{can't she?} \tabularnewline
\en{We should leave,} & \en{shouldn't we?} \tabularnewline
\en{You've finished,} & \en{haven't you?} \tabularnewline
\en{They'll come,} & \en{won't they?} \tabularnewline
\en{I'm late,} & \en{\textbf{aren't I}?} (irrégulier !) \tabularnewline
\en{Let's go,} & \en{\textbf{shall we}?} \tabularnewline
\en{Don't forget,} & \en{\textbf{will you}?} \tabularnewline
\en{Open the door,} & \en{\textbf{will you / can you}?} \tabularnewline
\en{Nobody called,} & \en{\textbf{did they}?} (sujet négatif \ra tag affirmatif, pronom \en{they}) \tabularnewline
\en{Everything is fine,} & \en{\textbf{isn't it}?} \tabularnewline
\end{tabularx}
\end{tableau}

\begin{astuce}[Deux intonations, deux sens]
Tag \textbf{descendant} \searrow\ = je suis presque sûr, je cherche simplement l'accord :
\en{It's hot today, isn't it?} \searrow\ (conversation normale).\\
Tag \textbf{montant} \nearrow\ = vraie question, je ne sais pas :
\en{You haven't seen my keys, have you?} \nearrow\\
C'est l'outil de \textbf{small talk} par excellence : \en{Nice weather, isn't it?}
\en{You're new here, aren't you?} \en{That was a great match, wasn't it?}
\end{astuce}

\section{Unité 5 --- \en{so}, \en{neither}, \en{either} : répondre en écho}

\begin{tableau}[Marquer l'accord et le désaccord]
\begin{tabularx}{\linewidth}{@{}p{5.0cm}X@{}}
\entete{Situation} & \entete{Réponse} \tabularnewline
\hline
\en{I like coffee.} (affirmatif) & \en{\textbf{So do I.}} --- \textit{Moi aussi.} \tabularnewline
\en{I'm tired.} (avec \en{be}) & \en{\textbf{So am I.}} \tabularnewline
\en{I can swim.} (avec modal) & \en{\textbf{So can I.}} \tabularnewline
\en{I've been there.} (avec \en{have}) & \en{\textbf{So have I.}} \tabularnewline
\en{I don't like tea.} (négatif) & \en{\textbf{Neither do I.}} / \en{Nor do I.} --- \textit{Moi non plus.} \tabularnewline
\en{I'm not hungry.} & \en{\textbf{Neither am I.}} \tabularnewline
\en{I can't come.} & \en{\textbf{Neither can I.}} \tabularnewline
Désaccord sur un affirmatif & \en{\textbf{I don't.}} / \en{Really? I do.} \tabularnewline
Désaccord sur un négatif & \en{\textbf{I do.}} / \en{Well, I am.} \tabularnewline
\en{either} en fin de négation & \en{I don't like it \textbf{either}.} --- \textit{Je ne l'aime pas non plus.} \tabularnewline
\end{tabularx}
\end{tableau}

\begin{piege}[\en{*Me too} après une négation]
\en{``I don't like fish.''} --- \en{*Me too.} \ko \ra \en{\textbf{Neither do I.}} ou
\en{\textbf{Me neither.}} (familier, très courant à l'oral).
\end{piege}

\section{Unité 6 --- Les quantifieurs avancés}

\begin{tableau}[Tous les quantifieurs, avec leur construction]
\footnotesize
\begin{longtable}{@{}p{3.0cm}p{2.6cm}p{4.4cm}p{3.6cm}@{}}
\entete{Quantifieur} & \entete{Avec} & \entete{Sens} & \entete{Exemple} \tabularnewline
\endfirsthead
\entete{Quantifieur} & \entete{Avec} & \entete{Sens} & \entete{Exemple} \tabularnewline
\endhead
\hline
\en{a few} & dénombrables & quelques (positif) & \en{I have a few friends here.} \tabularnewline
\en{few} & dénombrables & très peu (négatif) & \en{Few people know this.} \tabularnewline
\en{a little} & indénombrables & un peu (positif) & \en{I speak a little English.} \tabularnewline
\en{little} & indénombrables & très peu (négatif) & \en{There is little hope.} \tabularnewline
\en{several} & dénombrables & plusieurs & \en{several reasons} \tabularnewline
\en{many / much} & dénombr. / indénombr. & beaucoup (surtout négation, question) & \en{How many? How much?} \tabularnewline
\en{a lot of / lots of} & les deux & beaucoup (affirmatif naturel) & \en{a lot of money, a lot of cars} \tabularnewline
\en{plenty of} & les deux & largement assez & \en{There's plenty of time.} \tabularnewline
\en{enough} & les deux & assez & \en{enough food, enough chairs} \tabularnewline
\en{too much / too many} & indénombr. / dénombr. & trop & \en{too much noise, too many errors} \tabularnewline
\en{some / any} & les deux & du, des / (négation, question) & \en{some water, any idea?} \tabularnewline
\en{no / none} & --- & aucun & \en{I have no time.} / \en{None of us knew.} \tabularnewline
\en{all (of)} & les deux & tout, tous & \en{all the students, all of it} \tabularnewline
\en{both (of)} & deux & les deux & \en{both of my brothers} \tabularnewline
\en{either / neither} & deux & l'un ou l'autre / ni l'un ni l'autre & \en{either day is fine; neither is true} \tabularnewline
\en{each / every} & singulier & chaque / tous (sans exception) & \en{each student, every day} \tabularnewline
\en{most (of)} & les deux & la plupart & \en{most people, most of the time} \tabularnewline
\en{half of, a third of} & les deux & la moitié, un tiers & \en{half of the cake} \tabularnewline
\end{longtable}
\end{tableau}

\begin{memo}[\en{every} contre \en{each} contre \en{all}]
\en{Every day} = tous les jours, sans exception (vision d'ensemble).\\
\en{Each student} = chaque étudiant pris un par un (vision individuelle).\\
\en{All the students} = tous les étudiants (le groupe).\\
Tous trois se construisent avec un verbe au \textbf{singulier} pour \en{every} et \en{each} :
\en{Every child \textbf{has} a book.}
\end{memo}

\section{Unité 7 --- Les adverbes : place et formation}

\begin{tableau}[Formation et exceptions]
\begin{tabularx}{\linewidth}{@{}p{5.0cm}X@{}}
\entete{Règle} & \entete{Exemples} \tabularnewline
\hline
adjectif + \en{-ly} & \en{quick} \ra \en{quickly}, \en{careful} \ra \en{carefully}, \en{quiet} \ra \en{quietly} \tabularnewline
\en{-y} \ra \en{-ily} & \en{happy} \ra \en{happily}, \en{easy} \ra \en{easily}, \en{angry} \ra \en{angrily} \tabularnewline
\en{-le} \ra \en{-bly} & \en{terrible} \ra \en{terribly}, \en{simple} \ra \en{simply} \tabularnewline
\en{-ic} \ra \en{-ically} & \en{basic} \ra \en{basically}, \en{economic} \ra \en{economically} \tabularnewline
Irréguliers & \en{good} \ra \en{well} ; \en{fast} \ra \en{fast} ; \en{hard} \ra \en{hard} ; \en{late} \ra \en{late} \tabularnewline
Deux formes, deux sens & \en{hard} (avec effort) / \en{hardly} (presque pas) ; \en{late} (en retard) / \en{lately} (récemment) ; \en{near} (près) / \en{nearly} (presque) \tabularnewline
\end{tabularx}
\end{tableau}

\begin{piege}[\en{hardly} ne veut pas dire « durement »]
\en{He works \textbf{hard}.} = il travaille dur.\\
\en{He \textbf{hardly} works.} = il ne travaille presque pas.\\
Même piège : \en{lately} = récemment (pas « tardivement »), \en{nearly} = presque,
\en{eventually} = finalement (pas « éventuellement » !), \en{actually} = en fait (pas
« actuellement » !). Ces quatre adverbes sont des \textbf{faux amis} --- voir chapitre 13.
\end{piege}

\begin{tableau}[Place de l'adverbe]
\begin{tabularx}{\linewidth}{@{}p{3.4cm}p{4.4cm}X@{}}
\entete{Type} & \entete{Place} & \entete{Exemple} \tabularnewline
\hline
Fréquence & avant le verbe principal, après \en{be} & \en{She \textbf{often} travels. He is \textbf{always} late.} \tabularnewline
Manière & après le verbe ou l'objet & \en{He drove \textbf{carefully}. She speaks English \textbf{well}.} \tabularnewline
Lieu & après le verbe, avant le temps & \en{They waited \textbf{outside}.} \tabularnewline
Temps & en fin de phrase (ou en début pour insister) & \en{I saw him \textbf{yesterday}. \textbf{Yesterday} I saw him.} \tabularnewline
Degré & avant l'adjectif ou l'adverbe modifié & \en{\textbf{very} good, \textbf{really} quickly, \textbf{quite} interesting} \tabularnewline
Commentaire & en début de phrase & \en{\textbf{Unfortunately}, we lost. \textbf{Honestly}, I don't know.} \tabularnewline
\end{tabularx}
\end{tableau}

\begin{astuce}[L'ordre quand plusieurs adverbes coexistent]
\textbf{Manière \ra Lieu \ra Temps} (\en{MPT}) : \en{She sang \textbf{beautifully}
\textbf{at the concert} \textbf{last night}.} Mémorisez « MPT » et vous serez correct à
chaque fois.
\end{astuce}

\section{Unité 8 --- Mettre en relief : emphase, inversion, clivage}

\subsection{\en{do} emphatique}
\ex{I \textbf{do} like your dress! You \textbf{did} call me, didn't you? Do sit down.}{J'aime \emph{vraiment} ta robe ! Tu m'as \emph{bien} appelé, non ? Assieds-toi donc.}

\subsection{Les phrases clivées (\en{cleft sentences})}
\begin{tableau}[Trois structures pour insister]
\begin{tabularx}{\linewidth}{@{}p{4.4cm}X@{}}
\entete{Structure} & \entete{Exemple} \tabularnewline
\hline
\en{It is / was X that\ldots} & \en{\textbf{It was} Ali \textbf{that} broke the window.} --- \textit{C'est Ali qui a cassé la vitre.} \tabularnewline
\en{What I need is\ldots} & \en{\textbf{What I need} is a good rest.} --- \textit{Ce dont j'ai besoin, c'est de repos.} \tabularnewline
\en{The thing that\ldots / All I want is\ldots} & \en{\textbf{All I want} is some peace.} --- \textit{Tout ce que je veux, c'est la paix.} \tabularnewline
\en{What happened was that\ldots} & \en{\textbf{What happened was that} the power went off.} --- \textit{Ce qui s'est passé, c'est que le courant a sauté.} \tabularnewline
\end{tabularx}
\end{tableau}

\subsection{L'inversion après un adverbe négatif ou restrictif (registre soutenu)}
\begin{tableau}[Inversions à reconnaître et à oser]
\begin{tabularx}{\linewidth}{@{}p{4.6cm}X@{}}
\entete{Amorce} & \entete{Phrase} \tabularnewline
\hline
\en{Never} & \en{\textbf{Never have I} seen such a thing.} \tabularnewline
\en{Rarely / Seldom} & \en{\textbf{Rarely does} he arrive on time.} \tabularnewline
\en{Hardly\ldots when} & \en{\textbf{Hardly had} we left \textbf{when} it started to rain.} \tabularnewline
\en{No sooner\ldots than} & \en{\textbf{No sooner had} I sat down \textbf{than} the phone rang.} \tabularnewline
\en{Not only\ldots but also} & \en{\textbf{Not only is} she intelligent, \textbf{but} she \textbf{is also} kind.} \tabularnewline
\en{Under no circumstances} & \en{\textbf{Under no circumstances should} you open this door.} \tabularnewline
\en{Only after / Only when} & \en{\textbf{Only after} the meeting \textbf{did I} understand.} \tabularnewline
\en{Little} & \en{\textbf{Little did} he know what was coming.} \tabularnewline
\end{tabularx}
\end{tableau}

\begin{astuce}[Utilisez-en une seule, mais utilisez-la]
Dans un exposé ou un entretien, une inversion bien placée fait un effet remarquable :
\en{``Not only did we finish on time, but we also came in under budget.''}
(Nous avons non seulement fini à temps, mais en plus sous le budget.)
Une par discours suffit --- au-delà, cela devient affecté.
\end{astuce}

\section{Unité 9 --- Les articles de niveau avancé}

\begin{tableau}[Cas délicats, résolus]
\begin{tabularx}{\linewidth}{@{}p{5.4cm}X@{}}
\entete{Situation} & \entete{Article} \tabularnewline
\hline
Instruments de musique & \en{the} : \en{play \textbf{the} piano, the guitar} \tabularnewline
Sports et jeux & rien : \en{play football, chess, tennis} \tabularnewline
Repas & rien en général : \en{have breakfast} ; mais \en{a} si on le décrit : \en{We had \textbf{a} lovely dinner} \tabularnewline
Institutions au sens fonctionnel & rien : \en{go to school / to bed / to prison / to church / to hospital (BrE) / to work / home} \tabularnewline
Mêmes lieux, sens bâtiment & \en{the} : \en{I went to \textbf{the} school to meet the teacher.} \tabularnewline
Moyens de transport avec \en{by} & rien : \en{by car, by bus, by train, by plane} --- mais \en{on foot} \tabularnewline
Moyens avec un déterminant & \en{the / my} : \en{in \textbf{the} car, on \textbf{the} bus} \tabularnewline
Noms géographiques & rien : \en{Senegal, Africa, Lake Victoria, Mount Everest} --- \en{the} : \en{the Gambia, the USA, the Sahara, the Atlantic, the Alps, the Nile} \tabularnewline
Rues, places, aéroports, gares & rien : \en{Oxford Street, Times Square, Heathrow Airport, Victoria Station} \tabularnewline
\en{the} + adjectif = groupe & \en{the poor, the rich, the young, the unemployed} \tabularnewline
Noms abstraits généralisés & rien : \en{Life is hard. Love matters. Education counts.} \tabularnewline
\en{a} avec un sens de « par » & \en{twice \textbf{a} week, 50 kilometres \textbf{an} hour} \tabularnewline
\end{tabularx}
\end{tableau}

\section{Unité 10 --- Verbe + préposition : les combinaisons obligatoires}
En anglais, le choix de la préposition dépend du verbe, pas de la logique française.
Apprenez chaque bloc d'un seul tenant.

\begin{tableau}[Les 60 combinaisons verbe/préposition les plus utiles]
\footnotesize
\begin{longtable}{@{}p{4.6cm}p{4.6cm}p{4.6cm}@{}}
\entete{+ \en{about}} & \entete{+ \en{for}} & \entete{+ \en{to}} \tabularnewline
\endfirsthead
\entete{+ \en{about}} & \entete{+ \en{for}} & \entete{+ \en{to}} \tabularnewline
\endhead
\en{think about} & \en{wait for} & \en{listen to} \tabularnewline
\en{talk about} & \en{look for} & \en{belong to} \tabularnewline
\en{dream about / of} & \en{ask for} & \en{happen to} \tabularnewline
\en{worry about} & \en{pay for} & \en{speak to / with} \tabularnewline
\en{care about} & \en{apologise for} & \en{write to} \tabularnewline
\en{know about} & \en{hope for} & \en{reply to} \tabularnewline
\en{complain about} & \en{thank someone for} & \en{agree to / with} \tabularnewline
\en{argue about} & \en{blame someone for} & \en{object to} \tabularnewline
\entete{+ \en{at}} & \entete{+ \en{of}} & \en{look forward to} \tabularnewline
\en{look at} & \en{think of} & \en{used to} \tabularnewline
\en{laugh at} & \en{dream of} & \en{prefer A to B} \tabularnewline
\en{arrive at (lieu précis)} & \en{consist of} & \en{get used to} \tabularnewline
\en{be good at} & \en{be afraid of} & \en{introduce someone to} \tabularnewline
\en{be bad at} & \en{be proud of} & \entete{+ \en{with}} \tabularnewline
\en{be surprised at} & \en{be tired of} & \en{agree with (personne)} \tabularnewline
\en{shout at} & \en{be full of} & \en{deal with} \tabularnewline
\entete{+ \en{in}} & \en{be made of} & \en{argue with} \tabularnewline
\en{believe in} & \en{be aware of} & \en{live with} \tabularnewline
\en{arrive in (ville, pays)} & \en{die of} & \en{fill in / out} \tabularnewline
\en{succeed in} & \en{think of / about} & \en{compare with} \tabularnewline
\en{be interested in} & \en{remind someone of} & \en{help someone with} \tabularnewline
\en{take part in} & \entete{+ \en{from}} & \en{provide someone with} \tabularnewline
\en{spend money in / on} & \en{come from} & \en{be satisfied with} \tabularnewline
\en{result in} & \en{suffer from} & \en{be angry with (qqn)} \tabularnewline
\en{specialise in} & \en{differ from} & \en{be busy with} \tabularnewline
\en{be involved in} & \en{recover from} & \en{begin with} \tabularnewline
\en{depend on} & \en{protect someone from} & \en{stay with} \tabularnewline
\en{rely on} & \en{prevent someone from} & \en{share with} \tabularnewline
\en{insist on} & \en{graduate from} & \en{be pleased with} \tabularnewline
\en{congratulate someone on} & \en{hide from} & \en{work with} \tabularnewline
\en{work on} & \en{learn from} & \en{be careful with} \tabularnewline
\en{be based on} & \en{escape from} & \en{be strict with} \tabularnewline
\en{spend time on} & \en{be different from} & \en{be familiar with} \tabularnewline
\en{focus on} & \en{be absent from} & \en{be bored with} \tabularnewline
\en{be keen on} & \en{take something from} & \en{be honest with} \tabularnewline
\en{count on} & \en{be far from} & \en{be popular with} \tabularnewline
\end{longtable}
\end{tableau}

\begin{piege}[Trois prépositions que le français fait rater]
\en{*depend of} \ra \en{depend \textbf{on}} \qquad
\en{*listen something} \ra \en{listen \textbf{to}} something \qquad
\en{*discuss about} \ra \en{discuss something} (pas de préposition !)\\
Ajoutez : \en{*married with} \ra \en{married \textbf{to}} ; \en{*angry against} \ra
\en{angry \textbf{with}} (personne) / \en{angry \textbf{about}} (situation) ;
\en{*search something} \ra \en{search \textbf{for}} (ou \en{look for}).
\end{piege}

\begin{exo}[10.1 --- Prépositions]
Complétez : \en{about, at, for, from, in, of, on, to, with} (parfois rien) :
\begin{enumerate}
  \item It depends \trou[1.2cm] the weather.
  \item I'm waiting \trou[1.2cm] the bus.
  \item She's married \trou[1.2cm] a doctor.
  \item Let's discuss \trou[1.2cm] the problem tomorrow.
  \item I'm looking forward \trou[1.2cm] hearing from you.
  \item He's very good \trou[1.2cm] mathematics.
  \item I arrived \trou[1.2cm] Dakar \trou[1.2cm] Monday.
  \item Listen \trou[1.2cm] me, please.
  \item She's afraid \trou[1.2cm] spiders.
  \item I agree \trou[1.2cm] you completely.
  \item He was born \trou[1.2cm] a small village.
  \item Thanks \trou[1.2cm] your help.
\end{enumerate}
\end{exo}
\begin{corrige}[10.1]
1. \en{on} 2. \en{for} 3. \en{to} 4. rien 5. \en{to} 6. \en{at}
7. \en{in} --- \en{on} 8. \en{to} 9. \en{of} 10. \en{with} 11. \en{in} 12. \en{for}
\end{corrige}

\begin{exo}[Synthèse du chapitre 11 --- reformulez]
\begin{enumerate}
  \item \en{The woman is my aunt. She called you.} (avec \en{who})
  \item \en{I want to thank you. You helped me.} (avec \en{for})
  \item \en{It's too difficult. I can't understand it.} (avec \en{too\ldots to})
  \item \en{You practise a lot. You improve.} (avec \en{the more\ldots the more})
  \item \en{Ali broke the window.} (mise en relief)
  \item \en{I have never seen such a beautiful beach.} (inversion)
  \item \en{``I'm hungry.''} (rapporté : \en{she said})
  \item \en{She is tall. Her sister is tall too.} (avec \en{as\ldots as})
  \item \en{I don't like tea.} --- accord négatif
  \item \en{You are tired.} (avec un tag)
\end{enumerate}
\end{exo}
\begin{corrige}[Synthèse]
1. \en{The woman who called you is my aunt.}
2. \en{I want to thank you for helping me.}
3. \en{It's too difficult for me to understand.}
4. \en{The more you practise, the more you improve.}
5. \en{It was Ali who broke the window.}
6. \en{Never have I seen such a beautiful beach.}
7. \en{She said (that) she was hungry.}
8. \en{Her sister is as tall as she is.}
9. \en{Neither do I.} / \en{Me neither.}
10. \en{You're tired, aren't you?}
\end{corrige}

\begin{memo}[Les dix réflexes qui marquent le niveau B2]
\begin{enumerate}
  \item Une relative au lieu de deux phrases : \en{The man who called\ldots}
  \item \en{for + -ing} après un remerciement : \en{Thanks for coming.}
  \item \en{look forward to + -ing} dans chaque conclusion d'e-mail.
  \item \en{used to} pour les souvenirs : \en{I used to live\ldots}
  \item \en{should have} pour les regrets : \en{I should have told you.}
  \item Un tag en fin de phrase pour créer du lien : \en{It's nice, isn't it?}
  \item \en{So do I / Neither do I} au lieu de \en{me too / me neither} systématiques.
  \item \en{as\ldots as} et \en{much better than} pour nuancer.
  \item \en{It was X who\ldots} pour mettre en relief.
  \item Le passif quand l'auteur importe peu : \en{The meeting was cancelled.}
\end{enumerate}
\end{memo}


% ===== PARTIE IV — VOCABULAIRE =====================================
\part{Vocabulaire : les mots qui servent vraiment}
% =====================================================================
%  CHAPITRE 12 — VOCABULAIRE : STRATÉGIE ET THÈMES
% =====================================================================
\chapter{Vocabulaire : apprendre les mots qui servent}

\section{La stratégie : 3\,000 mots suffisent pour tout dire}
Les études de fréquence sont formelles :

\begin{tableau}[Couverture d'un texte ordinaire selon la taille du vocabulaire]
\begin{tabularx}{\linewidth}{@{}p{3.6cm}p{3.0cm}X@{}}
\entete{Mots connus} & \entete{Couverture} & \entete{Ce que vous pouvez faire} \tabularnewline
\hline
100 mots les plus fréquents & 50\,\% & Presque rien, mais c'est la structure de toute phrase. \tabularnewline
500 mots & 75\,\% & Phrases de survie, se débrouiller dans un magasin. \tabularnewline
1\,000 mots & 85\,\% & Conversation quotidienne simple. \tabularnewline
2\,000 mots & 90\,\% & Comprendre un journal, suivre une série. \tabularnewline
\textbf{3\,000 mots} & \textbf{95\,\%} & \textbf{Conversation courante complète. C'est l'objectif du mois 9.} \tabularnewline
5\,000 mots & 98\,\% & Lire un roman, travailler en anglais. \tabularnewline
20\,000 mots & 99\,\% & Niveau d'un natif cultivé --- inutile à viser. \tabularnewline
\end{tabularx}
\end{tableau}

\begin{astuce}[N'apprenez pas « des mots », apprenez des blocs]
\en{decision} seul ne sert à rien. \en{make a decision} sert immédiatement.
Apprenez toujours le mot \textbf{avec son compagnon} :\\
\en{make a decision, make a mistake, make progress, take a photo, take a break,
do your best, do the shopping, have a look, have a rest, get married, get lost,
pay attention, keep in touch, catch a cold, spend time, feel like doing}.
\end{astuce}

\section{Votre avantage secret : les mots que vous connaissez déjà}
Entre 30 et 45\,\% du vocabulaire anglais vient du français (souvent via le normand après
1066). Voici 120 mots anglais que vous connaissez \textbf{déjà}, à condition de les
prononcer à l'anglaise.

\begin{tableau}[120 mots transparents --- lisez-les à voix haute]
\footnotesize
\begin{multicols}{3}
\begin{anglais}
\begin{itemize}\setlength{\itemsep}{1pt}
\item able --- capable
\item accept --- accepter
\item accident --- accident
\item action --- action
\item active --- actif
\item address --- adresse
\item admirable
\item adventure --- aventure
\item advice --- conseil
\item age --- âge
\item agree --- être d'accord
\item almost --- presque
\item already --- déjà
\item animal --- animal
\item announce --- annoncer
\item answer --- réponse
\item apart --- à part
\item apple --- pomme (piège !)
\item argument --- argument
\item arrive --- arriver
\item article --- article
\item attention --- attention
\item attitude
\item automatic
\item available --- disponible
\item average --- moyen
\item avoid --- éviter
\item balance --- équilibre
\item bank --- banque
\item base --- base
\item beautiful --- beau
\item because --- parce que
\item begin --- commencer
\item believe --- croire
\item benefit --- bénéfice
\item bicycle --- bicyclette
\item brave --- courageux
\item cable
\item calendar --- calendrier
\item calm --- calme
\item camera --- appareil photo
\item capital --- capitale
\item careful --- prudent
\item central --- central
\item certain --- certain
\item chance --- chance
\item change --- changer
\item chapter --- chapitre
\item charity --- charité
\item circle --- cercle
\item city --- cité, ville
\item civil --- civil
\item clear --- clair
\item climate --- climat
\item collection
\item colour / color
\item comfortable
\item common --- commun
\item company --- entreprise
\item compare --- comparer
\item complete --- complet
\item complicated
\item condition
\item confident
\item consequence
\item contact
\item content --- contenu
\item continue --- continuer
\item control --- contrôler
\item conversation
\item copy --- copie
\item correct --- correct
\item country --- pays
\item courage
\item create --- créer
\item crisis --- crise
\item critic --- critique
\item culture
\item curious --- curieux
\item dance --- danse
\item danger
\item decide --- décider
\item declare --- déclarer
\item decorate --- décorer
\item defence / defense
\item delicate
\item demand --- demande
\item democracy
\item demonstrate
\item department
\item depend --- dépendre
\item describe --- décrire
\item desert --- désert
\item desire --- désir
\item detail
\item develop
\item dictionary
\item difference
\item difficult
\item digest
\item direct --- direct
\item disappear
\item discover
\item discuss
\item disease --- maladie (piège !)
\item distance
\item distribute
\item document
\item domestic
\item dominant
\item donate
\item double
\item dramatic
\item economy
\item effect --- effet
\item efficient
\item effort
\item elegant
\item element
\item emotion
\item employ --- employer
\item empty --- vide (piège !)
\item energy
\item enormous
\item enter --- entrer
\item environment
\item equal --- égal
\item essential
\item establish
\item estimate
\item exact --- exact
\item examine
\item excellent
\item except --- sauf
\item excuse --- excuse
\item exercise --- exercice
\item exist
\item experience
\item explain
\item express
\item extraordinary
\item extreme
\item fabric --- tissu (piège !)
\item face --- visage
\item fact --- fait
\item factor
\item family
\item famous --- célèbre
\item favourite
\item festival
\item figure
\item final --- final
\item finance
\item firm --- entreprise
\item flexible
\item fluent
\item focus
\item force --- force
\item forest --- forêt
\item forget --- oublier (piège !)
\item formal --- formel
\item fortune
\item fragile
\item free --- libre / gratuit
\item frequent
\item fresh --- frais
\item front --- devant
\item fruit --- fruit
\item function
\item fundamental
\item furniture --- meubles
\item future --- futur
\item gallery
\item garden --- jardin
\item general
\item generous
\item gentle --- doux
\item gesture --- geste
\item global
\item government
\item grateful --- reconnaissant
\item group --- groupe
\item guarantee
\item guide
\item habit --- habitude
\item harmonious
\item honest --- honnête
\item honour --- honneur
\item horrible
\item hospital
\item hotel
\item humour
\item idea --- idée
\item identical
\item ignore --- ignorer
\item illegal
\item illustrate
\item imagine
\item immediate
\item immense
\item immigrant
\item immune
\item impact
\item impatient
\item important
\item impossible
\item impression
\item improve --- améliorer
\item include
\item increase
\item incredible
\item independent
\item individual
\item industrial
\item influence
\item inform
\item information
\item ingredient
\item inhabit --- habiter (piège !)
\item initial
\item innocent
\item insist
\item inspiration
\item install
\item instant
\item instinct
\item institution
\item instruct
\item instrument
\item insurance --- assurance
\item intelligent
\item intense
\item intention
\item interest
\item international
\item interpret
\item interrupt
\item interview
\item intimate
\item introduce
\item invent
\item invest
\item invisible
\item invitation
\item irregular
\item isolate
\item issue --- problème
\item January --- janvier
\item journey --- voyage
\item joy --- joie
\item judge --- juge
\item juice --- jus
\item July --- juillet
\item June --- juin
\item junior
\item justice
\item justify
\item keen --- passionné (piège !)
\end{itemize}
\end{anglais}
\end{multicols}
\end{tableau}

\begin{piege}[Les transparents trompeurs]
Dans cette liste, certains mots sont des \textbf{faux amis} partiels : \en{disease}
(maladie, pas « désaise »), \en{empty} (vide), \en{fabric} (tissu), \en{forget} (oublier),
\en{inhabit} (habiter un lieu, pas « habiter » au sens de vivre : on dit \en{live in}),
\en{keen} (passionné), \en{free} (libre \emph{et} gratuit), \en{issue} (problème, pas
« issue » de secours --- \en{exit}). Le chapitre 13 liste les 150 faux amis complets.
\end{piege}

\section{Les familles de mots : un mot appris = cinq mots gagnés}
\begin{tableau}[Exemples de familles]
\footnotesize
\begin{longtable}{@{}p{2.6cm}p{3.2cm}p{3.2cm}p{3.0cm}p{2.0cm}@{}}
\entete{Verbe} & \entete{Nom (personne)} & \entete{Nom (chose/action)} & \entete{Adjectif} & \entete{Adverbe} \tabularnewline
\endfirsthead
\entete{Verbe} & \entete{Nom (pers.)} & \entete{Nom (chose)} & \entete{Adjectif} & \entete{Adverbe} \tabularnewline
\endhead
\hline
\en{decide} & \en{decision-maker} & \en{decision} & \en{decisive, decided} & \en{decisively} \tabularnewline
\en{develop} & \en{developer} & \en{development} & \en{developed, developing} & --- \tabularnewline
\en{employ} & \en{employer, employee} & \en{employment} & \en{employable, unemployed} & --- \tabularnewline
\en{explain} & --- & \en{explanation} & \en{explanatory} & --- \tabularnewline
\en{govern} & \en{governor} & \en{government} & \en{governmental} & --- \tabularnewline
\en{inform} & --- & \en{information} & \en{informative, informed} & --- \tabularnewline
\en{manage} & \en{manager} & \en{management} & \en{manageable} & --- \tabularnewline
\en{organise} & \en{organiser} & \en{organisation} & \en{organised} & --- \tabularnewline
\en{produce} & \en{producer} & \en{production, product} & \en{productive} & \en{productively} \tabularnewline
\en{succeed} & \en{successor} & \en{success} & \en{successful} & \en{successfully} \tabularnewline
\en{teach} & \en{teacher} & \en{teaching} & --- & --- \tabularnewline
\en{understand} & --- & \en{understanding} & \en{understandable} & \en{understandably} \tabularnewline
\end{longtable}
\end{tableau}

\begin{tableau}[Préfixes utiles (ils doublent votre vocabulaire)]
\begin{tabularx}{\linewidth}{@{}p{2.6cm}p{3.0cm}X@{}}
\entete{Préfixe} & \entete{Sens} & \entete{Exemples} \tabularnewline
\hline
\en{un-} & contraire & \en{unhappy, unsure, unfair, unusual, unemployed} \tabularnewline
\en{in- / im- / il- / ir-} & contraire & \en{impossible, incorrect, illegal, irregular, informal} \tabularnewline
\en{dis-} & contraire & \en{disagree, disappear, dislike, dishonest} \tabularnewline
\en{mis-} & mal & \en{mistake, misunderstand, misspell, mislead} \tabularnewline
\en{re-} & à nouveau & \en{rewrite, rebuild, redo, return, recycle} \tabularnewline
\en{pre-} & avant & \en{prepare, prevent, predict, preview} \tabularnewline
\en{over-} & trop & \en{overwork, overeat, overestimate, overtime} \tabularnewline
\en{under-} & pas assez, dessous & \en{underpaid, underestimate, underwear, underground} \tabularnewline
\en{out-} & dépasser & \en{outrun, outnumber, output} \tabularnewline
\en{co-} & ensemble & \en{cooperate, coexist, co-author} \tabularnewline
\en{self-} & soi-même & \en{self-confidence, self-employed, selfish} \tabularnewline
\en{anti-} & contre & \en{antisocial, antibiotic} \tabularnewline
\en{non-} & non & \en{non-stop, nonsense, non-smoker} \tabularnewline
\end{tabularx}
\end{tableau}

\begin{tableau}[Suffixes utiles]
\begin{tabularx}{\linewidth}{@{}p{2.8cm}p{3.4cm}X@{}}
\entete{Suffixe} & \entete{Crée} & \entete{Exemples} \tabularnewline
\hline
\en{-er / -or} & personne qui fait & \en{teacher, worker, actor, doctor, writer} \tabularnewline
\en{-ist} & spécialiste & \en{scientist, artist, dentist, economist} \tabularnewline
\en{-ian} & métier, origine & \en{musician, politician, Nigerian, Ghanaian} \tabularnewline
\en{-ee} & personne qui reçoit & \en{employee, trainee, refugee} \tabularnewline
\en{-ment} & nom d'action & \en{development, agreement, movement} \tabularnewline
\en{-tion / -sion} & nom d'action & \en{education, decision, information} \tabularnewline
\en{-ness} & nom d'état & \en{happiness, kindness, darkness, illness} \tabularnewline
\en{-ity} & nom de qualité & \en{ability, possibility, university} \tabularnewline
\en{-ance / -ence} & nom d'état & \en{importance, difference, patience} \tabularnewline
\en{-ful} & plein de & \en{beautiful, useful, careful, helpful} \tabularnewline
\en{-less} & sans & \en{useless, careless, homeless, wireless} \tabularnewline
\en{-able / -ible} & qui peut être & \en{possible, comfortable, available, reliable} \tabularnewline
\en{-ous} & qui a la qualité & \en{dangerous, famous, serious, generous} \tabularnewline
\en{-y} & caractérisé par & \en{rainy, sunny, dirty, happy, tasty} \tabularnewline
\en{-al} & relatif à & \en{national, personal, natural, musical} \tabularnewline
\en{-ish} & approximatif, nationalité & \en{reddish, thirtyish, British, childish} \tabularnewline
\en{-ly} & adverbe & \en{quickly, carefully, usually} \tabularnewline
\end{tabularx}
\end{tableau}

\section{Vocabulaire thématique}
Vingt thèmes, les mots dont vous aurez réellement besoin. Chaque entrée donne la
\textbf{prononciation à la française} --- dites-la à voix haute en lisant le tableau.
Créez vos cartes Anki au fur et à mesure, 15 mots par jour maximum.

\subsection{Thème 1 --- Se présenter, identité, état civil}
\begin{tabvocab}
\mot{first name}{feust né-im}{prénom --- \en{What's your first name?}}
\mot{last name / surname}{laast né-im}{nom de famille}
\mot{full name}{foul né-im}{nom complet}
\mot{date of birth}{dé-it euv beuss}{date de naissance}
\mot{place of birth}{pléiss euv beuss}{lieu de naissance}
\mot{age}{é-idj}{âge --- \en{How old are you?} / \en{I'm 30.}}
\mot{nationality}{na-cheu-na-LI-ti}{nationalité --- \en{Senegalese, French, Nigerian}}
\mot{address}{e-DRESS}{adresse (accent sur la 2\textsuperscript{e} syllabe !)}
\mot{phone number}{foun NEUM-beu}{numéro de téléphone}
\mot{email}{ii-mé-il}{e-mail --- \en{email address}}
\mot{marital status}{MA-ri-tl STÉ-i-teuss}{situation matrimoniale}
\mot{single}{si-ngl}{célibataire}
\mot{married}{MA-rid}{marié(e)}
\mot{divorced}{di-VOOSST}{divorcé(e)}
\mot{widow / widower}{ouii-dou}{veuve / veuf}
\mot{husband}{HOUZ-beund}{mari}
\mot{wife}{ouaïf}{femme (épouse) --- pluriel : \en{wives}}
\mot{job}{djob}{emploi, métier --- \en{What's your job?}}
\mot{occupation}{o-kieu-PÉ-i-cheun}{profession (formel)}
\mot{student}{STiou-d'nt}{étudiant, élève}
\mot{pupil}{piou-pl}{élève (école primaire, BrE)}
\mot{ID card}{aï-di kaad}{carte d'identité}
\mot{passport}{PAASS-poort}{passeport}
\mot{signature}{SIG-neu-tcheu}{signature}
\mot{to spell}{spèl}{épeler --- \en{How do you spell your name?}}
\end{tabvocab}

\subsection{Thème 2 --- La famille et les relations}
\begin{tabvocab}
\mot{parents}{PÉ-e-reunts}{parents --- \en{my mother and father}}
\mot{mother / mom / mum}{MO-dheu}{mère, maman}
\mot{father / dad}{FAA-dheu}{père, papa}
\mot{son}{seun}{fils}
\mot{daughter}{DOO-teu}{fille (enfant) --- le \en{gh} est muet}
\mot{child / children}{tchaïld / TCHIL-dren}{enfant / enfants}
\mot{brother}{BRA-dheu}{frère}
\mot{sister}{SI-steu}{s\oe ur}
\mot{uncle}{EUN-kl}{oncle}
\mot{aunt}{aant}{tante}
\mot{cousin}{KEU-z'n}{cousin, cousine}
\mot{nephew}{NÉ-fiou}{neveu}
\mot{niece}{niiss}{nièce}
\mot{grandmother}{GRAND-mo-dheu}{grand-mère}
\mot{grandfather}{GRAND-faa-dheu}{grand-père}
\mot{grandson / granddaughter}{GRAND-seun}{petit-fils / petite-fille}
\mot{in-laws}{in-looz}{la belle-famille --- \en{mother-in-law, brother-in-law}}
\mot{stepfather / stepmother}{STÈP-faa-dheu}{beau-père / belle-mère (remariage)}
\mot{half-brother}{haaf-BRA-dheu}{demi-frère}
\mot{twin}{touin}{jumeau --- \en{twin brother}}
\mot{relative}{RÉ-leu-tiv}{membre de la famille, parent}
\mot{friend}{frènd}{ami --- \en{best friend, close friend}}
\mot{neighbour}{NÉ-i-beu}{voisin}
\mot{colleague}{KO-liig}{collègue}
\mot{boyfriend / girlfriend}{BOÏ-frend}{petit ami / petite amie}
\mot{fiancé(e)}{fi-ON-sé-i}{fiancé(e)}
\mot{to get married}{guèt MA-rid}{se marier}
\mot{to give birth}{guiv beuss}{accoucher, donner naissance}
\mot{to bring up / raise}{bri-ng UP}{élever (des enfants)}
\mot{to take care of}{té-ik KÉ-eu euv}{s'occuper de}
\end{tabvocab}

\subsection{Thème 3 --- Le corps, la santé, le médecin}
\begin{tabvocab}
\mot{head}{hèd}{tête --- \en{I have a headache}}
\mot{hair}{hè-eu}{cheveux (indénombrable) --- \en{long hair}}
\mot{face}{fé-iss}{visage}
\mot{eye / eyes}{aï}{\oe il / yeux --- \en{brown eyes}}
\mot{ear}{i-eu}{oreille}
\mot{nose}{nôouz}{nez}
\mot{mouth}{maouth}{bouche --- pluriel \en{mouths} \pron{maou-dhz}}
\mot{tooth / teeth}{touuth / tiith}{dent / dents}
\mot{tongue}{teung}{langue (organe)}
\mot{lip}{lip}{lèvre}
\mot{neck}{nèk}{cou}
\mot{shoulder}{CHOUL-deu}{épaule}
\mot{arm}{aam}{bras}
\mot{hand}{hand}{main}
\mot{finger}{FI-ngheu}{doigt}
\mot{chest}{tchèst}{poitrine}
\mot{back}{bak}{dos}
\mot{stomach}{STO-meuk}{ventre, estomac}
\mot{leg}{lèg}{jambe}
\mot{knee}{nii}{genou --- le \en{k} est muet}
\mot{foot / feet}{fout / fiit}{pied / pieds}
\mot{skin}{skine}{peau}
\mot{blood}{bleud}{sang}
\mot{bone}{bôoun}{os}
\mot{heart}{haat}{c\oe ur (organe) --- \en{heart attack}}
\mot{ill / sick}{il / sik}{malade}
\mot{pain}{pé-in}{douleur --- \en{a pain in my back}}
\mot{headache}{HÈ-dé-ik}{mal de tête}
\mot{stomachache}{STO-meuk-é-ik}{mal de ventre}
\mot{fever / temperature}{FII-veu}{fièvre --- \en{I have a fever}}
\mot{cough}{kof}{toux --- \en{to cough} \pron{kof}}
\mot{cold}{kôould}{rhume --- \en{to catch a cold} ; aussi : froid}
\mot{flu}{flou}{grippe}
\mot{sore throat}{soo THRO-OUT}{mal de gorge}
\mot{runny nose}{RA-ni nôouz}{nez qui coule}
\mot{dizzy}{DI-zi}{qui tourne la tête --- \en{I feel dizzy}}
\mot{tired}{TAÏ-eud}{fatigué}
\mot{injury / wound}{IN-djeu-ri}{blessure}
\mot{to bleed}{bliid}{saigner}
\mot{medicine}{MÉ-deu-sine}{médicament, médecine}
\mot{pill / tablet}{pil}{comprimé}
\mot{prescription}{pri-SKRIP-cheun}{ordonnance}
\mot{pharmacy / chemist}{FAA-meu-si}{pharmacie}
\mot{doctor / GP}{DOOK-teu}{médecin --- \en{general practitioner}}
\mot{nurse}{neuss}{infirmier, infirmière}
\mot{hospital}{HOSS-pi-tl}{hôpital}
\mot{appointment}{e-POÏNT-m'nt}{rendez-vous --- \en{to make an appointment}}
\mot{surgery}{SEU-djeu-ri}{cabinet médical (BrE) ; chirurgie}
\mot{treatment}{TRIIT-m'nt}{traitement}
\mot{vaccine}{VA-siinn}{vaccin}
\mot{health insurance}{hèlth in-Chou-reuns}{assurance maladie}
\mot{to recover}{ri-KEU-veu}{guérir, se rétablir}
\end{tabvocab}

\begin{astuce}[Les dix phrases médicales qui sauvent]
\en{I don't feel well.} (Je ne me sens pas bien.) \quad
\en{I have a headache / a fever / a sore throat.} \quad
\en{It hurts here.} (Ça fait mal ici.) \quad
\en{I'm allergic to\ldots} \quad
\en{I've been feeling ill for three days.} (Je suis malade depuis trois jours.) \quad
\en{I need to see a doctor.} \quad
\en{Where is the nearest pharmacy?} \quad
\en{Is it serious?} \quad
\en{What should I do?}
\end{astuce}

\subsection{Thème 4 --- La maison et les objets du quotidien}
\begin{tabvocab}
\mot{house / home}{haouss / hâam}{maison (bâtiment) / chez soi}
\mot{flat / apartment}{flatt / e-PAAUT-m'nt}{appartement (BrE / AmE)}
\mot{room}{rouum}{pièce}
\mot{living room}{LI-ving rouum}{salon}
\mot{bedroom}{BÈD-roum}{chambre}
\mot{kitchen}{KI-tchin}{cuisine}
\mot{bathroom}{BAAth-roum}{salle de bains}
\mot{toilet}{TO-i-l't}{toilettes --- \en{Where is the toilet?}}
\mot{hall / entrance}{hooll}{entrée, couloir}
\mot{balcony}{BA-l-keu-ni}{balcon}
\mot{yard / garden}{iaad / GAA-d'n}{cour / jardin}
\mot{roof}{rouuf}{toit}
\mot{wall}{ouool}{mur}
\mot{floor}{floo}{sol, étage --- \en{on the ground floor}}
\mot{ceiling}{SII-ling}{plafond}
\mot{door}{doo}{porte}
\mot{window}{OUIN-dou}{fenêtre}
\mot{stairs}{stè-euz}{escalier --- \en{go up the stairs}}
\mot{lift / elevator}{lift / É-leu-vé-i-teu}{ascenseur (BrE / AmE)}
\mot{key}{kii}{clé}
\mot{lock}{lok}{serrure ; \en{to lock} verrouiller}
\mot{furniture}{FEU-ni-tcheu}{meubles (indénombrable)}
\mot{table}{TÉ-i-bl}{table}
\mot{chair}{tchè-eu}{chaise}
\mot{armchair}{AArm-tchè-eu}{fauteuil}
\mot{sofa / couch}{SÔ-fa}{canapé}
\mot{bed}{bèd}{lit --- \en{to make the bed}}
\mot{cupboard}{KEU-beud}{placard --- le \en{p} est muet}
\mot{shelf / shelves}{chèlf}{étagère}
\mot{desk}{dèsk}{bureau (meuble)}
\mot{mirror}{MI-reu}{miroir}
\mot{lamp}{lamp}{lampe}
\mot{light / lamp}{laït}{lumière / luminaire}
\mot{fridge}{fridj}{réfrigérateur --- \en{refrigerator}}
\mot{freezer}{FRII-zeu}{congélateur}
\mot{oven}{EUV'n}{four}
\mot{stove / cooker}{stôouv}{cuisinière}
\mot{kettle}{KÈ-tl}{bouilloire}
\mot{sink}{sink}{évier}
\mot{tap / faucet}{tap}{robinet (BrE / AmE)}
\mot{washing machine}{OUO-ching me-chiin}{machine à laver}
\mot{iron}{AÏ-eun}{fer à repasser --- aussi : fer (métal)}
\mot{broom}{brouum}{balai}
\mot{bucket}{BEU-k't}{seau}
\mot{towel}{TA-ouel}{serviette}
\mot{soap}{sôoup}{savon}
\mot{blanket}{BLANG-k't}{couverture}
\mot{pillow}{PI-lou}{oreiller}
\mot{sheet}{chiit}{drap ; feuille}
\mot{electricity}{i-lèk-TRI-si-ti}{électricité}
\mot{water}{OUO-teu}{eau}
\mot{power cut}{PA-ou-eu keut}{coupure de courant}
\end{tabvocab}

\subsection{Thème 5 --- La nourriture et les boissons}
\begin{tabvocab}
\mot{food}{fouud}{nourriture}
\mot{meal}{miil}{repas --- \en{three meals a day}}
\mot{breakfast}{BRÈK-feust}{petit-déjeuner --- \en{to have breakfast}}
\mot{lunch}{leuntch}{déjeuner}
\mot{dinner}{DI-neu}{dîner}
\mot{snack}{snak}{en-cas, casse-croûte}
\mot{bread}{brèd}{pain --- \en{a loaf of bread}}
\mot{rice}{raïss}{riz}
\mot{meat}{miit}{viande}
\mot{beef}{biif}{b\oe uf}
\mot{chicken}{TCHI-k'n}{poulet}
\mot{fish}{fich}{poisson --- pluriel \en{fish}}
\mot{egg}{èg}{\oe uf}
\mot{vegetable}{VEDJ-teu-bl}{légume}
\mot{onion}{EUN-ieun}{oignon}
\mot{tomato}{teu-MÉ-i-tou}{tomate --- pluriel \en{tomatoes}}
\mot{potato}{peu-TÉ-i-tou}{pomme de terre}
\mot{carrot}{KA-rot}{carotte}
\mot{cabbage}{KA-bidj}{chou}
\mot{salad}{SA-lad}{salade}
\mot{fruit}{frouut}{fruit}
\mot{banana}{be-NA-neu}{banane}
\mot{orange}{O-rindj}{orange}
\mot{mango}{MANG-gou}{mangue}
\mot{lemon}{LÉ-meun}{citron}
\mot{apple}{A-pl}{pomme}
\mot{sugar}{CHOU-geu}{sucre}
\mot{salt}{soolt}{sel}
\mot{pepper}{PÉ-peu}{poivre}
\mot{oil}{oïl}{huile}
\mot{butter}{BEU-teu}{beurre}
\mot{cheese}{tchiiz}{fromage}
\mot{milk}{milk}{lait}
\mot{yoghurt}{IO-gheuut}{yaourt}
\mot{water}{OUO-teu}{eau --- \en{still / sparkling}}
\mot{tea / coffee}{tii / KO-fi}{thé / café}
\mot{juice}{djouuss}{jus}
\mot{soda / soft drink}{SÔ-da}{boisson gazeuse}
\mot{sauce}{sooss}{sauce}
\mot{soup}{souup}{soupe}
\mot{dish}{dich}{plat --- \en{a traditional dish}}
\mot{recipe}{RÉ-si-pi}{recette}
\mot{ingredients}{in-GRI-di-eunts}{ingrédients}
\mot{to cook}{kouk}{cuisiner}
\mot{to boil}{boïl}{bouillir}
\mot{to fry}{fraï}{frire}
\mot{to grill}{gril}{griller}
\mot{to bake}{bé-ik}{cuire au four}
\mot{to roast}{rouust}{rôtir}
\mot{to cut}{keut}{couper}
\mot{to mix}{miks}{mélanger}
\mot{to peel}{piil}{éplucher}
\mot{to taste}{té-ist}{goûter ; avoir le goût de}
\mot{delicious}{di-LI-cheus}{délicieux}
\mot{tasty}{TÉ-i-sti}{savoureux}
\mot{spicy}{SPA-i-si}{épicé}
\mot{sweet}{souit}{sucré}
\mot{sour}{SA-ou-eu}{acide, aigre}
\mot{bitter}{BI-teu}{amer}
\mot{fresh}{frèch}{frais}
\mot{raw}{roo}{cru}
\mot{ripe}{raïp}{mûr}
\mot{hungry}{HEU-ngri}{avoir faim --- \en{I'm hungry}}
\mot{thirsty}{THEU-sti}{avoir soif}
\mot{full}{foul}{plein, repu --- \en{I'm full}}
\mot{leftovers}{LÈFT-ou-veuz}{restes (du repas)}
\end{tabvocab}

\subsection{Thème 6 --- Au restaurant et au marché}
\begin{tabvocab}
\mot{restaurant}{RÈSS-teu-ront}{restaurant}
\mot{waiter / waitress}{OUÉ-i-teu}{serveur / serveuse}
\mot{menu}{MÉ-niou}{carte, menu}
\mot{bill / check}{bil / tchèk}{addition (BrE / AmE)}
\mot{tip}{tip}{pourboire}
\mot{reservation / booking}{rè-zeu-VÉ-i-cheun}{réservation --- \en{to book a table}}
\mot{order}{OO-deu}{commande ; \en{to order} commander}
\mot{dish of the day}{dich euv dhe DÉ-i}{plat du jour}
\mot{starter / appetizer}{STAA-teu}{entrée}
\mot{main course}{mé-in kooss}{plat principal}
\mot{dessert}{di-ZEUT}{dessert}
\mot{drink}{drink}{boisson}
\mot{glass / bottle}{glaass / BO-tl}{verre / bouteille}
\mot{plate}{plé-it}{assiette}
\mot{bowl}{bôoul}{bol}
\mot{fork / knife / spoon}{fook / naïf / spouun}{fourchette / couteau / cuillère}
\mot{napkin}{NA-pkine}{serviette de table}
\mot{market}{MAA-k't}{marché}
\mot{shop / store}{chop / stoo}{magasin}
\mot{supermarket}{SOU-peu-maa-k't}{supermarché}
\mot{customer}{KEUSS-teu-meu}{client}
\mot{seller}{SÉ-leu}{vendeur}
\mot{price}{praïss}{prix}
\mot{cheap / expensive}{tchiip / ik-SPÈN-siv}{bon marché / cher}
\mot{discount}{DIS-kaount}{remise, réduction}
\mot{bargain}{BAA-g'n}{bonne affaire}
\mot{receipt}{ri-SIIT}{reçu, ticket --- le \en{p} est muet}
\mot{cash}{kach}{espèces --- \en{to pay in cash}}
\mot{change}{tché-indj}{monnaie ; changement}
\mot{shopping}{CHO-ping}{courses --- \en{to do the shopping}}
\mot{basket / trolley}{BAAS-k't}{panier / chariot}
\mot{queue}{kiou}{file d'attente --- \en{to queue up}}
\mot{size}{saïz}{taille, pointure}
\mot{weigh}{oué-i}{peser}
\mot{kilo / gram}{KII-lou}{kilo / gramme}
\end{tabvocab}

\begin{astuce}[Les 12 phrases du restaurant]
\en{A table for two, please.} \quad \en{Could I see the menu?} \quad
\en{What do you recommend?} \quad \en{I'd like the grilled fish.} \quad
\en{I'll have the same.} \quad \en{Could I have some water, please?} \quad
\en{It's delicious, thank you.} \quad \en{I'm vegetarian / I don't eat pork.} \quad
\en{Could we have the bill, please?} \quad \en{Is service included?} \quad
\en{Keep the change.} \quad \en{Everything was lovely, thanks.}
\end{astuce}

\subsection{Thème 7 --- Les vêtements et les couleurs}
\begin{tabvocab}
\mot{clothes}{klôoudhz}{vêtements (toujours pluriel)}
\mot{shirt}{cheut}{chemise}
\mot{T-shirt}{TII-cheut}{t-shirt}
\mot{trousers / pants}{TRA-ou-zeuz}{pantalon (BrE / AmE)}
\mot{jeans}{djiinz}{jean}
\mot{shorts}{choo-ts}{short}
\mot{skirt}{skeut}{jupe}
\mot{dress}{drèss}{robe}
\mot{suit}{sout}{costume --- aussi : \en{to suit} aller à (qqn)}
\mot{jacket}{DJA-k't}{veste}
\mot{coat}{kôout}{manteau}
\mot{sweater / jumper}{SOÉ-teu}{pull}
\mot{hoodie}{HOU-di}{sweat à capuche}
\mot{shoes}{chouuz}{chaussures}
\mot{trainers / sneakers}{TRÉ-i-neuz}{baskets}
\mot{sandals}{SAN-deulz}{sandales}
\mot{boots}{bouuts}{bottes}
\mot{slippers}{SLI-peuz}{pantoufles}
\mot{socks}{soks}{chaussettes}
\mot{hat / cap}{hatt / kap}{chapeau / casquette}
\mot{scarf}{skaaf}{écharpe --- pluriel \en{scarves}}
\mot{gloves}{gleuvz}{gants}
\mot{belt}{bèlt}{ceinture}
\mot{bag}{bag}{sac}
\mot{purse / wallet}{peuss / OUO-l't}{porte-monnaie / portefeuille}
\mot{glasses}{glaa-siz}{lunettes}
\mot{watch}{ouotch}{montre}
\mot{jewellery / jewelry}{DJO-ou-el-ri}{bijoux}
\mot{ring}{ring}{bague ; anneau}
\mot{umbrella}{em-BRÉ-la}{parapluie}
\mot{to wear}{oué-eu}{porter (vêtement) --- \en{She wears a dress}}
\mot{to put on}{pout ON}{mettre (action)}
\mot{to take off}{té-ik OF}{enlever}
\mot{to try on}{traï ON}{essayer}
\mot{to fit}{fit}{aller à la bonne taille}
\mot{size}{saïz}{taille}
\mot{red, blue, green, yellow}{rèd, blou, griin, IÉ-lou}{rouge, bleu, vert, jaune}
\mot{black, white, grey / gray}{blak, ouaït, gré-i}{noir, blanc, gris}
\mot{brown, orange, pink, purple}{braoun, O-rindj, pink, PEU-pl}{marron, orange, rose, violet}
\mot{light blue / dark blue}{laït blou}{bleu clair / bleu foncé}
\mot{colourful}{KEU-leu-ful}{coloré}
\end{tabvocab}

\subsection{Thème 8 --- La ville, les transports, se déplacer}
\begin{tabvocab}
\mot{city / town}{SI-ti / taoun}{grande ville / ville}
\mot{village}{VI-lidj}{village}
\mot{street}{striit}{rue --- \en{main street}}
\mot{road}{rôoud}{route}
\mot{avenue / boulevard}{A-veu-niou}{avenue}
\mot{square}{skoué-eu}{place}
\mot{corner}{KOO-neu}{coin --- \en{at the corner}}
\mot{crossroads}{KROSS-rôoudz}{carrefour}
\mot{traffic lights}{TRA-fik laïts}{feux de circulation}
\mot{roundabout}{RAUND-eu-baout}{rond-point}
\mot{bridge}{bridj}{pont}
\mot{market}{MAA-k't}{marché}
\mot{mosque / church}{mosk / tcheutch}{mosquée / église}
\mot{school}{skouul}{école}
\mot{hospital}{HOSS-pi-tl}{hôpital}
\mot{bank}{bangk}{banque}
\mot{post office}{pôoust O-fiss}{bureau de poste}
\mot{police station}{peu-LIISS STÉ-i-cheun}{commissariat}
\mot{petrol station / gas station}{PÉ-tr'l}{station-service}
\mot{bus stop}{beuss stop}{arrêt de bus}
\mot{bus / coach}{beuss}{bus / autocar}
\mot{car}{kaa}{voiture}
\mot{taxi}{TA-k-si}{taxi}
\mot{train}{tré-in}{train}
\mot{plane / airplane}{plé-in}{avion}
\mot{boat / ship}{bôout / chip}{bateau / navire}
\mot{bicycle / motorbike}{BA-i-sik'l}{vélo / moto}
\mot{ticket}{TI-k't}{billet, ticket}
\mot{fare}{fé-eu}{tarif, prix du trajet}
\mot{platform}{PLAT-foom}{quai, plateforme}
\mot{station / airport}{STÉ-i-cheun}{gare / aéroport}
\mot{traffic}{TRA-fik}{circulation}
\mot{traffic jam}{TRA-fik djam}{embouteillage}
\mot{driver}{DRAÏ-veu}{conducteur}
\mot{passenger}{PA-ss'n-djeu}{passager}
\mot{to drive}{draïv}{conduire}
\mot{to ride}{raïd}{rouler, monter (à vélo, à cheval)}
\mot{to walk}{ouook}{marcher --- \en{on foot}}
\mot{to take the bus}{té-ik dhe beuss}{prendre le bus}
\mot{to get on / off}{guèt ON / OF}{monter / descendre (bus, train)}
\mot{to get in / out of}{guèt IN / AOUT euv}{monter / sortir (voiture)}
\mot{to turn left / right}{teun lèft / raït}{tourner à gauche / à droite}
\mot{to go straight on}{gou stré-it ON}{aller tout droit}
\mot{to cross}{kross}{traverser}
\mot{opposite}{O-ppo-zit}{en face de}
\mot{near / far from}{ni-eu / faa}{près / loin de}
\mot{next to}{nèkst tou}{à côté de}
\mot{between}{bi-TOUINN}{entre}
\mot{around / about}{e-RAOUND}{autour de, environ}
\mot{direction}{di-RÈK-cheun}{direction}
\mot{map}{map}{carte, plan}
\mot{distance}{DI-ss'teuns}{distance --- \en{It's two kilometres away.}}
\end{tabvocab}

\subsection{Thème 9 --- Le temps, les routines, le calendrier}
\begin{tabvocab}
\mot{time}{taïm}{temps, heure, fois}
\mot{clock / watch}{klok / ouotch}{horloge / montre}
\mot{hour / minute / second}{A-ou-eu / MI-nit / SÉ-keund}{heure / minute / seconde}
\mot{day / week / month / year}{dé-i / ouii-k / meunth / i-eu}{jour / semaine / mois / année}
\mot{today / tomorrow / yesterday}{teu-DÉ-i / te-MO-rou / IÈS-teu-dé-i}{aujourd'hui / demain / hier}
\mot{morning / afternoon / evening}{MOO-ning / aaf-teu-NOOUN / II-v'-ning}{matin / après-midi / soir}
\mot{night / midnight / noon}{naït / MID-naït / nouun}{nuit / minuit / midi}
\mot{Monday \ldots Sunday}{MEUN-di \ldots SEUN-di}{lundi \ldots dimanche}
\mot{January \ldots December}{DJA-niou-é-ri \ldots di-SSÈM-beu}{janvier \ldots décembre}
\mot{season}{SII-z'n}{saison --- \en{dry season, rainy season}}
\mot{weekend}{ouii-k'ènd}{week-end}
\mot{holiday / vacation}{HO-li-dé-i}{vacances, jour férié}
\mot{daily routine}{DÉ-i-li rou-TIIN}{routine quotidienne}
\mot{to wake up / get up}{oué-ik EP / guèt EP}{se réveiller / se lever}
\mot{to wash / shower}{ouoch / CHA-ou-eu}{se laver / se doucher}
\mot{to get dressed}{guèt drèst}{s'habiller}
\mot{to pray}{pré-i}{prier}
\mot{to commute}{keu-MIOUT}{faire le trajet domicile-travail}
\mot{to work / study}{oueu-k / STOU-di}{travailler / étudier}
\mot{to rest / relax}{rèst / ri-LAKS}{se reposer / se détendre}
\mot{to go to bed}{gou tou bèd}{aller se coucher}
\mot{to sleep}{sliip}{dormir --- \en{slept, slept}}
\mot{early / late}{EU-li / lé-it}{tôt / tard ; en retard}
\mot{on time / in time}{on taïm}{à l'heure / à temps}
\mot{busy / free}{BI-zi / frii}{occupé / libre}
\mot{always / never}{AAL-oué-iz / NÉ-veu}{toujours / jamais}
\end{tabvocab}

\subsection{Thème 10 --- Le travail et les métiers}
\begin{tabvocab}
\mot{work / job}{oueu-k / djob}{travail (activité) / emploi (poste)}
\mot{company / firm}{KEUM-peu-ni}{entreprise}
\mot{office}{O-fiss}{bureau (lieu)}
\mot{boss / manager}{boss / MA-n'e-djeu}{patron / directeur}
\mot{employee / employer}{em-PLOÏ-ii / em-PLOÏ-eu}{employé / employeur}
\mot{staff / team}{staaf / tiim}{personnel / équipe}
\mot{colleague}{KO-liig}{collègue}
\mot{customer / client}{KEUSS-teu-meu}{client}
\mot{salary / wage}{SA-leu-ri / oué-idj}{salaire (mensuel) / salaire (horaire)}
\mot{income}{IN-keum}{revenu}
\mot{contract}{KON-trakt}{contrat}
\mot{interview}{IN-teu-viou}{entretien d'embauche}
\mot{CV / resume}{si-vi / RÈ-ziou-mé-i}{curriculum vitæ}
\mot{experience}{ik-SPÍ-e-ri-euns}{expérience}
\mot{skill}{skil}{compétence}
\mot{training}{TRÉ-i-ning}{formation}
\mot{meeting}{MII-ting}{réunion}
\mot{deadline}{DÈD-laïn}{date limite}
\mot{task / project}{taask / PRO-djèkt}{tâche / projet}
\mot{shift}{chift}{poste, horaire --- \en{night shift}}
\mot{overtime}{Ô-veu-taïm}{heures supplémentaires}
\mot{full-time / part-time}{foul-taïm}{à temps plein / partiel}
\mot{unemployed}{eun-im-PLOÏD}{au chômage}
\mot{to hire / employ}{haï-eu}{embaucher}
\mot{to fire / sack}{faï-eu}{licencier}
\mot{to resign}{ri-ZAÏN}{démissionner}
\mot{to retire}{ri-TAÏ-eu}{prendre sa retraite}
\mot{to earn}{eun}{gagner (de l'argent)}
\mot{to apply for}{e-PLAÏ fo}{postuler à}
\mot{teacher, doctor, nurse}{TII-tcheu}{enseignant, médecin, infirmier}
\mot{engineer, technician}{en-dji-NI-eu}{ingénieur, technicien}
\mot{driver, mechanic}{DRAÏ-veu, me-KA-nik}{chauffeur, mécanicien}
\mot{farmer, fisherman}{FAA-meu}{agriculteur, pêcheur}
\mot{trader, shopkeeper}{TRÉ-i-deu}{commerçant}
\mot{tailor, carpenter, mason}{TÉ-i-lou}{tailleur, menuisier, maçon}
\mot{civil servant}{SI-v'l SEU-veunt}{fonctionnaire}
\mot{lawyer, judge}{LOÏ-eu}{avocat, juge}
\mot{police officer}{peu-LIISS O-fi-seu}{policier}
\mot{soldier}{SÔUL-djeu}{soldat}
\mot{accountant}{e-KAUN-teunt}{comptable}
\mot{secretary}{SÉ-creu-teu-ri}{secrétaire}
\mot{chef / cook}{chèf / kouk}{chef cuisinier / cuisinier}
\mot{hairdresser}{HÉ-eu-drè-seu}{coiffeur}
\mot{journalist}{DJEU-neu-list}{journaliste}
\end{tabvocab}

\subsection{Thème 11 --- L'école, les études, l'apprentissage}
\begin{tabvocab}
\mot{school}{skouul}{école}
\mot{primary / secondary school}{PRAÏ-me-ri}{école primaire / secondaire}
\mot{university}{iou-ni-VEU-si-ti}{université}
\mot{class / classroom}{klaass}{classe, cours}
\mot{lesson / course}{LÈ-ss'n / kooss}{leçon / cours (AmE : cours ; BrE : \en{course} = formation)}
\mot{subject}{SEUB-djèkt}{matière}
\mot{teacher / lecturer}{TII-tcheu}{enseignant / professeur d'université}
\mot{pupil / student}{piou-pl / STiou-d'nt}{élève / étudiant}
\mot{exam / test}{ig-ZAM / tèst}{examen / contrôle}
\mot{result / grade / mark}{ri-ZEULT}{résultat / note}
\mot{to pass / fail an exam}{paass / fé-il}{réussir / échouer à un examen}
\mot{homework}{HÔUM-oueu-k}{devoirs (à la maison)}
\mot{exercise}{ÉK-seu-saïz}{exercice}
\mot{question / answer}{KWÈS-tchoun / AA-n'seu}{question / réponse}
\mot{dictionary}{DIK-cheu-neu-ri}{dictionnaire}
\mot{notebook / textbook}{NÔUT-bouk}{cahier / manuel}
\mot{pen, pencil, eraser}{pèn, PÈN-sl, i-RÉ-i-zeu}{stylo, crayon, gomme}
\mot{board}{bood}{tableau --- \en{whiteboard, blackboard}}
\mot{to learn}{leun}{apprendre (acquérir)}
\mot{to study}{STOU-di}{étudier (travailler)}
\mot{to teach}{tiitch}{enseigner --- \en{taught, taught}}
\mot{to explain}{ik-SPLÉ-IN}{expliquer}
\mot{to understand}{en-deu-STAND}{comprendre --- \en{understood}}
\mot{to remember / forget}{ri-MÈM-beu / feu-GÈT}{se souvenir / oublier}
\mot{to practise}{PRAK-tiss}{s'entraîner}
\mot{to improve}{im-PROUUV}{améliorer}
\mot{to graduate}{GRA-djou-é-it}{obtenir son diplôme}
\mot{degree / diploma}{di-GRIi}{diplôme universitaire / diplôme}
\mot{scholarship}{SKO-leu-chip}{bourse d'études}
\mot{knowledge}{NO-lidj}{connaissances, savoir}
\mot{skill}{skil}{compétence}
\mot{progress}{PRO-grès}{progrès --- \en{to make progress}}
\end{tabvocab}

\subsection{Thème 12 --- L'argent, la banque, les achats en ligne}
\begin{tabvocab}
\mot{money}{MEU-ni}{argent (indénombrable)}
\mot{cash}{kach}{espèces}
\mot{coin / banknote}{koïn / BANGK-nôut}{pièce / billet}
\mot{currency}{KEU-ren-si}{monnaie --- \en{the CFA franc, the dollar, the euro}}
\mot{exchange rate}{iks-TCHÉ-indj ré-it}{taux de change}
\mot{bank}{bangk}{banque}
\mot{bank account}{bang e-KAOUNT}{compte bancaire}
\mot{credit / debit card}{KRÉ-dit kaad}{carte de crédit / débit}
\mot{ATM / cash machine}{é-i-tii-èm}{distributeur automatique}
\mot{transfer}{TRANS-feu}{virement}
\mot{loan}{lôoun}{prêt}
\mot{debt}{dèt}{dette --- le \en{b} est muet}
\mot{insurance}{in-CHOU-reuns}{assurance}
\mot{tax / duty}{taks / DIOUTI}{impôt / taxe}
\mot{receipt}{ri-SIIT}{reçu}
\mot{invoice}{IN-voïs}{facture}
\mot{budget}{BEU-djit}{budget}
\mot{savings}{SÉ-i-vingz}{économies}
\mot{to pay}{pé-i}{payer --- \en{paid, paid}}
\mot{to spend}{spènd}{dépenser --- \en{spent, spent}}
\mot{to save}{sé-iv}{économiser ; sauver}
\mot{to cost}{kost}{coûter --- \en{cost, cost}}
\mot{to borrow / lend}{BO-rou / lènd}{emprunter / prêter}
\mot{to owe}{ôou}{devoir (de l'argent)}
\mot{to refund}{ri-FEUND}{rembourser}
\mot{to afford}{e-FOOD}{avoir les moyens de}
\mot{online shopping}{ON-laïn CHO-ping}{achats en ligne}
\mot{to order / deliver}{OO-deu / di-LI-veu}{commander / livrer}
\mot{delivery fee}{di-LI-ve-ri fii}{frais de livraison}
\end{tabvocab}

\subsection{Thème 13 --- Le voyage, l'aéroport, l'hôtel}
\begin{tabvocab}
\mot{trip / journey}{trip / DJEU-ni}{voyage (court) / trajet}
\mot{travel}{TRA-v'l}{voyager (verbe)}
\mot{flight}{flaït}{vol --- \en{flight number}}
\mot{airport}{É-eu-poort}{aéroport}
\mot{check-in}{TCHÈK-in}{enregistrement}
\mot{boarding pass}{BOO-ding paass}{carte d'embarquement}
\mot{luggage / baggage}{LEU-gidj}{bagages}
\mot{suitcase}{SOUT-ké-iss}{valise}
\mot{hand luggage}{hand LEU-gidj}{bagage à main}
\mot{passport / visa}{PAASS-poort / VII-za}{passeport / visa}
\mot{customs}{KEUSS-teumz}{douane}
\mot{security check}{si-KIO-ou-ri-ti}{contrôle de sécurité}
\mot{gate}{gé-it}{porte d'embarquement}
\mot{departure / arrival}{di-PAA-tcheu / e-RAÏ-veul}{départ / arrivée}
\mot{delay}{di-LÉ-i}{retard}
\mot{to book}{bouk}{réserver}
\mot{hotel}{hÔ-TÈL}{hôtel}
\mot{reception / front desk}{ri-SÈP-cheun}{réception}
\mot{single / double room}{SI-ngl / DOU-bl roum}{chambre simple / double}
\mot{key / key card}{kii}{clé / carte-clé}
\mot{floor}{floo}{étage}
\mot{air conditioning}{É-eu keun-di-CHEU-ning}{climatisation}
\mot{breakfast included}{BRÈK-feust in-KLOU-did}{petit-déjeuner compris}
\mot{tourist / guide}{TO-ou-rist / gaïd}{touriste / guide}
\mot{sightseeing}{SAÏT-sii-ing}{visites touristiques}
\mot{beach}{biitch}{plage}
\mot{embassy}{ÈM-beu-si}{ambassade}
\end{tabvocab}

\subsection{Thème 14 --- La nature, la météo, les animaux}
\begin{tabvocab}
\mot{weather}{OUÉ-dheu}{temps (météo) --- \en{What's the weather like?}}
\mot{sunny / cloudy}{SEU-ni / KLA-ou-di}{ensoleillé / nuageux}
\mot{hot / warm / cool / cold}{hot / ouoom / kouul / kôould}{chaud / tiède / frais / froid}
\mot{rain / rainy}{ré-in}{pluie --- \en{It's raining}}
\mot{storm / thunder}{stoom / THEU-n'deu}{tempête / tonnerre}
\mot{wind / windy}{ouind}{vent --- \en{It's windy}}
\mot{humidity}{hiou-MI-di-ti}{humidité}
\mot{heat / drought}{hiit / draout}{chaleur / sécheresse}
\mot{flood}{fleud}{inondation}
\mot{sky / sun / moon / star}{skaï / seun / mouun / staa}{ciel / soleil / lune / étoile}
\mot{sea / ocean / river / lake}{sii / Ô-cheun / RI-veu / lé-ik}{mer / océan / fleuve / lac}
\mot{mountain / hill}{MAUN-t'n / hil}{montagne / colline}
\mot{forest / tree / plant}{FO-rist / trii / plaant}{forêt / arbre / plante}
\mot{sand / stone / soil}{sand / stôoun / soïl}{sable / pierre / sol}
\mot{animal}{A-ni-ml}{animal --- pluriel \en{animals}}
\mot{dog / cat / horse / cow / sheep / goat}{dog / katt / hooss / kaou / chiip / gôut}{chien / chat / cheval / vache / mouton / chèvre}
\mot{chicken / duck / bird}{TCHI-k'n / deuk / beud}{poule / canard / oiseau}
\mot{lion / elephant / monkey / snake}{LAÏ-eun}{lion / éléphant / singe / serpent}
\mot{fish / mosquito / fly / spider}{fich / meuss-KII-tou / flaï / SPAÏ-deu}{poisson / moustique / mouche / araignée}
\mot{tail / wing / paw}{té-il / ouing / poo}{queue / aile / patte}
\end{tabvocab}

\subsection{Thème 15 --- Loisirs, sport, musique, cinéma}
\begin{tabvocab}
\mot{hobby}{HO-bi}{passe-temps --- pluriel \en{hobbies}}
\mot{free time}{frii taïm}{temps libre}
\mot{sport / match / game}{spoort / mattch / gé-im}{sport / match / jeu}
\mot{football / basketball}{FOOT-booll}{football / basketball}
\mot{team / player / coach}{tiim / PLÉ-i-eu / kôoutch}{équipe / joueur / entraîneur}
\mot{to win / lose / draw}{ouine / louz / droo}{gagner / perdre / faire match nul}
\mot{goal / score}{gôoul / skoo}{but / score}
\mot{championship / tournament}{TCHAM-pieun-chip}{championnat / tournoi}
\mot{gym}{djiim}{salle de sport}
\mot{to work out}{oueu-k AOUT}{faire de l'exercice}
\mot{running / swimming}{RA-n'ning / SOI-m'ning}{course / natation}
\mot{music}{MIou-zik}{musique}
\mot{song / singer / band}{song / SI-ngheu / band}{chanson / chanteur / groupe}
\mot{to sing / dance / play}{sing / daans / plé-i}{chanter / danser / jouer}
\mot{film / movie}{film / MOU-vi}{film (BrE / AmE)}
\mot{cinema / theatre}{SI-neu-ma}{cinéma / théâtre}
\mot{actor / actress}{AK-teu}{acteur / actrice}
\mot{series / episode}{SÍ-eu-ri-z}{série / épisode}
\mot{documentary}{do-kieu-MÈN-t'ri}{documentaire --- \en{watch a documentary}}
\mot{subtitles}{SEUB-taï-tlz}{sous-titres}
\mot{screen / sound / picture}{skriin / saound / PIK-tcheu}{écran / son / image}
\mot{book / novel / story}{bouk / NO-v'l / STOO-ri}{livre / roman / histoire}
\mot{to read / write}{riid / raït}{lire / écrire --- \en{read, read ; wrote, written}}
\mot{to paint / draw}{pé-int / droo}{peindre / dessiner}
\mot{to travel / visit}{TRA-v'l / VI-zit}{voyager / visiter}
\mot{to cook / bake}{kouk / bé-ik}{cuisiner / pâtisser}
\mot{to take photos}{té-ik FO-tôuz}{prendre des photos --- \en{take pictures}}
\end{tabvocab}

\subsection{Thème 16 --- Les émotions et le caractère}
\begin{tabvocab}
\mot{happy / glad}{HA-pi / glad}{heureux, content}
\mot{unhappy / sad}{en-HA-pi / sad}{malheureux, triste}
\mot{angry}{AEN-gri}{en colère --- \en{to be angry with somebody}}
\mot{annoyed}{e-NOÏD}{agacé}
\mot{worried}{OUA-rid}{inquiet}
\mot{afraid / scared}{e-FRÉ-ID / ské-eud}{effrayé --- \en{afraid of}}
\mot{surprised}{seu-PRAÏZD}{surpris}
\mot{excited}{ik-SAÏ-tid}{enthousiaste --- \en{I'm excited about it}}
\mot{bored}{bood}{qui s'ennuie}
\mot{tired / exhausted}{TAÏ-eud / ig-ZOOS-tid}{fatigué / épuisé}
\mot{proud}{praoud}{fier --- \en{proud of}}
\mot{ashamed}{e-CHÉ-IMD}{honteux}
\mot{jealous}{DJÉ-leuss}{jaloux}
\mot{disappointed}{di-ss'e-POÏN-tid}{déçu}
\mot{confident}{KON-fi-d'nt}{confiant}
\mot{nervous}{NEU-veuss}{nerveux, stressé}
\mot{calm}{kaam}{calme}
\mot{lonely}{LÔUN-li}{seul (sentiment)}
\mot{grateful / thankful}{GRÉ-i-tful}{reconnaissant}
\mot{kind / generous}{kaïnd / DJÉ-neu-reus}{gentil / généreux}
\mot{polite / rude}{peu-LAÏT / rood}{poli / impoli}
\mot{honest / dishonest}{O-nist}{honnête / malhonnête}
\mot{hard-working / lazy}{haad-OUOE-king / LÉ-i-zi}{travailleur / paresseux}
\mot{serious / funny}{SÍ-eu-ri-eus / FA-ni}{sérieux / drôle}
\mot{friendly / shy}{FRÈND-li / chaï}{sympathique / timide}
\mot{patient / impatient}{PÉ-i-cheunt}{patient / impatient}
\mot{clever / smart / stupid}{KLÉ-veu / smaart / STiou-pid}{intelligent / futé / stupide}
\mot{brave / cowardly}{bré-iv}{courageux / lâche}
\mot{talkative / quiet}{TOO-keu-tiv / KOUAÏ-et}{bavard / calme}
\mot{ambitious}{am-BI-cheuss}{ambitieux}
\mot{feelings}{FII-lingz}{sentiments}
\mot{mood}{mouud}{humeur --- \en{in a good mood}}
\end{tabvocab}

\subsection{Thème 17 --- Technologie, téléphone, Internet}
\begin{tabvocab}
\mot{phone / mobile / cell}{fôoun / MÔ-baïl}{téléphone / portable}
\mot{screen}{skriin}{écran}
\mot{battery}{BA-teu-ri}{batterie}
\mot{charger}{TCHAAD-jeu}{chargeur}
\mot{to charge / unplug}{tchaadj / eun-PLEUG}{charger / débrancher}
\mot{to call / ring}{kooll / ring}{appeler --- \en{to give someone a call}}
\mot{text message / SMS}{tèkst MÈ-ssidj}{message texte}
\mot{to send / receive}{sènd / ri-SII-uv}{envoyer / recevoir}
\mot{call}{kooll}{appel --- \en{missed call}}
\mot{network / signal}{NÈT-oueu-k / SIG-neul}{réseau / signal}
\mot{data / Wi-Fi}{DÉ-i-ta}{données / wi-fi}
\mot{computer / laptop}{keum-Piou-teu}{ordinateur / portable}
\mot{keyboard / mouse}{KII-bood / maouss}{clavier / souris}
\mot{file / folder}{faïl / FÔUL-deu}{fichier / dossier}
\mot{to save / delete / download}{sé-iv / di-LIIT / DAOUN-lôd}{enregistrer / supprimer / télécharger}
\mot{to log in / out}{log IN / AOUT}{se connecter / se déconnecter}
\mot{password}{PAASS-oueu-d}{mot de passe}
\mot{username}{Iou-zeu-NÉ-im}{nom d'utilisateur}
\mot{website}{OUÈB-saït}{site web}
\mot{search engine}{seutch ÈN-djin}{moteur de recherche}
\mot{app / software / hardware}{ap / SOFT-oué-eu}{application / logiciel / matériel}
\mot{email / inbox}{ii-mé-il / IN-boks}{courriel / boîte de réception}
\mot{attachment}{e-TATCH-m'nt}{pièce jointe}
\mot{social media}{SÔ-cheul MII-di-a}{réseaux sociaux}
\mot{account}{e-KAOUNT}{compte}
\mot{to post / share / like / follow}{pôst / ché-eu}{publier / partager / aimer / suivre}
\mot{to click / tap / scroll}{klik / tap / skroul}{cliquer / taper / faire défiler}
\mot{video call}{VI-di-ou kooll}{appel vidéo}
\mot{to take a screenshot}{ské-in-chot}{faire une capture d'écran}
\mot{to fix / repair}{fiks / ri-PÉ-eu}{réparer}
\end{tabvocab}

\subsection{Thème 18 --- Médias, actualité, société}
\begin{tabvocab}
\mot{news}{niouz}{actualités (singulier !) --- \en{The news is bad}}
\mot{newspaper}{NIouz-pé-i-peu}{journal}
\mot{TV / channel}{téli-vi / TCHA-n'l}{télévision / chaîne}
\mot{radio}{RÉ-i-di-ou}{radio}
\mot{press}{prèss}{presse}
\mot{journalist / reporter}{DJEU-neu-list}{journaliste / reporter}
\mot{headline}{HÈD-laïn}{gros titre}
\mot{article}{AA-ti-kl}{article}
\mot{interview}{IN-teu-viou}{entretien, interview}
\mot{advertisement / advert}{ed-VEU-tiss-m'nt}{publicité}
\mot{information}{in-feu-MÉ-i-cheun}{informations (indénombrable)}
\mot{rumour / fake news}{ROU-meu}{rumeur / fausse nouvelle}
\mot{government}{GOU-veun-m'nt}{gouvernement}
\mot{president / minister}{PRÉ-zi-d'nt}{président / ministre}
\mot{election / vote}{i-LÈK-cheun / vôt}{élection / vote}
\mot{law / rule}{loo / roul}{loi / règle}
\mot{rights}{raïts}{droits --- \en{human rights}}
\mot{freedom}{FRII-deum}{liberté}
\mot{citizen}{SI-ti-z'n}{citoyen}
\mot{society}{seu-SAÏ-eu-ti}{société}
\mot{community}{keu-MIou-ni-ti}{communauté, quartier}
\mot{population}{po-piou-LÉ-i-cheun}{population}
\mot{development}{di-VÉ-leup-m'nt}{développement}
\mot{economy / market}{i-KO-no-mi}{économie / marché}
\mot{poverty / wealth}{PO-veu-ti / ouèlth}{pauvreté / richesse}
\mot{unemployment}{eun-im-PLOÏ-m'nt}{chômage}
\mot{education / health}{é-djou-KÉ-i-cheun / hèlth}{éducation / santé}
\mot{environment}{in-VAÏ-reun-m'nt}{environnement}
\mot{crime / theft}{kraïm / thèft}{crime / vol}
\mot{police / court / prison}{peu-LIISS / koort / PRÍ-z'n}{police / tribunal / prison}
\end{tabvocab}

\subsection{Thème 19 --- Religion et vie spirituelle}
\begin{tabvocab}
\mot{religion}{ri-LI-djeun}{religion}
\mot{God}{god}{Dieu}
\mot{faith}{fé-ith}{foi}
\mot{to believe in}{bi-LII-uv in}{croire en}
\mot{prayer}{pré-eu}{prière --- \en{to say a prayer}}
\mot{to pray}{pré-i}{prier}
\mot{worship}{OUOE-chip}{culte, adoration}
\mot{mosque / church / temple}{mosk / tcheutch / TÈM-pl}{mosquée / église / temple}
\mot{imam / priest / pastor}{i-MAAM / priist / PAASS-teu}{imam / prêtre / pasteur}
\mot{holy book}{HÔ-li bouk}{livre saint}
\mot{verse / chapter}{veuss / TCHA-p-teu}{verset / chapitre (sourate : \en{surah})}
\mot{recitation}{rè-si-TÉ-i-cheun}{récitation}
\mot{pilgrimage}{PI-lgri-midj}{pèlerinage}
\mot{fasting}{FAAS-ting}{jeûne --- \en{to fast}}
\mot{charity / almsgiving}{TCHA-ri-ti}{aumône, charité}
\mot{blessing}{BLÈ-ssing}{bénédiction}
\mot{to bless}{blèss}{bénir}
\mot{soul}{sôoul}{âme}
\mot{sin / forgiveness}{sine / feu-GIV-n's}{péché / pardon}
\mot{to repent / forgive}{ri-PÈNT / feu-GIV}{se repentir / pardonner}
\mot{mercy}{MEU-si}{miséricorde}
\mot{peace}{piiss}{paix --- \en{Peace be upon you}}
\mot{religious / secular}{ri-LI-djeus}{religieux / laïque}
\mot{celebration / festival}{sè-leu-BRÉ-i-cheun}{célébration / fête}
\end{tabvocab}

\begin{astuce}[Formules de respect interculturel]
\en{Peace be upon you.} (Que la paix soit sur vous.) \quad
\en{God willing.} / \en{If God wills it.} (Si Dieu le veut.) \quad
\en{May God bless you.} (Que Dieu vous bénisse.) \quad
\en{Congratulations on your wedding.} (Félicitations pour votre mariage.) \quad
\en{My condolences.} / \en{I'm sorry for your loss.} (Mes condoléances.) \quad
\en{Ramadan Mubarak.} / \en{Happy Tabaski!} / \en{Merry Christmas!}
\end{astuce}

\subsection{Thème 20 --- Décrire, expliquer, nuancer (mots-outils)}
\begin{tabvocab}
\mot{thing}{thing}{chose --- \en{the most important thing}}
\mot{way}{oué-i}{façon, manière --- \en{in this way, by the way}}
\mot{reason}{RII-z'n}{raison --- \en{the reason why}}
\mot{idea}{aï-DÍ-eu}{idée}
\mot{example}{ig-ZAM-pl}{exemple --- \en{for example}}
\mot{fact}{fakt}{fait --- \en{in fact}}
\mot{opinion}{e-PÍ-nieun}{opinion --- \en{in my opinion}}
\mot{point}{poïnt}{point, argument --- \en{good point}}
\mot{meaning}{MII-ning}{sens, signification}
\mot{difference}{DI-feu-reuns}{différence --- \en{make a difference}}
\mot{problem / solution}{PRO-b'l / se-LIOU-cheun}{problème / solution}
\mot{question / answer}{KWÈS-tchoun}{question / réponse}
\mot{result}{ri-ZEULT}{résultat}
\mot{chance}{tchaans}{chance, occasion --- \en{give me a chance}}
\mot{choice}{tchoïs}{choix}
\mot{decision}{di-SÍ-jeun}{décision}
\mot{advice}{ed-VAÏS}{conseil (indénombrable) --- \en{a piece of advice}}
\mot{news}{niouz}{nouvelle (singulier)}
\mot{experience}{ik-SPÍ-eu-ri-euns}{expérience}
\mot{knowledge}{NO-lidj}{connaissances}
\mot{purpose}{PEU-peuss}{but --- \en{on purpose} exprès}
\mot{advantage / disadvantage}{ed-VAAN-tidj}{avantage / inconvénient}
\mot{as soon as possible}{euss SOUN euss}{dès que possible --- \en{ASAP}}
\mot{at least / at most}{at liist}{au moins / au plus}
\mot{for instance}{fo IN-ss'teuns}{par exemple}
\mot{by the way}{baï dhe oué-i}{au fait, à propos}
\mot{in other words}{in EUDheu oueu-dz}{autrement dit}
\mot{that's why}{dhats ouaï}{c'est pourquoi}
\mot{it depends}{it di-PÈNDZ}{cela dépend}
\mot{I mean}{aï miin}{je veux dire}
\mot{actually}{AK-tchou-a-li}{en fait (faux ami : pas « actuellement »)}
\mot{probably / maybe / perhaps}{PRO-beu-bli}{probablement / peut-être / peut-être}
\end{tabvocab}

\section{Les 98 verbes essentiels (à savoir conjuguer par c\oe ur)}
\begin{tableau}[98 verbes essentiels avec leurs trois formes]
\footnotesize
\begin{longtable}{@{}p{2.4cm}p{2.2cm}p{2.4cm}p{3.0cm}p{2.4cm}p{2.2cm}@{}}
\entete{Base} & \entete{Passé} & \entete{Participe} & \entete{Sens} & \entete{Base} & \entete{Passé / p.p.} \tabularnewline
\endfirsthead
\entete{Base} & \entete{Passé} & \entete{Participe} & \entete{Sens} & \entete{Base} & \entete{Passé / p.p.} \tabularnewline
\endhead
\hline
\en{be} & \en{was/were} & \en{been} & être & \en{make} & \en{made / made} \tabularnewline
\en{have} & \en{had} & \en{had} & avoir & \en{take} & \en{took / taken} \tabularnewline
\en{do} & \en{did} & \en{done} & faire & \en{come} & \en{came / come} \tabularnewline
\en{say} & \en{said} & \en{said} & dire & \en{go} & \en{went / gone} \tabularnewline
\en{get} & \en{got} & \en{got(ten)} & obtenir, devenir & \en{see} & \en{saw / seen} \tabularnewline
\en{know} & \en{knew} & \en{known} & savoir & \en{want} & \en{wanted / wanted} \tabularnewline
\en{think} & \en{thought} & \en{thought} & penser & \en{need} & \en{needed / needed} \tabularnewline
\en{tell} & \en{told} & \en{told} & dire à qqn & \en{feel} & \en{felt / felt} \tabularnewline
\en{find} & \en{found} & \en{found} & trouver & \en{try} & \en{tried / tried} \tabularnewline
\en{give} & \en{gave} & \en{given} & donner & \en{leave} & \en{left / left} \tabularnewline
\en{work} & \en{worked} & \en{worked} & travailler & \en{become} & \en{became / become} \tabularnewline
\en{call} & \en{called} & \en{called} & appeler & \en{put} & \en{put / put} \tabularnewline
\en{ask} & \en{asked} & \en{asked} & demander & \en{mean} & \en{meant / meant} \tabularnewline
\en{use} & \en{used} & \en{used} & utiliser & \en{keep} & \en{kept / kept} \tabularnewline
\en{look} & \en{looked} & \en{looked} & regarder & \en{let} & \en{let / let} \tabularnewline
\en{seem} & \en{seemed} & \en{seemed} & sembler & \en{begin} & \en{began / begun} \tabularnewline
\en{help} & \en{helped} & \en{helped} & aider & \en{show} & \en{showed / shown} \tabularnewline
\en{talk} & \en{talked} & \en{talked} & parler & \en{hear} & \en{heard / heard} \tabularnewline
\en{turn} & \en{turned} & \en{turned} & tourner & \en{play} & \en{played / played} \tabularnewline
\en{start} & \en{started} & \en{started} & commencer & \en{run} & \en{ran / run} \tabularnewline
\en{move} & \en{moved} & \en{moved} & déménager, bouger & \en{answer} & \en{answered / answered} \tabularnewline
\en{like} & \en{liked} & \en{liked} & aimer & \en{believe} & \en{believed / believed} \tabularnewline
\en{live} & \en{lived} & \en{lived} & habiter & \en{bring} & \en{brought / brought} \tabularnewline
\en{love} & \en{loved} & \en{loved} & aimer (fort) & \en{happen} & \en{happened / happened} \tabularnewline
\en{eat} & \en{ate} & \en{eaten} & manger & \en{write} & \en{wrote / written} \tabularnewline
\en{drink} & \en{drank} & \en{drunk} & boire & \en{sit} & \en{sat / sat} \tabularnewline
\en{sleep} & \en{slept} & \en{slept} & dormir & \en{stand} & \en{stood / stood} \tabularnewline
\en{read} & \en{read} & \en{read} & lire & \en{lose} & \en{lost / lost} \tabularnewline
\en{buy} & \en{bought} & \en{bought} & acheter & \en{pay} & \en{paid / paid} \tabularnewline
\en{sell} & \en{sold} & \en{sold} & vendre & \en{meet} & \en{met / met} \tabularnewline
\en{open} & \en{opened} & \en{opened} & ouvrir & \en{include} & \en{included / included} \tabularnewline
\en{close} & \en{closed} & \en{closed} & fermer & \en{continue} & \en{continued / continued} \tabularnewline
\en{learn} & \en{learnt / learned} & \en{learnt} & apprendre & \en{set} & \en{set / set} \tabularnewline
\en{teach} & \en{taught} & \en{taught} & enseigner & \en{change} & \en{changed / changed} \tabularnewline
\en{understand} & \en{understood} & \en{understood} & comprendre & \en{lead} & \en{led / led} \tabularnewline
\en{remember} & \en{remembered} & \en{remembered} & se souvenir & \en{explain} & \en{explained / explained} \tabularnewline
\en{forget} & \en{forgot} & \en{forgotten} & oublier & \en{watch} & \en{watched / watched} \tabularnewline
\en{wait} & \en{waited} & \en{waited} & attendre & \en{follow} & \en{followed / followed} \tabularnewline
\en{stop} & \en{stopped} & \en{stopped} & arrêter & \en{finish} & \en{finished / finished} \tabularnewline
\en{walk} & \en{walked} & \en{walked} & marcher & \en{create} & \en{created / created} \tabularnewline
\en{drive} & \en{drove} & \en{driven} & conduire & \en{speak} & \en{spoke / spoken} \tabularnewline
\en{fly} & \en{flew} & \en{flown} & voler & \en{carry} & \en{carried / carried} \tabularnewline
\en{grow} & \en{grew} & \en{grown} & grandir, cultiver & \en{allow} & \en{allowed / allowed} \tabularnewline
\en{build} & \en{built} & \en{built} & construire & \en{add} & \en{added / added} \tabularnewline
\en{send} & \en{sent} & \en{sent} & envoyer & \en{spend} & \en{spent / spent} \tabularnewline
\en{break} & \en{broke} & \en{broken} & casser & \en{cut} & \en{cut / cut} \tabularnewline
\en{catch} & \en{caught} & \en{caught} & attraper & \en{raise} & \en{raised / raised} \tabularnewline
\en{choose} & \en{chose} & \en{chosen} & choisir & \en{kill} & \en{killed / killed} \tabularnewline
\en{cost} & \en{cost} & \en{cost} & coûter & \en{die} & \en{died / died} \tabularnewline
\end{longtable}
\end{tableau}

\begin{mission}[Chapitre 12 --- votre premier millier]
Objectif chiffré : \textbf{1\,000 mots en 100 jours}. Procédure :\\
\textbf{1.} Choisissez un thème par jour (20 thèmes ci-dessus, puis répétez).\\
\textbf{2.} Sélectionnez 15 mots du thème --- ceux dont vous avez vraiment besoin.\\
\textbf{3.} Créez 15 cartes Anki, chacune avec \textbf{une phrase personnelle} et la
prononciation.\\
\textbf{4.} Dites les 15 phrases à voix haute, deux fois.\\
\textbf{5.} Le soir, écrivez 5 phrases de journal en utilisant au moins 5 mots du jour.\\
En 100 jours, vous aurez 1\,000 mots actifs --- c'est-à-dire 85\,\% de la conversation
quotidienne.
\end{mission}

% =====================================================================
%  CHAPITRE 13 — PHRASAL VERBS, IDIOMES, FAUX AMIS
% =====================================================================
\chapter{\en{Phrasal verbs}, expressions idiomatiques et faux amis}

C'est le chapitre qui fait la différence entre un anglais « correct » et un anglais
« vivant ». Les anglophones utilisent des \en{phrasal verbs} en permanence : ils disent
\en{give up}, \en{find out}, \en{turn out}, \en{come up with}, là où un manuel écrit
\en{abandon}, \en{discover}, \en{result}, \en{invent}.

\section{Unité 1 --- Comprendre les \en{phrasal verbs}}

\begin{regle}[Structure]
\textbf{Verbe + particule} (adverbe ou préposition). Le sens est souvent
\textbf{imprévisible} : \en{look} (regarder) + \en{after} (après) = \en{look after}
(s'occuper de). Il faut donc les apprendre \textbf{comme des mots nouveaux}, avec un exemple.
\end{regle}

\begin{tableau}[Les quatre types --- savoir si l'on peut séparer]
\begin{tabularx}{\linewidth}{@{}p{3.2cm}p{4.0cm}X@{}}
\entete{Type} & \entete{Règle} & \entete{Exemples} \tabularnewline
\hline
\textbf{Séparable} & L'objet peut se placer entre le verbe et la particule. Avec un pronom (\en{it, them, him}), la séparation est \textbf{obligatoire}. & \en{turn off the light} = \en{turn the light off} = \en{turn \textbf{it} off} (\en{*turn off it} est faux) \tabularnewline
\textbf{Inséparable} & La particule reste collée au verbe. & \en{look after \textbf{her}} (jamais \en{*look her after}) ; \en{get over it} \tabularnewline
\textbf{Intransitif} & Pas d'objet du tout. & \en{He grew up in Dakar.} \en{The plane took off.} \tabularnewline
\textbf{À trois mots} & Verbe + particule + préposition, toujours inséparable. & \en{look forward to, put up with, get on with, run out of} \tabularnewline
\end{tabularx}
\end{tableau}

\begin{astuce}[La stratégie qui marche : apprendre par particule]
Au lieu d'apprendre des listes alphabétiques, apprenez le \textbf{sens général} de chaque
particule, puis regroupez :\\
\en{up} = achèvement, augmentation (\en{eat up, finish up, speed up, turn up, give up, use up})\\
\en{out} = sortir, apparaître, révéler (\en{go out, find out, work out, run out, take out, look out})\\
\en{off} = séparation, départ, arrêt (\en{take off, turn off, set off, cut off, show off, call off})\\
\en{on} = continuation (\en{go on, carry on, keep on, move on, put on})\\
\en{back} = retour (\en{come back, go back, give back, look back, call back})\\
\en{down} = diminution, arrêt (\en{slow down, calm down, break down, turn down, write down})\\
\en{away} = éloignement, disparition (\en{go away, run away, throw away, put away, take away})\\
\en{over} = au-dessus, réexamen (\en{get over, think over, look over, take over})\\
\en{through} = traversée, achèvement (\en{get through, look through, go through, come through})
\end{astuce}

\subsection{Les 150 \en{phrasal verbs} les plus utiles}
\begin{tableau}[150 \en{phrasal verbs} --- sens et exemple]
\footnotesize
\begin{longtable}{@{}p{3.0cm}p{4.4cm}p{6.4cm}@{}}
\entete{\en{Phrasal verb}} & \entete{Sens français} & \entete{Exemple} \tabularnewline
\endfirsthead
\entete{\en{Phrasal verb}} & \entete{Sens} & \entete{Exemple} \tabularnewline
\endhead
\hline
\en{ask out} & inviter (qqn) à sortir & \en{He asked her out.} \tabularnewline
\en{back up} & sauvegarder ; soutenir & \en{Back up your files. I'll back you up.} \tabularnewline
\en{blow up} & exploser ; gonfler & \en{The car blew up.} \tabularnewline
\en{break down} & tomber en panne ; s'effondrer & \en{My car broke down. She broke down in tears.} \tabularnewline
\en{break into} & s'introduire par effraction & \en{Someone broke into our house.} \tabularnewline
\en{break up} & se séparer ; se terminer & \en{They broke up last month.} \tabularnewline
\en{bring up} & élever ; aborder un sujet & \en{She brought up three kids. Don't bring that up again.} \tabularnewline
\en{build up} & développer progressivement & \en{Build up your vocabulary.} \tabularnewline
\en{call back} & rappeler (téléphone) & \en{I'll call you back.} \tabularnewline
\en{call off} & annuler & \en{They called off the meeting.} \tabularnewline
\en{calm down} & se calmer & \en{Calm down, it's not serious.} \tabularnewline
\en{carry on} & continuer & \en{Carry on working.} \tabularnewline
\en{catch up} & rattraper (son retard) & \en{Let's catch up soon!} \tabularnewline
\en{check in / out} & s'enregistrer / régler sa note & \en{We checked in at the hotel.} \tabularnewline
\en{check out} & examiner ; regarder & \en{Check out this video!} \tabularnewline
\en{cheer up} & remonter le moral & \en{Cheer up, it'll be fine.} \tabularnewline
\en{clean up} & nettoyer, ranger & \en{Clean up your room.} \tabularnewline
\en{come across} & trouver par hasard & \en{I came across an old photo.} \tabularnewline
\en{come from} & venir de & \en{Where do you come from?} \tabularnewline
\en{come in} & entrer & \en{Come in and sit down.} \tabularnewline
\en{come on} & allez ! dépêche-toi & \en{Come on, we're late!} \tabularnewline
\en{come out} & sortir ; paraître & \en{The sun came out. The book comes out in May.} \tabularnewline
\en{come over} & passer chez quelqu'un & \en{Come over for dinner.} \tabularnewline
\en{come up with} & trouver (une idée) & \en{She came up with a great idea.} \tabularnewline
\en{count on} & compter sur & \en{You can count on me.} \tabularnewline
\en{cut down on} & réduire & \en{Cut down on sugar.} \tabularnewline
\en{cut off} & couper ; interrompre & \en{The electricity was cut off. Don't cut me off.} \tabularnewline
\en{deal with} & s'occuper de, traiter & \en{I'll deal with it.} \tabularnewline
\en{depend on} & dépendre de & \en{It depends on the weather.} \tabularnewline
\en{dress up} & s'habiller élégamment & \en{We need to dress up.} \tabularnewline
\en{drop off} & déposer (qqn/qqch) & \en{I'll drop you off at school.} \tabularnewline
\en{eat out} & manger au restaurant & \en{Let's eat out tonight.} \tabularnewline
\en{end up} & finir par & \en{We ended up staying home.} \tabularnewline
\en{fall asleep} & s'endormir & \en{I fell asleep during the film.} \tabularnewline
\en{fall down} & tomber & \en{He fell down the stairs.} \tabularnewline
\en{figure out} & comprendre, résoudre & \en{I can't figure it out.} \tabularnewline
\en{fill in / out} & remplir (un formulaire) & \en{Fill in this form.} \tabularnewline
\en{find out} & découvrir, apprendre & \en{I found out the truth.} \tabularnewline
\en{get along / on with} & bien s'entendre avec & \en{I get along with my colleagues.} \tabularnewline
\en{get away} & s'échapper, partir en vacances & \en{We're getting away for the weekend.} \tabularnewline
\en{get back} & revenir ; récupérer & \en{When did you get back?} \tabularnewline
\en{get in / out of} & entrer / sortir (véhicule) & \en{Get in the car!} \tabularnewline
\en{get off / on} & descendre / monter (bus, train) & \en{Get off at the next stop.} \tabularnewline
\en{get over} & surmonter, se remettre & \en{She got over the illness quickly.} \tabularnewline
\en{get rid of} & se débarrasser de & \en{I got rid of my old phone.} \tabularnewline
\en{get up} & se lever & \en{I get up at six.} \tabularnewline
\en{give away} & donner, offrir ; trahir & \en{Don't give away the ending!} \tabularnewline
\en{give back} & rendre & \en{Give me back my pen.} \tabularnewline
\en{give up} & abandonner ; arrêter & \en{Never give up. He gave up smoking.} \tabularnewline
\en{go ahead} & continuer, vas-y & \en{Go ahead, I'm listening.} \tabularnewline
\en{go back} & retourner & \en{I want to go back home.} \tabularnewline
\en{go on} & continuer ; se passer & \en{Go on! What's going on?} \tabularnewline
\en{go out} & sortir ; s'éteindre & \en{Let's go out. The lights went out.} \tabularnewline
\en{go over} & revoir, vérifier & \en{Let's go over the plan.} \tabularnewline
\en{grow up} & grandir & \en{I grew up in Thiès.} \tabularnewline
\en{hand in} & rendre (un devoir) & \en{Hand in your papers.} \tabularnewline
\en{hang out} & traîner, passer du temps & \en{We hang out at the café.} \tabularnewline
\en{hang up} & raccrocher & \en{Don't hang up!} \tabularnewline
\en{hold on} & attendre, tenir & \en{Hold on a second.} \tabularnewline
\en{hurry up} & se dépêcher & \en{Hurry up or we'll be late.} \tabularnewline
\en{keep on} & continuer & \en{Keep on trying.} \tabularnewline
\en{keep up with} & suivre le rythme & \en{I can't keep up with you.} \tabularnewline
\en{leave out} & omettre, exclure & \en{Don't leave out the details.} \tabularnewline
\en{let down} & décevoir & \en{I won't let you down.} \tabularnewline
\en{listen to} & écouter & \en{Listen to me.} \tabularnewline
\en{live on} & vivre de & \en{They live on rice and fish.} \tabularnewline
\en{look after} & s'occuper de & \en{She looks after her grandmother.} \tabularnewline
\en{look for} & chercher & \en{I'm looking for my keys.} \tabularnewline
\en{look forward to} & attendre avec impatience & \en{I'm looking forward to it.} \tabularnewline
\en{look into} & examiner, enquêter sur & \en{We'll look into the problem.} \tabularnewline
\en{look out} & attention ! & \en{Look out! There's a car!} \tabularnewline
\en{look up} & chercher (un mot) & \en{Look it up in the dictionary.} \tabularnewline
\en{make up} & inventer ; se réconcilier & \en{Don't make up stories. They made up after the fight.} \tabularnewline
\en{move in / out} & emménager / déménager & \en{They moved in last week.} \tabularnewline
\en{pass away} & décéder (poli) & \en{His father passed away in 2020.} \tabularnewline
\en{pay back} & rembourser & \en{I'll pay you back tomorrow.} \tabularnewline
\en{pick up} & ramasser ; aller chercher ; apprendre & \en{Pick it up. I'll pick you up at 8. I picked up some English.} \tabularnewline
\en{point out} & faire remarquer & \en{He pointed out my mistake.} \tabularnewline
\en{put away} & ranger & \en{Put your phone away.} \tabularnewline
\en{put off} & reporter, remettre & \en{Don't put it off until tomorrow.} \tabularnewline
\en{put on} & mettre (vêtement) ; allumer & \en{Put on your coat. Put on the light.} \tabularnewline
\en{put out} & éteindre (feu) & \en{They put out the fire.} \tabularnewline
\en{put up with} & supporter, tolérer & \en{I can't put up with the noise.} \tabularnewline
\en{run into} & rencontrer par hasard & \en{I ran into Ali at the market.} \tabularnewline
\en{run out of} & être à court de & \en{We've run out of sugar.} \tabularnewline
\en{set off} & partir (en voyage) & \en{We set off at dawn.} \tabularnewline
\en{set up} & installer, fonder & \en{They set up a business.} \tabularnewline
\en{show off} & se vanter, frimer & \en{Stop showing off!} \tabularnewline
\en{show up} & arriver, se pointer & \en{He showed up late.} \tabularnewline
\en{shut down} & fermer, éteindre & \en{Shut down your computer.} \tabularnewline
\en{shut up} & se taire (familier, impoli) & \en{Shut up!} --- à éviter avec des inconnus. \tabularnewline
\en{sit down} & s'asseoir & \en{Please sit down.} \tabularnewline
\en{slow down} & ralentir & \en{Slow down, please, I can't follow.} \tabularnewline
\en{speak up} & parler plus fort & \en{Could you speak up?} \tabularnewline
\en{stand up} & se lever & \en{Everybody stood up.} \tabularnewline
\en{switch on / off} & allumer / éteindre & \en{Switch off the TV.} \tabularnewline
\en{take after} & ressembler à (un parent) & \en{She takes after her mother.} \tabularnewline
\en{take away} & emporter & \en{Let's get a takeaway.} \tabularnewline
\en{take care of} & prendre soin de & \en{Take care of yourself.} \tabularnewline
\en{take off} & décoller ; enlever & \en{The plane took off. Take off your shoes.} \tabularnewline
\en{take over} & prendre le contrôle & \en{She took over the company.} \tabularnewline
\en{take up} & commencer (un loisir) ; prendre (de la place) & \en{I took up football. It takes up too much space.} \tabularnewline
\en{talk about} & parler de & \en{Let's talk about your plans.} \tabularnewline
\en{think about / of} & penser à & \en{I'm thinking about moving.} \tabularnewline
\en{think over} & réfléchir à & \en{I'll think it over.} \tabularnewline
\en{throw away} & jeter & \en{Don't throw it away.} \tabularnewline
\en{try on} & essayer (vêtement) & \en{Can I try it on?} \tabularnewline
\en{turn down} & baisser (le son) ; refuser & \en{Turn down the music. He turned down the offer.} \tabularnewline
\en{turn into} & se transformer en & \en{The dream turned into reality.} \tabularnewline
\en{turn off / on} & éteindre / allumer & \en{Turn off your phone.} \tabularnewline
\en{turn out} & s'avérer & \en{It turned out to be true.} \tabularnewline
\en{turn up} & arriver (inattendu) ; monter le son & \en{He turned up at midnight. Turn up the volume.} \tabularnewline
\en{use up} & utiliser jusqu'au bout & \en{We've used up all the money.} \tabularnewline
\en{wake up} & se réveiller & \en{I wake up at five.} \tabularnewline
\en{warm up} & se réchauffer ; s'échauffer & \en{Let's warm up before the match.} \tabularnewline
\en{watch out} & faire attention & \en{Watch out for the dog!} \tabularnewline
\en{wear out} & user ; épuiser & \en{These shoes are worn out. I'm worn out.} \tabularnewline
\en{work out} & fonctionner ; s'entraîner & \en{It worked out fine. I work out twice a week.} \tabularnewline
\en{write down} & noter & \en{Write down the address.} \tabularnewline
\end{longtable}
\end{tableau}

\begin{exo}[1.1 --- Choisissez le \en{phrasal verb}]
\begin{enumerate}
  \item My car (broke down / broke up) on the road.
  \item Please (look after / look for) my bag while I'm away.
  \item I can't (put up with / put on) this noise.
  \item She (takes after / takes off) her father --- they look the same.
  \item We (ran out of / ran into) milk, I need to buy some.
  \item He (made up / made out) an excuse and left.
  \item Can you (speak up / speak out)? I can't hear you.
  \item Don't (give up / give in) --- you're almost there.
  \item I (came across / came up) this photo while cleaning.
  \item They (called off / called back) the match because of the rain.
\end{enumerate}
\end{exo}
\begin{corrige}[1.1]
1. broke down 2. look after 3. put up with 4. takes after 5. ran out of
6. made up 7. speak up 8. give up 9. came across 10. called off
\end{corrige}

\section{Unité 2 --- Les expressions idiomatiques}

Un idiome ne se traduit pas mot à mot : \en{it's raining cats and dogs} ne parle ni de
chats ni de chiens. Apprenez-en \textbf{100} --- cela suffit pour comprendre l'humour,
les films et les conversations normales.

\subsection{Les 40 idiomes les plus fréquents (à connaître absolument)}
\begin{tableau}[40 expressions indispensables]
\footnotesize
\begin{longtable}{@{}p{5.2cm}p{4.8cm}p{3.8cm}@{}}
\entete{Expression} & \entete{Sens français} & \entete{Exemple} \tabularnewline
\endfirsthead
\entete{Expression} & \entete{Sens} & \entete{Exemple} \tabularnewline
\endhead
\hline
\en{a piece of cake} & très facile & \en{The exam was a piece of cake.} \tabularnewline
\en{break a leg!} & bonne chance ! & \en{Break a leg on stage!} \tabularnewline
\en{call it a day} & arrêter pour aujourd'hui & \en{Let's call it a day.} \tabularnewline
\en{cost an arm and a leg} & coûter très cher & \en{That car cost an arm and a leg.} \tabularnewline
\en{don't judge a book by its cover} & ne pas se fier aux apparences & \en{Don't judge a book by its cover.} \tabularnewline
\en{every cloud has a silver lining} & à quelque chose malheur est bon & \en{You lost your job, but every cloud has a silver lining.} \tabularnewline
\en{get cold feet} & se dégonfler & \en{He got cold feet before the wedding.} \tabularnewline
\en{hit the nail on the head} & mettre le doigt dessus & \en{You hit the nail on the head.} \tabularnewline
\en{in the same boat} & dans le même bateau & \en{We're all in the same boat.} \tabularnewline
\en{it's not my cup of tea} & ce n'est pas mon truc & \en{Football is not my cup of tea.} \tabularnewline
\en{keep an eye on} & garder un \oe il sur & \en{Keep an eye on my bag.} \tabularnewline
\en{kill two birds with one stone} & faire d'une pierre deux coups & \en{I killed two birds with one stone.} \tabularnewline
\en{let the cat out of the bag} & vendre la mèche & \en{She let the cat out of the bag.} \tabularnewline
\en{make ends meet} & joindre les deux bouts & \en{It's hard to make ends meet.} \tabularnewline
\en{miss the boat} & rater l'occasion & \en{Don't miss the boat --- apply now.} \tabularnewline
\en{no pain, no gain} & on n'a rien sans rien & \en{No pain, no gain!} \tabularnewline
\en{once in a blue moon} & très rarement & \en{I see him once in a blue moon.} \tabularnewline
\en{on the one hand\ldots on the other hand} & d'un côté\ldots de l'autre & \en{On the one hand, it's cheap\ldots} \tabularnewline
\en{out of the blue} & à l'improviste & \en{He called me out of the blue.} \tabularnewline
\en{pull someone's leg} & faire marcher qqn & \en{Are you pulling my leg?} \tabularnewline
\en{put all your eggs in one basket} & mettre tous ses \oe ufs dans le même panier & \en{Don't put all your eggs in one basket.} \tabularnewline
\en{rain cats and dogs} & pleuvoir des cordes & \en{It's raining cats and dogs.} \tabularnewline
\en{see eye to eye} & être d'accord & \en{We don't see eye to eye on politics.} \tabularnewline
\en{so far, so good} & jusqu'ici, tout va bien & \en{So far, so good.} \tabularnewline
\en{take it easy} & détends-toi, vas-y doucement & \en{Take it easy, you'll be fine.} \tabularnewline
\en{the best of both worlds} & le meilleur des deux mondes & \en{Working from home gives you the best of both worlds.} \tabularnewline
\en{the early bird catches the worm} & le monde appartient à ceux qui se lèvent tôt & \en{Come at 7 --- the early bird catches the worm.} \tabularnewline
\en{there's no such thing as a free lunch} & rien n'est gratuit & \en{There's no such thing as a free lunch.} \tabularnewline
\en{time flies} & le temps passe vite & \en{Time flies when you're busy.} \tabularnewline
\en{to make a long story short} & bref, pour faire court & \en{Long story short, I got the job.} \tabularnewline
\en{under the weather} & pas dans son assiette & \en{I feel a bit under the weather.} \tabularnewline
\en{up in the air} & incertain, en suspens & \en{The plan is still up in the air.} \tabularnewline
\en{when pigs fly} & quand les poules auront des dents & \en{He'll apologise when pigs fly.} \tabularnewline
\en{you can say that again!} & tu l'as dit ! & \en{``It's hot.'' --- ``You can say that again!''} \tabularnewline
\en{your guess is as good as mine} & je n'en sais pas plus que toi & \en{Where is he? Your guess is as good as mine.} \tabularnewline
\en{bite the bullet} & se résigner, s'y mettre & \en{I bit the bullet and called her.} \tabularnewline
\en{break the ice} & briser la glace & \en{A joke helps break the ice.} \tabularnewline
\en{get the hang of it} & prendre le coup de main & \en{You'll get the hang of it.} \tabularnewline
\en{it takes two to tango} & il faut être deux & \en{It takes two to tango.} \tabularnewline
\en{the ball is in your court} & c'est à toi de jouer & \en{I've done my part --- the ball is in your court.} \tabularnewline
\end{longtable}
\end{tableau}

\subsection{80 expressions supplémentaires, par thème}
\begin{tableau}[Expressions de la vie quotidienne]
\footnotesize
\begin{longtable}{@{}p{5.4cm}p{8.4cm}@{}}
\entete{Expression} & \entete{Sens} \tabularnewline
\endfirsthead
\entete{Expression} & \entete{Sens} \tabularnewline
\endhead
\hline
\en{by the way} & au fait, à propos \tabularnewline
\en{that's up to you} & c'est toi qui décides \tabularnewline
\en{it doesn't matter} & ça n'a pas d'importance \tabularnewline
\en{never mind} & laisse tomber, peu importe \tabularnewline
\en{don't worry about it} & ne t'inquiète pas pour ça \tabularnewline
\en{take your time} & prends ton temps \tabularnewline
\en{hurry up, we're late} & dépêche-toi, on est en retard \tabularnewline
\en{I'm on my way} & j'arrive, je suis en route \tabularnewline
\en{I'm running late} & je vais être en retard \tabularnewline
\en{long time no see!} & ça fait longtemps ! \tabularnewline
\en{what's up?} & quoi de neuf ? \tabularnewline
\en{not much / nothing much} & pas grand-chose \tabularnewline
\en{how's it going?} & comment ça va ? \tabularnewline
\en{I can't complain} & je ne peux pas me plaindre \tabularnewline
\en{same old, same old} & toujours la même chose \tabularnewline
\en{I'm stuffed} & je n'en peux plus (de manger) \tabularnewline
\en{I'm starving} & je meurs de faim \tabularnewline
\en{I'm exhausted / knackered} & je suis épuisé \tabularnewline
\en{I'm fed up with it} & j'en ai marre \tabularnewline
\en{it's a pain} & c'est une corvée \tabularnewline
\en{no big deal} & pas grave \tabularnewline
\en{that's a shame / too bad} & c'est dommage \tabularnewline
\en{what a pity!} & quel dommage ! \tabularnewline
\en{congratulations! well done!} & félicitations ! bravo ! \tabularnewline
\en{good luck / all the best} & bonne chance / tous mes v\oe ux \tabularnewline
\en{take care} & prends soin de toi \tabularnewline
\en{see you later / soon} & à plus tard / à bientôt \tabularnewline
\en{have a good one} & bonne journée (AmE) \tabularnewline
\en{my pleasure / you're welcome} & avec plaisir / de rien \tabularnewline
\en{no problem / don't mention it} & pas de souci / il n'y a pas de quoi \tabularnewline
\en{I really appreciate it} & j'apprécie vraiment \tabularnewline
\en{it's my treat} & c'est moi qui invite \tabularnewline
\en{can I get you something?} & je vous sers quelque chose ? \tabularnewline
\en{help yourself} & servez-vous \tabularnewline
\en{make yourself at home} & faites comme chez vous \tabularnewline
\en{I'd better go} & je ferais mieux d'y aller \tabularnewline
\en{I'm afraid I can't} & je crains de ne pas pouvoir \tabularnewline
\en{that's fine by me} & ça me va \tabularnewline
\en{sounds good / sounds great} & ça me paraît bien \tabularnewline
\en{I'm not sure about that} & je n'en suis pas sûr \tabularnewline
\en{I have no idea} & je n'en ai aucune idée \tabularnewline
\en{as far as I know} & pour autant que je sache \tabularnewline
\en{it seems to me that} & il me semble que \tabularnewline
\en{to be honest} & pour être honnête \tabularnewline
\en{frankly speaking} & franchement \tabularnewline
\en{in my opinion / I reckon} & à mon avis / je crois \tabularnewline
\en{I couldn't agree more} & je suis tout à fait d'accord \tabularnewline
\en{I see your point, but} & je comprends votre point de vue, mais \tabularnewline
\en{that's not what I meant} & ce n'est pas ce que je voulais dire \tabularnewline
\en{what do you mean by that?} & qu'entendez-vous par là ? \tabularnewline
\en{you've got a point there} & tu marques un point \tabularnewline
\en{let's agree to disagree} & restons chacun sur nos positions \tabularnewline
\en{it depends on the situation} & cela dépend de la situation \tabularnewline
\en{at the end of the day} & au bout du compte \tabularnewline
\en{the thing is\ldots} & le problème, c'est que\ldots \tabularnewline
\en{to cut a long story short} & pour faire court \tabularnewline
\en{in a nutshell} & en résumé \tabularnewline
\en{that reminds me} & cela me fait penser \tabularnewline
\en{speaking of which} & en parlant de ça \tabularnewline
\en{by any chance} & par hasard \tabularnewline
\en{out of order} & en panne, hors service \tabularnewline
\en{sold out} & épuisé, tout vendu \tabularnewline
\en{on sale} & en solde \tabularnewline
\en{for free / free of charge} & gratuitement \tabularnewline
\en{in advance} & à l'avance \tabularnewline
\en{on time / right on time} & à l'heure / pile à l'heure \tabularnewline
\en{from time to time} & de temps en temps \tabularnewline
\en{all of a sudden} & tout à coup \tabularnewline
\en{as usual} & comme d'habitude \tabularnewline
\en{more or less} & plus ou moins \tabularnewline
\en{pretty much} & à peu près \tabularnewline
\en{kind of / sort of} & un peu, genre \tabularnewline
\en{a lot of / loads of} & beaucoup de \tabularnewline
\en{a bit of a problem} & un petit problème \tabularnewline
\en{the worst part is} & le pire, c'est \tabularnewline
\en{to be on the safe side} & par précaution \tabularnewline
\en{it's worth a try} & ça vaut le coup d'essayer \tabularnewline
\en{give it a go / a shot} & tente le coup \tabularnewline
\en{make the most of it} & profites-en au maximum \tabularnewline
\en{it's now or never} & c'est maintenant ou jamais \tabularnewline
\en{better late than never} & mieux vaut tard que jamais \tabularnewline
\end{longtable}
\end{tableau}

\subsection{L'argot courant (à comprendre, à employer avec modération)}
\begin{tableau}[30 mots familiers très entendus]
\footnotesize
\begin{longtable}{@{}p{3.4cm}p{4.0cm}p{6.4cm}@{}}
\entete{Mot familier} & \entete{Équivalent standard} & \entete{Exemple} \tabularnewline
\endfirsthead
\entete{Familier} & \entete{Standard} & \entete{Exemple} \tabularnewline
\endhead
\hline
\en{guy / dude / mate / bro} & \en{man, friend} & \en{Hey dude, what's up?} \tabularnewline
\en{kids} & \en{children} & \en{She has three kids.} \tabularnewline
\en{stuff / things} & \en{belongings} & \en{Put your stuff here.} \tabularnewline
\en{a lot / tons / loads} & \en{many} & \en{There were loads of people.} \tabularnewline
\en{awesome / amazing / cool} & \en{very good} & \en{That's awesome!} \tabularnewline
\en{awful / terrible / rubbish} & \en{very bad} & \en{The film was rubbish.} \tabularnewline
\en{cheap / pricey} & \en{inexpensive / expensive} & \en{It's a bit pricey.} \tabularnewline
\en{broke} & \en{without money} & \en{I'm broke until payday.} \tabularnewline
\en{bucks / quid} & \en{dollars / pounds} & \en{It costs ten bucks.} \tabularnewline
\en{gig} & \en{job, concert} & \en{I got a new gig.} \tabularnewline
\en{chill / hang out} & \en{relax} & \en{Let's just chill at home.} \tabularnewline
\en{no way!} & \en{impossible!} & \en{No way! You're joking.} \tabularnewline
\en{for real?} & \en{really?} & \en{For real? That's crazy.} \tabularnewline
\en{I'm done} & \en{I have finished} & \en{I'm done with this.} \tabularnewline
\en{my bad} & \en{my mistake, sorry} & \en{My bad, I forgot.} \tabularnewline
\en{no worries} & \en{it's fine} & \en{``Sorry!'' --- ``No worries.''} \tabularnewline
\en{I'm knackered / beat} & \en{I'm exhausted} & \en{I'm beat after work.} \tabularnewline
\en{to be broke / skint} & \en{to have no money} & \en{I'm skint this month.} \tabularnewline
\en{a rip-off} & \en{overpriced} & \en{That price is a rip-off.} \tabularnewline
\en{to be into something} & \en{to like} & \en{I'm really into football.} \tabularnewline
\en{to ghost someone} & \en{to stop replying} & \en{He ghosted me after two dates.} \tabularnewline
\en{to flake} & \en{to cancel at the last minute} & \en{Don't flake on me.} \tabularnewline
\en{to screw up} & \en{to make a mistake} & \en{I screwed up the interview.} \tabularnewline
\en{to be fed up} & \en{to be annoyed} & \en{I'm fed up with waiting.} \tabularnewline
\en{to hang in there} & \en{to persevere} & \en{Hang in there, it'll get better.} \tabularnewline
\en{to play it by ear} & \en{to improvise} & \en{Let's play it by ear.} \tabularnewline
\en{to be on the same page} & \en{to agree} & \en{Are we on the same page?} \tabularnewline
\en{to hit the sack} & \en{to go to bed} & \en{I'm hitting the sack.} \tabularnewline
\en{what's the catch?} & \en{what's the hidden problem?} & \en{Free? What's the catch?} \tabularnewline
\en{it's a no-brainer} & \en{it's an obvious choice} & \en{Take the job --- it's a no-brainer.} \tabularnewline
\end{longtable}
\end{tableau}

\section{Unité 3 --- Les faux amis : 150 mots qui vous trahissent}

Un faux ami est un mot qui ressemble au français mais ne veut pas dire la même chose.
Ce sont les fautes les plus dangereuses : personne ne vous corrige, on vous comprend mal
--- et vous ne le savez pas.

\begin{tableau}[150 faux amis français \ra anglais]
\footnotesize
\begin{longtable}{@{}p{3.0cm}p{3.6cm}p{3.6cm}p{3.6cm}@{}}
\entete{Mot français} & \entete{Ce que dit l'anglais} & \entete{Le mot anglais qui ressemble} & \entete{Ce qu'il veut dire vraiment} \tabularnewline
\endfirsthead
\entete{Français} & \entete{Anglais correct} & \entete{Faux ami} & \entete{Sens réel} \tabularnewline
\endhead
\hline
actuellement & \en{currently, at present} & \en{actually} & en fait, en réalité \tabularnewline
éventuellement & \en{possibly} & \en{eventually} & finalement, à la fin \tabularnewline
attendre & \en{to wait (for)} & \en{to attend} & assister à, être présent \tabularnewline
assister à & \en{to attend} & \en{to assist} & aider \tabularnewline
blesser & \en{to injure, hurt} & \en{to bless} & bénir \tabularnewline
librairie & \en{bookshop} & \en{library} & bibliothèque \tabularnewline
bibliothèque & \en{library} & \en{librairie} (rare) & --- \tabularnewline
blessé & \en{injured} & \en{blessed} & béni \tabularnewline
caméra & \en{video camera} & \en{camera} & appareil photo \tabularnewline
appareil photo & \en{camera} & --- & --- \tabularnewline
monnaie (pièces) & \en{change, coins} & \en{money} & argent en général \tabularnewline
monnaie (devise) & \en{currency} & \en{money} & argent \tabularnewline
déjeuner & \en{lunch} & \en{dinner} & dîner \tabularnewline
dîner & \en{dinner} & \en{diner} & personne qui dîne ; petit restaurant \tabularnewline
pain & \en{bread} & \en{pain} & douleur \tabularnewline
travail & \en{work, job} & \en{travail} (rare) & --- \tabularnewline
travailler & \en{to work} & \en{to travel} & voyager \tabularnewline
voyage & \en{trip, journey} & \en{voyage} & traversée maritime \tabularnewline
magasin & \en{shop, store} & \en{magazine} & revue \tabularnewline
revue & \en{magazine} & --- & --- \tabularnewline
blessure & \en{wound, injury} & --- & --- \tabularnewline
envie & \en{desire, wish} & \en{envy} & jalousie \tabularnewline
jalousie & \en{jealousy} & --- & --- \tabularnewline
sensible & \en{sensitive} & \en{sensible} & raisonnable, sensé \tabularnewline
raisonnable & \en{reasonable} & \en{sensible} & sensé \tabularnewline
coin & \en{corner} & \en{coin} & pièce de monnaie \tabularnewline
pièce & \en{room, coin} & --- & --- \tabularnewline
chance & \en{luck} & \en{chance} & occasion, hasard \tabularnewline
hasard & \en{chance, randomness} & \en{hazard} & danger \tabularnewline
danger & \en{danger} & \en{hazard} & risque \tabularnewline
médecine & \en{medicine} & \en{medication} & traitement médicamenteux \tabularnewline
médicament & \en{medicine, drug, pill} & \en{medicine} & aussi : la médecine \tabularnewline
drogue & \en{drug} & \en{drug} & médicament ET drogue \tabularnewline
pharmacie & \en{pharmacy, chemist's} & \en{pharmacy} & idem \tabularnewline
préservatif & \en{condom} & \en{preservative} & conservateur (alimentaire) \tabularnewline
conservateur & \en{preservative, conservative} & \en{preservative} & additif alimentaire \tabularnewline
expérience & \en{experience} & \en{experiment} & expérience scientifique \tabularnewline
expérience (labo) & \en{experiment} & --- & --- \tabularnewline
essayer & \en{to try} & \en{to essay} (rare) & --- \tabularnewline
essai & \en{attempt, essay} & \en{essay} & dissertation \tabularnewline
sujet & \en{subject, topic} & \en{subject} & aussi : sujet grammatical \tabularnewline
objet & \en{object, item} & \en{objective} & objectif \tabularnewline
objectif & \en{goal, objective} & \en{objective} & objectif (impartial) \tabularnewline
réellement & \en{really} & \en{actually} & en fait \tabularnewline
actuel & \en{current, present} & \en{actual} & réel \tabularnewline
issue & \en{exit} & \en{issue} & problème, question \tabularnewline
sortie & \en{exit, way out} & --- & --- \tabularnewline
blessé & \en{injured} & --- & --- \tabularnewline
rester & \en{to stay} & \en{to rest} & se reposer \tabularnewline
se reposer & \en{to rest} & --- & --- \tabularnewline
demander & \en{to ask} & \en{to demand} & exiger \tabularnewline
exiger & \en{to demand, require} & --- & --- \tabularnewline
répondre & \en{to answer, reply} & \en{to respond} & réagir \tabularnewline
proposer & \en{to suggest, offer} & \en{to propose} & demander en mariage ; proposer un plan \tabularnewline
soutenir & \en{to support} & \en{to sustain} & maintenir \tabularnewline
supporter (tolérer) & \en{to stand, put up with} & \en{to support} & soutenir \tabularnewline
introduire & \en{to introduce (people)} & \en{to introduce} & présenter quelqu'un \tabularnewline
présenter & \en{to introduce, present} & --- & --- \tabularnewline
connaître & \en{to know} & \en{to recognise} & reconnaître \tabularnewline
reconnaître & \en{to recognise} & --- & --- \tabularnewline
manquer & \en{to miss} & \en{to lack} & manquer de, être dépourvu \tabularnewline
apprendre & \en{to learn} & \en{to teach} & enseigner \tabularnewline
enseigner & \en{to teach} & --- & --- \tabularnewline
emprunter & \en{to borrow} & \en{to lend} & prêter \tabularnewline
prêter & \en{to lend} & \en{to borrow} & emprunter \tabularnewline
apporter & \en{to bring} & \en{to take} & emporter \tabularnewline
emporter & \en{to take (away)} & \en{to bring} & apporter \tabularnewline
visiter (un lieu) & \en{to visit} & \en{to visit} & idem (mais on ne « visite » pas une personne : \en{to visit somebody} existe, \en{to see somebody} est plus naturel) \tabularnewline
rendre visite & \en{to visit, to see} & --- & --- \tabularnewline
occasion & \en{opportunity, chance} & \en{occasion} & événement, circonstance \tabularnewline
parfois & \en{sometimes} & \en{often} & souvent \tabularnewline
souvent & \en{often} & --- & --- \tabularnewline
peut-être & \en{maybe, perhaps} & \en{possibly} & éventuellement \tabularnewline
possible & \en{possible} & \en{feasible} & réalisable \tabularnewline
important & \en{important} & \en{impressive} & impressionnant \tabularnewline
impressionnant & \en{impressive} & --- & --- \tabularnewline
intéressant & \en{interesting} & \en{interested} & intéressé \tabularnewline
intéressé & \en{interested} & --- & --- \tabularnewline
ennuyeux & \en{boring} & \en{bored} & ennuyé \tabularnewline
ennuyé & \en{annoyed, bothered} & \en{bored} & qui s'ennuie \tabularnewline
fatigué & \en{tired} & \en{fatigued} & épuisé (médical) \tabularnewline
malade & \en{ill, sick} & \en{malade} (inexistant) & --- \tabularnewline
médecin & \en{doctor} & \en{medical} & médical \tabularnewline
hôpital & \en{hospital} & \en{hotel} & hôtel \tabularnewline
hôtel & \en{hotel} & \en{hospital} & hôpital \tabularnewline
hôpital/hôtel & --- & \en{hostel} & auberge \tabularnewline
monnaie/pièce & \en{coin} & --- & --- \tabularnewline
billet & \en{ticket, banknote} & \en{bill} & facture, addition (AmE) \tabularnewline
facture & \en{invoice, bill} & --- & --- \tabularnewline
recette & \en{recipe} & \en{receipt} & reçu, ticket de caisse \tabularnewline
reçu & \en{receipt} & \en{recipe} & recette de cuisine \tabularnewline
prévenir & \en{to warn} & \en{to prevent} & empêcher \tabularnewline
empêcher & \en{to prevent, stop} & --- & --- \tabularnewline
réussir & \en{to succeed} & \en{to success} (n'existe pas) & --- \tabularnewline
réussir un examen & \en{to pass} & \en{to succeed} & réussir (dans la vie) \tabularnewline
échouer & \en{to fail} & \en{to fall} & tomber \tabularnewline
tomber & \en{to fall} & --- & --- \tabularnewline
monter & \en{to go up, climb} & \en{to mount} & assembler, monter (technique) \tabularnewline
descendre & \en{to go down} & \en{to descend} & descendre (littéraire) \tabularnewline
partir & \en{to leave} & \en{to part} & se séparer \tabularnewline
quitter & \en{to leave} & \en{to quit} & arrêter, démissionner \tabularnewline
arrêter & \en{to stop} & \en{to arrest} & arrêter (police) \tabularnewline
se marier & \en{to get married} & \en{to marry} & épouser \tabularnewline
divorcer & \en{to get divorced} & \en{to divorce} & idem \tabularnewline
enceinte & \en{pregnant} & \en{enceinte} (n'existe pas) & --- \tabularnewline
grossesse & \en{pregnancy} & --- & --- \tabularnewline
accoucher & \en{to give birth} & \en{to accouche} (n'existe pas) & --- \tabularnewline
couche & \en{diaper / nappy} & --- & --- \tabularnewline
couche (strate) & \en{layer} & --- & --- \tabularnewline
élever (un enfant) & \en{to raise, bring up} & \en{to elevate} & soulever, hausser \tabularnewline
monter (prix) & \en{to rise, increase} & \en{to raise} & augmenter (qqch) \tabularnewline
augmenter & \en{to increase} & \en{to augment} (rare) & --- \tabularnewline
diminuer & \en{to decrease} & \en{to diminish} & diminuer (soutenu) \tabularnewline
changer & \en{to change} & \en{to exchange} & échanger \tabularnewline
échanger & \en{to exchange} & --- & --- \tabularnewline
échanger (monnaie) & \en{to change, exchange} & --- & --- \tabularnewline
dépenser & \en{to spend} & \en{to dispense} & distribuer \tabularnewline
dépens & --- & \en{expense} & frais, dépense \tabularnewline
coûter & \en{to cost} & \en{to count} & compter \tabularnewline
compter & \en{to count} & \en{to compute} & calculer \tabularnewline
payer & \en{to pay} & \en{to pay for} & payer pour (au sens figuré) \tabularnewline
vendre & \en{to sell} & \en{to sale} (n'existe pas) & --- \tabularnewline
vente & \en{sale} & --- & --- \tabularnewline
acheter & \en{to buy} & \en{to book} & réserver \tabularnewline
réserver & \en{to book, reserve} & --- & --- \tabularnewline
livrer & \en{to deliver} & \en{to delivery} (n'existe pas) & --- \tabularnewline
livraison & \en{delivery} & --- & --- \tabularnewline
envoyer & \en{to send} & \en{to envy} & envier \tabularnewline
envier & \en{to envy} & --- & --- \tabularnewline
recevoir & \en{to receive} & \en{to perceive} & percevoir \tabularnewline
apercevoir & \en{to notice, glimpse} & --- & --- \tabularnewline
écouter & \en{to listen to} & \en{to hear} & entendre \tabularnewline
entendre & \en{to hear} & --- & --- \tabularnewline
regarder & \en{to watch, look at} & \en{to see} & voir \tabularnewline
voir & \en{to see} & --- & --- \tabularnewline
sentir (odeur) & \en{to smell} & \en{to feel} & ressentir, toucher \tabularnewline
sentir (ressentir) & \en{to feel} & \en{to sense} & pressentir \tabularnewline
toucher & \en{to touch} & \en{to feel} & palper, ressentir \tabularnewline
goûter (n.) & \en{snack} & \en{taste} & goût \tabularnewline
goûter (v.) & \en{to taste} & --- & --- \tabularnewline
savourer & \en{to savour, enjoy} & --- & --- \tabularnewline
assaisonnement & \en{seasoning} & \en{season} & saison \tabularnewline
saison & \en{season} & --- & --- \tabularnewline
temps (météo) & \en{weather} & \en{time} & temps (durée) \tabularnewline
heure & \en{hour, time} & --- & --- \tabularnewline
année & \en{year} & \en{yearly} & annuel \tabularnewline
journée & \en{day} & \en{journey} & voyage \tabularnewline
voyage & \en{trip, journey} & --- & --- \tabularnewline
promenade & \en{walk} & \en{promenade} & front de mer \tabularnewline
aéroport & \en{airport} & \en{port} & port maritime \tabularnewline
port & \en{port, harbour} & --- & --- \tabularnewline
aéroport & --- & \en{airfield} & aérodrome \tabularnewline
voiture & \en{car} & \en{vehicle} & véhicule \tabularnewline
camion & \en{truck / lorry} & \en{camera} & appareil photo \tabularnewline
car (bus) & \en{coach} & \en{car} & voiture \tabularnewline
train & \en{train} & \en{tram} & tramway \tabularnewline
tramway & \en{tram / streetcar} & --- & --- \tabularnewline
métro & \en{underground / subway / metro} & \en{metro} & idem (AmE : \en{subway}) \tabularnewline
billet & \en{ticket} & \en{billet} (rare) & --- \tabularnewline
rue & \en{street} & \en{road} & route \tabularnewline
route & \en{road} & \en{route} & itinéraire \tabularnewline
itinéraire & \en{route, itinerary} & --- & --- \tabularnewline
adresse & \en{address} & \en{advice} & conseil \tabularnewline
conseil & \en{advice} & \en{advert} & publicité \tabularnewline
publicité & \en{advertising, advert} & \en{publicity} & notoriété \tabularnewline
réclame & --- & \en{publicity} & notoriété, battage \tabularnewline
information & \en{information} & \en{inform} & informer \tabularnewline
nouvelles & \en{news} & --- & --- \tabularnewline
journal & \en{newspaper} & \en{journal} & revue spécialisée ; journal intime \tabularnewline
journal intime & \en{diary} & \en{diary} & agenda (BrE) \tabularnewline
agenda & \en{diary / planner} & \en{agenda} & ordre du jour \tabularnewline
ordre du jour & \en{agenda} & --- & --- \tabularnewline
réunion & \en{meeting} & \en{union} & syndicat, union \tabularnewline
syndicat & \en{union, trade union} & --- & --- \tabularnewline
bureau & \en{desk / office} & \en{bureau} & agence, service \tabularnewline
agence & \en{agency} & --- & --- \tabularnewline
entreprise & \en{company, business} & \en{enterprise} & esprit d'entreprise \tabularnewline
usine & \en{factory, plant} & \en{enterprise} & --- \tabularnewline
patron & \en{boss} & \en{patron} & mécène, client régulier \tabularnewline
client & \en{customer, client} & --- & --- \tabularnewline
personne (individu) & \en{person} & \en{people} & les gens \tabularnewline
gens & \en{people} & \en{persons} & personnes (formel) \tabularnewline
homme & \en{man} & \en{human} & humain \tabularnewline
humain & \en{human} & --- & --- \tabularnewline
monsieur & \en{Mr / sir} & \en{monsieur} (n'existe pas) & --- \tabularnewline
madame & \en{Mrs / Ms / madam} & --- & --- \tabularnewline
demoiselle & \en{Miss} & --- & --- \tabularnewline
nom & \en{name} & \en{noun} & nom (grammaire) \tabularnewline
verbe & \en{verb} & \en{adverb} & adverbe \tabularnewline
phrase & \en{sentence} & \en{phrase} & expression, groupe de mots \tabularnewline
peine (sanction) & \en{sentence} & \en{penalty} & pénalité \tabularnewline
erreur & \en{mistake, error} & \en{error} & erreur (technique) \tabularnewline
faute & \en{mistake, fault} & \en{fault} & responsabilité, défaut \tabularnewline
défaut & \en{flaw, fault} & \en{default} & défaut (informatique) \tabularnewline
succès & \en{success} & \en{successful} & réussi \tabularnewline
chance & \en{luck, opportunity} & \en{lucky} & chanceux \tabularnewline
malchance & \en{bad luck} & \en{unlucky} & malchanceux \tabularnewline
\end{longtable}
\end{tableau}

\begin{piege}[Les cinq faux amis les plus coûteux]
\begin{enumerate}
  \item \en{actually} = \textbf{en fait} (pas « actuellement » \ra \en{currently}).
  \item \en{eventually} = \textbf{finalement} (pas « éventuellement » \ra \en{possibly}).
  \item \en{library} = \textbf{bibliothèque} (pas « librairie » \ra \en{bookshop}).
  \item \en{to attend} = \textbf{assister à} (pas « attendre » \ra \en{to wait}).
  \item \en{sensible} = \textbf{raisonnable} (pas « sensible » \ra \en{sensitive}).
\end{enumerate}
\end{piege}

\section{Unité 4 --- Les calques du français à éviter}
\begin{tableau}[Traductions littérales fautives et formes correctes]
\footnotesize
\begin{longtable}{@{}p{5.4cm}p{4.0cm}p{4.4cm}@{}}
\entete{\ko Calque du français} & \entete{\ok Anglais correct} & \entete{Français} \tabularnewline
\endfirsthead
\entete{\ko Calque} & \entete{\ok Correct} & \entete{Français} \tabularnewline
\endhead
\hline
\en{*I have 30 years.} & \en{I'm 30 (years old).} & J'ai 30 ans. \tabularnewline
\en{*I am agree.} & \en{I agree.} & Je suis d'accord. \tabularnewline
\en{*I have cold / hunger.} & \en{I'm cold / hungry.} & J'ai froid / faim. \tabularnewline
\en{*It depends of you.} & \en{It depends on you.} & Cela dépend de toi. \tabularnewline
\en{*I listen music.} & \en{I listen to music.} & J'écoute de la musique. \tabularnewline
\en{*I search my keys.} & \en{I'm looking for my keys.} & Je cherche mes clés. \tabularnewline
\en{*I explain you the problem.} & \en{I'll explain the problem to you.} & Je t'explique le problème. \tabularnewline
\en{*I said you something.} & \en{I told you something.} & Je t'ai dit quelque chose. \tabularnewline
\en{*He is married with a doctor.} & \en{He is married to a doctor.} & Il est marié à une médecin. \tabularnewline
\en{*I will see you to tomorrow.} & \en{I'll see you tomorrow.} & Je te verrai demain. \tabularnewline
\en{*In the morning I go in the office.} & \en{In the morning I go to the office.} & Le matin je vais au bureau. \tabularnewline
\en{*I am in Dakar since 2015.} & \en{I have been in Dakar since 2015.} & Je suis à Dakar depuis 2015. \tabularnewline
\en{*It is difficult for me of understand.} & \en{It is difficult for me to understand.} & C'est difficile pour moi de comprendre. \tabularnewline
\en{*I have fear.} & \en{I'm afraid.} & J'ai peur. \tabularnewline
\en{*She makes a mistake.} \emph{(correct, mais)} & \en{She made a mistake.} (passé simple, pas présent) & Elle a fait une erreur. \tabularnewline
\en{*I demand you a question.} & \en{I'll ask you a question.} & Je te pose une question. \tabularnewline
\en{*This movie passes on TV.} & \en{This movie is on TV.} & Ce film passe à la télé. \tabularnewline
\en{*I have the possibility of coming.} & \en{I can come.} / \en{I have the chance to come.} & J'ai la possibilité de venir. \tabularnewline
\en{*Make attention!} & \en{Be careful!} / \en{Watch out!} & Fais attention ! \tabularnewline
\en{*I am a Senegalese.} & \en{I'm Senegalese.} / \en{I'm from Senegal.} & Je suis sénégalais. \tabularnewline
\en{*It is not possible of doing that.} & \en{It isn't possible to do that.} & Ce n'est pas possible de faire ça. \tabularnewline
\en{*I wait your answer.} & \en{I'm waiting for your answer.} / \en{I look forward to your reply.} & J'attends votre réponse. \tabularnewline
\en{*The weather is very hot.} \emph{(correct)} & \en{It's very hot.} (plus naturel) & Il fait très chaud. \tabularnewline
\en{*I want that you come.} & \en{I want you to come.} & Je veux que tu viennes. \tabularnewline
\en{*I think that it is true.} \emph{(correct)} & \en{I think it's true.} (\en{that} omis) & Je pense que c'est vrai. \tabularnewline
\end{longtable}
\end{tableau}

\begin{mission}[Chapitre 13 --- 21 jours d'expressions]
\textbf{Semaine 1 :} 30 \en{phrasal verbs} (les plus fréquents de la liste) --- 5 par jour,
chacun dans une phrase personnelle, à voix haute.\\
\textbf{Semaine 2 :} 25 idiomes --- 5 par jour, chacun dans une mini-situation jouée à voix haute.\\
\textbf{Semaine 3 :} 25 faux amis --- 5 par jour, avec la phrase correcte.\\
À la fin, vous emploierez \en{figure out, get rid of, look forward to, it depends on,
make ends meet, by the way} --- et votre anglais cessera de sonner traduit.
\end{mission}


% ===== PARTIE V — PARLER ===========================================
\part{Parler : 40 situations, des dialogues modèles, la fluidité}
% =====================================================================
%  CHAPITRE 14 — PARLER : SITUATIONS ET DIALOGUES
% =====================================================================
\chapter{40 situations : dialogues, blocs et jeux de rôle}

Ce chapitre est le c\oe ur pratique du livre. Chaque situation donne :
\textbf{(1)} un dialogue complet, traduit et prononcé ; \textbf{(2)} les blocs réutilisables ;
\textbf{(3)} une variation à jouer. Travaillez une situation par jour pendant 40 jours ---
à voix haute, en jouant les deux rôles.

\begin{astuce}[Le protocole « 5 passes » pour chaque dialogue]
\textbf{Passe 1 :} lisez le dialogue en français seulement --- comprenez la situation.
\textbf{Passe 2 :} lisez-le en anglais à voix haute, lentement, en regardant les
prononciations. \textbf{Passe 3 :} écoutez-le (application de synthèse vocale, ou lisez-le
vous-même en vous enregistrant). \textbf{Passe 4 :} jouez les deux rôles, le livre ouvert.
\textbf{Passe 5 :} jouez-le \textbf{sans le livre}, en changeant trois détails (nom, lieu,
objet). C'est la passe qui compte : elle transforme un texte appris en langage disponible.
\end{astuce}

\section{Situation 1 --- Se présenter (formel)}
\begin{dialoguebox}[1 --- lors d'un rendez-vous professionnel]
\dial{M. Sow}{Good morning. I'm Moussa Sow. Nice to meet you.}{Bonjour. Je suis Moussa Sow. Ravi de vous rencontrer.}
\dial{Mme Clarke}{Good morning, Mr Sow. I'm Helen Clarke. Please, have a seat.}{Bonjour, Monsieur Sow. Je suis Helen Clarke. Asseyez-vous, je vous prie.}
\dial{M. Sow}{Thank you very much.}{Merci beaucoup.}
\dial{Mme Clarke}{So, what do you do, Mr Sow?}{Alors, que faites-vous dans la vie, Monsieur Sow ?}
\dial{M. Sow}{I'm an accountant. I work for a trading company in Dakar.}{Je suis comptable. Je travaille pour une société de négoce à Dakar.}
\dial{Mme Clarke}{How long have you been working there?}{Depuis combien de temps y travaillez-vous ?}
\dial{M. Sow}{For about six years now. I started as an assistant.}{Depuis environ six ans. J'ai commencé comme assistant.}
\dial{Mme Clarke}{And where are you from originally?}{Et d'où venez-vous à l'origine ?}
\dial{M. Sow}{I'm from Thiès, but I've been living in Dakar since 2018.}{Je viens de Thiès, mais j'habite Dakar depuis 2018.}
\dial{Mme Clarke}{Lovely. Shall we start, then?}{Parfait. Commençons-nous alors ?}
\dial{M. Sow}{Of course. I'm ready.}{Bien sûr. Je suis prêt.}
\end{dialoguebox}

\begin{tableau}[Blocs réutilisables]
\begin{tabularx}{\linewidth}{@{}p{6.2cm}X@{}}
\entete{Anglais} & \entete{Français} \tabularnewline
\hline
\en{I'm\ldots\ / My name is\ldots} & Je suis\ldots\ / Je m'appelle\ldots \tabularnewline
\en{Nice to meet you.} & Ravi de vous rencontrer. \tabularnewline
\en{Likewise. / Same here.} & Moi de même. \tabularnewline
\en{I'm a\ldots\ (job). I work for / as / in\ldots} & Je suis\ldots\ Je travaille pour / comme / dans\ldots \tabularnewline
\en{I've been working there for six years.} & J'y travaille depuis six ans. \tabularnewline
\en{I'm from\ldots\ but I live in\ldots} & Je viens de\ldots\ mais j'habite à\ldots \tabularnewline
\en{How do you do?} (très formel) & Comment allez-vous ? (première rencontre) \tabularnewline
\en{Shall we start?} & Commençons-nous ? \tabularnewline
\end{tabularx}
\end{tableau}

\begin{mission}[Variation]
Rejouez le dialogue en vous présentant \textbf{vous}, avec votre vrai métier, votre vraie
ville, votre vraie durée d'expérience. Trois fois de suite, chaque fois plus vite.
\end{mission}

\section{Situation 2 --- Se présenter (informel) et \en{small talk}}
\begin{dialoguebox}[2 --- à une fête de famille]
\dial{Ali}{Hi! I don't think we've met. I'm Ali.}{Salut ! Je ne crois pas qu'on se connaisse. Je suis Ali.}
\dial{Sarah}{Hi Ali, I'm Sarah. How do you know the family?}{Salut Ali, je suis Sarah. Comment connais-tu la famille ?}
\dial{Ali}{Ousmane is my cousin. And you?}{Ousmane est mon cousin. Et toi ?}
\dial{Sarah}{I work with his sister. So, what do you do?}{Je travaille avec sa s\oe ur. Alors, que fais-tu ?}
\dial{Ali}{I'm a teacher --- I teach mathematics at a high school.}{Je suis enseignant --- j'enseigne les mathématiques au lycée.}
\dial{Sarah}{Oh, that must be interesting. How long have you been teaching?}{Oh, ça doit être intéressant. Depuis combien de temps enseignes-tu ?}
\dial{Ali}{About five years. It's hard work, but I love it. And what about you?}{Environ cinq ans. C'est un travail difficile, mais j'adore. Et toi ?}
\dial{Sarah}{I'm a nurse at the regional hospital.}{Je suis infirmière à l'hôpital régional.}
\dial{Ali}{Nice! Do you live around here?}{Super ! Tu habites dans le coin ?}
\dial{Sarah}{Yes, just ten minutes away. By the way, the food is amazing, isn't it?}{Oui, à dix minutes d'ici. Au fait, la nourriture est délicieuse, non ?}
\dial{Ali}{It really is. Have you tried the fish?}{Vraiment. Tu as goûté le poisson ?}
\dial{Sarah}{Not yet --- I will!}{Pas encore --- je vais le faire !}
\end{dialoguebox}

\begin{tableau}[Le \en{small talk} : 12 amorces universelles]
\begin{tabularx}{\linewidth}{@{}p{6.4cm}X@{}}
\entete{Amorce} & \entete{Emploi} \tabularnewline
\hline
\en{Lovely day, isn't it?} & Le temps (sujet sans risque) \tabularnewline
\en{How do you know everyone here?} & Fête, mariage, réunion de famille \tabularnewline
\en{So, what do you do?} & La question standard --- répondez puis \textbf{retournez-la} \tabularnewline
\en{How long have you lived here?} & Résidence \tabularnewline
\en{Did you watch the match last night?} & Sport (le meilleur sujet universel) \tabularnewline
\en{The food is great, isn't it?} & Repas, réception \tabularnewline
\en{Have you been here before?} & Lieu nouveau \tabularnewline
\en{How was your trip?} & Voyage \tabularnewline
\en{What brings you here?} & Raison de la présence (professionnel) \tabularnewline
\en{Do you have any plans for the weekend?} & Vendredi \tabularnewline
\en{How's your week going?} & Lundi ou mardi \tabularnewline
\en{Is this seat taken?} & Transport, café, salle d'attente \tabularnewline
\end{tabularx}
\end{tableau}

\begin{astuce}[La règle des trois relances]
Une conversation meurt quand on répond par une phrase fermée. Entraînez-vous à toujours
ajouter : \textbf{(1)} une réponse, \textbf{(2)} un détail, \textbf{(3)} une question en
retour.\\
\en{``Do you like football?''} --- \en{``Yes, I love it. I play every Saturday with my
brothers. Do you play?''}\\
Trois éléments au lieu d'un mot : la conversation vit.
\end{astuce}

\section{Situation 3 --- Salutations, au revoir, et formules de politesse}
\begin{tableau}[Toutes les salutations, du plus formel au plus familier]
\footnotesize
\begin{longtable}{@{}p{5.2cm}p{4.4cm}p{4.2cm}@{}}
\entete{Rencontrer} & \entete{Prendre congé} & \entete{Contexte} \tabularnewline
\endfirsthead
\entete{Rencontrer} & \entete{Prendre congé} & \entete{Contexte} \tabularnewline
\endhead
\hline
\en{Good morning.} (jusqu'à midi) & \en{Goodbye.} & Formel, universel \tabularnewline
\en{Good afternoon.} (12 h - 18 h) & \en{Have a nice day.} & Formel, commerce \tabularnewline
\en{Good evening.} (après 18 h) & \en{Have a good evening.} & Formel \tabularnewline
\en{Hello. / Hi.} & \en{See you. / See you later.} & Neutre \tabularnewline
\en{Hi there!} & \en{See you tomorrow / on Monday.} & Amical \tabularnewline
\en{Hey!} & \en{Take care.} & Familier, chaleureux \tabularnewline
\en{How are you?} & \en{Catch you later!} & Familier \tabularnewline
\en{How are you doing?} & \en{I'm off --- see you!} & Familier \tabularnewline
\en{How's it going?} & \en{Bye for now.} & Neutre \tabularnewline
\en{What's up? / How come?} & \en{Have a good one.} & Très familier (AmE) \tabularnewline
\en{Long time no see!} & \en{It was great seeing you.} & Retrouvailles \tabularnewline
\en{Pleased to meet you.} & \en{It was a pleasure.} & Première rencontre formelle \tabularnewline
\end{longtable}
\end{tableau}

\begin{tableau}[Répondre à « comment ça va ? »]
\begin{tabularx}{\linewidth}{@{}p{5.4cm}X@{}}
\entete{Réponse} & \entete{Nuance} \tabularnewline
\hline
\en{I'm fine, thanks. And you?} & Standard, poli, toujours correct. \tabularnewline
\en{Pretty good, thanks!} & Positif, naturel. \tabularnewline
\en{Not bad. You?} & Neutre, courant. \tabularnewline
\en{Can't complain.} & « Je ne peux pas me plaindre » --- très idiomatique. \tabularnewline
\en{Same old, same old.} & « Comme d'habitude ». \tabularnewline
\en{I'm doing well, thank you.} & Plus formel. \tabularnewline
\en{Great, thanks! I'm really happy today.} & Enthousiaste. \tabularnewline
\en{To be honest, I'm exhausted.} & Sincère --- ouvre une vraie conversation. \tabularnewline
\en{I've been better, but it's fine.} & Pudeur sur un moment difficile. \tabularnewline
\end{tabularx}
\end{tableau}

\begin{tableau}[Politesse : 20 formules]
\begin{tabularx}{\linewidth}{@{}X X@{}}
\en{Thank you very much.} --- merci beaucoup & \en{Thanks a lot.} --- merci bien \\
\en{I really appreciate it.} --- j'apprécie vraiment & \en{You're welcome.} --- de rien \\
\en{My pleasure.} --- avec plaisir & \en{No problem. / Don't mention it.} --- pas de souci \\
\en{I'm sorry.} --- je suis désolé & \en{Excuse me.} --- pardon (pour attirer l'attention) \\
\en{Sorry to bother you.} --- pardon de vous déranger & \en{Pardon? / Sorry?} --- pardon ? (faire répéter) \\
\en{Could you repeat that, please?} --- pourriez-vous répéter ? & \en{Could you speak more slowly?} --- pourriez-vous parler plus lentement ? \\
\en{Would you mind\ldots?} --- cela vous ennuierait-il de\ldots ? & \en{Could I possibly\ldots?} --- pourrais-je éventuellement\ldots ? \\
\en{I'd appreciate it if you could\ldots} --- je vous serais reconnaissant de\ldots & \en{That's very kind of you.} --- c'est très aimable à vous \\
\en{After you.} --- je vous en prie (passer devant) & \en{Go ahead.} --- allez-y \\
\en{Make yourself at home.} --- faites comme chez vous & \en{Help yourself.} --- servez-vous \\
\en{Congratulations!} --- félicitations ! & \en{Best wishes!} --- meilleurs v\oe ux ! \\
\en{Good luck!} --- bonne chance ! & \en{Take care of yourself.} --- prenez soin de vous \\
\en{Get well soon.} --- bon rétablissement & \en{My condolences.} --- mes condoléances \\
\en{I'm sorry for your loss.} --- je suis désolé pour votre perte & \en{Safe travels!} --- bon voyage ! \\
\end{tabularx}
\end{tableau}

\section{Situation 4 --- Au café et au restaurant}
\begin{dialoguebox}[4 --- commander, demander l'addition]
\dial{Serveur}{Good afternoon! Table for how many?}{Bonjour ! Une table pour combien ?}
\dial{Vous}{Two, please. Do you have a table outside?}{Deux, s'il vous plaît. Avez-vous une table en terrasse ?}
\dial{Serveur}{Yes, right this way. Here are the menus.}{Oui, par ici. Voici les cartes.}
\dial{Vous}{Thank you. What do you recommend?}{Merci. Que recommandez-vous ?}
\dial{Serveur}{The grilled fish is very popular today.}{Le poisson grillé est très demandé aujourd'hui.}
\dial{Vous}{That sounds good. I'll have the grilled fish, please.}{Cela semble bon. Je vais prendre le poisson grillé, s'il vous plaît.}
\dial{Ami}{And I'll take the chicken with rice.}{Et je prendrai le poulet avec du riz.}
\dial{Serveur}{Anything to drink?}{Quelque chose à boire ?}
\dial{Vous}{A large bottle of water and two orange juices, please.}{Une grande bouteille d'eau et deux jus d'orange, s'il vous plaît.}
\dial{Serveur}{Right away. Do you have any allergies?}{Tout de suite. Avez-vous des allergies ?}
\dial{Vous}{No, we're fine, thank you.}{Non, tout va bien, merci.}
\variante{(vingt minutes plus tard)}
\dial{Serveur}{Is everything all right?}{Tout se passe bien ?}
\dial{Vous}{Yes, it's delicious, thank you. Could we have some bread, please?}{Oui, c'est délicieux, merci. Pourrions-nous avoir du pain ?}
\dial{Serveur}{Of course. Enjoy your meal!}{Bien sûr. Bon appétit !}
\dial{Vous}{Could we have the bill, please?}{Pourrions-nous avoir l'addition, s'il vous plaît ?}
\dial{Serveur}{Here you are. Service is not included.}{Voici. Le service n'est pas compris.}
\dial{Vous}{No problem --- we'll leave a tip. Keep the change.}{Pas de problème --- nous laisserons un pourboire. Gardez la monnaie.}
\end{dialoguebox}

\begin{tableau}[Les 15 blocs du restaurant]
\begin{tabularx}{\linewidth}{@{}p{6.4cm}X@{}}
\entete{Bloc} & \entete{Emploi} \tabularnewline
\hline
\en{A table for two, please.} & En arrivant \tabularnewline
\en{Do you have a table by the window?} & Préférence \tabularnewline
\en{I've booked a table under the name\ldots} & Réservation \tabularnewline
\en{Could I see the menu, please?} & Demander la carte \tabularnewline
\en{What's the soup of the day?} & Demander le plat du jour \tabularnewline
\en{What do you recommend?} & Demander un conseil \tabularnewline
\en{I'll have the\ldots, please.} & \textbf{La formule de commande} --- jamais \en{*I want} \tabularnewline
\en{I'd like the\ldots, please.} & Variante polie \tabularnewline
\en{Could I have\ldots\ without onions?} & Exclusion d'un ingrédient \tabularnewline
\en{I'm vegetarian / I don't eat meat.} & Régime \tabularnewline
\en{Could we have some water, please?} & Demande en cours de repas \tabularnewline
\en{It was delicious, thank you.} & Compliment \tabularnewline
\en{Could we have the bill / check, please?} & Addition \tabularnewline
\en{Is service included?} & Question pourboire \tabularnewline
\en{Can I pay by card? / Keep the change.} & Paiement \tabularnewline
\end{tabularx}
\end{tableau}

\begin{piege}[\en{*I want a coffee}]
Grammaticalement correct, socialement brutal. En anglais, \en{I want} sonne comme un ordre
d'enfant. Utilisez \en{I'd like\ldots}, \en{I'll have\ldots}, \en{Could I have\ldots?}
--- trois secondes de politesse qui changent tout le ton d'un échange.
\end{piege}

\section{Situation 5 --- Au marché et dans les magasins}
\begin{dialoguebox}[5 --- acheter, demander le prix, négocier]
\dial{Vendeur}{Hello! Can I help you find something?}{Bonjour ! Je peux vous aider à trouver quelque chose ?}
\dial{Vous}{Yes, please. I'm looking for a shirt --- a blue one.}{Oui, s'il vous plaît. Je cherche une chemise --- une bleue.}
\dial{Vendeur}{We have this one, and this one in a lighter blue.}{Nous avons celle-ci, et celle-là en bleu plus clair.}
\dial{Vous}{How much is this one?}{Combien coûte celle-ci ?}
\dial{Vendeur}{It's 12,000 francs.}{C'est 12\,000 francs.}
\dial{Vous}{That's a bit expensive. Can you do 10,000?}{C'est un peu cher. Pouvez-vous la faire à 10\,000 ?}
\dial{Vendeur}{Tell you what --- 11,000, and it's a very good price.}{Voici ce que je vous propose --- 11\,000, et c'est un très bon prix.}
\dial{Vous}{Deal. Can I try it on first?}{D'accord. Puis-je l'essayer d'abord ?}
\dial{Vendeur}{Of course, the fitting room is over there.}{Bien sûr, la cabine d'essayage est par là.}
\dial{Vous}{It fits perfectly. I'll take it.}{Elle me va parfaitement. Je la prends.}
\dial{Vendeur}{Would you like a bag? Do you need a receipt?}{Voulez-vous un sac ? Avez-vous besoin d'un reçu ?}
\dial{Vous}{A bag, please. No receipt --- I'm paying in cash.}{Un sac, s'il vous plaît. Pas de reçu --- je paie en espèces.}
\end{dialoguebox}

\begin{tableau}[Achats : 18 blocs]
\begin{tabularx}{\linewidth}{@{}p{6.4cm}X@{}}
\entete{Bloc} & \entete{Emploi} \tabularnewline
\hline
\en{I'm just looking, thanks.} & Refuser poliment de l'aide \tabularnewline
\en{I'm looking for\ldots} & Annoncer sa recherche \tabularnewline
\en{Do you have this in a bigger / smaller size?} & Tailles \tabularnewline
\en{Do you have it in another colour?} & Couleur \tabularnewline
\en{Where can I find\ldots?} & Localisation dans le magasin \tabularnewline
\en{How much is it? / How much does it cost?} & Prix \tabularnewline
\en{That's too expensive for me.} & Négocier \tabularnewline
\en{Can you give me a discount?} & Demander une remise \tabularnewline
\en{What's your best price?} & Négocier \tabularnewline
\en{Can I try it on?} & Essayer \tabularnewline
\en{Where is the fitting room?} & Cabine \tabularnewline
\en{It fits well / It's too tight / too loose.} & Essayage \tabularnewline
\en{It suits you!} & Compliment (aller à la personne) \tabularnewline
\en{I'll take it.} & Acheter \tabularnewline
\en{Can I pay by card? / Is cash okay?} & Paiement \tabularnewline
\en{Could I have a receipt, please?} & Reçu \tabularnewline
\en{Can I get a refund / exchange it?} & Retour, échange \tabularnewline
\en{Is it still under guarantee / warranty?} & Garantie \tabularnewline
\end{tabularx}
\end{tableau}

\section{Situation 6 --- Demander son chemin}
\begin{dialoguebox}[6 --- dans la rue]
\dial{Vous}{Excuse me, sorry to bother you. How do I get to the central market?}{Pardon, excusez-moi de vous déranger. Comment aller au marché central ?}
\dial{Passant}{The central market? It's about a fifteen-minute walk from here.}{Le marché central ? C'est à environ quinze minutes à pied d'ici.}
\dial{Vous}{Oh, I see. Which way should I go?}{Ah, d'accord. Par où dois-je aller ?}
\dial{Passant}{Go straight on for two blocks, then turn left at the traffic lights.}{Allez tout droit pendant deux rues, puis tournez à gauche au feu.}
\dial{Vous}{Straight on, then left at the lights. Is it far from the mosque?}{Tout droit, puis à gauche au feu. C'est loin de la mosquée ?}
\dial{Passant}{No, it's just opposite the mosque, next to the pharmacy.}{Non, c'est juste en face de la mosquée, à côté de la pharmacie.}
\dial{Vous}{Great. And is there a bus that goes there?}{Parfait. Et y a-t-il un bus qui y va ?}
\dial{Passant}{Yes, bus number 4 stops right in front of it.}{Oui, le bus numéro 4 s'arrête juste devant.}
\dial{Vous}{Thank you very much for your help.}{Merci beaucoup pour votre aide.}
\dial{Passant}{No problem. You can't miss it.}{Pas de souci. Vous ne pouvez pas le rater.}
\end{dialoguebox}

\begin{tableau}[Les indications : tout le vocabulaire du déplacement]
\begin{tabularx}{\linewidth}{@{}p{5.4cm}X@{}}
\entete{Instruction} & \entete{Sens} \tabularnewline
\hline
\en{Go straight on / ahead.} & Allez tout droit. \tabularnewline
\en{Turn left / right.} & Tournez à gauche / à droite. \tabularnewline
\en{Take the first / second turning.} & Prenez la première / deuxième rue. \tabularnewline
\en{At the roundabout, take the second exit.} & Au rond-point, prenez la deuxième sortie. \tabularnewline
\en{Cross the street / the bridge.} & Traversez la rue / le pont. \tabularnewline
\en{Go past the bank.} & Passez devant la banque. \tabularnewline
\en{It's on your left / right.} & C'est sur votre gauche / droite. \tabularnewline
\en{It's next to / opposite / behind / in front of\ldots} & C'est à côté de / en face de / derrière / devant\ldots \tabularnewline
\en{It's at the end of the street.} & C'est au bout de la rue. \tabularnewline
\en{It's about five minutes' walk.} & C'est à cinq minutes à pied. \tabularnewline
\en{You can't miss it.} & Vous ne pouvez pas le rater. \tabularnewline
\en{I'm afraid I'm not from around here.} & Je ne suis pas du coin. \tabularnewline
\en{Let me check on my phone.} & Laissez-moi vérifier sur mon téléphone. \tabularnewline
\end{tabularx}
\end{tableau}

\section{Situation 7 --- À l'hôtel}
\begin{dialoguebox}[7 --- arrivée et problème de chambre]
\dial{Réceptionniste}{Good evening, sir. How can I help you?}{Bonsoir, monsieur. Comment puis-je vous aider ?}
\dial{Vous}{Good evening. I have a reservation under the name Diop, for three nights.}{Bonsoir. J'ai une réservation au nom de Diop, pour trois nuits.}
\dial{Réceptionniste}{Let me check\ldots\ Yes, here it is. A single room with breakfast. Could I see your ID, please?}{Je vérifie\ldots\ Oui, la voici. Une chambre simple avec petit-déjeuner. Puis-je voir votre pièce d'identité ?}
\dial{Vous}{Of course, here's my passport.}{Bien sûr, voici mon passeport.}
\dial{Réceptionniste}{Thank you. You're in room 304, on the third floor. Here's your key. Breakfast is from 7 to 10 am.}{Merci. Vous êtes en chambre 304, au troisième étage. Voici votre clé. Le petit-déjeuner est servi de 7 h à 10 h.}
\dial{Vous}{Is there Wi-Fi in the rooms?}{Y a-t-il du wi-fi dans les chambres ?}
\dial{Réceptionniste}{Yes, it's free. The password is on your key card holder.}{Oui, il est gratuit. Le mot de passe est sur l'étui de votre carte.}
\dial{Vous}{Perfect, thank you.}{Parfait, merci.}
\variante{(le lendemain matin, à la réception)}
\dial{Vous}{Excuse me, I'm sorry, but there's a problem with my room.}{Excusez-moi, je suis désolé, mais il y a un problème avec ma chambre.}
\dial{Réceptionniste}{I'm sorry to hear that. What seems to be the problem?}{Je suis navrée de l'apprendre. Quel est le problème ?}
\dial{Vous}{The air conditioning isn't working and there's no hot water.}{La climatisation ne fonctionne pas et il n'y a pas d'eau chaude.}
\dial{Réceptionniste}{I do apologise. I'll send someone up right away. If we can't fix it, we'll move you to another room.}{Je vous présente mes excuses. J'envoie quelqu'un tout de suite. Si nous ne pouvons pas réparer, nous vous changerons de chambre.}
\dial{Vous}{That would be great. Thank you for your help.}{Ce serait parfait. Merci pour votre aide.}
\end{dialoguebox}

\begin{tableau}[Se plaindre sans agressivité --- la technique du sandwich]
\begin{tabularx}{\linewidth}{@{}p{6.2cm}X@{}}
\entete{Étape} & \entete{Formule} \tabularnewline
\hline
1. Excuse polie & \en{Excuse me, I'm sorry, but\ldots} / \en{I'm afraid there's a problem with\ldots} \tabularnewline
2. Fait précis & \en{The air conditioning isn't working.} / \en{I've been waiting for 40 minutes.} / \en{This isn't what I ordered.} \tabularnewline
3. Demande claire & \en{Could you send someone, please?} / \en{Is there anything you can do?} / \en{Could I have a refund?} \tabularnewline
4. Remerciement & \en{Thank you, I really appreciate it.} \tabularnewline
\end{tabularx}
\end{tableau}

\begin{astuce}[\en{I'm afraid} et \en{I'm sorry} adoucissent tout]
\en{I'm afraid there's a mistake.} (Je crains qu'il y ait une erreur.) sonne bien mieux que
\en{There's a mistake.} En anglais, la plainte commence presque toujours par une excuse ---
même quand c'est vous qui avez raison. C'est culturel, pas logique.
\end{astuce}

\section{Situation 8 --- À l'aéroport}
\begin{dialoguebox}[8 --- enregistrement et contrôle]
\dial{Agent}{Good morning. Where are you flying to today?}{Bonjour. Où vous envolez-vous aujourd'hui ?}
\dial{Vous}{To London, via Casablanca.}{Vers Londres, via Casablanca.}
\dial{Agent}{May I see your passport and ticket, please? Do you have any luggage to check in?}{Puis-je voir votre passeport et votre billet ? Avez-vous des bagages à enregistrer ?}
\dial{Vous}{Here you are. Just one suitcase --- this is my hand luggage.}{Voici. Une seule valise --- ceci est mon bagage à main.}
\dial{Agent}{Thank you. Do you have anything sharp or liquid in your hand luggage?}{Merci. Avez-vous des objets tranchants ou des liquides dans votre bagage à main ?}
\dial{Vous}{No, nothing like that.}{Non, rien de ce genre.}
\dial{Agent}{Great. Window or aisle seat?}{Parfait. Siège côté hublot ou côté couloir ?}
\dial{Vous}{A window seat, please, if possible.}{Un siège côté hublot, s'il vous plaît, si c'est possible.}
\dial{Agent}{You're on seat 24A. Boarding starts at gate 12 at 10:40.}{Vous êtes au siège 24A. L'embarquement commence à la porte 12 à 10 h 40.}
\dial{Vous}{Thank you. How long is the flight?}{Merci. Combien de temps dure le vol ?}
\dial{Agent}{About six hours. Have a pleasant flight!}{Environ six heures. Bon vol !}
\variante{(au contrôle des passeports)}
\dial{Officier}{Passport, please. What's the purpose of your visit?}{Votre passeport, s'il vous plaît. Quel est le motif de votre visite ?}
\dial{Vous}{It's for a training course. I'll be here for two weeks.}{C'est pour une formation. Je resterai deux semaines.}
\dial{Officier}{Where will you be staying?}{Où logerez-vous ?}
\dial{Vous}{At the Hilton Hotel, in the city centre.}{À l'hôtel Hilton, au centre-ville.}
\dial{Officier}{Thank you. Enjoy your stay.}{Merci. Bon séjour.}
\end{dialoguebox}

\section{Situation 9 --- Chez le médecin}
\begin{dialoguebox}[9 --- consultation]
\dial{Médecin}{Good morning. What seems to be the problem?}{Bonjour. Quel est le problème ?}
\dial{Vous}{I've been feeling ill for three days. I have a headache and a sore throat.}{Je me sens malade depuis trois jours. J'ai mal à la tête et mal à la gorge.}
\dial{Médecin}{Do you have a fever?}{Avez-vous de la fièvre ?}
\dial{Vous}{Yes, I think so. I felt hot last night.}{Oui, je crois. J'avais chaud hier soir.}
\dial{Médecin}{Let me take your temperature\ldots\ 38.5. Open your mouth and say ``ah''. Any cough?}{Je prends votre température\ldots\ 38,5. Ouvrez la bouche et dites « ah ». De la toux ?}
\dial{Vous}{A dry cough, mostly at night.}{Une toux sèche, surtout la nuit.}
\dial{Médecin}{Are you allergic to any medication? Do you take any other medicines?}{Êtes-vous allergique à un médicament ? Prenez-vous d'autres traitements ?}
\dial{Vous}{No allergies, and I only take vitamins.}{Aucune allergie, et je ne prends que des vitamines.}
\dial{Médecin}{It looks like a bad cold, nothing serious. I'll prescribe something for the fever and the cough. Rest for three days and drink plenty of water.}{On dirait un gros rhume, rien de grave. Je vais vous prescrire quelque chose pour la fièvre et la toux. Reposez-vous trois jours et buvez beaucoup d'eau.}
\dial{Vous}{Thank you, doctor. Should I come back if it doesn't get better?}{Merci, docteur. Dois-je revenir si cela ne s'améliore pas ?}
\dial{Médecin}{Yes, come back in a week if you're not better. Take care.}{Oui, revenez dans une semaine si vous n'allez pas mieux. Prenez soin de vous.}
\end{dialoguebox}

\begin{tableau}[Décrire ses symptômes --- 20 blocs]
\begin{tabularx}{\linewidth}{@{}p{6.4cm}X@{}}
\entete{Bloc} & \entete{Sens} \tabularnewline
\hline
\en{I don't feel well.} & Je ne me sens pas bien. \tabularnewline
\en{I've been feeling ill for three days.} & Je suis malade depuis trois jours. \tabularnewline
\en{I have a headache / stomachache / backache.} & J'ai mal à la tête / au ventre / au dos. \tabularnewline
\en{I have a sore throat.} & J'ai mal à la gorge. \tabularnewline
\en{I have a runny nose and a cough.} & J'ai le nez qui coule et je tousse. \tabularnewline
\en{I have a fever / a temperature.} & J'ai de la fièvre. \tabularnewline
\en{I feel dizzy / sick.} & J'ai la tête qui tourne / je me sens mal. \tabularnewline
\en{It hurts when I swallow.} & Ça fait mal quand j'avale. \tabularnewline
\en{My leg is swollen.} & Ma jambe est enflée. \tabularnewline
\en{I've been vomiting.} & J'ai vomi. \tabularnewline
\en{I can't sleep.} & Je n'arrive pas à dormir. \tabularnewline
\en{I'm always tired.} & Je suis toujours fatigué. \tabularnewline
\en{I have an allergy to\ldots} & Je suis allergique à\ldots \tabularnewline
\en{I take medicine for\ldots} & Je prends un traitement pour\ldots \tabularnewline
\en{I need a prescription.} & J'ai besoin d'une ordonnance. \tabularnewline
\en{I need a medical certificate.} & J'ai besoin d'un certificat médical. \tabularnewline
\en{How long will it take to recover?} & Combien de temps pour guérir ? \tabularnewline
\en{Is it serious?} & Est-ce grave ? \tabularnewline
\en{Should I see a specialist?} & Dois-je voir un spécialiste ? \tabularnewline
\en{What should I avoid?} & Que dois-je éviter ? \tabularnewline
\end{tabularx}
\end{tableau}

\section{Situation 10 --- Au téléphone}
\begin{dialoguebox}[10 --- appel professionnel]
\dial{Vous}{Good morning, Diop and Sons, how can I help you?}{Bonjour, Diop et Fils, comment puis-je vous aider ?}
\dial{Appelant}{Hello, this is John Mensah from Bright Ltd. May I speak to Mr Diop, please?}{Bonjour, ici John Mensah de Bright Ltd. Puis-je parler à Monsieur Diop ?}
\dial{Vous}{I'm afraid he's in a meeting at the moment. Can I take a message?}{Je crains qu'il ne soit en réunion pour le moment. Puis-je prendre un message ?}
\dial{Appelant}{Yes, please. Could you ask him to call me back about the delivery?}{Oui, s'il vous plaît. Pourriez-vous lui demander de me rappeler au sujet de la livraison ?}
\dial{Vous}{Of course. Could you spell your surname for me, please?}{Bien sûr. Pourriez-vous épeler votre nom de famille ?}
\dial{Appelant}{It's M-E-N-S-A-H.}{C'est M-E-N-S-A-H.}
\dial{Vous}{Thank you. And what's the best number to reach you on?}{Merci. Et quel est le meilleur numéro pour vous joindre ?}
\dial{Appelant}{It's 77 123 45 67. I'm available after 2 pm.}{C'est le 77 123 45 67. Je suis disponible après 14 h.}
\dial{Vous}{Right, I'll make sure he gets the message. Anything else?}{D'accord, je m'assure qu'il reçoive le message. Autre chose ?}
\dial{Appelant}{No, that's all. Thanks for your help.}{Non, c'est tout. Merci de votre aide.}
\dial{Vous}{You're welcome. Goodbye, Mr Mensah.}{De rien. Au revoir, Monsieur Mensah.}
\end{dialoguebox}

\begin{tableau}[Le téléphone : 20 blocs]
\begin{tabularx}{\linewidth}{@{}p{6.4cm}X@{}}
\entete{Bloc} & \entete{Emploi} \tabularnewline
\hline
\en{Hello, this is Moussa speaking.} & Se présenter au téléphone (\textbf{jamais} \en{*I am Moussa}) \tabularnewline
\en{May I speak to\ldots, please?} & Demander quelqu'un \tabularnewline
\en{Could I ask who's calling?} & Demander qui appelle \tabularnewline
\en{Hold on, please. / One moment, please.} & Faire patienter \tabularnewline
\en{I'm afraid he's not available.} & Personne n'est joignable \tabularnewline
\en{Can I take a message?} & Prendre un message \tabularnewline
\en{Could you ask him to call me back?} & Demander un rappel \tabularnewline
\en{I'll call back later.} & Je rappellerai plus tard. \tabularnewline
\en{Could you spell that for me?} & Faire épeler \tabularnewline
\en{Sorry, you're breaking up.} & Vous grésillez / ça coupe. \tabularnewline
\en{The line is very bad. Can you hear me?} & La ligne est mauvaise. \tabularnewline
\en{I'm sorry, could you repeat that?} & Faire répéter \tabularnewline
\en{Could you speak a bit louder?} & Plus fort \tabularnewline
\en{Let me put you through.} & Je vous passe quelqu'un. \tabularnewline
\en{I think we've been cut off.} & Nous avons été coupés. \tabularnewline
\en{I'll send you an email to confirm.} & Je confirme par écrit. \tabularnewline
\en{Thanks for calling. Goodbye.} & Conclure \tabularnewline
\en{Is now a good time to talk?} & Vérifier la disponibilité \tabularnewline
\en{I'm calling about\ldots} & Annoncer l'objet de l'appel \tabularnewline
\en{Sorry to call so late.} & S'excuser de l'heure \tabularnewline
\end{tabularx}
\end{tableau}

\section{Situation 11 --- L'entretien d'embauche}
\begin{dialoguebox}[11 --- les six questions incontournables]
\dial{Recruteur}{Thanks for coming in. Could you tell me a little about yourself?}{Merci d'être venu. Pouvez-vous vous présenter brièvement ?}
\dial{Vous}{Of course. I'm an accountant with six years' experience. I've been working for a trading company in Dakar since 2020, where I handle invoicing and monthly reports.}{Bien sûr. Je suis comptable avec six ans d'expérience. Je travaille pour une société de négoce à Dakar depuis 2020, où je m'occupe de la facturation et des rapports mensuels.}
\dial{Recruteur}{Why do you want to work for us?}{Pourquoi souhaitez-vous travailler chez nous ?}
\dial{Vous}{Your company is growing fast in West Africa, and I'd like to work on larger accounts. I also value your training programmes.}{Votre entreprise se développe rapidement en Afrique de l'Ouest, et j'aimerais travailler sur des comptes plus importants. J'apprécie aussi vos programmes de formation.}
\dial{Recruteur}{What are your strengths and weaknesses?}{Quelles sont vos forces et vos faiblesses ?}
\dial{Vous}{My strength is accuracy --- I rarely make errors in reporting. My weakness used to be public speaking, but I've been taking courses, and I've improved a lot.}{Ma force est la rigueur --- je fais rarement des erreurs dans les rapports. Ma faiblesse était la prise de parole en public, mais je suis des formations et je me suis beaucoup amélioré.}
\dial{Recruteur}{Where do you see yourself in five years?}{Où vous voyez-vous dans cinq ans ?}
\dial{Vous}{Ideally, leading a small team and helping the company expand into new markets.}{Idéalement, à la tête d'une petite équipe, en aidant l'entreprise à s'étendre vers de nouveaux marchés.}
\dial{Recruteur}{Why did you leave your last job?}{Pourquoi avez-vous quitté votre dernier emploi ?}
\dial{Vous}{There were no opportunities for promotion, so I'm looking for a new challenge. I left on very good terms.}{Il n'y avait pas de perspectives de promotion, je cherche donc un nouveau défi. Je suis parti en très bons termes.}
\dial{Recruteur}{Do you have any questions for us?}{Avez-vous des questions à nous poser ?}
\dial{Vous}{Yes --- what does a typical week look like in this role, and what would success look like after six months?}{Oui --- à quoi ressemble une semaine type dans ce poste, et à quoi ressemblerait la réussite après six mois ?}
\end{dialoguebox}

\begin{astuce}[Les trois réponses à préparer par écrit]
Avant tout entretien, rédigez et mémorisez \textbf{60 secondes} sur : (1) qui vous êtes
professionnellement ; (2) pourquoi cette entreprise ; (3) votre plus belle réussite, avec un
chiffre. Ces trois blocs couvrent 60\,\% de l'entretien. Entraînez-les à voix haute dix fois
--- et enregistrez-vous.
\end{astuce}

\section{Situation 12 --- Donner son opinion et débattre}
\begin{dialoguebox}[12 --- discussion entre amis]
\dial{Ami}{What do you think about working abroad? Everyone wants to leave.}{Que penses-tu de travailler à l'étranger ? Tout le monde veut partir.}
\dial{Vous}{Well, it depends. On the one hand, it's a great opportunity to learn and earn more.}{Eh bien, cela dépend. D'un côté, c'est une grande occasion d'apprendre et de gagner plus.}
\dial{Ami}{Exactly! And you can support your family back home.}{Exactement ! Et tu peux aider ta famille restée au pays.}
\dial{Vous}{That's true. But on the other hand, I think we also need skilled people here. If everyone leaves, who builds the country?}{C'est vrai. Mais d'un autre côté, je pense que nous avons aussi besoin de gens qualifiés ici. Si tout le monde part, qui construit le pays ?}
\dial{Ami}{I see your point, but honestly, without opportunities here, people have no choice.}{Je comprends ton point de vue, mais honnêtement, sans opportunités ici, les gens n'ont pas le choix.}
\dial{Vous}{I couldn't agree more on that. That's why I think the real question is how to create opportunities here.}{Sur ce point, je suis tout à fait d'accord. C'est pourquoi je pense que la vraie question est de savoir comment créer des opportunités ici.}
\dial{Ami}{Fair enough. We're not going to solve that today, though!}{C'est juste. Mais nous n'allons pas résoudre ça aujourd'hui !}
\dial{Vous}{No, we're not --- but it was a good conversation.}{Non, en effet --- mais c'était une bonne conversation.}
\end{dialoguebox}

\begin{tableau}[Exprimer, nuancer, contredire --- 25 blocs]
\begin{tabularx}{\linewidth}{@{}p{6.4cm}X@{}}
\entete{Bloc} & \entete{Fonction} \tabularnewline
\hline
\en{In my opinion, \ldots} & donner un avis neutre \tabularnewline
\en{I think / I believe / I feel that\ldots} & avis courant \tabularnewline
\en{As far as I'm concerned, \ldots} & avis personnel marqué \tabularnewline
\en{From my point of view, \ldots} & point de vue \tabularnewline
\en{I'd say that\ldots} & je dirais que \tabularnewline
\en{It seems to me that\ldots} & il me semble que \tabularnewline
\en{I'm convinced that\ldots} & je suis convaincu que \tabularnewline
\en{I couldn't agree more.} & je suis totalement d'accord \tabularnewline
\en{That's exactly what I think.} & c'est exactement ce que je pense \tabularnewline
\en{You've got a point there.} & tu marques un point \tabularnewline
\en{Fair enough.} & c'est juste (accord partiel) \tabularnewline
\en{I agree up to a point, but\ldots} & je suis d'accord jusqu'à un certain point \tabularnewline
\en{I see what you mean, but\ldots} & je vois ce que tu veux dire, mais \tabularnewline
\en{That's not quite how I see it.} & ce n'est pas tout à fait ainsi que je vois les choses \tabularnewline
\en{I'm afraid I have to disagree.} & je crains de devoir être en désaccord \tabularnewline
\en{With all due respect, \ldots} & avec tout le respect que je vous dois \tabularnewline
\en{On the one hand\ldots on the other hand\ldots} & d'un côté\ldots de l'autre \tabularnewline
\en{It depends on the situation.} & cela dépend de la situation \tabularnewline
\en{That's a good question --- let me think.} & bonne question --- laisse-moi réfléchir \tabularnewline
\en{I'm not sure, but I'd guess that\ldots} & je ne suis pas sûr, mais je suppose que \tabularnewline
\en{To be honest, \ldots} & pour être honnête \tabularnewline
\en{In other words, \ldots} & autrement dit \tabularnewline
\en{Let's agree to disagree.} & restons sur nos positions \tabularnewline
\en{That reminds me of something.} & cela me fait penser à quelque chose \tabularnewline
\en{Going back to what you said earlier\ldots} & pour revenir à ce que tu disais \tabularnewline
\end{tabularx}
\end{tableau}

\section{Situation 13 --- Raconter une histoire}
\begin{dialoguebox}[13 --- une anecdote (passé simple, continu et perfect)]
\dial{Ami}{Have you ever had a really embarrassing moment?}{T'est-il déjà arrivé un moment vraiment embarrassant ?}
\dial{Vous}{Oh yes! Last year, I was travelling to a conference in Banjul.}{Oh oui ! L'année dernière, je voyageais pour une conférence à Banjul.}
\dial{Ami}{What happened?}{Que s'est-il passé ?}
\dial{Vous}{Well, I had arrived early, so I was waiting at the gate. Suddenly, a man started waving at me and shouting my name.}{Eh bien, j'étais arrivé tôt, donc j'attendais à la porte. Soudain, un homme a commencé à me faire signe et à crier mon nom.}
\dial{Ami}{Did you know him?}{Le connaissais-tu ?}
\dial{Vous}{Not at all! It turned out he was waving at his brother, who was standing right behind me.}{Pas du tout ! Il s'est avéré qu'il faisait signe à son frère, qui se tenait juste derrière moi.}
\dial{Ami}{Ha! What did you do?}{Ha ! Qu'as-tu fait ?}
\dial{Vous}{I waved back for about ten seconds before I realised. I've never been so embarrassed in my life.}{J'ai répondu en faisant signe pendant environ dix secondes avant de réaliser. Je n'ai jamais été aussi embarrassé de ma vie.}
\dial{Ami}{That's a great story. Did you ever meet him again?}{C'est une bonne histoire. L'as-tu revu ?}
\dial{Vous}{Actually, yes --- we were on the same plane, and we ended up sitting next to each other!}{En fait oui --- nous étions dans le même avion, et nous avons fini par nous asseoir l'un à côté de l'autre !}
\end{dialoguebox}

\begin{astuce}[La structure d'une bonne anecdote --- 5 étapes]
\textbf{1. Cadre} (passé continu / \en{past perfect}) : \en{It was raining and I was
waiting for the bus.}\\
\textbf{2. Élément déclencheur} (\en{suddenly}, passé simple) : \en{Suddenly, a dog ran
towards me.}\\
\textbf{3. Réaction} : \en{I didn't know what to do, so I climbed onto a bench.}\\
\textbf{4. Résultat} : \en{In the end, the owner came and took it away.}\\
\textbf{5. Commentaire} (\en{present perfect}) : \en{I've never been so scared!}\\
Apprenez cette structure et vous pourrez raconter n'importe quelle histoire de votre vie
--- c'est exactement ce que font les natifs.
\end{astuce}

\section{Situation 14 --- Parler de son pays et de sa culture}
\begin{dialoguebox}[14 --- avec un étranger curieux]
\dial{Étranger}{So you're from Senegal --- what's it like?}{Tu viens donc du Sénégal --- comment c'est ?}
\dial{Vous}{It's a beautiful country on the west coast of Africa. The people are very welcoming, and we have a long tradition of hospitality --- we call it \en{teranga}.}{C'est un beau pays sur la côte ouest de l'Afrique. Les gens sont très accueillants, et nous avons une longue tradition d'hospitalité --- nous l'appelons la teranga.}
\dial{Étranger}{What's the food like? I've never tried Senegalese food.}{Comment est la nourriture ? Je n'ai jamais goûté la cuisine sénégalaise.}
\dial{Vous}{You should try \en{thieboudienne} --- it's fish with rice and vegetables. It's our national dish. We also eat a lot of millet, peanuts and fresh fish.}{Tu devrais essayer le thieboudienne --- c'est du poisson avec du riz et des légumes. C'est notre plat national. Nous mangeons aussi beaucoup de mil, d'arachide et de poisson frais.}
\dial{Étranger}{That sounds delicious. And what languages do people speak?}{Cela semble délicieux. Et quelles langues les gens parlent-ils ?}
\dial{Vous}{French is the official language, but most people speak Wolof at home. There are also Pulaar, Serer, Diola and others.}{Le français est la langue officielle, mais la plupart des gens parlent wolof à la maison. Il y a aussi le pulaar, le sérère, le diola et d'autres.}
\dial{Étranger}{Wow. Is it hot all year round?}{Ouah. Il fait chaud toute l'année ?}
\dial{Vous}{Pretty much. We have a dry season from November to May and a rainy season from June to October. It can get very hot in April.}{À peu près. Nous avons une saison sèche de novembre à mai et une saison des pluies de juin à octobre. Il peut faire très chaud en avril.}
\dial{Étranger}{I'd love to visit one day.}{J'aimerais y aller un jour.}
\dial{Vous}{You'd be very welcome. I'd be happy to show you around.}{Tu serais le bienvenu. Je serais heureux de te faire visiter.}
\end{dialoguebox}

\begin{tableau}[Parler de son pays : 20 blocs]
\begin{tabularx}{\linewidth}{@{}p{6.6cm}X@{}}
\entete{Bloc} & \entete{Sens} \tabularnewline
\hline
\en{I'm from\ldots, in West Africa.} & Je viens de\ldots, en Afrique de l'Ouest. \tabularnewline
\en{It's a country of about 17 million people.} & C'est un pays d'environ 17 millions d'habitants. \tabularnewline
\en{The capital is\ldots, but I'm from\ldots} & La capitale est\ldots, mais je viens de\ldots \tabularnewline
\en{We speak\ldots\ at home and\ldots\ at school.} & Nous parlons\ldots à la maison et\ldots à l'école. \tabularnewline
\en{Our national dish is called\ldots} & Notre plat national s'appelle\ldots \tabularnewline
\en{It's made with rice, fish and vegetables.} & C'est fait avec du riz, du poisson et des légumes. \tabularnewline
\en{The climate is hot and dry for most of the year.} & Le climat est chaud et sec la majeure partie de l'année. \tabularnewline
\en{We have a dry season and a rainy season.} & Nous avons une saison sèche et une saison des pluies. \tabularnewline
\en{Football is by far the most popular sport.} & Le football est de loin le sport le plus populaire. \tabularnewline
\en{Traditional wrestling is also very popular.} & La lutte traditionnelle est aussi très populaire. \tabularnewline
\en{Family is at the centre of our culture.} & La famille est au c\oe ur de notre culture. \tabularnewline
\en{People are known for their hospitality.} & Les gens sont connus pour leur hospitalité. \tabularnewline
\en{Most people are Muslim, and there are also many Christians.} & La plupart des gens sont musulmans, et il y a aussi beaucoup de chrétiens. \tabularnewline
\en{We celebrate\ldots, which is a very important festival.} & Nous célébrons\ldots, qui est une fête très importante. \tabularnewline
\en{The economy is based on agriculture, fishing and services.} & L'économie repose sur l'agriculture, la pêche et les services. \tabularnewline
\en{Tourism is growing, especially around\ldots} & Le tourisme se développe, surtout autour de\ldots \tabularnewline
\en{It's about two hours from the capital by road.} & C'est à environ deux heures de la capitale par la route. \tabularnewline
\en{If you visit, you should definitely see\ldots} & Si tu visites, tu dois absolument voir\ldots \tabularnewline
\en{The best time to come is between November and February.} & La meilleure période pour venir est entre novembre et février. \tabularnewline
\en{I'd be happy to show you around.} & Je serais heureux de te faire visiter. \tabularnewline
\end{tabularx}
\end{tableau}

\section{Situation 15 --- La réunion de travail}
\begin{dialoguebox}[15 --- animer et participer]
\dial{Animatrice}{Good morning, everyone. Thanks for coming. The purpose of today's meeting is to review the second-quarter results. Let's get started.}{Bonjour à tous. Merci d'être venus. L'objet de la réunion d'aujourd'hui est de passer en revue les résultats du deuxième trimestre. Commençons.}
\dial{Vous}{Before we start, could I just clarify the agenda? Are we discussing the budget today as well?}{Avant de commencer, puis-je préciser l'ordre du jour ? Discutons-nous aussi du budget aujourd'hui ?}
\dial{Animatrice}{Good question. We'll come to the budget at the end, if we have time. First, Moussa, could you give us an overview of sales?}{Bonne question. Nous aborderons le budget à la fin, si nous avons le temps. D'abord, Moussa, pourriez-vous nous donner un aperçu des ventes ?}
\dial{Vous}{Sure. Sales increased by 12\,\% compared to last year, which is above our target. However, costs also rose, so the margin is slightly lower.}{Bien sûr. Les ventes ont augmenté de 12\,\% par rapport à l'année dernière, ce qui dépasse notre objectif. Cependant, les coûts ont aussi augmenté, donc la marge est légèrement plus basse.}
\dial{Collègue}{Sorry to interrupt, but does that include the new market?}{Pardon de vous interrompre, mais cela inclut-il le nouveau marché ?}
\dial{Vous}{No, it doesn't. If we include it, the increase is closer to 18\,\%.}{Non. Si on l'inclut, l'augmentation est plus proche de 18\,\%.}
\dial{Animatrice}{That's very encouraging. Let's move on to the next point.}{C'est très encourageant. Passons au point suivant.}
\dial{Collègue}{Could you go back to the costs for a moment? I'd like to understand where they come from.}{Pourriez-vous revenir sur les coûts un instant ? J'aimerais comprendre d'où ils viennent.}
\dial{Vous}{Of course. I'll send you the detailed breakdown after the meeting.}{Bien sûr. Je vous enverrai le détail après la réunion.}
\dial{Animatrice}{Right, to sum up: sales are up, costs need attention, and Moussa will send the details. Any other business?}{Bien, pour résumer : les ventes sont en hausse, les coûts demandent attention, et Moussa enverra le détail. Autres questions ?}
\dial{Vous}{Just one thing --- could we schedule a follow-up meeting next week?}{Une seule chose --- pourrions-nous prévoir une réunion de suivi la semaine prochaine ?}
\dial{Animatrice}{Good idea. I'll send an invitation. Thanks everyone --- meeting adjourned.}{Bonne idée. J'enverrai une invitation. Merci à tous --- la séance est levée.}
\end{dialoguebox}

\begin{tableau}[Animer et participer : 25 blocs]
\begin{tabularx}{\linewidth}{@{}p{6.6cm}X@{}}
\entete{Bloc} & \entete{Fonction} \tabularnewline
\hline
\en{Let's get started.} & Ouvrir \tabularnewline
\en{The purpose of this meeting is to\ldots} & Annoncer l'objet \tabularnewline
\en{There are three items on the agenda.} & Annoncer le plan \tabularnewline
\en{First of all, I'd like to\ldots} & Premier point \tabularnewline
\en{Let's move on to the next point.} & Passer au suivant \tabularnewline
\en{Could you give us an overview of\ldots?} & Demander un aperçu \tabularnewline
\en{To put it simply, \ldots} & Simplifier \tabularnewline
\en{The figures show that\ldots} & Présenter des chiffres \tabularnewline
\en{Sales rose / fell by 12\,\%.} & Hausse / baisse \tabularnewline
\en{Compared to last year / last quarter\ldots} & Comparaison \tabularnewline
\en{Sorry to interrupt, but\ldots} & Interrompre poliment \tabularnewline
\en{Could I come in here?} & Prendre la parole \tabularnewline
\en{Can I just add something?} & Ajouter \tabularnewline
\en{Could you go back to that point?} & Revenir en arrière \tabularnewline
\en{Could you clarify what you mean by\ldots?} & Faire préciser \tabularnewline
\en{If I understand correctly, you're saying that\ldots} & Vérifier sa compréhension \tabularnewline
\en{That's a valid point, but\ldots} & Nuancer \tabularnewline
\en{I'd like to make a suggestion.} & Proposer \tabularnewline
\en{What do you think about\ldots?} & Solliciter un avis \tabularnewline
\en{Shall we take a vote?} & Voter \tabularnewline
\en{Let's put that aside for now.} & Mettre de côté \tabularnewline
\en{We're running out of time.} & Le temps presse \tabularnewline
\en{To sum up, \ldots} & Résumer \tabularnewline
\en{Any other business?} & Questions diverses \tabularnewline
\en{I'll send you the minutes.} & J'enverrai le compte rendu. \tabularnewline
\end{tabularx}
\end{tableau}

\section{Situation 16 --- S'excuser, remercier, demander un service}
\begin{tableau}[Trois moments clés de la vie sociale]
\begin{tabularx}{\linewidth}{@{}p{5.0cm}p{4.6cm}X@{}}
\entete{Situation} & \entete{Bloc} & \entete{Réponse attendue} \tabularnewline
\hline
S'excuser (léger) & \en{Sorry! / I'm sorry.} & \en{That's okay. / No problem.} \tabularnewline
S'excuser (sérieux) & \en{I'm really sorry about what happened. It was my fault.} & \en{Don't worry about it. / I understand.} \tabularnewline
S'excuser (retard) & \en{Sorry I'm late --- the traffic was terrible.} & \en{That's fine, we've just started.} \tabularnewline
S'excuser et réparer & \en{I apologise. How can I make it right?} & \en{Just\ldots\ and that'll be fine.} \tabularnewline
Remercier (simple) & \en{Thanks! / Thank you.} & \en{You're welcome.} \tabularnewline
Remercier (sincère) & \en{Thank you so much, I really appreciate it.} & \en{It was my pleasure.} \tabularnewline
Remercier (écrit) & \en{Thanks a lot for your help with the report.} & \en{Any time.} \tabularnewline
Demander un service & \en{Could you do me a favour?} & \en{Sure, what do you need?} \tabularnewline
Demander directement & \en{Would you mind helping me with this?} & \en{Of course not --- here you are.} \tabularnewline
Refuser poliment & \en{I'd love to, but I'm afraid I can't.} & \en{No problem, maybe next time.} \tabularnewline
Reporter & \en{Could we do it another time? I'm swamped today.} & \en{Sure, let me know when.} \tabularnewline
\end{tabularx}
\end{tableau}

\begin{piege}[\en{I'm sorry} a deux sens]
\en{I'm sorry} = « je m'excuse » \emph{et} « je suis désolé pour toi » (compassion).
\en{``My father died.''} --- \en{``I'm so sorry.''} est la seule bonne réponse --- jamais
\en{*It's not a problem} ou \en{*You're welcome}. Cette confusion fait commettre des
impairs énormes ; retenez : devant une mauvaise nouvelle, \en{I'm sorry} (ou
\en{That's awful, I'm so sorry for your loss}).
\end{piege}

\section{Situation 17 --- Ne pas comprendre : les phrases de survie}
\begin{memo}[Les 12 phrases qui vous sauvent dans chaque conversation]
\begin{enumerate}
  \item \en{Sorry, could you repeat that, please?} --- Pardon, pourriez-vous répéter ?
  \item \en{Could you speak more slowly, please?} --- Pourriez-vous parler plus lentement ?
  \item \en{Could you say that again more simply?} --- Pourriez-vous le redire plus simplement ?
  \item \en{What does ``\ldots'' mean?} --- Que veut dire « \ldots » ?
  \item \en{How do you say ``\ldots'' in English?} --- Comment dit-on « \ldots » en anglais ?
  \item \en{Sorry, I didn't catch that.} --- Pardon, je n'ai pas saisi.
  \item \en{Do you mean\ldots?} --- Voulez-vous dire\ldots ?
  \item \en{Let me check I understood: you're saying that\ldots} --- Je vérifie : vous dites que\ldots
  \item \en{I'm still learning English, so please bear with me.} --- J'apprends encore l'anglais, soyez indulgent.
  \item \en{Could you write it down for me?} --- Pourriez-vous me l'écrire ?
  \item \en{Could you spell that, please?} --- Pourriez-vous l'épeler ?
  \item \en{I understand most of it, but I'm not sure about\ldots} --- Je comprends l'essentiel, mais pas sûr pour\ldots
\end{enumerate}
\end{memo}

\begin{astuce}[La phrase qui détend tout le monde]
\en{My English isn't perfect yet, but I'm working on it every day.}
(Mon anglais n'est pas encore parfait, mais je travaille dessus chaque jour.)
Cette phrase, dite avec le sourire, produit trois effets : elle excuse d'avance vos erreurs,
elle montre votre sérieux, et elle pousse l'interlocuteur à ralentir et à vous aider.
Apprenez-la par c\oe ur --- elle vaut de l'or dans les six premiers mois.
\end{astuce}

\section{Les 23 autres situations à préparer vous-même}
Vous avez maintenant la mécanique. Les situations suivantes sont listées avec leurs trois
blocs essentiels. Construisez vous-même un dialogue de 8 répliques pour chacune --- c'est
l'exercice le plus rentable du livre.

\begin{tableau}[Situations 18 à 40 --- blocs de départ]
\footnotesize
\begin{longtable}{@{}p{3.4cm}p{5.6cm}p{4.8cm}@{}}
\entete{Situation} & \entete{Blocs essentiels} & \entete{Votre mission} \tabularnewline
\endfirsthead
\entete{Situation} & \entete{Blocs essentiels} & \entete{Mission} \tabularnewline
\endhead
\hline
\textbf{18. À la pharmacie} & \en{Do you have something for\ldots?} / \en{Do I need a prescription?} / \en{How often should I take it?} & Jouer la scène complète \tabularnewline
\textbf{19. À la banque} & \en{I'd like to open an account.} / \en{What's the exchange rate?} / \en{I'd like to transfer money abroad.} & Ajouter les frais et délais \tabularnewline
\textbf{20. À la poste} & \en{I'd like to send this to France.} / \en{How long will it take?} / \en{Registered, please.} & Peser, payer, confirmer \tabularnewline
\textbf{21. Chez le coiffeur} & \en{Just a trim, please.} / \en{How short would you like it?} / \en{Not too short at the back.} & Négocier la coupe \tabularnewline
\textbf{22. En taxi} & \en{Could you take me to\ldots?} / \en{How much will it be?} / \en{Could you drop me here, please?} & Discuter le prix avant \tabularnewline
\textbf{23. Louer une voiture} & \en{I'd like to rent a car for three days.} / \en{Is insurance included?} / \en{Do I need to return it with a full tank?} & Ajouter l'état des lieux \tabularnewline
\textbf{24. Acheter en ligne} & \en{When will my order arrive?} / \en{Can I change the delivery address?} / \en{I'd like a refund.} & Réclamation complète \tabularnewline
\textbf{25. Panne technique} & \en{My phone won't turn on.} / \en{It's still under warranty.} / \en{Can you fix it today?} & Décrire le problème précisément \tabularnewline
\textbf{26. À la salle de sport} & \en{I'd like to sign up.} / \en{How much is the monthly fee?} / \en{Do you have a personal trainer?} & Poser 5 questions \tabularnewline
\textbf{27. Parler football} & \en{Did you watch the match?} / \en{We won two-nil!} / \en{The referee was terrible.} & Débat de 5 minutes \tabularnewline
\textbf{28. À l'école (parent)} & \en{How is my child doing?} / \en{He needs to improve his English.} / \en{When are the exams?} & Rendez-vous complet \tabularnewline
\textbf{29. Avec un professeur} & \en{Could you explain this again?} / \en{When is the deadline?} / \en{Can I have an extension?} & Demander un délai poliment \tabularnewline
\textbf{30. Au commissariat} & \en{I'd like to report a theft.} / \en{Someone stole my phone.} / \en{I need a police report.} & Déclaration détaillée \tabularnewline
\textbf{31. Urgence} & \en{Call an ambulance!} / \en{There's been an accident.} / \en{He's not breathing.} & Message d'alerte complet \tabularnewline
\textbf{32. Inviter} & \en{Would you like to\ldots?} / \en{Are you free on Saturday?} / \en{Let's meet at 7.} & Trois invitations différentes \tabularnewline
\textbf{33. Refuser} & \en{I'd love to, but I can't.} / \en{Something came up.} / \en{Can we do it another time?} & Refus + proposition \tabularnewline
\textbf{34. Faire connaissance en ligne} & \en{How long have you been learning English?} / \en{Would you like to practise together?} / \en{Shall we do a voice call?} & Premier appel de 5 min \tabularnewline
\textbf{35. Parler de religion} & \en{I pray five times a day.} / \en{We're fasting this month.} / \en{Everyone is welcome to join us.} & Échange respectueux \tabularnewline
\textbf{36. Parler de sa famille} & \en{I come from a big family.} / \en{I have four brothers and two sisters.} / \en{My parents live in\ldots} & Description de 2 min \tabularnewline
\textbf{37. Décrire une photo} & \en{In the foreground, there's\ldots} / \en{In the background, I can see\ldots} / \en{It looks like they're\ldots} & Décrire 3 photos réelles \tabularnewline
\textbf{38. Présenter un projet} & \en{The aim of this project is to\ldots} / \en{There are three stages.} / \en{I'd be happy to take questions.} & Exposé de 3 minutes \tabularnewline
\textbf{39. Passer un appel vidéo} & \en{Can you see me? Can you hear me?} / \en{You're on mute.} / \en{Let me share my screen.} & Gérer les problèmes techniques \tabularnewline
\textbf{40. Faire un discours / toast} & \en{I'd like to say a few words.} / \en{Thank you all for coming.} / \en{Let's raise our glasses.} & Discours de 90 secondes \tabularnewline
\end{longtable}
\end{tableau}

\begin{mission}[Chapitre 14 --- les 40 jours]
Une situation par jour, pendant 40 jours. Protocole quotidien (25 minutes) :\\
\textbf{5 min} --- lire le dialogue et ses blocs, à voix haute.\\
\textbf{5 min} --- écouter (synthèse vocale ou lecture enregistrée par vous).\\
\textbf{10 min} --- jouer les deux rôles sans le livre, en changeant trois détails,
et en vous enregistrant.\\
\textbf{5 min} --- créer 5 cartes Anki avec les blocs les plus utiles du jour.\\
Au jour 40, vous posséderez environ \textbf{500 blocs de conversation actifs} --- c'est
précisément ce qui permet de tenir une conversation de 30 minutes.
\end{mission}

% =====================================================================
%  CHAPITRE 15 — TECHNIQUES DE L'ORAL : PARLER SANS TRADUIRE
% =====================================================================
\chapter{Parler sans traduire : les techniques de fluidité}

Connaître des mots et des règles ne suffit pas : il faut \textbf{produire en temps réel}.
C'est un sport, pas un savoir. Ce chapitre donne les techniques qui transforment un savoir
passif en parole fluide.

\section{Le problème réel : le délai de traduction}
Quand vous traduisez du français dans votre tête, votre cerveau exécute cinq opérations :
comprendre le français \ra choisir les mots \ra les traduire \ra conjuguer \ra prononcer.
Résultat : trois à cinq secondes de silence, une intonation française, et des calques.
L'objectif de ce chapitre est de réduire ce circuit à \textbf{une seule opération} :
penser directement en anglais.

\begin{tableau}[Les cinq niveaux de la pensée en anglais --- progressez dans l'ordre]
\begin{tabularx}{\linewidth}{@{}p{1.4cm}p{4.0cm}X@{}}
\entete{Niv.} & \entete{Exercice} & \entete{Détail} \tabularnewline
\hline
\textbf{1} & Nommer le monde & 2 minutes, trois fois par jour : nommez en anglais tout ce
que vous voyez. \en{table, phone, window, a red car, two children, a man walking}. Aucune
traduction autorisée --- si vous ne savez pas, contournez (voir p. suivante). \tabularnewline
\textbf{2} & Phrases simples sur soi & \en{I'm tired. I want tea. It's hot. I have work at
two.} --- 20 phrases par jour, à voix haute, en marchant. \tabularnewline
\textbf{3} & Narration de l'action en cours & Commentez ce que vous faites :
\en{I'm cutting the onions. Now I'm adding oil. It smells good.} \tabularnewline
\textbf{4} & Récit du passé et projets & Le soir : 3 minutes sur votre journée. Le matin :
3 minutes sur ce que vous allez faire. \tabularnewline
\textbf{5} & Pensée abstraite et débat & Pensez directement en anglais quand vous réfléchissez
à un problème, préparez un argument, imaginez une réponse à une objection. \tabularnewline
\end{tabularx}
\end{tableau}

\begin{memo}[Le déclencheur visuel]
Choisissez un objet que vous voyez cent fois par jour --- votre téléphone, une porte, votre
montre. Chaque fois que vous le voyez, dites une phrase en anglais (même dans votre tête).
C'est un « rappel » automatique qui installe l'anglais dans votre journée sans effort
d'organisation.
\end{memo}

\section{Le contournement : parler sans connaître le mot}
La compétence la plus utile d'un locuteur intermédiaire n'est pas le vocabulaire : c'est
\textbf{la capacité à expliquer un mot manquant}. C'est ce que font les natifs eux-mêmes
--- ils disent \en{thingy}, \en{whatchamacallit}, \en{you know, the thing\ldots}

\begin{tableau}[Les sept stratégies de contournement]
\begin{tabularx}{\linewidth}{@{}p{3.6cm}p{4.2cm}X@{}}
\entete{Stratégie} & \entete{Formule} & \entete{Exemple (mot manquant : \en{corkscrew})} \tabularnewline
\hline
Définition fonctionnelle & \en{It's a thing you use to\ldots} & \en{It's a thing you use to open a wine bottle.} \tabularnewline
Description & \en{It's like a small metal tool with a spiral.} & \en{It's made of metal, it's small, it has a spiral.} \tabularnewline
Catégorie + précision & \en{It's a kind of\ldots} & \en{It's a kind of kitchen tool.} \tabularnewline
Synonyme simple & Remplacez le mot difficile par un mot facile & \en{a bottle opener} (pour du vin) \tabularnewline
Exemple d'usage & \en{You need it when you\ldots} & \en{You need it when you buy a bottle of wine.} \tabularnewline
Geste et mimique & Mimez, montrez, dessinez & (tourner le poignet) \tabularnewline
Demande d'aide & \en{How do you say\ldots?} / \en{What's the word for\ldots?} & \en{What's the word for the thing you open wine with?} \tabularnewline
\end{tabularx}
\end{tableau}

\begin{astuce}[Entraînement « mot interdit » --- 5 min/jour]
Choisissez 5 mots courants (\en{chair, water, phone, car, bread}) et \textbf{interdisez-vous
de les prononcer} pendant 3 minutes de description. Expliquez chacun sans le nommer :
\en{``It's a thing with four legs, you sit on it at the table.''} C'est exactement
l'exercice que font les interprètes. En trois semaines, vous ne serez plus jamais bloqué.
\end{astuce}

\begin{tableau}[12 phrases qui vous donnent du temps pour réfléchir]
\begin{tabularx}{\linewidth}{@{}p{6.4cm}X@{}}
\entete{Phrase} & \entete{Effet} \tabularnewline
\hline
\en{That's a good question --- let me think.} & Gagne 3 secondes, flatte l'interlocuteur. \tabularnewline
\en{Well, it depends\ldots} & Gagne du temps en ouvrant une nuance. \tabularnewline
\en{How can I put this\ldots} & Signale que vous cherchez la formulation juste. \tabularnewline
\en{What I mean is\ldots} & Reformule après une hésitation. \tabularnewline
\en{Let me see\ldots} & Neutre, très courant. \tabularnewline
\en{To be honest, I'm not sure, but I'd say\ldots} & Répond sans savoir --- bien mieux que le silence. \tabularnewline
\en{Sorry, I've lost my thread. Where was I?} & Reprend après un trou de mémoire. \tabularnewline
\en{It's on the tip of my tongue.} & Le mot est presque là. \tabularnewline
\en{Give me a second, I'll find the word.} & Annonce la recherche. \tabularnewline
\en{In French we say\ldots\ --- how do you say it in English?} & Transforme le trou en échange. \tabularnewline
\en{Anyway, the main point is\ldots} & Reprend le fil. \tabularnewline
\en{Sorry, could I start again?} & Recommencer proprement. \tabularnewline
\end{tabularx}
\end{tableau}

\section{Les \emph{fillers} naturels : remplir sans bégayer}
Les natifs ne parlent pas sans pause : ils remplissent leurs pauses avec des mots-outils.
Les connaître vous évite le « euh » français (qui sonne étranger) et vous donne l'air naturel.

\begin{tableau}[25 \emph{fillers} et connecteurs oraux]
\begin{tabularx}{\linewidth}{@{}p{3.4cm}p{3.0cm}X@{}}
\entete{\en{Filler}} & \entete{Équivalent} & \entete{Emploi et exemple} \tabularnewline
\hline
\en{well} & « eh bien » & Pour commencer une réponse : \en{Well, I think it's fine.} \tabularnewline
\en{you know} & « tu vois » & Recherche d'accord : \en{It's hard, you know?} \tabularnewline
\en{I mean} & « je veux dire » & Reformulation : \en{I mean, it's not that bad.} \tabularnewline
\en{kind of / sorta} & « un peu, genre » & Adoucissement : \en{I'm kind of tired.} \tabularnewline
\en{actually} & « en fait » & Correction d'attente : \en{Actually, I've never been there.} \tabularnewline
\en{basically} & « en gros » & Résumé : \en{Basically, we need more time.} \tabularnewline
\en{literally} & « littéralement » & Insistance (très familier, à doser) : \en{I literally waited two hours.} \tabularnewline
\en{to be honest} & « honnêtement » & Franchise : \en{To be honest, I didn't like it.} \tabularnewline
\en{I guess / I suppose} & « je suppose » & Hésitation polie : \en{I guess you're right.} \tabularnewline
\en{sort of} & « en quelque sorte » & Approximation. \tabularnewline
\en{like} & « genre » (très familier) & \en{It was, like, ten metres high.} --- à éviter en contexte formel. \tabularnewline
\en{anyway} & « bref, enfin » & Reprendre le fil : \en{Anyway, let's move on.} \tabularnewline
\en{so} & « donc, alors » & Enchaînement : \en{So, what happened next?} \tabularnewline
\en{right} & « d'accord, bon » & Validation : \en{Right, I understand now.} \tabularnewline
\en{OK} & « d'accord » & Transition. \tabularnewline
\en{er / erm} & « euh » & Le vrai « euh » anglais --- bref et rare. \tabularnewline
\en{let's see} & « voyons » & Réflexion. \tabularnewline
\en{you see} & « tu comprends » & Explication. \tabularnewline
\en{as I was saying} & « comme je disais » & Reprendre après interruption. \tabularnewline
\en{by the way} & « au fait » & Digression. \tabularnewline
\en{that said} & « cela dit » & Nuance. \tabularnewline
\en{on the other hand} & « d'un autre côté » & Contraste. \tabularnewline
\en{in fact} & « en fait » & Précision (attention : pas \en{actually} partout). \tabularnewline
\en{for example / for instance} & « par exemple » & Illustration. \tabularnewline
\en{at the end of the day} & « au bout du compte » & Conclusion. \tabularnewline
\end{tabularx}
\end{tableau}

\begin{piege}[Trois \en{fillers} maximum]
Trop de \en{you know} et de \en{like} donne l'impression inverse de la fluidité : un discours
hésitant. Choisissez \textbf{deux ou trois} marqueurs qui vous plaisent (\en{well},
\en{I mean}, \en{actually}, \en{to be honest}) et utilisez-les proprement. Le silence bref
vaut mieux qu'un \en{like} répété dix fois.
\end{piege}

\section{Les sept drills de fluidité}

\subsection{Drill 1 --- Les 60 secondes}
Choisissez un mot (\en{water}, \en{family}, \en{work}, \en{rain}) et parlez \textbf{60
secondes sans vous arrêter} sur ce mot : ce que c'est, ce que vous aimez, un souvenir,
une opinion. Chronométrez. Trois fois par jour. Au bout d'un mois, vous tiendrez
90 secondes sans blanc.

\subsection{Drill 2 --- La chaîne de mots}
Partez d'un mot, associez librement, et justifiez chaque enchaînement en une phrase :
\en{rain \ra umbrella \ra my grandmother \ra her house \ra the village\ldots}
Cinq minutes, à voix haute. Ce drill casse la peur de la phrase suivante.

\subsection{Drill 3 --- Le récit en trois temps}
Racontez la même anecdote \textbf{trois fois} : la première en 3 minutes, la deuxième en
2 minutes, la troisième en 1 minute. Vous apprenez à \textbf{condenser} et à éliminer le
superflu --- c'est l'exercice préféré des interprètes.

\subsection{Drill 4 --- Le double rôle}
Jouez seul un dialogue de ce livre, en faisant les deux voix, debout, avec gestes.
Changez une variable à chaque tour (vous êtes pressé, vous êtes poli, vous êtes agacé).

\subsection{Drill 5 --- La description d'image}
Prenez une photo (téléphone, journal, Internet) et décrivez-la pendant 2 minutes :
\en{In the foreground\ldots\ In the background\ldots\ It looks like\ldots\ They might be\ldots\
Maybe he's feeling\ldots} --- l'usage de \en{might} et \en{looks like} fait passer votre
anglais au niveau B2.

\subsection{Drill 6 --- La traduction différée}
Écoutez une phrase anglaise (podcast), mettez pause, \textbf{résumez-la en anglais} avec vos
propres mots. Jamais de traduction française : c'est ce qui installe la pensée directe.

\subsection{Drill 7 --- Le miroir à questions}
Préparez 30 questions sur votre vie (voir liste ci-dessous) et répondez-y à voix haute,
chacune en 45 secondes, sans préparation. Enregistrez. Répétez deux fois par semaine.

\begin{tableau}[30 questions sur vous-même --- le kit de l'entretien oral]
\footnotesize
\begin{multicols}{2}
\begin{anglais}
\begin{enumerate}\setlength{\itemsep}{1pt}
\item Can you tell me about yourself?
\item Where did you grow up?
\item What does your name mean?
\item Describe your family.
\item What do you do for a living?
\item Why did you choose that job?
\item What's a typical day like for you?
\item What did you do last weekend?
\item What are you planning to do next year?
\item Why are you learning English?
\item How do you learn best?
\item What was your favourite subject at school?
\item Describe your best friend.
\item What do you do in your free time?
\item What kind of music do you like?
\item Do you play any sport?
\item What's your favourite food, and how is it made?
\item Describe your city or village.
\item What's the best thing about your country?
\item What would you change about it?
\item Tell me about a memorable trip.
\item What was the happiest moment of your life?
\item Describe a difficulty you overcame.
\item Who do you admire, and why?
\item What are your three strengths?
\item What are you working on improving?
\item Where do you see yourself in five years?
\item What's your opinion on social media?
\item If you could learn anything, what would it be?
\item What advice would you give to a beginner learning English?
\end{enumerate}
\end{anglais}
\end{multicols}
\end{tableau}

\section{Fluidité contre précision : l'équilibre}
\begin{tableau}[Quand privilégier quoi]
\begin{tabularx}{\linewidth}{@{}p{4.6cm}p{4.6cm}X@{}}
\entete{Objectif} & \entete{Quand} & \entete{Comment} \tabularnewline
\hline
\textbf{Fluidité} & Conversation réelle, appels, discussions de groupe & Parlez sans vous arrêter, ne corrigez pas en chemin, acceptez les fautes, restez dans le rythme. \tabularnewline
\textbf{Précision} & Exercices écrits, préparation d'un exposé, enregistrement de travail & Corrigez chaque phrase, vérifiez les temps, les articles, la prononciation. \tabularnewline
\textbf{Règle 80/20} & Toute l'année & 80\,\% du temps de parole en mode fluide, 20\,\% en mode précis. L'inverse est l'erreur classique des francophones scolarisés. \tabularnewline
\end{tabularx}
\end{tableau}

\begin{astuce}[Ne vous corrigez jamais en pleine phrase]
S'interrompre pour corriger un \en{-s} oublié détruit le rythme, la confiance et le plaisir
de l'interlocuteur. Terminez votre idée, puis --- si la faute compte vraiment --- reformulez
la phrase entière correctement, une fois. Mieux : notez-la et corrigez-la le soir, hors
conversation.
\end{astuce}

\section{Le débit, les pauses, la voix}
\begin{itemize}
  \item \textbf{Débit :} visez 120 à 140 mots par minute en conversation. Trop lent =
  incompréhension chez l'autre ; trop rapide = fautes et finales mangées.
  \item \textbf{Pauses :} une pause courte après chaque groupe de sens. La pause n'est pas
  un échec, c'est de la ponctuation orale.
  \item \textbf{Volume :} parlez légèrement plus fort qu'en français. L'anglais porte.
  \item \textbf{Regard :} regardez votre interlocuteur, pas le plafond --- le regard fuyant
  fait croire que vous ne comprenez pas.
  \item \textbf{Respiration :} respirez par le ventre avant de commencer une phrase longue.
  Une phrase dite sans souffle se termine en murmure --- et l'information importante est perdue.
\end{itemize}

\section{Le plan anti-blocage en 6 réflexes}
\begin{enumerate}
  \item \textbf{Le mot manque ?} Contournez (\en{it's a thing you use to\ldots}).
  \item \textbf{La grammaire bloque ?} Simplifiez : sujet + verbe + complément, au présent.
  Une phrase simple et juste vaut mieux qu'une phrase complexe et fausse.
  \item \textbf{Vous n'avez pas compris ?} \en{Sorry, could you say that again more slowly?}
  --- c'est un réflexe à automatiser, pas une honte.
  \item \textbf{Vous paniquez ?} Buvez une gorgée d'eau, respirez, dites
  \en{Let me think for a second.} Le silence de trois secondes paraît long pour vous,
  court pour l'autre.
  \item \textbf{La conversation dérape vers un sujet que vous ne maîtrisez pas ?}
  Ramenez-la : \en{That reminds me of something I know well\ldots}
  \item \textbf{Vous avez fait une faute grave ?} Corrigez en souriant :
  \en{Sorry --- I mean\ldots} et continuez. Personne ne vous en tiendra rigueur.
\end{enumerate}

\begin{mission}[Chapitre 15 --- le défi des 21 jours de parole]
Pendant 21 jours, faites chaque jour, sans exception :\\
\textbf{1.} 3 minutes de description d'image (drill 5).\\
\textbf{2.} 3 minutes de « 60 secondes » sur trois mots différents (drill 1).\\
\textbf{3.} 2 minutes de réponse à trois questions de la liste des 30 (drill 7).\\
\textbf{4.} Enregistrez-vous deux fois dans la semaine et notez trois progrès.\\
Total : 8 minutes par jour. C'est le plus petit investissement du livre pour le plus grand
bénéfice : après 21 jours, les mots arriveront \textbf{avant} que vous ayez fini de les
chercher.
\end{mission}


% ===== PARTIE VI — ÉCOUTER, LIRE, ÉCRIRE ===========================
\part{Écouter, lire, écrire}
% =====================================================================
%  CHAPITRE 16 — COMPRÉHENSION ORALE
% =====================================================================
\chapter{Comprendre l'anglais parlé : la méthode d'écoute}

\begin{quote}\itshape
« Je comprends quand je lis, mais dès qu'on me parle, tout disparaît. »
C'est la plainte la plus fréquente --- et elle a une cause précise, donc une solution précise.
\end{quote}

\section{Pourquoi vous ne comprenez pas (les cinq causes)}
\begin{enumerate}
  \item \textbf{Vous cherchez des mots isolés.} Le natif ne dit pas les mots séparément :
  \en{What do you want to do?} devient \pron{oua-da-ya-wou-na-dou}. Votre oreille cherche
  cinq mots français-style ; il n'y en a qu'un seul, fondu. (Voir le chapitre 6.)
  \item \textbf{Votre oreille n'est pas entraînée aux sons manquants.} \ipatheta, \ipadh, /h/,
  /w/, \ipang, \ipeschwa : tant que vous ne les distinguez pas, des dizaines de mots
  courants vous échappent.
  \item \textbf{Le niveau est trop élevé.} Écouter CNN au mois 2 ne fait pas progresser :
  cela démoralise. Il faut 90\,\% de compréhension pour apprendre.
  \item \textbf{Vous écoutez passivement.} Le fond sonore aide au rythme, mais ne crée pas
  de compétence. Seule l'écoute \textbf{active} (avec tâche) progresse.
  \item \textbf{Vous traduisez en français.} Traduire prend 2 secondes ; la phrase suivante
  arrive avant. Il faut comprendre \textbf{directement} --- c'est un entraînement, pas un don.
\end{enumerate}

\section{Les deux écoutes : intensive et extensive}
\begin{tableau}[Les deux régimes, tous deux indispensables]
\begin{tabularx}{\linewidth}{@{}p{3.2cm}p{5.0cm}X@{}}
\entete{} & \entete{Écoute intensive} & \entete{Écoute extensive} \tabularnewline
\hline
Durée & 15 à 25 min/jour & 30 à 60 min/jour (temps mort) \tabularnewline
Matériau & 1 à 3 minutes d'audio, niveau 90\,\% compris & Contenu long, facile, agréable \tabularnewline
Méthode & 5 passes, dictée, shadowing, analyse & On écoute en faisant autre chose, sans pause \tabularnewline
Objectif & Précision : sons, liaisons, mots nouveaux & Volume : rythme, familiarité, endurance \tabularnewline
Effet & Progresse vite, fatigue & Progresse lentement, ne fatigue pas \tabularnewline
\end{tabularx}
\end{tableau}

\subsection{Le protocole des 5 passes (écoute intensive, 20 minutes)}
\begin{enumerate}
  \item \textbf{Passe 1 --- écoute aveugle (2 min).} Écoutez l'audio sans texte, deux fois.
  Notez seulement : quel est le sujet ? qui parle ? quel est le ton ? Ne cherchez pas à tout
  comprendre.
  \item \textbf{Passe 2 --- écoute avec le texte (3 min).} Lisez la transcription en écoutant.
  Soulignez les mots inconnus et les endroits où le son ne correspondait pas à ce que vous
  attendiez.
  \item \textbf{Passe 3 --- écoute sans le texte (3 min).} Deux fois. Vous devez maintenant
  comprendre 95\,\%. Sinon, revenez à la passe 2.
  \item \textbf{Passe 4 --- dictée (7 min).} Écrivez une phrase sur cinq, phrase par phrase,
  avec pauses. Comparez avec le texte. \textbf{Analysez chaque erreur} : mot inconnu ?
  liaison non perçue ? son non reconnu ? C'est l'étape qui fait le plus progresser.
  \item \textbf{Passe 5 --- shadowing (5 min).} Répétez en même temps que l'audio, sans le
  texte, en imitant tout.
\end{enumerate}

\begin{astuce}[Le même audio trois jours de suite]
Jour 1 : protocole complet. Jour 2 : passe 3 + passe 5 seulement (7 min). Jour 3 : écoute
simple en marchant. Ce triple contact ancre l'audio dans votre mémoire --- vous retiendrez
80\,\% des expressions, contre 20\,\% avec un audio écouté une fois.
\end{astuce}

\section{La méthode « série » : 3 passes par épisode}
\begin{tableau}[Comment travailler une série sans se décourager]
\begin{tabularx}{\linewidth}{@{}p{2.4cm}p{3.6cm}X@{}}
\entete{Passe} & \entete{Sous-titres} & \entete{Objectif} \tabularnewline
\hline
1 & Français (mois 1-3) puis anglais (mois 4+) & Comprendre l'histoire, prendre du plaisir. Sans plaisir, on abandonne. \tabularnewline
2 & \textbf{Anglais} & Relier le son à l'écrit. Vous découvrirez que vous « saviez » la moitié des mots sans les reconnaître à l'oreille. \tabularnewline
3 & Aucun & Tester votre oreille réelle. 10 à 15 minutes suffisent --- le reste de l'épisode en passe 1 ou 2. \tabularnewline
\end{tabularx}
\end{tableau}

\begin{memo}[Le choix des séries]
Préférez les \textbf{sitcoms du quotidien} (dialogues courts, situations répétitives, humour
simple) aux séries policières ou fantastiques (vocabulaire technique, débit rapide, bruit de
fond). Regardez \textbf{deux fois le même épisode} : la seconde fois, tout devient audible,
et c'est cette sensation qui vous prouve vos progrès.
\end{memo}

\section{La méthode « chanson » (15 minutes, une fois par semaine)}
\begin{enumerate}
  \item Choisissez une chanson au débit clair (pop, folk, soul --- évitez le rap au début).
  \item Écoutez deux fois sans paroles, en notant les mots que vous reconnaissez.
  \item Lisez les paroles en écoutant : soulignez le vocabulaire nouveau et les liaisons.
  \item Cherchez le sens global (de quoi parle la chanson ? quelle émotion ?).
  \item Chantez --- oui, vraiment. Chanter installe le rythme, la longueur des voyelles et
  les enchaînements mieux que n'importe quel exercice.
  \item Ajoutez 5 expressions de la chanson à votre Anki.
\end{enumerate}

\section{Dix textes d'écoute à travailler}
Les textes ci-dessous sont conçus pour être \textbf{écoutés} : faites-les lire par une
application de synthèse vocale (votre téléphone, un dictionnaire en ligne, un traducteur
vocal), par un partenaire, ou lisez-les à voix haute en vous enregistrant, puis travaillez-les
sans le texte. Chaque texte est suivi de questions ; les réponses sont en fin de section.

\subsection{Texte 1 --- niveau A1 : se présenter}
\begin{dialoguebox}[Texte 1 --- à écouter 5 fois]
\begin{anglais}
Hello! My name is Awa Diop. I'm twenty-four years old and I come from Senegal. I live in
Thiès with my parents and my two brothers. I work in a small shop near the market. I open
the shop at eight in the morning and I close it at six in the evening. In my free time, I
like listening to music and playing basketball with my friends. I'm learning English because
I want to work with customers from other countries. My favourite subject is English, but I
find pronunciation difficult. On Sundays, I usually visit my grandmother and we eat rice
with fish together.
\end{anglais}
\end{dialoguebox}
\textbf{Questions :} 1. Quel âge a Awa ? 2. Où habite-t-elle et avec qui ?
3. Quelles sont ses heures de travail ? 4. Pourquoi apprend-elle l'anglais ?
5. Que fait-elle le dimanche ?

\subsection{Texte 2 --- niveau A2 : une journée de travail}
\begin{dialoguebox}[Texte 2 --- à écouter 5 fois]
\begin{anglais}
I usually get up at six, but yesterday was different. I woke up at nine because I had worked
very late the night before. When I looked at my phone, I saw five missed calls from my boss.
I called him back immediately and he wasn't angry --- he just needed the report for the
morning meeting. I took a taxi because I didn't have time to wait for the bus, and the
traffic was terrible, so I arrived at eleven. In the end, I finished the report at two
o'clock and my boss thanked me. That evening, I went to bed at nine. I've never been so
tired in my life.
\end{anglais}
\end{dialoguebox}
\textbf{Questions :} 1. Pourquoi s'est-il réveillé à neuf heures ? 2. Qu'a-t-il vu sur son
téléphone ? 3. Son patron était-il fâché ? 4. Pourquoi a-t-il pris un taxi ?
5. À quelle heure a-t-il fini le rapport ? 6. Repérez les quatre temps utilisés.

\subsection{Texte 3 --- niveau A2 : au restaurant}
\begin{dialoguebox}[Texte 3 --- à écouter 5 fois]
\begin{anglais}
Good evening! Do you have a reservation? No? That's not a problem --- we have a table by the
window. Here are the menus. Today's special is grilled fish with vegetables, and we also
have chicken with rice. Would you like something to drink while you decide? We have fresh
orange juice, mineral water, and mint tea. I'll give you a few minutes. By the way, the fish
is very fresh --- it came from the market this morning. Let me know when you're ready to
order, and if you have any allergies, please tell me now.
\end{anglais}
\end{dialoguebox}
\textbf{Questions :} 1. Le client avait-il réservé ? 2. Où est placée la table ?
3. Quels sont les deux plats proposés ? 4. Quelles boissons sont citées ?
5. Pourquoi le poisson est-il « very fresh » ?

\subsection{Texte 4 --- niveau B1 : un appel téléphonique professionnel}
\begin{dialoguebox}[Texte 4 --- à écouter 5 fois]
\begin{anglais}
Good morning, West Africa Logistics, Fatou speaking. How can I help you? \ldots\ I see. So
you're calling about the container that was supposed to arrive on Tuesday? Let me check that
for you. Could you give me the reference number, please? Thank you. Ah, yes --- I can see
the problem. Unfortunately, the shipment has been delayed because of bad weather at the port,
and it's now expected on Friday. I do apologise for the inconvenience. What I can do is send
you a confirmation by email within the hour, and I'll also ask my colleague to call you if
there's any further change. Would that be all right? \ldots\ Perfect. Is there anything else
I can help you with today? Thank you for calling, Mr Ba. Have a good day.
\end{anglais}
\end{dialoguebox}
\textbf{Questions :} 1. Pourquoi le client appelle-t-il ? 2. Quelle est la cause du retard ?
3. Quand le conteneur doit-il arriver ? 4. Quelles sont les deux promesses de Fatou ?
5. Relevez trois formules de politesse professionnelle.

\subsection{Texte 5 --- niveau B1 : donner son opinion}
\begin{dialoguebox}[Texte 5 --- à écouter 5 fois]
\begin{anglais}
In my opinion, social media has changed the way young people learn English, and most of the
change is positive. Ten years ago, if you wanted to hear real English, you had to find a
teacher or a cassette. Today, you can watch a film, follow a footballer, or listen to a
podcast on your phone --- for free, and as often as you like. On the other hand, I think
there's a danger. Many young people believe that watching videos is enough, and they never
open their mouths. Understanding is not speaking. So my advice is simple: use social media
every day, but spend at least twenty minutes a day producing English --- out loud, alone if
necessary. That's how you turn input into ability.
\end{anglais}
\end{dialoguebox}
\textbf{Questions :} 1. Quelle est l'opinion générale du locuteur ? 2. Quelle comparaison
fait-il avec il y a dix ans ? 3. Quel danger identifie-t-il ? 4. Que signifie pour lui
« Understanding is not speaking » ? 5. Quel conseil donne-t-il précisément ?

\subsection{Texte 6 --- niveau B1 : raconter un voyage}
\begin{dialoguebox}[Texte 6 --- à écouter 5 fois]
\begin{anglais}
Last year I travelled to Ghana for the first time, and honestly, I wasn't expecting to enjoy
it so much. I'd planned everything: the flight, the hotel, the places I wanted to visit.
But of course, nothing went as planned. On the second day, while I was waiting for a bus to
Cape Coast, I started talking to a man called Kwame. He told me he was going to the same
place, and that his cousin had a car. Two hours later, I was sitting in the front seat,
eating bread with eggs, and listening to stories about his family. We spent three days
together. I've kept in touch with him ever since --- we speak on the phone almost every
month. If I hadn't missed that bus, I'd never have met him. Travelling isn't about the
places; it's about the people you meet in them.
\end{anglais}
\end{dialoguebox}
\textbf{Questions :} 1. Où le narrateur est-il allé ? 2. Qu'est-ce qui ne s'est pas passé
comme prévu ? 3. Comment a-t-il rencontré Kwame ? 4. Sont-ils encore en contact ?
5. Repérez le conditionnel 3 et le \en{present perfect} dans le texte.
6. Quelle est la conclusion du narrateur ?

\subsection{Texte 7 --- niveau B2 : les actualités}
\begin{dialoguebox}[Texte 7 --- à écouter 5 fois]
\begin{anglais}
Good evening. The government announced this afternoon that a new programme will provide
free digital skills training to fifty thousand young people over the next three years.
Speaking at a press conference in the capital, the Minister of Youth said that the scheme
had already been tested successfully in two regions, where more than seventy per cent of
participants found work within six months. Critics, however, have questioned whether the
funding will be sufficient. One economist told our correspondent that similar initiatives
in the past had failed to reach rural areas. The programme is expected to begin in January,
and registration will open online next week. In other news, the national team has qualified
for the Africa Cup of Nations after a two-nil victory last night.
\end{anglais}
\end{dialoguebox}
\textbf{Questions :} 1. Combien de jeunes seront formés, et en combien de temps ?
2. Quel résultat a été observé dans les deux régions test ? 3. Quelle critique est formulée ?
4. Quand le programme commencera-t-il ? 5. Quelle est la seconde nouvelle ?
6. Relevez deux exemples de voix passive.

\subsection{Texte 8 --- niveau B2 : un entretien}
\begin{dialoguebox}[Texte 8 --- à écouter 5 fois]
\begin{anglais}
So, tell me --- what would you say is your greatest professional achievement? \ldots\ Hmm.
I think it would have to be the reorganisation of our accounting system two years ago.
Before that, we were closing the monthly accounts by hand, and it took almost two weeks.
I suggested that we move to a new software, but the idea wasn't popular at first --- people
were used to the old way, and nobody wanted extra work. So I did two things: I trained the
team myself, in the evenings, and I showed the management how much time we would save.
Six months later, we were closing the accounts in three days instead of fourteen. To be
honest, I'm prouder of the way I convinced people than of the result itself. Getting people
to change is harder than changing a system.
\end{anglais}
\end{dialoguebox}
\textbf{Questions :} 1. Quelle question est posée au début ? 2. Quel problème existait avant ?
3. Pourquoi l'idée n'était-elle pas populaire ? 4. Quelles sont les deux choses faites par
le candidat ? 5. Quel est le résultat chiffré ? 6. De quoi est-il le plus fier, et pourquoi ?

\subsection{Texte 9 --- niveau A2/B1 : conversation familière}
\begin{dialoguebox}[Texte 9 --- à écouter 5 fois (débit rapide, réductions)]
\begin{anglais}
Hey, what's up? Long time no see! \ldots\ Yeah, I know, I've been really busy lately. Work's
been crazy --- we've got a new manager and he wants everything done yesterday. Anyway, how
about you? \ldots\ Oh, that's great! Congrats! When do you start? \ldots\ Nice. Hey, we
should catch up properly sometime --- maybe grab a coffee this weekend? I'm free Saturday
afternoon. \ldots\ Cool, let's say three o'clock at the place near the beach? Perfect.
I gotta run now, but I'll text you the details. Take care!
\end{anglais}
\end{dialoguebox}
\textbf{Questions :} 1. Quel est le registre (formel ou familier) ? 2. Que veut dire
\en{everything done yesterday} ? 3. Quelle bonne nouvelle l'autre personne annonce-t-elle ?
4. Que proposent-ils de faire ? 5. Où et quand se retrouveront-ils ?
6. Relevez quatre réductions familières.

\subsection{Texte 10 --- niveau B2/C1 : conférence courte}
\begin{dialoguebox}[Texte 10 --- à écouter 5 fois]
\begin{anglais}
Here's a question I'd like you to sit with for a moment: why do most adults find language
learning so much harder than children do? The usual answer is that children's brains are
more flexible. That's partly true, but I'd like to suggest something else. Children are
allowed to be terrible. A three-year-old says ``I goed to school'' and everybody smiles and
repeats the correct form. Nobody laughs at them; nobody thinks they're stupid. Adults, on
the other hand, have a reputation to protect. We're afraid of sounding foolish, so we stay
silent --- and silence is the one thing that guarantees you'll never improve. So my advice
isn't about grammar or vocabulary. It's much simpler: give yourself permission to be bad at
something, publicly, every single day, for a year. That's the whole secret.
\end{anglais}
\end{dialoguebox}
\textbf{Questions :} 1. Quelle est la question posée ? 2. Quelle est la réponse habituelle,
et pourquoi est-elle insuffisante ? 3. Quel avantage les enfants ont-ils selon la conférencière ?
4. Que craignent les adultes ? 5. Quel est le conseil final ? 6. Expliquez
\en{sit with a moment} et \en{give yourself permission}.

\subsection{Corrigés des dix textes}
\begin{corrige}[Textes 1 à 10 --- réponses essentielles]
\textbf{1.} 24 ans ; Thiès, avec ses parents et ses deux frères ; 8 h - 18 h ;
pour travailler avec des clients étrangers ; elle visite sa grand-mère et ils mangent du
riz au poisson.\\
\textbf{2.} Il avait travaillé très tard la veille ; cinq appels manqués de son patron ;
non, il avait seulement besoin du rapport ; pas le temps d'attendre le bus (et circulation
terrible) ; à 14 h ; passés utilisés : \en{present simple} (\en{I usually get up}),
\en{past simple} (\en{I woke up}), \en{past perfect} (\en{I had worked}),
\en{present perfect} (\en{I've never been}).\\
\textbf{3.} Non ; près de la fenêtre ; poisson grillé avec légumes, poulet avec riz ;
jus d'orange frais, eau minérale, thé à la menthe ; il vient du marché ce matin.\\
\textbf{4.} Un conteneur qui devait arriver mardi ; le mauvais temps au port ; vendredi ;
confirmation par e-mail dans l'heure + appel du collègue en cas de changement ;
\en{I do apologise for the inconvenience}, \en{Would that be all right?},
\en{Is there anything else I can help you with today?}\\
\textbf{5.} Plutôt positive ; avant, il fallait un professeur ou une cassette ;
croire que regarder suffit et ne jamais parler ; comprendre ne remplace pas produire ;
utiliser les réseaux sociaux chaque jour \textbf{et} produire 20 minutes à voix haute.\\
\textbf{6.} Au Ghana ; rien ne s'est passé comme prévu / il a raté le bus ; en attendant le
bus pour Cape Coast ; oui, ils se parlent presque chaque mois ;
\en{If I hadn't missed that bus, I'd never have met him} (conditionnel 3) et
\en{I've kept in touch with him ever since} (\en{present perfect}) ;
le voyage compte par les personnes rencontrées.\\
\textbf{7.} 50\,000 jeunes en trois ans ; plus de 70\,\% ont trouvé un emploi en six mois ;
le financement pourrait être insuffisant et les zones rurales oubliées ; en janvier,
inscriptions en ligne la semaine prochaine ; l'équipe nationale s'est qualifiée
(victoire 2-0) ; passifs : \en{the scheme had already been tested}, \en{the programme is expected to begin}.\\
\textbf{8.} « Quelle est votre plus grande réalisation professionnelle ? » ; la clôture des
comptes prenait deux semaines, à la main ; les gens étaient habitués et ne voulaient pas de
travail supplémentaire ; il a formé l'équipe lui-même le soir + a montré les gains de temps
à la direction ; trois jours au lieu de quatorze ; il est plus fier d'avoir convaincu les
gens, car changer les habitudes est plus dur que changer un système.\\
\textbf{9.} Familier ; tout doit être fait immédiatement (hier = déjà trop tard) ;
l'autre a un nouveau travail (ou une bonne nouvelle personnelle) ; prendre un café samedi ;
à 15 h, près de la plage ; réductions : \en{what's up}, \en{I've been} \pron{aïv-bin}, \en{we've got}, \en{gotta run}, \en{Let's say}, \en{I'll text you}.\\
\textbf{10.} Pourquoi les adultes apprennent-ils moins facilement que les enfants ?
La flexibilité cérébrale --- vraie mais incomplète ; les enfants ont le droit d'être
mauvais, personne ne se moque ; les adultes craignent d'avoir l'air idiots et se taisent ;
se donner la permission d'être mauvais, en public, chaque jour, pendant un an ;
\en{sit with} = garder une question en tête, y réfléchir ; \en{give yourself permission}
= s'autoriser.
\end{corrige}

\section{Le plan d'écoute sur 12 mois}
\begin{tableau}[Ce que vous écoutez, mois par mois]
\begin{tabularx}{\linewidth}{@{}p{1.8cm}p{4.6cm}p{4.2cm}X@{}}
\entete{Mois} & \entete{Intensif (20 min/jour)} & \entete{Extensif (temps mort)} & \entete{Objectif} \tabularnewline
\hline
1-2 & Dialogues de ce livre lus à voix haute + audios pour débutants (vitesse 0,75×) & Chansons lentes, radio en fond & Reconnaître les 44 sons et le rythme \tabularnewline
3-4 & Textes 1 à 5 ci-dessus, audios A2, dessins animés & Podcasts pour apprenants, séries avec sous-titres français & Comprendre un dialogue simple sans texte \tabularnewline
5-6 & Textes 6 à 10, podcasts A2/B1, interviews courtes & Sitcoms avec sous-titres anglais & Comprendre 70\,\% d'une sitcom \tabularnewline
7-8 & Actualités lentes, conférences courtes, TED-Ed & Séries en sous-titres anglais, documentaires & Suivre une conférence simple \tabularnewline
9-10 & Journaux télévisés, podcasts natifs, débats & Films, humour, interviews longues & Comprendre les informations à 80\,\% \tabularnewline
11-12 & Contenu natif non adapté, accents variés & Films sans sous-titres, radio, humour & Comprendre un film sans sous-titres \tabularnewline
\end{tabularx}
\end{tableau}

\begin{mission}[Chapitre 16 --- la semaine d'écoute type]
\textbf{Lundi :} texte 1, protocole des 5 passes. \textbf{Mardi :} texte 1, passes 3 et 5
seulement + texte 2, passes 1-2. \textbf{Mercredi :} texte 2 complet.
\textbf{Jeudi :} dictée de 8 phrases d'un texte déjà travaillé. \textbf{Vendredi :} chanson
(méthode chanson complète). \textbf{Samedi :} un épisode de série en 3 passes.
\textbf{Dimanche :} réécoute passive de tout le matériel de la semaine en marchant.
\end{mission}

% =====================================================================
%  CHAPITRE 17 — LECTURE ET ÉCRITURE
% =====================================================================
\chapter{Lire et écrire : orthographe, ponctuation, modèles}

\section{Partie A --- La lecture}

\subsection{Pourquoi lire est l'accélérateur le moins fatigant}
Lire apporte le vocabulaire en contexte, la grammaire par imprégnation, et la structure des
phrases anglaises naturelles. Vingt minutes de lecture agréable par jour valent une heure
d'exercices pénibles. Mais il faut lire \textbf{au bon niveau} --- c'est là que tout se joue.

\begin{regle}[La règle des 98\,\%]
Prenez une page au hasard. Comptez les mots que vous ne connaissez pas :\\
\textbf{0 à 2 mots inconnus sur 100} \ra niveau idéal, lisez pour le plaisir.\\
\textbf{3 à 5 mots} \ra niveau utile mais fatigant : travaillez 10 minutes, pas plus.\\
\textbf{Plus de 5 mots} \ra trop difficile : changez de livre, ce n'est pas une question
de volonté.
\end{regle}

\begin{tableau}[Les deux lectures, à alterner]
\begin{tabularx}{\linewidth}{@{}p{3.2cm}p{5.2cm}X@{}}
\entete{} & \entete{Lecture extensive} & \entete{Lecture intensive} \tabularnewline
\hline
Volume & 20 à 40 pages faciles par jour & 1 à 2 pages difficiles par jour \tabularnewline
But & Plaisir, fluidité, volume de mots rencontrés & Précision, vocabulaire actif, structures \tabularnewline
Dictionnaire & \textbf{Interdit} (ou une fois par page) & Autorisé, avec notes \tabularnewline
À voix haute & Non & Oui, en imitant le rythme \tabularnewline
Ce qu'on retient & Le sens général & Chaque mot, chaque tournure \tabularnewline
\end{tabularx}
\end{tableau}

\subsection{Que lire, mois par mois}
\begin{tableau}[Parcours de lecture sur 12 mois]
\begin{tabularx}{\linewidth}{@{}p{1.8cm}p{4.6cm}X@{}}
\entete{Mois} & \entete{Matériau} & \entete{Conseil} \tabularnewline
\hline
1-2 & Textes de ce livre ; dialogues ; listes de phrases ; panneaux, étiquettes, notices d'appareils & Relisez ce que vous connaissez déjà : la relecture crée la vitesse. \tabularnewline
3-4 & Livres simplifiés (\en{graded readers}) niveau A2 ; livres pour enfants illustrés ; articles courts en ligne & 10 pages par jour, sans dictionnaire. Un livre entier en une semaine. \tabularnewline
5-6 & \en{Graded readers} niveau B1 ; nouvelles courtes ; articles de presse simples ; bandes dessinées & Notez 10 expressions par semaine, pas 50 mots. \tabularnewline
7-8 & Premiers romans faciles (jeunesse, policier simple) ; articles d'actualité ; blogs sur vos centres d'intérêt & Changez de genre chaque semaine. \tabularnewline
9-10 & Romans pour adultes, presse (actualités, sport, économie), essais courts & Lisez le même sujet dans deux sources différentes. \tabularnewline
11-12 & Littérature, presse d'opinion, documentation professionnelle de votre métier & Objectif : lire comme vous lisez en français --- vite, et pour le contenu. \tabularnewline
\end{tabularx}
\end{tableau}

\begin{astuce}[Choisissez des livres que vous avez déjà lus en français]
Relire en anglais un roman ou une bande dessinée que vous connaissez par c\oe ur est la
méthode la plus efficace qui existe : vous comprenez l'histoire, donc votre cerveau se
concentre sur la \textbf{forme}. Même chose avec les films : revoir un film déjà vu avec
sous-titres anglais fait des miracles.
\end{astuce}

\subsection{Trois textes de lecture commentés}

\subsubsection*{Texte A --- niveau A2 : une lettre d'un ami}
\begin{dialoguebox}[Texte A]
\begin{anglais}
Dear Sam,

Thanks for your message --- it was great to hear from you. Sorry I haven't written for so
long, but work has been really busy these last few weeks.

Let me tell you my news. In March, I moved to a new flat in the centre of town. It's small,
but it has a balcony, and I can walk to work in fifteen minutes. I've been living there for
two months now and I'm very happy with it.

I've also started playing football again, twice a week with some colleagues. I'm not as fit
as I used to be, so the first games were hard! Last Saturday we won three-two, and I scored
--- the first goal of my life, believe me.

What about you? Are you still working at the hospital? How are your parents? I'm planning to
visit Dakar in August, so maybe we could meet up and have dinner.

Write soon and send me some photos.

Best wishes,\\
Moussa
\end{anglais}
\end{dialoguebox}

\textbf{Questions :} 1. Pourquoi Moussa n'a-t-il pas écrit plus tôt ? 2. Depuis combien de
temps habite-t-il son nouvel appartement ? 3. Quel avantage principal présente cet
appartement ? 4. Pourquoi les premiers matchs ont-ils été difficiles ?
5. Que propose-t-il pour août ? 6. Relevez trois temps différents et expliquez leur usage.

\subsubsection*{Texte B --- niveau B1 : un article}
\begin{dialoguebox}[Texte B]
\begin{anglais}
\textbf{Why walking is the most underrated exercise}

Most people think that to get fit, they need a gym, expensive equipment and several hours a
week. In fact, one of the best things you can do for your health costs nothing and requires
no training: walking.

Studies have shown that thirty minutes of brisk walking a day reduces the risk of heart
disease, improves sleep and lowers stress. Unlike running, walking can be continued into old
age, and it rarely causes injuries. It also has an advantage that is often forgotten: you
can do it with other people, which makes it easier to keep going.

So how do you start? Forget about the idea of a daily hour-long walk --- it will last two
weeks. Instead, add walking to what you already do. Get off the bus one stop earlier. Walk
while you are on the phone. Take the stairs instead of the lift. After a month, you will
have added forty minutes a day to your life without ever finding the time.

The hardest part of exercise is not the exercise. It is starting, and continuing.
\end{anglais}
\end{dialoguebox}

\textbf{Questions :} 1. Quelle idée reçue l'article conteste-t-il ? 2. Citez trois bénéfices
de la marche. 3. Pourquoi la marche est-elle préférable à la course selon l'auteur ?
4. Quel conseil concret donne-t-il pour commencer ? 5. Que signifie la dernière phrase ?
6. Repérez deux exemples de voix passive et un comparatif.

\subsubsection*{Texte C --- niveau B2 : un extrait argumentatif}
\begin{dialoguebox}[Texte C]
\begin{anglais}
\textbf{Should English replace French as our second language?}

It is a question that comes up every few years, usually after a diplomatic meeting or a
football match, and it deserves a calmer answer than the one it usually receives.

The argument for English is, on the surface, irresistible. It is the language of science, of
aviation, of computing, of most international business and of a very large share of the
entertainment the world consumes. A young person who speaks English well has access to
information that a young person who does not simply cannot reach.

But the argument is often presented as though languages were tools in a box, to be kept or
discarded according to usefulness. They are not. A language is also a way of thinking, a
literature, a set of relationships, and in our region, a history. French connects us to more
than twenty countries and to a diplomatic tradition in which we have invested for decades.
Dropping it would not make us more competitive; it would make us poorer.

The sensible answer is therefore not either/or, but both --- and, increasingly, more. The
real problem is not that we speak too many languages, but that we speak them badly. Improving
English teaching should not mean abandoning French; it should mean being honest about how
little of either we currently master.
\end{anglais}
\end{dialoguebox}

\textbf{Questions :} 1. Quand cette question resurgit-elle habituellement ? 2. Quels sont les
domaines où l'anglais domine ? 3. Quelle métaphore l'auteur rejette-t-il, et pourquoi ?
4. Que perdrait-on en abandonnant le français ? 5. Quelle est la conclusion ?
6. Que critique l'auteur dans la dernière phrase ? 7. Relevez trois connecteurs logiques et
expliquez leur fonction.

\begin{corrige}[Textes A, B, C]
\textbf{A.} 1. Le travail a été très chargé. 2. Deux mois. 3. Il est à quinze minutes à pied
du travail (et il a un balcon). 4. Il n'est plus aussi en forme qu'avant (\en{not as fit as
I used to be}). 5. Se rencontrer et dîner ensemble à Dakar. 6. \en{I haven't written}
(\en{present perfect} : lien avec le présent), \en{I moved} (passé simple : date précise),
\en{I've been living there for two months} (\en{present perfect continuous} : durée en cours),
\en{I used to be} (habitude passée révolue).\\[4pt]
\textbf{B.} 1. L'idée qu'il faut une salle de sport, du matériel coûteux et beaucoup de temps.
2. Réduit le risque de maladie cardiaque, améliore le sommeil, diminue le stress.
3. On peut la continuer en vieillissant et elle cause rarement des blessures ; on peut aussi
la pratiquer à plusieurs. 4. Ajouter la marche à ce qu'on fait déjà : descendre un arrêt plus
tôt, marcher au téléphone, prendre les escaliers. 5. Le plus dur n'est pas l'effort, c'est de
commencer et de continuer. 6. Passifs : \en{that is often forgotten}, \en{it will be
continued} (ou \en{can be continued}) ; comparatif : \en{Unlike running\ldots} /
\en{easier to keep going}.\\[4pt]
\textbf{C.} 1. Après une réunion diplomatique ou un match de football. 2. Science, aviation,
informatique, affaires internationales, divertissement. 3. La métaphore des langues comme
outils dans une boîte : une langue est aussi une pensée, une littérature, des relations, une
histoire. 4. Des relations avec plus de vingt pays et une tradition diplomatique ; on
deviendrait plus pauvre, pas plus compétitif. 5. Non pas l'un ou l'autre, mais les deux ---
et même davantage. 6. Nous parlons mal les deux langues, plutôt que trop de langues.
7. \en{therefore} (conséquence), \en{but} (opposition), \en{so} (conséquence),
\en{unlike} (contraste), \en{increasingly} (progression).
\end{corrige}


\section{Partie B --- L'orthographe}

\subsection{Les sept règles qui résolvent 90\,\% des problèmes}
\begin{tableau}[Règles d'orthographe]
\begin{tabularx}{\linewidth}{@{}p{4.2cm}X@{}}
\entete{Règle} & \entete{Application} \tabularnewline
\hline
\textbf{1. Consonne doublée} & Verbe d'une syllabe, voyelle + consonne : on double avant
\en{-ing} et \en{-ed} : \en{stop} \ra \en{stopped, stopping} ; \en{plan} \ra \en{planned,
planning} ; \en{travel} \ra \en{travelled} (BrE) / \en{traveled} (AmE). \tabularnewline
\textbf{2. \en{e} final supprimé} & Devant une voyelle : \en{make} \ra \en{making} ;
\en{write} \ra \en{writing} ; mais conservé devant consonne : \en{making} / \en{likes}. \tabularnewline
\textbf{3. \en{y} \ra \en{i}} & Après consonne : \en{study} \ra \en{studies, studied} ;
\en{happy} \ra \en{happier, happiness}. Après voyelle, on garde : \en{play} \ra \en{plays,
playing}. \tabularnewline
\textbf{4. \en{i} avant \en{e}} & Sauf après \en{c} : \en{believe, receive, piece, ceiling}.
Exceptions : \en{weird, height, either, their}. \tabularnewline
\textbf{5. Pluriels en \en{-es}} & Après \en{-s, -sh, -ch, -x, -o} : \en{boxes, brushes,
watches, buses, tomatoes}. \tabularnewline
\textbf{6. \en{-f} \ra \en{-ves}} & \en{leaf} \ra \en{leaves}, \en{life} \ra \en{lives},
\en{wife} \ra \en{wives}, \en{knife} \ra \en{knives} ; mais \en{roofs, chiefs, beliefs}. \tabularnewline
\textbf{7. Adverbes en \en{-ly}} & \en{happy} \ra \en{happily} ; \en{true} \ra \en{truly} ;
\en{terrible} \ra \en{terribly} ; mais \en{public} \ra \en{publicly}. \tabularnewline
\end{tabularx}
\end{tableau}

\subsection{Orthographe britannique ou américaine ?}
\begin{memo}[Choisissez-en une et gardez-la]
Mélanger les deux dans un même document est la seule vraie faute. En Afrique de l'Ouest,
l'enseignement suit généralement le britannique ; les médias et Internet sont majoritairement
américains. Recommandation : \textbf{britannique à l'écrit formel} (école, administration,
examens), \textbf{américain accepté} dans les échanges informels et sur Internet.
\end{memo}

\begin{tableau}[Les 40 différences les plus courantes]
\footnotesize
\begin{longtable}{@{}p{3.4cm}p{3.4cm}p{3.4cm}p{3.4cm}@{}}
\entete{Britannique} & \entete{Américain} & \entete{Britannique} & \entete{Américain} \tabularnewline
\endfirsthead
\entete{BrE} & \entete{AmE} & \entete{BrE} & \entete{AmE} \tabularnewline
\endhead
\hline
\en{colour} & \en{color} & \en{organise} & \en{organize} \tabularnewline
\en{favourite} & \en{favorite} & \en{realise} & \en{realize} \tabularnewline
\en{honour} & \en{honor} & \en{apologise} & \en{apologize} \tabularnewline
\en{behaviour} & \en{behavior} & \en{analyse} & \en{analyze} \tabularnewline
\en{neighbour} & \en{neighbor} & \en{practise} (v.) & \en{practice} (v.) \tabularnewline
\en{centre} & \en{center} & \en{licence} (n.) & \en{license} (n.) \tabularnewline
\en{theatre} & \en{theater} & \en{defence} & \en{defense} \tabularnewline
\en{metre} & \en{meter} & \en{offence} & \en{offense} \tabularnewline
\en{litre} & \en{liter} & \en{programme} & \en{program} \tabularnewline
\en{travelled} & \en{traveled} & \en{catalogue} & \en{catalog} \tabularnewline
\en{cancelled} & \en{canceled} & \en{dialogue} & \en{dialog} \tabularnewline
\en{labelled} & \en{labeled} & \en{cheque} & \en{check} (bancaire) \tabularnewline
\en{fulfil} & \en{fulfill} & \en{grey} & \en{gray} \tabularnewline
\en{enrol} & \en{enroll} & \en{tyre} & \en{tire} (pneu) \tabularnewline
\en{jewellery} & \en{jewelry} & \en{aeroplane} & \en{airplane} \tabularnewline
\en{kerb} & \en{curb} & \en{lorry} & \en{truck} \tabularnewline
\en{pavement} & \en{sidewalk} & \en{petrol} & \en{gas / gasoline} \tabularnewline
\en{lift} & \en{elevator} & \en{flat} & \en{apartment} \tabularnewline
\en{biscuit} & \en{cookie} & \en{sweets} & \en{candy} \tabularnewline
\en{holiday} & \en{vacation} & \en{autumn} & \en{fall} \tabularnewline
\end{longtable}
\end{tableau}

\section{Partie C --- La ponctuation et la typographie anglaises}
\begin{tableau}[Douze règles de ponctuation --- différentes du français]
\begin{tabularx}{\linewidth}{@{}p{4.6cm}X@{}}
\entete{Règle} & \entete{Détail} \tabularnewline
\hline
\textbf{Pas d'espace avant} \en{, ; : ! ?} & Contrairement au français : \en{Hello! How are you?}
(aucune espace avant \en{!} ni \en{?}). \tabularnewline
\textbf{Guillemets} & \en{``double quotes''} en anglais américain et souvent britannique ;
\en{`single quotes'} pour une citation dans la citation. Pas de « ». \tabularnewline
\textbf{Apostrophe} & \en{'} pour la possession (\en{Ali's book}) et les contractions
(\en{don't, I'm}). \textbf{Jamais} de pluriel avec apostrophe : \en{the 1990s} (pas \en{1990's}). \tabularnewline
\textbf{Majuscules} & Jours, mois, langues, nationalités prennent une majuscule :
\en{Monday, January, English, Senegalese}. Pas les adjectifs de matière ni les titres
complets (\en{the president} sauf devant un nom : \en{President Diop}). \tabularnewline
\textbf{Virgule} & Avant \en{and} dans une énumération (virgule d'Oxford) :
\en{bread, milk, and eggs}. Après un connecteur en début de phrase : \en{However, I disagree.} \tabularnewline
\textbf{Point-virgule} & Plus rare qu'en français ; préférer deux phrases. \tabularnewline
\textbf{Deux-points} & Pour introduire une liste ou une explication : \en{I need three things:
time, patience and coffee.} \tabularnewline
\textbf{Tiret long} & \en{---} (em dash) pour une incise, sans espaces en américain :
\en{He arrived --- late, as usual --- and apologised.} \tabularnewline
\textbf{Nombres} & Virgule pour les milliers, point pour les décimales :
\en{12,500 francs} ; \en{3.5 kilometres} --- \textbf{l'inverse du français}. \tabularnewline
\textbf{Heures} & \en{6.30 pm} (BrE) ou \en{6:30 pm} (AmE) ; \en{a.m. / p.m.} \tabularnewline
\textbf{Dates} & \en{5 May 2026} (BrE) ou \en{May 5, 2026} (AmE). \tabularnewline
\textbf{Titres} & Majuscule aux mots importants : \en{How to Learn English in Twelve Months}. \tabularnewline
\end{tabularx}
\end{tableau}

\begin{piege}[Les cinq erreurs de ponctuation les plus visibles]
\en{*Hello !} \ra \en{Hello!} \qquad
\en{*« I'm fine », he said.} \ra \en{``I'm fine,'' he said.} \qquad
\en{*12.500 francs} \ra \en{12,500 francs} \qquad
\en{*the 1990's} \ra \en{the 1990s} \qquad
\en{*in monday} \ra \en{on Monday}
\end{piege}

\section{Partie D --- Écrire : huit modèles prêts à l'emploi}

\subsection{Modèle 1 --- E-mail informel}
\begin{dialoguebox}[E-mail à un ami ou un partenaire de langue]
\begin{anglais}
\textbf{Subject: Let's practise this week?}

Hi John,

Hope you're doing well! Thanks for your last message --- it really helped.

I've been working on pronunciation this month, and I think I've improved a bit, but I still
struggle with the ``th'' sound. Would you like to do a voice call this weekend? I'm free on
Saturday afternoon or Sunday morning, your time.

In exchange, I can help you with your French --- I've prepared a short list of expressions
that people actually use.

Let me know what works for you.

Take care,\\
Moussa
\end{anglais}
\end{dialoguebox}
\textbf{Structure :} salutation courte \ra une phrase de lien \ra l'objet de l'e-mail
\ra la proposition précise \ra l'offre en retour \ra formule de fin chaleureuse.

\subsection{Modèle 2 --- E-mail professionnel (demande)}
\begin{dialoguebox}[Demande d'information à une entreprise]
\begin{anglais}
\textbf{Subject: Enquiry about the training programme --- ref. 2026/EN}

Dear Mr Mensah,

I am writing to enquire about the digital marketing training programme advertised on your
website last week.

I would be grateful if you could provide me with the following information:
\begin{itemize}
  \item the dates and duration of the next session;
  \item the total cost, and whether payment can be made in instalments;
  \item whether the course is available online for participants based in Thiès.
\end{itemize}

I have attached my CV for your reference. I look forward to hearing from you.

Yours sincerely,\\
Moussa Diop\\
\textit{Accountant --- Tel: +221 77 000 00 00}
\end{anglais}
\end{dialoguebox}
\textbf{Les blocs fixes à retenir :} \en{I am writing to\ldots} /
\en{I would be grateful if you could\ldots} / \en{Please find attached\ldots} /
\en{I look forward to hearing from you} (+ \en{-ing} !) /
\en{Yours sincerely} (si vous connaissez le nom) --- \en{Yours faithfully} (si
\en{Dear Sir or Madam}).

\subsection{Modèle 3 --- Réclamation polie mais ferme}
\begin{dialoguebox}[Produit défectueux]
\begin{anglais}
\textbf{Subject: Faulty product --- order \#45892}

Dear Sir or Madam,

On 12 May I ordered a wireless printer from your online shop (order number 45892). It was
delivered on 18 May, but unfortunately it does not work properly: the printer turns off after
a few minutes and cannot connect to my computer.

I have already contacted your support team twice by phone, but the problem has not been
resolved. I am therefore writing to request either a replacement or a full refund.

I would appreciate a reply within seven days. I have attached a copy of the invoice and a
short video showing the problem.

Yours faithfully,\\
Fatou Ndiaye
\end{anglais}
\end{dialoguebox}
\textbf{La technique du sandwich appliquée à l'écrit :} faits datés \ra problème précis
\ra démarches déjà faites \ra demande claire (\en{either\ldots or\ldots}) \ra délai
\ra pièces jointes \ra formule polie. Aucune émotion, aucun reproche personnel : c'est ce
qui rend une réclamation efficace.

\subsection{Modèle 4 --- Lettre de motivation (structure en 4 paragraphes)}
\begin{dialoguebox}[Squelette à personnaliser]
\begin{anglais}
\textbf{Paragraph 1 --- Why I am writing.}\\
\en{I am writing to apply for the position of \ldots, which I saw advertised on \ldots\ I am
very interested in this role because \ldots}

\textbf{Paragraph 2 --- Why I am suitable (2 proofs + figures).}\\
\en{I have \ldots years' experience in \ldots\ In my current position at \ldots, I am
responsible for \ldots\ Last year, I managed to \ldots, which resulted in a \ldots per cent
increase in \ldots}

\textbf{Paragraph 3 --- Why this company.}\\
\en{I am particularly attracted to your company because of \ldots\ I believe my experience in
\ldots would be a valuable addition to your team.}

\textbf{Paragraph 4 --- Availability and closing.}\\
\en{I would welcome the opportunity to discuss my application in an interview. I am available
from \ldots\ Thank you for considering my application. Yours sincerely, \ldots}
\end{anglais}
\end{dialoguebox}

\subsection{Modèle 5 --- CV anglais (les différences à connaître)}
\begin{tableau}[CV britannique/américain : ce qu'on met et ce qu'on ne met pas]
\begin{tabularx}{\linewidth}{@{}p{5.0cm}X@{}}
\entete{Élément} & \entete{Règle} \tabularnewline
\hline
Photo & \textbf{Pas de photo} dans les pays anglophones (sauf demande explicite). \tabularnewline
Âge, situation familiale, religion & \textbf{À ne jamais mentionner} --- c'est illégal ou très mal vu. \tabularnewline
Titre & \en{Accountant with 6 years' experience in the retail sector} --- un titre avec une valeur, pas seulement un métier. \tabularnewline
Profil (\en{Personal statement}) & 3 lignes en haut du CV : qui vous êtes, ce que vous savez faire, ce que vous cherchez. \tabularnewline
Expérience & Ordre \textbf{anti-chronologique} ; puces commençant par un \textbf{verbe d'action} : \en{Managed, Developed, Increased, Organised, Trained, Reduced}. \tabularnewline
Résultats & Toujours chiffrés : \en{increased sales by 20\,\%}, \en{trained a team of 12}, \en{reduced costs by 3 million}. \tabularnewline
Formation & \en{Education} --- diplômes, dates, établissement. \tabularnewline
Compétences & \en{Skills} --- langues avec niveau CECRL (\en{French: native; English: B2}), informatique, permis. \tabularnewline
Références & \en{References available on request} --- ne listez pas de noms. \tabularnewline
Longueur & 1 à 2 pages maximum. \tabularnewline
\end{tabularx}
\end{tableau}

\subsection{Modèle 6 --- Essai argumentatif en 5 paragraphes}
\begin{dialoguebox}[Structure classique exigée dans les examens]
\begin{anglais}
\textbf{1. Introduction :} amorce + reformulation du sujet + annonce du plan.\\
\en{In recent years, \ldots has become a topic of debate. While some argue that \ldots,
others believe that \ldots\ This essay will examine both views before reaching a conclusion.}

\textbf{2. Argument 1 :} idée + explication + exemple.\\
\en{Firstly, \ldots\ This is because \ldots\ For example, \ldots}

\textbf{3. Argument 2 :} même structure.\\
\en{Secondly, \ldots\ In addition, \ldots\ A clear example of this is \ldots}

\textbf{4. Contre-argument et réfutation :}\\
\en{Admittedly, some people claim that \ldots\ However, this argument overlooks the fact
that \ldots}

\textbf{5. Conclusion :} synthèse + opinion + ouverture.\\
\en{In conclusion, although \ldots, it seems clear that \ldots\ Ultimately, the question is
not whether \ldots but how \ldots}
\end{anglais}
\end{dialoguebox}

\subsection{Modèle 7 --- Compte rendu de réunion}
\begin{dialoguebox}[Structure]
\begin{anglais}
\textbf{Meeting minutes --- 12 May 2026, 10:00, Head Office}\\
\textbf{Present:} \ldots \quad \textbf{Apologies:} \ldots \quad \textbf{Chair:} \ldots

\textbf{1. Review of Q1 results.} Sales increased by 12\,\% compared with the same period
last year. Costs also rose, so the margin is slightly lower. Action: M. Diop to send the
detailed breakdown by 19 May.

\textbf{2. New market.} The team agreed to postpone the launch until September. Action:
F. Ndiaye to prepare a revised budget.

\textbf{3. Any other business.} A follow-up meeting will be held next week.

\textbf{Next meeting:} 26 May 2026, 10:00.
\end{anglais}
\end{dialoguebox}
\textbf{Formules fixes :} \en{Action: [nom] to [verbe] by [date]} ;
\en{The team agreed to\ldots} ; \en{It was decided that\ldots} ;
\en{Apologies:} (absents excusés) ; \en{Any other business} (questions diverses).

\subsection{Modèle 8 --- Message court (WhatsApp, SMS professionnel)}
\begin{tableau}[Registre court mais correct]
\begin{tabularx}{\linewidth}{@{}p{6.4cm}X@{}}
\entete{Message} & \entete{Emploi} \tabularnewline
\hline
\en{Hi Ali, are you free at 3? Need to talk about the report.} & Court, neutre, efficace \tabularnewline
\en{Running 10 min late, sorry!} & Retard --- acceptable même au travail \tabularnewline
\en{Got it, thanks!} & Accusé de réception \tabularnewline
\en{Could you send me the file when you have a moment?} & Demande polie sans urgence \tabularnewline
\en{Thanks again for today --- much appreciated.} & Remerciement après une réunion \tabularnewline
\en{Just a quick reminder about tomorrow's meeting at 10.} & Rappel poli \tabularnewline
\en{No worries, we'll sort it out tomorrow.} & Dédramatiser \tabularnewline
\end{tabularx}
\end{tableau}

\section{Partie E --- Le journal écrit : 60 invites}
Écrire cinq à dix phrases par jour est l'exercice le plus rentable de tout le livre : il
révèle immédiatement quels mots et quelles structures vous manquent.

\begin{tableau}[60 invites de journal (deux par jour pendant un mois)]
\footnotesize
\begin{multicols}{2}
\begin{anglais}
\begin{enumerate}\setlength{\itemsep}{1pt}
\item Describe your morning in detail.
\item What made you smile today?
\item Describe your bedroom to a blind person.
\item Write about someone who helped you this week.
\item What are you looking forward to?
\item Describe the best meal you have ever eaten.
\item What would you do with a free week?
\item Write about a fear you would like to overcome.
\item Describe your favourite place in your city.
\item What did you learn this week?
\item Write a letter to yourself in five years.
\item Describe the weather and how it affects you.
\item What habit would you like to change?
\item Write about your first job.
\item Describe a stranger you remember.
\item What are you grateful for today?
\item Write about a mistake that taught you something.
\item Describe your daily route to work.
\item What does success mean to you?
\item Write about a song that moves you.
\item Describe your family in five sentences each.
\item What would you change about your country?
\item Write about a time you were very tired.
\item Describe something you own and treasure.
\item What do you want to be remembered for?
\item Write about an argument you had.
\item Describe a market or a busy street.
\item What is the hardest thing about learning English?
\item Write about a person you admire from afar.
\item Describe a journey that went wrong.
\item What do you do when you are angry?
\item Write about a tradition in your family.
\item Describe your phone and what it does for you.
\item What would you tell a younger version of yourself?
\item Write about a book or film that changed you.
\item Describe the sound of your neighbourhood.
\item What are you avoiding doing? Why?
\item Write about a kindness you received.
\item Describe your ideal working day.
\item What is overrated in your opinion?
\item Write about a time you were lost.
\item Describe your hands and what they have done.
\item What do you believe but cannot prove?
\item Write about a celebration you attended.
\item Describe your relationship with money.
\item What skill would you like to have?
\item Write about a teacher who mattered.
\item Describe something you built or fixed.
\item What makes you laugh?
\item Write about a promise you made.
\item Describe the view from your window right now.
\item What would you do if you had no fear?
\item Write about a conversation you regret.
\item Describe your favourite time of day.
\item What do you need to let go of?
\item Write about a place you will never return to.
\item Describe your English level honestly.
\item What are your three goals for next month?
\item Write about someone you miss.
\item Summarise your month in ten sentences.
\end{enumerate}
\end{anglais}
\end{multicols}
\end{tableau}

\begin{astuce}[La méthode des trois jets]
\textbf{Jet 1 (5 min) :} écrivez sans vous arrêter, sans dictionnaire, sans corriger.
\textbf{Jet 2 (5 min) :} relisez et corrigez la grammaire, les temps, les articles ---
c'est là que vous apprenez. \textbf{Jet 3 (3 min, deux fois par semaine) :} faites corriger
par un outil ou un partenaire, et notez \textbf{une seule} erreur récurrente à surveiller
la semaine suivante. Une erreur corrigée à la fois : c'est la seule méthode qui fonctionne.
\end{astuce}

\begin{exo}[17.1 --- Orthographe]
Écrivez la forme correcte :
\en{stop (past), plan (-ing), study (3rd person), happy (comparative), leaf (plural),
life (plural), travel (past, BrE), colour (AmE), centre (AmE), practise (noun, BrE),
organise (AmE), programme (AmE)}.
\end{exo}
\begin{corrige}[17.1]
\en{stopped, planning, studies, happier, leaves, lives, travelled, color, center, practice,
organize, program}.
\end{corrige}

\begin{exo}[17.2 --- Ponctuation]
Corrigez ces cinq phrases (ponctuation et majuscules) :
\begin{enumerate}
  \item \en{*hello ! how are you today ?}
  \item \en{*i was born in 12.03.1995 in thies}
  \item \en{*he said ``i'm tired'' , and then he left}
  \item \en{*the meeting is on monday at 14.30}
  \item \en{*i bought apples , bananas and oranges}
\end{enumerate}
\end{exo}
\begin{corrige}[17.2]
1. \en{Hello! How are you today?} 2. \en{I was born on 12 March 1995 in Thiès.}
3. \en{He said, ``I'm tired,'' and then he left.} 4. \en{The meeting is on Monday at 2.30 pm.}
5. \en{I bought apples, bananas, and oranges.}
\end{corrige}

\begin{exo}[17.3 --- Écrivez]
Rédigez (150 à 200 mots) : un e-mail formel à une école de langue pour demander les dates,
les tarifs et les modalités d'inscription. Utilisez au moins :
\en{I am writing to\ldots}, \en{I would be grateful if you could\ldots},
\en{Please find attached\ldots}, \en{I look forward to hearing from you},
\en{Yours faithfully}.
\end{exo}
\begin{corrige}[17.3 --- exemple de réponse attendue]
\begin{anglais}
\textbf{Subject: Enquiry about English courses --- September intake}

Dear Sir or Madam,

I am writing to enquire about the intensive English courses advertised on your website.

I would be grateful if you could let me know the start dates for the September intake, the
weekly timetable, and the total cost of the course. I would also like to know whether an
entrance test is required, and whether payment can be made in two instalments.

I have been studying English on my own for one year, and I believe my level is
approximately B1. Please find attached a short description of my goals, in case it is useful
for the placement test.

I look forward to hearing from you.

Yours faithfully,\\
Awa Sarr
\end{anglais}
\end{corrige}

\begin{mission}[Chapitre 17 --- lire et écrire chaque jour]
\textbf{Tous les jours :} 15 minutes de lecture extensive (facile, sans dictionnaire) +
10 minutes de journal (5 phrases minimum).\\
\textbf{Deux fois par semaine :} un texte de ce livre relu à voix haute, puis réécrit de
mémoire --- exercice puissant qui installe les structures.\\
\textbf{Une fois par semaine :} un e-mail complet (modèle 2, 3 ou 4) envoyé à un vrai
destinataire --- partenaire de langue, école, entreprise. L'écrit envoyé progresse trois
fois plus vite que l'écrit dans un cahier.
\end{mission}

% =====================================================================
%  CHAPITRE 18 — LES 100 ERREURS TYPOIQUES DES FRANCOPHONES
% =====================================================================
\chapter{Les 100 erreurs qui trahissent un francophone}

Ce chapitre est un outil de diagnostic. Lisez-le une première fois en entier (30 minutes),
puis revenez-y une fois par mois. Chaque erreur est présentée avec sa forme fautive,
sa forme correcte et l'explication. \textbf{Ne travaillez qu'une erreur à la fois} :
choisissez-en trois par semaine, affichez-les, et corrigez-les dans tous vos écrits et
enregistrements jusqu'à l'automatisme.

\begin{tableau}[Erreurs 1 à 25 --- le verbe et ses temps]
\footnotesize
\begin{longtable}{@{}p{5.0cm}p{4.6cm}p{4.2cm}@{}}
\entete{\ko Forme fautive} & \entete{\ok Forme correcte} & \entete{Explication} \tabularnewline
\endfirsthead
\entete{\ko Faux} & \entete{\ok Juste} & \entete{Pourquoi} \tabularnewline
\endhead
\hline
\en{*I have 25 years.} & \en{I'm 25 (years old).} & L'âge se dit avec \en{be}. \tabularnewline
\en{*I have cold / hunger / fear.} & \en{I'm cold / hungry / afraid.} & Sensations avec \en{be} + adjectif. \tabularnewline
\en{*I am agree.} & \en{I agree.} & \en{agree} est un verbe, pas un adjectif. \tabularnewline
\en{*I am knowing him.} & \en{I know him.} & Verbes d'état jamais au \en{-ing}. \tabularnewline
\en{*I want that you come.} & \en{I want you to come.} & \en{want + personne + to}. \tabularnewline
\en{*He don't work.} & \en{He doesn't work.} & 3\textsuperscript{e} personne : \en{doesn't}. \tabularnewline
\en{*He doesn't has a car.} & \en{He doesn't have a car.} & Après \en{does}, base verbale. \tabularnewline
\en{*Did you went?} & \en{Did you go?} & Après \en{did}, base verbale. \tabularnewline
\en{*She don't like it.} & \en{She doesn't like it.} & Idem 3\textsuperscript{e} personne. \tabularnewline
\en{*I have seen him yesterday.} & \en{I saw him yesterday.} & Date précise \ra passé simple. \tabularnewline
\en{*I live here since 2015.} & \en{I've lived here since 2015.} & « Depuis » \ra \en{present perfect}. \tabularnewline
\en{*I am here since Monday.} & \en{I've been here since Monday.} & Idem avec \en{be}. \tabularnewline
\en{*I will call you when I will arrive.} & \en{I'll call you when I arrive.} & Présent après \en{when}. \tabularnewline
\en{*If I will have time\ldots} & \en{If I have time\ldots} & Jamais \en{will} après \en{if}. \tabularnewline
\en{*I used to go there when I was young} \emph{(correct)} \en{*I am used to go} & \en{I'm used to going.} & \en{be used to} + \en{-ing}. \tabularnewline
\en{*He can to swim.} & \en{He can swim.} & Pas de \en{to} après un modal. \tabularnewline
\en{*She cans / *he musts.} & \en{She can / he must.} & Les modaux ne prennent jamais \en{-s}. \tabularnewline
\en{*I mustn't come} \emph{(= interdit)} & \en{I don't have to come.} & « Pas obligé » \ra \en{don't have to}. \tabularnewline
\en{*I have gone to Paris twice.} & \en{I have been to Paris twice.} & Expérience \ra \en{been}. \tabularnewline
\en{*She's gone to the market} \emph{(elle est au marché)} & \en{She's gone to the market.} & Correct --- mais vérifiez le sens voulu. \tabularnewline
\en{*I didn't nothing say.} & \en{I didn't say anything.} & Une seule négation. \tabularnewline
\en{*I don't never go.} & \en{I never go.} & \en{never} suffit. \tabularnewline
\en{*I have make a mistake.} & \en{I have made a mistake.} & Participe passé, pas base. \tabularnewline
\en{*She has wrote it.} & \en{She has written it.} & 3\textsuperscript{e} forme. \tabularnewline
\en{*I didn't knew.} & \en{I didn't know.} & Base après \en{didn't}. \tabularnewline
\end{longtable}
\end{tableau}

\begin{tableau}[Erreurs 26 à 50 --- articles, noms, prépositions]
\footnotesize
\begin{longtable}{@{}p{5.0cm}p{4.6cm}p{4.2cm}@{}}
\entete{\ko Forme fautive} & \entete{\ok Forme correcte} & \entete{Explication} \tabularnewline
\endfirsthead
\entete{\ko Faux} & \entete{\ok Juste} & \entete{Pourquoi} \tabularnewline
\endhead
\hline
\en{*I like the music.} & \en{I like music.} & Généralité \ra pas d'article. \tabularnewline
\en{*The life is beautiful.} & \en{Life is beautiful.} & Concept abstrait général. \tabularnewline
\en{*I live in the Senegal.} & \en{I live in Senegal.} & Pays sans \en{the} (sauf pluriel/États). \tabularnewline
\en{*I go to the school} \emph{(étudier)} & \en{I go to school.} & Fonction \ra pas d'article. \tabularnewline
\en{*I had the breakfast at 7.} & \en{I had breakfast at 7.} & Repas sans article. \tabularnewline
\en{*She speaks the English.} & \en{She speaks English.} & Langues sans article. \tabularnewline
\en{*I play the football.} & \en{I play football.} & Sports sans article (mais \en{play \textbf{the} piano}). \tabularnewline
\en{*He is in the Monday.} & \en{He arrives on Monday.} & \en{on} + jour, sans article. \tabularnewline
\en{*I will see you in tomorrow.} & \en{I'll see you tomorrow.} & Jamais de préposition avant \en{tomorrow}. \tabularnewline
\en{*in the night} \emph{(en général)} & \en{at night} & \en{at night} mais \en{in the morning}. \tabularnewline
\en{*an university} & \en{a university} & \en{u} = \pron{iou} (consonne). \tabularnewline
\en{*a hour} & \en{an hour} & \en{h} muet. \tabularnewline
\en{*two informations} & \en{two pieces of information} & Indénombrable. \tabularnewline
\en{*many advices} & \en{a lot of advice} & Indénombrable. \tabularnewline
\en{*a news} & \en{a piece of news / some news} & \en{news} est singulier. \tabularnewline
\en{*the news are good} & \en{the news is good} & Accord au singulier. \tabularnewline
\en{*a hairs} & \en{hair} (chevelure) & \en{hair} est indénombrable ici. \tabularnewline
\en{*peoples} & \en{people} & Déjà pluriel. \tabularnewline
\en{*childrens} & \en{children} & Déjà pluriel. \tabularnewline
\en{*informations about} & \en{information about} & Indénombrable. \tabularnewline
\en{*depend of} & \en{depend on} & Préposition imposée. \tabularnewline
\en{*listen music} & \en{listen to music} & Préposition imposée. \tabularnewline
\en{*discuss about} & \en{discuss (something)} & Pas de préposition. \tabularnewline
\en{*married with} & \en{married to} & Préposition imposée. \tabularnewline
\en{*I have fear of dogs.} & \en{I'm afraid of dogs.} & Structure anglaise. \tabularnewline
\end{longtable}
\end{tableau}

\begin{tableau}[Erreurs 51 à 75 --- structure de phrase et vocabulaire]
\footnotesize
\begin{longtable}{@{}p{5.0cm}p{4.6cm}p{4.2cm}@{}}
\entete{\ko Forme fautive} & \entete{\ok Forme correcte} & \entete{Explication} \tabularnewline
\endfirsthead
\entete{\ko Faux} & \entete{\ok Juste} & \entete{Pourquoi} \tabularnewline
\endhead
\hline
\en{*Is cold today.} & \en{It's cold today.} & Sujet \en{it} obligatoire. \tabularnewline
\en{*Have a problem.} & \en{There's a problem.} & \en{il y a} = \en{there is}. \tabularnewline
\en{*It has two rooms.} & \en{It has got / There are two rooms.} & \en{il y a} $\neq$ \en{it has}. \tabularnewline
\en{*I am a Senegalese.} & \en{I'm Senegalese.} & Nationalité : adjectif sans article. \tabularnewline
\en{*She is a doctor woman.} & \en{She is a doctor.} & Pas de redondance. \tabularnewline
\en{*I said you something.} & \en{I told you something.} & \en{tell + personne}, \en{say + chose}. \tabularnewline
\en{*He explained me the problem.} & \en{He explained the problem to me.} & \en{explain} n'accepte pas le double complément direct. \tabularnewline
\en{*I demand you a question.} & \en{I'll ask you a question.} & \en{demand} = exiger. \tabularnewline
\en{*Make attention!} & \en{Be careful! / Watch out!} & Calque de « faire attention ». \tabularnewline
\en{*Make a photo.} & \en{Take a photo.} & Collocation imposée. \tabularnewline
\en{*Do a mistake.} & \en{Make a mistake.} & \en{make} pour créer, \en{do} pour exécuter. \tabularnewline
\en{*Do a party.} & \en{Have / throw a party.} & Collocation. \tabularnewline
\en{*Make sport.} & \en{Do sport / play sports.} & Collocation. \tabularnewline
\en{*I have a good relation with him.} & \en{I get on well with him.} & \en{relation} = relation amoureuse ou diplomatique. \tabularnewline
\en{*The people is nice here.} & \en{The people are nice here.} & \en{people} est pluriel. \tabularnewline
\en{*I am boring.} \emph{(je m'ennuie)} & \en{I'm bored.} & \en{-ed} = ressenti. \tabularnewline
\en{*The film was bored.} & \en{The film was boring.} & \en{-ing} = cause. \tabularnewline
\en{*I have 32 years old.} & \en{I'm 32.} & Redondance. \tabularnewline
\en{*It's difficult for me of understanding.} & \en{It's difficult for me to understand.} & \en{for + personne + to}. \tabularnewline
\en{*I think to go.} & \en{I'm thinking of going.} & \en{think of + -ing} (projeter). \tabularnewline
\en{*I look my phone.} & \en{I look at my phone.} & \en{look at} ; \en{watch} pour ce qui bouge. \tabularnewline
\en{*I see TV every night.} & \en{I watch TV every night.} & \en{watch} = attention prolongée. \tabularnewline
\en{*The weather is very hot} \emph{(correct, mais)} & \en{It's very hot.} & Plus naturel. \tabularnewline
\en{*How is it called?} & \en{What is it called?} & \en{what} pour un nom. \tabularnewline
\en{*How do you call this?} & \en{What do you call this?} & Même règle. \tabularnewline
\end{longtable}
\end{tableau}

\begin{tableau}[Erreurs 76 à 100 --- discours, prononciation, style]
\footnotesize
\begin{longtable}{@{}p{5.0cm}p{4.6cm}p{4.2cm}@{}}
\entete{\ko Forme fautive} & \entete{\ok Forme correcte} & \entete{Explication} \tabularnewline
\endfirsthead
\entete{\ko Faux} & \entete{\ok Juste} & \entete{Pourquoi} \tabularnewline
\endhead
\hline
\en{*He asked me where did I live.} & \en{He asked me where I lived.} & Question indirecte : ordre affirmatif. \tabularnewline
\en{*I don't know what does it mean.} & \en{I don't know what it means.} & Idem. \tabularnewline
\en{*The man who I met him.} & \en{The man (who) I met.} & Ne pas répéter le pronom. \tabularnewline
\en{*My brother, that lives in Dakar\ldots} & \en{My brother, who lives in Dakar\ldots} & \en{that} interdit en relative non définissante. \tabularnewline
\en{*I am living here since 2010.} & \en{I have been living here since 2010.} & Durée \ra \en{perfect}. \tabularnewline
\en{*She said me that\ldots} & \en{She told me that\ldots} & \en{say} ne prend pas de personne directement. \tabularnewline
\en{*I am since one hour here.} & \en{I have been here for an hour.} & \en{for} + durée. \tabularnewline
\en{*I search my keys.} & \en{I'm looking for my keys.} & \en{search} = fouiller un lieu. \tabularnewline
\en{*I will see you to tomorrow.} & \en{I'll see you tomorrow.} & Préposition inutile. \tabularnewline
\en{*Actually I'm working} \emph{(= actuellement)} & \en{Currently I'm working.} & \en{actually} = en fait. \tabularnewline
\en{*Eventually we'll meet} \emph{(= éventuellement)} & \en{We'll possibly meet.} & \en{eventually} = finalement. \tabularnewline
\en{*I visited my grandmother} \emph{(rendre visite)} & \en{I visited my grandmother.} \ok & Correct ici --- mais \en{visit a place} est plus fréquent ; pour les personnes, \en{see} est souvent plus naturel. \tabularnewline
\en{*I want a coffee} \emph{(au café)} & \en{I'd like a coffee, please.} & Registre : \en{want} est brutal. \tabularnewline
\en{*Give me the salt.} \emph{(à table, adulte)} & \en{Could you pass me the salt, please?} & Politesse obligatoire entre adultes. \tabularnewline
\en{*I have the possibility to come.} & \en{I can come.} / \en{I have the chance to come.} & \en{possibility} = faisabilité technique. \tabularnewline
\en{*It is necessary that you come.} & \en{You need to come.} / \en{You must come.} & Éviter les structures lourdes. \tabularnewline
\en{*I suggest you to come.} & \en{I suggest coming.} / \en{I suggest that you come.} & \en{suggest} refuse \en{suggest sb to do}. \tabularnewline
\en{*He explained me everything} & \en{He explained everything to me.} & Structure imposée. \tabularnewline
\en{*I prefer coffee than tea.} & \en{I prefer coffee to tea.} & \en{prefer A to B}. \tabularnewline
\en{*I am agree with you} & \en{I agree with you.} & Déjà vu, mais faute tenace. \tabularnewline
\en{*This is more better.} & \en{This is better.} & Jamais \en{more} + comparatif. \tabularnewline
\en{*He runs more fast than me.} & \en{He runs faster than me.} & Comparatif en \en{-er}. \tabularnewline
\en{*I have so much friends.} & \en{I have so many friends.} & \en{many} + dénombrable. \tabularnewline
\en{*There is many people.} & \en{There are many people.} & Accord avec le pluriel. \tabularnewline
\en{*I look forward to hear from you.} & \en{I look forward to hearing from you.} & \en{to} est une préposition ici. \tabularnewline
\end{longtable}
\end{tableau}

\section{Les cinq fautes de prononciation à corriger en premier}
\begin{enumerate}
  \item \textbf{Le \en{h} avalé} : \en{he, him, house, happy} --- soufflez.
  \item \textbf{Le \en{th} transformé} : \en{think} $\neq$ \en{sink} ; \en{this} $\neq$ \en{dis}.
  \item \textbf{Le \en{w} dit \en{v}} : \en{we} $\neq$ \en{ve} ; \en{water} $\neq$ \en{vater}.
  \item \textbf{Le \en{e} final ajouté} : \en{book} \pron{bouk}, pas \pron{bou-keu}.
  \item \textbf{L'accent tonique plat} : \en{baNAna}, \en{inforMAtion}, \en{toDAY}.
\end{enumerate}

\section{Auto-diagnostic : 40 phrases à corriger}
Couvrez la colonne de droite, corrigez les 40 phrases, puis comparez. Score sur 40 :
\textbf{35 et plus} = niveau B1 solide sur les structures de base ; \textbf{25-34} =
niveau A2, travaillez trois erreurs par semaine ; \textbf{moins de 25} = reprenez les
chapitres 8 et 9 en priorité.

\begin{tableau}[Test de diagnostic]
\footnotesize
\begin{longtable}{@{}p{0.9cm}p{6.2cm}p{6.7cm}@{}}
\entete{\#} & \entete{Phrase à corriger} & \entete{Correction} \tabularnewline
\endfirsthead
\entete{\#} & \entete{Phrase à corriger} & \entete{Correction} \tabularnewline
\endhead
1 & \en{I have 40 years.} & \en{I'm 40 (years old).} \tabularnewline
2 & \en{She don't like fish.} & \en{She doesn't like fish.} \tabularnewline
3 & \en{I am agree with you.} & \en{I agree with you.} \tabularnewline
4 & \en{He can to help us.} & \en{He can help us.} \tabularnewline
5 & \en{I live here since five years.} & \en{I've lived here for five years.} \tabularnewline
6 & \en{Yesterday I have seen Ali.} & \en{Yesterday I saw Ali.} \tabularnewline
7 & \en{I will call you when I will arrive.} & \en{I'll call you when I arrive.} \tabularnewline
8 & \en{If it will rain, we'll stay.} & \en{If it rains, we'll stay.} \tabularnewline
9 & \en{I have cold.} & \en{I'm cold.} \tabularnewline
10 & \en{The life is hard.} & \en{Life is hard.} \tabularnewline
11 & \en{I like the football.} & \en{I like football.} \tabularnewline
12 & \en{He explained me the rule.} & \en{He explained the rule to me.} \tabularnewline
13 & \en{She said me she was tired.} & \en{She told me she was tired.} \tabularnewline
14 & \en{I am boring today.} & \en{I'm bored today.} \tabularnewline
15 & \en{Make attention to the step.} & \en{Watch your step.} / \en{Be careful.} \tabularnewline
16 & \en{I did a mistake.} & \en{I made a mistake.} \tabularnewline
17 & \en{He asked me where do I live.} & \en{He asked me where I live(d).} \tabularnewline
18 & \en{The man who I met him yesterday.} & \en{The man (who) I met yesterday.} \tabularnewline
19 & \en{My sister, that works in a bank\ldots} & \en{My sister, who works in a bank\ldots} \tabularnewline
20 & \en{I search my glasses.} & \en{I'm looking for my glasses.} \tabularnewline
21 & \en{She is married with a teacher.} & \en{She is married to a teacher.} \tabularnewline
22 & \en{I look TV every evening.} & \en{I watch TV every evening.} \tabularnewline
23 & \en{Actually I'm busy.} \emph{(= actuellement)} & \en{Currently / At the moment I'm busy.} \tabularnewline
24 & \en{I prefer tea than coffee.} & \en{I prefer tea to coffee.} \tabularnewline
25 & \en{This is more easy.} & \en{This is easier.} \tabularnewline
26 & \en{I have so much work.} \emph{(correct !)} & \en{I have so much work.} --- \en{work} est indénombrable, donc \en{much} est juste. \tabularnewline
27 & \en{There is many problems.} & \en{There are many problems.} \tabularnewline
28 & \en{I look forward to meet you.} & \en{I look forward to meeting you.} \tabularnewline
29 & \en{I want that you help me.} & \en{I want you to help me.} \tabularnewline
30 & \en{I am knowing the answer.} & \en{I know the answer.} \tabularnewline
31 & \en{I didn't went there.} & \en{I didn't go there.} \tabularnewline
32 & \en{She has wrote a letter.} & \en{She has written a letter.} \tabularnewline
33 & \en{You mustn't pay} \emph{(= pas obligé)} & \en{You don't have to pay.} \tabularnewline
34 & \en{I used to go there} \emph{(= j'ai l'habitude d'y aller)} & \en{I usually go there.} \tabularnewline
35 & \en{I'm used to wake up early.} & \en{I'm used to waking up early.} \tabularnewline
36 & \en{I have been in Paris two times.} & \en{I have been to Paris twice.} \tabularnewline
37 & \en{The news are bad.} & \en{The news is bad.} \tabularnewline
38 & \en{He gave me an advice.} & \en{He gave me a piece of advice.} \tabularnewline
39 & \en{I'll see you in tomorrow.} & \en{I'll see you tomorrow.} \tabularnewline
40 & \en{How is it called in English?} & \en{What is it called in English?} \tabularnewline
\end{longtable}
\end{tableau}

\begin{mission}[Chapitre 18 --- trois erreurs par semaine]
Chaque lundi, choisissez \textbf{trois} erreurs de ce chapitre (pas plus). Écrivez-les sur
un papier que vous gardez dans votre poche ou sur votre bureau. Pendant toute la semaine,
chaque fois que vous parlez ou écrivez en anglais, vérifiez ces trois points. Le dimanche,
faites un test : écrivez 10 phrases en utilisant les trois structures surveillées.
En 33 semaines, vous aurez éliminé les 100 erreurs --- et personne ne devinera plus que
vous êtes francophone en vous lisant.
\end{mission}


% ===== PARTIE VII — TESTS ==========================================
\part{Tests, suivi et certification}
% =====================================================================
%  CHAPITRE 19 — TESTS, SUIVI ET CERTIFICATION
% =====================================================================
\chapter{Tests, suivi et certification}

Sans mesure, pas de progrès --- seulement une impression. Ce chapitre vous donne les outils
pour savoir \textbf{où vous en êtes}, à tout moment, sans payer un examen.

\section{Test de positionnement (40 questions, 40 minutes)}
Faites ce test \textbf{avant} de commencer le livre, puis au mois 6 et au mois 12.
Ne regardez aucun corrigé, n'utilisez aucun dictionnaire. Notez vos réponses sur une feuille.

\subsection{Partie 1 --- Grammaire (questions 1 à 20)}
Choisissez la bonne réponse.

\begin{tableau}[Q1 à Q20]
\footnotesize
\begin{longtable}{@{}p{0.9cm}p{7.0cm}p{5.9cm}@{}}
\entete{\#} & \entete{Phrase} & \entete{Choix} \tabularnewline
\endfirsthead
\entete{\#} & \entete{Phrase} & \entete{Choix} \tabularnewline
\endhead
1 & \en{She \trou[2cm] in a hospital.} & a) work \ b) works \ c) working \ d) is work \tabularnewline
2 & \en{I \trou[2cm] 28 years old.} & a) have \ b) am \ c) has \ d) be \tabularnewline
3 & \en{\trou[2cm] you speak English?} & a) Are \ b) Do \ c) Does \ d) Is \tabularnewline
4 & \en{There aren't \trou[2cm] chairs.} & a) some \ b) much \ c) any \ d) a \tabularnewline
5 & \en{This is \trou[2cm] best restaurant in town.} & a) the \ b) a \ c) --- \ d) an \tabularnewline
6 & \en{I \trou[2cm] him yesterday.} & a) have seen \ b) see \ c) saw \ d) seen \tabularnewline
7 & \en{She \trou[2cm] here since 2019.} & a) lives \ b) is living \ c) has lived \ d) lived \tabularnewline
8 & \en{Look! It \trou[2cm].} & a) rains \ b) is raining \ c) rained \ d) rain \tabularnewline
9 & \en{I \trou[2cm] to the gym twice a week.} & a) am going \ b) go \ c) goes \ d) have gone \tabularnewline
10 & \en{If it rains, we \trou[2cm] at home.} & a) will stay \ b) would stay \ c) stayed \ d) stay will \tabularnewline
11 & \en{If I \trou[2cm] rich, I would travel.} & a) am \ b) was \ c) were \ d) will be \tabularnewline
12 & \en{You \trou[2cm] smoke here --- it's forbidden.} & a) don't have to \ b) mustn't \ c) shouldn't \ d) couldn't \tabularnewline
13 & \en{That's the man \trou[2cm] car was stolen.} & a) who \ b) which \ c) whose \ d) that \tabularnewline
14 & \en{I enjoy \trou[2cm] to music.} & a) listen \ b) to listen \ c) listening \ d) listened \tabularnewline
15 & \en{She promised \trou[2cm] early.} & a) coming \ b) to come \ c) come \ d) came \tabularnewline
16 & \en{The meeting \trou[2cm] cancelled last week.} & a) is \ b) has been \ c) was \ d) were \tabularnewline
17 & \en{He asked me where \trou[2cm].} & a) did I live \ b) I lived \ c) do I live \ d) lived I \tabularnewline
18 & \en{I should \trou[2cm] harder at school.} & a) have worked \ b) has worked \ c) working \ d) to work \tabularnewline
19 & \en{It's time we \trou[2cm] home.} & a) go \ b) went \ c) will go \ d) gone \tabularnewline
20 & \en{You're tired, \trou[2cm]?} & a) aren't you \ b) don't you \ c) isn't it \ d) are you \tabularnewline
\end{longtable}
\end{tableau}

\subsection{Partie 2 --- Vocabulaire et expressions (questions 21 à 30)}
\begin{tableau}[Q21 à Q30]
\footnotesize
\begin{longtable}{@{}p{0.9cm}p{7.0cm}p{5.9cm}@{}}
\entete{\#} & \entete{Phrase} & \entete{Choix} \tabularnewline
\endfirsthead
\entete{\#} & \entete{Phrase} & \entete{Choix} \tabularnewline
\endhead
21 & \en{I'm looking \trou[2cm] my keys.} & a) at \ b) for \ c) after \ d) up \tabularnewline
22 & \en{It depends \trou[2cm] the weather.} & a) of \ b) on \ c) from \ d) to \tabularnewline
23 & \en{Can you \trou[2cm] me a favour?} & a) make \ b) do \ c) give \ d) take \tabularnewline
24 & \en{He \trou[2cm] up smoking last year.} & a) gave \ b) took \ c) put \ d) got \tabularnewline
25 & \en{``Actually'' means:} & a) actuellement \ b) en fait \ c) activement \ d) éventuellement \tabularnewline
26 & \en{A ``library'' is:} & a) une librairie \ b) une bibliothèque \ c) un laboratoire \ d) un libre \tabularnewline
27 & \en{``To attend a meeting'' means:} & a) attendre \ b) assister à \ c) organiser \ d) annuler \tabularnewline
28 & \en{Choose the polite form at a café:} & a) \en{Give me a coffee.} \ b) \en{I want a coffee.} \ c) \en{I'd like a coffee, please.} \ d) \en{Coffee!} \tabularnewline
29 & \en{``It's a piece of cake'' means:} & a) c'est un gâteau \ b) c'est très facile \ c) c'est cher \ d) c'est fini \tabularnewline
30 & \en{``I'm fed up'' means:} & a) j'ai mangé \ b) j'en ai marre \ c) je suis repu \ d) je suis fatigué \tabularnewline
\end{longtable}
\end{tableau}

\subsection{Partie 3 --- Correction d'erreurs (questions 31 à 40)}
Chaque phrase contient \textbf{une} erreur. Trouvez-la et corrigez-la.
\begin{enumerate}\setlength{\itemsep}{3pt}
  \item \en{I am agree with you.}
  \item \en{She don't like tea.}
  \item \en{I live in Dakar since 2012.}
  \item \en{He can to drive a car.}
  \item \en{I have seen him last week.}
  \item \en{The people is very friendly here.}
  \item \en{He explained me the problem.}
  \item \en{I'm boring this evening.}
  \item \en{This book is more interesting that the other one.}
  \item \en{I look forward to hear from you.}
\end{enumerate}

\subsection{Corrigé du test de positionnement}
\begin{corrige}[Parties 1 et 2 --- réponses]
\footnotesize
1\,b \quad 2\,b \quad 3\,b \quad 4\,c \quad 5\,a \quad 6\,c \quad 7\,c \quad 8\,b \quad
9\,b \quad 10\,a \quad 11\,c \quad 12\,b \quad 13\,c \quad 14\,c \quad 15\,b \quad
16\,c \quad 17\,b \quad 18\,a \quad 19\,b \quad 20\,a \quad 21\,b \quad 22\,b \quad
23\,b \quad 24\,a \quad 25\,b \quad 26\,b \quad 27\,b \quad 28\,c \quad 29\,b \quad 30\,b
\end{corrige}

\begin{corrige}[Partie 3 --- corrections]
1. \en{I agree with you.} 2. \en{She doesn't like tea.}
3. \en{I have lived / have been living in Dakar since 2012.}
4. \en{He can drive a car.} 5. \en{I saw him last week.}
6. \en{The people are very friendly here.} 7. \en{He explained the problem to me.}
8. \en{I'm bored this evening.} 9. \en{\ldots more interesting \textbf{than} the other one.}
10. \en{I look forward to \textbf{hearing} from you.}
\end{corrige}

\begin{tableau}[Interprétation de votre score sur 40]
\begin{tabularx}{\linewidth}{@{}p{3.0cm}p{3.0cm}X@{}}
\entete{Score} & \entete{Niveau estimé} & \entete{Point de départ dans ce livre} \tabularnewline
\hline
0 à 12 & A1 débutant & Commencez au mois 1, chapitre 5. Ne sautez rien. \tabularnewline
13 à 20 & A1 confirmé & Commencez au mois 1 ou 2. Insistez sur les chapitres 8 et 9. \tabularnewline
21 à 28 & A2 & Commencez au mois 3. Faites le chapitre 8 en révision rapide. \tabularnewline
29 à 34 & B1 faible & Commencez au mois 4-5. Priorité : \en{present perfect}, articles, prépositions. \tabularnewline
35 à 40 & B1 confirmé / B2 & Commencez au mois 6. Votre priorité n'est plus la grammaire : c'est l'oral (chapitres 15 et 16) et le vocabulaire (chapitre 13). \tabularnewline
\end{tabularx}
\end{tableau}

\section{Les tests mensuels : modèle reproductible}
Plutôt que douze tests tout faits --- que vous finiriez par connaître par c\oe ur --- voici
un \textbf{modèle de 45 minutes} à reproduire chaque dernier dimanche du mois, avec le
contenu du mois.

\begin{tableau}[Structure du test mensuel]
\begin{tabularx}{\linewidth}{@{}p{2.4cm}p{1.6cm}X@{}}
\entete{Épreuve} & \entete{Durée} & \entete{Contenu} \tabularnewline
\hline
\textbf{1. Grammaire} & 10 min & 20 phrases à compléter tirées des unités du mois. Barème : 1 point par phrase. \textbf{Objectif : 16/20.} \tabularnewline
\textbf{2. Vocabulaire} & 10 min & 30 mots du mois : donnez le sens, ou donnez le mot anglais à partir du français. \textbf{Objectif : 24/30.} \tabularnewline
\textbf{3. Écoute} & 10 min & Un audio neuf de 2 minutes (niveau du mois), écouté deux fois, avec 8 questions de compréhension. \textbf{Objectif : 6/8.} \tabularnewline
\textbf{4. Écrit} & 10 min & 120 à 150 mots sur un sujet du mois. Comptez vos erreurs : \textbf{objectif : moins de 6 erreurs.} \tabularnewline
\textbf{5. Oral} & 5 min & 3 minutes de parole libre enregistrée sur un sujet tiré au hasard parmi les 30 questions du chapitre 15, puis auto-évaluation avec la grille ci-dessous. \tabularnewline
\end{tabularx}
\end{tableau}

\subsection{Test du mois 6 (exemple complet)}
\begin{exo}[Mois 6 --- grammaire, 20 phrases]
\begin{enumerate}\setlength{\itemsep}{2pt}
  \item \en{I (never / be) \trou[3cm] to Europe.}
  \item \en{She (work) \trou[3cm] here since March.}
  \item \en{While I (walk) \trou[3cm] home, it started to rain.}
  \item \en{If I (know) \trou[3cm] the answer, I would tell you.}
  \item \en{The bridge (build) \trou[3cm] in 1994.}
  \item \en{He said he (be) \trou[3cm] tired.}
  \item \en{You (should / call) \trou[3cm] me --- I was worried.}
  \item \en{That's the girl \trou[2cm] brother won the prize.}
  \item \en{I'm looking forward to \trou[3cm] (see) you.}
  \item \en{We (visit) \trou[3cm] Morocco three times so far.}
  \item \en{By the time we arrived, the film (already / start) \trou[3.5cm].}
  \item \en{She asked me where \trou[3.5cm] (I / buy) it.}
  \item \en{I'd rather you (not / smoke) \trou[3cm] inside.}
  \item \en{The more you practise, the (good) \trou[3cm] you get.}
  \item \en{He must (be) \trou[3cm] at home --- his car is outside.}
  \item \en{My phone (steal) \trou[3cm] on the bus last week.}
  \item \en{I (used to / play) \trou[3cm] football every day.}
  \item \en{There's not (much / many) \trou[2cm] time left.}
  \item \en{She's the most (experience) \trou[3.5cm] person in the team.}
  \item \en{I'll call you as soon as I (get) \trou[3cm] home.}
\end{enumerate}
\end{exo}
\begin{corrige}[Mois 6 --- grammaire]
1. \en{have never been} 2. \en{has been working} 3. \en{was walking} 4. \en{knew}
5. \en{was built} 6. \en{was} 7. \en{should have called} 8. \en{whose}
9. \en{seeing} 10. \en{have visited} 11. \en{had already started} 12. \en{I had bought}
13. \en{didn't smoke} 14. \en{better} 15. \en{be} 16. \en{was stolen}
17. \en{used to play} 18. \en{much} 19. \en{experienced} 20. \en{get}
\end{corrige}

\subsection{Test du mois 12 (exemple complet)}
\begin{exo}[Mois 12 --- production écrite (250 mots, 40 minutes)]
Sujet au choix :\\
\textbf{A.} \en{Some people believe that learning English online is as effective as taking
classes with a teacher. To what extent do you agree?}\\
\textbf{B.} \en{Write a formal email to a company complaining about a service you received,
and explaining what you expect them to do.}\\
\textbf{C.} \en{Describe a person who has influenced your life, and explain how.}
\end{exo}
\begin{corrige}[Mois 12 --- grille de correction (à utiliser soi-même)]
\textbf{Sur 20 points :}\\
\textbf{Structure (5)} : introduction claire, paragraphes distincts, conclusion.
Connecteurs variés (\en{however, therefore, in addition, admittedly, ultimately}).\\
\textbf{Grammaire (5)} : au moins trois temps différents utilisés correctement ;
subordonnées relatives ; un passif ; un conditionnel.\\
\textbf{Vocabulaire (5)} : au moins dix mots ou expressions de niveau B2
(\en{a wide range of, to keep in touch, it depends on, as far as I'm concerned,
to be worth doing, a piece of advice}).\\
\textbf{Précision (5)} : moins de huit erreurs au total ; aucune erreur sur les cinq
points surveillés (articles, \en{-s} de la 3\textsuperscript{e} personne, temps,
prépositions, \en{make/do}).
\end{corrige}

\begin{exo}[Mois 12 --- production orale (5 minutes enregistrées)]
Tirez un sujet au hasard et parlez 3 minutes sans préparation, puis répondez à deux
questions improvisées :
\begin{enumerate}
  \item \en{Is technology making us more or less connected?}
  \item \en{What would you change about the education system in your country?}
  \item \en{Describe a skill everyone should learn, and explain why.}
  \item \en{Should governments do more to protect the environment?}
  \item \en{What is the most important quality in a friend?}
  \item \en{Is it better to work for a large company or a small one?}
\end{enumerate}
\end{exo}

\section{Grille d'auto-évaluation orale}
À remplir après chaque enregistrement. Soyez sévère mais juste : c'est votre seul miroir.

\begin{tableau}[Grille de notation orale --- sur 5 par critère]
\begin{tabularx}{\linewidth}{@{}p{3.2cm}p{4.6cm}X@{}}
\entete{Critère} & \entete{5 = excellent} & \entete{Ce qu'il faut vérifier} \tabularnewline
\hline
\textbf{Fluidité} & Aucune pause de plus de 2 secondes ; pas de retour au français. &
Comptez les pauses longues et les mots français. \tabularnewline
\textbf{Étendue} & Vocabulaire varié, synonymes, expressions idiomatiques. &
Avez-vous répété 5 fois le même mot ? \tabularnewline
\textbf{Précision} & Erreurs rares et non gênantes ; temps corrects. &
Comptez les erreurs de temps, d'article, de \en{-s}. \tabularnewline
\textbf{Prononciation} & Sons corrects, accents toniques placés, intonation vivante. &
Un anglophone comprendrait-il du premier coup ? \tabularnewline
\textbf{Interaction} & Répond, relance, reformule, pose des questions. &
Avez-vous posé au moins trois questions ? \tabularnewline
\end{tabularx}
\end{tableau}

\begin{tableau}[Équivalence approximative avec le CECRL]
\begin{tabularx}{\linewidth}{@{}p{2.4cm}p{4.0cm}X@{}}
\entete{Total /25} & \entete{Niveau} & \entete{Ce que cela signifie} \tabularnewline
\hline
0 à 6 & A1 & Phrases isolées, gros efforts de l'interlocuteur. \tabularnewline
7 à 11 & A2 & Échanges simples sur des sujets connus, pauses fréquentes. \tabularnewline
12 à 16 & B1 & Conversation possible, avec des trous et des reformulations. \tabularnewline
17 à 21 & B2 & Conversation fluide et spontanée, nuances possibles. \tabularnewline
22 à 25 & C1 & Aisance complète, humour, abstraction, peu d'effort pour l'autre. \tabularnewline
\end{tabularx}
\end{tableau}

\section{Les certifications : faut-il passer un examen ?}
\begin{tableau}[Les quatre examens utiles]
\begin{tabularx}{\linewidth}{@{}p{2.6cm}p{2.0cm}p{4.0cm}X@{}}
\entete{Examen} & \entete{Score} & \entete{Ce qu'il évalue} & \entete{Pour qui} \tabularnewline
\hline
\textbf{TOEIC} (Listening \& Reading) & 10 à 990 & Compréhension orale et écrite en contexte professionnel. & Emploi, entreprises, recrutement. Le plus demandé en Afrique francophone. \tabularnewline
\textbf{Cambridge B2 First (FCE) / C1 Advanced (CAE)} & niveaux A2 à C2 & Les quatre compétences, dont l'oral en face à face. & Diplôme permanent, études, crédibilité académique. \tabularnewline
\textbf{IELTS} & 0 à 9 & Quatre compétences, très orienté études et immigration. & Université à l'étranger, visa, bourse. \tabularnewline
\textbf{TOEFL iBT} & 0 à 120 & Anglais académique nord-américain. & Universités américaines et canadiennes. \tabularnewline
\end{tabularx}
\end{tableau}

\begin{astuce}[Faut-il s'inscrire à un examen ?]
\textbf{Oui}, si vous avez besoin d'un papier pour un emploi, une bourse ou un visa ---
et l'échéance est alors un excellent moteur. \textbf{Peut-être}, si vous voulez une mesure
objective : un examen passé dans de vraies conditions dit la vérité sur votre niveau, ce
qu'aucune auto-évaluation ne fait. \textbf{Non}, si cela vous pousse à préparer des QCM
plutôt qu'à parler : le but reste la conversation, pas le score. Un bon compromis :
passez un examen blanc gratuit en ligne au mois 6 et au mois 12, et décidez ensuite.
\end{astuce}

\subsection{Stratégies d'examen (valables pour tous)}
\begin{enumerate}
  \item \textbf{Gestion du temps avant tout.} En QCM, ne passez jamais plus de 45 secondes
  sur une question : marquez, avancez, revenez. Une question non répondue vaut zéro ;
  une question devinée vaut parfois un point.
  \item \textbf{Écoute : lisez les questions avant l'audio.} Vous saurez quoi chercher.
  Les réponses apparaissent presque toujours dans l'ordre des questions.
  \item \textbf{Lecture : cherchez les mots-clés, ne traduisez pas tout.} Le \emph{skimming}
  (survol) puis le \emph{scanning} (recherche ciblée) font gagner 40\,\% du temps.
  \item \textbf{Méfiez-vous des synonymes.} Le texte dit \en{inexpensive}, la bonne réponse
  dit \en{cheap}. Ce sont presque toujours des reformulations, jamais les mêmes mots.
  \item \textbf{Écrit : respectez la consigne de longueur} et la structure attendue.
  Un texte de 180 mots demandé, écrit en 300, perd des points. Comptez vos mots à
  l'entraînement pour développer un sens visuel de la longueur.
  \item \textbf{Oral : parlez, même quand vous cherchez.} Utilisez les \emph{fillers} du
  chapitre 15. Le silence coûte plus cher que l'approximation.
  \item \textbf{Préparez-vous avec de vrais sujets officiels} (les organismes publient des
  exemples gratuits) --- pas avec des résumés trouvés au hasard.
\end{enumerate}

\section{Le bilan de fin d'année}
Le dernier jour de votre douzième mois, faites ce bilan écrit --- en anglais, évidemment.

\begin{tableau}[Le bilan des 365 jours]
\begin{tabularx}{\linewidth}{@{}p{7.6cm}X@{}}
\entete{Question} & \entete{Votre réponse} \tabularnewline
\hline
Combien de jours ai-je étudié ? (journal de progression) & \trou[4cm] \tabularnewline
Combien de mots sont actifs dans mon Anki ? & \trou[4cm] \tabularnewline
Combien d'heures de conversation réelle ai-je eues ? & \trou[4cm] \tabularnewline
Quel score au test de positionnement (sur 40), au mois 0 ? & \trou[4cm] \tabularnewline
Quel score aujourd'hui ? & \trou[4cm] \tabularnewline
Quel total à la grille orale (sur 25) ? & \trou[4cm] \tabularnewline
Quelles sont mes trois forces ? & \trou[4cm] \tabularnewline
Quelles sont mes trois faiblesses restantes ? & \trou[4cm] \tabularnewline
Quel est mon objectif pour l'année 2 ? & \trou[4cm] \tabularnewline
\end{tabularx}
\end{tableau}

\begin{mission}[Et après ? L'année 2]
Si vous avez atteint B2, votre année 2 ne doit plus ressembler à de l'étude : elle doit
ressembler à de la \textbf{vie en anglais}. Voici le programme :\\
\textbf{1.} Un film ou une série par semaine, sans aucun sous-titre.\\
\textbf{2.} Deux conversations d'une heure par semaine avec des natifs ou des personnes
plus avancées que vous.\\
\textbf{3.} Un livre par mois, en anglais, sur un sujet qui vous passionne vraiment.\\
\textbf{4.} Un domaine de spécialité à conquérir : votre métier, une discipline académique,
un loisir --- devenez capable de parler technique pendant 30 minutes.\\
\textbf{5.} Un accent secondaire à comprendre : indien, écossais, australien, nigérian.\\
\textbf{6.} De l'écrit ambitieux : un texte de 500 mots par semaine, corrigé par un natif.\\
\textbf{7.} Enseigner : expliquer l'anglais à quelqu'un de moins avancé est la méthode la
plus rapide pour consolider le C1.\\
\textbf{8.} Passer une certification si cela vous ouvre une porte.
\end{mission}

\begin{memo}[Le seul indicateur qui compte]
À la fin de l'année, la question n'est pas « quel est mon niveau CECRL ? ». C'est :
\textbf{est-ce que je peux parler trente minutes avec un anglophone, sans préparation,
sans peur, et sans repasser au français ?} Si la réponse est oui, vous avez réussi ---
quel que soit le chiffre. Écoutez votre enregistrement du jour 1, puis celui du jour 365.
C'est la seule preuve dont vous ayez besoin.
\end{memo}


% ------------------------------------------------------------------
%  ANNEXES
% ------------------------------------------------------------------
\appendix
\part*{Annexes \mdseries\small--- tout ce qu'il faut avoir sous la main}
\addcontentsline{toc}{part}{Annexes}
% =====================================================================
%  ANNEXE A — LES VERBES IRRÉGULIERS (150 verbes, classés par modèle)
% =====================================================================
\chapter{Les verbes irréguliers : 150 verbes, quatre modèles}

Il existe environ 200 verbes irréguliers en anglais courant ; les 150 qui suivent couvrent la quasi-totalité de ceux que vous rencontrerez. Appris en vrac, c'est une
corvée. Appris \textbf{par modèle}, c'est une semaine de travail : les verbes se regroupent
par sonorité, et votre mémoire auditive fait le reste.

\begin{memo}[Comment les apprendre --- la méthode des quatre semaines]
\textbf{Semaine 1 :} groupe 1 (formes identiques) + les 25 verbes du groupe 2 les plus
fréquents. \textbf{Semaine 2 :} le reste du groupe 2 (modèle \en{-t} : \en{felt, slept,
kept}). \textbf{Semaine 3 :} groupe 3, modèle \en{i} \ra \en{a} \ra \en{u}
(\en{sing, sang, sung}). \textbf{Semaine 4 :} groupe 3, modèles restants + révision
générale. Chaque jour : 15 verbes à voix haute, \textbf{dans une phrase} --- jamais en
listes nues. Exemple : \en{``begin, began, begun --- I began my work at eight.''}
\end{memo}

\section{Groupe 1 --- Les trois formes identiques (le plus facile)}
\begin{tableau}[Verbes invariables --- 20 verbes]
\footnotesize
\begin{longtable}{@{}p{3.0cm}p{3.0cm}p{3.0cm}p{5.2cm}@{}}
\entete{Base} & \entete{Passé simple} & \entete{Participe passé} & \entete{Sens} \tabularnewline
\endfirsthead
\entete{Base} & \entete{Passé} & \entete{Participe} & \entete{Sens} \tabularnewline
\endhead
\hline
\en{bet} & \en{bet} & \en{bet} & parier \tabularnewline
\en{burst} & \en{burst} & \en{burst} & éclater \tabularnewline
\en{cost} & \en{cost} & \en{cost} & coûter \tabularnewline
\en{cut} & \en{cut} & \en{cut} & couper \tabularnewline
\en{hit} & \en{hit} & \en{hit} & frapper \tabularnewline
\en{hurt} & \en{hurt} & \en{hurt} & blesser, faire mal \tabularnewline
\en{let} & \en{let} & \en{let} & laisser, permettre \tabularnewline
\en{put} & \en{put} & \en{put} & mettre \tabularnewline
\en{read} & \en{read} \pron{rèd} & \en{read} \pron{rèd} & lire (orthographe identique, son différent !) \tabularnewline
\en{set} & \en{set} & \en{set} & poser, régler \tabularnewline
\en{shut} & \en{shut} & \en{shut} & fermer \tabularnewline
\en{spread} & \en{spread} & \en{spread} & étaler, se répandre \tabularnewline
\en{cast} & \en{cast} & \en{cast} & lancer, distribuer (rôles) \tabularnewline
\en{fit} & \en{fit} & \en{fit} & aller (taille), adapter \tabularnewline
\en{quit} & \en{quit} & \en{quit} & quitter, arrêter \tabularnewline
\en{split} & \en{split} & \en{split} & fendre, diviser \tabularnewline
\en{upset} & \en{upset} & \en{upset} & bouleverser, renverser \tabularnewline
\en{broadcast} & \en{broadcast} & \en{broadcast} & diffuser \tabularnewline
\en{forecast} & \en{forecast} & \en{forecast} & prévoir (météo) \tabularnewline
\en{rid} & \en{rid} & \en{rid} & débarrasser (\en{get rid of}) \tabularnewline
\end{longtable}
\end{tableau}

\section{Groupe 2 --- Passé et participe identiques}

\subsection{Modèle 2A --- terminaison en \en{-t} (voyelle courte)}
\begin{tableau}[Modèle \en{-t} --- 30 verbes]
\footnotesize
\begin{longtable}{@{}p{3.0cm}p{3.0cm}p{3.0cm}p{5.2cm}@{}}
\entete{Base} & \entete{Passé} & \en{= Participe} & \entete{Sens} \tabularnewline
\endfirsthead
\entete{Base} & \entete{Passé / Participe} & \entete{Prononciation} & \entete{Sens} \tabularnewline
\endhead
\hline
\en{feel} & \en{felt} & \pron{fèlt} & sentir \tabularnewline
\en{keep} & \en{kept} & \pron{kèpt} & garder \tabularnewline
\en{sleep} & \en{slept} & \pron{slèpt} & dormir \tabularnewline
\en{sweep} & \en{swept} & \pron{souèpt} & balayer \tabularnewline
\en{weep} & \en{wept} & \pron{ouèpt} & pleurer \tabularnewline
\en{leave} & \en{left} & \pron{lèft} & quitter, laisser \tabularnewline
\en{lose} & \en{lost} & \pron{lost} & perdre \tabularnewline
\en{mean} & \en{meant} & \pron{mènt} & signifier, vouloir dire \tabularnewline
\en{meet} & \en{met} & \pron{mèt} & rencontrer \tabularnewline
\en{deal} & \en{dealt} & \pron{dèlt} & traiter, distribuer \tabularnewline
\en{dream} & \en{dreamt / dreamed} & \pron{drèmt} & rêver \tabularnewline
\en{learn} & \en{learnt / learned} & \pron{leunt} & apprendre \tabularnewline
\en{burn} & \en{burnt / burned} & \pron{beunt} & brûler \tabularnewline
\en{spell} & \en{spelt / spelled} & \pron{spèlt} & épeler \tabularnewline
\en{spill} & \en{spilt / spilled} & \pron{spilt} & renverser \tabularnewline
\en{spoil} & \en{spoilt / spoiled} & \pron{spoïlt} & gâcher, gâter \tabularnewline
\en{smell} & \en{smelt / smelled} & \pron{smèlt} & sentir \tabularnewline
\en{build} & \en{built} & \pron{bilt} & construire \tabularnewline
\en{spend} & \en{spent} & \pron{spènt} & dépenser, passer (temps) \tabularnewline
\en{send} & \en{sent} & \pron{sènt} & envoyer \tabularnewline
\en{lend} & \en{lent} & \pron{lènt} & prêter \tabularnewline
\en{bend} & \en{bent} & \pron{bènt} & plier \tabularnewline
\en{buy} & \en{bought} & \pron{boot} & acheter \tabularnewline
\en{bring} & \en{brought} & \pron{broot} & apporter \tabularnewline
\en{think} & \en{thought} & \pron{thoot} & penser \tabularnewline
\en{fight} & \en{fought} & \pron{foot} & combattre \tabularnewline
\en{seek} & \en{sought} & \pron{soot} & chercher (soutenu) \tabularnewline
\en{teach} & \en{taught} & \pron{toot} & enseigner \tabularnewline
\en{catch} & \en{caught} & \pron{koot} & attraper \tabularnewline
\en{sell} & \en{sold} & \pron{sôld} & vendre \tabularnewline
\end{longtable}
\end{tableau}

\subsection{Modèle 2B --- terminaison en \en{-d} ou \en{-n}}
\begin{tableau}[Modèle \en{-d/-n} --- 45 verbes]
\footnotesize
\begin{longtable}{@{}p{3.0cm}p{3.2cm}p{3.0cm}p{5.0cm}@{}}
\entete{Base} & \entete{Passé / Participe} & \entete{Prononciation} & \entete{Sens} \tabularnewline
\endfirsthead
\entete{Base} & \entete{Passé / Participe} & \entete{Prononciation} & \entete{Sens} \tabularnewline
\endhead
\hline
\en{find} & \en{found} & \pron{faound} & trouver \tabularnewline
\en{have} & \en{had} & \pron{had} & avoir \tabularnewline
\en{hear} & \en{heard} & \pron{heud} & entendre \tabularnewline
\en{hold} & \en{held} & \pron{hèld} & tenir \tabularnewline
\en{make} & \en{made} & \pron{mé-id} & faire, fabriquer \tabularnewline
\en{pay} & \en{paid} & \pron{pé-id} & payer \tabularnewline
\en{say} & \en{said} & \pron{sèd} & dire \tabularnewline
\en{sit} & \en{sat} & \pron{satt} & s'asseoir \tabularnewline
\en{stand} & \en{stood} & \pron{stoud} & se tenir debout \tabularnewline
\en{understand} & \en{understood} & \pron{en-deu-stoud} & comprendre \tabularnewline
\en{win} & \en{won} & \pron{ouone} & gagner \tabularnewline
\en{bind} & \en{bound} & \pron{baound} & lier \tabularnewline
\en{feed} & \en{fed} & \pron{fèd} & nourrir \tabularnewline
\en{flee} & \en{fled} & \pron{flèd} & fuir \tabularnewline
\en{lead} & \en{led} & \pron{lèd} & mener \tabularnewline
\en{shed} & \en{shed} & \pron{chèd} & verser, perdre \tabularnewline
\en{shoot} & \en{shot} & \pron{chot} & tirer \tabularnewline
\en{slide} & \en{slid} & \pron{slid} & glisser \tabularnewline
\en{stick} & \en{stuck} & \pron{steuk} & coller, planter \tabularnewline
\en{strike} & \en{struck} & \pron{streu-k} & frapper \tabularnewline
\en{dig} & \en{dug} & \pron{deug} & creuser \tabularnewline
\en{swing} & \en{swung} & \pron{souang} & balancer \tabularnewline
\en{hang} & \en{hung} & \pron{hang} & suspendre \tabularnewline
\en{spin} & \en{spun} & \pron{speun} & filer, tourner \tabularnewline
\en{sting} & \en{stung} & \pron{steung} & piquer \tabularnewline
\en{string} & \en{strung} & \pron{streung} & enfiler, tendre \tabularnewline
\en{cling} & \en{clung} & \pron{kleung} & s'accrocher \tabularnewline
\en{fling} & \en{flung} & \pron{fleung} & jeter \tabularnewline
\en{spring} & \en{sprang / sprung} & \pron{spraang} & jaillir \tabularnewline
\en{lay} & \en{laid} & \pron{lé-id} & poser, pondre \tabularnewline
\en{light} & \en{lit / lighted} & \pron{lit} & allumer \tabularnewline
\en{shine} & \en{shone / shined} & \pron{chôn} & briller \tabularnewline
\en{weave} & \en{wove} & \en{woven} & tisser \tabularnewline
\end{longtable}
\end{tableau}

\section{Groupe 3 --- Les trois formes différentes}

\subsection{Modèle 3A --- \en{i} \ra \en{a} \ra \en{u}}
\begin{tableau}[Le modèle le plus régulier des irréguliers]
\footnotesize
\begin{longtable}{@{}p{3.0cm}p{2.8cm}p{3.0cm}p{5.4cm}@{}}
\entete{Base} & \entete{Passé} & \entete{Participe} & \entete{Sens} \tabularnewline
\endfirsthead
\entete{Base} & \entete{Passé} & \entete{Participe} & \entete{Sens} \tabularnewline
\endhead
\hline
\en{begin} & \en{began} & \en{begun} & commencer \tabularnewline
\en{drink} & \en{drank} & \en{drunk} & boire \tabularnewline
\en{ring} & \en{rang} & \en{rung} & sonner \tabularnewline
\en{sing} & \en{sang} & \en{sung} & chanter \tabularnewline
\en{sink} & \en{sank} & \en{sunk} & couler \tabularnewline
\en{swim} & \en{swam} & \en{swum} & nager \tabularnewline
\en{shrink} & \en{shrank} & \en{shrunk} & rétrécir \tabularnewline
\en{stink} & \en{stank} & \en{stunk} & puer \tabularnewline
\en{wring} & \en{wrung} & \en{wrung} & tordre \tabularnewline
\end{longtable}
\end{tableau}

\subsection{Modèle 3B --- \en{i} \ra \en{o} \ra \en{i}}
\begin{tableau}[Voyelle centrale puis retour]
\footnotesize
\begin{longtable}{@{}p{3.0cm}p{2.8cm}p{3.0cm}p{5.4cm}@{}}
\entete{Base} & \entete{Passé} & \entete{Participe} & \entete{Sens} \tabularnewline
\endfirsthead
\entete{Base} & \entete{Passé} & \entete{Participe} & \entete{Sens} \tabularnewline
\endhead
\hline
\en{give} & \en{gave} & \en{given} & donner \tabularnewline
\en{drive} & \en{drove} & \en{driven} & conduire \tabularnewline
\en{ride} & \en{rode} & \en{ridden} & monter, chevaucher \tabularnewline
\en{rise} & \en{rose} & \en{risen} & se lever \tabularnewline
\en{write} & \en{wrote} & \en{written} & écrire \tabularnewline
\en{hide} & \en{hid} & \en{hidden} & cacher \tabularnewline
\en{bite} & \en{bit} & \en{bitten} & mordre \tabularnewline
\en{choose} & \en{chose} & \en{chosen} & choisir \tabularnewline
\en{freeze} & \en{froze} & \en{frozen} & geler \tabularnewline
\en{speak} & \en{spoke} & \en{spoken} & parler \tabularnewline
\en{break} & \en{broke} & \en{broken} & casser \tabularnewline
\en{steal} & \en{stole} & \en{stolen} & voler \tabularnewline
\en{wake} & \en{woke} & \en{woken} & (se) réveiller \tabularnewline
\en{forget} & \en{forgot} & \en{forgotten} & oublier \tabularnewline
\en{forgive} & \en{forgave} & \en{forgiven} & pardonner \tabularnewline
\en{get} & \en{got} & \en{gotten} (AmE) & obtenir \tabularnewline
\en{beget} & \en{begot} & \en{begotten} & engendrer \tabularnewline
\end{longtable}
\end{tableau}

\subsection{Modèle 3C --- \en{ew} / \en{own} / \en{grown}}
\begin{tableau}[Verbes en \en{-ow} et \en{-aw}]
\footnotesize
\begin{longtable}{@{}p{3.0cm}p{2.8cm}p{3.0cm}p{5.4cm}@{}}
\entete{Base} & \entete{Passé} & \entete{Participe} & \entete{Sens} \tabularnewline
\endfirsthead
\entete{Base} & \entete{Passé} & \entete{Participe} & \entete{Sens} \tabularnewline
\endhead
\hline
\en{know} & \en{knew} & \en{known} & savoir, connaître \tabularnewline
\en{grow} & \en{grew} & \en{grown} & grandir, cultiver \tabularnewline
\en{throw} & \en{threw} & \en{thrown} & jeter \tabularnewline
\en{blow} & \en{blew} & \en{blown} & souffler \tabularnewline
\en{fly} & \en{flew} & \en{flown} & voler \tabularnewline
\en{draw} & \en{drew} & \en{drawn} & dessiner, tirer \tabularnewline
\en{show} & \en{showed} & \en{shown} & montrer \tabularnewline
\en{sew} & \en{sewed} & \en{sewn} & coudre \tabularnewline
\en{withdraw} & \en{withdrew} & \en{withdrawn} & retirer \tabularnewline
\end{longtable}
\end{tableau}

\subsection{Modèle 3D --- Les verbes à connaître absolument (haute fréquence)}
\begin{tableau}[Les 50 irréguliers les plus utilisés]
\footnotesize
\begin{longtable}{@{}p{2.6cm}p{2.4cm}p{2.6cm}p{6.4cm}@{}}
\entete{Base} & \entete{Passé} & \entete{Participe} & \entete{Sens et exemple} \tabularnewline
\endfirsthead
\entete{Base} & \entete{Passé} & \entete{Participe} & \entete{Sens et exemple} \tabularnewline
\endhead
\hline
\en{be} & \en{was / were} & \en{been} & être --- \en{I have been there.} \tabularnewline
\en{become} & \en{became} & \en{become} & devenir --- \en{She has become a doctor.} \tabularnewline
\en{come} & \en{came} & \en{come} & venir --- \en{He has come home.} \tabularnewline
\en{do} & \en{did} & \en{done} & faire --- \en{Have you done it?} \tabularnewline
\en{eat} & \en{ate} & \en{eaten} & manger --- \en{I have eaten already.} \tabularnewline
\en{fall} & \en{fell} & \en{fallen} & tomber --- \en{He has fallen.} \tabularnewline
\en{go} & \en{went} & \en{gone} & aller --- \en{She has gone out.} \tabularnewline
\en{run} & \en{ran} & \en{run} & courir --- \en{I have run 5 km.} \tabularnewline
\en{see} & \en{saw} & \en{seen} & voir --- \en{I have seen that film.} \tabularnewline
\en{take} & \en{took} & \en{taken} & prendre --- \en{They have taken the bus.} \tabularnewline
\en{tell} & \en{told} & \en{told} & dire à --- \en{He told me the truth.} \tabularnewline
\en{wear} & \en{wore} & \en{worn} & porter --- \en{She has worn it twice.} \tabularnewline
\en{know} & \en{knew} & \en{known} & savoir --- \en{I have known him for years.} \tabularnewline
\en{think} & \en{thought} & \en{thought} & penser --- \en{I thought about it.} \tabularnewline
\en{say} & \en{said} & \en{said} & dire --- \en{He said he was tired.} \tabularnewline
\en{make} & \en{made} & \en{made} & faire --- \en{She made a mistake.} \tabularnewline
\en{find} & \en{found} & \en{found} & trouver --- \en{I found my keys.} \tabularnewline
\en{give} & \en{gave} & \en{given} & donner --- \en{They gave us food.} \tabularnewline
\en{get} & \en{got} & \en{got / gotten} & obtenir --- \en{I got a new job.} \tabularnewline
\en{leave} & \en{left} & \en{left} & partir --- \en{He left at six.} \tabularnewline
\en{buy} & \en{bought} & \en{bought} & acheter --- \en{We bought a car.} \tabularnewline
\en{bring} & \en{brought} & \en{brought} & apporter --- \en{She brought tea.} \tabularnewline
\en{teach} & \en{taught} & \en{taught} & enseigner --- \en{He taught me English.} \tabularnewline
\en{learn} & \en{learnt / learned} & \en{learnt / learned} & apprendre --- \en{I have learnt a lot.} \tabularnewline
\en{lose} & \en{lost} & \en{lost} & perdre --- \en{They lost the match.} \tabularnewline
\en{win} & \en{won} & \en{won} & gagner --- \en{We won two-nil.} \tabularnewline
\en{meet} & \en{met} & \en{met} & rencontrer --- \en{I met her in 2019.} \tabularnewline
\en{pay} & \en{paid} & \en{paid} & payer --- \en{I paid in cash.} \tabularnewline
\en{send} & \en{sent} & \en{sent} & envoyer --- \en{She sent an email.} \tabularnewline
\en{spend} & \en{spent} & \en{spent} & dépenser --- \en{We spent too much.} \tabularnewline
\en{build} & \en{built} & \en{built} & construire --- \en{They built a house.} \tabularnewline
\en{keep} & \en{kept} & \en{kept} & garder --- \en{He kept the change.} \tabularnewline
\en{feel} & \en{felt} & \en{felt} & sentir --- \en{I felt ill.} \tabularnewline
\en{sleep} & \en{slept} & \en{slept} & dormir --- \en{I slept eight hours.} \tabularnewline
\en{speak} & \en{spoke} & \en{spoken} & parler --- \en{She has spoken to him.} \tabularnewline
\en{break} & \en{broke} & \en{broken} & casser --- \en{It's broken.} \tabularnewline
\en{choose} & \en{chose} & \en{chosen} & choisir --- \en{They chose the blue one.} \tabularnewline
\en{write} & \en{wrote} & \en{written} & écrire --- \en{I have written to her.} \tabularnewline
\en{read} & \en{read} \pron{rèd} & \en{read} \pron{rèd} & lire --- \en{I have read it twice.} \tabularnewline
\en{understand} & \en{understood} & \en{understood} & comprendre --- \en{I understood everything.} \tabularnewline
\en{stand} & \en{stood} & \en{stood} & se tenir --- \en{He stood up.} \tabularnewline
\en{sit} & \en{sat} & \en{sat} & s'asseoir --- \en{We sat outside.} \tabularnewline
\en{hold} & \en{held} & \en{held} & tenir --- \en{She held my hand.} \tabularnewline
\en{hear} & \en{heard} & \en{heard} & entendre --- \en{I heard the news.} \tabularnewline
\en{hide} & \en{hid} & \en{hidden} & cacher --- \en{They hid the gift.} \tabularnewline
\en{begin} & \en{began} & \en{begun} & commencer --- \en{The film has begun.} \tabularnewline
\en{drink} & \en{drank} & \en{drunk} & boire --- \en{He drank too much.} \tabularnewline
\en{sing} & \en{sang} & \en{sung} & chanter --- \en{She has sung beautifully.} \tabularnewline
\en{swim} & \en{swam} & \en{swum} & nager --- \en{We swam in the sea.} \tabularnewline
\en{throw} & \en{threw} & \en{thrown} & jeter --- \en{He threw the ball.} \tabularnewline
\end{longtable}
\end{tableau}

\subsection{Modèle 3E --- Verbes irréguliers restants (à reconnaître)}
\begin{tableau}[Complément --- 45 verbes moins fréquents]
\footnotesize
\begin{longtable}{@{}p{2.8cm}p{2.8cm}p{2.8cm}p{5.4cm}@{}}
\entete{Base} & \entete{Passé} & \entete{Participe} & \entete{Sens} \tabularnewline
\endfirsthead
\entete{Base} & \entete{Passé} & \entete{Participe} & \entete{Sens} \tabularnewline
\endhead
\hline
\en{arise} & \en{arose} & \en{arisen} & surgir \tabularnewline
\en{awake} & \en{awoke} & \en{awoken} & (se) réveiller \tabularnewline
\en{bear} & \en{bore} & \en{borne / born} & supporter ; naître (\en{born}) \tabularnewline
\en{beat} & \en{beat} & \en{beaten} & battre \tabularnewline
\en{behold} & \en{beheld} & \en{beheld} & contempler \tabularnewline
\en{bleed} & \en{bled} & \en{bled} & saigner \tabularnewline
\en{breed} & \en{bred} & \en{bred} & élever (animaux) \tabularnewline
\en{creep} & \en{crept} & \en{crept} & ramper \tabularnewline
\en{forbid} & \en{forbade} & \en{forbidden} & interdire \tabularnewline
\en{foresee} & \en{foresaw} & \en{foreseen} & prévoir \tabularnewline
\en{grind} & \en{ground} & \en{ground} & moudre \tabularnewline
\en{kneel} & \en{knelt} & \en{knelt} & s'agenouiller \tabularnewline
\en{knit} & \en{knit / knitted} & \en{knit} & tricoter \tabularnewline
\en{leap} & \en{leapt / leaped} & \en{leapt} & bondir \tabularnewline
\en{lie} (s'allonger) & \en{lay} & \en{lain} & être couché \tabularnewline
\en{lie} (mentir) & \en{lied} & \en{lied} & mentir (régulier) \tabularnewline
\en{mistake} & \en{mistook} & \en{mistaken} & se tromper \tabularnewline
\en{overcome} & \en{overcame} & \en{overcome} & surmonter \tabularnewline
\en{prove} & \en{proved} & \en{proved / proven} & prouver \tabularnewline
\en{saw} & \en{sawed} & \en{sawn} & scier \tabularnewline
\en{shake} & \en{shook} & \en{shaken} & secouer \tabularnewline
\en{shave} & \en{shaved} & \en{shaven} & raser \tabularnewline
\en{shear} & \en{sheared} & \en{shorn} & tondre \tabularnewline
\en{slay} & \en{slew} & \en{slain} & tuer (littéraire) \tabularnewline
\en{sow} & \en{sowed} & \en{sown} & semer \tabularnewline
\en{stride} & \en{strode} & \en{stridden} & enjamber \tabularnewline
\en{swear} & \en{swore} & \en{sworn} & jurer \tabularnewline
\en{swell} & \en{swelled} & \en{swollen} & enfler \tabularnewline
\en{tear} & \en{tore} & \en{torn} & déchirer \tabularnewline
\en{undertake} & \en{undertook} & \en{undertaken} & entreprendre \tabularnewline
\end{longtable}
\end{tableau}

\begin{piege}[Les cinq paires qu'on confond tout le temps]
\en{lie} (s'allonger) : \en{lie --- lay --- lain}. \quad
\en{lie} (mentir) : \en{lie --- lied --- lied} (régulier).\\
\en{lay} (poser) : \en{lay --- laid --- laid}. \quad
\en{rise} (monter, intransitif) : \en{rise --- rose --- risen}. \quad
\en{raise} (lever, transitif) : \en{raise --- raised --- raised} (régulier).\\
Exemple : \en{Yesterday I \textbf{lay} on the bed} (je me suis allongé) ;
\en{I \textbf{laid} the book on the table} (j'ai posé le livre).
\end{piege}

\begin{exo}[A.1 --- Les trois formes, de mémoire]
Écrivez les trois formes des vingt verbes suivants :
\en{take, see, go, come, do, make, find, know, think, tell, buy, bring, teach, meet,
win, break, choose, write, speak, understand}.
\end{exo}
\begin{corrige}[A.1]
\en{take/took/taken ; see/saw/seen ; go/went/gone ; come/came/come ; do/did/done ;
make/made/made ; find/found/found ; know/knew/known ; think/thought/thought ;
tell/told/told ; buy/bought/bought ; bring/brought/brought ; teach/taught/taught ;
meet/met/met ; win/won/won ; break/broke/broken ; choose/chose/chosen ;
write/wrote/written ; speak/spoke/spoken ; understand/understood/understood}.
\end{corrige}

% =====================================================================
%  ANNEXE B — 600 PHRASES PRÊTES À L'EMPLOI
% =====================================================================
\chapter{600 phrases prêtes à l'emploi}

Voici votre \textbf{boîte à outils}. Ces phrases ne s'étudient pas comme du vocabulaire :
elles se \textbf{mémorisent comme des blocs}, entiers, avec leur son. Une phrase apprise
par c\oe ur est disponible instantanément à l'oral --- alors qu'une phrase construite mot à
mot coûte trois secondes de réflexion.

\begin{regle}[Comment utiliser ce répertoire]
\textbf{1.} Choisissez \textbf{dix phrases par jour}, pas plus.\\
\textbf{2.} Lisez-les à voix haute trois fois, en exagérant l'intonation.\\
\textbf{3.} Enregistrez-vous, écoutez, corrigez.\\
\textbf{4.} Utilisez chacune d'elles \textbf{au moins une fois dans la journée} --- même en
vous parlant à vous-même sous la douche.\\
\textbf{5.} Le week-end, test : cachez la colonne française, lisez la colonne anglaise.
Puis cachez l'anglais, dites la phrase en partant du français.\\
Une phrase est « acquise » quand vous pouvez la dire en moins de deux secondes, sans
traduire. Comptez environ \textbf{40 jours} pour maîtriser les 600.
\end{regle}

\section{1. Saluer et prendre congé}
\begin{tableau}[30 formules de salutation]
\footnotesize
\begin{longtable}{@{}p{8.0cm}p{6.6cm}@{}}
\entete{Anglais} & \entete{Français} \tabularnewline
\endfirsthead
\entete{Anglais} & \entete{Français} \tabularnewline
\endhead
\phr{Hello! How are you?}{Bonjour ! Comment allez-vous ?}
\phr{Hi there! Long time no see.}{Salut ! Ça fait longtemps.}
\phr{Good morning, everyone.}{Bonjour à tous.}
\phr{Good afternoon. How's it going?}{Bonjour. Comment ça va ?}
\phr{Good evening. Nice to see you again.}{Bonsoir. Content de vous revoir.}
\phr{Hey! What's up?}{Salut ! Quoi de neuf ?}
\phr{Hi! How have you been?}{Salut ! Comment vas-tu depuis la dernière fois ?}
\phr{Nice to meet you.}{Enchanté de faire votre connaissance.}
\phr{It's a pleasure to meet you.}{C'est un plaisir de vous rencontrer.}
\phr{I've heard so much about you.}{J'ai beaucoup entendu parler de vous.}
\phr{Sorry I'm late.}{Désolé, je suis en retard.}
\phr{Welcome! Come in.}{Bienvenue ! Entrez.}
\phr{How do you do?}{Comment allez-vous ? (très formel, première rencontre)}
\phr{Good to see you.}{Content de te voir.}
\phr{How's life treating you?}{Comment va la vie ?}
\phr{See you later!}{À plus tard !}
\phr{See you tomorrow morning.}{À demain matin.}
\phr{Take care!}{Prends soin de toi !}
\phr{Have a nice day!}{Bonne journée !}
\phr{Have a good weekend.}{Bon week-end.}
\phr{Have a safe trip.}{Bon voyage / bonne route.}
\phr{I'd better get going.}{Il faut que j'y aille.}
\phr{I must be off --- time flies.}{Je dois y aller --- le temps passe vite.}
\phr{It was lovely talking to you.}{C'était vraiment agréable de parler avec vous.}
\phr{Thanks for having me.}{Merci de m'avoir reçu.}
\phr{Give my regards to your family.}{Transmettez mes amitiés à votre famille.}
\phr{Say hello to Aïcha for me.}{Dis bonjour à Aïcha de ma part.}
\phr{Let's keep in touch.}{Restons en contact.}
\phr{Catch you later!}{On se capte plus tard ! (familier)}
\phr{Goodbye, and good luck with everything.}{Au revoir, et bonne chance pour tout.}
\end{longtable}
\end{tableau}

\section{2. Se présenter}
\begin{tableau}[28 phrases pour parler de soi]
\footnotesize
\begin{longtable}{@{}p{8.0cm}p{6.6cm}@{}}
\entete{Anglais} & \entete{Français} \tabularnewline
\endfirsthead
\entete{Anglais} & \entete{Français} \tabularnewline
\endhead
\phr{My name is Moussa --- but everyone calls me Mouss.}{Je m'appelle Moussa --- mais tout le monde m'appelle Mouss.}
\phr{I'm from Senegal, from a town called Thiès.}{Je viens du Sénégal, d'une ville qui s'appelle Thiès.}
\phr{I was born in 1995.}{Je suis né en 1995.}
\phr{I'm in my late twenties.}{J'ai la fin de la vingtaine.}
\phr{I live in Dakar with my family.}{J'habite à Dakar avec ma famille.}
\phr{I've been living here for five years.}{J'habite ici depuis cinq ans.}
\phr{I work as an accountant.}{Je travaille comme comptable.}
\phr{I've been working in this company since 2020.}{Je travaille dans cette entreprise depuis 2020.}
\phr{I'm a student at university.}{Je suis étudiant à l'université.}
\phr{I'm studying business administration.}{J'étudie la gestion.}
\phr{I'm currently looking for a job.}{Je cherche un emploi en ce moment.}
\phr{I'm self-employed.}{Je travaille à mon compte.}
\phr{In my free time, I play football and read.}{Pendant mon temps libre, je joue au football et je lis.}
\phr{I'm really into photography.}{Je suis passionné de photographie.}
\phr{I've been learning English for a year.}{J'apprends l'anglais depuis un an.}
\phr{I also speak French and Wolof.}{Je parle aussi français et wolof.}
\phr{I'm married and I have two children.}{Je suis marié et j'ai deux enfants.}
\phr{I'm single at the moment.}{Je suis célibataire en ce moment.}
\phr{I come from a big family --- we are seven.}{Je viens d'une grande famille --- nous sommes sept.}
\phr{My dream is to work abroad.}{Mon rêve est de travailler à l'étranger.}
\phr{I'd describe myself as calm and curious.}{Je me décrirais comme calme et curieux.}
\phr{People say I'm a good listener.}{On dit que je sais écouter.}
\phr{What about you? Tell me about yourself.}{Et vous ? Parlez-moi de vous.}
\phr{And what do you do for a living?}{Et que faites-vous dans la vie ?}
\phr{How long have you lived here?}{Depuis combien de temps habitez-vous ici ?}
\phr{Where are you from originally?}{D'où venez-vous à l'origine ?}
\phr{Sorry, how do you pronounce your name?}{Pardon, comment prononce-t-on votre nom ?}
\phr{Could you spell that for me, please?}{Pouvez-vous me l'épeler, s'il vous plaît ?}
\end{longtable}
\end{tableau}

\section{3. Répondre, réagir, montrer qu'on écoute}
\begin{tableau}[30 réactions --- indispensables pour ne pas paraître froid]
\footnotesize
\begin{longtable}{@{}p{8.0cm}p{6.6cm}@{}}
\entete{Anglais} & \entete{Français} \tabularnewline
\endfirsthead
\entete{Anglais} & \entete{Français} \tabularnewline
\endhead
\phr{I'm fine, thanks. And you?}{Je vais bien, merci. Et vous ?}
\phr{Not bad, thanks --- you?}{Pas mal, merci --- et toi ?}
\phr{Pretty good, actually.}{Plutôt bien, en fait.}
\phr{Can't complain.}{Je ne peux pas me plaindre.}
\phr{I'm exhausted --- it's been a long week.}{Je suis épuisé --- la semaine a été longue.}
\phr{Things are going well, thanks.}{Les choses se passent bien, merci.}
\phr{Really? Tell me more.}{Vraiment ? Racontez-moi.}
\phr{That's interesting.}{C'est intéressant.}
\phr{Wow, that's amazing!}{Waouh, c'est génial !}
\phr{That sounds great.}{Ça a l'air super.}
\phr{I'm so happy for you.}{Je suis vraiment content pour vous.}
\phr{Congratulations! Well done.}{Félicitations ! Bravo.}
\phr{That's terrible --- I'm so sorry.}{C'est terrible --- je suis désolé.}
\phr{Oh no, what happened?}{Oh non, qu'est-ce qui s'est passé ?}
\phr{That must have been difficult.}{Ça a dû être difficile.}
\phr{I understand how you feel.}{Je comprends ce que vous ressentez.}
\phr{Are you all right?}{Vous allez bien ?}
\phr{Take your time.}{Prenez votre temps.}
\phr{No worries.}{Pas de souci.}
\phr{It doesn't matter.}{Ça ne fait rien.}
\phr{Fair enough.}{C'est juste / je l'admets.}
\phr{I see what you mean.}{Je vois ce que vous voulez dire.}
\phr{That makes sense.}{C'est logique.}
\phr{Right, I get it.}{D'accord, je comprends.}
\phr{You're joking!}{Vous plaisantez !}
\phr{No way!}{Pas possible !}
\phr{That's unbelievable.}{C'est incroyable.}
\phr{I can imagine.}{J'imagine.}
\phr{Same here.}{Pareil pour moi.}
\phr{Me too --- exactly the same.}{Moi aussi --- exactement pareil.}
\end{longtable}
\end{tableau}

\section{4. Phrases de survie : quand vous ne comprenez pas}
\begin{tableau}[28 phrases --- les plus importantes de tout ce répertoire]
\footnotesize
\begin{longtable}{@{}p{8.0cm}p{6.6cm}@{}}
\entete{Anglais} & \entete{Français} \tabularnewline
\endfirsthead
\entete{Anglais} & \entete{Français} \tabularnewline
\endhead
\phr{Sorry, could you say that again, please?}{Pardon, pouvez-vous répéter, s'il vous plaît ?}
\phr{Could you speak more slowly, please?}{Pourriez-vous parler plus lentement ?}
\phr{I'm sorry, I didn't catch that.}{Désolé, je n'ai pas saisi.}
\phr{What does that mean?}{Qu'est-ce que ça veut dire ?}
\phr{How do you say ``vélo'' in English?}{Comment dit-on « vélo » en anglais ?}
\phr{What's the word for this?}{Comment appelle-t-on cela ?}
\phr{Could you spell it for me, please?}{Pouvez-vous me l'épeler, s'il vous plaît ?}
\phr{Sorry, could you write it down?}{Pardon, pourriez-vous l'écrire ?}
\phr{Do you mean\ldots{} ?}{Vous voulez dire\ldots{} ?}
\phr{So you're saying that\ldots{} ?}{Donc vous dites que\ldots{} ?}
\phr{Let me see if I understand correctly.}{Voyons si j'ai bien compris.}
\phr{I understand the words, but not the meaning.}{Je comprends les mots, mais pas le sens.}
\phr{My English isn't perfect yet, but I'm working on it.}{Mon anglais n'est pas encore parfait, mais je progresse.}
\phr{Could you use a simpler word, please?}{Pourriez-vous employer un mot plus simple ?}
\phr{I'm not sure I follow you.}{Je ne suis pas sûr de vous suivre.}
\phr{Sorry, I've lost the thread.}{Désolé, j'ai perdu le fil.}
\phr{Could you give me an example?}{Pourriez-vous me donner un exemple ?}
\phr{Is that a common expression?}{Est-ce une expression courante ?}
\phr{Would you mind explaining that again?}{Cela vous ennuierait-il d'expliquer encore ?}
\phr{I know the word but I can't pronounce it.}{Je connais le mot mais je n'arrive pas à le prononcer.}
\phr{Please stop me if I make a mistake.}{Arrêtez-moi si je fais une erreur.}
\phr{Can we speak a bit more slowly? I'm still learning.}{Pouvons-nous parler un peu plus lentement ? J'apprends encore.}
\phr{Just a second --- let me think.}{Un instant --- laissez-moi réfléchir.}
\phr{Give me a moment, please.}{Donnez-moi un instant, s'il vous plaît.}
\phr{What did you say the name was?}{Quel était le nom, déjà ?}
\phr{Sorry --- the line is breaking up.}{Pardon --- la ligne est mauvaise.}
\phr{Could you repeat the number, please?}{Pouvez-vous répéter le numéro ?}
\phr{Thanks for your patience with my English.}{Merci de votre patience avec mon anglais.}
\end{longtable}
\end{tableau}

\section{5. Donner et demander une opinion}
\begin{tableau}[30 phrases]
\footnotesize
\begin{longtable}{@{}p{8.0cm}p{6.6cm}@{}}
\entete{Anglais} & \entete{Français} \tabularnewline
\endfirsthead
\entete{Anglais} & \entete{Français} \tabularnewline
\endhead
\phr{In my opinion, it's not a good idea.}{À mon avis, ce n'est pas une bonne idée.}
\phr{As far as I'm concerned, it works fine.}{En ce qui me concerne, ça marche très bien.}
\phr{I think we should start earlier.}{Je pense que nous devrions commencer plus tôt.}
\phr{I honestly believe that\ldots{}}{Je crois sincèrement que\ldots{}}
\phr{If you ask me, it's too expensive.}{Si vous voulez mon avis, c'est trop cher.}
\phr{To my mind, the second option is better.}{Selon moi, la deuxième option est meilleure.}
\phr{I'd say it depends on the situation.}{Je dirais que ça dépend de la situation.}
\phr{I'm not convinced, to be honest.}{Je ne suis pas convaincu, honnêtement.}
\phr{I have mixed feelings about it.}{J'ai des sentiments mitigés à ce sujet.}
\phr{I strongly disagree.}{Je suis fortement en désaccord.}
\phr{I couldn't agree more.}{Je ne pourrais pas être plus d'accord.}
\phr{That's exactly what I think.}{C'est exactement ce que je pense.}
\phr{You've got a point there.}{Vous marquez un point / vous avez raison sur ce point.}
\phr{I see it differently.}{Je vois les choses différemment.}
\phr{That's one way of looking at it.}{C'est une façon de voir les choses.}
\phr{What do you think about it?}{Qu'en pensez-vous ?}
\phr{What's your opinion on this?}{Quelle est votre opinion là-dessus ?}
\phr{How do you feel about the new rules?}{Que pensez-vous des nouvelles règles ?}
\phr{Do you agree with me?}{Êtes-vous d'accord avec moi ?}
\phr{Would you say it's worth it?}{Diriez-vous que ça vaut le coup ?}
\phr{Let me put it this way.}{Laissez-moi le formuler ainsi.}
\phr{The way I see it, we have two options.}{À mes yeux, nous avons deux options.}
\phr{It seems to me that nothing has changed.}{Il me semble que rien n'a changé.}
\phr{I might be wrong, but\ldots{}}{J'ai peut-être tort, mais\ldots{}}
\phr{Correct me if I'm wrong.}{Corrigez-moi si je me trompe.}
\phr{There's no doubt that education matters.}{Il ne fait aucun doute que l'éducation compte.}
\phr{It's obvious that prices are rising.}{Il est évident que les prix augmentent.}
\phr{Some people think that\ldots{}, but I disagree.}{Certains pensent que\ldots{}, mais je ne suis pas d'accord.}
\phr{On the one hand\ldots{}, on the other hand\ldots{}}{D'un côté\ldots{}, de l'autre\ldots{}}
\phr{All in all, I'd do it again.}{Tout bien considéré, je le referais.}
\end{longtable}
\end{tableau}

\section{6. Demander, offrir, proposer}
\begin{tableau}[30 phrases]
\footnotesize
\begin{longtable}{@{}p{8.0cm}p{6.6cm}@{}}
\entete{Anglais} & \entete{Français} \tabularnewline
\endfirsthead
\entete{Anglais} & \entete{Français} \tabularnewline
\endhead
\phr{Could you help me with this, please?}{Pourriez-vous m'aider avec ceci ?}
\phr{Would you mind opening the window?}{Cela vous ennuierait-il d'ouvrir la fenêtre ?}
\phr{Could you do me a favour?}{Pourriez-vous me rendre un service ?}
\phr{Do you have a minute?}{Avez-vous une minute ?}
\phr{I was wondering if you could\ldots{}}{Je me demandais si vous pouviez\ldots{}}
\phr{Would it be possible to change the date?}{Serait-il possible de changer la date ?}
\phr{Can I ask you something?}{Puis-je vous demander quelque chose ?}
\phr{Where can I find a pharmacy?}{Où puis-je trouver une pharmacie ?}
\phr{How do I get to the train station?}{Comment aller à la gare ?}
\phr{Is it far from here?}{Est-ce loin d'ici ?}
\phr{Could you show me on the map?}{Pouvez-vous me montrer sur la carte ?}
\phr{What time does it open?}{À quelle heure ça ouvre ?}
\phr{Can I have the menu, please?}{Puis-je avoir la carte, s'il vous plaît ?}
\phr{May I use your phone?}{Puis-je utiliser votre téléphone ?}
\phr{Can I pay by card?}{Puis-je payer par carte ?}
\phr{Could I have the bill, please?}{Pourrais-je avoir l'addition ?}
\phr{Let me help you with your bags.}{Laissez-moi vous aider avec vos sacs.}
\phr{Would you like a cup of tea?}{Voulez-vous une tasse de thé ?}
\phr{Can I get you anything?}{Puis-je vous apporter quelque chose ?}
\phr{Shall I close the door?}{Voulez-vous que je ferme la porte ?}
\phr{Do you want me to drive?}{Voulez-vous que je conduise ?}
\phr{I'll take care of it.}{Je m'en occupe.}
\phr{Why don't we meet on Friday?}{Et si on se voyait vendredi ?}
\phr{How about going for a walk?}{Que diriez-vous d'aller se promener ?}
\phr{Would you like to join us?}{Voulez-vous vous joindre à nous ?}
\phr{I'd like to invite you to dinner.}{J'aimerais vous inviter à dîner.}
\phr{Are you free this evening?}{Êtes-vous libre ce soir ?}
\phr{That's very kind of you, thanks.}{C'est très aimable à vous, merci.}
\phr{Don't mention it.}{Je vous en prie / il n'y a pas de quoi.}
\phr{Any time!}{Quand vous voulez !}
\end{longtable}
\end{tableau}

\section{7. Accepter, refuser, hésiter}
\begin{tableau}[28 phrases --- refuser poliment est une compétence]
\footnotesize
\begin{longtable}{@{}p{8.0cm}p{6.6cm}@{}}
\entete{Anglais} & \entete{Français} \tabularnewline
\endfirsthead
\entete{Anglais} & \entete{Français} \tabularnewline
\endhead
\phr{Yes, of course.}{Oui, bien sûr.}
\phr{That sounds perfect.}{Ça me paraît parfait.}
\phr{I'd love to, thanks!}{Avec plaisir, merci !}
\phr{Sure, why not.}{Bien sûr, pourquoi pas.}
\phr{Count me in.}{Comptez sur moi / je suis partant.}
\phr{That works for me.}{Cela me convient.}
\phr{Let's do it.}{Faisons-le.}
\phr{I'm in.}{J'en suis.}
\phr{I'm afraid I can't.}{J'ai bien peur de ne pas pouvoir.}
\phr{I'd love to, but I already have plans.}{J'adorerais, mais j'ai déjà quelque chose de prévu.}
\phr{Sorry, that won't be possible.}{Désolé, ce ne sera pas possible.}
\phr{I'd rather not, if you don't mind.}{Je préférerais ne pas, si cela ne vous dérange pas.}
\phr{Maybe another time.}{Une autre fois, peut-être.}
\phr{I'm not sure --- let me think about it.}{Je ne suis pas sûr --- laissez-moi y réfléchir.}
\phr{I need to check my schedule first.}{Je dois d'abord vérifier mon emploi du temps.}
\phr{Could we postpone it to next week?}{Pourrions-nous reporter à la semaine prochaine ?}
\phr{I'll let you know as soon as I can.}{Je vous préviens dès que possible.}
\phr{It depends --- what time exactly?}{Ça dépend --- à quelle heure exactement ?}
\phr{To be honest, I'm not very keen on it.}{Honnêtement, ça ne m'enthousiasme pas.}
\phr{I'm not really into that kind of thing.}{Ce n'est pas vraiment mon truc.}
\phr{Thanks for asking, but no.}{Merci de proposer, mais non.}
\phr{I appreciate the offer, but I'll pass.}{J'apprécie l'offre, mais je passe mon tour.}
\phr{No, thank you --- I'm fine.}{Non, merci --- ça va.}
\phr{Sorry, I can't make it.}{Désolé, je ne peux pas venir.}
\phr{Something has come up.}{J'ai un imprévu.}
\phr{I have to cancel --- I'm really sorry.}{Je dois annuler --- je suis vraiment désolé.}
\phr{Let's say yes for now and confirm later.}{Disons oui pour l'instant et confirmons plus tard.}
\phr{If it's not too much trouble\ldots{}}{Si ce n'est pas trop de peine\ldots{}}
\end{longtable}
\end{tableau}

\section{8. Remercier et s'excuser}
\begin{tableau}[24 phrases]
\footnotesize
\begin{longtable}{@{}p{8.0cm}p{6.6cm}@{}}
\entete{Anglais} & \entete{Français} \tabularnewline
\endfirsthead
\entete{Anglais} & \entete{Français} \tabularnewline
\endhead
\phr{Thank you very much.}{Merci beaucoup.}
\phr{Thanks a lot for your help.}{Merci beaucoup pour votre aide.}
\phr{I really appreciate it.}{J'apprécie vraiment.}
\phr{That's very kind of you.}{C'est très gentil de votre part.}
\phr{I can't thank you enough.}{Je ne vous remercierai jamais assez.}
\phr{You've saved my life.}{Vous m'avez sauvé la vie.}
\phr{You're welcome.}{De rien / je vous en prie.}
\phr{My pleasure.}{Avec plaisir.}
\phr{Not at all.}{Pas du tout.}
\phr{No problem.}{Pas de problème.}
\phr{It was nothing.}{Ce n'était rien.}
\phr{I'm sorry --- it's my fault.}{Je suis désolé --- c'est ma faute.}
\phr{I apologise for the delay.}{Je vous présente mes excuses pour le retard.}
\phr{Sorry to bother you.}{Désolé de vous déranger.}
\phr{I shouldn't have said that.}{Je n'aurais pas dû dire cela.}
\phr{Please forgive me.}{Pardonnez-moi, je vous prie.}
\phr{It won't happen again.}{Cela ne se reproduira pas.}
\phr{I made a mistake and I take responsibility.}{J'ai fait une erreur et j'en assume la responsabilité.}
\phr{No need to apologise.}{Pas besoin de vous excuser.}
\phr{Forget about it --- it's fine.}{Oubliez --- ce n'est rien.}
\phr{I hope I didn't offend you.}{J'espère ne pas vous avoir offensé.}
\phr{That was rude of me --- sorry.}{C'était impoli de ma part --- pardon.}
\phr{Thanks for understanding.}{Merci de votre compréhension.}
\phr{I'm grateful for your support.}{Je suis reconnaissant de votre soutien.}
\end{longtable}
\end{tableau}

\section{9. Goûts, sentiments, émotions}
\begin{tableau}[26 phrases]
\footnotesize
\begin{longtable}{@{}p{8.0cm}p{6.6cm}@{}}
\entete{Anglais} & \entete{Français} \tabularnewline
\endfirsthead
\entete{Anglais} & \entete{Français} \tabularnewline
\endhead
\phr{I love this song.}{J'adore cette chanson.}
\phr{I'm really into African literature.}{Je suis passionné de littérature africaine.}
\phr{I'm crazy about football.}{Je suis fou de football.}
\phr{I quite like cooking.}{J'aime assez cuisiner.}
\phr{I don't really like horror films.}{Je n'aime pas trop les films d'horreur.}
\phr{I can't stand the noise.}{Je ne supporte pas le bruit.}
\phr{I hate being late.}{Je déteste être en retard.}
\phr{It's not really my thing.}{Ce n'est pas vraiment mon truc.}
\phr{I've never been interested in cars.}{Les voitures ne m'ont jamais intéressé.}
\phr{It grows on you.}{On s'y attache avec le temps.}
\phr{I'm so excited about this trip!}{Je suis tellement impatient pour ce voyage !}
\phr{I'm really looking forward to it.}{J'attends cela avec impatience.}
\phr{I'm a bit nervous, actually.}{Je suis un peu nerveux, en fait.}
\phr{I'm worried about the exam.}{Je m'inquiète pour l'examen.}
\phr{I'm afraid of flying.}{J'ai peur de l'avion.}
\phr{I feel so much better today.}{Je me sens beaucoup mieux aujourd'hui.}
\phr{I'm tired of waiting.}{J'en ai assez d'attendre.}
\phr{I'm fed up with this weather.}{J'en ai marre de ce temps.}
\phr{I'm proud of you.}{Je suis fier de toi.}
\phr{It made me laugh so much.}{Ça m'a beaucoup fait rire.}
\phr{I was really moved by the film.}{J'ai été très ému par le film.}
\phr{It's a bit disappointing.}{C'est un peu décevant.}
\phr{I'm relieved --- everything went well.}{Je suis soulagé --- tout s'est bien passé.}
\phr{I don't mind waiting.}{Cela ne me dérange pas d'attendre.}
\phr{It doesn't bother me at all.}{Cela ne me gêne pas du tout.}
\phr{I'm fine with that.}{Ça me va.}
\end{longtable}
\end{tableau}

\section{10. Le temps, les dates, les rendez-vous}
\begin{tableau}[28 phrases]
\footnotesize
\begin{longtable}{@{}p{8.0cm}p{6.6cm}@{}}
\entete{Anglais} & \entete{Français} \tabularnewline
\endfirsthead
\entete{Anglais} & \entete{Français} \tabularnewline
\endhead
\phr{What time is it?}{Quelle heure est-il ?}
\phr{It's half past nine.}{Il est neuf heures et demie.}
\phr{It's a quarter to eight.}{Il est huit heures moins le quart.}
\phr{What day is it today?}{Quel jour sommes-nous ?}
\phr{Today is Wednesday, the twelfth of March.}{Nous sommes mercredi douze mars.}
\phr{What's the date today?}{Quelle est la date d'aujourd'hui ?}
\phr{I was born on the fifth of May, 1995.}{Je suis né le 5 mai 1995.}
\phr{The meeting starts at ten sharp.}{La réunion commence à dix heures précises.}
\phr{Let's meet at eight in the evening.}{Rendez-vous à vingt heures.}
\phr{Can we move it to Friday?}{Pouvons-nous le déplacer à vendredi ?}
\phr{I'm busy on Monday afternoon.}{Je suis occupé lundi après-midi.}
\phr{Are you free next week?}{Êtes-vous libre la semaine prochaine ?}
\phr{I'll be there in ten minutes.}{J'arrive dans dix minutes.}
\phr{It takes about half an hour by bus.}{Cela prend environ une demi-heure en bus.}
\phr{I've been waiting for twenty minutes.}{J'attends depuis vingt minutes.}
\phr{We're running late.}{Nous sommes en retard.}
\phr{The deadline is next Monday.}{La date limite est lundi prochain.}
\phr{I do this twice a week.}{Je fais cela deux fois par semaine.}
\phr{Every other day, I go running.}{Un jour sur deux, je vais courir.}
\phr{I work from Monday to Friday.}{Je travaille du lundi au vendredi.}
\phr{It happens once in a while.}{Cela arrive de temps en temps.}
\phr{I haven't seen him for ages.}{Je ne l'ai pas vu depuis des lustres.}
\phr{A long time ago, I lived abroad.}{Il y a longtemps, j'habitais à l'étranger.}
\phr{These days, I study every evening.}{En ce moment, j'étudie tous les soirs.}
\phr{Soon, I hope.}{Bientôt, j'espère.}
\phr{Not yet --- maybe later.}{Pas encore --- peut-être plus tard.}
\phr{It's time to go.}{Il est temps de partir.}
\phr{The shop is open until nine.}{Le magasin est ouvert jusqu'à neuf heures.}
\end{longtable}
\end{tableau}

\section{11. Au téléphone et en visioconférence}
\begin{tableau}[26 phrases]
\footnotesize
\begin{longtable}{@{}p{8.0cm}p{6.6cm}@{}}
\entete{Anglais} & \entete{Français} \tabularnewline
\endfirsthead
\entete{Anglais} & \entete{Français} \tabularnewline
\endhead
\phr{Hello, this is Moussa speaking.}{Bonjour, Moussa à l'appareil.}
\phr{May I speak to Mrs Diop, please?}{Puis-je parler à Mme Diop, s'il vous plaît ?}
\phr{Who's calling, please?}{C'est de la part de qui ?}
\phr{Could you hold on a second?}{Pouvez-vous patienter un instant ?}
\phr{I'm afraid she's not available right now.}{Je crains qu'elle ne soit pas disponible actuellement.}
\phr{Can I take a message?}{Puis-je prendre un message ?}
\phr{Could you ask her to call me back?}{Pourriez-vous lui demander de me rappeler ?}
\phr{I'm calling about my order.}{J'appelle au sujet de ma commande.}
\phr{I'm returning your call.}{Je réponds à votre appel.}
\phr{Sorry, you've got the wrong number.}{Désolé, vous avez fait erreur de numéro.}
\phr{The line is bad --- can you hear me?}{La ligne est mauvaise --- vous m'entendez ?}
\phr{You're breaking up.}{Votre voix se coupe.}
\phr{Can you hear me now?}{Vous m'entendez maintenant ?}
\phr{I think we've been cut off.}{Je crois que nous avons été coupés.}
\phr{Let me put you through.}{Je vous passe quelqu'un.}
\phr{I'll call you back in five minutes.}{Je vous rappelle dans cinq minutes.}
\phr{Can everyone hear me all right?}{Tout le monde m'entend bien ?}
\phr{You're on mute.}{Vous êtes en mode muet.}
\phr{Could you share your screen, please?}{Pourriez-vous partager votre écran ?}
\phr{Sorry, go ahead --- you were speaking.}{Pardon, allez-y --- vous parliez.}
\phr{Let's schedule a call for Thursday.}{Programmons un appel pour jeudi.}
\phr{I'll send you the link by email.}{Je vous envoie le lien par courriel.}
\phr{My connection is unstable today.}{Ma connexion est instable aujourd'hui.}
\phr{Thanks for calling. Goodbye.}{Merci d'avoir appelé. Au revoir.}
\phr{I'll text you the details.}{Je vous envoie les détails par message.}
\phr{Speak to you soon!}{À bientôt au téléphone !}
\end{longtable}
\end{tableau}

\section{12. Au restaurant}
\begin{tableau}[26 phrases]
\footnotesize
\begin{longtable}{@{}p{8.0cm}p{6.6cm}@{}}
\entete{Anglais} & \entete{Français} \tabularnewline
\endfirsthead
\entete{Anglais} & \entete{Français} \tabularnewline
\endhead
\phr{A table for two, please.}{Une table pour deux, s'il vous plaît.}
\phr{Do you have a table by the window?}{Avez-vous une table près de la fenêtre ?}
\phr{What do you recommend?}{Que recommandez-vous ?}
\phr{What's in this dish?}{Qu'y a-t-il dans ce plat ?}
\phr{Is it spicy?}{Est-ce épicé ?}
\phr{I'm allergic to peanuts.}{Je suis allergique aux arachides.}
\phr{I don't eat meat --- I'm vegetarian.}{Je ne mange pas de viande --- je suis végétarien.}
\phr{Could I have this without onions?}{Pourrais-je l'avoir sans oignons ?}
\phr{I'll have the grilled fish, please.}{Je prendrai le poisson grillé, s'il vous plaît.}
\phr{Same for me.}{Pareil pour moi.}
\phr{Could we have some water, please?}{Pourrions-nous avoir de l'eau ?}
\phr{Still or sparkling?}{Plate ou gazeuse ?}
\phr{A glass of red wine, please.}{Un verre de vin rouge, s'il vous plaît.}
\phr{This is delicious!}{C'est délicieux !}
\phr{It's a bit cold, I'm afraid.}{C'est un peu froid, j'ai peur.}
\phr{I think there's been a mistake --- this isn't what I ordered.}{Je crois qu'il y a une erreur --- ce n'est pas ce que j'ai commandé.}
\phr{Could we have some more bread?}{Pourrions-nous avoir encore du pain ?}
\phr{Is service included?}{Le service est-il compris ?}
\phr{Could we have the bill, please?}{L'addition, s'il vous plaît.}
\phr{Can we pay separately?}{Pouvons-nous payer séparément ?}
\phr{Keep the change.}{Gardez la monnaie.}
\phr{The meal was excellent --- thank you.}{Le repas était excellent --- merci.}
\phr{Could I get this to take away?}{Puis-je l'emporter ?}
\phr{We're in a hurry --- is it quick to prepare?}{Nous sommes pressés --- est-ce rapide à préparer ?}
\phr{Do you accept credit cards?}{Acceptez-vous les cartes bancaires ?}
\phr{That was lovely --- we'll come back.}{C'était très agréable --- nous reviendrons.}
\end{longtable}
\end{tableau}

\section{13. À l'hôtel, en voyage, dans les transports}
\begin{tableau}[30 phrases]
\footnotesize
\begin{longtable}{@{}p{8.0cm}p{6.6cm}@{}}
\entete{Anglais} & \entete{Français} \tabularnewline
\endfirsthead
\entete{Anglais} & \entete{Français} \tabularnewline
\endhead
\phr{I have a reservation under the name Diop.}{J'ai une réservation au nom de Diop.}
\phr{I'd like to book a double room for two nights.}{Je voudrais réserver une chambre double pour deux nuits.}
\phr{Is breakfast included?}{Le petit-déjeuner est-il inclus ?}
\phr{What time is check-out?}{À quelle heure faut-il libérer la chambre ?}
\phr{Could I have a wake-up call at six?}{Pourrais-je avoir un réveil téléphonique à six heures ?}
\phr{The air conditioning isn't working.}{La climatisation ne fonctionne pas.}
\phr{Could I change rooms, please?}{Pourrais-je changer de chambre ?}
\phr{Is there Wi-Fi in the hotel?}{Y a-t-il du Wi-Fi dans l'hôtel ?}
\phr{What's the password, please?}{Quel est le mot de passe ?}
\phr{Could you call a taxi for me?}{Pourriez-vous m'appeler un taxi ?}
\phr{Where is the nearest bus stop?}{Où est l'arrêt de bus le plus proche ?}
\phr{Which platform does the train leave from?}{De quel quai part le train ?}
\phr{Is this seat taken?}{Cette place est-elle libre ?}
\phr{One ticket to London, please.}{Un billet pour Londres, s'il vous plaît.}
\phr{A return ticket --- how much is it?}{Un aller-retour --- combien cela coûte-t-il ?}
\phr{Is there a discount for students?}{Y a-t-il une réduction pour les étudiants ?}
\phr{What time does the flight leave?}{À quelle heure décolle le vol ?}
\phr{I'd like to check in, please.}{Je voudrais m'enregistrer, s'il vous plaît.}
\phr{Here's my passport and my ticket.}{Voici mon passeport et mon billet.}
\phr{I have one suitcase to check in.}{J'ai une valise à enregistrer.}
\phr{Is this hand luggage allowed?}{Ce bagage à main est-il autorisé ?}
\phr{My flight has been delayed.}{Mon vol a été retardé.}
\phr{I've missed my connection.}{J'ai raté ma correspondance.}
\phr{Where can I find the baggage claim?}{Où se trouve la livraison des bagages ?}
\phr{My suitcase hasn't arrived.}{Ma valise n'est pas arrivée.}
\phr{Could you help me, please? I'm lost.}{Pouvez-vous m'aider ? Je suis perdu.}
\phr{How long does it take to get there?}{Combien de temps faut-il pour y arriver ?}
\phr{Turn left at the second street.}{Tournez à gauche à la deuxième rue.}
\phr{Is it within walking distance?}{Est-ce accessible à pied ?}
\phr{Thanks --- you've been very helpful.}{Merci --- vous m'avez été très utile.}
\end{longtable}
\end{tableau}

\section{14. Achats et argent}
\begin{tableau}[24 phrases]
\footnotesize
\begin{longtable}{@{}p{8.0cm}p{6.6cm}@{}}
\entete{Anglais} & \entete{Français} \tabularnewline
\endfirsthead
\entete{Anglais} & \entete{Français} \tabularnewline
\endhead
\phr{I'm just looking, thanks.}{Je regarde seulement, merci.}
\phr{Do you have this in a larger size?}{Avez-vous ceci en taille plus grande ?}
\phr{Can I try it on?}{Puis-je l'essayer ?}
\phr{Where are the fitting rooms?}{Où sont les cabines d'essayage ?}
\phr{It doesn't fit --- it's too tight.}{Ça ne va pas --- c'est trop serré.}
\phr{How much is this?}{Combien coûte ceci ?}
\phr{That's a bit expensive for me.}{C'est un peu cher pour moi.}
\phr{Is there any discount?}{Y a-t-il une réduction ?}
\phr{Can you give me a better price?}{Pouvez-vous me faire un meilleur prix ?}
\phr{It's a good deal.}{C'est une bonne affaire.}
\phr{I'll take it.}{Je le prends.}
\phr{Could I have a receipt, please?}{Pourrais-je avoir un reçu ?}
\phr{Do you deliver?}{Livrez-vous ?}
\phr{Is it in stock?}{Est-ce en stock ?}
\phr{When will it be available again?}{Quand sera-t-il de nouveau disponible ?}
\phr{I'd like to return this --- it's faulty.}{Je voudrais retourner ceci --- il est défectueux.}
\phr{Can I get a refund or an exchange?}{Puis-je obtenir un remboursement ou un échange ?}
\phr{Here's my receipt.}{Voici mon reçu.}
\phr{Cash or card?}{En espèces ou par carte ?}
\phr{Could you break this note for me?}{Pouvez-vous me faire de la monnaie sur ce billet ?}
\phr{It costs fifty dollars.}{Cela coûte cinquante dollars.}
\phr{That's ten pounds each.}{C'est dix livres pièce.}
\phr{It's half price today.}{C'est à moitié prix aujourd'hui.}
\phr{Thanks for your help.}{Merci pour votre aide.}
\end{longtable}
\end{tableau}

\section{15. Santé et urgences}
\begin{tableau}[22 phrases]
\footnotesize
\begin{longtable}{@{}p{8.0cm}p{6.6cm}@{}}
\entete{Anglais} & \entete{Français} \tabularnewline
\endfirsthead
\entete{Anglais} & \entete{Français} \tabularnewline
\endhead
\phr{I don't feel well.}{Je ne me sens pas bien.}
\phr{I have a headache and a fever.}{J'ai mal à la tête et de la fièvre.}
\phr{I've had a sore throat for three days.}{J'ai mal à la gorge depuis trois jours.}
\phr{I've been coughing since Monday.}{Je tousse depuis lundi.}
\phr{My stomach hurts.}{J'ai mal au ventre.}
\phr{I think I've caught a cold.}{Je crois que j'ai attrapé un rhume.}
\phr{I'm allergic to penicillin.}{Je suis allergique à la pénicilline.}
\phr{I'm taking medication for my heart.}{Je prends un traitement pour le c\oe ur.}
\phr{I need to see a doctor.}{J'ai besoin de voir un médecin.}
\phr{Where is the nearest hospital?}{Où est l'hôpital le plus proche ?}
\phr{I'd like to make an appointment.}{Je voudrais prendre rendez-vous.}
\phr{Is it serious?}{Est-ce grave ?}
\phr{How long have you had these symptoms?}{Depuis combien de temps avez-vous ces symptômes ?}
\phr{Take this twice a day after meals.}{Prenez ceci deux fois par jour après les repas.}
\phr{You should rest for a few days.}{Vous devriez vous reposer quelques jours.}
\phr{Call an ambulance!}{Appelez une ambulance !}
\phr{Help! Somebody call the police!}{Au secours ! Que quelqu'un appelle la police !}
\phr{There's been an accident.}{Il y a eu un accident.}
\phr{I've lost my passport.}{J'ai perdu mon passeport.}
\phr{My bag has been stolen.}{On m'a volé mon sac.}
\phr{Where is the police station?}{Où est le commissariat ?}
\phr{I need to contact my embassy.}{Je dois contacter mon ambassade.}
\end{longtable}
\end{tableau}

\section{16. Au travail et en réunion}
\begin{tableau}[30 phrases]
\footnotesize
\begin{longtable}{@{}p{8.0cm}p{6.6cm}@{}}
\entete{Anglais} & \entete{Français} \tabularnewline
\endfirsthead
\entete{Anglais} & \entete{Français} \tabularnewline
\endhead
\phr{Let's get started, shall we?}{Commençons, voulez-vous ?}
\phr{The purpose of this meeting is to\ldots{}}{L'objet de cette réunion est de\ldots{}}
\phr{I'd like to go over three points.}{Je voudrais passer en revue trois points.}
\phr{First of all, let's look at the figures.}{Tout d'abord, regardons les chiffres.}
\phr{Could you give us an update on the project?}{Pouvez-vous nous faire un point sur le projet ?}
\phr{What's the current status?}{Où en sommes-nous ?}
\phr{I'd like to add something.}{J'aimerais ajouter quelque chose.}
\phr{Sorry to interrupt, but\ldots{}}{Pardon d'interrompre, mais\ldots{}}
\phr{If I may come back to what you said\ldots{}}{Si je puis revenir sur ce que vous avez dit\ldots{}}
\phr{Could you clarify that point?}{Pourriez-vous clarifier ce point ?}
\phr{What do you mean by ``efficiency''?}{Qu'entendez-vous par « efficacité » ?}
\phr{That's a good question --- let me think.}{C'est une bonne question --- laissez-moi réfléchir.}
\phr{I'll get back to you on that.}{Je reviens vers vous à ce sujet.}
\phr{Let's not go off topic.}{Ne nous écartons pas du sujet.}
\phr{Can we come back to the main point?}{Pouvons-nous revenir au point principal ?}
\phr{I'm afraid I don't have the figures with me.}{Je crains de ne pas avoir les chiffres avec moi.}
\phr{We're running out of time.}{Nous manquons de temps.}
\phr{Let's leave that for the next meeting.}{Laissons cela pour la prochaine réunion.}
\phr{To sum up, we agree on two points.}{Pour résumer, nous sommes d'accord sur deux points.}
\phr{So the next step is to contact the client.}{La prochaine étape est donc de contacter le client.}
\phr{Who's going to be in charge of this?}{Qui va s'en charger ?}
\phr{I'll take care of the report.}{Je m'occupe du rapport.}
\phr{Could you send me the minutes?}{Pourriez-vous m'envoyer le compte rendu ?}
\phr{The deadline is the end of the month.}{L'échéance est la fin du mois.}
\phr{We're behind schedule.}{Nous sommes en retard sur le planning.}
\phr{We've made good progress.}{Nous avons bien progressé.}
\phr{I'll keep you informed.}{Je vous tiendrai informé.}
\phr{Thanks everyone --- see you next week.}{Merci à tous --- à la semaine prochaine.}
\phr{I'm applying for the position of manager.}{Je postule au poste de responsable.}
\phr{Could I take a day off on Friday?}{Puis-je prendre un jour de congé vendredi ?}
\end{longtable}
\end{tableau}

\section{17. Apprendre une langue et parler de soi en tant qu'apprenant}
\begin{tableau}[24 phrases]
\footnotesize
\begin{longtable}{@{}p{8.0cm}p{6.6cm}@{}}
\entete{Anglais} & \entete{Français} \tabularnewline
\endfirsthead
\entete{Anglais} & \entete{Français} \tabularnewline
\endhead
\phr{I've been learning English for one year.}{J'apprends l'anglais depuis un an.}
\phr{My English is getting better every day.}{Mon anglais s'améliore chaque jour.}
\phr{I still translate in my head sometimes.}{Il m'arrive encore de traduire dans ma tête.}
\phr{Speaking is harder for me than writing.}{Parler est plus difficile pour moi qu'écrire.}
\phr{I understand much more than I can say.}{Je comprends bien plus que je ne peux dire.}
\phr{Could you correct my pronunciation?}{Pouvez-vous corriger ma prononciation ?}
\phr{Please correct me when I make mistakes.}{Corrigez-moi quand je fais des erreurs.}
\phr{Is this sentence natural?}{Cette phrase est-elle naturelle ?}
\phr{Would a native speaker say it like that?}{Un natif le dirait-il ainsi ?}
\phr{What's the difference between ``say'' and ``tell''?}{Quelle est la différence entre « say » et « tell » ?}
\phr{How do you pronounce this word?}{Comment prononce-t-on ce mot ?}
\phr{Is ``borrow'' the same as ``lend''?}{« Borrow » est-il identique à « lend » ?}
\phr{I need more practice with listening.}{J'ai besoin de plus d'entraînement à l'écoute.}
\phr{I watch one episode a day in English.}{Je regarde un épisode par jour en anglais.}
\phr{I keep a vocabulary notebook.}{Je tiens un carnet de vocabulaire.}
\phr{I review my flashcards every morning.}{Je révise mes cartes mémoire chaque matin.}
\phr{I'd like to find a conversation partner.}{J'aimerais trouver un partenaire de conversation.}
\phr{Do you know anyone who wants to practise?}{Connaissez-vous quelqu'un qui veut s'entraîner ?}
\phr{Could we speak English for half an hour?}{Pourrions-nous parler anglais pendant une demi-heure ?}
\phr{Let's switch to English, if that's all right.}{Passons à l'anglais, si cela vous convient.}
\phr{I'm preparing for the TOEIC exam.}{Je prépare l'examen TOEIC.}
\phr{My goal is to work in an international company.}{Mon objectif est de travailler dans une entreprise internationale.}
\phr{I'm proud of the progress I've made.}{Je suis fier des progrès que j'ai faits.}
\phr{One day, I'll think in English.}{Un jour, je penserai en anglais.}
\end{longtable}
\end{tableau}

\section{18. Conversation sociale (\emph{small talk})}
\begin{tableau}[28 phrases --- l'art de remplir les silences]
\footnotesize
\begin{longtable}{@{}p{8.0cm}p{6.6cm}@{}}
\entete{Anglais} & \entete{Français} \tabularnewline
\endfirsthead
\entete{Anglais} & \entete{Français} \tabularnewline
\endhead
\phr{Lovely weather today, isn't it?}{Il fait beau aujourd'hui, n'est-ce pas ?}
\phr{It's been raining for a week.}{Il pleut depuis une semaine.}
\phr{Is it always this hot in July?}{Fait-il toujours aussi chaud en juillet ?}
\phr{Do you live around here?}{Vous habitez dans le coin ?}
\phr{How long have you been in this city?}{Depuis combien de temps êtes-vous dans cette ville ?}
\phr{What do you do for a living?}{Que faites-vous dans la vie ?}
\phr{How did you get into that field?}{Comment en êtes-vous arrivé à ce métier ?}
\phr{Do you have any brothers or sisters?}{Avez-vous des frères ou s\oe urs ?}
\phr{What do you do in your free time?}{Que faites-vous pendant votre temps libre ?}
\phr{Have you seen any good films recently?}{Avez-vous vu de bons films récemment ?}
\phr{I'd recommend it --- it's excellent.}{Je vous le recommande --- c'est excellent.}
\phr{What kind of music do you like?}{Quel genre de musique aimez-vous ?}
\phr{Do you play any sports?}{Faites-vous du sport ?}
\phr{Did you support any team at the World Cup?}{Avez-vous soutenu une équipe à la Coupe du monde ?}
\phr{Have you ever been abroad?}{Êtes-vous déjà allé à l'étranger ?}
\phr{Where would you love to travel?}{Où aimeriez-vous voyager ?}
\phr{I've never been to Europe, but I'd love to go.}{Je ne suis jamais allé en Europe, mais j'adorerais y aller.}
\phr{What do you think of the food here?}{Que pensez-vous de la nourriture ici ?}
\phr{Have you tried the local dishes?}{Avez-vous goûté les plats locaux ?}
\phr{How was your weekend?}{Comment s'est passé votre week-end ?}
\phr{Nothing special --- I stayed at home.}{Rien de particulier --- je suis resté chez moi.}
\phr{What are your plans for the holidays?}{Quels sont vos projets pour les vacances ?}
\phr{That sounds like fun.}{Ça a l'air amusant.}
\phr{Anyway, it was nice chatting with you.}{En tout cas, c'était agréable de bavarder avec vous.}
\phr{Let's exchange numbers.}{Échangeons nos numéros.}
\phr{Are you on WhatsApp?}{Êtes-vous sur WhatsApp ?}
\phr{We should do this again.}{Il faudrait recommencer.}
\phr{Take care of yourself.}{Prenez soin de vous.}
\end{longtable}
\end{tableau}

\section{19. Raconter une histoire}
\begin{tableau}[24 phrases et connecteurs de récit]
\footnotesize
\begin{longtable}{@{}p{8.0cm}p{6.6cm}@{}}
\entete{Anglais} & \entete{Français} \tabularnewline
\endfirsthead
\entete{Anglais} & \entete{Français} \tabularnewline
\endhead
\phr{Let me tell you what happened.}{Laissez-moi vous raconter ce qui s'est passé.}
\phr{It all started last summer.}{Tout a commencé l'été dernier.}
\phr{One day, I was walking home when\ldots{}}{Un jour, je rentrais chez moi quand\ldots{}}
\phr{At first, I didn't understand.}{Au début, je n'ai pas compris.}
\phr{Then something strange happened.}{Puis quelque chose d'étrange s'est produit.}
\phr{Suddenly, the lights went out.}{Soudain, les lumières se sont éteintes.}
\phr{A few minutes later, he arrived.}{Quelques minutes plus tard, il est arrivé.}
\phr{Meanwhile, my sister was waiting outside.}{Pendant ce temps, ma s\oe ur attendait dehors.}
\phr{In the end, everything was fine.}{Finalement, tout s'est bien passé.}
\phr{It turned out that he knew her.}{Il s'est avéré qu'il la connaissait.}
\phr{To cut a long story short\ldots{}}{Pour faire court\ldots{}}
\phr{The funny thing is that\ldots{}}{Le plus drôle, c'est que\ldots{}}
\phr{Believe it or not, I won.}{Croyez-le ou non, j'ai gagné.}
\phr{I'll never forget that day.}{Je n'oublierai jamais ce jour-là.}
\phr{Looking back, it was a great experience.}{Avec le recul, ce fut une belle expérience.}
\phr{That's when I realised something important.}{C'est là que j'ai compris quelque chose d'important.}
\phr{If it hadn't been for him, I would have failed.}{Sans lui, j'aurais échoué.}
\phr{I used to do this every Sunday.}{Je faisais cela tous les dimanches.}
\phr{We were about to leave when it started.}{Nous étions sur le point de partir quand cela a commencé.}
\phr{He kept talking for an hour.}{Il n'a pas arrêté de parler pendant une heure.}
\phr{Nobody expected that.}{Personne ne s'y attendait.}
\phr{And that's the story.}{Et voilà l'histoire.}
\phr{What happened next?}{Que s'est-il passé ensuite ?}
\phr{How did you feel at that moment?}{Comment vous êtes-vous senti à ce moment-là ?}
\end{longtable}
\end{tableau}

\section{20. Parler de ses projets et de l'avenir}
\begin{tableau}[22 phrases]
\footnotesize
\begin{longtable}{@{}p{8.0cm}p{6.6cm}@{}}
\entete{Anglais} & \entete{Français} \tabularnewline
\endfirsthead
\entete{Anglais} & \entete{Français} \tabularnewline
\endhead
\phr{I'm going to start a new job in June.}{Je vais commencer un nouveau travail en juin.}
\phr{I'm planning to travel next year.}{Je prévois de voyager l'an prochain.}
\phr{I hope to improve my English quickly.}{J'espère améliorer mon anglais rapidement.}
\phr{My goal is to become fluent by December.}{Mon objectif est de devenir à l'aise d'ici décembre.}
\phr{I'd like to work for an international company.}{J'aimerais travailler pour une entreprise internationale.}
\phr{I'm thinking of going back to university.}{J'envisage de retourner à l'université.}
\phr{One day, I'd love to live abroad.}{Un jour, j'adorerais vivre à l'étranger.}
\phr{If everything goes well, I'll move in the spring.}{Si tout va bien, je déménagerai au printemps.}
\phr{I might change career --- nothing is decided.}{Je changerai peut-être de carrière --- rien n'est décidé.}
\phr{I've just started my own business.}{Je viens de créer ma propre entreprise.}
\phr{I'm saving money for a project.}{J'économise pour un projet.}
\phr{In five years, I see myself\ldots{}}{Dans cinq ans, je me vois\ldots{}}
\phr{I don't know yet, but I'll figure it out.}{Je ne sais pas encore, mais je trouverai.}
\phr{It depends on how things go.}{Cela dépend de comment les choses évoluent.}
\phr{I've made up my mind.}{J'ai pris ma décision.}
\phr{I'm still weighing up the options.}{Je pèse encore les options.}
\phr{Let's see how it goes.}{On verra comment ça se passe.}
\phr{I'm looking forward to the next step.}{J'attends la prochaine étape avec impatience.}
\phr{Whatever happens, I'll keep learning.}{Quoi qu'il arrive, je continuerai à apprendre.}
\phr{This is only the beginning.}{Ce n'est que le début.}
\phr{I believe hard work always pays off.}{Je crois que le travail finit toujours par payer.}
\phr{Thank you --- you've been very kind.}{Merci --- vous avez été très aimable.}
\end{longtable}
\end{tableau}

\section{21. Phrases pour l'écrit : e-mails et messages}
\begin{tableau}[32 formules d'e-mail --- du plus formel au plus courant]
\footnotesize
\begin{longtable}{@{}p{8.0cm}p{6.6cm}@{}}
\entete{Anglais} & \entete{Français} \tabularnewline
\endfirsthead
\entete{Anglais} & \entete{Français} \tabularnewline
\endhead
\phr{Dear Sir or Madam,}{Madame, Monsieur,}
\phr{Dear Mr Diop,}{Cher Monsieur Diop,}
\phr{Dear Aïcha,}{Chère Aïcha,}
\phr{Hi Sam --- hope you're doing well.}{Salut Sam --- j'espère que tu vas bien.}
\phr{I hope this email finds you well.}{J'espère que ce courriel vous trouvera en bonne santé.}
\phr{I'm writing to enquire about\ldots{}}{Je vous écris pour me renseigner au sujet de\ldots{}}
\phr{I'm writing in connection with your advertisement.}{Je vous écris au sujet de votre annonce.}
\phr{I'm writing to follow up on our conversation.}{Je vous écris pour faire suite à notre échange.}
\phr{Further to our phone call this morning\ldots{}}{Suite à notre appel téléphonique de ce matin\ldots{}}
\phr{With reference to your email of 3 May\ldots{}}{En référence à votre courriel du 3 mai\ldots{}}
\phr{I would like to request more information.}{Je souhaiterais obtenir davantage d'informations.}
\phr{Could you please send me the document?}{Pourriez-vous m'envoyer le document ?}
\phr{I would be grateful if you could reply by Friday.}{Je vous serais reconnaissant de bien vouloir répondre avant vendredi.}
\phr{Please find attached my CV.}{Veuillez trouver ci-joint mon CV.}
\phr{I've attached the report to this email.}{J'ai joint le rapport à ce courriel.}
\phr{Just a quick note to say thanks.}{Un petit mot pour vous remercier.}
\phr{Sorry for the late reply.}{Désolé pour la réponse tardive.}
\phr{Apologies for the delay in getting back to you.}{Mes excuses pour le délai de ma réponse.}
\phr{Thank you for your prompt reply.}{Merci pour votre réponse rapide.}
\phr{That's very clear --- thank you.}{C'est très clair --- merci.}
\phr{I'm afraid there's been a misunderstanding.}{Je crains qu'il y ait eu un malentendu.}
\phr{Could you clarify one point, please?}{Pourriez-vous clarifier un point, s'il vous plaît ?}
\phr{Unfortunately, we are unable to accept your request.}{Malheureusement, nous ne sommes pas en mesure d'accepter votre demande.}
\phr{We will do our best to solve the problem.}{Nous ferons de notre mieux pour résoudre le problème.}
\phr{Please don't hesitate to contact me.}{N'hésitez pas à me contacter.}
\phr{Let me know if you need anything else.}{Dites-moi si vous avez besoin d'autre chose.}
\phr{I look forward to hearing from you.}{Dans l'attente de votre réponse.}
\phr{Looking forward to meeting you.}{Au plaisir de vous rencontrer.}
\phr{Kind regards,}{Cordialement,}
\phr{Best regards,}{Bien cordialement,}
\phr{Yours sincerely,}{Veuillez agréer, Madame, Monsieur, mes salutations distinguées,}
\phr{See you soon! Talk later.}{À bientôt ! On se parle plus tard.}
\end{longtable}
\end{tableau}

\section{22. Tournures naturelles qui font « anglais courant »}
\begin{tableau}[30 expressions --- celles qui font la différence]
\footnotesize
\begin{longtable}{@{}p{8.0cm}p{6.6cm}@{}}
\entete{Anglais} & \entete{Français} \tabularnewline
\endfirsthead
\entete{Anglais} & \entete{Français} \tabularnewline
\endhead
\phr{It's up to you.}{C'est à toi de décider.}
\phr{That's not my cup of tea.}{Ce n'est pas mon truc.}
\phr{Let's play it by ear.}{Improvisons, on verra bien.}
\phr{It's not a big deal.}{Ce n'est pas grave.}
\phr{No biggie.}{Pas de souci. (familier)}
\phr{Long story short, I lost my job.}{Bref, j'ai perdu mon travail.}
\phr{Make yourself at home.}{Faites comme chez vous.}
\phr{It's on me --- my treat.}{C'est pour moi --- je vous invite.}
\phr{Let's split the bill.}{Partageons l'addition.}
\phr{I'm on my way.}{J'arrive / je suis en route.}
\phr{I'll be right there.}{J'arrive tout de suite.}
\phr{Give me a hand, will you?}{Donne-moi un coup de main, tu veux ?}
\phr{I'm swamped at work this week.}{Je suis débordé au travail cette semaine.}
\phr{Take it easy.}{Vas-y doucement / détends-toi.}
\phr{Don't worry about it.}{Ne t'inquiète pas pour ça.}
\phr{It doesn't make any sense.}{Cela n'a aucun sens.}
\phr{I've had enough of this.}{J'en ai assez de ça.}
\phr{Let's call it a day.}{Arrêtons-nous là pour aujourd'hui.}
\phr{It's worth a try.}{Cela vaut la peine d'essayer.}
\phr{There's no point in arguing.}{Il est inutile de se disputer.}
\phr{I can't afford it right now.}{Je ne peux pas me le permettre actuellement.}
\phr{It cost me an arm and a leg.}{Cela m'a coûté les yeux de la tête.}
\phr{That's a rip-off!}{C'est du vol !}
\phr{I'll get to it later.}{Je m'en occuperai plus tard.}
\phr{It slipped my mind.}{Cela m'est sorti de la tête.}
\phr{I'm about to leave.}{Je suis sur le point de partir.}
\phr{It turned out better than expected.}{Cela s'est avéré mieux que prévu.}
\phr{Everything went smoothly.}{Tout s'est déroulé sans accroc.}
\phr{I'm getting the hang of it.}{Je commence à piger.}
\phr{Practice makes perfect.}{C'est en forgeant qu'on devient forgeron.}
\end{longtable}
\end{tableau}

\begin{mission}[Annexe B --- le plan des 40 jours]
600 phrases, 15 par jour, 40 jours. Chaque matin : 15 phrases à voix haute (10 minutes).
Chaque soir : utilisez-en au moins trois dans la journée --- en conversation réelle, en
enregistrement, ou simplement en vous parlant. Le dimanche : test des 105 phrases de la
semaine, colonne française cachée. Au jour 40, vous possédez un répertoire qui couvre
90\,\% des situations de la vie quotidienne. C'est le plus grand raccourci vers la fluidité :
\textbf{parler en blocs tout faits pendant que votre grammaire, elle, continue de se
construire.}
\end{mission}

% =====================================================================
%  ANNEXE C — 14 FICHES DE RÉVISION (une page chacune)
% =====================================================================
\chapter{Quatorze fiches de révision à imprimer}

Ces fiches résument tout ce que le livre développe. Leur usage est simple :
\textbf{imprimez-les} (ou recopiez-les à la main, ce qui est encore plus efficace),
\textbf{affichez-les} là où vous passez du temps --- miroir de la salle de bain, porte du
frigo, mur du bureau, écran de veille du téléphone --- et \textbf{lisez-les à voix haute}
une minute par jour. En trois mois, vous connaîtrez la grammaire anglaise sans jamais
ouvrir un livre de grammaire.

\section*{Fiche 1 --- Les douze temps en une page}
\begin{tableau}[Les temps : forme, emploi, signal]
\footnotesize
\begin{longtable}{@{}p{3.4cm}p{4.6cm}p{3.2cm}p{3.0cm}@{}}
\entete{Temps} & \entete{Emploi} & \entete{Signal} & \entete{Exemple} \tabularnewline
\endfirsthead
\entete{Temps} & \entete{Emploi} & \entete{Signal} & \entete{Exemple} \tabularnewline
\endhead
\hline
\en{present simple} & habitude, vérité générale & \en{every day, always, usually} & \en{I work in Dakar.} \tabularnewline
\en{present continuous} & en cours maintenant, temporaire & \en{now, at the moment, these days} & \en{I'm working now.} \tabularnewline
\en{present perfect} & passé lié au présent, expérience & \en{already, yet, ever, since, for} & \en{I have finished.} \tabularnewline
\en{present perfect cont.} & activité commencée et encore en cours & \en{for, since, all day, how long} & \en{I have been waiting for an hour.} \tabularnewline
\en{past simple} & action terminée à un moment précis & \en{yesterday, last week, in 2019, ago} & \en{I worked yesterday.} \tabularnewline
\en{past continuous} & décor, action interrompue & \en{while, when, at 8 o'clock} & \en{I was working when he called.} \tabularnewline
\en{past perfect} & antériorité dans le passé & \en{before, after, already, by the time} & \en{I had worked before he arrived.} \tabularnewline
\en{past perfect cont.} & durée avant un moment passé & \en{for, since, all morning} & \en{I had been working for hours.} \tabularnewline
\en{future (will)} & décision immédiate, prédiction, promesse & \en{tomorrow, soon, I think} & \en{I'll work tomorrow.} \tabularnewline
\en{future (going to)} & intention, plan, évidence & \en{tonight, next week, look!} & \en{I'm going to work late.} \tabularnewline
\en{future continuous} & en cours à un moment futur & \en{this time tomorrow, at 9} & \en{I'll be working at nine.} \tabularnewline
\en{future perfect} & achevé avant un moment futur & \en{by, by the time, before} & \en{I'll have finished by Friday.} \tabularnewline
\end{longtable}
\end{tableau}

\begin{regle}[La règle d'or des temps]
Trois questions suffisent : \textbf{(1)} l'action est-elle finie et datée ? \ra passé simple.
\textbf{(2)} est-elle finie mais encore liée à maintenant ? \ra \en{present perfect}.
\textbf{(3)} est-elle en train de se dérouler ? \ra un temps en \en{-ing}.
\end{regle}

\section*{Fiche 2 --- Modaux et auxiliaires}
\begin{tableau}[Ce que chaque modal exprime]
\footnotesize
\begin{longtable}{@{}p{3.0cm}p{5.2cm}p{6.0cm}@{}}
\entete{Modal} & \entete{Sens} & \entete{Exemple} \tabularnewline
\endfirsthead
\entete{Modal} & \entete{Sens} & \entete{Exemple} \tabularnewline
\endhead
\hline
\en{can / could} & capacité, possibilité, demande & \en{I can swim. Could you help me?} \tabularnewline
\en{may / might} & permission, probabilité faible & \en{It might rain. May I come in?} \tabularnewline
\en{must} & obligation forte, déduction certaine & \en{You must stop. He must be tired.} \tabularnewline
\en{have to} & obligation extérieure, factuelle & \en{I have to work on Saturdays.} \tabularnewline
\en{should / ought to} & conseil, ce qui serait bien & \en{You should see a doctor.} \tabularnewline
\en{will / would} & futur, volonté, habitude passée, politesse & \en{I'll come. Would you like tea?} \tabularnewline
\en{shall} & proposition, suggestion (surtout \en{I/we}) & \en{Shall we go?} \tabularnewline
\en{used to} & habitude passée révolue & \en{I used to smoke.} \tabularnewline
\en{be used to} & être habitué à (+ \en{-ing}) & \en{I'm used to getting up early.} \tabularnewline
\en{needn't / don't have to} & absence d'obligation & \en{You needn't come.} \tabularnewline
\en{mustn't} & interdiction (\ko pas « pas obligé ») & \en{You mustn't smoke here.} \tabularnewline
\en{can't have / must have} & déduction sur le passé & \en{He can't have left. She must have forgotten.} \tabularnewline
\en{should have} & regret, reproche & \en{You should have told me.} \tabularnewline
\end{longtable}
\end{tableau}

\section*{Fiche 3 --- Poser toutes les questions}
\begin{tableau}[Les six structures interrogatives]
\footnotesize
\begin{longtable}{@{}p{4.4cm}p{9.8cm}@{}}
\entete{Type} & \entete{Structure et exemples} \tabularnewline
\endfirsthead
\entete{Type} & \entete{Structure et exemples} \tabularnewline
\endhead
\hline
\textbf{Oui/non} & Auxiliaire + sujet + verbe. \en{Do you like tea? Are you ready? Can she drive? Have they left? Did he call?} \tabularnewline
\textbf{Mot interrogatif} & Mot en \en{Wh-} + auxiliaire + sujet + verbe. \en{Where do you live? Why did she cry? How long have you been waiting?} \tabularnewline
\textbf{Sujet} (\en{who/what}) & Pas d'auxiliaire \en{do} : \en{Who called you? What happened? Who wants coffee?} \tabularnewline
\textbf{Question tag} & Affirmation + tag négatif (et inversement). \en{You're tired, aren't you? She doesn't smoke, does she? He can swim, can't he?} \tabularnewline
\textbf{Indirecte} & Formule polie + ordre affirmatif. \en{Could you tell me where the station is? Do you know what time it starts?} \tabularnewline
\textbf{Choix} & \en{or} entre deux options. \en{Do you want tea or coffee? Are you coming or staying?} \tabularnewline
\end{longtable}
\end{tableau}

\begin{piege}[Le \en{-s} de la question]
\en{*Does she works?} \ko \ra \en{Does she work?} \ok. L'auxiliaire \en{does} porte déjà la
marque de la 3\textsuperscript{e} personne : le verbe principal reste à la base.
\end{piege}

\section*{Fiche 4 --- Les articles}
\begin{tableau}[Quand utiliser \en{a/an}, \en{the}, ou rien]
\footnotesize
\begin{longtable}{@{}p{3.2cm}p{11.0cm}@{}}
\entete{Cas} & \entete{Règle et exemples} \tabularnewline
\endfirsthead
\entete{Cas} & \entete{Règle et exemples} \tabularnewline
\endhead
\hline
\en{a / an} & Singulier dénombrable, première mention, non identifié. \en{I saw \textbf{a} dog. She's \textbf{an} engineer.} \tabularnewline
\en{the} & Chose identifiée, unique, déjà mentionnée, ou avec superlatif. \en{\textbf{The} dog was black. \textbf{The} sun. \textbf{The} best film.} \tabularnewline
\textbf{rien} & Généralités au pluriel ou indénombrables. \en{\textbf{Dogs} are friendly. I love \textbf{music}. \textbf{Water} is essential.} \tabularnewline
\textbf{rien} & Repas, langues, sports, matières, jours, mois. \en{I had \textbf{breakfast}. She speaks \textbf{English}. I play \textbf{tennis}.} \tabularnewline
\textbf{rien} & Institutions dans leur fonction. \en{go to \textbf{school} / \textbf{bed} / \textbf{church} / \textbf{work} / \textbf{hospital} (BrE)} \tabularnewline
\en{the} & Pays pluriels ou composés, fleuves, mers, déserts, chaînes. \en{\textbf{the} USA, \textbf{the} Netherlands, \textbf{the} Nile, \textbf{the} Sahara, \textbf{the} Alps} \tabularnewline
\textbf{rien} & Pays, villes, continents, lacs, monts isolés. \en{Senegal, Dakar, Africa, Lake Victoria, Mount Kilimanjaro} \tabularnewline
\end{longtable}
\end{tableau}

\section*{Fiche 5 --- Prépositions de temps et de lieu}
\begin{tableau}[\en{at / on / in} --- la règle des trois échelles]
\footnotesize
\begin{longtable}{@{}p{2.2cm}p{4.4cm}p{7.6cm}@{}}
\entete{Prép.} & \entete{Échelle} & \entete{Exemples} \tabularnewline
\endfirsthead
\entete{Prép.} & \entete{Échelle} & \entete{Exemples} \tabularnewline
\endhead
\hline
\en{at} & point précis & \en{at 7 o'clock, at noon, at night, at the weekend (BrE), at Christmas, at the bus stop, at home, at work} \tabularnewline
\en{on} & jour, surface & \en{on Monday, on 3 May, on Monday morning, on my birthday, on the table, on the wall, on the bus, on the plane} \tabularnewline
\en{in} & période longue, contenant & \en{in May, in 2026, in the morning, in summer, in a week, in the box, in the car, in Dakar, in Africa} \tabularnewline
\en{by} & avant ou au plus tard & \en{by Friday, by 6 p.m.} \tabularnewline
\en{until / till} & jusque'à & \en{until midnight, till Friday} \tabularnewline
\en{for / since} & durée / point de départ & \en{for two hours, since 2019} \tabularnewline
\en{during / while} & pendant & \en{during the film, while I was sleeping} \tabularnewline
\en{from\ldots{} to} & de\ldots{} à & \en{from Monday to Friday, from 9 to 5} \tabularnewline
\en{between} & entre deux & \en{between 2019 and 2022, between the two buildings} \tabularnewline
\en{to / into / onto} & mouvement vers & \en{go to school, come into the room, climb onto the roof} \tabularnewline
\en{from / out of / off} & mouvement depuis & \en{from Dakar, out of the car, off the bus} \tabularnewline
\end{longtable}
\end{tableau}

\begin{astuce}[Le moyen mnémotechnique des échelles]
\en{at} = un point (heure, adresse précise). \en{on} = un jour ou une surface.
\en{in} = une boîte (mois, année, ville, pays, contenant). Plus l'échelle est grande,
plus la préposition est « large ».
\end{astuce}

\section*{Fiche 6 --- Comparatifs et superlatifs}
\begin{tableau}[Formation]
\footnotesize
\begin{longtable}{@{}p{3.6cm}p{4.0cm}p{6.6cm}@{}}
\entete{Adjectif} & \entete{Comparatif} & \entete{Superlatif} \tabularnewline
\endfirsthead
\entete{Adjectif} & \entete{Comparatif} & \entete{Superlatif} \tabularnewline
\endhead
\hline
1 syllabe & \en{+ -er} & \en{the + -est} --- \en{tall, taller, the tallest} \tabularnewline
1 syllabe en \en{-e} & \en{+ -r} & \en{the + -st} --- \en{large, larger, the largest} \tabularnewline
consonne + voyelle + consonne & double consonne & \en{big, bigger, the biggest} \tabularnewline
en \en{-y} & \en{-ier} & \en{the -iest} --- \en{happy, happier, the happiest} \tabularnewline
2 syllabes et plus & \en{more +} & \en{the most +} --- \en{careful, more careful, the most careful} \tabularnewline
irréguliers & \en{good/better/best} & \en{bad/worse/worst ; far/further/furthest ; little/less/least ; much/many/more/most} \tabularnewline
égalité & \en{as + adj + as} & \en{She's as tall as me.} \tabularnewline
inégalité & \en{not as + adj + as} & \en{It's not as cold as yesterday.} \tabularnewline
progression & \en{comparative + and + comparative} & \en{It's getting colder and colder.} \tabularnewline
proportion & \en{the + comparative, the + comparative} & \en{The more you practise, the better you get.} \tabularnewline
\end{longtable}
\end{tableau}

\section*{Fiche 7 --- Le discours rapporté}
\begin{tableau}[Le décalage des temps]
\footnotesize
\begin{longtable}{@{}p{5.4cm}p{4.4cm}p{4.4cm}@{}}
\entete{Discours direct} & \entete{Discours rapporté} & \entete{Exemple} \tabularnewline
\endfirsthead
\entete{Direct} & \entete{Rapporté} & \entete{Exemple} \tabularnewline
\endhead
\hline
\en{present simple} & \en{past simple} & \en{``I like it'' \ra She said she liked it.} \tabularnewline
\en{present continuous} & \en{past continuous} & \en{``I'm working'' \ra He said he was working.} \tabularnewline
\en{past simple} & \en{past perfect} & \en{``I saw it'' \ra She said she had seen it.} \tabularnewline
\en{present perfect} & \en{past perfect} & \en{``I've done it'' \ra He said he had done it.} \tabularnewline
\en{will} & \en{would} & \en{``I'll call'' \ra She said she would call.} \tabularnewline
\en{can} & \en{could} & \en{``I can help'' \ra He said he could help.} \tabularnewline
\en{must} & \en{had to} & \en{``I must go'' \ra She said she had to go.} \tabularnewline
\en{may} & \en{might} & \en{``It may rain'' \ra He said it might rain.} \tabularnewline
\en{today} & \en{that day} & --- \tabularnewline
\en{tomorrow} & \en{the next day / the following day} & --- \tabularnewline
\en{yesterday} & \en{the day before / the previous day} & --- \tabularnewline
\en{here} & \en{there} & --- \tabularnewline
\en{this / these} & \en{that / those} & --- \tabularnewline
\en{ago} & \en{before / earlier} & --- \tabularnewline
\end{longtable}
\end{tableau}

\begin{regle}[\en{say} ou \en{tell} ?]
\en{tell} + personne : \en{She \textbf{told me} she was tired.}\\
\en{say} + chose (+ \en{to} personne) : \en{She \textbf{said} she was tired.} ---
\en{She \textbf{said to me}\ldots{}} (plus rare).\\
Jamais \en{*She said me}. Jamais \en{*She told that\ldots{}} sans complément de personne.
\end{regle}

\section*{Fiche 8 --- Les conditionnels}
\begin{tableau}[Les cinq conditionnels]
\footnotesize
\begin{longtable}{@{}p{3.0cm}p{5.0cm}p{6.2cm}@{}}
\entete{Type} & \entete{Structure} & \entete{Emploi et exemple} \tabularnewline
\endfirsthead
\entete{Type} & \entete{Structure} & \entete{Emploi et exemple} \tabularnewline
\endhead
\hline
\textbf{Zéro} & \en{if + present, present} & Vérité générale. \en{If you heat water, it boils.} \tabularnewline
\textbf{1} & \en{if + present, will} & Futur réel ou probable. \en{If it rains, I'll stay home.} \tabularnewline
\textbf{2} & \en{if + past, would} & Imaginaire présent, hypothèse. \en{If I had time, I'd learn Spanish.} \tabularnewline
\textbf{3} & \en{if + past perfect, would have + p.p.} & Regret sur le passé. \en{If I had studied, I would have passed.} \tabularnewline
\textbf{Mixte} & \en{if + past perfect, would + base} & Cause passée, effet présent. \en{If I had taken that job, I'd be rich now.} \tabularnewline
\end{longtable}
\end{tableau}

\begin{piege}[Jamais \en{will} après \en{if}]
\en{*If it will rain\ldots{}} \ko \ra \en{If it rains\ldots{}} \ok.
Exception rare : \en{if} = « si jamais » dans une demande polie ---
\en{If you \textbf{will} wait a moment, I'll check.}
\end{piege}

\section*{Fiche 9 --- Le passif}
\begin{tableau}[Formation à tous les temps]
\footnotesize
\begin{longtable}{@{}p{4.0cm}p{4.4cm}p{5.8cm}@{}}
\entete{Temps} & \entete{Forme passive} & \entete{Exemple} \tabularnewline
\endfirsthead
\entete{Temps} & \entete{Forme passive} & \entete{Exemple} \tabularnewline
\endhead
\hline
\en{present simple} & \en{am/is/are + p.p.} & \en{English is spoken here.} \tabularnewline
\en{present continuous} & \en{am/is/are being + p.p.} & \en{The road is being repaired.} \tabularnewline
\en{past simple} & \en{was/were + p.p.} & \en{The bridge was built in 1990.} \tabularnewline
\en{past continuous} & \en{was/were being + p.p.} & \en{The car was being washed.} \tabularnewline
\en{present perfect} & \en{has/have been + p.p.} & \en{The work has been finished.} \tabularnewline
\en{past perfect} & \en{had been + p.p.} & \en{It had been stolen.} \tabularnewline
\en{future (will)} & \en{will be + p.p.} & \en{The results will be published.} \tabularnewline
\en{future (going to)} & \en{is/are going to be + p.p.} & \en{A school is going to be built.} \tabularnewline
\en{modals} & \en{modal + be + p.p.} & \en{It must be signed. It can be done.} \tabularnewline
\en{modals (past)} & \en{modal + have been + p.p.} & \en{It should have been checked.} \tabularnewline
\end{longtable}
\end{tableau}

\section*{Fiche 10 --- Quantifieurs}
\begin{tableau}[Dénombrable ou indénombrable ?]
\footnotesize
\begin{longtable}{@{}p{3.2cm}p{4.0cm}p{3.6cm}p{3.4cm}@{}}
\entete{Quantifieur} & \entete{S'emploie avec} & \entete{Sens} & \entete{Exemple} \tabularnewline
\endfirsthead
\entete{Quantifieur} & \entete{Avec} & \entete{Sens} & \entete{Exemple} \tabularnewline
\endhead
\hline
\en{many} & dénombrables (surtout nég./interr.) & beaucoup & \en{How many books?} \tabularnewline
\en{much} & indénombrables (surtout nég./interr.) & beaucoup & \en{How much time?} \tabularnewline
\en{a lot of / lots of} & les deux (affirmatif) & beaucoup & \en{a lot of friends / a lot of water} \tabularnewline
\en{(a) few} & dénombrables & quelques / peu & \en{a few days} \tabularnewline
\en{(a) little} & indénombrables & un peu & \en{a little sugar} \tabularnewline
\en{several} & dénombrables & plusieurs & \en{several times} \tabularnewline
\en{some} & les deux (affirmatif) & un peu de, des & \en{some money, some people} \tabularnewline
\en{any} & les deux (nég./interr.) & aucun, du & \en{I don't have any money.} \tabularnewline
\en{enough} & les deux & assez & \en{enough time, enough chairs} \tabularnewline
\en{all / most / no} & les deux & tout, la plupart, aucun & \en{all the students, most people} \tabularnewline
\en{each / every} & dénombrables singuliers & chaque & \en{every morning} \tabularnewline
\en{both / either / neither} & dénombrables (deux) & les deux / l'un ou l'autre / ni l'un ni l'autre & \en{both of us, neither option} \tabularnewline
\end{longtable}
\end{tableau}

\begin{piege}[\en{a few} $\neq$ \en{few}]
\en{a few friends} = quelques amis (positif, j'en ai). \en{few friends} = peu d'amis
(négatif, presque pas). Idem : \en{a little} (un peu) vs \en{little} (presque pas).
\end{piege}

\section*{Fiche 11 --- Prononciation : les dix sons à surveiller}
\begin{tableau}[Les sons qui trahissent un francophone]
\footnotesize
\begin{longtable}{@{}p{2.4cm}p{3.4cm}p{8.4cm}@{}}
\entete{Son} & \entete{Mot modèle} & \entete{Conseil} \tabularnewline
\endfirsthead
\entete{Son} & \entete{Mot modèle} & \entete{Conseil} \tabularnewline
\endhead
\hline
\en{h} aspiré & \en{house, he, hand} & Soufflez vraiment --- le \en{h} anglais existe. \ko \en{*ouse}. \tabularnewline
\en{th} sourd & \en{think, three, thing} & Langue entre les dents, sans voix. Pas \en{sink}. \tabularnewline
\en{th} sonore & \en{this, that, mother} & Même position, avec vibration. Pas \en{dis}. \tabularnewline
\en{w} & \en{we, water, window} & Lèvres arrondies comme pour \en{ou}, puis \en{wé}. Pas \en{v}. \tabularnewline
\en{v} & \en{very, vote} & Dents sur la lèvre --- le vrai \en{v} français. \tabularnewline
\en{r} anglais & \en{red, really} & Langue en arrière, jamais roulée ni grasseyée. \tabularnewline
\en{i} bref & \en{ship, sit, six} & Court et relâché --- différent de \en{sheep, seat}. \tabularnewline
\en{a} de \en{cat} & \en{cat, man, bad} & Bouche ouverte, entre \en{a} et \en{è}. \tabularnewline
\en{ng} final & \en{sing, thing} & \pron{si-ng} --- jamais de \en{g} explosif, jamais de \en{n} simple. \tabularnewline
\en{schwa} & \en{about, teacher} & Le son le plus fréquent de l'anglais : un \en{e} très court et neutre. \tabularnewline
\end{longtable}
\end{tableau}

\begin{regle}[Accent tonique --- trois réflexes]
\textbf{1.} Dans les mots de deux syllabes, l'accent tombe presque toujours sur la
\textbf{première} : \en{TAble, WAter, HAPpy, PREsent} (nom).\\
\textbf{2.} Les suffixes \en{-tion, -sion, -ic, -ity} attirent l'accent sur la syllabe
\textbf{précédente} : \en{naTION, deCIsion, ecoNOmic, uniVERsity}.\\
\textbf{3.} Verbe et nom de même orthographe changent d'accent :
\en{PREsent} (nom) / \en{preSENT} (verbe) ; \en{REcord} / \en{reCORD} ;
\en{OBject} / \en{obJECT}.
\end{regle}

\section*{Fiche 12 --- Orthographe : les 25 mots les plus ratés}
\begin{tableau}[À recopier trois fois chacun]
\footnotesize
\begin{multicols}{2}
\begin{itemize}\setlength{\itemsep}{1pt}
\item \en{because}
\item \en{necessary}
\item \en{separate}
\item \en{definitely}
\item \en{occurred}
\item \en{beginning}
\item \en{believe}
\item \en{receive}
\item \en{foreign}
\item \en{environment}
\item \en{government}
\item \en{different}
\item \en{business}
\item \en{friend}
\item \en{tomorrow}
\item \en{Wednesday}
\item \en{February}
\item \en{restaurant}
\item \en{beautiful}
\item \en{interesting}
\item \en{important}
\item \en{experience}
\item \en{knowledge}
\item \en{language}
\item \en{successful}
\end{itemize}
\end{multicols}
\end{tableau}

\section*{Fiche 13 --- Faux amis éclair (les 30 plus dangereux)}
\begin{tableau}[À relire avant chaque conversation]
\footnotesize
\begin{longtable}{@{}p{3.6cm}p{5.0cm}p{5.6cm}@{}}
\entete{Mot anglais} & \entete{Vrai sens} & \entete{Ce qu'on croit} \tabularnewline
\endfirsthead
\entete{Mot anglais} & \entete{Vrai sens} & \entete{Ce qu'on croit} \tabularnewline
\endhead
\hline
\en{actually} & en fait & actuellement \tabularnewline
\en{eventually} & finalement & éventuellement \tabularnewline
\en{library} & bibliothèque & librairie \tabularnewline
\en{shop} & magasin & shop (informatique) \tabularnewline
\en{coin} & pièce de monnaie & coin (angle) \tabularnewline
\en{chair} & chaise, présidence & chair (viande) \tabularnewline
\en{money} & argent (financier) & monnaie \tabularnewline
\en{currency} & monnaie (système) & courance \tabularnewline
\en{attend} & assister à & attendre \tabularnewline
\en{wait for} & attendre & --- \tabularnewline
\en{afford} & avoir les moyens & affermir \tabularnewline
\en{assist} & aider & assister à \tabularnewline
\en{injury} & blessure & injure \tabularnewline
\en{insult} & insulte & --- (attention : verbe \pron{inSEULT}) \tabularnewline
\en{sensible} & raisonnable & sensible \tabularnewline
\en{sensitive} & sensible & --- \tabularnewline
\en{comprehensive} & complet, exhaustif & compréhensif \tabularnewline
\en{understanding} & compréhensif & --- \tabularnewline
\en{large} & grand (dimensions) & large \tabularnewline
\en{wide} & large & --- \tabularnewline
\en{to travel} & voyager & travailler \tabularnewline
\en{to work} & travailler & --- \tabularnewline
\en{journey} & trajet, voyage & journée \tabularnewline
\en{day} & journée & --- \tabularnewline
\en{car} & voiture & car (autocar) \tabularnewline
\en{coach} & autocar, entraîneur & --- \tabularnewline
\en{to bless} & bénir & blesser \tabularnewline
\en{to hurt / to injure} & blesser & --- \tabularnewline
\en{preservative} & conservateur (alimentaire) & préservatif \tabularnewline
\en{condom} & préservatif & --- \tabularnewline
\end{longtable}
\end{tableau}

\section*{Fiche 14 --- La trousse de secours orale}
\begin{tableau}[Dix phrases qui sauvent une conversation]
\footnotesize
\begin{longtable}{@{}p{7.6cm}p{6.6cm}@{}}
\entete{Phrase} & \entete{À quoi elle sert} \tabularnewline
\endfirsthead
\entete{Phrase} & \entete{À quoi elle sert} \tabularnewline
\endhead
\hline
\en{Sorry, could you say that again?} & Gagner une seconde chance de comprendre. \tabularnewline
\en{Could you speak more slowly, please?} & Ralentir l'interlocuteur sans l'offenser. \tabularnewline
\en{How do you say ``\ldots{}'' in English?} & Combler un mot manquant. \tabularnewline
\en{You mean the thing you use to\ldots{} ?} & Décrire un objet dont on ignore le nom. \tabularnewline
\en{Let me put it another way.} & Reformuler au lieu de se bloquer. \tabularnewline
\en{What I'm trying to say is\ldots{}} & Rattraper une phrase mal partie. \tabularnewline
\en{Just a second --- let me think.} & Acheter du temps légitimement. \tabularnewline
\en{I'm not sure I follow you.} & Signaler une incompréhension sans s'excuser. \tabularnewline
\en{Could you give me an example?} & Transformer un mot abstrait en concret. \tabularnewline
\en{My English isn't perfect yet, but I'm working on it.} & Désamorcer la pression, en une phrase. \tabularnewline
\end{longtable}
\end{tableau}

\begin{mission}[Annexe C --- le rituel des fiches]
Imprimez les quatorze fiches. Collez-en \textbf{deux} sur le mur de votre espace de travail.
Chaque matin, lisez à voix haute les deux fiches affichées : une minute trente suffit.
Tous les sept jours, remplacez-les par les deux suivantes. En sept semaines, les quatorze
fiches sont passées sept fois devant vos yeux --- et la grammaire anglaise cesse d'être une
liste de règles pour devenir un automatisme.
\end{mission}

% =====================================================================
%  ANNEXE D — CONNECTEURS ET PRÉPOSITIONS
% =====================================================================
\chapter{Connecteurs, prépositions et particules}

Ce sont les deux éléments qui font passer un discours de « scolaire » à « naturel ».
Les francophones utilisent en anglais cinq ou six connecteurs (\en{and, but, because,
so, then}) --- alors qu'il en existe plus d'une centaine, chacun avec sa nuance.
Quant aux prépositions, elles ne se devinent jamais : elles s'apprennent \textbf{avec} le
mot, comme une seule unité.

\section{Les connecteurs par fonction}
\textbf{Registre :} \cible = formel/écrit ; \etoile = courant, oral et écrit ;
les autres sont neutres. Apprenez-en trois nouveaux par semaine.

\subsection{Ajouter une idée}
\begin{tableau}[Addition --- 14 connecteurs]
\footnotesize
\begin{longtable}{@{}p{4.2cm}p{4.2cm}p{5.8cm}@{}}
\entete{Connecteur} & \entete{Équivalent français} & \entete{Exemple} \tabularnewline
\endfirsthead
\entete{Connecteur} & \entete{Français} & \entete{Exemple} \tabularnewline
\endhead
\en{and} & et & \en{I work and I study.} \tabularnewline
\en{also} & aussi & \en{She also speaks Wolof.} \tabularnewline
\en{too / as well} & aussi (fin de phrase) & \en{I'd like to come, too.} \tabularnewline
\en{in addition} & de plus & \en{In addition, the price includes breakfast.} \tabularnewline
\en{furthermore} \cible & en outre & \en{Furthermore, the results were late.} \tabularnewline
\en{moreover} \cible & de plus, qui plus est & \en{Moreover, no one was informed.} \tabularnewline
\en{what is more} & qui plus est & \en{What is more, it costs nothing.} \tabularnewline
\en{besides} & d'ailleurs, en plus & \en{Besides, I don't have time.} \tabularnewline
\en{apart from that} & à part cela & \en{Apart from that, everything went well.} \tabularnewline
\en{not only\ldots{} but also} & non seulement\ldots{} mais aussi & \en{He's not only kind but also clever.} \tabularnewline
\en{both\ldots{} and} & à la fois\ldots{} et & \en{Both the film and the book are great.} \tabularnewline
\en{on top of that} & en plus de cela & \en{On top of that, it started to rain.} \tabularnewline
\en{to make matters worse} & pour ne rien arranger & \en{To make matters worse, I lost my phone.} \tabularnewline
\en{plus} (oral) & et puis & \en{Plus, I was really tired.} \tabularnewline
\end{longtable}
\end{tableau}

\subsection{Opposer et concéder}
\begin{tableau}[Contraste --- 16 connecteurs]
\footnotesize
\begin{longtable}{@{}p{4.2cm}p{4.2cm}p{5.8cm}@{}}
\entete{Connecteur} & \entete{Équivalent français} & \entete{Exemple} \tabularnewline
\endfirsthead
\entete{Connecteur} & \entete{Français} & \entete{Exemple} \tabularnewline
\endhead
\en{but} & mais & \en{I tried, but I failed.} \tabularnewline
\en{however} & cependant & \en{It's expensive. However, it's excellent.} \tabularnewline
\en{although / though} & bien que & \en{Although it rained, we went out.} \tabularnewline
\en{even though} & même si (réel) & \en{Even though he was ill, he came.} \tabularnewline
\en{even if} & même si (hypothétique) & \en{Even if you ask me, I won't agree.} \tabularnewline
\en{despite / in spite of} + nom & malgré & \en{Despite the rain, we played.} \tabularnewline
\en{nevertheless} \cible & néanmoins & \en{Nevertheless, the plan worked.} \tabularnewline
\en{nonetheless} \cible & malgré tout & \en{Nonetheless, doubts remain.} \tabularnewline
\en{on the other hand} & d'un autre côté & \en{On the other hand, it takes time.} \tabularnewline
\en{whereas / while} & alors que & \en{He's calm, whereas she's nervous.} \tabularnewline
\en{yet} & pourtant & \en{It's simple, yet effective.} \tabularnewline
\en{still} & tout de même & \en{Still, I wouldn't do it again.} \tabularnewline
\en{while it is true that} \cible & s'il est vrai que & \en{While it is true that prices rise, wages stagnate.} \tabularnewline
\en{admittedly} \cible & certes & \en{Admittedly, I was late.} \tabularnewline
\en{on the contrary} & au contraire & \en{On the contrary, I agree completely.} \tabularnewline
\en{unlike} & contrairement à & \en{Unlike his brother, he's shy.} \tabularnewline
\end{longtable}
\end{tableau}

\begin{piege}[\en{although} + proposition, \en{despite} + nom]
\en{Although \textbf{it was} raining} (sujet + verbe) \ok\\
\en{Despite \textbf{the rain}} (nom ou \en{-ing}) \ok\\
\en{*Despite of the rain} \ko --- jamais \en{of} après \en{despite}.\\
\en{*Although of the rain} \ko.
\end{piege}

\subsection{Expliquer la cause et la conséquence}
\begin{tableau}[Cause, conséquence, but --- 18 connecteurs]
\footnotesize
\begin{longtable}{@{}p{4.2cm}p{4.2cm}p{5.8cm}@{}}
\entete{Connecteur} & \entete{Équivalent français} & \entete{Exemple} \tabularnewline
\endfirsthead
\entete{Connecteur} & \entete{Français} & \entete{Exemple} \tabularnewline
\endhead
\en{because} & parce que & \en{I stayed because I was tired.} \tabularnewline
\en{because of} + nom & à cause de & \en{Because of the traffic, I was late.} \tabularnewline
\en{since} & puisque, comme & \en{Since you're here, let's start.} \tabularnewline
\en{as} & comme & \en{As it was late, we left.} \tabularnewline
\en{due to / owing to} \cible & en raison de & \en{Due to a technical problem\ldots{}} \tabularnewline
\en{thanks to} & grâce à & \en{Thanks to your help, I passed.} \tabularnewline
\en{the reason why is that} & la raison est que & \en{The reason why I called is that\ldots{}} \tabularnewline
\en{so} & donc & \en{It was late, so I went home.} \tabularnewline
\en{so\ldots{} that} & tellement\ldots{} que & \en{It was so hot that we left.} \tabularnewline
\en{such\ldots{} that} & un tel\ldots{} que & \en{It was such a good film that I saw it twice.} \tabularnewline
\en{therefore} \cible & par conséquent & \en{Therefore, we must act.} \tabularnewline
\en{thus} \cible & ainsi & \en{Thus, the problem was solved.} \tabularnewline
\en{consequently} \cible & par conséquent & \en{Consequently, sales fell.} \tabularnewline
\en{as a result} & en conséquence & \en{As a result, he lost his job.} \tabularnewline
\en{that's why} & c'est pourquoi & \en{That's why I'm learning English.} \tabularnewline
\en{in order to / so as to} & afin de & \en{In order to pass, I studied hard.} \tabularnewline
\en{so that} + proposition & pour que & \en{I spoke slowly so that they could understand.} \tabularnewline
\en{with the aim of} \cible & dans le but de & \en{With the aim of improving quality\ldots{}} \tabularnewline
\end{longtable}
\end{tableau}

\begin{piege}[\en{so that} $\neq$ \en{so}]
\en{It was late, \textbf{so} I went home} (conséquence, virgule avant).\\
\en{I left early \textbf{so that} I wouldn't be late} (but, suivi de sujet + verbe).\\
\en{*I left early so that not to be late} \ko \ra \en{so as not to be late} \ok.
\end{piege}

\subsection{Donner un exemple, illustrer, préciser}
\begin{tableau}[Exemple et reformulation --- 12 connecteurs]
\footnotesize
\begin{longtable}{@{}p{4.2cm}p{4.2cm}p{5.8cm}@{}}
\entete{Connecteur} & \entete{Équivalent français} & \entete{Exemple} \tabularnewline
\endfirsthead
\entete{Connecteur} & \entete{Français} & \entete{Exemple} \tabularnewline
\endhead
\en{for example / for instance} & par exemple & \en{For instance, Senegal and Ghana.} \tabularnewline
\en{such as} + nom & tel que & \en{Fruits such as mangoes and oranges.} \tabularnewline
\en{like} + nom (oral) & comme & \en{Countries like Mali and Niger.} \tabularnewline
\en{in particular} & en particulier & \en{In particular, the second point.} \tabularnewline
\en{especially} & surtout & \en{Especially in the morning.} \tabularnewline
\en{namely} \cible & à savoir & \en{Two cities, namely Dakar and Thiès.} \tabularnewline
\en{in other words} & autrement dit & \en{In other words, we have no choice.} \tabularnewline
\en{that is (to say)} & c'est-à-dire & \en{He's retired, that is to say, he doesn't work.} \tabularnewline
\en{what I mean is} & ce que je veux dire, c'est & \en{What I mean is that it's urgent.} \tabularnewline
\en{to put it simply} & pour le dire simplement & \en{To put it simply, it failed.} \tabularnewline
\en{take\ldots{}, for example} & prenons\ldots{}, par exemple & \en{Take my brother, for example.} \tabularnewline
\en{such as, among others} \cible & entre autres & \en{Several issues, among others.} \tabularnewline
\end{longtable}
\end{tableau}

\subsection{Organiser le temps et le récit}
\begin{tableau}[Chronologie --- 20 connecteurs]
\footnotesize
\begin{longtable}{@{}p{4.2cm}p{4.2cm}p{5.8cm}@{}}
\entete{Connecteur} & \entete{Équivalent français} & \entete{Exemple} \tabularnewline
\endfirsthead
\entete{Connecteur} & \entete{Français} & \entete{Exemple} \tabularnewline
\endhead
\en{first / first of all} & d'abord & \en{First of all, thank you all.} \tabularnewline
\en{to begin with} & pour commencer & \en{To begin with, let me explain.} \tabularnewline
\en{then} & ensuite & \en{Then we went home.} \tabularnewline
\en{next} & puis & \en{Next, add the flour.} \tabularnewline
\en{after that / afterwards} & après cela & \en{After that, nothing happened.} \tabularnewline
\en{later} & plus tard & \en{Two hours later, he arrived.} \tabularnewline
\en{soon} & bientôt & \en{See you soon.} \tabularnewline
\en{before} & avant & \en{Before leaving, check the door.} \tabularnewline
\en{while} & pendant que & \en{While I was cooking, he called.} \tabularnewline
\en{as} & au moment où & \en{As I opened the door, the phone rang.} \tabularnewline
\en{during} + nom & pendant & \en{During the meeting, nobody spoke.} \tabularnewline
\en{at the same time} & en même temps & \en{At the same time, she was studying.} \tabularnewline
\en{meanwhile / in the meantime} & pendant ce temps & \en{Meanwhile, the water was boiling.} \tabularnewline
\en{since then} & depuis lors & \en{Since then, I've changed my mind.} \tabularnewline
\en{until} & jusqu'à ce que & \en{Wait until I come back.} \tabularnewline
\en{as soon as} & dès que & \en{I'll call you as soon as I arrive.} \tabularnewline
\en{eventually} & finalement (après longtemps) & \en{Eventually, he agreed.} \tabularnewline
\en{in the end} & en fin de compte & \en{In the end, we stayed home.} \tabularnewline
\en{at last} & enfin (avec soulagement) & \en{At last, the bus came!} \tabularnewline
\en{finally} & enfin (dernier point) & \en{Finally, I'd like to thank my family.} \tabularnewline
\end{longtable}
\end{tableau}

\begin{piege}[\en{eventually}, \en{finally}, \en{at last}, \en{in the end}]
\en{Eventually} = après une longue attente ou beaucoup d'efforts (\textbf{pas} « éventuellement »).\\
\en{Finally} = introduit le dernier élément d'une liste.\\
\en{At last} = soulagement, souvent exclamatif.\\
\en{In the end} = bilan global, après hésitations.
\end{piege}

\subsection{Conclure, résumer, nuancer}
\begin{tableau}[Conclusion et nuance --- 14 connecteurs]
\footnotesize
\begin{longtable}{@{}p{4.6cm}p{4.0cm}p{5.6cm}@{}}
\entete{Connecteur} & \entete{Équivalent français} & \entete{Exemple} \tabularnewline
\endfirsthead
\entete{Connecteur} & \entete{Français} & \entete{Exemple} \tabularnewline
\endhead
\en{in conclusion} \cible & en conclusion & \en{In conclusion, the project is viable.} \tabularnewline
\en{to sum up} & pour résumer & \en{To sum up, we need more time.} \tabularnewline
\en{in short / in brief} & bref & \en{In short, it was a disaster.} \tabularnewline
\en{all in all} & tout bien considéré & \en{All in all, it was a good year.} \tabularnewline
\en{on the whole} & dans l'ensemble & \en{On the whole, I'm satisfied.} \tabularnewline
\en{overall} \cible & globalement & \en{Overall, sales increased by 8\%.} \tabularnewline
\en{to conclude} \cible & pour conclure & \en{To conclude, I thank you all.} \tabularnewline
\en{basically} (oral) & en gros & \en{Basically, we agree.} \tabularnewline
\en{in a nutshell} (oral) & en deux mots & \en{In a nutshell, it's too expensive.} \tabularnewline
\en{generally speaking} & de manière générale & \en{Generally speaking, people are kind.} \tabularnewline
\en{to some extent} & dans une certaine mesure & \en{To some extent, you're right.} \tabularnewline
\en{more or less} & plus ou moins & \en{It's more or less finished.} \tabularnewline
\en{rather than} & plutôt que & \en{Rather than complaining, let's act.} \tabularnewline
\en{instead of} + \en{-ing} & au lieu de & \en{Instead of waiting, call them.} \tabularnewline
\end{longtable}
\end{tableau}

\subsection{Structurer un essai ou une présentation}
\begin{tableau}[Le squelette d'un discours de cinq minutes]
\footnotesize
\begin{longtable}{@{}p{3.4cm}p{10.8cm}@{}}
\entete{Étape} & \entete{Formules prêtes à l'emploi} \tabularnewline
\endfirsthead
\entete{Étape} & \entete{Formules} \tabularnewline
\endhead
\textbf{Introduction} & \en{Today I'd like to talk about\ldots{}} / \en{The purpose of this presentation is to\ldots{}} / \en{I've divided my talk into three parts.} \tabularnewline
\textbf{Annonce du plan} & \en{First, I'll look at\ldots{} Then I'll move on to\ldots{} Finally, I'll conclude with\ldots{}} \tabularnewline
\textbf{Transition} & \en{Let's move on to the second point.} / \en{This brings me to\ldots{}} / \en{Turning now to\ldots{}} \tabularnewline
\textbf{Développement} & \en{As I mentioned earlier\ldots{}} / \en{To give you an example\ldots{}} / \en{The key point here is that\ldots{}} \tabularnewline
\textbf{Nuance} & \en{Of course, there are other views.} / \en{It's worth noting that\ldots{}} / \en{However, we must remember\ldots{}} \tabularnewline
\textbf{Conclusion} & \en{To sum up\ldots{}} / \en{The main takeaway is that\ldots{}} / \en{If you remember one thing, let it be this\ldots{}} \tabularnewline
\textbf{Questions} & \en{Thank you for listening. I'd be happy to take your questions.} / \en{Does that make sense?} / \en{Would anyone like to add something?} \tabularnewline
\end{longtable}
\end{tableau}

\section{Verbe + préposition : 100 combinaisons à mémoriser en bloc}
Ne traduisez jamais la préposition --- apprenez le verbe \textbf{avec} elle, comme un seul mot.

\begin{tableau}[Verbes suivis de \en{at}]
\footnotesize
\begin{longtable}{@{}p{4.4cm}p{9.8cm}@{}}
\entete{Verbe} & \entete{Sens et exemple} \tabularnewline
\endfirsthead
\entete{Verbe} & \entete{Sens et exemple} \tabularnewline
\endhead
\en{look at} & regarder --- \en{Look at this photo.} \tabularnewline
\en{listen to} & écouter --- \en{Listen to me.} (attention : \en{to}, pas \en{at}) \tabularnewline
\en{laugh at} & rire de --- \en{Don't laugh at him.} \tabularnewline
\en{smile at} & sourire à --- \en{She smiled at the baby.} \tabularnewline
\en{shout at} & crier après --- \en{He shouted at the driver.} \tabularnewline
\en{stare at} & fixer du regard --- \en{Stop staring at me.} \tabularnewline
\en{point at / to} & pointer vers --- \en{He pointed at the map.} \tabularnewline
\en{aim at} & viser --- \en{We're aiming at a 20\% increase.} \tabularnewline
\en{arrive at} (lieu précis) & arriver à --- \en{We arrived at the airport.} \tabularnewline
\en{be good at} & être bon en --- \en{She's good at maths.} \tabularnewline
\en{be bad at} & être mauvais en --- \en{I'm bad at drawing.} \tabularnewline
\en{be surprised at} & être surpris de --- \en{I was surprised at his answer.} \tabularnewline
\end{longtable}
\end{tableau}

\begin{tableau}[Verbes suivis de \en{for}]
\footnotesize
\begin{longtable}{@{}p{4.4cm}p{9.8cm}@{}}
\entete{Verbe} & \entete{Sens et exemple} \tabularnewline
\endfirsthead
\entete{Verbe} & \entete{Sens et exemple} \tabularnewline
\endhead
\en{look for} & chercher --- \en{I'm looking for my keys.} \tabularnewline
\en{wait for} & attendre --- \en{Wait for me!} \tabularnewline
\en{ask for} & demander --- \en{He asked for help.} \tabularnewline
\en{pay for} & payer --- \en{I'll pay for the tickets.} \tabularnewline
\en{search for} & rechercher --- \en{They searched for hours.} \tabularnewline
\en{apologise for} & s'excuser de --- \en{I apologise for the delay.} \tabularnewline
\en{thank (sb) for} & remercier de --- \en{Thanks for coming.} \tabularnewline
\en{blame (sb) for} & reprocher à --- \en{Don't blame me for this.} \tabularnewline
\en{hope for} & espérer --- \en{We hope for a better future.} \tabularnewline
\en{care for} & s'occuper de / aimer --- \en{She cares for her mother.} \tabularnewline
\en{apply for} & postuler à --- \en{I applied for the job.} \tabularnewline
\en{prepare for} & se préparer à --- \en{We're preparing for the exam.} \tabularnewline
\end{longtable}
\end{tableau}

\begin{tableau}[Verbes suivis de \en{to}]
\footnotesize
\begin{longtable}{@{}p{4.4cm}p{9.8cm}@{}}
\entete{Verbe} & \entete{Sens et exemple} \tabularnewline
\endfirsthead
\entete{Verbe} & \entete{Sens et exemple} \tabularnewline
\endhead
\en{listen to} & écouter --- \en{I listen to music every day.} \tabularnewline
\en{speak to / talk to} & parler à --- \en{Can I speak to you?} \tabularnewline
\en{write to} & écrire à --- \en{Write to me soon.} \tabularnewline
\en{reply to} & répondre à --- \en{He never replied to my email.} \tabularnewline
\en{explain to} & expliquer à --- \en{Explain it to me slowly.} \tabularnewline
\en{belong to} & appartenir à --- \en{This bag belongs to her.} \tabularnewline
\en{happen to} & arriver à --- \en{What happened to you?} \tabularnewline
\en{agree with / agree to} & être d'accord avec / consentir à --- \en{I agree with you. I agree to the terms.} \tabularnewline
\en{object to} + \en{-ing} & s'opposer à --- \en{I object to paying twice.} \tabularnewline
\en{look forward to} + \en{-ing} & attendre avec impatience --- \en{I look forward to meeting you.} \tabularnewline
\en{be used to} + \en{-ing} & être habitué à --- \en{I'm used to getting up early.} \tabularnewline
\en{prefer A to B} & préférer A à B --- \en{I prefer tea to coffee.} \tabularnewline
\end{longtable}
\end{tableau}

\begin{tableau}[Verbes suivis de \en{on / of / about / with / from}]
\footnotesize
\begin{longtable}{@{}p{4.4cm}p{9.8cm}@{}}
\entete{Verbe} & \entete{Sens et exemple} \tabularnewline
\endfirsthead
\entete{Verbe} & \entete{Sens et exemple} \tabularnewline
\endhead
\en{depend on} & dépendre de --- \en{It depends on the weather.} \tabularnewline
\en{rely on} & compter sur --- \en{You can rely on me.} \tabularnewline
\en{congratulate on} & féliciter pour --- \en{I congratulated him on his success.} \tabularnewline
\en{insist on} + \en{-ing} & insister pour --- \en{She insisted on paying.} \tabularnewline
\en{focus on} & se concentrer sur --- \en{Let's focus on the main issue.} \tabularnewline
\en{work on} & travailler sur --- \en{I'm working on a new project.} \tabularnewline
\en{spend (money/time) on} & dépenser en --- \en{I spent too much on books.} \tabularnewline
\en{think of / about} & penser à --- \en{I often think of you.} \tabularnewline
\en{dream of / about} & rêver de --- \en{I dreamt about the sea.} \tabularnewline
\en{hear of / about} & entendre parler de --- \en{I've never heard of him.} \tabularnewline
\en{talk about} & parler de --- \en{What are you talking about?} \tabularnewline
\en{complain about} & se plaindre de --- \en{He complained about the noise.} \tabularnewline
\en{argue about / with} & se disputer au sujet de / avec --- \en{They argued about money.} \tabularnewline
\en{agree with (sb)} & être d'accord avec --- \en{I agree with you.} \tabularnewline
\en{deal with} & traiter, gérer --- \en{I'll deal with it.} \tabularnewline
\en{cope with} & faire face à --- \en{She copes well with stress.} \tabularnewline
\en{fill in / fill out} & remplir (formulaire) --- \en{Fill in this form.} \tabularnewline
\en{believe in} & croire en --- \en{I believe in you.} \tabularnewline
\en{consist of} & consister en --- \en{The team consists of five people.} \tabularnewline
\en{succeed in} + \en{-ing} & réussir à --- \en{He succeeded in passing.} \tabularnewline
\en{prevent from} + \en{-ing} & empêcher de --- \en{Nothing prevented us from going.} \tabularnewline
\en{stop (sb) from} & empêcher --- \en{The rain stopped us from playing.} \tabularnewline
\en{protect from / against} & protéger de --- \en{Protect your skin from the sun.} \tabularnewline
\en{suffer from} & souffrir de --- \en{He suffers from asthma.} \tabularnewline
\en{die of / from} & mourir de --- \en{She died of malaria.} \tabularnewline
\en{recover from} & se remettre de --- \en{He recovered from his illness.} \tabularnewline
\en{escape from} & s'échapper de --- \en{They escaped from the building.} \tabularnewline
\en{differ from} & différer de --- \en{English differs from French.} \tabularnewline
\en{benefit from} & bénéficier de --- \en{You'll benefit from this course.} \tabularnewline
\en{learn from} & apprendre de --- \en{Learn from your mistakes.} \tabularnewline
\end{longtable}
\end{tableau}

\section{Adjectif + préposition : 50 combinaisons}
\begin{tableau}[Les adjectifs et leur préposition]
\footnotesize
\begin{longtable}{@{}p{4.4cm}p{9.8cm}@{}}
\entete{Adjectif} & \entete{Construction et exemple} \tabularnewline
\endfirsthead
\entete{Adjectif} & \entete{Construction et exemple} \tabularnewline
\endhead
\en{afraid of} & peur de --- \en{I'm afraid of spiders.} \tabularnewline
\en{good at / bad at} & bon / mauvais en --- \en{She's good at languages.} \tabularnewline
\en{interested in} & intéressé par --- \en{I'm interested in history.} \tabularnewline
\en{keen on} & passionné de --- \en{He's keen on football.} \tabularnewline
\en{fond of} & attaché à --- \en{I'm fond of my grandmother.} \tabularnewline
\en{proud of} & fier de --- \en{We're proud of you.} \tabularnewline
\en{ashamed of} & honteux de --- \en{He was ashamed of his mistake.} \tabularnewline
\en{tired of} & las de --- \en{I'm tired of waiting.} \tabularnewline
\en{bored with / by} & ennuyé par --- \en{She's bored with her job.} \tabularnewline
\en{fed up with} & exaspéré par --- \en{I'm fed up with this weather.} \tabularnewline
\en{angry with (sb) / about (sth)} & fâché contre / au sujet de --- \en{She's angry with me about the delay.} \tabularnewline
\en{annoyed with / about} & agacé par --- \en{He's annoyed about the noise.} \tabularnewline
\en{worried about} & inquiet au sujet de --- \en{I'm worried about the exam.} \tabularnewline
\en{concerned about} & préoccupé par --- \en{We're concerned about climate change.} \tabularnewline
\en{excited about} & enthousiaste à l'idée de --- \en{They're excited about the trip.} \tabularnewline
\en{surprised at / by} & surpris de --- \en{I was surprised at the result.} \tabularnewline
\en{amazed at / by} & émerveillé par --- \en{We were amazed by the view.} \tabularnewline
\en{shocked by / at} & choqué par --- \en{She was shocked by the news.} \tabularnewline
\en{disappointed with / by / in} & déçu par --- \en{I'm disappointed with the service.} \tabularnewline
\en{satisfied with} & satisfait de --- \en{Are you satisfied with your work?} \tabularnewline
\en{pleased with / about} & content de --- \en{He's pleased with his results.} \tabularnewline
\en{happy with / about} & heureux de --- \en{I'm happy with my new flat.} \tabularnewline
\en{full of} & plein de --- \en{The room was full of people.} \tabularnewline
\en{empty of} & vide de --- \en{A street empty of cars.} \tabularnewline
\en{covered with / in} & couvert de --- \en{Covered in dust.} \tabularnewline
\en{made of / from} & fait de / à partir de --- \en{Made of wood / made from grapes.} \tabularnewline
\en{similar to} & semblable à --- \en{Similar to the French system.} \tabularnewline
\en{different from} & différent de --- \en{Different from what I expected.} \tabularnewline
\en{the same as} & le même que --- \en{The same as yesterday.} \tabularnewline
\en{married to} & marié à --- \en{She's married to a doctor.} \tabularnewline
\en{engaged to} & fiancé à --- \en{He's engaged to Aïcha.} \tabularnewline
\en{divorced from} & divorcé de --- \en{Divorced from his first wife.} \tabularnewline
\en{related to} & lié à / parent de --- \en{Is he related to you?} \tabularnewline
\en{connected to / with} & relié à --- \en{Connected to the internet.} \tabularnewline
\en{addicted to} & dépendant de --- \en{Addicted to coffee.} \tabularnewline
\en{used to} + \en{-ing} & habitué à --- \en{I'm used to the heat.} \tabularnewline
\en{ready for / to} & prêt pour / à --- \en{Ready for the meeting / ready to go.} \tabularnewline
\en{available for / to} & disponible pour --- \en{Available for questions.} \tabularnewline
\en{responsible for} & responsable de --- \en{She's responsible for the budget.} \tabularnewline
\en{famous for} & célèbre pour --- \en{Famous for its beaches.} \tabularnewline
\en{capable of} + \en{-ing} & capable de --- \en{Capable of running 10 km.} \tabularnewline
\en{aware of} & conscient de --- \en{Are you aware of the risk?} \tabularnewline
\en{conscious of} & conscient de --- \en{Conscious of the time.} \tabularnewline
\en{guilty of} & coupable de --- \en{Guilty of theft.} \tabularnewline
\en{innocent of} & innocent de --- \en{Innocent of all charges.} \tabularnewline
\en{jealous of} & jaloux de --- \en{Jealous of his brother.} \tabularnewline
\en{envious of} & envieux de --- \en{Envious of her success.} \tabularnewline
\en{grateful to (sb) for (sth)} & reconnaissant envers --- \en{Grateful to you for your help.} \tabularnewline
\en{sorry for / about} & désolé de --- \en{Sorry for the noise / about yesterday.} \tabularnewline
\en{typical of} & typique de --- \en{That's typical of him.} \tabularnewline
\end{longtable}
\end{tableau}

\section{La logique des particules}
Les particules des \emph{phrasal verbs} ne sont pas arbitraires : elles conservent presque
toujours leur sens spatial d'origine. Comprendre cette logique divise par deux l'effort de
mémorisation des 150 verbes du chapitre 13.

\begin{tableau}[Les dix particules et leur image mentale]
\footnotesize
\begin{longtable}{@{}p{2.0cm}p{5.2cm}p{7.0cm}@{}}
\entete{Particule} & \entete{Image d'origine} & \entete{Sens dérivé et exemples} \tabularnewline
\endfirsthead
\entete{Particule} & \entete{Image} & \entete{Sens dérivé} \tabularnewline
\endhead
\en{up} & vers le haut, puis \textbf{jusqu'au bout} & Complétude, augmentation, amélioration. \en{eat up} (tout manger), \en{clean up} (nettoyer à fond), \en{speed up} (accélérer), \en{cheer up} (remonter le moral), \en{give up} (abandonner), \en{look up} (chercher). \tabularnewline
\en{down} & vers le bas & Réduction, arrêt, échec. \en{slow down}, \en{calm down}, \en{turn down} (refuser / baisser), \en{break down} (tomber en panne), \en{write down}, \en{shut down}. \tabularnewline
\en{off} & séparation, éloignement & Départ, coupure, éloignement. \en{set off} (partir), \en{take off} (décoller / enlever), \en{turn off}, \en{put off} (reporter), \en{call off} (annuler), \en{show off} (frimer), \en{drop off} (déposer). \tabularnewline
\en{out} & vers l'extérieur, puis \textbf{visible} & Apparition, épuisement, révélation. \en{go out}, \en{find out} (découvrir), \en{work out} (fonctionner / s'entraîner), \en{run out of} (ne plus avoir), \en{pass out} (s'évanouir), \en{figure out} (comprendre), \en{come out}. \tabularnewline
\en{on} & contact, continuité & Poursuite, mise en marche, dépendance. \en{go on} (continuer), \en{carry on}, \en{turn on}, \en{put on} (mettre un vêtement), \en{get on with} (s'entendre avec), \en{rely on}, \en{hold on} (attendre). \tabularnewline
\en{back} & en arrière & Retour, restitution. \en{come back}, \en{go back}, \en{give back} (rendre), \en{call back}, \en{look back}, \en{pay back} (rembourser), \en{get back to} (recontacter). \tabularnewline
\en{over} & au-dessus, d'un bout à l'autre & Répétition, transfert, fin. \en{think over} (réfléchir à), \en{go over} (revoir), \en{take over} (reprendre le contrôle), \en{get over} (surmonter), \en{start over} (recommencer). \tabularnewline
\en{away} & au loin, disparition & Éloignement, disparition progressive. \en{go away}, \en{run away}, \en{throw away} (jeter), \en{fade away}, \en{put away} (ranger), \en{give away} (offrir / révéler). \tabularnewline
\en{through} & à travers, du début à la fin & Traversée, achèvement. \en{get through} (joindre / traverser), \en{go through} (parcourir), \en{look through} (feuilleter), \en{see through} (voir clair dans), \en{come through} (survivre). \tabularnewline
\en{in / into} & à l'intérieur & Entrée, implication, transformation. \en{get in}, \en{check in}, \en{drop in} (passer voir), \en{break into} (entrer par effraction), \en{turn into} (se transformer en), \en{look into} (examiner). \tabularnewline
\end{longtable}
\end{tableau}

\begin{astuce}[Comment deviner un \emph{phrasal verb} inconnu]
Décomposez : \textbf{verbe de base} (sens concret) + \textbf{particule} (image du tableau
ci-dessus). \en{To \textbf{hand} \textbf{in}} = donner « vers l'intérieur » \ra remettre
(une copie, un dossier). \en{To \textbf{sleep} \textbf{off}} = dormir en « séparant » \ra
cuver. Cette gymnastique fonctionne environ une fois sur deux --- et vous donne une
hypothèse plausible plutôt qu'un blocage total.
\end{astuce}

\begin{exo}[D.1 --- Complétez avec la préposition]
\begin{enumerate}\setlength{\itemsep}{2pt}
  \item \en{It depends \trou[1.5cm] the weather.}
  \item \en{I'm really interested \trou[1.5cm] African history.}
  \item \en{She's married \trou[1.5cm] an engineer.}
  \item \en{He's very good \trou[1.5cm] telling stories.}
  \item \en{I'm looking forward \trou[1.5cm] seeing you again.}
  \item \en{They congratulated me \trou[1.5cm] my promotion.}
  \item \en{This bag belongs \trou[1.5cm] my brother.}
  \item \en{We're proud \trou[1.5cm] our results.}
  \item \en{He suffers \trou[1.5cm] back pain.}
  \item \en{The film is similar \trou[1.5cm] the one we saw last week.}
  \item \en{I apologise \trou[1.5cm] being late.}
  \item \en{She's responsible \trou[1.5cm] the whole project.}
\end{enumerate}
\end{exo}
\begin{corrige}[D.1]
1. \en{on} 2. \en{in} 3. \en{to} 4. \en{at} 5. \en{to} 6. \en{on} 7. \en{to}
8. \en{of} 9. \en{from} 10. \en{to} 11. \en{for} 12. \en{for}
\end{corrige}

\begin{exo}[D.2 --- Choisissez le bon connecteur]
\begin{enumerate}\setlength{\itemsep}{2pt}
  \item (\en{Although / Despite}) the rain, we went to the beach.
  \item (\en{Although / Despite}) it was raining, we went to the beach.
  \item He studied hard; (\en{therefore / however}), he passed the exam.
  \item (\en{Because / Because of}) I was tired, I went to bed early.
  \item (\en{So that / So as to}) not be late, I left at seven.
  \item It was (\en{so / such}) a difficult exam that nobody passed.
  \item (\en{Eventually / Finally}), I'd like to thank my teachers.
  \item We waited three hours, but (\en{eventually / at last}) he arrived.
  \item She's fluent (\en{thanks to / due to}) her mother, who is British.
  \item (\en{In spite of / Even though}) he's rich, he isn't happy.
\end{enumerate}
\end{exo}
\begin{corrige}[D.2]
1. \en{Despite} (suivi d'un nom). 2. \en{Although} (suivi d'une proposition).
3. \en{therefore}. 4. \en{Because} (proposition) --- \en{Because of my tiredness} sinon.
5. \en{So as not to}. 6. \en{such} (\en{such + a/an + nom}) --- \en{so difficult that} sinon.
7. \en{Finally}. 8. \en{eventually}. 9. \en{thanks to} (connotation positive).
10. \en{Even though} (proposition).
\end{corrige}

\begin{memo}[Révision express --- Annexe D]
Dix connecteurs à connaître absolument, pour commencer : \en{however, although, because of,
therefore, as a result, in addition, for instance, on the other hand, as soon as, to sum up}.
Dix régimes prépositionnels à ne jamais rater : \en{depend \textbf{on}, listen \textbf{to},
look \textbf{for}, married \textbf{to}, good \textbf{at}, interested \textbf{in},
afraid \textbf{of}, wait \textbf{for}, arrive \textbf{at}, talk \textbf{about}}.
\end{memo}

% =====================================================================
%  ANNEXE E — RESSOURCES
% =====================================================================
\chapter{Ressources : quoi utiliser, et surtout quoi ne pas utiliser}

Le plus grand danger de l'apprenant moderne n'est pas le manque de ressources : c'est leur
abondance. Passer trois heures à comparer dix applications, c'est trois heures où vous ne
parlez pas. Cette annexe vous donne une sélection courte et justifiée --- puis une règle :
\textbf{choisissez une ressource par besoin, et gardez-la douze mois.}

\begin{piege}[Le piège du collectionneur]
Symptômes : vous avez installé sept applications, ouvert trois chaînes YouTube, créé cinq
comptes, téléchargé deux méthodes\ldots{} et votre niveau n'a pas bougé.\\
Remède : une application de vocabulaire, une source d'écoute, un partenaire de parole,
un dictionnaire. \textbf{Quatre outils, pas quarante.} La variété ne fait pas progresser ;
la répétition, si.
\end{piege}

\section{Le kit minimal : quatre outils, zéro franc}
\begin{tableau}[Ce dont vous avez réellement besoin]
\begin{tabularx}{\linewidth}{@{}p{3.0cm}p{4.0cm}X@{}}
\entete{Besoin} & \entete{Outil} & \entete{Comment l'utiliser} \tabularnewline
\hline
\textbf{Vocabulaire} & Anki (gratuit sur ordinateur et Android) &
15 cartes par jour, créées \textbf{par vous}, avec une phrase entière au recto. Révision quotidienne de 10 minutes. \tabularnewline
\textbf{Dictionnaire} & WordReference ou Cambridge Dictionary (gratuits) &
Cherchez un mot, \textbf{lisez la phrase d'exemple}, notez-la. Ne cherchez jamais plus de cinq mots à la suite. \tabularnewline
\textbf{Écoute} & Un podcast ou une chaîne pour apprenants (gratuit) &
Un épisode de 5 à 10 minutes, trois fois, selon le protocole des 5 passes du chapitre 16. \tabularnewline
\textbf{Parole} & Un humain : ami, collègue, club, application d'échange &
Trois fois par semaine, 20 minutes minimum. Rien ne remplace cette ligne du tableau. \tabularnewline
\end{tabularx}
\end{tableau}

\section{Dictionnaires et outils de référence}
\begin{tableau}[Les six références utiles]
\footnotesize
\begin{longtable}{@{}p{4.6cm}p{9.6cm}@{}}
\entete{Outil} & \entete{À quoi il sert} \tabularnewline
\endfirsthead
\entete{Outil} & \entete{À quoi il sert} \tabularnewline
\endhead
\en{WordReference} & Dictionnaire anglais-français très riche, avec forums où des milliers de nuances sont discutées. Idéal au début. \tabularnewline
\en{Cambridge Dictionary} & Dictionnaire anglais-anglais pour apprenants, avec transcription phonétique britannique \textbf{et} américaine, et audio. À privilégier dès le niveau A2. \tabularnewline
\en{Oxford Learner's Dictionaries} & Même usage, très fiable, avec l'\en{Oxford 3000} : la liste des 3\,000 mots essentiels --- exactement le volume visé par ce livre. \tabularnewline
\en{Linguee / DeepL Write} & Traducteur de \textbf{contextes} : vous voyez comment un mot est réellement employé dans des documents bilingues. Excellent pour vérifier une tournure. \tabularnewline
\en{Forvo} & Prononciation enregistrée par des natifs du monde entier. Indispensable pour un mot dont vous doutez. \tabularnewline
\en{YouGlish} & Tapez un mot ou une expression : le site trouve des vidéos YouTube où des natifs la prononcent en contexte réel. L'outil le plus utile pour l'accent. \tabularnewline
\end{longtable}
\end{tableau}

\begin{astuce}[Le réflexe anglais-anglais]
Dès le troisième mois, cessez de chercher \en{book} = « livre ». Cherchez la
\textbf{définition anglaise} et la \textbf{phrase d'exemple}. Cela double le temps de
consultation mais triple la mémorisation : vous apprenez le mot dans son habitat naturel,
et vous lisez de l'anglais au passage.
\end{astuce}

\section{Écoute : où trouver de l'anglais adapté}
\begin{tableau}[Podcasts et chaînes, par niveau]
\footnotesize
\begin{longtable}{@{}p{2.4cm}p{4.6cm}p{7.2cm}@{}}
\entete{Niveau} & \entete{Type de ressource} & \entete{Ce qu'il faut chercher} \tabularnewline
\endfirsthead
\entete{Niveau} & \entete{Type} & \entete{Critères de choix} \tabularnewline
\endhead
\textbf{A1--A2} & Podcasts pour apprenants (débit lent, vocabulaire contrôlé) &
Épisodes de 3 à 8 minutes, texte disponible. Exemples de formats à chercher : « 6 Minute English », « Learning English Broadcast », « English for beginners ». \tabularnewline
\textbf{A2--B1} & Chaînes pédagogiques, histoires racontées, actualités simplifiées &
Épisodes de 10 minutes avec transcript. Les sites d'actualités à trois niveaux (le même article en facile, moyen, difficile) sont parfaits. \tabularnewline
\textbf{B1--B2} & Podcasts conversationnels pour apprenants avancés, séries, interviews &
Débit naturel mais articulé ; sujets variés. Commencez \textbf{sans} sous-titres, puis vérifiez avec. \tabularnewline
\textbf{B2+} & Podcasts natifs : actualité, sciences, sport, humour, débats &
Ce que les Anglophones écoutent eux-mêmes. Objectif : comprendre sans aucun support. \tabularnewline
\end{longtable}
\end{tableau}

\begin{regle}[Comment choisir sa source d'écoute]
Trois critères, dans cet ordre : \textbf{(1) un transcript existe} (sinon le protocole des
5 passes est impossible) ; \textbf{(2) le sujet vous intéresse vraiment} --- on supporte
mieux la difficulté quand le contenu plaît ; \textbf{(3) une voix que vous aimez entendre}.
Vous écouterez cette personne 150 heures cette année : choisissez bien.
\end{regle}

\section{Lecture : des textes à votre taille}
\begin{tableau}[Progression de lecture sur 12 mois]
\begin{tabularx}{\linewidth}{@{}p{2.6cm}p{3.6cm}X@{}}
\entete{Période} & \entete{Type de texte} & \entete{Volume conseillé} \tabularnewline
\hline
Mois 1--3 & Textes écrits pour les apprenants, livres gradués niveau 1-2, articles très courts & 10 minutes par jour, 300 mots par texte. \tabularnewline
Mois 4--6 & Livres gradués niveau 3-4, articles d'actualité simplifiés, blogs personnels simples & 15 minutes par jour, un chapitre par semaine. \tabularnewline
Mois 7--9 & Romans jeunesse, premières pages de Wikipédia en anglais simple, articles de presse sur des sujets connus & 20 minutes par jour, 30 pages par semaine. \tabularnewline
Mois 10--12 & Romans pour adultes, presse générale, essais, blogs spécialisés & 30 minutes par jour, un livre par mois. \tabularnewline
\end{tabularx}
\end{tableau}

\begin{astuce}[Où trouver des livres gratuits et légaux]
Le projet Gutenberg (gutenberg.org) propose des dizaines de milliers de livres tombés dans
le domaine public, en anglais, téléchargeables en une fois. Commencez par des textes courts :
nouvelles, contes, récits de voyage. La version simplifiée de Wikipédia en anglais
(simple.wikipedia.org) est une mine : des milliers d'articles écrits en anglais facile.
Enfin, les bibliothèques municipales et les centres culturels étrangers prêtent souvent des
livres en anglais --- et organisent des clubs de conversation gratuits.
\end{astuce}

\section{Parler : trouver des partenaires}
\begin{tableau}[Les six voies possibles, de la plus accessible à la plus efficace]
\begin{tabularx}{\linewidth}{@{}p{3.2cm}p{2.0cm}X@{}}
\entete{Voie} & \entete{Coût} & \entete{Comment faire} \tabularnewline
\hline
\textbf{Se parler à soi-même} & 0 &
Narration continue de vos gestes, monologue sous la douche, réponses aux 30 questions du chapitre 15. C'est la voie la plus sous-estimée : elle crée l'automatisme sans peur du jugement. \tabularnewline
\textbf{Applications d'échange linguistique} & 0 &
Des applications de type « tandem » mettent en relation des francophones et des anglophones : 15 minutes en anglais, 15 minutes en français. Cherchez un partenaire sérieux, fixez un rendez-vous hebdomadaire, gardez-le un an. \tabularnewline
\textbf{Clubs et associations} & 0 à faible &
Clubs d'anglais universitaires, associations culturelles, centres culturels britanniques ou américains, groupes de discussion en ligne. L'avantage : la régularité et le collectif. \tabularnewline
\textbf{Amis et collègues anglophones} & 0 &
Proposez explicitement : « Pendant nos pauses, parlons anglais 15 minutes. » Beaucoup acceptent, personne n'ose demander. \tabularnewline
\textbf{Réseaux sociaux et messagerie} & 0 &
Commentez en anglais sous des vidéos, participez à des forums sur vos centres d'intérêt, échangez des messages vocaux. L'écrit spontané est un excellent pont vers l'oral. \tabularnewline
\textbf{Tuteurs en ligne} & payant &
Plateformes de cours particuliers à la demande. Même deux heures par mois suffisent à débloquer l'oral et à faire corriger vos erreurs récurrentes. Cherchez un tuteur qui \textbf{vous fait parler} 80\,\% du temps, pas un professeur qui explique. \tabularnewline
\end{tabularx}
\end{tableau}

\begin{mission}[La phrase à envoyer aujourd'hui]
Copiez, adaptez et envoyez ce message à trois personnes anglophones de votre entourage ---
ou à trois contacts d'une application d'échange :\\[4pt]
\en{Hello! My name is \ldots{} and I'm from Senegal. I've been learning English for about a
year and I'd love to practise speaking with a real person. Would you be interested in a
language exchange? We could talk for 20 minutes twice a week --- half in English, half in
French. No pressure at all, just let me know if you'd like to try!}\\[4pt]
Sur trois messages, vous obtiendrez statistiquement une réponse enthousiaste. C'est le
meilleur retour sur investissement de toute cette annexe.
\end{mission}

\section{Écrire et se faire corriger}
\begin{tableau}[Outils d'écriture]
\begin{tabularx}{\linewidth}{@{}p{4.0cm}X@{}}
\entete{Outil} & \entete{Usage} \tabularnewline
\hline
\textbf{Correcteurs automatiques} & Utiles pour l'orthographe et les erreurs de base. Mais ne les laissez \textbf{pas} corriger à votre place : lisez la suggestion, comprenez-la, puis écrivez-la vous-même. Sinon vous ne progresserez jamais. \tabularnewline
\textbf{Sites de correction par des natifs} & Certains sites permettent de soumettre un texte et de recevoir des corrections commentées par des anglophones. Renvoyez l'ascenseur en corrigeant des textes français. \tabularnewline
\textbf{Le journal personnel} & L'outil le plus puissant et le moins cher : 60 invites au chapitre 17, une par jour, 100 mots. Personne ne vous juge, donc vous osez des phrases ambitieuses. \tabularnewline
\textbf{Les échanges écrits réels} & Messages, e-mails, commentaires, forums. Écrire à un humain motive infiniment plus qu'écrire un exercice. \tabularnewline
\end{tabularx}
\end{tableau}

\begin{astuce}[La règle des deux couleurs]
Quand un texte vous est corrigé, recopiez-le : en \textbf{noir} ce que vous aviez écrit,
en \textbf{rouge} chaque correction. Relisez cette version une semaine plus tard. Vous
constaterez que 80\,\% de vos erreurs sont toujours les mêmes --- ce sont elles qu'il faut
inscrire dans le plan « trois erreurs par semaine » du chapitre 18.
\end{astuce}

\section{Grammaire : trois livres suffisent}
\begin{tableau}[Les ouvrages de référence classiques]
\begin{tabularx}{\linewidth}{@{}p{5.0cm}X@{}}
\entete{Ouvrage} & \entete{Pourquoi} \tabularnewline
\hline
\textbf{Une grammaire pratique pour apprenants} (type \emph{English Grammar in Use}) &
Une page de règle, une page d'exercices, corrigés inclus. Le format le plus efficace jamais inventé. Prenez le niveau intermédiaire. \tabularnewline
\textbf{Un ouvrage d'usage pour les questions fines} (type \emph{Practical English Usage}) &
Organisé comme un dictionnaire des difficultés. À consulter quand vous butez sur une nuance --- pas à lire d'une traite. \tabularnewline
\textbf{Un dictionnaire anglais-anglais pour apprenants} &
Voir la section précédente. C'est le livre que vous ouvrirez le plus souvent de toute votre année. \tabularnewline
\end{tabularx}
\end{tableau}

\begin{piege}[N'achetez pas trois méthodes]
Une méthode de grammaire, un dictionnaire, et ce livre : c'est tout. Ajouter une deuxième
méthode crée le sentiment de progresser sans produire le moindre progrès, parce que vous
passez votre temps à recommencer au chapitre 1 d'un autre ouvrage.
\end{piege}

\section{Applications : ce qu'elles font bien, ce qu'elles ne font pas}
\begin{tableau}[Verdict honnête]
\begin{tabularx}{\linewidth}{@{}p{4.4cm}p{4.6cm}X@{}}
\entete{Type d'application} & \entete{Ce qu'elle fait bien} & \entete{Ce qu'elle ne fait pas} \tabularnewline
\hline
\textbf{Applications ludiques de type « cours du jour »} &
Créer une habitude quotidienne, entretenir la motivation, donner les 800 premiers mots. &
Elles ne vous feront \textbf{jamais} parler. Leurs phrases sont trop courtes, trop propres, trop faciles. Utiles les trois premiers mois, puis insuffisantes. \tabularnewline
\textbf{Répétition espacée (cartes mémoire)} &
Mémoriser 3\,000 mots de façon durable --- c'est l'outil le plus efficace qui existe pour le vocabulaire. &
Elles n'enseignent ni la prononciation, ni la syntaxe, ni la spontanéité. \tabularnewline
\textbf{Reconnaissance vocale} &
Vérifier qu'un son est correct, donner un retour immédiat sans honte. &
Elle comprend trop bien : elle devine votre intention. Un humain, lui, ne devine pas. \tabularnewline
\textbf{Traduction automatique} &
Dépanner, vérifier un doute, donner un point de départ. &
Elle vous empêche de faire l'effort qui construit la compétence. À utiliser \textbf{après} avoir essayé, jamais avant. \tabularnewline
\end{tabularx}
\end{tableau}

\section{Sans internet : huit ressources que vous avez déjà}
Si votre connexion est limitée, rien n'est perdu. Voici ce qui fonctionne hors ligne.

\begin{enumerate}
  \item \textbf{Votre téléphone.} Changez la langue du système en anglais. Vous apprendrez
  200 mots d'interface en trois semaines, sans effort : \en{settings, share, save, delete,
  search, download, mute, brightness}.
  \item \textbf{La radio.} Les stations internationales anglophones émettent en ondes
  courtes ou en FM dans de nombreuses régions d'Afrique de l'Ouest. Vingt minutes par jour
  suffisent --- même en bruit de fond, votre oreille se calibre.
  \item \textbf{La télévision.} Un film que vous connaissez déjà, en version originale avec
  sous-titres anglais. Vous connaissez l'histoire : votre cerveau est libre de se concentrer
  sur la langue.
  \item \textbf{Les fichiers téléchargés.} Téléchargez podcasts et vidéos quand vous avez du
  réseau, puis écoutez-les hors ligne --- dans le bus, en marchant, en cuisinant. C'est du
  temps gratuit.
  \item \textbf{Les livres papier.} Un seul roman anglais lu trois fois vaut mieux que dix
  romans jamais ouverts. Les bibliothèques, librairies d'occasion et centres culturels en
  fournissent gratuitement.
  \item \textbf{Les humains autour de vous.} Un collègue, un voisin, un commerçant, un
  étudiant en anglais : créez un club de deux personnes qui se parlent en anglais une heure
  par semaine. C'est la ressource la plus puissante de cette liste.
  \item \textbf{Les produits de consommation.} Lisez les étiquettes, les notices, les
  emballages, les modes d'emploi en anglais. Vocabulaire concret et immédiatement réutilisable.
  \item \textbf{Ce livre.} Les dix textes d'écoute du chapitre 16, les trois textes de
  lecture du chapitre 17, les 600 phrases de l'annexe B, les 17 dialogues du chapitre 14 :
  vous avez déjà 20 heures de matériel hors ligne, corrigé et gradué.
\end{enumerate}

\begin{tableau}[Budget-type sur douze mois]
\begin{tabularx}{\linewidth}{@{}p{5.4cm}p{2.6cm}X@{}}
\entete{Poste} & \entete{Coût} & \entete{Commentaire} \tabularnewline
\hline
Anki + dictionnaires en ligne & 0 & Gratuit, sans publicité, sans limite. \tabularnewline
Podcasts et vidéos & 0 & Illimité. Le seul coût : vos données mobiles. \tabularnewline
Livres gradués d'occasion & faible & Deux ou trois suffisent pour l'année. \tabularnewline
Partenaire d'échange & 0 & Vous payez en français, eux en anglais. \tabularnewline
Tuteur en ligne (facultatif) & variable & 2 heures par mois, à partir du mois 4. \tabularnewline
Examen blanc en ligne (facultatif) & 0 à faible & Au mois 6 et au mois 12, pour une mesure objective. \tabularnewline
\end{tabularx}
\end{tableau}

\begin{memo}[Révision express --- Annexe E]
Quatre outils, pas quarante : \textbf{Anki} pour les mots, \textbf{un dictionnaire
anglais-anglais} pour le sens, \textbf{un podcast avec transcript} pour l'oreille,
\textbf{un humain} pour la bouche. Tout le reste est du bruit. Et souvenez-vous : la
meilleure ressource du monde ne fait rien à votre place. Vingt minutes par jour, pendant
365 jours, valent mieux que n'importe quelle collection d'applications.
\end{memo}

% =====================================================================
%  ANNEXE F — JOURNAL DE PROGRESSION (à imprimer)
% =====================================================================
\chapter{Votre journal de progression}

Cette annexe est faite pour être \textbf{imprimée, remplie à la main, et conservée}.
Un journal de progression remplit trois fonctions : il rend l'effort visible, il
transforme une intention en rendez-vous, et il fournit --- dans les moments de doute ---
la preuve que vous avancez. Au mois 8, quand la motivation fléchira, vous relirez les
pages des mois 2 et 3.

\begin{regle}[La seule règle du journal]
\textbf{On remplit le journal le jour même, jamais le lendemain.} Deux minutes suffisent :
trois cases à cocher, un chiffre à écrire, une phrase à noter. Un journal tenu une fois par
semaine devient un journal abandonné en un mois.
\end{regle}

\section{Le contrat personnel}
À remplir le jour 1, à signer, et à relire chaque premier dimanche du mois.

\begin{tableau}[Mon contrat d'apprentissage]
\begin{tabularx}{\linewidth}{@{}p{7.4cm}X@{}}
\entete{Question} & \entete{Ma réponse} \tabularnewline
\hline
Pourquoi est-ce que je veux parler anglais ? (la vraie raison, pas la raison polie) & \rule{0pt}{2.2cm}\trou[6.6cm] \tabularnewline
Qu'est-ce que cela changera concrètement dans ma vie ? & \rule{0pt}{2.0cm}\trou[6.6cm] \tabularnewline
Combien de minutes par jour est-ce que je m'engage à y consacrer ? & \rule{0pt}{1.2cm}\trou[6.6cm] \tabularnewline
À quelle heure exactement ? (créez le rendez-vous) & \rule{0pt}{1.2cm}\trou[6.6cm] \tabularnewline
Où ? (le lieu compte autant que l'heure) & \rule{0pt}{1.2cm}\trou[6.6cm] \tabularnewline
Quelle est la plus petite version possible de ma séance, les jours difficiles ? & \rule{0pt}{1.6cm}\trou[6.6cm] \tabularnewline
Qui sera au courant de mon engagement ? & \rule{0pt}{1.2cm}\trou[6.6cm] \tabularnewline
Quelle récompense si je tiens 90 jours ? & \rule{0pt}{1.2cm}\trou[6.6cm] \tabularnewline
Ma date de début : & \rule{0pt}{1.0cm}\trou[4.0cm] \tabularnewline
Ma date de fin (365 jours plus tard) : & \rule{0pt}{1.0cm}\trou[4.0cm] \tabularnewline
Signature : & \rule{0pt}{1.2cm}\trou[6.6cm] \tabularnewline
\end{tabularx}
\end{tableau}

\begin{astuce}[Le seuil minimal des mauvais jours]
Écrivez dans votre contrat une « plus petite version possible » ridiculement facile :
\emph{5 minutes --- 5 cartes Anki et une phrase à voix haute}. Le but n'est pas de
progresser ces jours-là, mais de \textbf{ne pas casser la chaîne}. Une chaîne de 300 jours
avec vingt séances de 5 minutes vaut infiniment mieux qu'une chaîne brisée trois fois.
\end{astuce}

\section{La grille hebdomadaire}
À photocopier 52 fois. Une ligne par jour, une case par activité. Une case remplie vaut
mieux qu'une case parfaite.

\newcommand{\case}{\fbox{\hspace{1.0em}\rule[-0.35em]{0pt}{1.05em}}}
\begin{tableau}[Semaine du \rule{2cm}{0.4pt} au \rule{2cm}{0.4pt}]
\footnotesize
\begin{tabularx}{\linewidth}{@{}p{1.5cm}p{1.0cm}p{1.0cm}p{1.0cm}p{1.0cm}p{1.0cm}p{1.0cm}p{1.0cm}X@{}}
\entete{Jour} & \entete{15 min Anki} & \entete{15 min grammaire} & \entete{20 min écoute} & \entete{15 min lecture} & \entete{10 min écrit} & \entete{20 min oral} & \case & \entete{Note du jour} \tabularnewline
\hline
Lundi & \case & \case & \case & \case & \case & \case & \case & \tabularnewline
Mardi & \case & \case & \case & \case & \case & \case & \case & \tabularnewline
Mercredi & \case & \case & \case & \case & \case & \case & \case & \tabularnewline
Jeudi & \case & \case & \case & \case & \case & \case & \case & \tabularnewline
Vendredi & \case & \case & \case & \case & \case & \case & \case & \tabularnewline
Samedi & \case & \case & \case & \case & \case & \case & \case & \tabularnewline
Dimanche & \multicolumn{6}{@{}p{11.0cm}}{\footnotesize Bilan et test de la semaine (30 min)} & \case & \tabularnewline
\hline
\multicolumn{9}{@{}p{16.2cm}@{}}{\footnotesize\textbf{Total de la semaine} \quad Jours étudiés : \trou[1.2cm] / 7 \quad Minutes : \trou[1.6cm] \quad Mots nouveaux : \trou[1.2cm] \quad Minutes de parole : \trou[1.4cm]} \tabularnewline
\end{tabularx}
\end{tableau}

\begin{tableau}[Les sept activités, en clair]
\footnotesize
\begin{tabularx}{\linewidth}{@{}p{2.4cm}p{1.4cm}X@{}}
\entete{Activité} & \entete{Durée} & \entete{Ce qu'elle contient} \tabularnewline
\hline
\textbf{Anki} & 15 min & 15 cartes nouvelles + révision du jour. Toujours avec une phrase entière. \tabularnewline
\textbf{Grammaire} & 15 min & Une unité de ce livre : règle, exemples lus à voix haute, un exercice corrigé. \tabularnewline
\textbf{Écoute} & 20 min & Un même audio, protocole des 5 passes (chapitre 16). Le même texte trois jours de suite. \tabularnewline
\textbf{Lecture} & 15 min & Un texte à 98\,\% de mots connus. Plaisir d'abord, dictionnaire ensuite. \tabularnewline
\textbf{Écrit} & 10 min & Journal : 80 à 100 mots sur une des 60 invites du chapitre 17. \tabularnewline
\textbf{Oral} & 20 min & Conversation réelle, ou monologue enregistré, ou un des 7 drills du chapitre 15. \tabularnewline
\textbf{Bonus} & --- & Chanson, série, film, message vocal à un ami, téléphone en anglais. \tabularnewline
\end{tabularx}
\end{tableau}

\section{Le registre annuel}
Une ligne par semaine. Ce tableau tient sur une seule page : affichez-le.

\begin{tableau}[52 semaines --- remplissez le dimanche soir]
\footnotesize
\begin{longtable}{@{}p{1.4cm}p{1.6cm}p{2.2cm}p{2.2cm}p{2.4cm}p{2.2cm}p{2.0cm}@{}}
\entete{Sem.} & \entete{Dates} & \entete{Jours /7} & \entete{Minutes totales} & \entete{Minutes de parole} & \entete{Mots actifs} & \entete{Humeur} \tabularnewline
\endfirsthead
\entete{Sem.} & \entete{Dates} & \entete{Jours} & \entete{Minutes} & \entete{Parole} & \entete{Mots} & \entete{Humeur} \tabularnewline
\endhead
\hline
1 & \trou[1.4cm] & \trou[2.0cm] & \trou[2.0cm] & \trou[2.2cm] & \trou[2.0cm] & \trou[1.8cm] \tabularnewline
2 & \trou[1.4cm] & \trou[2.0cm] & \trou[2.0cm] & \trou[2.2cm] & \trou[2.0cm] & \trou[1.8cm] \tabularnewline
3 & \trou[1.4cm] & \trou[2.0cm] & \trou[2.0cm] & \trou[2.2cm] & \trou[2.0cm] & \trou[1.8cm] \tabularnewline
4 & \trou[1.4cm] & \trou[2.0cm] & \trou[2.0cm] & \trou[2.2cm] & \trou[2.0cm] & \trou[1.8cm] \tabularnewline
5 & \trou[1.4cm] & \trou[2.0cm] & \trou[2.0cm] & \trou[2.2cm] & \trou[2.0cm] & \trou[1.8cm] \tabularnewline
6 & \trou[1.4cm] & \trou[2.0cm] & \trou[2.0cm] & \trou[2.2cm] & \trou[2.0cm] & \trou[1.8cm] \tabularnewline
7 & \trou[1.4cm] & \trou[2.0cm] & \trou[2.0cm] & \trou[2.2cm] & \trou[2.0cm] & \trou[1.8cm] \tabularnewline
8 & \trou[1.4cm] & \trou[2.0cm] & \trou[2.0cm] & \trou[2.2cm] & \trou[2.0cm] & \trou[1.8cm] \tabularnewline
9 & \trou[1.4cm] & \trou[2.0cm] & \trou[2.0cm] & \trou[2.2cm] & \trou[2.0cm] & \trou[1.8cm] \tabularnewline
10 & \trou[1.4cm] & \trou[2.0cm] & \trou[2.0cm] & \trou[2.2cm] & \trou[2.0cm] & \trou[1.8cm] \tabularnewline
11 & \trou[1.4cm] & \trou[2.0cm] & \trou[2.0cm] & \trou[2.2cm] & \trou[2.0cm] & \trou[1.8cm] \tabularnewline
12 & \trou[1.4cm] & \trou[2.0cm] & \trou[2.0cm] & \trou[2.2cm] & \trou[2.0cm] & \trou[1.8cm] \tabularnewline
13 & \trou[1.4cm] & \trou[2.0cm] & \trou[2.0cm] & \trou[2.2cm] & \trou[2.0cm] & \trou[1.8cm] \tabularnewline
14 & \trou[1.4cm] & \trou[2.0cm] & \trou[2.0cm] & \trou[2.2cm] & \trou[2.0cm] & \trou[1.8cm] \tabularnewline
15 & \trou[1.4cm] & \trou[2.0cm] & \trou[2.0cm] & \trou[2.2cm] & \trou[2.0cm] & \trou[1.8cm] \tabularnewline
16 & \trou[1.4cm] & \trou[2.0cm] & \trou[2.0cm] & \trou[2.2cm] & \trou[2.0cm] & \trou[1.8cm] \tabularnewline
17 & \trou[1.4cm] & \trou[2.0cm] & \trou[2.0cm] & \trou[2.2cm] & \trou[2.0cm] & \trou[1.8cm] \tabularnewline
18 & \trou[1.4cm] & \trou[2.0cm] & \trou[2.0cm] & \trou[2.2cm] & \trou[2.0cm] & \trou[1.8cm] \tabularnewline
19 & \trou[1.4cm] & \trou[2.0cm] & \trou[2.0cm] & \trou[2.2cm] & \trou[2.0cm] & \trou[1.8cm] \tabularnewline
20 & \trou[1.4cm] & \trou[2.0cm] & \trou[2.0cm] & \trou[2.2cm] & \trou[2.0cm] & \trou[1.8cm] \tabularnewline
21 & \trou[1.4cm] & \trou[2.0cm] & \trou[2.0cm] & \trou[2.2cm] & \trou[2.0cm] & \trou[1.8cm] \tabularnewline
22 & \trou[1.4cm] & \trou[2.0cm] & \trou[2.0cm] & \trou[2.2cm] & \trou[2.0cm] & \trou[1.8cm] \tabularnewline
23 & \trou[1.4cm] & \trou[2.0cm] & \trou[2.0cm] & \trou[2.2cm] & \trou[2.0cm] & \trou[1.8cm] \tabularnewline
24 & \trou[1.4cm] & \trou[2.0cm] & \trou[2.0cm] & \trou[2.2cm] & \trou[2.0cm] & \trou[1.8cm] \tabularnewline
25 & \trou[1.4cm] & \trou[2.0cm] & \trou[2.0cm] & \trou[2.2cm] & \trou[2.0cm] & \trou[1.8cm] \tabularnewline
26 & \trou[1.4cm] & \trou[2.0cm] & \trou[2.0cm] & \trou[2.2cm] & \trou[2.0cm] & \trou[1.8cm] \tabularnewline
27 & \trou[1.4cm] & \trou[2.0cm] & \trou[2.0cm] & \trou[2.2cm] & \trou[2.0cm] & \trou[1.8cm] \tabularnewline
28 & \trou[1.4cm] & \trou[2.0cm] & \trou[2.0cm] & \trou[2.2cm] & \trou[2.0cm] & \trou[1.8cm] \tabularnewline
29 & \trou[1.4cm] & \trou[2.0cm] & \trou[2.0cm] & \trou[2.2cm] & \trou[2.0cm] & \trou[1.8cm] \tabularnewline
30 & \trou[1.4cm] & \trou[2.0cm] & \trou[2.0cm] & \trou[2.2cm] & \trou[2.0cm] & \trou[1.8cm] \tabularnewline
31 & \trou[1.4cm] & \trou[2.0cm] & \trou[2.0cm] & \trou[2.2cm] & \trou[2.0cm] & \trou[1.8cm] \tabularnewline
32 & \trou[1.4cm] & \trou[2.0cm] & \trou[2.0cm] & \trou[2.2cm] & \trou[2.0cm] & \trou[1.8cm] \tabularnewline
33 & \trou[1.4cm] & \trou[2.0cm] & \trou[2.0cm] & \trou[2.2cm] & \trou[2.0cm] & \trou[1.8cm] \tabularnewline
34 & \trou[1.4cm] & \trou[2.0cm] & \trou[2.0cm] & \trou[2.2cm] & \trou[2.0cm] & \trou[1.8cm] \tabularnewline
35 & \trou[1.4cm] & \trou[2.0cm] & \trou[2.0cm] & \trou[2.2cm] & \trou[2.0cm] & \trou[1.8cm] \tabularnewline
36 & \trou[1.4cm] & \trou[2.0cm] & \trou[2.0cm] & \trou[2.2cm] & \trou[2.0cm] & \trou[1.8cm] \tabularnewline
37 & \trou[1.4cm] & \trou[2.0cm] & \trou[2.0cm] & \trou[2.2cm] & \trou[2.0cm] & \trou[1.8cm] \tabularnewline
38 & \trou[1.4cm] & \trou[2.0cm] & \trou[2.0cm] & \trou[2.2cm] & \trou[2.0cm] & \trou[1.8cm] \tabularnewline
39 & \trou[1.4cm] & \trou[2.0cm] & \trou[2.0cm] & \trou[2.2cm] & \trou[2.0cm] & \trou[1.8cm] \tabularnewline
40 & \trou[1.4cm] & \trou[2.0cm] & \trou[2.0cm] & \trou[2.2cm] & \trou[2.0cm] & \trou[1.8cm] \tabularnewline
41 & \trou[1.4cm] & \trou[2.0cm] & \trou[2.0cm] & \trou[2.2cm] & \trou[2.0cm] & \trou[1.8cm] \tabularnewline
42 & \trou[1.4cm] & \trou[2.0cm] & \trou[2.0cm] & \trou[2.2cm] & \trou[2.0cm] & \trou[1.8cm] \tabularnewline
43 & \trou[1.4cm] & \trou[2.0cm] & \trou[2.0cm] & \trou[2.2cm] & \trou[2.0cm] & \trou[1.8cm] \tabularnewline
44 & \trou[1.4cm] & \trou[2.0cm] & \trou[2.0cm] & \trou[2.2cm] & \trou[2.0cm] & \trou[1.8cm] \tabularnewline
45 & \trou[1.4cm] & \trou[2.0cm] & \trou[2.0cm] & \trou[2.2cm] & \trou[2.0cm] & \trou[1.8cm] \tabularnewline
46 & \trou[1.4cm] & \trou[2.0cm] & \trou[2.0cm] & \trou[2.2cm] & \trou[2.0cm] & \trou[1.8cm] \tabularnewline
47 & \trou[1.4cm] & \trou[2.0cm] & \trou[2.0cm] & \trou[2.2cm] & \trou[2.0cm] & \trou[1.8cm] \tabularnewline
48 & \trou[1.4cm] & \trou[2.0cm] & \trou[2.0cm] & \trou[2.2cm] & \trou[2.0cm] & \trou[1.8cm] \tabularnewline
49 & \trou[1.4cm] & \trou[2.0cm] & \trou[2.0cm] & \trou[2.2cm] & \trou[2.0cm] & \trou[1.8cm] \tabularnewline
50 & \trou[1.4cm] & \trou[2.0cm] & \trou[2.0cm] & \trou[2.2cm] & \trou[2.0cm] & \trou[1.8cm] \tabularnewline
51 & \trou[1.4cm] & \trou[2.0cm] & \trou[2.0cm] & \trou[2.2cm] & \trou[2.0cm] & \trou[1.8cm] \tabularnewline
52 & \trou[1.4cm] & \trou[2.0cm] & \trou[2.0cm] & \trou[2.2cm] & \trou[2.0cm] & \trou[1.8cm] \tabularnewline
\hline
\multicolumn{2}{@{}l}{\textbf{Total de l'année}} & \trou[1.4cm] & \trou[1.6cm] & \trou[1.8cm] & \trou[1.6cm] & \tabularnewline
\end{longtable}
\end{tableau}

\section{Le registre des conversations}
À remplir après chaque échange oral --- réel ou enregistré. C'est le tableau le plus
important de cette annexe : c'est celui qui mesure \textbf{ce que vous voulez vraiment}.

\begin{tableau}[Registre de conversation --- à photocopier]
\footnotesize
\begin{longtable}{@{}p{1.4cm}p{2.0cm}p{1.2cm}p{2.5cm}p{2.5cm}p{2.3cm}p{1.7cm}@{}}
\entete{Date} & \entete{Partenaire} & \entete{Durée} & \entete{Sujet} & \entete{Ce qui m'a bloqué} & \entete{3 phrases à retenir} & \entete{Note /5} \tabularnewline
\endfirsthead
\entete{Date} & \entete{Qui} & \entete{min} & \entete{Sujet} & \entete{Blocage} & \entete{À retenir} & \entete{Note} \tabularnewline
\endhead
\hline
 & & & & & & \tabularnewline
 & & & & & & \tabularnewline
 & & & & & & \tabularnewline
 & & & & & & \tabularnewline
 & & & & & & \tabularnewline
 & & & & & & \tabularnewline
 & & & & & & \tabularnewline
 & & & & & & \tabularnewline
 & & & & & & \tabularnewline
 & & & & & & \tabularnewline
 & & & & & & \tabularnewline
 & & & & & & \tabularnewline
 & & & & & & \tabularnewline
 & & & & & & \tabularnewline
 & & & & & & \tabularnewline
 & & & & & & \tabularnewline
 & & & & & & \tabularnewline
 & & & & & & \tabularnewline
 & & & & & & \tabularnewline
 & & & & & & \tabularnewline
\end{longtable}
\end{tableau}

\begin{astuce}[La note sur 5]
Utilisez la grille du chapitre 19, réduite à une seule note globale :
\textbf{1} = j'ai subi la conversation ; \textbf{2} = j'ai répondu, avec effort ;
\textbf{3} = j'ai tenu, en simplifiant ; \textbf{4} = j'ai été à l'aise ;
\textbf{5} = j'ai oublié que je parlais anglais. La moyenne de vos dix dernières notes est
votre vrai niveau --- bien plus fiable que n'importe quel test.
\end{astuce}

\section{Le registre des enregistrements}
Enregistrez-vous au moins une fois par semaine. Ce tableau garde la trace de vos progrès
auditifs --- les seuls que vous puissiez réécouter six mois plus tard avec stupéfaction.

\begin{tableau}[Mes enregistrements]
\footnotesize
\begin{tabularx}{\linewidth}{@{}p{1.6cm}p{2.2cm}p{3.0cm}p{1.6cm}p{2.2cm}X@{}}
\entete{\#} & \entete{Date} & \entete{Sujet / exercice} & \entete{Durée} & \entete{Pauses longues} & \entete{Ce que j'ai remarqué} \tabularnewline
\hline
1 & & Jour 1 --- présentation de soi & & & \tabularnewline
2 & & Semaine 4 & & & \tabularnewline
3 & & Semaine 8 & & & \tabularnewline
4 & & Semaine 12 & & & \tabularnewline
5 & & Semaine 16 & & & \tabularnewline
6 & & Semaine 20 & & & \tabularnewline
7 & & Semaine 24 & & & \tabularnewline
8 & & Semaine 28 & & & \tabularnewline
9 & & Semaine 32 & & & \tabularnewline
10 & & Semaine 36 & & & \tabularnewline
11 & & Semaine 40 & & & \tabularnewline
12 & & Semaine 44 & & & \tabularnewline
13 & & Semaine 48 & & & \tabularnewline
14 & & Jour 365 --- présentation de soi (même sujet qu'au jour 1) & & & \tabularnewline
\end{tabularx}
\end{tableau}

\begin{regle}[Le test du jour 1 et du jour 365]
Le jour 1, enregistrez-vous en train de vous présenter pendant trois minutes --- même sujet,
mêmes questions. Le jour 365, refaites exactement le même enregistrement. Écoutez-les à la
suite. Cette comparaison vaut tous les diplômes du monde : c'est la seule preuve irréfutable
de ce que 900 heures de travail ont produit.
\end{regle}

\section{Le suivi des 100 erreurs}
Une case par erreur du chapitre 18. Cochez-la quand vous l'avez \textbf{éliminée}, c'est-à-dire
quand vous ne la commettez plus pendant trois semaines consécutives.

\begin{tableau}[Tableau de chasse des erreurs 1 à 100]
\footnotesize
\begin{tabular}{@{}p{0.72cm}p{0.72cm}p{0.72cm}p{0.72cm}p{0.72cm}p{0.72cm}p{0.72cm}p{0.72cm}p{0.72cm}p{0.72cm}p{0.72cm}p{0.72cm}p{0.72cm}p{0.72cm}@{}}
\case & \case & \case & \case & \case & \case & \case & \case & \case & \case & \case & \case & \case & \case \\
1 & 2 & 3 & 4 & 5 & 6 & 7 & 8 & 9 & 10 & 11 & 12 & 13 & 14 \\[2pt]
\case & \case & \case & \case & \case & \case & \case & \case & \case & \case & \case & \case & \case & \case \\
15 & 16 & 17 & 18 & 19 & 20 & 21 & 22 & 23 & 24 & 25 & 26 & 27 & 28 \\[2pt]
\case & \case & \case & \case & \case & \case & \case & \case & \case & \case & \case & \case & \case & \case \\
29 & 30 & 31 & 32 & 33 & 34 & 35 & 36 & 37 & 38 & 39 & 40 & 41 & 42 \\[2pt]
\case & \case & \case & \case & \case & \case & \case & \case & \case & \case & \case & \case & \case & \case \\
43 & 44 & 45 & 46 & 47 & 48 & 49 & 50 & 51 & 52 & 53 & 54 & 55 & 56 \\[2pt]
\case & \case & \case & \case & \case & \case & \case & \case & \case & \case & \case & \case & \case & \case \\
57 & 58 & 59 & 60 & 61 & 62 & 63 & 64 & 65 & 66 & 67 & 68 & 69 & 70 \\[2pt]
\case & \case & \case & \case & \case & \case & \case & \case & \case & \case & \case & \case & \case & \case \\
71 & 72 & 73 & 74 & 75 & 76 & 77 & 78 & 79 & 80 & 81 & 82 & 83 & 84 \\[2pt]
\case & \case & \case & \case & \case & \case & \case & \case & \case & \case & \case & \case & \case & \case \\
85 & 86 & 87 & 88 & 89 & 90 & 91 & 92 & 93 & 94 & 95 & 96 & 97 & 98 \\[2pt]
\case & \case & & & & & & & & & & & & \\
99 & 100 & & & & & & & & & & & & \\
\end{tabular}
\end{tableau}

\section{Le bilan mensuel}
À photocopier douze fois. Dix minutes, le dernier dimanche de chaque mois.

\begin{tableau}[Mois \rule{1cm}{0.4pt} --- du \rule{2cm}{0.4pt} au \rule{2cm}{0.4pt}]
\footnotesize
\begin{tabularx}{\linewidth}{@{}p{7.6cm}X@{}}
\entete{Question} & \entete{Réponse} \tabularnewline
\hline
Nombre de jours étudiés sur 30 & \rule{0pt}{0.9cm}\trou[5.0cm] \tabularnewline
Nombre total de minutes & \rule{0pt}{0.9cm}\trou[5.0cm] \tabularnewline
Nombre de minutes de parole réelle & \rule{0pt}{0.9cm}\trou[5.0cm] \tabularnewline
Nombre de conversations tenues & \rule{0pt}{0.9cm}\trou[5.0cm] \tabularnewline
Nombre de mots actifs dans Anki & \rule{0pt}{0.9cm}\trou[5.0cm] \tabularnewline
Chapitres de ce livre terminés & \rule{0pt}{0.9cm}\trou[5.0cm] \tabularnewline
Test mensuel --- score de grammaire & \rule{0pt}{0.9cm}\trou[2.0cm] / 20 \tabularnewline
Test mensuel --- score de vocabulaire & \rule{0pt}{0.9cm}\trou[2.0cm] / 30 \tabularnewline
Grille orale --- total & \rule{0pt}{0.9cm}\trou[2.0cm] / 25 \tabularnewline
Erreurs éliminées ce mois-ci (numéros) & \rule{0pt}{0.9cm}\trou[5.0cm] \tabularnewline
Ce qui a le mieux marché & \rule{0pt}{1.4cm}\trou[6.0cm] \tabularnewline
Ce qui m'a fait perdre du temps & \rule{0pt}{1.4cm}\trou[6.0cm] \tabularnewline
Ce que je change le mois prochain & \rule{0pt}{1.4cm}\trou[6.0cm] \tabularnewline
Mon moment de fierté du mois & \rule{0pt}{1.4cm}\trou[6.0cm] \tabularnewline
Un mot d'encouragement pour moi-même & \rule{0pt}{1.4cm}\trou[6.0cm] \tabularnewline
\end{tabularx}
\end{tableau}

\section{La lettre du jour 365}
Le dernier jour, écrivez-vous une lettre --- \textbf{en anglais}, évidemment. Elle répond
à quatre questions : qu'est-ce que je sais faire aujourd'hui que je ne savais pas faire il y
a un an ? Quelle est la conversation dont je suis le plus fier ? Qu'est-ce qui m'a fait
vouloir abandonner, et comment ai-je continué ? Que vais-je faire de cette compétence ?
Puis cachetez-la, et ouvrez-la dans cinq ans.

\begin{tableau}[Mes repères de fin d'année]
\begin{tabularx}{\linewidth}{@{}p{8.2cm}X@{}}
\entete{Indicateur} & \entete{Jour 1 \qquad\qquad Jour 365} \tabularnewline
\hline
Je peux me présenter et parler de ma vie pendant\ldots{} & \trou[2.2cm] min \qquad \trou[2.2cm] min \tabularnewline
Je comprends un anglophone qui parle normalement\ldots{} & \trou[2.2cm] \% \qquad \trou[2.2cm] \% \tabularnewline
Je connais activement\ldots{} & \trou[2.2cm] mots \quad \trou[2.2cm] mots \tabularnewline
Je lis un article de presse\ldots{} & \trou[2.2cm] \% \qquad \trou[2.2cm] \% \tabularnewline
J'écris un e-mail correct en\ldots{} & \trou[2.2cm] min \qquad \trou[2.2cm] min \tabularnewline
Je traduis encore dans ma tête\ldots{} & \trou[2.2cm] fois/min \quad \trou[2.2cm] fois/min \tabularnewline
Test de positionnement (chapitre 19)\ldots{} & \trou[2.2cm] /40 \qquad \trou[2.2cm] /40 \tabularnewline
Grille d'auto-évaluation orale\ldots{} & \trou[2.2cm] /25 \qquad \trou[2.2cm] /25 \tabularnewline
\end{tabularx}
\end{tableau}

\begin{mission}[Commencez maintenant, pas demain]
Imprimez cette annexe aujourd'hui. Remplissez le contrat personnel ce soir --- même
imparfaitement, même en dix minutes. Fixez l'heure de votre première séance, et mettez une
alarme sur votre téléphone pour demain matin. Le jour 1 n'est pas une date dans un livre :
c'est le moment où vous prenez le stylo.
\end{mission}

% =====================================================================
%  ANNEXE G — GLOSSAIRE, TABLE DE REPÉRAGE ET MOT FINAL
% =====================================================================
\chapter{Glossaire, repérage et mot final}

\section{Glossaire des termes grammaticaux}
Tous les termes techniques employés dans ce livre, avec leur équivalent anglais, une
explication en une ligne et un exemple. Classés par ordre alphabétique français.

\begin{tableau}[Glossaire --- 54 termes]
\footnotesize
\begin{longtable}{@{}p{3.6cm}p{3.2cm}p{7.4cm}@{}}
\entete{Terme français} & \entete{Terme anglais} & \entete{Explication et exemple} \tabularnewline
\endfirsthead
\entete{Français} & \entete{Anglais} & \entete{Explication} \tabularnewline
\endhead
\hline
adjectif & \en{adjective} & Mot qui qualifie un nom, placé \textbf{avant} en anglais. \en{a \textbf{beautiful} house}. \tabularnewline
adverbe & \en{adverb} & Modifie un verbe ou un adjectif. \en{She sings \textbf{well}.} \tabularnewline
auxiliaire & \en{auxiliary verb} & Verbe d'aide : \en{be, have, do} et les modaux. \en{She \textbf{has} left.} \tabularnewline
base verbale & \en{base form / bare infinitive} & Forme nue du verbe, sans \en{to} ni \en{-s}. \en{to go \ra go}. \tabularnewline
complément & \en{complement} & Élément qui complète le verbe ou le nom. \en{She is \textbf{a doctor}.} \tabularnewline
comparatif & \en{comparative} & Forme de comparaison entre deux éléments. \en{taller than, more interesting than}. \tabularnewline
compte rendu & \en{minutes / report} & Texte qui résume une réunion ou un événement. \tabularnewline
conditionnel & \en{conditional} & Structure en \en{if} exprimant une hypothèse. \en{If I had time, I'd come.} \tabularnewline
conjonction & \en{conjunction / linking word} & Mot qui relie deux idées. \en{and, but, because, although}. \tabularnewline
contractée (forme) & \en{contraction} & Forme raccourcie de l'oral. \en{I'm, don't, she's, won't, I'd}. \tabularnewline
coordination & \en{coordination} & Jonction de deux éléments de même niveau. \en{tea \textbf{and} coffee}. \tabularnewline
dénombrable & \en{countable} & Nom qui se compte, avec un pluriel. \en{one book, two books}. \tabularnewline
discours rapporté & \en{reported / indirect speech} & Façon de rapporter les paroles d'autrui. \en{She said she was tired.} \tabularnewline
ellipse & \en{ellipsis} & Omission d'un élément sous-entendu. \en{``Who wants tea?'' --- ``Me!''} \tabularnewline
gérondif & \en{gerund} & Forme en \en{-ing} employée comme nom. \en{\textbf{Swimming} is fun. I enjoy \textbf{reading}.} \tabularnewline
impératif & \en{imperative} & Ordre, conseil, invitation. \en{\textbf{Sit} down. \textbf{Don't} worry.} \tabularnewline
indénombrable & \en{uncountable} & Nom sans pluriel, mesuré et non compté. \en{water, information, advice, money}. \tabularnewline
infinitif & \en{infinitive} & \en{to} + base verbale. \en{I want \textbf{to leave}.} \tabularnewline
interrogatif (mot) & \en{Wh-word} & \en{who, what, where, when, why, how, which, whose}. \tabularnewline
intransitif & \en{intransitive} & Verbe sans complément d'objet. \en{He \textbf{arrived}. She \textbf{slept}.} \tabularnewline
marqueur temporel & \en{time marker / signal word} & Mot qui indique le temps à employer. \en{yesterday, since, every day}. \tabularnewline
modal & \en{modal verb} & Auxiliaire de nuance : \en{can, could, may, might, must, should, will, would}. \tabularnewline
nom & \en{noun} & \en{house, freedom, Senegal}. \tabularnewline
négation & \en{negation} & \en{not}, ou les auxiliaires négatifs. \en{I do\textbf{n't} know.} \tabularnewline
objet (complément d') & \en{(direct) object} & Ce qui subit l'action. \en{I read \textbf{a book}.} \tabularnewline
participe passé & \en{past participle} & Troisième forme du verbe. \en{gone, written, eaten, played}. \tabularnewline
participe présent & \en{present participle} & Forme en \en{-ing} des temps continus. \en{I am \textbf{working}.} \tabularnewline
passé simple & \en{past simple / preterite} & Temps du récit achevé. \en{I \textbf{went} home.} \tabularnewline
passif & \en{passive voice} & Sujet qui subit l'action. \en{The bridge \textbf{was built} in 1990.} \tabularnewline
phrase complexe & \en{complex sentence} & Une principale + une ou plusieurs subordonnées. \tabularnewline
phrasal verb & \en{phrasal verb} & Verbe + particule au sens propre. \en{give up, look after, turn off}. \tabularnewline
préposition & \en{preposition} & Mot de relation. \en{in, on, at, for, with, about}. \tabularnewline
présent continu & \en{present continuous} & \en{be + -ing}, action en cours. \en{I \textbf{am reading}.} \tabularnewline
présent parfait & \en{present perfect} & \en{have + participe passé}, passé lié au présent. \en{I \textbf{have seen} it.} \tabularnewline
pronom & \en{pronoun} & Remplace un nom. \en{I, you, he, it, we, they, mine, this}. \tabularnewline
pronom relatif & \en{relative pronoun} & \en{who, which, that, whose, where}. \en{The man \textbf{who} called.} \tabularnewline
proposition & \en{clause} & Groupe sujet + verbe. \en{[I think] [that she is right]}. \tabularnewline
quantifieur & \en{quantifier} & \en{some, any, much, many, a lot of, few, several}. \tabularnewline
question tag & \en{question tag} & Mini-question en fin de phrase. \en{It's cold, \textbf{isn't it}?} \tabularnewline
réfléchi (pronom) & \en{reflexive pronoun} & \en{myself, yourself, himself, ourselves, themselves}. \tabularnewline
registre & \en{register} & Niveau de langue : soutenu, courant, familier. \tabularnewline
relative (proposition) & \en{relative clause} & Subordonnée qui précise un nom. \en{the book \textbf{that I read}}. \tabularnewline
schwa & \en{schwa} & Son neutre très court, le plus fréquent de l'anglais. \en{about \pron{eubaout}}. \tabularnewline
subjonctif & \en{subjunctive} & Rare en anglais, presque invisible. \en{I suggest that he \textbf{be} informed.} \tabularnewline
subordonnée & \en{subordinate clause} & Proposition dépendante. \en{\textbf{When I arrived}, she left.} \tabularnewline
superlatif & \en{superlative} & Le plus haut degré. \en{the tallest, the most interesting}. \tabularnewline
sujet & \en{subject} & Celui qui fait l'action --- \textbf{toujours exprimé} en anglais. \en{\textbf{It} is raining.} \tabularnewline
suffixe & \en{suffix} & Terminaison qui change la catégorie. \en{teach \ra teach\textbf{er}, kind \ra kind\textbf{ness}}. \tabularnewline
syllabe tonique & \en{stressed syllable} & La syllabe prononcée plus fort. \en{ba\textbf{NA}na}. \tabularnewline
temps & \en{tense} & Forme verbale qui situe l'action dans le temps. \tabularnewline
transitif & \en{transitive} & Verbe qui exige un complément d'objet. \en{I \textbf{bought} a car.} \tabularnewline
verbe d'état & \en{stative verb} & Verbe de sensation, opinion, possession --- rarement au \en{-ing}. \en{know, like, want, have, believe}. \tabularnewline
verbe irrégulier & \en{irregular verb} & Verbe aux trois formes non standard. \en{go/went/gone}. \tabularnewline
voix active & \en{active voice} & Sujet qui agit. \en{The dog bit the man.} \tabularnewline
\end{longtable}
\end{tableau}

\section{Table de repérage : où trouver ce que je cherche ?}
\begin{tableau}[Vos questions, nos réponses]
\footnotesize
\begin{longtable}{@{}p{8.2cm}p{5.9cm}@{}}
\entete{Je veux\ldots{}} & \entete{Voir\ldots{}} \tabularnewline
\endfirsthead
\entete{Je veux\ldots{}} & \entete{Voir\ldots{}} \tabularnewline
\endhead
\hline
comprendre comment organiser mon année & chapitre 2 (plan des 12 mois) \tabularnewline
apprendre à travailler efficacement & chapitre 3 (huit techniques) \tabularnewline
trouver des ressources gratuites & chapitre 4 et annexe E \tabularnewline
prononcer un son correctement & chapitres 5 et 6 \tabularnewline
corriger mon accent français & chapitre 7 \tabularnewline
comprendre le \en{present perfect} & chapitre 9 \tabularnewline
maîtriser les articles \en{a/an/the} & chapitre 8 \tabularnewline
poser des questions & chapitre 8 et fiche 3 \tabularnewline
utiliser les modaux & chapitre 10 et fiche 2 \tabularnewline
construire le passif & chapitre 10 et fiche 9 \tabularnewline
employer le discours rapporté & chapitre 10 et fiche 7 \tabularnewline
comprendre les conditionnels & chapitre 11 et fiche 8 \tabularnewline
apprendre du vocabulaire par thème & chapitre 12 \tabularnewline
conjuguer un verbe irrégulier & annexe A \tabularnewline
comprendre un \emph{phrasal verb} & chapitre 13 et annexe D (particules) \tabularnewline
éviter un faux ami & chapitre 13 et fiche 13 \tabularnewline
apprendre les idiomes & chapitre 13 \tabularnewline
dialoguer dans une situation précise & chapitre 14 \tabularnewline
avoir une phrase toute faite & annexe B \tabularnewline
parler sans traduire dans ma tête & chapitre 15 \tabularnewline
me débloquer quand je cherche un mot & chapitre 15 (stratégies de contournement) \tabularnewline
améliorer ma compréhension orale & chapitre 16 \tabularnewline
trouver un texte d'écoute avec corrigé & chapitre 16 (dix textes) \tabularnewline
lire un texte adapté à mon niveau & chapitre 17, partie A \tabularnewline
écrire un e-mail correct & chapitre 17, partie D \tabularnewline
tenir un journal en anglais & chapitre 17, partie E \tabularnewline
connaître l'orthographe britannique ou américaine & chapitre 17, partie B \tabularnewline
corriger mes erreurs typiques & chapitre 18 \tabularnewline
choisir la bonne préposition & annexe D \tabularnewline
varier mes connecteurs & annexe D \tabularnewline
évaluer mon niveau & chapitre 19 \tabularnewline
préparer un examen & chapitre 19 \tabularnewline
réviser toute la grammaire en 14 pages & annexe C \tabularnewline
suivre ma progression au quotidien & annexe F \tabularnewline
\end{longtable}
\end{tableau}

\section{Le mot final}

Vous êtes arrivé au bout de ce livre --- ce qui, statistiquement, est déjà un exploit.
La plupart des gens qui achètent une méthode la feuillettent trois jours. Vous, vous l'avez
lue jusqu'à la dernière annexe.

Mais vous savez maintenant une chose essentielle : \textbf{ce livre n'était pas le but}.
Il était une carte. La carte n'est pas le territoire, et connaître le chemin par c\oe ur ne
fait pas avancer d'un mètre. Ce qui vous fera parler anglais, ce ne sont pas les 3\,000 mots
de ces pages : ce sont les 300 heures pendant lesquelles votre bouche les prononcera, les
50 conversations où vous aurez eu peur et où vous serez resté, les 200 enregistrements que
vous aurez réécoutés en grimaçant, les 400 jours où vous vous serez assis à la même heure
alors que vous n'en aviez aucune envie.

Vous allez connaître trois moments difficiles. Prévoyons-les.

\textbf{Le premier, vers la sixième semaine.} L'enthousiasme retombe, la nouveauté s'use,
et vous aurez l'impression de ne rien apprendre. C'est exactement le moment où les autres
abandonnent. La seule chose à faire : réduire, mais ne pas arrêter. Vingt minutes par jour
plutôt que zéro.

\textbf{Le deuxième, vers le cinquième mois.} Vous comprendrez presque tout, mais vous
parlerez mal, et cet écart entre votre oreille et votre bouche est douloureux. C'est bon
signe : cela signifie que votre compréhension a progressé plus vite que votre production,
ce qui est l'ordre naturel des choses. La parole rattrapera. Continuez les drills du
chapitre 15 --- ils sont faits précisément pour ce passage.

\textbf{Le troisième, vers le neuvième mois.} Vous aurez une conversation ratée. Un mot
manquera au mauvais moment, vous rougirez, et vous vous direz que tout cela ne sert à rien.
Le lendemain, relisez votre journal de progression (annexe F). Regardez la ligne de la
semaine 1. Puis écoutez votre enregistrement du jour 1. Vous ne reconnaîtrez pas cette
personne --- et c'est précisément la preuve que vous êtes devenu quelqu'un d'autre.

Une dernière chose, la plus importante de toutes. Vous n'apprenez pas l'anglais pour passer
un examen, ni pour impressionner quelqu'un, ni pour cocher une case sur un CV. Vous
l'apprenez pour \textbf{rencontrer des gens}. Chaque conversation réelle --- même ratée,
même boiteuse, même avec vingt silences --- vaut plus que dix heures d'étude solitaire.
Allez vers les gens dès le premier mois. Parlez mal, tôt, souvent. C'est le secret que tous
les polyglottes finissent par avouer : ils n'ont jamais été bons avant d'avoir été mauvais
longtemps, en public.

Dans douze mois, quelqu'un vous demandera : « Mais vous parlez anglais ? » Et vous répondrez
--- peut-être en haussant les épaules, parce que vous n'y penserez même plus :
\en{``Yes, more or less. I've been practising for a year.''}

Ce jour-là, pensez à revenir écrire une ligne dans votre journal.

\begin{center}
\vspace{0.6cm}
\textit{\en{``The best time to plant a tree was twenty years ago.\\ The second best time is now.''}}\\[6pt]
\footnotesize\selectlanguage{french}\textit{Le meilleur moment pour planter un arbre,
c'était il y a vingt ans. Le deuxième meilleur moment, c'est maintenant.}
\vspace{0.6cm}
\end{center}

\begin{mission}[Votre première action, dans les dix prochaines minutes]
\textbf{1.} Ouvrez votre téléphone et réglez une alarme quotidienne à l'heure que vous avez
choisie dans votre contrat personnel (annexe F).\\
\textbf{2.} Enregistrez-vous pendant trois minutes : présentez-vous en anglais, comme vous
pouvez, maintenant. Gardez ce fichier. Ne le supprimez jamais.\\
\textbf{3.} Envoyez le message de l'annexe E à une personne.\\
\textbf{4.} Commencez le chapitre 5 demain matin.\\
C'est tout. Quatre gestes, dix minutes, et votre année a commencé.
\end{mission}


% ------------------------------------------------------------------
%  FIN
% ------------------------------------------------------------------
\backmatter

\end{document}
